\documentclass[10pt]{article}
\usepackage[margin=0.78in]{geometry}
\usepackage[T1]{fontenc}
\usepackage{amsmath,amssymb,amsthm,mathtools,booktabs,longtable,array,tabularx,microtype,lmodern}
\usepackage{needspace,etoolbox,xurl}
\usepackage[dvipsnames]{xcolor}
\usepackage[colorlinks=true,linkcolor=NavyBlue,urlcolor=BrickRed,hypertexnames=false]{hyperref}
\usepackage{bookmark}
\usepackage{fancyhdr,enumitem}
\usepackage{titlesec}

\definecolor{ink}{HTML}{17212B}
\definecolor{navy}{HTML}{17365D}
\definecolor{teal}{HTML}{147D78}
\definecolor{pale}{HTML}{EEF4F7}
\color{ink}

\hypersetup{
  pdftitle={All New Matrix Groups: Formulas, Derivations, and Proofs},
  pdfauthor={Lambda Distributions workspace},
  pdfsubject={Consolidated proof anthology for matrix-group sigma-MGF formulas}
}

\pagestyle{fancy}
\fancyhf{}
\lhead{\nouppercase{\leftmark}}
\rhead{Proof anthology}
\cfoot{\thepage}
\renewcommand{\sectionmark}[1]{\markboth{\thesection\quad #1}{}}
\setlength{\headheight}{14pt}
\setlength{\parindent}{0pt}
\setlength{\parskip}{5pt}
\raggedbottom
\clubpenalty=10000
\widowpenalty=10000
\displaywidowpenalty=10000
\allowdisplaybreaks[3]
\setlength{\emergencystretch}{2.5em}
\urlstyle{same}
\setlist[itemize]{leftmargin=1.35em,itemsep=2pt,topsep=3pt}
\setlist[enumerate]{leftmargin=1.55em,itemsep=2pt,topsep=3pt}
% Keep headings and short literal command blocks with the material that
% follows them.  This prevents orphaned headings and one-line spill pages.
\AddToHook{cmd/subsection/before}{\Needspace{5\baselineskip}}
\AddToHook{cmd/subsubsection/before}{\Needspace{4\baselineskip}}
\BeforeBeginEnvironment{verbatim}{\Needspace{6\baselineskip}}
% Keep the anthology navigation compact: only proof-level sections are
% numbered and listed.  Internal derivation headings remain visible in the
% text, but do not produce entries such as 12.1 or 12.1.1.
\setcounter{secnumdepth}{3}
\setcounter{tocdepth}{1}
\titleformat{\part}[display]{\centering\color{navy}\bfseries\Large}
{Part \thepart}{0.4em}{\LARGE}
\titlespacing*{\part}{0pt}{0pt}{1.5em}
\titleformat{\section}{\color{navy}\bfseries\Large}{\thesection}{0.65em}{}
\titleformat{\subsection}{\color{teal}\bfseries\large}{\thesubsection}{0.6em}{}
\titleformat{\subsubsection}{\bfseries}{\thesubsubsection}{0.55em}{}

\newcommand{\Sym}{\operatorname{Sym}}
\newcommand{\Tr}{\operatorname{Tr}}
\newcommand{\Fix}{\operatorname{Fix}}
\newcommand{\CT}{\operatorname{CT}}
\newcommand{\ExpS}{\operatorname{Exp}_{\sigma}}
\newcommand{\SMG}{\operatorname{SMG}}
\newcommand{\MGF}{\SMG}
\newcommand{\smg}[1]{\mathbb E\!\left[\ExpS\!\left(X_{#1}h_1\right)\right]}
\newcommand{\smgv}[2]{\mathbb E\!\left[\ExpS\!\left(X_{#1,#2}h_1\right)\right]}
\newcommand{\F}{\mathbb F}
\newcommand{\C}{\mathbb C}
\newcommand{\E}{\mathbb E}
\newcommand{\Z}{\mathbb Z}
\newcommand{\Aut}{\operatorname{Aut}}
\newcommand{\End}{\operatorname{End}}
\newcommand{\Hom}{\operatorname{Hom}}
\newcommand{\GL}{\operatorname{GL}}
\newcommand{\SL}{\operatorname{SL}}
\newcommand{\Sp}{\operatorname{Sp}}
\newcommand{\Gr}{\operatorname{Gr}}
\newcommand{\Ad}{\operatorname{Ad}}
\newcommand{\Cl}{\operatorname{Cl}}
\newcommand{\Cyc}{\operatorname{Cyc}}
\newcommand{\Irr}{\operatorname{Irr}}
\newcommand{\Ind}{\operatorname{Ind}}
\newcommand{\Lie}{\operatorname{Lie}}
\newcommand{\MSet}{\operatorname{MSet}}
\newcommand{\cwe}{\operatorname{cwe}}
\newcommand{\im}{\operatorname{im}}
\newcommand{\ch}{\operatorname{ch}}
\newcommand{\rank}{\operatorname{rank}}
\newcommand{\calO}{\mathcal O}
\newcommand{\one}{\mathbf 1}
\newcommand{\BMGF}{\mathcal B}
\newcommand{\BSMG}{\mathcal B}
\newcommand{\PP}{\mathcal P}
\newcommand{\qbinom}[2]{\genfrac{[}{]}{0pt}{}{#1}{#2}_q}
\newcommand{\resultbox}[1]{%
  \begin{center}\fcolorbox{navy}{pale}{\parbox{0.93\linewidth}{#1}}\end{center}}

\newcommand{\proofstatus}[2]{%
  \begin{center}\fcolorbox{teal}{pale}{\parbox{0.93\linewidth}{%
  \textbf{Status: #1.} #2}}\end{center}}
\newenvironment{intuition}{\begin{quote}\small\textbf{Intuition.}\ }{\end{quote}}
\newenvironment{proofaudit}{\begin{quote}\small\textbf{Proof audit.}\ }{\end{quote}}

% Shared theorem environments for the inlined proof bodies.
\newtheorem{theorem}{Theorem}
\newtheorem{proposition}{Proposition}
\newtheorem{lemma}{Lemma}
\newtheorem{corollary}{Corollary}
\newtheorem{remark}{Remark}

\title{\color{navy}\bfseries All New Matrix Groups:\\[3pt]
Formulas, Derivations, and Proofs\\[6pt]
\large A consolidated proof anthology organized around established tools}
\author{Lambda Distributions workspace}
\date{21 July 2026 --- integrated with the 14 July matrix-group paper}

\begin{document}
\maketitle

\begin{abstract}
This volume is the proof-level companion to
\emph{All New Matrix Groups and Formulas}.  It does not merely list the final
identities: it begins with the common tools that produce them and then bundles
the complete proof-and-verification notes for every family, in the same family
order as the formula compendium.  The source notes are included without
abridgment so that hypotheses, finite-rank qualifications, corrections,
computational checks, and references remain attached to the claims they
support.  Superseded flat exports and byte-level duplicate copies are omitted;
their expanded canonical editions appear once in the sections below.

This revised edition treats the foundational identities of
\emph{Matrix group $\Lambda$-distributions} (draft of 14 July 2026) as
imported tools rather than reproving them in each family. It corrects the
normalized image/kernel convention, separates earlier theorems from new
extensions, adds explicit stable-range and central-charge hypotheses, and
moves unresolved proof obligations to a final provisional part.
\end{abstract}

\resultbox{\textbf{How to use this volume.}  Start with the derivation toolkit
below to identify the relevant reduction: Reynolds/Molien averaging, cycle
index, Fourier projection, Weyl constant term, central-character selection,
or a two-sided balanced series.  Then follow the table of contents to the
family-specific note, where the reduction is specialized and independently
verified.  Exact finite formulas, stable formulas, limiting laws, and numerical
diagnostics are deliberately kept distinct.}



\section*{Conventions and claim labels}
\addcontentsline{toc}{section}{Conventions and claim labels}
\begin{center}
\small
\renewcommand{\arraystretch}{1.2}
\begin{tabularx}{0.97\linewidth}{@{}>{\bfseries}p{0.25\linewidth}X@{}}
\toprule
Object & Convention used throughout this edition \\
\midrule
Representation & A specified continuous finite-dimensional complex representation
$\rho:G\to GL(V)$; the represented image $\rho(G)$ is part of the data. \\
Averaging & Normalized counting measure for finite groups and normalized Haar
probability for compact groups. Non-Haar laws are always named explicitly. \\
$\sigma$-MGF & The coefficientwise formal series
$\MGF_{G,V}:=\smgv{G}{V}$, equivalently
$\int_G\prod_i\det(I-t_i\rho(g))^{-1}\,dg$. \\
Coefficient & For a partition $\tau=(\tau_1\ge\cdots\ge\tau_s>0)$,
$[m_\tau]\MGF_{G,V}=\dim(\bigotimes_i\Sym^{\tau_i}V)^G$
under Haar averaging; this is imported from Lemma 3.3.1 of the 14 July paper. \\
Image versus kernel & Normalized Haar/counting averages depend only on the
represented image. Kernels matter only for unnormalized sums or non-Haar
pushforwards. \\
One-sided versus balanced & Scalar-rich representations often have trivial
one-sided series. The balanced series pairs $V$ with $V^*$ and is graded by
central charge. \\
Validity labels & \emph{Exact}: all stated finite parameters. \emph{Stable}:
a proved threshold for a fixed coefficient. \emph{Limit}: a specified topology
and normalization. \emph{Conditional}: an explicitly named missing structural
input. \emph{Diagnostic}: finite computation only. \\
\bottomrule
\end{tabularx}
\end{center}

\resultbox{\textbf{Global correction.} If $H=\rho(G)$ and the average is
normalized Haar/counting measure, then integration over $G$ equals integration
over $H$. Earlier warnings that coincident matrices must be counted with
``abstract multiplicity'' have been removed.}

\tableofcontents
\clearpage

\part{Derivation toolkit: how the formulas are obtained}

\section{Imported tools from \emph{Matrix group $\Lambda$-distributions}}
\label{sec:imported-tools}

We use the following results from M.~Bertucci, S.~Van Blarcom,
B.~Glancy, S.~Howe, T.~Madden, B.~Miller, S.~Pal, G.~Papagrigoriou,
N.~Raikman, and T.~Tran, \emph{Matrix group $\Lambda$-distributions},
draft of 14 July 2026. They are not reproved in the family notes below.

Let $G$ be finite or compact, let $V$ be a specified finite-dimensional
complex representation, and put
\[
W_\tau(V)=\bigotimes_{a=1}^{\ell(\tau)}\Sym^{\tau_a}V.
\]
Throughout,
\[
\boxed{\MGF_{G,V}:=\smgv{G}{V}}
\]
(or $\MGF_G:=\smg{G}$ when the matrix realization is understood). Thus
$\MGF$ is only a compact abbreviation for the displayed expectation.

\begin{theorem}[Imported coefficient identity; Lemma 3.3.1]
\label{tool:coefficient}
For normalized Haar or counting measure,
\[
\boxed{\MGF_{G,V}=\sum_\tau \dim W_\tau(V)^G\,m_\tau,\qquad
[m_\tau]\MGF_{G,V}=\dim W_\tau(V)^G.}
\]
\end{theorem}

\begin{theorem}[Imported generalized Molien identity; Lemma 3.3.2]
\label{tool:molien}
Coefficientwise as a formal symmetric series,
\[
\boxed{\MGF_{G,V}=\int_G\prod_{i\ge1}\det(I-t_i\rho(g))^{-1}\,dg.}\label{eq:tool-finite-molien}
\]
For a finite group the integral is the normalized sum over $G$.
\end{theorem}

\begin{theorem}[Imported permutation-orbit identity; Lemma 3.3.5]
\label{tool:permutation-orbits}
If $V=\C[X]$ is a permutation representation, then
\[
[m_\tau]\MGF_{G,X}=\#\!\left(G\backslash
\prod_a\Sym^{\tau_a}X\right).
\]
\end{theorem}

\begin{theorem}[Imported independence criterion; Theorem C]
\label{tool:independence}
The power-sum trace variables are independent precisely when the
$\sigma$-MGF factors into one-variable moment-generating factors in the
coordinates $p_r/r$. The analogous statement for cycle-counting functions
uses the falling $\sigma$-MGF.
\end{theorem}

Theorem A and Theorem B of the same paper are cited rather than reproved for
the cyclic, dihedral, Pauli, alternating, and classical Weyl cases. Later
sections retain only a correction, an exact finite-rank strengthening, a more
general representation, or an independent verification. In particular, the
normalized one-dimensional cyclic coefficient is $1$, not $n$, when
$n\mid|\tau|$; that section is explicitly presented as a correction and
extension.

\subsection{Proof pipeline after importing the tools}
A family-specific proof now begins at the first structural reduction:
weights, cycles, conjugacy classes, a maximal torus, a normal subgroup, or a
representation-theoretic decomposition. It then supplies a genuinely
independent coefficient extraction and states the exact range of validity.
The Reynolds projection and determinant expansion are not repeated merely to
identify $[m_\tau]$ with an invariant dimension.

\begin{theorem}[Normalized image invariance]\label{thm:image-invariance}\label{thm:normalized-image-invariance}
Let $\rho:G\to H:=\rho(G)$ be a continuous surjective homomorphism of finite
or compact groups. The pushforward of normalized Haar measure on $G$ is
normalized Haar measure on $H$. Consequently
\[
\MGF_{G,V}=\MGF_{H,V},\qquad W^G=W^H
\]
for every module $W$ obtained functorially from $V$.
\end{theorem}
\begin{proof}
The pushforward is a left-invariant probability measure on $H$, hence is Haar
measure by uniqueness. The fixed spaces agree because the two actions have
the same image.
\end{proof}

\begin{intuition}
A kernel repeats every image matrix uniformly in a normalized finite average;
Haar pushforward is the compact analogue. Central extensions still matter
when their central scalars act nontrivially in the chosen linear lift.
\end{intuition}

\section{Plethystic exponentials and trace coordinates}

The logarithmic determinant expansion
\[
-\log\det(I-tA)=\sum_{r\ge1}\frac{t^r}{r}\Tr(A^r)
\]
rewrites every Molien integrand as
\begin{equation}\label{eq:tool-plethystic}
\prod_i\det(I-t_iA)^{-1}
=\exp\!\left(\sum_{r\ge1}\frac{p_r(\mathbf t)}r\Tr(A^r)\right).
\end{equation}
This is the bridge between eigenvalue calculations and symmetric functions.
It explains why a permutation cycle of length $r$ contributes
$\ExpS(p_r)$, why wreath products exponentiate connected cycles, and why
Schur-functor problems reduce to polynomial substitutions in trace variables.

Two expansions are repeatedly useful:
\[
\ExpS(F)=\exp\!\left(\sum_{r\ge1}\frac1r p_r[F]\right),
\qquad
h_d=\sum_{\lambda\vdash d}\frac{p_\lambda}{z_\lambda}.
\]
The first builds multisets from connected components; the second changes from
power sums (natural for cycles and traces) to complete symmetric functions
(natural for symmetric powers).

\section{Finite Fourier filters and zero-weight counting}

Suppose $A=\prod_{\ell=1}^rC_{m_\ell}$ is abelian and
$V=\bigoplus_{j=1}^d\C_{w_j}$ is decomposed into one-dimensional weight
spaces.  A monomial in $W_\tau(V)$ is described by occupancies $q_{bj}\ge0$
with row sums $\sum_jq_{bj}=\tau_b$.  Its total $A$-weight is
\[
\left(\sum_{b,j}w_{\ell j}q_{bj}\right)_{\ell=1}^r.
\]
It is invariant exactly when every component is zero modulo $m_\ell$.
Equivalently, use character orthogonality
\[
\frac1{|A|}\sum_{a\in A}\chi(a)=
\begin{cases}1,&\chi=1,\\0,&\chi\ne1.\end{cases}
\]
This derives the simultaneous-congruence formula and, after conditioning on a
permutation part, the dual-code filter for monomial groups.  The same idea
produces scalar-center selection rules: if $\zeta I$ lies in the group, a
one-sided coefficient vanishes unless $\zeta^{|\tau|}=1$ for every allowed
$\zeta$.

\begin{theorem}[Normal-abelian Fourier--Reynolds reduction]
Let $G$ be finite, let $A\triangleleft G$ be abelian, put $Q=G/A$, and let
$W$ be a complex $G$-module. With
$P_A=|A|^{-1}\sum_{a\in A}\rho_W(a)$ and any section $s:Q\to G$,
\[
 P_G=\frac1{|Q|}\sum_{q\in Q}\rho_W(s(q))P_A,
 \qquad
 \dim W^G=\frac1{|Q|}\sum_{q\in Q}
 \Tr\bigl(\rho_W(s(q))\mid W^A\bigr).
\]
The expression is independent of the section.
\end{theorem}
\begin{proof}
Normality makes $P_A$ commute with $G$. Expanding the displayed average gives
$|G|^{-1}\sum_{g\in G}\rho_W(g)$. On $W^A$ the subgroup $A$ acts
trivially, so changing a section representative does not change the trace.
\end{proof}
This single reduction contains the finite-abelian, code-monomial,
$G(r,p,n)$, generalized-dihedral, and several Pauli/extraspecial filters as
special cases.

\section{Cycle indices, fixed points, and M\"obius inversion}

If $G$ permutes a finite set $X$ and $c_d(g;X)$ is the number of $d$-cycles,
then a cycle contributes the determinant factor $1-t^d$.  Hence
\begin{equation}\label{eq:tool-permutation}
\boxed{\MGF_{G,X}=\frac1{|G|}\sum_{g\in G}
\prod_{i,d}(1-t_i^d)^{-c_d(g;X)}.}
\end{equation}
The coefficient has an independent orbit interpretation:
\[
[m_\tau]\MGF_{G,X}
=\#\left(G\backslash\prod_j\Sym^{\tau_j}X\right).
\]
This is often the cleanest verification route, since it avoids eigenvalues
entirely.

When cycles are difficult to enumerate directly but fixed points are easy,
use
\[
|\Fix_X(g^r)|=\sum_{d\mid r}d\,c_d(g;X),
\qquad
c_d(g;X)=\frac1d\sum_{r\mid d}\mu(d/r)|\Fix_X(g^r)|.
\]
This fixed-points-to-cycles transform is the principal tool for subset,
Hamming, Grassmann, polar, affine, and finite Lie-type actions.  The geometry
enters only in the fixed-point count: solve a linear system, count invariant
subspaces, or impose an isotropic/singular predicate.

\section{Wreath products and the block-cycle determinant}

Let $H\le GL(V)$ and consider $H\wr S_n$ on $V^{\oplus n}$.  Around a
permutation cycle $C=(i_1\cdots i_r)$, multiply the labels in cyclic order,
$h_C=h_{i_r}\cdots h_{i_1}$.  Direct block elimination gives
\[
\det(I-tg_C)=\det(I-t^rh_C).
\]
Different cycles are independent.  The exponential formula for labelled
cycles therefore yields
\begin{equation}\label{eq:tool-wreath}
\boxed{\sum_{n\ge0}u^n\MGF_{H\wr S_n}(\mathbf t)
=\exp\!\left(\sum_{r\ge1}\frac{u^r}{r}
\MGF_H(t_1^r,t_2^r,\ldots)\right).}
\end{equation}
This is also the recursive mechanism behind rooted-tree automorphism groups:
isomorphic branches are permuted, so every vertex layer contributes a wreath
product.

\section{Conjugacy-class and spectral compression}

If the integrand in \eqref{eq:tool-finite-molien} is a class function, replace
the element sum by
\[
\frac1{|G|}\sum_{C}|C|\prod_i\det(I-t_i\rho(g_C))^{-1}.
\]
The proof task becomes a classification of spectra.  This produces the exact
formulas for cyclic, dihedral, dicyclic, polyhedral, Pauli/extraspecial,
sporadic frame-shape, and many braid-image examples.  A reliable derivation
records three pieces separately:
\begin{itemize}
\item the number of elements in each spectral sector;
\item the eigenvalue multiset in the represented module; and
\item the represented image and scalar center. By
Theorem~\ref{thm:normalized-image-invariance}, normalized Haar/counting
averages may be taken directly over the image; the kernel affects only
unnormalized or non-Haar weights.
\end{itemize}
For a regular representation, an element of order $d$ acts as $|G|/d$
disjoint $d$-cycles.  Grouping only by element order immediately gives
\[
\MGF_{G,\mathrm{Reg}}=
\sum_{d\ge1}\frac{N_d(G)}{|G|}
\ExpS\!\left(\frac{|G|}{d}p_d\right).
\]

\section{Weyl integration and constant-term extraction}

For compact connected $G$, choose a maximal torus $T$, Weyl group $W$, and
positive roots $\Phi^+$.  If the weights of $V$ are $\mu$ with multiplicities
$m(\mu)$, Weyl integration converts the Haar integral into
\begin{equation}\label{eq:tool-weyl}
\boxed{\MGF_{G,V}=\frac1{|W|}\CT_z\left[
\prod_{\alpha\in\Phi^+}(1-z^\alpha)(1-z^{-\alpha})
\prod_{i,\mu}(1-t_i z^\mu)^{-m(\mu)}\right].}
\end{equation}
The constant term imposes weight zero; the Weyl denominator performs the
alternating projection required to pass from torus weights to $G$-invariants.
This is the exact engine for compact classical groups, exceptional groups, and
maximal compact corrections.  Low-degree checks can be made independently by
constructing invariant bilinear, alternating, cubic, or quartic tensors and by
computing Lie-algebra kernels.

\section{Schur functors and plethystic pushforward}

For a partition $\lambda\vdash d$, the Frobenius character formula gives the
trace of a Schur-functor representation:
\begin{equation}\label{eq:tool-schur}
\boxed{\Tr(S_\lambda(g)^r)=s_\lambda[X(g^r)]
=\sum_{\alpha\vdash d}\frac{\chi^\lambda(\alpha)}{z_\alpha}
\prod_{a\in\alpha}\Tr(g^{ra}).}
\end{equation}
Substitute \eqref{eq:tool-schur} into \eqref{eq:tool-plethystic}.  In stable
orthogonal and symplectic rank, the base invariant series are respectively
$\ExpS(h_2)$ and $\ExpS(e_2)$, so the Schur-functor coefficients become Hall
inner products against $h_\tau[s_\lambda]$.  A genuinely independent finite
check constructs $S_\lambda(V)$ with Young symmetrizers and averages the
resulting matrices.

\section{Central characters and balanced two-sided series}

A continuous scalar center often forces every positive-degree holomorphic
invariant to vanish.  For example, if $U(1)$ acts on $V$ with weight one, then
$z\in U(1)$ acts on $W_\tau(V)$ by $z^{|\tau|}$, so the one-sided series is
just $1$.  The correct nontrivial object pairs $V$ and $V^*$:
\[
b_{\alpha,\beta}
=\dim\left(\Sym^\alpha V\otimes\Sym^\beta V^*\right)^G,
\qquad b_{\alpha,\beta}=0\ \text{unless }|\alpha|=|\beta|.
\]
The corresponding integrand contains both
$\det(I-s_i\rho(g))^{-1}$ and
$\det(I-t_j\rho(g)^{-\!*})^{-1}$.  This correction is essential for Clifford
groups with all scalars, local unitary tensor products, unitary Schur
functors, and conjugation on endomorphisms.  Frame potentials arise as the
multilinear diagonal coefficients $\E|\Tr g|^{2k}$.

\begin{theorem}[Central-charge lattice filter]
Let a central subgroup $Z$ act on modules $V_1,\ldots,V_s$ by characters
$\omega_1,\ldots,\omega_s$. Then
\[
 \left(\bigotimes_{j=1}^s\Sym^{d_j}V_j\right)^G=0
 \quad\text{unless}\quad
 \prod_{j=1}^s\omega_j(z)^{d_j}=1\quad(z\in Z).
\]
For $V$ and $V^*$ this is the charge condition $|\alpha|-|\beta|$. Exact
charge zero is invariant under arbitrary projective rephasing; a finite
congruence sector records a chosen scalar lift.
\end{theorem}
\begin{proof}
Every vector in the displayed homogeneous component is an eigenvector for
$Z$ with the stated product character. A $G$-fixed vector must in particular
be fixed by $Z$.
\end{proof}

\section{Nonuniform measures, convolution, and heat flow}

For a probability measure $\nu$ that is not Haar, averaging is generally not
a projection.  The coefficient is
\[
[m_\tau]\MGF_{\nu,V}
=\Tr\!\left(\int_G\rho_{W_\tau}(g)\,d\nu(g)\right),
\]
so it need not be a nonnegative integer.  For a finite random walk
$\nu=\kappa^{*r}$, Fourier transforms turn convolution into a power:
$\widehat\nu(W)=\widehat\kappa(W)^r$.  For compact heat kernels, irreducible
Fourier modes are attenuated by their Casimir eigenvalues.  On $U(1)$ this is
the explicit weight filter
\[
[m_\tau]\MGF_t=\sum_{w\in\mathbb Z}M_\tau(w)e^{-w^2t/2},
\]
which interpolates from total dimension at $t=0$ to Haar invariant dimension
as $t\to\infty$.

\section{Central extensions, projective representations, and scalar filters}

The representation, not merely the abstract group, is part of the input.
For a central extension $1\to Z\to\widetilde G\to G\to1$ and an honest
representation $V$ of $\widetilde G$, a central element acting by
$\omega(z)$ acts on $W_\tau(V)$ by $\omega(z)^{|\tau|}$.  Hence the one-sided
coefficient vanishes unless that central character is trivial.  If the lift
is only projective, the intrinsic replacement is either the adjoint series or
the balanced two-sided series: scalar phases cancel in
$V\otimes V^*$.  This distinction is essential for spin covers, Weil lifts,
Clifford groups, braid images, and sporadic covers.

\section{Frame shapes and lattice automorphism groups}

For a finite-order integral matrix $g$ with frame shape
$\prod_{d\ge1}d^{k_d(g)}$, one has
\[
 \det(I-tg)=\prod_{d\ge1}(1-t^d)^{k_d(g)},\qquad
 \MGF_{G,V}=\frac1{|G|}\sum_g\prod_{i,d}(1-t_i^d)^{-k_d(g)}.
\]
Power traces and frame exponents determine one another by M\"obius inversion.
This gives an exact class-compressed route for root lattices, repeated
orthogonal sums, tensor products, and isolated lattices, while keeping
rank-tower limits separate from finite exceptional examples.

\section{Multiplicative wreath actions and tensor-induced modules}

The additive block representation of $H\wr S_n$ is governed by the
block-cycle determinant above.  Product actions and tensor-induced modules
behave differently: their fixed points or characters multiply over coordinate
cycles.  For $g=(h_1,\ldots,h_n;\pi)$, the data needed from a cycle $C$ are
its length and the conjugacy class of the cycle product $h_C$.  Powering splits
$C$ into $\gcd(|C|,r)$ cycles and powers the corresponding cycle product.
This class-power rule is the common engine for product actions, twisted wreath
actions, Fock modules, and multipartition representations.

\section{Fast asymptotic diagnostics}

Before asserting a coefficientwise limit, inspect the first orbit/moment
coefficients.  In a permutation module,
\[
 [m_{(1,1)}]\MGF_{G,X}=\#(G\backslash X^2).
\]
For a general complex module, $[m_{(1,1)}]=\dim(V\otimes V)^G$; it
measures self-duality rather than the usual frame potential. The balanced
second moment is $\mathbb E|\Tr\rho(g)|^2=\dim(V\otimes V^*)^G$. Fourth
moments often expose spinor or tensor-growth obstructions.  If one fixed coefficient
diverges, the raw sigma-MGF has no finite coefficientwise limit.  When all
bounded-degree data depend on finitely many cycle counts, Poisson or
compound-Poisson limits are plausible; otherwise the correct statement may
require normalization, fixed-tail shapes, fixed level, or no universal limit.

\section{Stable range, coefficientwise limits, and honest claim labels}

Three statements that look similar must not be conflated:
\begin{itemize}
\item an \emph{exact finite-rank identity}, valid for a specified $n$;
\item a \emph{stable-range identity}, whose coefficient of total degree $d$
agrees once $n$ exceeds a proved threshold; and
\item a \emph{distributional or coefficientwise limit}, which may have no
finite stage at which equality holds.
\end{itemize}
Stability is usually proved by showing that degree-$d$ data only sees at most
$d$ labels, blocks, or trace modes.  It can fail immediately when geometry
creates new low-degree orbits, as for complete flags or the binary-cube action.
Poisson cycle limits and Gaussian trace limits are therefore stated as limits,
not silently promoted to finite formulas.

\begin{proposition}[Integer convergence implies eventual stability]
For finite or compact Haar models, if a fixed coefficient
$a_n=[m_\tau]\MGF_{G_n,V_n}$ converges to a finite real limit, then $a_n$ is
eventually constant and the limit is a nonnegative integer.
\end{proposition}
\begin{proof}
Each $a_n$ is an invariant-space dimension. A convergent integer sequence is
eventually contained in an interval of length less than one.
\end{proof}

\begin{theorem}[Bounded-support orbit stability]\label{thm:bounded-support-orbit-stability}
Suppose every degree-$D$ configuration is contained in a supported subobject
of size or rank at most $R(D)$, and suppose that for $n\ge N(R)$ every
isomorphism between two supported configurations of rank at most $R$ extends
to an element of $G_n$. Then every coefficient of total degree $D$ in the
permutation $\sigma$-MGF is independent of $n$ for
$n\ge N(R(D))$.
\end{theorem}
\begin{proof}
The coefficient counts orbits of configurations. Bounded support leaves only
finitely many abstract support types, and the extension hypothesis makes
ambient equivalence exactly support isomorphism.
\end{proof}
A support bound alone is not a stability proof; the extension property must be
verified for each family.

\section{Verification hierarchy and proof-audit checklist}

The bundled notes use four distinct levels of evidence.  These labels are not
interchanged:
\begin{enumerate}
\item \textbf{Proof:} a symbolic argument valid for every parameter in the
stated range.
\item \textbf{Independent exact verification:} a second exact computation
based on a different mathematical object, such as orbit enumeration versus a
cycle formula, modular weight counting versus complex matrix averaging, or a
Lie-algebra kernel versus a character formula.
\item \textbf{Cross-check with different implementation failure modes:} two
computations share some mathematical input, such as an eigenvalue list or
class table, but use different coefficient algorithms.
\item \textbf{Numerical diagnostic:} floating-point evaluation, seeded Haar
sampling, or evaluation of an identity at finitely many parameter values.
\end{enumerate}
Every reported result should state which category it belongs to.  Finite
testing is never substituted for an all-parameter proof.  When a rank or
nullity is obtained numerically, the certificate reports the threshold, the
largest singular or eigenvalue classified as zero, and the smallest one
classified as positive.  If no clear spectral gap is present, the result is
labelled inconclusive rather than rounded to a rank.

\resultbox{\textbf{Fast routing guide.}  Diagonal abelian action: use a Fourier
zero-weight filter.  Permutation action: count fixed points, invert to cycles,
then use the cycle index.  Block permutation: use the wreath determinant.
Finite matrix group: compress by spectral classes.  Compact connected group:
use Weyl constant terms.  Continuous scalar center: switch to a balanced
two-sided series.  Non-Haar law: compute the averaging operator rather than an
invariant dimension.}

\clearpage
\part{Finite abelian, monomial, dihedral, and wreath families}
\resultbox{\textbf{Standing convention.}
The notation $\MGF_{G,V}$ means $\smgv{G}{V}$. The coefficient identity,
generalized Molien formula, and permutation-orbit formula are imported
Tools~\ref{tool:coefficient}--\ref{tool:permutation-orbits}; each section
starts at the first family-specific reduction.}
\section{Finite abelian zero-weight theorem}\label{proof:finite-abelian}
\begingroup
\definecolor{ink}{HTML}{17324D}
\definecolor{accent}{HTML}{0F766E}
\definecolor{soft}{HTML}{EEF6F5}
\providecommand{\CT}{\operatorname{CT}}
\renewcommand{\CT}{\operatorname{CT}}
\providecommand{\Sym}{\operatorname{Sym}}
\renewcommand{\Sym}{\operatorname{Sym}}
\providecommand{\Exp}{\operatorname{Exp}}
\renewcommand{\Exp}{\operatorname{Exp}}
\providecommand{\one}{\mathbf 1}
\renewcommand{\one}{\mathbf 1}
\setcounter{theorem}{0}
\setcounter{proposition}{0}
\setcounter{lemma}{0}
\setcounter{corollary}{0}
\setcounter{remark}{0}
\thispagestyle{plain}

\begin{abstract}
For a finite abelian group acting diagonally, the coefficient of the monomial
symmetric function $m_\tau$ in the sigma moment-generating function is the
number of monomial basis vectors whose total character is trivial.  This note
proves the resulting simultaneous-congruence formula, derives the diagonal
Molien average and finite constant-term form, and records a direct matrix test.
The accompanying marimo notebook constructs the represented group, builds the
full tensor product of symmetric powers, forms the Reynolds projector, and
compares its rank with an exact lattice count and a numerical character average.
\end{abstract}

\subsection{Setup and the represented matrix group}

Let
\[
 A=C_{m_1}\times\cdots\times C_{m_r},\qquad
 M=\widehat A\cong
 \mathbb Z/m_1\mathbb Z\oplus\cdots\oplus\mathbb Z/m_r\mathbb Z.
\]
Every complex representation of $A$ splits into one-dimensional characters.
Choose a weight basis and write
\[
 V=\bigoplus_{j=1}^{d}\mathbb C_{w_j},\qquad
 w_j=(w_{1j},\ldots,w_{rj})\in M.
\]
The weights are the columns of an integer matrix $W=(w_{\ell j})\in
\mathbb Z^{r\times d}$, understood row by row modulo $m_\ell$.  For
$a=(a_1,\ldots,a_r)$, the corresponding matrix is
\begin{equation}\label{finite-abelian-zero-weight-proof:eq:represented-group}
 \rho(a)=\operatorname{diag}_{1\le j\le d}
 \left(\exp\left(2\pi i\sum_{\ell=1}^{r}
 \frac{a_\ell w_{\ell j}}{m_\ell}\right)\right).
\end{equation}
Thus the concrete matrix group is obtained by letting
$0\le a_\ell<m_\ell$ in \eqref{finite-abelian-zero-weight-proof:eq:represented-group}. If $\rho$ has a
kernel, several abstract elements give the same matrix; every image matrix has
exactly $|\ker\rho|$ preimages. Hence the normalized average over $A$ equals
the normalized average over $\rho(A)$ by
Theorem~\ref{thm:normalized-image-invariance}.

For a partition or finite tuple $\tau=(\tau_1,\ldots,\tau_s)$, put
\[
 T_\tau(V)=\bigotimes_{b=1}^{s}\Sym^{\tau_b}V.
\]
The empty tuple gives $T_\varnothing(V)=\mathbb C$.

\subsection{The simultaneous-congruence formula}

\begin{theorem}[Zero-weight formula]\label{finite-abelian-zero-weight-proof:thm:main}
With the notation above,
\begin{equation}\label{finite-abelian-zero-weight-proof:eq:count}
 \dim T_\tau(V)^A=
 \#\left\{(q_{bj})\in\mathbb Z_{\ge0}^{s\times d}:
 \begin{array}{ll}
 \displaystyle\sum_{j=1}^{d}q_{bj}=\tau_b
   & (1\le b\le s),\\[4pt]
 \displaystyle\sum_{b=1}^{s}\sum_{j=1}^{d}
 w_{\ell j}q_{bj}\equiv0\pmod {m_\ell}
   & (1\le\ell\le r)
 \end{array}\right\}.
\end{equation}
Equivalently, for $q_b=(q_{b1},\ldots,q_{bd})^T$ and
$D=\operatorname{diag}(m_1,\ldots,m_r)$,
\begin{equation}\label{finite-abelian-zero-weight-proof:eq:lattice}
 \dim T_\tau(V)^A=
 \#\left\{(q_1,\ldots,q_s):
 q_b\in\mathbb Z_{\ge0}^{d},\quad
 \one^Tq_b=\tau_b,\quad
 W(q_1+\cdots+q_s)\in D\mathbb Z^r\right\}.
\end{equation}
\end{theorem}

\begin{proof}
The monomials
\[
 e_1^{q_{b1}}\cdots e_d^{q_{bd}},\qquad
 q_{bj}\ge0,\quad \sum_jq_{bj}=\tau_b,
\]
form a basis of $\Sym^{\tau_b}V$.  Tensoring these bases over $b$ gives a
basis of $T_\tau(V)$ indexed by the arrays in the first line of
\eqref{finite-abelian-zero-weight-proof:eq:count}.  Since $e_j$ has character $w_j$, the basis vector indexed
by $(q_{bj})$ has total character
\[
 \sum_{b=1}^{s}\sum_{j=1}^{d}q_{bj}w_j\in M.
\]
It is fixed by every element of $A$ exactly when that character is zero.  In
coordinates this is precisely the second line of \eqref{finite-abelian-zero-weight-proof:eq:count}.  Because
the displayed monomial basis is a weight basis, the invariant subspace is
spanned by exactly those basis vectors of weight zero.  Counting them proves
\eqref{finite-abelian-zero-weight-proof:eq:count}; collecting the $q_b$ and writing divisibility by each
$m_\ell$ as membership in $D\mathbb Z^r$ gives \eqref{finite-abelian-zero-weight-proof:eq:lattice}.
\end{proof}

\subsection{Fourier form of the imported sigma-MGF}

Tools~\ref{tool:coefficient} and~\ref{tool:molien} give, without a new
Reynolds--Molien proof,
\begin{align}
 \mathbb E_A\!\left[\Exp_\sigma(X_Ah_1)\right]
 &=\frac1{|A|}\sum_{a\in A}
 \prod_{i\ge1}\prod_{j=1}^{d}
 \frac{1}{1-t_i\chi_{w_j}(a)},
 \label{finite-abelian-zero-weight-proof:eq:mgf}\\
 [m_\tau]\,\mathbb E_A\!\left[\Exp_\sigma(X_Ah_1)\right]
 &=\frac1{|A|}\sum_{a\in A}
 \prod_{b=1}^{s}h_{\tau_b}
 \bigl(\chi_{w_1}(a),\ldots,\chi_{w_d}(a)\bigr)
 =\dim T_\tau(V)^A.
 \label{finite-abelian-zero-weight-proof:eq:molien-coefficient}
\end{align}
The family-specific step is to evaluate this imported average by finite
Fourier projection.  Work in the group algebra with
$z_\ell^{m_\ell}=1$ and put
$z^{w_j}=\prod_\ell z_\ell^{w_{\ell j}}$.  If $\CT_A$ extracts the coefficient
of the identity character, then  Work in the group algebra
with $z_\ell^{m_\ell}=1$ and put
$z^{w_j}=\prod_\ell z_\ell^{w_{\ell j}}$.  If $\CT_A$ extracts the coefficient
of the identity character, then
\begin{align}
 \mathbb E_A\!\left[\Exp_\sigma(X_Ah_1)\right]
 &=\CT_A\prod_{i\ge1}\prod_{j=1}^{d}(1-t_i z^{w_j})^{-1},\\
 [m_\tau]\,\mathbb E_A\!\left[\Exp_\sigma(X_Ah_1)\right]
 &=\CT_A\prod_{b=1}^{s}h_{\tau_b}
 (z^{w_1},\ldots,z^{w_d}).
\end{align}
Character orthogonality identifies this operator with the finite Fourier
projector
\[
 \CT_A F=\frac1{|A|}\sum_{a\in A}
 F\bigl(\zeta_{m_1}^{a_1},\ldots,\zeta_{m_r}^{a_r}\bigr).
\]

\subsection{Verification strategy}

The accompanying notebook uses three routes with different failure modes.

\begin{enumerate}[leftmargin=*,label=\textbf{\arabic*.}]
\item \textbf{Direct represented matrices.}
It constructs all $|A|=\prod_\ell m_\ell$ matrices in
\eqref{finite-abelian-zero-weight-proof:eq:represented-group}.  It first checks the identity, unitarity, and
homomorphism laws numerically for every pair of represented elements.  In the
monomial basis, the diagonal of
$\Sym^n(\rho(a))$ is indexed by weak compositions $q$ of $n$ and has entry
$\prod_j\rho(a)_{jj}^{q_j}$.  Kronecker products give the full matrices on
$T_\tau(V)$.  Their average is $P_A$; its numerical rank and trace give the
target quantity directly.  The residual $\lVert P_A^2-P_A\rVert_F$ checks
that the average is a projection.

\item \textbf{Exact zero-weight enumeration.}
It enumerates the same weak compositions using integer arithmetic only and
counts those for which every row of $W(q_1+\cdots+q_s)$ is divisible by the
corresponding $m_\ell$.  This is a literal implementation of
\eqref{finite-abelian-zero-weight-proof:eq:lattice}; no complex roots of unity are used.

\item \textbf{Character/Molien average.}
It computes
\[
 \frac1{|A|}\sum_{a\in A}
 \prod_b\operatorname{tr}\bigl(\Sym^{\tau_b}(\rho(a))\bigr)
\]
in floating-point complex arithmetic.  The result must be within $10^{-8}$ of
the exact count.  This route avoids the large Kronecker matrices while still
testing the Fourier projection independently of the modular test.
\end{enumerate}

A case passes when the projector rank, its rounded trace, the exact count, and
the rounded Molien average agree, and both the numerical discrepancy and
idempotence residual are below $10^{-8}$.

\subsection{Representative results and size choice}

The default case is
\[
 A=C_3\times C_4,\qquad
 W=\begin{pmatrix}1&0&2\\0&1&3\end{pmatrix},\qquad
 d=3,\qquad \tau=(3,2).
\]
Here $|A|=12$ and
\[
\dim T_{(3,2)}(V)=
\binom{3+3-1}{3}\binom{2+3-1}{2}=10\cdot6=60.
\]
\clearpage
\thispagestyle{fancy}
The notebook constructs twelve $60\times60$ target matrices.  It obtains
\[
 \operatorname{rank}P_A=2,\qquad
 \text{zero-weight count}=2,\qquad
 \text{Molien average}=2+3.8\times10^{-15}i,
\]
with maximum displayed numerical discrepancy $8.21\times10^{-15}$.

The automated suite broadens the coverage as follows.
\begin{center}
\small
\begin{tabular}{@{}lrrrrc@{}}
\toprule
Representation & $|A|$ & $d$ & largest $\tau$ & max. target dim. & checks \\
\midrule
$C_2$, trivial plus sign & 2 & 2 & $(3,2)$ & 12 & 5 \\
$C_4$, nonfaithful & 4 & 3 & $(2,1)$ & 18 & 4 \\
$C_2\times C_2$, all characters & 4 & 4 & $(2,1)$ & 40 & 4 \\
$C_3\times C_4$, mixed weights & 12 & 3 & $(3,2)$ & 60 & 5 \\
$C_2\times C_3$, repeated weights & 6 & 4 & $(2,2)$ & 100 & 4 \\
\midrule
\multicolumn{5}{r}{Total} & 22 \\
\bottomrule
\end{tabular}
\end{center}
All 22 checks pass.  The cases include the empty partition, trivial weights,
repeated weights, a nonfaithful representation, one and two cyclic factors,
and tensor products with repeated and unequal degrees.

In general the target dimension is
\begin{equation}\label{finite-abelian-zero-weight-proof:eq:target-dim}
 N=\prod_{b=1}^{s}\binom{\tau_b+d-1}{\tau_b}.
\end{equation}
Dense direct projection requires $O(N^2)$ memory and approximately
$O(|A|N^2)$ arithmetic, while literal congruence enumeration visits $N$
monomial basis vectors.  The notebook caps the dense target dimension at
$N=1000$ and uses $N=60$ by default.  For larger dimensions, the Molien trace
average or an exact convolution of residue distributions on
$\prod_\ell\mathbb Z/m_\ell\mathbb Z$ is the appropriate replacement for a
dense projector.

\subsection{Reproduction}

From the repository root, launch the notebook with
\begin{verbatim}
.venv/bin/marimo edit \
  marimo_notebooks/finite_abelian.py
\end{verbatim}
The controls accept comma-separated moduli and partition parts.  Rows of $W$
are comma-separated and separated from one another by semicolons.  For
example, the default matrix is entered as
\begin{verbatim}
1, 0, 2; 0, 1, 3
\end{verbatim}
The automated table is evaluated whenever the notebook runs.
\endgroup


\clearpage
\section{Cyclic-weight invariant distributions: correction and extension}\label{proof:cyclic-weight}
\begingroup
\definecolor{ink}{HTML}{17212B}
\definecolor{accent}{HTML}{2457A7}
\definecolor{soft}{HTML}{EEF4FC}
\providecommand{\C}{\mathbb C}
\renewcommand{\C}{\mathbb C}
\providecommand{\Z}{\mathbb Z}
\renewcommand{\Z}{\mathbb Z}
\providecommand{\Sym}{\operatorname{Sym}}
\renewcommand{\Sym}{\operatorname{Sym}}
\providecommand{\Exp}{\operatorname{Exp}}
\renewcommand{\Exp}{\operatorname{Exp}}
\setcounter{theorem}{0}
\setcounter{proposition}{0}
\setcounter{lemma}{0}
\setcounter{corollary}{0}
\setcounter{remark}{0}
\begin{abstract}
This note verifies the homogeneous-basis moments associated with a finite
cyclic representation. The answer depends on the multiset of weights of the
chosen representation, not only on the abstract cyclic group. We prove the
weight-congruence formula, derive its refined Molien series, specialize to the
cyclic permutation representation, and describe independent exact and
numerical tests implemented in the accompanying marimo notebook.
\end{abstract}

\subsection{Setup}

Let $C_m=\Z/m\Z=\langle R\rangle$ act on $V=\C^d$. Since $R^m=I$ and the
polynomial $x^m-1$ has distinct roots over $\C$, the operator $R$ is
diagonalizable. Fix an eigenbasis $u_1,\ldots,u_d$ such that
\[
  Ru_j=\omega^{a_j}u_j,
  \qquad \omega=e^{2\pi i/m},
  \qquad a_j\in\Z/m\Z.
\]
The multiset $(a_1,\ldots,a_d)$ is part of the representation data.

For a partition $\tau=(\tau_1,\ldots,\tau_s)$, define
\[
  W_\tau=\Sym^{\tau_1}(V)\otimes\cdots\otimes\Sym^{\tau_s}(V).
\]

\subsection{Invariant-count formula}

\begin{theorem}\label{cyclic-weight-distribution-proof:thm:count}
For the cyclic representation above,
\[
\dim(W_\tau^{C_m})=
\#\left\{
(q_{i,j})\in\Z_{\geq0}^{s\times d}:
\begin{array}{l}
\sum_{j=1}^{d}q_{i,j}=\tau_i\quad(1\leq i\leq s),\\[2pt]
\sum_{i=1}^{s}\sum_{j=1}^{d}a_jq_{i,j}\equiv0\pmod m
\end{array}
\right\}.
\]
\end{theorem}

\begin{proof}
A monomial basis of $\Sym^{\tau_i}(V)$ is indexed by the weak compositions
$(q_{i,1},\ldots,q_{i,d})$ of $\tau_i$. Hence a basis vector of $W_\tau$ has
the form
\[
  b_q=\bigotimes_{i=1}^{s}u_1^{q_{i,1}}\cdots u_d^{q_{i,d}}.
\]
Applying $R$ gives
\[
  Rb_q
  =\omega^{\sum_{i,j}a_jq_{i,j}}b_q.
\]
Thus every $b_q$ is an eigenvector, and it is fixed by the generator precisely
when its exponent is zero modulo $m$. A vector is fixed by $C_m$ if and only if
it is fixed by $R$. Counting the basis vectors with eigenvalue $1$ proves the
claim.
\end{proof}

By Tool~\ref{tool:coefficient},
$[m_\tau]\,\smg{C_m}=\dim(W_\tau^{C_m})$.  Thus
Theorem~\ref{cyclic-weight-distribution-proof:thm:count} is the
representation-specific coefficient formula; it extends the one-dimensional
cyclic calculation of Theorem~A(1) in the 14 July paper.

\subsection{Fourier evaluation of the imported sigma-MGF}

Tool~\ref{tool:molien} gives directly
\begin{align}
 \smg{C_m}(T)
 &=\mathbb E\!\left[\operatorname{Exp}_{\sigma}
 (X_{C_m}h_1)(T)\right]
 =\frac1m\sum_{k=0}^{m-1}
  \prod_{r\geq1}\prod_{j=1}^{d}
  \frac{1}{1-\omega^{ka_j}t_r},
 \label{cyclic-weight-distribution-proof:eq:sigma-mgf}\\
 \dim(W_\tau^{C_m})
 &=\frac1m\sum_{k=0}^{m-1}
   \prod_{i=1}^{s}h_{\tau_i}
   (\omega^{ka_1},\ldots,\omega^{ka_d}).
 \label{cyclic-weight-distribution-proof:eq:molien-coeff}
\end{align}
The only new reduction here is the roots-of-unity filter
\[
 \frac1m\sum_{k=0}^{m-1}\omega^{kL}
 =\begin{cases}1,&L\equiv0\pmod m,\\0,&L\not\equiv0\pmod m,\end{cases}
\]
which turns the imported average into the congruence count of
Theorem~\ref{cyclic-weight-distribution-proof:thm:count}.
In particular, extracting one monomial-symmetric coefficient gives the proposed
formula itself:
\[
\boxed{
[m_\tau]\,\smg{C_m}
=\mu_{X_{C_m}}(h_\tau)
=\dim(W_\tau^{C_m})
=\#\left\{(q_{i,j}):
\sum_jq_{i,j}=\tau_i,\ 
\sum_{i,j}a_jq_{i,j}\equiv0\pmod m
\right\}.
}
\]
All infinite products above are interpreted coefficientwise as formal symmetric
series, so no analytic convergence assertion is required.

\subsection{Cyclic permutation representation}

For the cyclic shift on $\C^n$,
\[
 R(e_j)=e_{j+1}\ (j<n),\qquad R(e_n)=e_1,
\]
the Fourier eigenbasis has weights $0,1,\ldots,n-1$. Thus, with $m=d=n$,
\[
 \dim(W_\tau^{C_n})=
 \#\left\{(q_{i,a}):
 \sum_{a=0}^{n-1}q_{i,a}=\tau_i,\quad
 \sum_{i=1}^{s}\sum_{a=0}^{n-1}a q_{i,a}\equiv0\pmod n
 \right\}.
\]

For $d=1$ and weight $a_1=1$, every row is forced to be
$q_{i,1}=\tau_i$, so the value is $1$ when $|\tau|\equiv0\pmod m$ and $0$
otherwise. The surviving value is $1$, not $m$, because the group average is
normalized by $1/m$.

\subsection{Verification strategy and tested range}

The accompanying notebook uses three computational paths.

\begin{enumerate}
\item \textbf{Direct enumeration.} Generate every weak composition of each
$\tau_i$ into $d$ parts, form their Cartesian product, and count arrays whose
total weight is zero modulo $m$.
\item \textbf{Exact residue convolution.} For each block, count monomials in
each residue class of $\Z/m\Z$ and cyclically convolve the block distributions.
The residue-zero entry is the answer. This path uses integer arithmetic.
\item \textbf{Molien average.} For every $k$, compute each
$h_{\tau_i}(\omega^{ka_1},\ldots,\omega^{ka_d})$ as the coefficient of
$z^{\tau_i}$ in $\prod_j(1-\omega^{ka_j}z)^{-1}$, multiply across blocks, and
average over $k$. This path uses complex arithmetic and is compared with the
exact integer count.
\end{enumerate}

The automated suite checks all $2\leq m\leq8$, all partitions of total degree
at most $7$ with at most three blocks, and four representation families at
each order: the faithful one-dimensional character, the $m$-dimensional
cyclic shift, a repeated zero-weight representation, and a nonuniform
repeated-weight pattern. These cases include dimensions from $1$ through $8$
and deliberately test representations that are not determined by the group
order alone.

\medskip
\noindent\fcolorbox{accent}{soft}{\parbox{0.92\linewidth}{
\textbf{Executed result.} All $868$ generated cases passed. The exact direct
count and exact residue convolution agreed in every case. The largest absolute
discrepancy between the numerical Molien average and the exact integer answer
was $2.2\times10^{-12}$.
}}
\medskip

\begin{remark}
Finite testing is evidence, not a substitute for Theorem~\ref{cyclic-weight-distribution-proof:thm:count}.
Its value is to detect normalization errors, mishandling of repeated weights,
incorrect tensor-block multiplication, and mistakes in extracting complete
homogeneous coefficients.
\end{remark}

\subsection{Conclusion}

The proposed weight-congruence formula is correct for arbitrary complex
representations of $C_m$. The cyclic-shift specialization follows by using the
full set of residues as weights. The refined Molien expression is also correct
when read coefficientwise, and normalized expectation is essential in the
one-dimensional specialization.
\endgroup


\clearpage
\section{Restricted monomial groups: dual-code filter}\label{proof:restricted-monomial}
\begingroup
\definecolor{ink}{HTML}{162235}
\definecolor{accent}{HTML}{245B78}
\definecolor{pale}{HTML}{EEF5F7}
\definecolor{passgreen}{HTML}{16794A}
\providecommand{\Sfun}{\SMG}
\renewcommand{\Sfun}{\SMG}
\providecommand{\Cyc}{\operatorname{Cyc}}
\renewcommand{\Cyc}{\operatorname{Cyc}}
\providecommand{\Sym}{\operatorname{Sym}}
\renewcommand{\Sym}{\operatorname{Sym}}
\providecommand{\PASS}{\textcolor{passgreen}{\textbf{PASS}}}
\renewcommand{\PASS}{\textcolor{passgreen}{\textbf{PASS}}}
\providecommand{\summarybox}[1]{\begin{center}\fcolorbox{accent}{pale}{\parbox{0.91\linewidth}{#1}}\end{center}}
\renewcommand{\summarybox}[1]{\begin{center}\fcolorbox{accent}{pale}{\parbox{0.91\linewidth}{#1}}\end{center}}
\setcounter{theorem}{0}
\setcounter{proposition}{0}
\setcounter{lemma}{0}
\setcounter{corollary}{0}
\setcounter{remark}{0}
\summarybox{\textbf{Outcome.}  The general dual-code formula and the specialized
$G(r,p,n)$ coefficient formula were compared with explicit matrix averages in
six representative cases.  All 180 coefficients through total degree six
agree.  A direct Reynolds projector for $G(2,2,3)$ on
$\Sym^2(V)\otimes V$ has rank and trace $1$, with idempotence error
$3.21\times10^{-16}$.}

\subsection{Reduction to the code-specific calculation}
By Tools~\ref{tool:coefficient} and~\ref{tool:molien},
\[
\MGF_{G,V}=\smgv{G}{V},\qquad [m_\tau]\MGF_{G,V}=\dim W_\tau(V)^G.
\]
The new work begins with the cycle spectrum of a code-constrained monomial
matrix and the resulting dual-code filter.

\subsection{Code-constrained monomial groups}

Fix $\zeta=e^{2\pi i/r}$, an $H$-stable additive code
$A\leq(\mathbb Z/r\mathbb Z)^n$, and
\[
G(A,H)=D(A)\rtimes H,\qquad
D(A)=\{\operatorname{diag}(\zeta^{x_1},\ldots,\zeta^{x_n}):x\in A\}.
\]
We use the convention $g=D(x)P_\pi$.  If $C$ is a cycle of $\pi$, the
restriction of $g$ to its coordinate subspace is a weighted cyclic shift.
Only the full traversal contributes to its characteristic polynomial, hence
\begin{equation}
\det(I-tg)=\prod_{C\in\Cyc(\pi)}
\left(1-\zeta^{x_C}t^{|C|}\right),\qquad
x_C=\sum_{i\in C}x_i.                                         \tag{3}
\end{equation}
Substitution in (1) proves the finite exact formula
\begin{equation}
\Sfun_{A,H}(\mathbf t)=\frac1{|H||A|}
\sum_{\pi\in H}\sum_{x\in A}
\prod_{a,C}\left(1-\zeta^{x_C}t_a^{|C|}\right)^{-1}.         \tag{4}
\end{equation}

\subsubsection{The dual-code filter}

Expand each factor geometrically and let $m_{C,a}\geq0$ be its selected
multiplicity.  Put $M_C=\sum_am_{C,a}$ and set $e_i=M_C$ for $i\in C$.  The
phase of a term is $\zeta^{x\cdot e}$.  For
\[
A^\perp=\{e\in(\mathbb Z/r)^n:x\cdot e=0\pmod r\text{ for all }x\in A\},
\]
finite character orthogonality gives
\[
\frac1{|A|}\sum_{x\in A}\zeta^{x\cdot e}=\mathbf1_{e\in A^\perp}.
\]
Therefore the coefficient indexed by $\tau=(\tau_1,\ldots,\tau_\ell)$ is
\begin{equation}
\frac1{|H|}\sum_{\pi\in H}
\#\left\{(m_{C,a}):
\begin{array}{l}
\sum_C|C|m_{C,a}=\tau_a\quad\text{for every }a,\\
e(\pi,m)\in A^\perp
\end{array}\right\}.                                        \tag{5}
\end{equation}
Equation (5) is the implemented general formula.  Notice that it contains no
matrix eigenvalue computation and no average over diagonal matrices.

For a residue $s\in\mathbb Z/r$, the equivalent root-of-unity filter is
\[
\Phi_{d,s}(\mathbf t)=\frac1r\sum_{\ell=0}^{r-1}
\zeta^{-\ell s}\prod_a(1-\zeta^\ell t_a^d)^{-1}.
\]
It collects those monomials for which $\sum_am_a\equiv s\pmod r$.

\subsection{Specialization to $G(r,p,n)$}

Assume $p\mid r$ and
\[
A_{r,p,n}=\{x\in(\mathbb Z/r)^n:\sum_i x_i=0\pmod p\}.
\]
The vectors $e_i-e_j\in A_{r,p,n}$ force a dual vector to have equal
coordinates.  The vectors $p e_i\in A_{r,p,n}$ then force the common coordinate
to be a multiple of $r/p$.  Thus
\begin{equation}
A_{r,p,n}^\perp
=\{q(r/p)(1,\ldots,1):0\leq q<p\}.                            \tag{6}
\end{equation}
The cycle-index theorem now yields
\begin{equation}
\Sfun_{G(r,p,n)}(\mathbf t)
=\sum_{q=0}^{p-1}[u^n]\exp\left(
\sum_{d\geq1}\frac{u^d}{d}\Phi_{d,qr/p}(\mathbf t)\right).  \tag{7}
\end{equation}
An equivalent coefficient algorithm regards each coordinate as carrying a
label $(m_1,\ldots,m_\ell)$ whose coordinate sum has residue $qr/p$.  An
$S_n$-orbit is a multiset of $n$ labels.  Bounded dynamic programming over
label multiplicities extracts exactly the coefficient of (7); this is the
second independent formula route in the software.

With only $t_1=t$ nonzero, coefficient extraction reduces to
\begin{equation}
\Sfun_{G(r,p,n)}(t,0,\ldots)=
\frac1{(1-t^{nr/p})\prod_{j=1}^{n-1}(1-t^{rj})}.              \tag{8}
\end{equation}
For $W(D_n)=G(2,2,n)$, the second residue sector requires every coordinate to
carry odd total multiplicity.  It cannot occur below degree $n$.  In
particular, it does not create a degree-one fixed vector; this is the key
zero-coordinate correction in the proposed discussion.

\subsection{Independent verification design}

\begin{enumerate}[leftmargin=*,itemsep=3pt]
\item Construct every matrix $D(x)P_\pi$.  Check the predicted order
$r^n n!/p$ in the $G(r,p,n)$ cases.
\item For every partition $\tau$ of every total degree from $0$ through $6$,
compute $h_{\tau_a}$ recursively from the eigenvalues of each matrix and
average $\prod_a h_{\tau_a}$ over the group.
\item Independently evaluate (5) for a general code case or the multiset
coefficient of (7) for $G(r,p,n)$.
\item For a selected target, explicitly embed normalized occupancy bases of
symmetric powers in tensor powers, form the full Reynolds matrix, and compare
its numerical rank and trace with the coefficient.
\end{enumerate}

These routes share only the mathematical input $(r,p,n)$ or $(A,H)$.  The
formula route does not call the matrix-average code.

\subsection{Representative cases and results}

There are 30 partitions through total degree six.  The following suite includes
an $H=C_3$ code group not covered by the full symmetric-group specialization,
the unrestricted wreath products, the $D_n$ parity sector, and two different
nontrivial divisors $p$ of $r$.

\begin{center}
\begin{tabularx}{\linewidth}{@{}Xrrrrc@{}}
\toprule
Case & $|G|$ & $\dim V$ & degree & checks & result\\
\midrule
even-sign code $\rtimes C_3$ & 12 & 3 & $0$--$6$ & 30 & \PASS\\
$G(2,1,2)$ & 8 & 2 & $0$--$6$ & 30 & \PASS\\
$G(2,2,3)=W(D_3)$ & 24 & 3 & $0$--$6$ & 30 & \PASS\\
$G(3,1,2)$ & 18 & 2 & $0$--$6$ & 30 & \PASS\\
$G(4,2,2)$ & 16 & 2 & $0$--$6$ & 30 & \PASS\\
$G(4,4,2)$ & 8 & 2 & $0$--$6$ & 30 & \PASS\\
\bottomrule
\end{tabularx}
\end{center}

Selected matched dimensions illustrate that the test is multivariate, not only
the ordinary Molien specialization.
\begin{center}
\begin{tabular}{@{}lrrrr@{}}
\toprule
Case & $(2)$ & $(3)$ & $(2,1)$ & $(3,2)$\\
\midrule
even-sign code $\rtimes C_3$ & 1 & 1 & 1 & 4\\
$G(2,1,2)$ & 1 & 0 & 0 & 0\\
$G(2,2,3)$ & 1 & 1 & 1 & 3\\
$G(3,1,2)$ & 0 & 1 & 1 & 0\\
$G(4,2,2)$ & 0 & 0 & 0 & 0\\
$G(4,4,2)$ & 1 & 0 & 0 & 0\\
\bottomrule
\end{tabular}
\end{center}

For $G(2,2,3)$ and $\tau=(2,1)$, the target dimension is
$\binom{4}{2}\binom{3}{1}=18$.  Direct averaging of the eighteen-dimensional
target matrices produced
\[
\operatorname{rank}P=1,\qquad \operatorname{tr}P=0.9999999999999998,
\qquad\|P^2-P\|_F=3.21\times10^{-16}.
\]
This agrees with both coefficient routes.

\subsection{Reproducibility and interpretation}

Run the command-line suite from the repository root with
{\scriptsize
\begin{verbatim}
.venv/bin/python -m for_this_guy.matrix_group_formula_verification.restricted_monomial.verification
\end{verbatim}
}
or open
\path{marimo_notebooks/restricted_monomial_verification.py}
with marimo.  Matrix
averages are accepted only when their imaginary part and distance from the
nearest integer are below $3\times10^{-7}$.  The observed residuals are much
smaller; the displayed projector error is at machine precision.

Finite tests do not replace the preceding proof.  Their role is to detect
indexing errors, wrong permutation conventions, missing residue sectors, and
incorrect treatment of zero-labeled coordinates.  Within the stated small
dimensions, no such discrepancy remains.
\endgroup


\clearpage
\section[Exact finite-n formula for G(r,p,n)]{Exact finite-$n$ formula for $G(r,p,n)$}\label{proof:grp}
\begingroup
\definecolor{Navy}{HTML}{17324D}
\definecolor{Teal}{HTML}{087E8B}
\definecolor{Pale}{HTML}{EDF6F7}
\definecolor{Rule}{HTML}{B8CDD2}
\providecommand{\Ms}{\SMG}
\renewcommand{\Ms}{\SMG}
\providecommand{\Sym}{\operatorname{Sym}}
\renewcommand{\Sym}{\operatorname{Sym}}
\providecommand{\ExpS}{\operatorname{Exp}_{\sigma}}
\renewcommand{\ExpS}{\operatorname{Exp}_{\sigma}}
\providecommand{\coeff}[1]{\left[#1\right]}
\renewcommand{\coeff}[1]{\left[#1\right]}
\providecommand{\C}{\mathbb C}
\renewcommand{\C}{\mathbb C}
\setcounter{theorem}{0}
\setcounter{proposition}{0}
\setcounter{lemma}{0}
\setcounter{corollary}{0}
\setcounter{remark}{0}
\colorbox{Pale}{\parbox{0.94\linewidth}{
\textbf{Result.} For $p\mid r$ and $n\geq1$, the auxiliary variable $u$ records the
number of coordinates exactly:
\[
\boxed{\displaystyle
\smg{G(r,p,n)}
=\coeff{u^n}\sum_{q=0}^{p-1}\ExpS\!\left(
u\sum_{j\geq0}h_{qr/p+jr}(\mathbf t)\right).}
\]
The supplied implementation constructs the matrix group, averages the target
character directly, and compares every partition through a chosen degree with an
independent bounded-multiset coefficient extraction.
}}

\subsection{Imported convention and the finite-size marker}
For $V=\C^n$, Tools~\ref{tool:coefficient} and~\ref{tool:molien} identify
\[
\MGF_{G(r,p,n)}=\smg{G(r,p,n)},\qquad
[m_\tau]\MGF_{G(r,p,n)}=\dim(\Sym^\tau V)^{G(r,p,n)}.
\]
The additional variable $u$ records the number of coordinates, with
\[
\ExpS\!\left(u\sum_{\mathbf v\in L}\mathbf t^{\mathbf v}\right)
=\prod_{\mathbf v\in L}(1-u\mathbf t^{\mathbf v})^{-1}.
\]
Thus the theorem below is an exact finite-$n$ strengthening, not a reproof of
the imported Molien identity.

\subsection{The monomial group and exact theorem}

Choose a primitive $r$th root $\omega$.  The group is
\[
G(r,p,n)=\left\{(\omega^{b_1},\ldots,\omega^{b_n};\sigma)
\in\mu_r^n\rtimes S_n:\ \sum_i b_i\equiv0\pmod p\right\},
\qquad p\mid r,
\]
acting by $e_i\mapsto\omega^{b_i}e_{\sigma(i)}$.  Its order is
$r^n n!/p$.  Set
\[
a_q=\frac{qr}{p},\qquad
H_q^{(r,p)}=\sum_{j\geq0}h_{a_q+jr},\qquad 0\leq q<p.
\]

\textbf{Theorem.}
For every $n\geq1$,
\[
\smg{G(r,p,n)}=\coeff{u^n}\sum_{q=0}^{p-1}
\ExpS\!\left(uH_q^{(r,p)}\right)
\tag{1}
\]
as a complete formal symmetric series.  Thus (1) has no stable-range
qualification.

\subsection{Proof}

Fix $\tau=(\tau_1,\ldots,\tau_s)$.  A monomial basis vector of
$\Sym^\tau V$ assigns to coordinate $i$ a label
\[
\mathbf v_i=(v_i(1),\ldots,v_i(s))\in\mathbb Z_{\geq0}^s,
\qquad \sum_{i=1}^n\mathbf v_i=\tau.
\]
Write $d_i=|\mathbf v_i|=\sum_jv_i(j)$.  The diagonal element with exponent
vector $\mathbf b$ acts by $\omega^{\sum_i d_i b_i}$.

\subsubsection*{The annihilator condition}

The character is trivial for every $\mathbf b$ satisfying
$\sum_i b_i\equiv0\pmod p$ if and only if there is a single
$q\in\{0,\ldots,p-1\}$ such that
\[
d_1\equiv\cdots\equiv d_n\equiv\frac{qr}{p}\pmod r.
\tag{2}
\]
Indeed, the allowed vector $e_i-e_j$ forces $d_i\equiv d_j\pmod r$.
Calling the common residue $c$, the allowed vector $p e_i$ forces
$pc\equiv0\pmod r$, so $c=qr/p$ modulo $r$.  Conversely, if (2) holds, then
\[
\sum_i d_i b_i\equiv\frac{qr}{p}\sum_i b_i\equiv0\pmod r.
\]
This proves both necessity and sufficiency, including $n=1$.

\subsubsection*{Permutation orbits become multisets}

After diagonal invariance is imposed, $S_n$ permutes the labels
$\mathbf v_1,\ldots,\mathbf v_n$.  The invariant dimension is therefore the
number of label multisets.  For a fixed $q$, their generating function is
\[
F_q(u,\mathbf t)=
\prod_{\substack{\mathbf v\in\mathbb Z_{\geq0}^s\\
|\mathbf v|\equiv a_q\ (\mathrm{mod}\ r)}}
\frac1{1-u\mathbf t^{\mathbf v}}.
\]
Here $u$ counts labels, and $\mathbf t$ records their componentwise sum.  Since
\[
\sum_{|\mathbf v|=d}\mathbf t^{\mathbf v}=h_d(\mathbf t),
\]
we have $F_q=\ExpS(uH_q^{(r,p)})$.  The residue classes for distinct $q$ are
disjoint, so summing $\coeff{u^n\mathbf t^\tau}F_q$ proves (1).

\subsection{Types \texorpdfstring{$B/C$ and $D$}{B/C and D}}

For $G(2,1,n)=W(B_n/C_n)$ only the even branch occurs:
\[
\Ms_{B_n/C_n}=\coeff{u^n}\ExpS\!\left(u(h_0+h_2+h_4+\cdots)\right).
\tag{3}
\]
The factor contributed by $h_0=1$ is $(1-u)^{-1}$; it supplies unused, or
zero-labelled, coordinates.

For $G(2,2,n)=W(D_n)$ both common residues occur:
\[
\Ms_{D_n}=\coeff{u^n}\left\{
\ExpS\!\left(u(h_0+h_2+h_4+\cdots)\right)
+\ExpS\!\left(u(h_1+h_3+h_5+\cdots)\right)\right\}.
\tag{4}
\]
The second branch has no zero label: all $n$ coordinates must have odd positive
degree.  Consequently it first contributes in total degree $n$.  This explains
why deleting $u$ gives a false degree-one term for $n\geq2$.

More generally, the first $q>0$ correction has every coordinate of degree at
least $r/p$, hence first appears in degree $nr/p$.  Before that threshold, the
$q=0$ branch has enough zero labels to stabilize, giving
\[
\Ms_{G(r,p,n)}\equiv\ExpS(h_r+h_{2r}+\cdots)
\quad\text{in total degrees }d<nr/p.
\tag{5}
\]

\subsection{Independent verification strategy}

The notebook does not verify (1) by evaluating its right-hand side twice.
It uses two computational paths.

\textbf{Direct matrix average.}
For every $\sigma\in S_n$ and exponent vector
$\mathbf b\in\{0,\ldots,r-1\}^n$ with $\sum b_i\equiv0\pmod p$, it constructs
the matrix
\[
g_{\sigma(i),i}=\omega^{b_i},\qquad g_{k,i}=0\quad(k\neq\sigma(i)).
\]
For eigenvalues $\lambda_1,\ldots,\lambda_n$, the character of $\Sym^dV$ is
computed from
\[
\sum_{d\geq0}\chi_{\Sym^dV}(g)z^d
=\prod_{i=1}^n(1-\lambda_i z)^{-1}.
\]
The observed integer is
\[
D_{\rm direct}(\tau)=\frac1{|G|}\sum_{g\in G}
\prod_j\chi_{\Sym^{\tau_j}V}(g).
\tag{6}
\]

\textbf{Bounded multiset extraction.}
For each $q$, the code lists only labels
$0\leq v(j)\leq\tau_j$ with $|\mathbf v|\equiv qr/p\pmod r$.  It processes each
label type once, selecting multiplicities from $0$ through $n$.  A dynamic
program indexed by ``labels used'' and ``componentwise degree'' extracts
$\coeff{u^n\mathbf t^\tau}F_q$ using exact Python integers.  This route never
constructs or diagonalizes a group matrix.

Floating-point arithmetic appears only in roots of unity and eigenvalues on the
direct side.  The group average is required to lie within $2\cdot10^{-8}$ of a
real integer before comparison; otherwise the verifier raises an error instead
of silently rounding.

\subsection{Representative results}

Every partition $\tau$ of every total degree $0\leq|\tau|\leq6$ was tested, for
30 coefficients per row of the table.

\begin{center}
\begin{tabular}{@{}lrrrl@{}}
\toprule
Group & Ambient dim. & Order & Coefficients & Result\\
\midrule
$G(2,1,2)=W(B_2/C_2)$ & 2 & 8 & 30 & pass\\
$G(2,2,4)=W(D_4)$ & 4 & 192 & 30 & pass\\
$G(3,1,2)$ & 2 & 18 & 30 & pass\\
$G(4,2,2)$ & 2 & 16 & 30 & pass\\
$G(4,4,2)$ & 2 & 8 & 30 & pass\\
\bottomrule
\end{tabular}
\end{center}

These cases exercise the zero-label mechanism, the odd $D_n$ branch, roots of
unity of orders $3$ and $4$, and several nontrivial values of $p$.  They are
finite checks rather than a proof, but they are structurally independent and
probe exactly the finite-$n$ issue addressed by $u$.

\subsection{Reproduction}

Open the marimo notebook
\href{run:marimo_notebooks/g_r_p_n.py}
{\texttt{marimo\_notebooks/g\_r\_p\_n.py}} and press
\emph{Construct matrices and verify}.  The standalone engine is
\texttt{verification.py}; running it from the repository root executes the five
representative cases and exits nonzero on any discrepancy.
\endgroup


\clearpage
\section[Exact finite-n wreath-product formula]{Exact finite-$n$ wreath-product formula}\label{proof:wreath-finite}
\begingroup
\definecolor{Navy}{HTML}{17324D}
\definecolor{Violet}{HTML}{6441A5}
\definecolor{Pale}{HTML}{F2EFF9}
\providecommand{\Ms}{\SMG}
\renewcommand{\Ms}{\SMG}
\providecommand{\Sym}{\operatorname{Sym}}
\renewcommand{\Sym}{\operatorname{Sym}}
\providecommand{\coeff}[1]{\left[#1\right]}
\renewcommand{\coeff}[1]{\left[#1\right]}
\providecommand{\C}{\mathbb C}
\renewcommand{\C}{\mathbb C}
\setcounter{theorem}{0}
\setcounter{proposition}{0}
\setcounter{lemma}{0}
\setcounter{corollary}{0}
\setcounter{remark}{0}
\colorbox{Pale}{\parbox{0.94\linewidth}{
\textbf{Result.} Let a finite matrix group $H\leq\mathrm{GL}(V)$ act blockwise
on $V^{\oplus n}$, with $S_n$ permuting the blocks.  Then
\[
\boxed{\displaystyle
\sum_{n\geq0}u^n\Ms_{H\wr S_n}(\mathbf t)
=\exp\!\left(\sum_{\ell\geq1}\frac{u^\ell}{\ell}
\Ms_H(t_1^\ell,t_2^\ell,\ldots)\right).}
\]
The notebook verifies fixed-$n$ coefficients by comparing explicit
block-monomial matrix averages with an exact rational cycle-index recurrence.
}}

\subsection{Set-up}

Let $H\leq\mathrm{GL}(V)$ be finite, with $d=\dim V$.  The wreath product
$H\wr S_n=H^n\rtimes S_n$ acts on
$V^{\oplus n}\cong V\otimes\C^n$.  For any finite matrix group $K$ define
\[
\Ms_K(\mathbf t)=\frac1{|K|}\sum_{k\in K}
\prod_{i\geq1}\det(I-t_i k)^{-1}.
\tag{1}
\]
If $\tau=(\tau_1,\ldots,\tau_s)$, then
\[
\coeff{m_\tau}\Ms_K
=\dim\!\left(\bigotimes_{j=1}^s\Sym^{\tau_j}W\right)^K,
\tag{2}
\]
where $W$ is the representation space of $K$.

For $\ell\geq1$ set
\[
A_\ell(\mathbf t)=\frac1{|H|}\sum_{h\in H}
\prod_{i\geq1}\det(I-t_i^\ell h)^{-1}
=\Ms_H(t_1^\ell,t_2^\ell,\ldots).
\tag{3}
\]

\subsection{Exact theorem}

\textbf{Theorem.}
As a formal power series in $u$ and the variables $\mathbf t$,
\[
\sum_{n\geq0}u^n\Ms_{H\wr S_n}(\mathbf t)
=\exp\!\left(\sum_{\ell\geq1}\frac{u^\ell}{\ell}A_\ell(\mathbf t)\right).
\tag{4}
\]
Consequently, at a prescribed number of blocks,
\[
\Ms_{H\wr S_n}(\mathbf t)=\coeff{u^n}
\exp\!\left(\sum_{\ell\geq1}\frac{u^\ell}{\ell}
\Ms_H(t_1^\ell,t_2^\ell,\ldots)\right).
\tag{5}
\]
No limit in $n$ and no stable-range assumption is used.

\subsection{Cycle-by-cycle proof}

Write $g=(h_1,\ldots,h_n;\sigma)$.  For a cycle
$c=(i_1\,i_2\,\cdots\,i_\ell)$ of $\sigma$, define its ordered cycle product
\[
h_c=h_{i_\ell}\cdots h_{i_2}h_{i_1}.
\]
On the corresponding $\ell$ summands of $V^{\oplus n}$, the block action gives
the standard determinant identity
\[
\det(I-zg_c)=\det(I-z^\ell h_c).
\tag{6}
\]
One way to see (6) is to solve the first $\ell-1$ block equations of
$(I-zg_c)x=0$ successively; closing the cycle leaves
$(I-z^\ell h_c)x_{i_1}=0$.  Thus the two determinants have the same zeros
with multiplicity.  Both have degree at most $\ell d$ and constant term one,
so their ratio is the constant one.  Equivalently, the identity follows by
block Gaussian elimination over $\mathbb C[z]$.

Since distinct cycles act on disjoint blocks, (6) yields
\[
\prod_{i\geq1}\det(I-t_i g)^{-1}
=\prod_{c}\prod_{i\geq1}\det(I-t_i^{|c|}h_c)^{-1}.
\tag{7}
\]
For fixed $\sigma$, averaging uniformly over $(h_1,\ldots,h_n)\in H^n$ makes
the cycle products for distinct cycles independent and uniform in $H$.  This
follows because, after all but one $h_i$ on a cycle are chosen, multiplication
by the remaining uniform factor makes $h_c$ uniform.  If $c_\ell(\sigma)$ is
the number of $\ell$-cycles, the conditional average of (7) is therefore
\[
\prod_{\ell\geq1}A_\ell(\mathbf t)^{c_\ell(\sigma)}.
\tag{8}
\]

Finally apply the cycle-index identity
\[
\sum_{n\geq0}\frac{u^n}{n!}\sum_{\sigma\in S_n}
\prod_{\ell\geq1}x_\ell^{c_\ell(\sigma)}
=\exp\!\left(\sum_{\ell\geq1}\frac{x_\ell u^\ell}{\ell}\right).
\tag{9}
\]
Substitution of $x_\ell=A_\ell$ proves (4).

\subsection{Roots of unity as a consistency check}

For the one-dimensional representation $H=\mu_r$,
\[
\Ms_{\mu_r}(\mathbf t)
=\frac1r\sum_{\zeta^r=1}\prod_i(1-\zeta t_i)^{-1}
=\sum_{j\geq0}h_{rj}(\mathbf t),
\tag{10}
\]
by the roots-of-unity filter.  Substituting (10) in (5) recovers the exact
$G(r,1,n)=\mu_r\wr S_n$ formula.  Thus the monomial-root case is the
one-dimensional specialization of the general block theorem.

\subsection{How the notebook verifies the theorem}

The two sides are implemented through different data structures and algorithms.

\subsubsection*{Direct block-matrix construction}

For each $\sigma\in S_n$ and $(h_1,\ldots,h_n)\in H^n$, the code creates an
$nd\times nd$ matrix whose only nonzero block in block column $i$ is
\[
g_{\sigma(i),i}=h_i.
\tag{11}
\]
It constructs all $|H|^n n!$ matrices.  For each matrix, complete symmetric
characters are generated by
\[
\sum_{a\geq0}\chi_{\Sym^a(V^{\oplus n})}(g)z^a
=\prod_{j=1}^{nd}(1-\lambda_jz)^{-1},
\tag{12}
\]
and the direct target is the normalized average of products of these characters,
as in (2).

\subsubsection*{Exact cycle-index recurrence}

Write the right side of (4) as $F(u)=\sum_{n\geq0}F_nu^n$.  Logarithmic
differentiation gives the useful fixed-$n$ recurrence
\[
nF_n=\sum_{\ell=1}^n A_\ell F_{n-\ell},
\qquad F_0=1.
\tag{13}
\]
This is the notebook's main coefficient tool.  For a target multi-degree
$\tau$, it stores only exponents bounded componentwise by $\tau$.  The
coefficient of $\mathbf t^\alpha$ in $A_\ell$ is zero unless every component
of $\alpha$ is divisible by $\ell$; when divisible, it is the direct base-group
coefficient at $\alpha/\ell$.  Truncated polynomial convolution and (13) are
performed with exact rational numbers.  The final denominator must cancel to
one, or the verifier raises an error.

This recurrence is useful beyond testing: it avoids enumerating partitions of
permutation cycle types and computes one finite-$n$ answer without expanding an
unbounded exponential.

\subsubsection*{Numerical guardrail}

Only the direct eigenvalue calculation uses floating point.  Each normalized
average must lie within $3\cdot10^{-8}$ of a real integer.  No comparison is
reported if that integrality check fails.  The formula side remains exact.

\subsection{Representative results}

All partitions through total degree five were checked: 19 coefficients for each
case.

\begin{center}
\begin{tabular}{@{}lrrrrl@{}}
\toprule
Base representation $H\curvearrowright V$ & $|H|$ & $\dim V$ & $n$ & $|H\wr S_n|$ & Result\\
\midrule
$C_2$ scalar on $\C$ & 2 & 1 & 3 & 48 & pass\\
$C_3$ scalar on $\C$ & 3 & 1 & 2 & 18 & pass\\
$C_2$ reflection on $\C^2$ & 2 & 2 & 2 & 8 & pass\\
$S_3$ standard on $\C^2$ & 6 & 2 & 2 & 72 & pass\\
\bottomrule
\end{tabular}
\end{center}

The last two cases test genuinely block-valued matrices.  The $C_2$ reflection
representation is reducible, while the standard $S_3$ representation is
two-dimensional, irreducible, and has a nonabelian base group.  Agreement in
these cases tests content not covered by the scalar roots-of-unity specialization.
\endgroup


\clearpage
\section{$\sigma$-MGF for arbitrary wreath products}\label{proof:wreath-general}
\begingroup
\definecolor{navy}{RGB}{24,55,95}
\providecommand{\Sym}{\operatorname{Sym}}
\renewcommand{\Sym}{\operatorname{Sym}}
\providecommand{\Exp}{\operatorname{Exp}}
\renewcommand{\Exp}{\operatorname{Exp}}
\providecommand{\cS}{\SMG}
\renewcommand{\cS}{\SMG}
\providecommand{\cW}{\SMG}
\renewcommand{\cW}{\SMG}
\setcounter{theorem}{0}
\setcounter{proposition}{0}
\setcounter{lemma}{0}
\setcounter{corollary}{0}
\setcounter{remark}{0}
\begin{abstract}
For a finite matrix group $\rho:H\to GL(V)$, this note proves the exact
generating function for the block-monomial representation of $H\wr S_n$ on
$V^{\oplus n}$ and its coefficientwise stable limit
$\Exp_\sigma(\cS_H-1)$.  A precise degreewise interpretation explains the
marked compound-Poisson language.  The companion Marimo notebook constructs
all wreath-product matrices and compares their generalized Molien
coefficients with two independent exact recurrences.  Tests for the sign
representation of $C_2$ and the rational two-dimensional representation of
$C_3$ checked 38 finite coefficients, 19 stable coefficients, and every
small cycle-product determinant identity; all passed exactly.
\end{abstract}

\subsection{Setup and target formula}
Let $\rho:H\to GL(V)$ be a finite matrix group, $\dim V=d$, and let
\[
 G_n=H\wr S_n=H^n\rtimes S_n
\]
act on $V_n=V^{\oplus n}$ by block-monomial matrices.  Write
\[
 \cW_{H,n}(\mathbf t)=\frac1{|H|^n n!}\sum_{g\in H\wr S_n}
 \prod_{a\ge1}\det(I-t_ag\mid V_n)^{-1}
\]
and
\[
 \cS_H(\mathbf t)=\frac1{|H|}\sum_{h\in H}
 \prod_{a\ge1}\det(I-t_ah\mid V)^{-1}.
\]
For the Adams substitution $\psi_rF(\mathbf t)=F(t_1^r,t_2^r,\ldots)$, the
claims under test are
\begin{equation}\label{wreath-products-proof-and-verification:eq:exact}
 \boxed{\displaystyle \sum_{n\ge0}\cW_{H,n}(\mathbf t)u^n
 =\exp\left(\sum_{r\ge1}\frac{u^r}{r}\psi_r\cS_H(\mathbf t)\right).}
\end{equation}
\begin{equation}\label{wreath-products-proof-and-verification:eq:stable}
 \boxed{\displaystyle \cW_{H,\infty}
 =\Exp_\sigma(\cS_H-1)
 =\exp\left(\sum_{r\ge1}\frac{\psi_r\cS_H-1}{r}\right).}
\end{equation}

\subsection{Cycle-product determinant lemma}
Write $g=(h_1,\ldots,h_n;\pi)$.  If
$C=(i_1\ i_2\ \cdots\ i_r)$ is a cycle of $\pi$, define
\[
 h_C=h_{i_r}\cdots h_{i_2}h_{i_1}.
\]
On $V_C=V_{i_1}\oplus\cdots\oplus V_{i_r}$, applying $g$ exactly $r$
times returns a vector to its original block and applies $h_C$.  If $\lambda$
is an eigenvalue of $h_C$, the corresponding block-cycle eigenvalues are the
$r$ roots of $z^r=\lambda$.  Their determinant factor is
\[
 \prod_{z^r=\lambda}(1-tz)=1-t^r\lambda.
\]
Multiplying over the eigenvalues of $h_C$ proves
\begin{equation}\label{wreath-products-proof-and-verification:eq:cycle-product}
 \boxed{\det(I-tg\mid V_C)=\det(I-t^rh_C\mid V).}
\end{equation}
The checker tests \eqref{wreath-products-proof-and-verification:eq:cycle-product} without eigenvalues: it constructs
both integer matrices and recovers each determinant polynomial from exact
traces via Newton identities.  Every labeling of cycles of lengths $1,2,3$
is tested for both base groups.

\subsection{Proof of the exact generating function}
Fix $\pi\in S_n$.  On an $r$-cycle, the product map $H^r\to H$ has exactly
$|H|^{r-1}$ preimages for each $h\in H$.  Thus the cycle product is uniform in
$H$, and disjoint cycles have independent products.  By
\eqref{wreath-products-proof-and-verification:eq:cycle-product}, an $r$-cycle contributes the average
\[
 B_r(\mathbf t)=\frac1{|H|}\sum_{h\in H}
 \prod_a\det(I-t_a^rh)^{-1}=\psi_r\cS_H(\mathbf t).
\]
It follows that
\[
 \cW_{H,n}=\frac1{n!}\sum_{\pi\in S_n}
 \prod_{C\in\operatorname{Cyc}(\pi)}B_{|C|}.
\]
Substitution in the symmetric-group cycle-index identity
\[
 \sum_{n\ge0}\frac{u^n}{n!}\sum_{\pi\in S_n}
 \prod_{r\ge1}x_r^{c_r(\pi)}
 =\exp\left(\sum_{r\ge1}\frac{x_ru^r}{r}\right)
\]
proves \eqref{wreath-products-proof-and-verification:eq:exact}.

\subsection{Coefficient recurrence used by the checker}
Write the right side of \eqref{wreath-products-proof-and-verification:eq:exact} as $F(u)=\sum_{n\ge0}F_nu^n$.
Logarithmic differentiation in $u$ gives
\begin{equation}\label{wreath-products-proof-and-verification:eq:finite-recurrence}
 \boxed{nF_n=\sum_{r=1}^n(\psi_r\cS_H)F_{n-r},\qquad F_0=1.}
\end{equation}
For $\tau=(\tau_1,\ldots,\tau_\ell)$ the code represents a truncated
multivariate series by an exact dictionary indexed by
$0\le\alpha_i\le\tau_i$.  The coefficient of $\mathbf t^\alpha$ in
$\psi_r\cS_H$ vanishes unless every $\alpha_i$ is divisible by $r$; in the
divisible case it is the direct base-group average
\[
 \frac1{|H|}\sum_{h\in H}\prod_i
 \chi_{\Sym^{\alpha_i/r}V}(h).
\]
This gives an exact rational implementation of \eqref{wreath-products-proof-and-verification:eq:finite-recurrence};
all denominators cancel in the final invariant dimensions.

The direct side does not use this recurrence.  It constructs every
$nd\times nd$ block-monomial matrix and computes its complete symmetric
characters from
\begin{equation}\label{wreath-products-proof-and-verification:eq:newton-h}
 m h_m(g)=\sum_{k=1}^m\operatorname{tr}(g^k)h_{m-k}(g),\qquad h_0=1.
\end{equation}
It then averages $\prod_i h_{\tau_i}(g)$ over all $|H|^nn!$ matrices.

\subsection{Stable range and an independent stable extractor}
Separate the constant term in \eqref{wreath-products-proof-and-verification:eq:exact}:
\[
 \sum_{n\ge0}\cW_{H,n}u^n
 =\frac1{1-u}\exp\left(\sum_{r\ge1}\frac{u^r}{r}
 (\psi_r\cS_H-1)\right).
\]
\newpage
Every nonconstant monomial in $\psi_r\cS_H-1$ has total $\mathbf t$-degree
at least $r$.  Products preserve the inequality ``$u$-degree at most
$\mathbf t$-degree.''  Therefore a coefficient of total degree $k$ in the
second factor is a polynomial in $u$ of degree at most $k$.  Multiplication
by $(1-u)^{-1}$ makes its $u^n$ coefficient constant once $n\ge k$.  Setting
$u=1$ in the finite polynomial gives \eqref{wreath-products-proof-and-verification:eq:stable}, and hence
\[
 \boxed{\cW_{H,n}\equiv\Exp_\sigma(\cS_H-1)\pmod{\deg>n}.}
\]

The stable checker does not call the finite recurrence at a large proxy
$n$.  If $L=\sum_{r\ge1}(\psi_r\cS_H-1)/r$ and $F=e^L$, the Euler operator
$D=\sum_i t_i\partial_{t_i}$ gives the multivariate recurrence
\begin{equation}\label{wreath-products-proof-and-verification:eq:stable-recurrence}
 |\alpha|F_\alpha=
 \sum_{0<\beta\le\alpha}|\beta|L_\beta F_{\alpha-\beta},
 \qquad
 L_\beta=\sum_{r\mid\gcd(\beta)}\frac1r
 [\mathbf t^{\beta/r}]\cS_H.
\end{equation}
This is a useful general-purpose coefficient extractor for plethystic
exponentials: it is finite, exact, and independent of the outer block count.

\subsection{Marked compound-Poisson interpretation}
Let $\mathcal C(H)$ be the conjugacy classes of $H$ and choose $h_C\in C$.
Put
\[
 F_{r,C}(\mathbf t)=\prod_a\det(I-t_a^rh_C)^{-1}.
\]
Then the logarithm of the stable series is
\[
 \log\cW_{H,\infty}=
 \sum_{r\ge1}\sum_{C\in\mathcal C(H)}
 \frac{|C|}{r|H|}\bigl(F_{r,C}-1\bigr).
\]
This is the probability-generating functional of independent marks
\[
 N_{r,C}\sim\operatorname{Poisson}\left(\frac{|C|}{r|H|}\right).
\]
The precise interpretation is degreewise: for a coefficient of total degree
$k$, only $r\le k$ can contribute.  Although the total intensity summed over
all $r$ diverges, every fixed-degree calculation sees a finite marked process.
This is exactly the local topology in which cycle counts of uniform
permutations converge to independent Poisson variables.

\subsection{Representations tested}
Two exact integral models exercise different dimensions.
\begin{enumerate}
\item $C_2\to GL_1(\mathbb Q)$, with matrices $[1]$ and $[-1]$.  Here a base
coefficient survives precisely when its total degree is even.
\item $C_3\to GL_2(\mathbb Q)$, generated by
\[
 R=\begin{pmatrix}0&-1\\1&-1\end{pmatrix},\qquad R^3=I.
\]
This is the rational form of the two nontrivial complex characters together;
it tests a genuinely two-dimensional block and non-scalar determinant.
\end{enumerate}
The largest enumerations construct all $2^4 4!=384$ matrices of dimension
$4$ and all $3^3 3!=162$ matrices of dimension $6$.

\newpage
\subsection{Exact verification evidence}
For each representation, every partition through total degree $5$ is checked.
The table shows selected rows; ``stable'' is asserted only in the proved range
$|\tau|\le n$.
\begin{center}
\small
\begin{tabular}{@{}llrrrr@{}}
\toprule
Base $H$ & $n,\tau$ & direct & finite formula & stable & result\\
\midrule
$C_2$ sign & $4,(1)$ & 0 & 0 & 0 & pass\\
$C_2$ sign & $4,(2)$ & 1 & 1 & 1 & pass\\
$C_2$ sign & $4,(1,1)$ & 1 & 1 & 1 & pass\\
$C_2$ sign & $4,(4)$ & 2 & 2 & 2 & pass\\
$C_3$ rational & $3,(2)$ & 1 & 1 & 1 & pass\\
$C_3$ rational & $3,(1,1)$ & 2 & 2 & 2 & pass\\
$C_3$ rational & $3,(3)$ & 2 & 2 & 2 & pass\\
$C_3$ rational & $3,(2,1)$ & 2 & 2 & 2 & pass\\
$C_3$ rational & $3,(4)$ & 2 & 2 & --- & pass\\
\bottomrule
\end{tabular}
\end{center}

Across the complete suite there are 38 direct-matrix versus finite-formula
comparisons and 19 direct-matrix versus stable-formula comparisons.  In
addition, 53 cycle-product determinant identities are checked: $2+4+8=14$
labelings for $C_2$ and $3+9+27=39$ for $C_3$; all pass.  The implementation uses integers and
\texttt{Fraction} throughout, so ``pass'' means exact equality.

\paragraph{Reproducibility and limitations.}
The Marimo notebook exposes the base representation, block count, and degree
cutoff, displays an actual block-monomial matrix, and marks every coefficient.
The standalone module runs the fixed suite and exits nonzero on failure.  The
proof is valid for any finite matrix group over characteristic zero.  The
included constructor focuses on two small faithful integral representations
so exhaustive matrix averaging remains immediate; adding a new base group
requires only its exact matrices.

\paragraph{Conclusion.}
The explicit matrix averages agree with the exact cycle-product generating
function and, in the uniform sufficient range $n\ge|\tau|$, with the
independently extracted stable plethystic exponential. This bound is sharp as
a uniform degree bound, although a fixed coefficient may stabilize earlier.  The marked compound-Poisson
description is therefore supported both formally and computationally.
\endgroup


\clearpage
\section{Multiplicative wreath representations}\label{proof:new-dist:multiplicative_wreath}

\resultbox{\textbf{Canonical testing artifacts.}\par\smallskip Exact backend: \path{lambda_distributions/dists/multiplicative_wreath_sigma_mgfs/verification.py}.\par Regression: \path{lambda_distributions/dists/multiplicative_wreath_sigma_mgfs/test_verification.py}.\par Marimo: \path{marimo_notebooks/multiplicative_wreath.py}.}

\subsection*{Scope and outcome}

Let $G_n=H\wr S_n$.  We prove the wreath class-power theorem and use it to
derive exact $\sigma$-MGFs for product actions on $\Omega^n$ and
tensor-induced modules $V^{\otimes n}$.  Their tensor moments are multiset
counts, which normally diverge polynomially; diagonal coset and twisted
regular actions give related polynomial or exponential rare-event laws.  We
also record the graded bosonic and fermionic Fock formulas, fixed-tail
multipartition limits, growing-multipartition obstructions, varying-base
criteria, and the contrasting compound-Poisson law for additive block
modules.  An executable marimo notebook constructs the relevant permutation
and signed-monomial matrix groups in dimensions up to $36$.  Direct orbit and
fixed-space computations agree with the proposed formulas in $13{,}183$
exact checks.

\subsection{Common target and wreath class powers}

Let $H$ be a finite group and $G_n=H^n\rtimes S_n$.  For a complex
$G_n$-module $W$ and a partition $\tau=(\tau_1,\ldots,\tau_s)$, put
\[
 a_\tau(G_n,W)=\dim\left(\bigotimes_{j=1}^s\Sym^{\tau_j}W\right)^{G_n},
 \qquad
 \MGF_{G_n,W}(\mathbf t)=\sum_\tau a_\tau m_\tau(\mathbf t).
\]
Reynolds averaging and
$\sum_{d\ge0}\Tr(\Sym^d A)t^d=\det(I-tA)^{-1}$ give
\begin{equation}\label{nd:multiplicative_wreath:eq:universal-molien}
 \MGF_{G_n,W}
 =\frac1{|G_n|}\sum_{g\in G_n}
 \exp\!\left(\sum_{r\ge1}\frac{p_r(\mathbf t)}r\chi_W(g^r)\right).
\end{equation}

Let $\mathcal C(H)$ be the conjugacy classes of $H$.  A wreath type is a
partition-valued function $\boldsymbol\mu:C\mapsto\mu(C)$ of total size $n$.
A part $\ell$ of $\mu(C)$ denotes an $\ell$-cycle of the top permutation whose
ordered label product lies in $C$.  For $c\in C$, set $z_C=|C_H(c)|$, and set
\[
 z_\lambda=\prod_{j\ge1}j^{m_j(\lambda)}m_j(\lambda)!,\qquad
 Z_{\boldsymbol\mu}=\prod_{C\in\mathcal C(H)}
 z_{\mu(C)}z_C^{\ell(\mu(C))}.
\]

\begin{theorem}[wreath class-power theorem]\label{nd:multiplicative_wreath:thm:class-power}
The proportion of elements of type $\boldsymbol\mu$ is
$Z_{\boldsymbol\mu}^{-1}$.  If an $\ell$-cycle has cycle product $c$ and
$d=\gcd(\ell,r)$, then the $r$th power has $d$ cycles of length $\ell/d$,
each with cycle product conjugate to $c^{r/d}$.  Denote the transformed type
by $\boldsymbol\mu^{[r]}$.  Then
\begin{equation}\label{nd:multiplicative_wreath:eq:universal-class}
 \boxed{\MGF_{G_n,W}=
 \sum_{\boldsymbol\mu\vdash_H n}\frac1{Z_{\boldsymbol\mu}}
 \exp\!\left[\sum_{r\ge1}\frac{p_r}{r}
 \chi_W(\boldsymbol\mu^{[r]})\right].}
\end{equation}
\end{theorem}

\begin{proof}
Fix the top cycle structure and one cycle of length $\ell$.  After choosing
$\ell-1$ labels freely, the last label has exactly $|C|$ choices that place
the cycle product in $C$.  Combining this with the ordinary permutation
cycle-index count gives the denominator $z_{\mu(C)}z_C^{\ell(\mu(C))}$;
multiplication over $C$ gives $Z_{\boldsymbol\mu}$.  Advancing by $r$ positions
on an $\ell$-cycle has $d=\gcd(\ell,r)$ orbits, each of length $\ell/d$.
Traversing one new orbit winds $r/d$ times around the old cycle, so its label
product is conjugate to $c^{r/d}$.  Substitution in
\eqref{nd:multiplicative_wreath:eq:universal-molien} proves \eqref{nd:multiplicative_wreath:eq:universal-class}.
\end{proof}

\begin{remark}[orientation]
Changing the starting point of a top cycle conjugates its label product.
Thus all formulas are independent of left-versus-right cycle notation once
the product is recorded only up to conjugacy.
\end{remark}

\subsection{Product-action permutation groups}

Suppose $H$ acts on a finite set $\Omega$ and $G_n$ acts on
$X_n=\Omega^n$.  Write $g=(h_1,\ldots,h_n;\pi)$.  For a cycle $Q$ of $\pi$,
let $a_Q$ be its ordered label product and put
$d_Q(r)=\gcd(|Q|,r)$ and $f_s(a)=|\Fix_\Omega(a^s)|$.

\begin{theorem}[product-action fixed powers and $\sigma$-MGF]
For every $r\ge1$,
\begin{equation}\label{nd:multiplicative_wreath:eq:product-fixed}
 \boxed{F_r(g):=|\Fix_{\Omega^n}(g^r)|
 =\prod_{Q\in\Cyc(\pi)}
 f_{r/d_Q(r)}(a_Q)^{d_Q(r)}.}
\end{equation}
If $c_j(g)$ is the number of $j$-cycles on $X_n$, then
\begin{equation}\label{nd:multiplicative_wreath:eq:mobius}
 c_j(g)=\frac1j\sum_{r\mid j}\mu(j/r)F_r(g),
\end{equation}
and therefore
\begin{equation}\label{nd:multiplicative_wreath:eq:product-smg}
 \boxed{\MGF^{\rm prod}_n
 =\sum_{\boldsymbol\mu\vdash_H n}\frac1{Z_{\boldsymbol\mu}}
 \prod_{i,j\ge1}(1-t_i^j)^{-c_j(\boldsymbol\mu)}.}
\end{equation}
\end{theorem}

\begin{proof}
The $r$th power splits $Q$ into $d_Q(r)$ coordinate cycles.  Going around one
new cycle applies an element conjugate to $a_Q^{r/d_Q(r)}$.  Each new cycle
may independently receive any point fixed by that element, proving
\eqref{nd:multiplicative_wreath:eq:product-fixed}.  Since
$F_r(g)=\sum_{j\mid r}j c_j(g)$, M\"obius inversion gives
\eqref{nd:multiplicative_wreath:eq:mobius}.  The represented permutation matrix has determinant
$\prod_j(1-t^j)^{c_j(g)}$; Reynolds averaging and
Theorem~\ref{nd:multiplicative_wreath:thm:class-power} give \eqref{nd:multiplicative_wreath:eq:product-smg}.
\end{proof}

Define $R_s(H,\Omega)=\#(H\backslash\Omega^s)$.

\begin{theorem}[exact product moments]\label{nd:multiplicative_wreath:thm:product-moment}
\begin{equation}\label{nd:multiplicative_wreath:eq:product-moment}
 \boxed{[m_{(1^s)}]\MGF^{\rm prod}_n
 =\binom{n+R_s(H,\Omega)-1}{R_s(H,\Omega)-1}.}
\end{equation}
\end{theorem}

\begin{proof}
An ordered $s$-tuple of words is an $s\times n$ array.  The base group $H^n$
replaces each column by its $H$-orbit in $\Omega^s$, leaving $R_s$ column
types.  The top group forgets their order.  The orbits are therefore
multisets of size $n$ from $R_s$ types, counted by \eqref{nd:multiplicative_wreath:eq:product-moment}.
The orbit count equals the invariant dimension in the permutation module.
\end{proof}

If the base action is nontrivial and transitive, then $R_2\ge2$, so the pair
coefficient grows as $n^{R_2-1}/(R_2-1)!$ and no finite coefficientwise limit
exists.  If $\delta_H$ is the derangement proportion of the base action, a
fixed word exists exactly when every top-cycle product is a nonderangement.
Since the cycle products are independent uniform elements of $H$ conditional
on $\pi$,
\begin{equation}\label{nd:multiplicative_wreath:eq:product-rare}
 \Pr(F_1>0)=\E(1-\delta_H)^{K(\pi)}
 =\frac{\Gamma(n+1-\delta_H)}
 {\Gamma(1-\delta_H)\Gamma(n+1)}
 \sim\frac{n^{-\delta_H}}{\Gamma(1-\delta_H)}.
\end{equation}
Thus $F_1\to0$ in probability while transitivity gives $\E F_1=1$ and higher
moments grow polynomially.

\subsection{Tensor-induced signed matrix groups}

Let $V$ be an $H$-module with character $\chi$, and let
$T_n(V)=V^{\otimes n}$ with the usual label action and permutation of tensor
factors.

\begin{theorem}[tensor cycle character and moments]\label{nd:multiplicative_wreath:thm:tensor}
For $g=(h_1,\ldots,h_n;\pi)$,
\begin{align}
 \chi_{T_n(V)}(g)&=\prod_{Q\in\Cyc(\pi)}\chi(a_Q),\label{nd:multiplicative_wreath:eq:tensor-char}\\
 \chi_{T_n(V)}(g^r)&=\prod_Q
 \chi(a_Q^{r/d_Q(r)})^{d_Q(r)}.\label{nd:multiplicative_wreath:eq:tensor-power}
\end{align}
Consequently,
\begin{equation}\label{nd:multiplicative_wreath:eq:tensor-smg}
 \boxed{\MGF_n^\otimes=
 \sum_{\boldsymbol\mu\vdash_Hn}\frac1{Z_{\boldsymbol\mu}}
 \exp\!\left[\sum_{r\ge1}\frac{p_r}{r}
 \prod_{C,\ell}\chi(c^{r/d_{\ell,r}})^{
 d_{\ell,r}m_\ell(\mu(C))}\right].}
\end{equation}
If $a_s(H,V)=\dim(V^{\otimes s})^H$, then
\begin{equation}\label{nd:multiplicative_wreath:eq:tensor-moment}
 \boxed{[m_{(1^s)}]\MGF_n^\otimes=
 \begin{cases}
 \binom{n+a_s-1}{a_s-1},&a_s>0,\\
 0,&a_s=0.
 \end{cases}}
\end{equation}
\end{theorem}

\begin{proof}
Contracting tensor indices around one top cycle produces the trace of its
label product, proving \eqref{nd:multiplicative_wreath:eq:tensor-char}; Theorem~\ref{nd:multiplicative_wreath:thm:class-power}
gives \eqref{nd:multiplicative_wreath:eq:tensor-power}.  Substitution in
\eqref{nd:multiplicative_wreath:eq:universal-class} proves \eqref{nd:multiplicative_wreath:eq:tensor-smg}.  For the moment,
rearrange tensor factors:
\[
 (V^{\otimes n})^{\otimes s}\cong(V^{\otimes s})^{\otimes n}.
\]
Taking $H^n$-invariants gives $((V^{\otimes s})^H)^{\otimes n}$; taking the
remaining $S_n$-invariants gives
$\Sym^n((V^{\otimes s})^H)$.  Its dimension is
\eqref{nd:multiplicative_wreath:eq:tensor-moment}.
\end{proof}

If $a_s\ge2$, the coefficient is asymptotic to
$n^{a_s-1}/(a_s-1)!$.  If $V$ is self-dual irreducible and $\dim V>1$, then
$a_2=1$ and
\[
 a_4=\dim\operatorname{End}_H(V\otimes V)\ge2,
\]
because $V\otimes V=\Sym^2V\oplus\Lambda^2V$ has two nonzero invariant
summands.  Hence the fourth tensor moment is at least $n+1$.  A
one-dimensional character $\theta$ of order $q$ is exceptional:
$a_s=1$ when $q\mid s$ and $a_s=0$ otherwise, so these obstruction
coefficients remain bounded.

\subsection{Bosonic and fermionic Fock modules}

For a finite $G$-module $W$, introduce the energy-graded modules
\[
 \mathcal F^+(W)=\bigoplus_{d\ge0}q^d\Sym^dW,
 \qquad
 \mathcal F^-(W)=\bigoplus_{d\ge0}q^d\Lambda^dW.
\]
The graded $\sigma$-MGF applies the outer symmetric-power construction to
these graded modules.

\begin{proposition}[graded Fock formulas]\label{nd:multiplicative_wreath:prop:fock}
\begin{align}
 \MGF^+_{G,W}(q;\mathbf t)
 &=\frac1{|G|}\sum_{g\in G}\exp\!\left[
 \sum_{r\ge1}\frac{p_r}{r}\det(I-q^rg^r\mid W)^{-1}\right],
 \label{nd:multiplicative_wreath:eq:fock-plus}\\
 \MGF^-_{G,W}(q;\mathbf t)
 &=\frac1{|G|}\sum_{g\in G}\exp\!\left[
 \sum_{r\ge1}\frac{p_r}{r}\det(I+q^rg^r\mid W)\right].
 \label{nd:multiplicative_wreath:eq:fock-minus}
\end{align}
Removing the vacuum replaces the determinant inverse in
\eqref{nd:multiplicative_wreath:eq:fock-plus} by that inverse minus $1$.
\end{proposition}

\begin{proof}
For each $r$, the graded characters are
\[
 \sum_{d\ge0}q^{dr}\Tr(g^r\mid\Sym^dW)
 =\det(I-q^rg^r\mid W)^{-1},
\]
and
\[
 \sum_{d\ge0}q^{dr}\Tr(g^r\mid\Lambda^dW)
 =\det(I+q^rg^r\mid W).
\]
Insert these identities into the plethystic form of Reynolds averaging.
\end{proof}

For the block module $W_n=V^{\oplus n}$, every total-energy-$E$ monomial has
support on at most $E$ blocks.  Therefore
\begin{equation}\label{nd:multiplicative_wreath:eq:fock-stability}
 [q^Em_\tau]\MGF^\pm_{H\wr S_n,W_n}
 \quad\text{is independent of $n$ for $n\ge E$.}
\end{equation}
The grading is essential: at $q=1$ the bosonic Fock space is
infinite-dimensional and ungraded coefficients are generally infinite.

\subsection{Irreducible multipartitions}

Irreducibles of $H\wr S_n$ are indexed by multipartitions
$\boldsymbol\lambda=(\lambda^\rho)_{\rho\in\Irr(H)}$.  The wreath Frobenius
characteristic sends their characters to
\begin{equation}\label{nd:multiplicative_wreath:eq:wreath-frob}
 \chi^{\boldsymbol\lambda}\longleftrightarrow
 \prod_{\rho\in\Irr(H)}s_{\lambda^\rho}(x^{(\rho)}).
\end{equation}
Theorem~\ref{nd:multiplicative_wreath:thm:class-power} immediately gives
\begin{equation}\label{nd:multiplicative_wreath:eq:irrep-smg}
 \boxed{\MGF_{\boldsymbol\lambda}
 =\sum_{\boldsymbol\mu\vdash_Hn}\frac1{Z_{\boldsymbol\mu}}
 \exp\!\left[\sum_{r\ge1}\frac{p_r}{r}
 \chi^{\boldsymbol\lambda}(\boldsymbol\mu^{[r]})\right].}
\end{equation}

For a fixed-tail family, only the trivial-color component acquires a long
first row and all other boxes form a fixed multipartition
$\boldsymbol\alpha$. Assume a wreath character-polynomial theorem supplies
a polynomial $Q_{\boldsymbol\alpha}(X_{\ell,C})$ and an explicit threshold
$N_{\boldsymbol\alpha}(D)$ valid for all power evaluations needed through
total outer degree $D$. This is the additional input not proved in this
section.  Under the
uniform measure,
\begin{equation}\label{nd:multiplicative_wreath:eq:marked-poisson}
 X_{\ell,C}\Longrightarrow Z_{\ell,C},\qquad
 Z_{\ell,C}\sim\operatorname{Poisson}
 \left(\frac1{\ell|C_H(c)|}\right),
\end{equation}
independently.  Define
\begin{equation}\label{nd:multiplicative_wreath:eq:power-poisson}
 Z_{j,D}^{[r]}=
 \sum_{\substack{\ell,C:\ \ell/\gcd(\ell,r)=j\\
 c^{r/\gcd(\ell,r)}\in D}}
 \gcd(\ell,r)Z_{\ell,C}.
\end{equation}
For each fixed degree $D$ and $n\ge N_{\boldsymbol\alpha}(D)$, the
character-polynomial substitution is valid; marked-cycle moment convergence
then yields the conditional marked-Poisson-chaos law
\begin{equation}\label{nd:multiplicative_wreath:eq:fixed-tail-limit}
 \boxed{\MGF_{\boldsymbol\alpha,\infty}
 =\E\exp\!\left[\sum_{r\ge1}\frac{p_r}{r}
 Q_{\boldsymbol\alpha}(\boldsymbol Z^{[r]})\right].}
\end{equation}

\proofstatus{Conditional}{Equation~\eqref{nd:multiplicative_wreath:eq:fixed-tail-limit}
is proved once the stated wreath character-polynomial theorem and threshold
$N_{\boldsymbol\alpha}(D)$ are supplied. A proof obligation is collected in
the final provisional part.}

There is no corresponding universal raw limit for growing multipartitions.
If a self-dual $W_{\boldsymbol\lambda}$ satisfies
\[
 W_{\boldsymbol\lambda}\otimes W_{\boldsymbol\lambda}
 =\bigoplus_{\boldsymbol\nu}
 g^H_{\boldsymbol\lambda,\boldsymbol\lambda,\boldsymbol\nu}
 W_{\boldsymbol\nu},
\]
then
\begin{equation}\label{nd:multiplicative_wreath:eq:kronecker-obstruction}
 [m_{(1,1,1,1)}]\MGF
 =\dim\operatorname{End}_{H\wr S_n}(W_{\boldsymbol\lambda}^{\otimes2})
 =\sum_{\boldsymbol\nu}
 (g^H_{\boldsymbol\lambda,\boldsymbol\lambda,\boldsymbol\nu})^2.
\end{equation}
Growth of this sum obstructs convergence.

\subsection{Diagonal coset matrix groups}

Let $N_n=H^n$, $\Delta H=\{(h,\ldots,h):h\in H\}$, and let $N_n$ act on
$X_n=N_n/\Delta H$.  For a class $C$ and $c\in C$, write
$z_C=|C_H(c)|$.

\begin{theorem}[diagonal fixed powers and moments]\label{nd:multiplicative_wreath:thm:diagonal}
For $g=(g_1,\ldots,g_n)$,
\begin{equation}\label{nd:multiplicative_wreath:eq:diagonal-fixed}
 F_r(g)=\begin{cases}
 z_C^{n-1},&g_1^r,\ldots,g_n^r\in C,\\
 0,&\text{otherwise}.
 \end{cases}
\end{equation}
If $X(g)=F_1(g)$, then
\begin{equation}\label{nd:multiplicative_wreath:eq:diagonal-moment}
 \boxed{\E X(g)^s=\sum_{C\in\mathcal C(H)}z_C^{(s-1)n-s}.}
\end{equation}
In particular, $[m_{(1,1)}]=\sum_Cz_C^{n-2}\ge|H|^{n-2}$.
\end{theorem}

\begin{proof}
A coset $x\Delta H$ is fixed by $g^r$ exactly when the components of
$x^{-1}g^rx$ are equal.  If the component powers do not lie in one class,
there is no solution.  If they lie in $C$, choose the common conjugate and
then choose a conjugator in each coordinate; division by $|\Delta H|$ leaves
$z_C^{n-1}$ cosets.  The event $g_1,\ldots,g_n\in C$ has probability
$z_C^{-n}$, and the fixed-point value there is $z_C^{n-1}$.  Summing the
$s$th powers proves \eqref{nd:multiplicative_wreath:eq:diagonal-moment}.
\end{proof}

The identity class alone gives exponential second-moment growth.  Meanwhile
$\Pr(X>0)=\sum_Cz_C^{-n}\to0$ for nontrivial $H$, although $\E X=1$.

\subsection{Twisted-wreath regular matrix groups}

Let $N_n=H^n$, let $P_n\le\operatorname{Aut}(N_n)$, and let
$G_n=N_n\rtimes P_n$ act affinely on $N_n$ by
$x^{(a,p)}=a\,p(x)$.  Put
\[
 a_r=a\,p(a)\cdots p^{r-1}(a),
 \qquad L_{p^r}(x)=x\,p^r(x)^{-1}.
\]

\begin{theorem}[Lang-map criterion and moments]\label{nd:multiplicative_wreath:thm:twisted}
\begin{equation}\label{nd:multiplicative_wreath:eq:lang}
 F_r(a,p)=\begin{cases}
 |\Fix_{N_n}(p^r)|,&a_r\in\im L_{p^r},\\
 0,&\text{otherwise}.
 \end{cases}
\end{equation}
Conditioned on $p$, $F_1$ equals $|\Fix_{N_n}(p)|$ with probability its
reciprocal, and otherwise equals zero.  Hence
\begin{equation}\label{nd:multiplicative_wreath:eq:twisted-moment}
 \boxed{[m_{(1^s)}]=\frac1{|P_n|}\sum_{p\in P_n}
 |\Fix_{N_n}(p)|^{s-1}=\#(P_n\backslash N_n^{s-1}).}
\end{equation}
If $P_n=S_n$ only permutes coordinates, then
\begin{equation}\label{nd:multiplicative_wreath:eq:pure-coordinate}
 \boxed{[m_{(1^s)}]=
 \binom{n+|H|^{s-1}-1}{|H|^{s-1}-1}.}
\end{equation}
\end{theorem}

\begin{proof}
The fixed equation $a_rp^r(x)=x$ is equivalent to
$a_r=L_{p^r}(x)$.  Every nonempty fiber of the Lang map is a coset of
$\Fix_{N_n}(p^r)$, proving \eqref{nd:multiplicative_wreath:eq:lang} and the conditional spike law.
Averaging its $s$th power gives the middle expression in
\eqref{nd:multiplicative_wreath:eq:twisted-moment}; Burnside's lemma gives the orbit count.  Under
coordinate permutations, $N_n^{s-1}$ is an array of $n$ elements from
$H^{s-1}$ modulo column permutation, hence a multiset of size $n$ from
$|H|^{s-1}$ types.  This proves \eqref{nd:multiplicative_wreath:eq:pure-coordinate}.
\end{proof}

For coordinate permutations,
\[
 \Pr(F_1>0)=\frac{(1/|H|)_n}{n!}
 \sim\frac{n^{1/|H|-1}}{\Gamma(1/|H|)},
\]
again giving vanishing fixed points in probability with divergent moments.

\subsection{Varying base groups and the additive contrast}

Let $H_N$ act on $\Omega_N$ and $V_N$, and put
$G_{N,k}=H_N\wr S_k$.  Define
\[
 R_{N,s}=\#(H_N\backslash\Omega_N^s),
 \qquad a_{N,s}=\dim(V_N^{\otimes s})^{H_N}.
\]
Theorems~\ref{nd:multiplicative_wreath:thm:product-moment} and \ref{nd:multiplicative_wreath:thm:tensor} give the exact
two-parameter laws
\begin{align}
 [m_{(1^s)}]\MGF_{G_{N,k},\Omega_N^k}
 &=\binom{k+R_{N,s}-1}{R_{N,s}-1},\label{nd:multiplicative_wreath:eq:vary-product}\\
 [m_{(1^s)}]\MGF_{G_{N,k},V_N^{\otimes k}}
 &=\begin{cases}
 \binom{k+a_{N,s}-1}{a_{N,s}-1},&a_{N,s}>0,\\
 0,&a_{N,s}=0.
 \end{cases}\label{nd:multiplicative_wreath:eq:vary-tensor}
\end{align}
These binomial coefficients govern every joint-growth path.

The direct-sum block module behaves differently.  If
$A_N(\mathbf t)=\MGF_{H_N,V_N}(\mathbf t)$, the labelled-cycle exponential
formula gives
\begin{equation}\label{nd:multiplicative_wreath:eq:block-series}
 \sum_{k\ge0}u^k\MGF_{H_N\wr S_k,V_N^{\oplus k}}
 =\exp\!\left[\sum_{\ell\ge1}\frac{u^\ell}{\ell}
 A_N(t_1^\ell,t_2^\ell,\ldots)\right].
\end{equation}
Separating its constant term $(1-u)^{-1}$ shows coefficientwise that
\begin{equation}\label{nd:multiplicative_wreath:eq:block-limit}
 \lim_{k\to\infty}\MGF_{H_N\wr S_k,V_N^{\oplus k}}
 =\exp\!\left[\sum_{\ell\ge1}\frac{
 A_N(t_1^\ell,t_2^\ell,\ldots)-1}{\ell}\right].
\end{equation}
Thus the additive block representation has a compound-Poisson limit, while
\eqref{nd:multiplicative_wreath:eq:vary-product} and \eqref{nd:multiplicative_wreath:eq:vary-tensor} usually diverge.

\begin{corollary}[multiplicative wreath obstruction theorem]
For product and tensor-induced wreath representations, the coefficient
$[m_{(1^s)}]$ is a multiset count
$\binom{n+r_s-1}{r_s-1}$.  If $r_s\ge2$, it grows as
$n^{r_s-1}/(r_s-1)!$, so the raw $\sigma$-MGF has no finite coefficientwise
limit.
\end{corollary}

\subsection{Exact matrix verification}

\subsubsection{Independent computational routes}

The companion marimo notebook and \texttt{verification.py} use the following
audit protocol.
\begin{enumerate}
\item Construct every group element with the semidirect-product law.  Store a
permutation matrix by its unique nonzero row in each column; store a signed
monomial matrix by that row and a sign.  This is the complete matrix, not a
character-only surrogate.
\item Form matrix powers directly and compare their fixed points or traces
with the structural cycle-product, class-power, diagonal-centralizer, or Lang
map formula.
\item Compute $\bigotimes_j\Sym^{\tau_j}W$ on its actual monomial basis.
Permutation modules use literal configuration orbits.  Signed modules
propagate signs around each basis orbit: an orbit contributes one invariant
unless a stabilizer sends a basis vector to its negative.
\item Independently compute the same coefficient by the full Reynolds--Molien
average.  Moment cases receive a third comparison with the closed binomial or
centralizer formula.
\end{enumerate}
All checks use integers or exact rational numbers.

\subsubsection{Wreath types and product actions}

\begin{center}
\begin{tabular}{lrrrr}
\toprule
group & order & wreath types & power checks & weight checks\\
\midrule
$C_2\wr S_3$ & 48 & 10 & 288 & 10\\
$S_3\wr S_2$ & 72 & 9 & 432 & 9\\
\bottomrule
\end{tabular}
\end{center}

For $S_3\curvearrowright\{0,1,2\}$, direct product-action matrices give:
\begin{center}
\begin{tabular}{rrrrrr}
\toprule
$n$ & matrix dimension & $s$ & $R_s$ & direct/Burnside & formula\\
\midrule
2 & 9  & 2 & 2 & 3  & 3\\
2 & 9  & 3 & 5 & 15 & 15\\
3 & 27 & 2 & 2 & 4  & 4\\
3 & 27 & 3 & 5 & 35 & 35\\
\bottomrule
\end{tabular}
\end{center}
The program additionally checks all first six powers of all $72+1296$
product-action matrices, for $8{,}208$ fixed-power comparisons.

\subsubsection{Tensor-induced and fixed-tail signed groups}

Take $H=C_2$ and $V=\mathbf1\oplus\mathrm{sgn}$.  Then
$a_s=2^{s-1}$.  The tensor-induced results are:
\begin{center}
\begin{tabular}{lrrrrrr}
\toprule
group & dim. & $a_{(1)}$ & $a_{(2)}$ & $a_{(1,1)}$ & $a_{(2,1)}$ & $a_{(1^4)}$\\
\midrule
$C_2\wr S_2$ & 4 & 1 & 3 & 3 & 7  & 36\\
$C_2\wr S_3$ & 8 & 1 & 4 & 4 & 13 & 120\\
\bottomrule
\end{tabular}
\end{center}
Every entry is simultaneously a direct signed-fixed-space count and a Molien
coefficient; the moment entries also equal \eqref{nd:multiplicative_wreath:eq:tensor-moment}.  Traces
of six powers give another $336$ checks.

For the fixed-tail bipartition $((n-1),(1))$ of $C_2\wr S_n$, the natural
signed module has character polynomial $X_{1,+}-X_{1,-}$.  For $n=2,3,4$,
both direct and class-formula routes give
\begin{center}
\begin{tabular}{c|rrrrrr}
\toprule
$n$ & $(1)$ & $(2)$ & $(1,1)$ & $(3)$ & $(2,1)$ & $(1^4)$\\
\midrule
2 & 0&1&1&0&0&4\\
3 & 0&1&1&0&0&4\\
4 & 0&1&1&0&0&4\\
\bottomrule
\end{tabular}
\end{center}
The matrix trace also agrees with the power-transformed character polynomial
in $1{,}760$ element-power comparisons.  Its limiting character is the
Skellam variable $Z_+-Z_-$ with independent
$Z_\pm\sim\operatorname{Poisson}(1/2)$, concretely illustrating
\eqref{nd:multiplicative_wreath:eq:fixed-tail-limit}.

\subsubsection{Diagonal, twisted, Fock, and varying-base checks}

For the diagonal coset actions, matrix averages and
\eqref{nd:multiplicative_wreath:eq:diagonal-moment} agree:
\begin{center}
\begin{tabular}{lrrrr}
\toprule
action & dimension & $\E X$ & $\E X^2$ & $\E X^3$\\
\midrule
$S_3^2\curvearrowright S_3^2/\Delta S_3$ & 6  & 1 & 3  & 11\\
$S_3^3\curvearrowright S_3^3/\Delta S_3$ & 36 & 1 & 11 & 251\\
\bottomrule
\end{tabular}
\end{center}
All first six powers of all $252$ matrices give $1{,}512$ further checks.
For $S_3^2\rtimes S_2\curvearrowright S_3^2$, the direct/formula moments are
$1,21,666$ for $s=1,2,3$; all $432$ fixed-power checks satisfy the Lang-map
criterion, and the direct pair-orbit count is $21$.

For $C_2$ on $\mathbf1\oplus\mathrm{sgn}$, explicit bosonic monomials give
invariant dimensions $1,1,2,2,3,3,4$ in energies $0$ through $6$.
Explicit exterior bases give $1,1,0,0,0,0,0$.  These values and all
unaveraged power characters agree with \eqref{nd:multiplicative_wreath:eq:fock-plus} and
\eqref{nd:multiplicative_wreath:eq:fock-minus} in $94$ checks.  Regular bases $C_2,C_3$, with
$k=2,3$ and $s=2,3$, give eight further exact specializations of
\eqref{nd:multiplicative_wreath:eq:vary-product}--\eqref{nd:multiplicative_wreath:eq:vary-tensor}.

\resultbox{\textbf{Verification result.}  The executable suite reports
\textbf{PASS: 13,183 exact checks}.  This includes group construction,
semidirect powers, class weights, fixed points, traces, orbit counts, signed
fixed spaces, Molien coefficients, moment theorems, and graded Fock
characters.  The asymptotic claims are deductions from the proved exact
formulas; they are not extrapolations from the small cases.}

\subsection{Remarks and asymptotic classification}

\begin{center}
\small
\begin{tabularx}{\linewidth}{>{\raggedright\arraybackslash}p{0.25\linewidth}
>{\raggedright\arraybackslash}p{0.29\linewidth}X}
\toprule
family & controlling statistic & behavior\\
\midrule
$V^{\oplus n}$ & cyclewise determinants & compound-Poisson limit\\
$\Omega^n$ & $R_s=\#(H\backslash\Omega^s)$ & polynomial divergence if $R_s\ge2$\\
$V^{\otimes n}$ & $a_s=\dim(V^{\otimes s})^H$ & polynomial divergence if $a_s\ge2$\\
bosonic/fermionic Fock & $\det(I\mp q^rg^r)^{\mp1}$ & fixed-energy stability\\
fixed-tail multipartition & marked character polynomial & marked-Poisson chaos\\
growing multipartition & wreath Kronecker multiplicities & no universal raw limit\\
$H^n/\Delta H$ & common conjugacy class & exponential moment divergence\\
$H^n\rtimes P_n\curvearrowright H^n$ & $|\Fix_{H^n}(p)|$ & Bernoulli-spike rare events\\
varying $H_N\wr S_k$ & $R_{N,s}$ or $a_{N,s}$ & exact two-parameter criterion\\
\bottomrule
\end{tabularx}
\end{center}

The computational examples are deliberately modest: their role is to test
the representation, conventions, power operation, and coefficient extraction
without relying on asymptotic numerics.  The formulas themselves are valid at
every finite size covered by their hypotheses.  The small groups already
display the structural distinction: cyclewise sums exponentiate into stable
laws; cyclewise products create rare, large contributions that force moments
to grow.

\subsection{Reproducibility and scope}
\thispagestyle{fancy}

The standalone source, compiled PDF, marimo notebook, exact backend, and
focused tests live together in
\texttt{new\_dists/multiplicative\_wreath\_sigma\_mgfs}.  From the repository
root, the complete audit is reproduced with
\begin{center}
\texttt{.venv/bin/python -m
new\_dists.multiplicative\_wreath\_sigma\_mgfs.verification},
\end{center}
and the interactive laboratory is opened with
\begin{center}
\texttt{.venv/bin/marimo run
marimo\_notebooks/multiplicative\_wreath.py}.
\end{center}
The PDF is compiled directly from its colocated \LaTeX{} source.  No combined
anthology source or PDF is modified by this package.

The exhaustive computations certify the listed finite instances and test the
implementation of each structural reduction.  Universality comes from the
proofs: Reynolds averaging, wreath class powers, orbit--multiset bijections,
and the Lang-map fiber argument.  Statements about growing multipartitions
remain criteria rather than unconditional convergence theorems, because the
relevant wreath Kronecker multiplicities depend on the chosen family.

\begin{thebibliography}{9}
\bibitem{nd:multiplicative_wreath:wreath-symmetric}
I.~G.~Macdonald-style wreath characteristic background is developed in
\emph{Wreath Product Symmetric Functions},
\href{https://arxiv.org/abs/0809.2439}{arXiv:0809.2439}.

\bibitem{nd:multiplicative_wreath:fock}
I.~B.~Frenkel, N.~Jing, and W.~Wang,
\emph{Vertex representations via finite groups and the McKay correspondence},
\href{https://arxiv.org/abs/math/9907166}{arXiv:math/9907166}.

\bibitem{nd:multiplicative_wreath:fig}
K.~Casto,
\emph{$\mathrm{FI}_G$-modules, orbit configuration spaces, and complex
reflection groups},
\href{https://arxiv.org/abs/1608.06317}{arXiv:1608.06317}.
\end{thebibliography}

\clearpage
\section{Generalized-dihedral master theorem}\label{proof:generalized-dihedral}
\begingroup
\definecolor{ink}{HTML}{17212B}
\definecolor{accent}{HTML}{2457A7}
\definecolor{soft}{HTML}{EEF4FC}
\providecommand{\C}{\mathbb C}
\renewcommand{\C}{\mathbb C}
\providecommand{\Z}{\mathbb Z}
\renewcommand{\Z}{\mathbb Z}
\providecommand{\Sym}{\operatorname{Sym}}
\renewcommand{\Sym}{\operatorname{Sym}}
\providecommand{\Dih}{\operatorname{Dih}}
\renewcommand{\Dih}{\operatorname{Dih}}
\providecommand{\Exp}{\operatorname{Exp}}
\renewcommand{\Exp}{\operatorname{Exp}}
\providecommand{\tr}{\operatorname{tr}}
\renewcommand{\tr}{\operatorname{tr}}
\providecommand{\diag}{\operatorname{diag}}
\renewcommand{\diag}{\operatorname{diag}}
\providecommand{\Mcal}{\SMG}
\renewcommand{\Mcal}{\SMG}
\providecommand{\Rcal}{\mathcal R}
\renewcommand{\Rcal}{\mathcal R}
\providecommand{\Scal}{\SMG}
\renewcommand{\Scal}{\SMG}
\setcounter{theorem}{0}
\setcounter{proposition}{0}
\setcounter{lemma}{0}
\setcounter{corollary}{0}
\setcounter{remark}{0}
\begin{abstract}
Let $G=\Dih(A)=A\rtimes C_2$ for a finite abelian group $A$, where the
nontrivial element acts on $A$ by inversion.  This note classifies the complex
irreducible representations of $G$ and proves a refined Molien formula for an
arbitrary finite-dimensional representation.  The formula is then checked by
an independent full-matrix computation: the accompanying marimo notebook
constructs every represented group element, constructs the action on tensor
products of symmetric powers, and compares the rank of the Reynolds projector
with coefficient extraction from the claimed formula.  All 33 representative
checks passed, covering cyclic and noncyclic $A$, odd and even orders,
one- and two-dimensional irreducible blocks, repeated blocks, and
representation dimensions from $2$ through $5$.
\end{abstract}

\subsection{Setup and imported target}
Let $A$ be finite abelian and
\[
G=\Dih(A)=A\rtimes\langle s\rangle,\qquad s^2=1,\qquad sas=a^{-1}.
\]
Thus $G=A\sqcup As$ and the usual dihedral group is the case $A=C_n$. For a
specified $\rho:G\to GL(V)$, put
\begin{equation}\label{generalized-dihedral-master-theorem:eq:molien-identity}
\Mcal_{G,V}:=\smgv{G}{V}
=\frac1{|G|}\sum_{g\in G}\prod_{r\ge1}\det(I-t_r\rho(g))^{-1}.
\end{equation}
By Tool~\ref{tool:coefficient},
\begin{equation}\label{generalized-dihedral-master-theorem:eq:invariant-identity}
\Mcal_{G,V}=\sum_\tau\dim(W_\tau^G)m_\tau.
\end{equation}
These are imported. The new work begins with the irreducible decomposition and
the two coset spectra.

\subsection{Classification of the irreducibles}

Write $\widehat A=\operatorname{Hom}(A,\C^\times)$.  Conjugation by $s$ acts
on $\widehat A$ by $\chi\mapsto\chi^{-1}$.

\begin{theorem}[Irreducibles of $\Dih(A)$]\label{generalized-dihedral-master-theorem:thm:classification}
The irreducible complex representations of $G=\Dih(A)$ are as follows.
\begin{enumerate}[label=\textup{(\roman*)}]
\item If $\eta\in\widehat A$ satisfies $\eta^2=1$, it has exactly two
one-dimensional extensions
\[
  L_{\eta,\epsilon}(a)=\eta(a),\qquad
  L_{\eta,\epsilon}(s)=\epsilon,\qquad \epsilon\in\{+1,-1\}.
\]
\item If $\chi\ne\chi^{-1}$, then
$\rho_\chi=\operatorname{Ind}_A^G\chi$ is irreducible of dimension two.  In a
suitable basis,
\[
  \rho_\chi(a)=
  \begin{pmatrix}\chi(a)&0\\0&\chi(a)^{-1}\end{pmatrix},
  \qquad
  \rho_\chi(s)=
  \begin{pmatrix}0&1\\1&0\end{pmatrix}.
\]
Moreover $\rho_\chi\cong\rho_{\chi^{-1}}$, and these are the only
duplications in this list.
\end{enumerate}
\end{theorem}

\begin{proof}
Restrict an irreducible $G$-module $U$ to $A$.  Since $A$ is finite abelian
over $\C$, the restriction is a direct sum of weight spaces
\[
  U|_A=\bigoplus_{\chi\in\widehat A}U_\chi.
\]
For $u\in U_\chi$ and $a\in A$,
\[
  a(su)=s(sas)u=s(a^{-1}u)=\chi(a)^{-1}su,
\]
so $sU_\chi=U_{\chi^{-1}}$.

Suppose first that $\chi\ne\chi^{-1}$ occurs.  For nonzero
$u\in U_\chi$, the span of $u$ and $su$ is $G$-stable: $A$ acts through the
two weights $\chi,\chi^{-1}$, while $s$ exchanges the two lines.
Irreducibility forces this span to be all of $U$.  In the basis $(u,su)$ the
matrices are those displayed above.  The two distinct $A$-weights show that
this module is irreducible, and exchanging the basis vectors identifies
$\rho_\chi$ with $\rho_{\chi^{-1}}$.

Suppose instead that every occurring weight satisfies
$\eta=\eta^{-1}$.  Then $s$ preserves each $U_\eta$, so irreducibility forces
a single weight.  On that space $A$ acts by scalars and $s^2=1$.  The operator
$s$ is diagonalizable with eigenvalues $\pm1$, and each eigenspace is
$G$-stable.  Thus irreducibility forces dimension one and gives
$L_{\eta,+}$ or $L_{\eta,-}$.  Conversely, the relation $sas=a^{-1}$ in a
one-dimensional representation requires $\eta(a)=\eta(a)^{-1}$ for every
$a$, exactly the condition $\eta^2=1$.
\end{proof}

Choose a set $\Omega$ containing one representative from each unordered pair
$\{\chi,\chi^{-1}\}$ with $\chi\ne\chi^{-1}$.  Complete reducibility and
Theorem~\ref{generalized-dihedral-master-theorem:thm:classification} give the unique multiplicity data
\begin{equation}\label{generalized-dihedral-master-theorem:eq:decomposition}
  V\cong
  \bigoplus_{\substack{\eta^2=1\\\epsilon=\pm1}}
    L_{\eta,\epsilon}^{\oplus c_{\eta,\epsilon}}
  \oplus
  \bigoplus_{\chi\in\Omega}
    \rho_\chi^{\oplus m_\chi}.
\end{equation}
Put
\[
  c_\eta=c_{\eta,+}+c_{\eta,-},
  \qquad M=\sum_{\chi\in\Omega}m_\chi.
\]

\subsection{The master formula}

\begin{theorem}[Generalized-dihedral refined Molien formula]
\label{generalized-dihedral-master-theorem:thm:master}
For the representation \eqref{generalized-dihedral-master-theorem:eq:decomposition},
\[
  \boxed{\Mcal_{G,V}=\frac12\Rcal_V+\frac12\Scal_V,}
\]
where
\begin{align}
\Rcal_V
&=\frac1{|A|}\sum_{a\in A}\prod_{r\geq1}
\left[
\prod_{\eta^2=1}(1-\eta(a)t_r)^{-c_\eta}
\prod_{\chi\in\Omega}
\frac{1}{
(1-\chi(a)t_r)^{m_\chi}
(1-\chi(a)^{-1}t_r)^{m_\chi}}
\right],\label{generalized-dihedral-master-theorem:eq:rotation}\\[3pt]
\Scal_V
&=\frac1{|A|}\sum_{a\in A}\prod_{r\geq1}
\left[
(1-t_r^2)^{-M}
\prod_{\substack{\eta^2=1\\\epsilon=\pm1}}
(1-\epsilon\eta(a)t_r)^{-c_{\eta,\epsilon}}
\right].\label{generalized-dihedral-master-theorem:eq:reflection}
\end{align}
Consequently, for every partition $\tau$,
\[
  [m_\tau]\Mcal_{G,V}
  =\dim\left(
    \bigotimes_b\Sym^{\tau_b}V
  \right)^G.
\]
\end{theorem}

\begin{proof}
Split the normalized sum \eqref{generalized-dihedral-master-theorem:eq:molien-identity} across
$G=A\sqcup As$.  This produces the factor $1/2$ and the two normalized
$A$-averages.

For $a\in A$, each copy of $L_{\eta,\epsilon}$ has eigenvalue $\eta(a)$;
the sign $\epsilon$ does not enter on $A$.  Each copy of $\rho_\chi$ has
eigenvalues $\chi(a)$ and $\chi(a)^{-1}$.  Multiplying
$(1-t_r\lambda)^{-1}$ over this eigenvalue multiset proves
\eqref{generalized-dihedral-master-theorem:eq:rotation}.

For $as\in As$, a copy of $L_{\eta,\epsilon}$ has eigenvalue
$\epsilon\eta(a)$.  On a copy of $\rho_\chi$,
\[
  \rho_\chi(as)=
  \begin{pmatrix}
  0&\chi(a)\\
  \chi(a)^{-1}&0
  \end{pmatrix}.
\]
This matrix squares to $I$ and has trace zero, hence eigenvalues $+1$ and
$-1$, independently of $a$.  Its determinant factor is therefore
\[
  \det(I-t_r\rho_\chi(as))^{-1}
  =\frac1{(1-t_r)(1+t_r)}
  =(1-t_r^2)^{-1}.
\]
There are $M$ such blocks.  Multiplication over the one- and two-dimensional
summands proves \eqref{generalized-dihedral-master-theorem:eq:reflection}.  The coefficient statement follows
from \eqref{generalized-dihedral-master-theorem:eq:invariant-identity}.
\end{proof}

\begin{remark}[Image invariance]
The formula may be averaged either over the abstract group or uniformly over
the represented image. The two normalized averages agree by
Theorem~\ref{thm:normalized-image-invariance}; only unnormalized sums or
non-Haar laws can distinguish the kernel.
\end{remark}

\subsection{Coefficientwise verification method}

For $X$ with eigenvalue multiset $\Lambda(X)$, define
\[
  H_q(X):=[z^q]\prod_{\lambda\in\Lambda(X)}(1-\lambda z)^{-1}
  =h_q(\Lambda(X)).
\]
Extracting the monomial of type $\tau$ from Theorem~\ref{generalized-dihedral-master-theorem:thm:master} gives the
finite prediction
\begin{equation}\label{generalized-dihedral-master-theorem:eq:finite-prediction}
\begin{split}
  F_\tau
  =\frac1{2|A|}\sum_{a\in A}
  \bigg\{&
  \prod_b h_{\tau_b}\!\left(
    \{\eta(a)^{[c_\eta]}\}_{\eta^2=1},
    \{\chi(a)^{[m_\chi]},
      (\chi(a)^{-1})^{[m_\chi]}\}_{\chi\in\Omega}
  \right)\\
  &+
  \prod_b h_{\tau_b}\!\left(
    \{(\epsilon\eta(a))^{[c_{\eta,\epsilon}]}\}_{\eta,\epsilon},
    \{1^{[M]},(-1)^{[M]}\}
  \right)
  \bigg\}.
\end{split}
\end{equation}
Here $x^{[m]}$ denotes $m$ repetitions of $x$.  The notebook evaluates each
$h_q$ by truncated coefficient multiplication in
$\prod_\lambda(1-\lambda z)^{-1}$.

The independent direct route does not diagonalize reflection matrices and
does not use \eqref{generalized-dihedral-master-theorem:eq:finite-prediction}.  It performs the following steps.
\begin{enumerate}
\item Represent $A=C_{n_1}\times\cdots\times C_{n_r}$ by residue vectors and
construct the matrices $\rho(a)$ and $\rho(as)$ block by block.
\item Verify $s^2=I$ and $s\rho(a)s=\rho(a^{-1})$ on generators of $A$.
\item In the weak-composition monomial basis, expand the image of every
degree-$q$ monomial to construct the full matrix
$\Sym^q(\rho(g))$.
\item Form
\[
  \rho_{W_\tau}(g)
  =\bigotimes_b\Sym^{\tau_b}(\rho(g))
  \quad\text{and}\quad
  P_\tau=\frac1{|G|}\sum_{g\in G}\rho_{W_\tau}(g).
\]
\item Compare $\operatorname{rank}P_\tau$ and $\tr P_\tau$ with $F_\tau$,
while also checking $P_\tau^2=P_\tau$.
\end{enumerate}
Thus agreement is stronger than comparing two evaluations of the same
determinant expression.

\subsection{Representative cases and executed results}

The direct matrix dimension is
\[
  \dim W_\tau
  =\prod_b\binom{\tau_b+d-1}{\tau_b},
  \qquad d=\dim V.
\]
The tested cases were selected to keep this dimension moderate while exercising
each structural feature of the theorem.

\begin{center}
\small
\renewcommand{\arraystretch}{1.15}
\begin{tabular}{@{}p{0.39\linewidth}rrrr@{}}
\toprule
Representation & $|G|$ & $\dim V$ & checks & largest $\dim W_\tau$\\
\midrule
$D_5$, one standard two-dimensional block
  & 10 & 2 & 6 & 9\\
$D_6$, trivial and order-two 1D blocks plus one 2D block
  & 12 & 4 & 7 & 100\\
$\Dih(C_2\times C_2)$, self-inverse characters only
  & 8 & 5 & 6 & 75\\
$\Dih(C_3\times C_4)$, mixed 1D and 2D blocks
  & 24 & 4 & 7 & 100\\
$D_5$, multiplicity two on the same 2D block
  & 10 & 4 & 7 & 100\\
\midrule
\textbf{Total} & & & \textbf{33} & \textbf{100}\\
\bottomrule
\end{tabular}
\end{center}

The partitions include the empty partition, single blocks of degrees
$1,2,3$, repeated tensor blocks such as $(1,1)$ and $(2,2)$, and the mixed
partition $(2,1)$.  The noncyclic examples test character weights with several
moduli.  The even cyclic and elementary 2-group examples test the extra
self-inverse characters.  The last example tests the exponent $M$ and repeated
two-dimensional summands.

\medskip
\noindent\fcolorbox{accent}{soft}{\parbox{0.92\linewidth}{
\textbf{Executed result.} All 33 checks passed.  Across the full suite, the
maximum of the group-relation error, Reynolds-projector idempotence error,
trace/formula discrepancy, and distance of the formula from its nearest
integer was \(4.48\times10^{-15}\).  The largest explicitly materialized
target was \(100\times100\).
}}
\medskip

\begin{remark}
Finite testing is not a proof of the theorem.  Its role is to expose common
implementation and transcription errors: a missing factor $1/2$, omission of
one coset, incorrect normalization when passing to a represented image, incorrect
$\epsilon$ signs, failure to pair $\chi$ with $\chi^{-1}$, or replacing the
reflection contribution $(1-t_r^2)^{-M}$ by an $a$-dependent expression.
\end{remark}

\subsection{Conclusion}

The proposed classification and master formula are correct for every
finite-dimensional complex representation of $\Dih(A)$ with $A$ finite
abelian.  The proof reduces the representation theory to inversion orbits in
$\widehat A$ and then reads the determinant factors from the two cosets
$A$ and $As$.  The matrix verification agrees coefficientwise with this
formula in 33 small but structurally varied cases, including reasonable
representation dimensions and full symmetric-power projectors.
\endgroup


\clearpage
\section{Exact sigma-MGFs for dicyclic groups}\label{proof:dicyclic}
\begingroup
\definecolor{navy}{RGB}{24,55,95}
\definecolor{green}{RGB}{35,112,75}
\providecommand{\Dic}{\operatorname{Dic}}
\renewcommand{\Dic}{\operatorname{Dic}}
\providecommand{\Sym}{\operatorname{Sym}}
\renewcommand{\Sym}{\operatorname{Sym}}
\providecommand{\E}{\mathbb E}
\renewcommand{\E}{\mathbb E}
\setcounter{theorem}{0}
\setcounter{proposition}{0}
\setcounter{lemma}{0}
\setcounter{corollary}{0}
\setcounter{remark}{0}
\begin{abstract}
For the dicyclic group $\Dic_m$ of order $4m$, this note proves the proposed
generalized Molien formulas for every complex representation, specializes them
to the two-dimensional irreducibles $\rho_k$, and derives the stated exact
coefficient count.  It also checks direct sums and the regular representation.
The accompanying marimo notebook constructs the matrices themselves and
compares the rank of an explicit Reynolds projector with three independent
formula evaluations.  All representative cases passed to numerical error below
$2.6\times10^{-15}$.
\end{abstract}

\subsection{Imported finite-group reduction}
The determinant and coefficient formulas are Tools~\ref{tool:molien} and
\ref{tool:coefficient}. The dicyclic argument starts with the representation
classification and the spectra on the two cosets.

\subsection{Representations of \texorpdfstring{$\Dic_m$}{Dic-m}}
Use
\[
 \Dic_m=\langle a,x:a^{2m}=1,\ x^2=a^m,\ xax^{-1}=a^{-1}\rangle,
 \qquad \zeta=e^{\pi i/m}.
\]
Its elements are $a^j$ and $xa^j$ for $0\le j<2m$.
The four one-dimensional representations are
\[
 L_{\alpha,\beta}:\quad a\mapsto\alpha,\quad x\mapsto\beta,
 \qquad \alpha^2=1,\quad\beta^2=\alpha^m.
\]
For $1\le k\le m-1$, the two-dimensional representations are
\[
 \rho_k(a)=\begin{pmatrix}\zeta^k&0\\0&\zeta^{-k}\end{pmatrix},\qquad
 \rho_k(x)=\begin{pmatrix}0&(-1)^k\\1&0\end{pmatrix}.
\]
The displayed matrices satisfy all three defining relations.  Their two
$a$-weights are distinct and are interchanged by $x$, hence $\rho_k$ is
irreducible.  The representations $\rho_k$, $1\leq k\leq m-1$, are pairwise
inequivalent: $\chi_k(a)=\zeta^k+\zeta^{-k}=2\cos(\pi k/m)$, and these values
are distinct because cosine is strictly decreasing on $(0,\pi)$.  They are
inequivalent to the one-dimensional representations by dimension.  The four
one-dimensional characters are pairwise distinct by their values on $a$ and
$x$.  Therefore the displayed irreducibles are pairwise inequivalent; now
$4\cdot1^2+(m-1)\cdot2^2=4m$ proves that the list is complete.

Write an arbitrary representation as
\[
 V\cong\bigoplus_{\alpha,\beta}L_{\alpha,\beta}^{\oplus c_{\alpha,\beta}}
 \oplus\bigoplus_{k=1}^{m-1}\rho_k^{\oplus d_k}.
\]
On $a^j$, the displayed summands have eigenvalues $\alpha^j$ and
$\zeta^{kj},\zeta^{-kj}$.  On $xa^j$, they have eigenvalue
$\beta\alpha^j$ and a pair whose square is $(-1)^k$.  Thus
\begin{align*}
\SMG_{\Dic_m,V}
={}&\frac1{4m}\sum_{j=0}^{2m-1}\prod_i
 \left[\prod_{\alpha,\beta}(1-\alpha^jt_i)^{-c_{\alpha,\beta}}
 \prod_k\frac1{(1-\zeta^{kj}t_i)^{d_k}(1-\zeta^{-kj}t_i)^{d_k}}\right]\\
&+\frac1{4m}\sum_{j=0}^{2m-1}\prod_i
 \left[\prod_{\alpha,\beta}(1-\beta\alpha^jt_i)^{-c_{\alpha,\beta}}
 \prod_k(1-(-1)^kt_i^2)^{-d_k}\right].
\end{align*}
This proves the arbitrary-representation formula directly from the spectrum.

\subsection{A single two-dimensional irreducible}
For $V=\rho_k$ the preceding expression becomes
\[
 \boxed{\SMG_{\Dic_m,\rho_k}=
 \frac1{4m}\sum_{j=0}^{2m-1}\prod_i
 \frac1{(1-\zeta^{kj}t_i)(1-\zeta^{-kj}t_i)}
 +\frac12\prod_i\frac1{1-(-1)^kt_i^2}.}
\]
Let $q_k=2m/\gcd(2m,k)$.  In the monomial basis
$u^{\tau_b-r_b}v^{r_b}$ of $\Sym^{\tau_b}V$, the $a$-weight is
$\zeta^{k(\tau_b-2r_b)}$.  Hence the cyclic-subgroup average selects exactly
\[
 N_{k,\tau}=\#\left\{(r_b):0\le r_b\le\tau_b,\quad
 q_k\mid |\tau|-2\sum_b r_b\right\}.
\]
For a coset element the eigenvalues are $q,-q$, where $q^2=(-1)^k$, and
\[
 h_d(q,-q)=\begin{cases}q^d,&d\text{ even},\\0,&d\text{ odd}.
 \end{cases}
\]
Half the group lies in each part, giving the exact integer formula
\[
 \boxed{[m_\tau]\SMG_{\Dic_m,\rho_k}
 =\frac12N_{k,\tau}+\frac12\mathbf1_{\text{all }\tau_b\text{ even}}
 (-1)^{k|\tau|/2}.}
\]
For $k=1$, the coset factor is $(1+t_i^2)^{-1}$, recording that lifted
reflections square to $-I$.

For $V=\bigoplus_k\rho_k^{\oplus d_k}$, the identical argument yields half a
zero-weight lattice count and half $\prod_bB_{\tau_b}$, where
\[
 B_d=[x^d]\prod_{k=1}^{m-1}(1-(-1)^kx^2)^{-d_k}.
\]

\subsection{Regular representation}
An element of order $d$ acts on the regular representation as $4m/d$ cycles
of length $d$.  There are $\varphi(d)$ elements of order $d$ in
$\langle a\rangle$, while every $xa^j$ has order four.  Therefore
\[
 \boxed{\SMG_{\Dic_m,\mathrm{reg}}=
 \frac1{4m}\sum_{d\mid2m}\varphi(d)
 \operatorname{Exp}_\sigma\!\left(\frac{4m}{d}p_d\right)
 +\frac12\operatorname{Exp}_\sigma(mp_4).}
\]

\subsection{Verification strategy and results}
The notebook uses an orthonormal occupancy basis for each symmetric power.
It embeds that basis in $V^{\otimes q}$, computes
$\Sym^q(\rho(g))$ by compression, forms the tensor product $W_\tau$, and then
forms the full Reynolds projector.  This route does not use eigenvalues or the
coefficient formula.  It is compared with (i) a character average, (ii) the
congruence count above, and (iii) coefficient extraction from the two spectral
sums.  A safety cap prevents accidental construction of oversized tensor
spaces.

\begin{center}
\small
\begin{tabular}{@{}lrrrr@{}}
\toprule
Case & $\dim W_\tau$ & rank & formula & error\\
\midrule
$\Dic_2,\rho_1,\tau=(2)$ & 3 & 0 & 0 & $6.8\cdot10^{-17}$\\
$\Dic_3,\rho_1,\tau=(2,2)$ & 9 & 2 & 2 & $6.9\cdot10^{-16}$\\
$\Dic_4,\rho_2,\tau=(3,1)$ & 8 & 2 & 2 & $3.4\cdot10^{-16}$\\
$\Dic_5,\rho_3,\tau=(4)$ & 5 & 1 & 1 & $5.5\cdot10^{-16}$\\
$\Dic_3,L_{1,-1}\oplus L_{-1,i}\oplus\rho_2,\tau=(2,1)$ & 40 & 4 & 4 & $<2.6\cdot10^{-15}$\\
$\Dic_4,\rho_1\oplus\rho_2,\tau=(2,1)$ & 40 & 2 & 2 & $<10^{-14}$\\
$\Dic_2,\mathrm{reg},\tau=(4)$ & 330 & 44 & 44 & $0$\\
\bottomrule
\end{tabular}
\end{center}
Here ``error'' is the projector idempotence residual or, for the two added
direct-sum rows, a conservative bound including relation and formula residuals.
The dimensions $2$, $4$, and $8$ exercise irreducible, reducible, and regular
representations without hiding the computation behind large numerical spaces.

\paragraph{Conclusion.}
Every tested route agrees.  More importantly, the proof identifies why the
formula must hold: the determinant series gives symmetric-power characters,
and averaging those characters is exactly the trace/rank of the invariant
projector.  The dicyclic sign is forced by $(xa^j)^2=a^m$.
\endgroup


\clearpage
\section{Regular representations of finite groups}\label{proof:regular-representations}
\begingroup
\providecommand{\Sym}{\operatorname{Sym}}
\renewcommand{\Sym}{\operatorname{Sym}}
\providecommand{\Exp}{\operatorname{Exp}_{\sigma}}
\renewcommand{\Exp}{\operatorname{Exp}_{\sigma}}
\providecommand{\E}{\mathbf E}
\renewcommand{\E}{\mathbf E}
\providecommand{\tr}{\operatorname{Tr}}
\renewcommand{\tr}{\operatorname{Tr}}
\providecommand{\ord}{\operatorname{ord}}
\renewcommand{\ord}{\operatorname{ord}}
\setcounter{theorem}{0}
\setcounter{proposition}{0}
\setcounter{lemma}{0}
\setcounter{corollary}{0}
\setcounter{remark}{0}
\begin{abstract}
Using the imported coefficient and Molien identities, we derive the
element-order formula for the regular $\sigma$-MGF and then test the regular
representations of $C_5$, $D_4$ (order $8$), $S_3$, and $A_4$.  Every group is
constructed as a concrete left-regular matrix group.  Coefficients are computed
by an explicit Reynolds projector and compared with an independent order-count
formula; the full rational functions are also compared at safe numerical
specializations.  All reported cases agree to at least $10^{-14}$.
\end{abstract}

\subsection{Imported reduction}
We use Tools~\ref{tool:coefficient} and~\ref{tool:molien}. The
family-specific datum is that an element of order $d$ acts on the regular
basis as $|G|/d$ cycles of length $d$.

\subsection{The regular representation}
Let $R_G=\mathbb C[G]$ with the left action.  If $g$ has order $d$, left
multiplication partitions $G$ into $|G|/d$ cycles of length $d$.  Hence
\[
 \det(I-t\rho_{R_G}(g))=(1-t^d)^{|G|/d}.
\]
Writing $N_d(G)=\#\{g\in G:\ord(g)=d\}$ gives
\begin{equation}
 \boxed{\displaystyle
 \smgv{G}{R_G}=
 \sum_{d\geq1}\frac{N_d(G)}{|G|}
 \Exp\!\left(\frac{|G|}{d}p_d\right).}
 \label{regular-representations-proof-and-verification:eq:regular}
\end{equation}
Equivalently, at finitely many scalar variables,
\[
 \smgv{G}{R_G}(t_1,\ldots,t_s)
 =\frac1{|G|}\sum_dN_d(G)
   \prod_{i=1}^s(1-t_i^d)^{-|G|/d}.
\]
Thus the regular $\sigma$-MGF depends only on the element-order distribution.

The principal families follow immediately:
\begin{align*}
 \smgv{C_n}{R}&=\frac1n\sum_{d\mid n}\varphi(d)
       \Exp\!\left(\frac ndp_d\right),\\
 \smgv{D_n}{R}&=\frac1{2n}\left[
   \sum_{d\mid n}\varphi(d)\Exp\!\left(\frac{2n}{d}p_d\right)
   +n\Exp(np_2)\right],\qquad |D_n|=2n,\\
 \smgv{S_n}{R}&=\sum_{\lambda\vdash n}\frac1{z_\lambda}
   \Exp\!\left(\frac{n!}{o(\lambda)}p_{o(\lambda)}\right),\\
 \smgv{A_n}{R}&=2\!\sum_{\substack{\lambda\vdash n\\n-\ell(\lambda)\ {
m even}}}
   \frac1{z_\lambda}
   \Exp\!\left(\frac{n!/2}{o(\lambda)}p_{o(\lambda)}\right),
\end{align*}
where $z_\lambda=\prod_jj^{m_j}m_j!$ and
$o(\lambda)=\operatorname{lcm}\{j:m_j>0\}$.

\subsection{Independent verification method}
The implementation uses two routes.

\paragraph{Direct matrix route.}
Index the standard basis by $x\in G$ and set
$(P_g)_{gx,x}=1$.  On
$W_\tau=\bigotimes_a\Sym^{\tau_a}R_G$ form
\[
 \Pi_\tau=\frac1{|G|}\sum_{g\in G}
     \bigotimes_a\Sym^{\tau_a}(P_g).
\]
Because $\Pi_\tau$ is the Reynolds projector,
$\operatorname{rank}\Pi_\tau=\dim(W_\tau)^G$.  The code separately checks
$\|\Pi_\tau^2-\Pi_\tau\|$ and averages the symmetric-power characters.

\paragraph{Order-formula route.}
Element orders are computed in the abstract multiplication table.  For an
element of order $d$, the coefficient of $t^k$ in
$(1-t^d)^{-|G|/d}$ is zero unless $d\mid k$, and otherwise is
$\binom{|G|/d+k/d-1}{k/d}$.  Multiplication over the parts of $\tau$ and
averaging with $N_d(G)$ extracts the coefficient in \eqref{regular-representations-proof-and-verification:eq:regular}.  This
route never uses the Reynolds matrix.

As an untruncated check, both determinant expressions are evaluated at
$(t_1,t_2)=(0.07,-0.04)$.  These small values are far from all poles.

\subsection{Representative results}
\begin{center}
\small
\begin{tabular}{@{}lrrrrr@{}}
\toprule
case & $\dim R_G$ & $\dim W_\tau$ & projector & formula & max. error\\
\midrule
$C_5$, $\tau=(2,1)$ & 5 & 75 & 15 & 15 & $9.4\cdot10^{-16}$\\
$D_4$, $\tau=(2)$ & 8 & 36 & 7 & 7 & $5.0\cdot10^{-16}$\\
$S_3$, $\tau=(1,1)$ & 6 & 36 & 6 & 6 & $4.1\cdot10^{-16}$\\
$A_4$, $\tau=(2)$ & 12 & 78 & 8 & 8 & $3.4\cdot10^{-16}$\\
\bottomrule
\end{tabular}
\end{center}
Here ``max. error'' is the larger of the determinant-specialization error and
the Reynolds idempotence error.  Character averages equal the same displayed
integers in every row.

\subsection{Dimension choice and a stability diagnostic}
The test dimensions $5,8,6,12$ are large enough to include several element
orders but keep the largest explicit target at dimension $78$.  This makes a
full matrix projector, rather than only a character computation, practical.
The implementation caps interactive symmetric-power targets to prevent an
accidental exponential allocation.

Finally, the regular family cannot have a finite coefficientwise limit when
$|G|\to\infty$: directly,
\[
 [m_{(1,1)}]\smgv{G}{R_G}=\dim(R_G\otimes R_G)^G=|G|.
\]
This is a proof-level obstruction, not merely a numerical observation.

\subsection*{Reproducibility}
The canonical marimo file is
\path{marimo_notebooks/regular_representations.py}.
The shared engine, README, and pytest command remain at the package root.
\endgroup


\clearpage
\section{Representation variants: exterior powers and regular modules}\label{proof:representation-variants}
\begingroup
\definecolor{ink}{HTML}{17302A}
\definecolor{accent}{HTML}{287A62}
\definecolor{pale}{HTML}{EDF7F3}
\definecolor{passgreen}{HTML}{16794A}
\providecommand{\Sfun}{\SMG}
\renewcommand{\Sfun}{\SMG}
\providecommand{\Sym}{\operatorname{Sym}}
\renewcommand{\Sym}{\operatorname{Sym}}
\providecommand{\Reg}{\operatorname{Reg}}
\renewcommand{\Reg}{\operatorname{Reg}}
\providecommand{\PASS}{\textcolor{passgreen}{\textbf{PASS}}}
\renewcommand{\PASS}{\textcolor{passgreen}{\textbf{PASS}}}
\providecommand{\summarybox}[1]{\begin{center}\fcolorbox{accent}{pale}{\parbox{0.91\linewidth}{#1}}\end{center}}
\renewcommand{\summarybox}[1]{\begin{center}\fcolorbox{accent}{pale}{\parbox{0.91\linewidth}{#1}}\end{center}}
\setcounter{theorem}{0}
\setcounter{proposition}{0}
\setcounter{lemma}{0}
\setcounter{corollary}{0}
\setcounter{remark}{0}
\summarybox{\textbf{Outcome.}  Two structurally different substitutions into
the finite-group master formula were tested.  For $\Lambda^2$ of the standard
three-dimensional $S_4$ representation, explicit compound matrices agree with
the pairwise-eigenvalue formula.  For the regular representation of $S_3$,
explicit $6\times6$ permutation matrices agree with the element-order formula.
All 38 multivariate coefficients through total degree five pass, as do direct
Reynolds projector checks.}

\subsection{Imported master formula}
The universal finite-group identities are Tools~\ref{tool:coefficient} and
\ref{tool:molien}. This section tests two nontrivial substitutions into those
tools: an exterior power and a regular module.

\subsection{Exterior powers}

If $g|_V$ has eigenvalues $\lambda_1,\ldots,\lambda_n$, then the wedge vectors
$v_{i_1}\wedge\cdots\wedge v_{i_k}$ show that $g|_{\Lambda^kV}$ has eigenvalues
\[
\lambda_I=\prod_{i\in I}\lambda_i,\qquad |I|=k.
\]
Therefore
\begin{equation}
\Sfun_{G,\Lambda^kV}(\mathbf t)=\frac1{|G|}\sum_{g\in G}
\prod_{a\geq1}\prod_{|I|=k}(1-t_a\lambda_I(g))^{-1}.         \tag{3}
\end{equation}
This proof needs no assumption that $G$ is a reflection group.

\subsubsection{Concrete test: $\Lambda^2$ of the $S_4$ reflection representation}

Start from the permutation representation of $S_4$ on $\mathbb C^4$ and
restrict it to
\[
V=\{(x_1,x_2,x_3,x_4):x_1+x_2+x_3+x_4=0\}.
\]
Numerical QR orthonormalization supplies a $4\times3$ basis matrix $B$, so the
explicit standard matrices are $B^*P_\pi B$.  The direct exterior-square
matrices are constructed without diagonalizing: in the ordered basis
$e_i\wedge e_j$, their entries are the $2\times2$ minors
\begin{equation}
(\Lambda^2g)_{(k,l),(i,j)}
=g_{ki}g_{lj}-g_{kj}g_{li}.                                  \tag{4}
\end{equation}
The formula route instead diagonalizes only the original $3\times3$ standard
matrices and forms the three pairwise eigenvalue products required by (3).
Thus an error in the compound-matrix construction is not reproduced in the
formula route.

\subsection{The regular representation}

Let $G$ act on its basis by left multiplication.  If $g$ has order $o(g)$,
then its permutation of the $|G|$ basis elements has $|G|/o(g)$ cycles, each of
length $o(g)$.  Hence
\begin{equation}
\det(I-tg|_{\Reg})=(1-t^{o(g)})^{|G|/o(g)}.                   \tag{5}
\end{equation}
If $N_m$ is the number of elements of order $m$, substitution into (1) gives
\begin{equation}
\Sfun_{G,\Reg}(\mathbf t)=\sum_m\frac{N_m}{|G|}
\prod_{a\geq1}(1-t_a^m)^{-|G|/m}.                            \tag{6}
\end{equation}

For coefficient extraction, use
\[
[t^d](1-t^m)^{-c}=
\begin{cases}
\binom{c+d/m-1}{d/m},&m\mid d,\\
0,&m\nmid d.
\end{cases}
\]
For $S_3$, the order counts are
\begin{equation}
N_1=1,\qquad N_2=3,\qquad N_3=2.                              \tag{7}
\end{equation}
Equations (6)--(7) therefore give an integer-arithmetic formula for every
multivariate coefficient.  The direct route constructs all six left-regular
$6\times6$ permutation matrices and never consults these order counts.

\subsection{Verification protocol}

For each representation and every partition $\tau$ through total degree five:
\begin{enumerate}[leftmargin=*,itemsep=3pt]
\item Construct all representation matrices explicitly.
\item Recursively compute $h_d$ from each matrix's eigenvalues and average
$\prod_a h_{\tau_a}$ over the group.
\item Compare with the representation-specific formula: pairwise spectra in
(3) for $S_4$, or binomial coefficient extraction from (6) for $S_3$.
\item Embed the normalized occupancy basis of $\Sym^2(W)$ in $W^{\otimes2}$,
form the explicit Reynolds matrix, and compare rank with the coefficient.
\end{enumerate}

\subsection{Results}

\begin{center}
\begin{tabularx}{\linewidth}{@{}Xrrrrc@{}}
\toprule
Representation & $|G|$ & $\dim W$ & degree & checks & result\\
\midrule
$S_4$ on $\Lambda^2(V_{\mathrm{std}})$ & 24 & 3 & $0$--$5$ & 19 & \PASS\\
$S_3$ on $\Reg$ & 6 & 6 & $0$--$5$ & 19 & \PASS\\
\bottomrule
\end{tabularx}
\end{center}

Selected matched dimensions are
\begin{center}
\begin{tabular}{@{}lrrrrrr@{}}
\toprule
Representation & $(1)$ & $(2)$ & $(3)$ & $(2,1)$ & $(2,2)$ & $(3,2)$\\
\midrule
$\Lambda^2(V_{\mathrm{std}})$ of $S_4$ & 0 & 1 & 0 & 0 & 3 & 1\\
$\Reg$ of $S_3$ & 1 & 5 & 10 & 21 & 78 & 196\\
\bottomrule
\end{tabular}
\end{center}
The large mixed coefficients in the regular case provide a useful stress test:
agreement is not caused merely by both routes rounding small numbers to zero.

For the exterior case, $\Sym^2(W)$ has dimension six and the Reynolds
projector has rank $1$, trace $0.9999999999999992$, and idempotence error
$8.31\times10^{-16}$.  For the regular case, $\Sym^2(W)$ has dimension $21$;
the projector has rank $5$, trace $4.999999999999999$, and idempotence error
$3.72\times10^{-16}$.  These ranks agree with the $(2)$ coefficient in each
row of the table.

\subsection{What these tests establish}

The exterior test exercises a nonlinear eigenvalue transformation and a
coordinate-level compound-matrix construction.  The regular test exercises a
combinatorial compression from individual group elements to element-order
counts.  Together they show that the master formula is representation-agnostic
in precisely the claimed way: only the spectra on the chosen representation
change.

The tests are finite and numerical; the all-degree statements (3) and (6)
follow from the preceding eigenvalue and cycle proofs.  Numerical acceptance
requires the direct average to be within $3\times10^{-7}$ of an integer.  The
projector residuals are at machine precision.

\subsection{Reproducibility}

Run from the repository root:
\begin{verbatim}
PYTHONPATH=for_this_guy/matrix_group_formula_verification \
  .venv/bin/python \
  for_this_guy/matrix_group_formula_verification/\
  representation_variants/verification.py
\end{verbatim}
The marimo notebook allows the maximum degree to be changed interactively and
displays every coefficient comparison and both projector diagnostics.
\endgroup


\clearpage
\part{Symmetric, alternating, cube, and permutation actions}
\section[Exact finite-n sigma-MGF for S-n]{Exact finite-$n$ sigma-MGF for $S_n$}\label{proof:symmetric-exact}
\begingroup
\definecolor{navy}{RGB}{24,55,95}
\providecommand{\Sym}{\operatorname{Sym}}
\renewcommand{\Sym}{\operatorname{Sym}}
\setcounter{theorem}{0}
\setcounter{proposition}{0}
\setcounter{lemma}{0}
\setcounter{corollary}{0}
\setcounter{remark}{0}
\begin{abstract}
For the permutation representation $V=\mathbb C^n$, this note proves the exact
cycle-type formula for $\smg{S_n}$, identifies its coefficients with vector
partitions, and explains the stable range.  The accompanying marimo notebook
constructs all permutation matrices and compares a full Reynolds-projector
rank with both formulas.  All representative tests passed.
\end{abstract}

\subsection{Imported coefficient interpretation}
The coefficient identity is Tool~\ref{tool:coefficient}; its orbit form is
Tool~\ref{tool:permutation-orbits}. The new content is the exact finite-$n$
cycle-type formula, the vector-partition size constraint, and stabilization.

\subsection{Exact cycle-type formula}
Write a cycle type as $\lambda=(1^{m_1}2^{m_2}\cdots)\vdash n$ and set
\[
 z_\lambda=\prod_{d\ge1}d^{m_d}m_d!.
\]
A $d$-cycle has all $d$th roots of unity as eigenvalues, so
$\det(I-tP_d)=1-t^d$.  A permutation of type $\lambda$ therefore contributes
\[
 \prod_i\prod_{d\ge1}(1-t_i^d)^{-m_d(\lambda)}.
\]
The conjugacy class has size $n!/z_\lambda$, proving
\[
 \boxed{\smg{S_n}=\sum_{\lambda\vdash n}\frac1{z_\lambda}
 \prod_i\prod_{d\ge1}(1-t_i^d)^{-m_d(\lambda)}.}
\]
This statement is exact for every $n$ and in every degree; it is not a stable
approximation.

\subsection{Exact coefficient formula}
A monomial-basis vector of $W_\tau$ is encoded by a labeling
\[
 \ell:\{1,\ldots,n\}\longrightarrow\mathbb Z_{\ge0}^s,
 \qquad \sum_{j=1}^n\ell(j)=\tau
\]
coordinatewise.  Two labelings lie in the same $S_n$-orbit exactly when the
multiplicities
$a_{\mathbf v}=\#\{j:\ell(j)=\mathbf v\}$ agree.  Since the dimension of the
invariant space of a permutation module equals its number of orbits,
\[
 \boxed{[m_\tau]\smg{S_n}=\#\left\{(a_{\mathbf v})_{\mathbf v\ge0}:
 a_{\mathbf v}\ge0,\quad \sum_{\mathbf v}a_{\mathbf v}=n,\quad
 \sum_{\mathbf v}a_{\mathbf v}\mathbf v=\tau\right\}.}
\]
Equivalently, this is the number of vector partitions of $\tau$ into at most
$n$ nonzero parts; zero parts fill the remaining slots.

Marking the number of labels with $u$ gives
\[
 \boxed{\sum_{n\ge0}u^n\smg{S_n}
 =\prod_{\mathbf v\in\mathbb Z_{\ge0}^{(\mathbb N)}}(1-ut^{\mathbf v})^{-1}.}
\]
The exponential formula for permutation cycles gives the equivalent form
\[
 \sum_{n\ge0}u^n\smg{S_n}
 =\exp\left(\sum_{d\ge1}\frac{u^d}{d}
 \prod_i\frac1{1-t_i^d}\right).
\]

\subsection{Stable range}
Separate the zero label.  Every nonzero vector has total degree at least one,
so a vector partition of $\tau$ uses at most $|\tau|$ nonzero parts.  If
$n\ge|\tau|$, zeros can always be added and the coefficient stabilizes:
\[
 [m_\tau]\smg{S_n}=[m_\tau]\prod_{\mathbf v\ne0}(1-t^{\mathbf v})^{-1}.
\]
When $n<|\tau|$, the exact finite-$n$ formula simply removes decompositions
requiring too many nonzero parts.

\subsection{Verification method and evidence}
For each $g\in S_n$, the notebook creates its permutation matrix, compresses
$g^{\otimes q}$ to an orthonormal occupancy basis of $\Sym^q(V)$, tensors the
selected blocks, and averages the resulting matrices.  The projector rank is
compared independently with the symmetric-power character average, cycle-type
coefficient extraction, and a dynamic-programming vector-partition count.

\begin{center}
\begin{tabular}{@{}lrrrrr@{}}
\toprule
Case & $\dim W_\tau$ & rank & character & cycle & orbit\\
\midrule
$S_2,\tau=(3)$ & 4 & 2 & 2 & 2 & 2\\
$S_3,\tau=(2,1)$ & 18 & 4 & 4 & 4 & 4\\
$S_4,\tau=(2,2)$ & 100 & 9 & 9 & 9 & 9\\
$S_5,\tau=(2,1)$ & 75 & 4 & 4 & 4 & 4\\
\bottomrule
\end{tabular}
\end{center}
The largest projector here is $100\times100$ and averages 24 explicitly
constructed matrices.  Idempotence residuals range from
$2.2\times10^{-16}$ to $1.1\times10^{-15}$.  These cases cover degree above
the stable threshold, multiblock tensor products, and two different values of
$n$ for the same $\tau$.

\paragraph{Conclusion.}
The cycle-index and orbit formulas are two descriptions of the same invariant
dimension, and the explicit matrix calculations agree with both in every
tested case.
\endgroup


\clearpage
\section[Exact finite-n sigma-MGF for A-n]{Exact finite-$n$ sigma-MGF for $A_n$}\label{proof:alternating-exact}
\begingroup
\definecolor{navy}{RGB}{24,55,95}
\providecommand{\Sym}{\operatorname{Sym}}
\renewcommand{\Sym}{\operatorname{Sym}}
\setcounter{theorem}{0}
\setcounter{proposition}{0}
\setcounter{lemma}{0}
\setcounter{corollary}{0}
\setcounter{remark}{0}
\begin{abstract}
The asymptotic agreement with $S_n$ is imported from Theorem B; for $n\ge2$, this section proves the exact cycle-type formula for the permutation
representation of $A_n$, derives the correction to the $S_n$ vector-partition
count from orbit splitting, and proves equality with $S_n$ in degree at most
$n-2$.  Explicit even permutation matrices and Reynolds projectors verify the
formula, including a boundary case with a nonzero splitting correction.
\end{abstract}

\subsection{From \texorpdfstring{$S_n$}{S-n} to \texorpdfstring{$A_n$}{A-n}}
For any class function $F$ on $S_n$, the indicator of an even permutation gives
\[
 \mathbb E_{A_n}[F]=\frac1{n!}\sum_{g\in S_n}(1+\operatorname{sgn}(g))F(g).
\]
By Tool~\ref{tool:coefficient}, the coefficient of $m_\tau$ is
$\dim(\bigotimes_b\Sym^{\tau_b}(\mathbb C^n))^{A_n}$.  No new
Reynolds argument is needed.

If $\lambda\vdash n$, then
\[
 \operatorname{sgn}(\lambda)=(-1)^{n-\ell(\lambda)},\qquad
 z_\lambda=\prod_{d\ge1}d^{m_d(\lambda)}m_d(\lambda)!.
\]
Combining the sign indicator with the determinant of each permutation cycle
proves the exact expression
\[
 \boxed{\smg{A_n}=\sum_{\lambda\vdash n}
 \frac{1+(-1)^{n-\ell(\lambda)}}{z_\lambda}
 \prod_i\prod_{d\ge1}(1-t_i^d)^{-m_d(\lambda)}.}
\]
Only even cycle types remain, each with coefficient $2/z_\lambda$.

\subsection{Generating function}
Set
\[
 G_+(u)=\prod_{\mathbf v\ge0}(1-ut^{\mathbf v})^{-1},\qquad
 G_-(u)=\prod_{\mathbf v\ge0}(1+ut^{\mathbf v}).
\]
The first product records multisets of labels.  The second records sets of
distinct labels and is the sign-weighted cycle-index contribution.  Hence, for
$n\ge2$,
\[
 \boxed{\smg{A_n}=[u^n]\bigl(G_+(u)+G_-(u)\bigr).}
\]
Equivalently, applying the exponential formula to both factors gives
\begin{align*}
\smg{A_n}=[u^n]\Bigg[&\exp\left(\sum_{d\ge1}\frac{u^d}{d}
 \prod_i\frac1{1-t_i^d}\right)\\
&+\exp\left(\sum_{d\ge1}\frac{(-1)^{d-1}u^d}{d}
 \prod_i\frac1{1-t_i^d}\right)\Bigg].
\end{align*}

\subsection{Exact orbit interpretation}
An $S_n$-orbit of monomial-basis labelings is determined by multiplicities
$(a_{\mathbf v})_{\mathbf v\ge0}$.  Its stabilizer is
$\prod_{\mathbf v}S_{a_{\mathbf v}}$.  If some $a_{\mathbf v}\ge2$, the
stabilizer contains the odd transposition of two equally labeled points, so the
$S_n$-orbit remains a single $A_n$-orbit.  If every $a_{\mathbf v}$ is zero or
one, the stabilizer lies in $A_n$, and the orbit splits into two $A_n$-orbits.
Thus
\begin{align*}
[m_\tau]\smg{A_n}
={}&\#\left\{(a_{\mathbf v})\ge0:\sum a_{\mathbf v}=n,\ 
 \sum a_{\mathbf v}\mathbf v=\tau\right\}\\
&+\#\left\{(a_{\mathbf v})\in\{0,1\}:\sum a_{\mathbf v}=n,\ 
 \sum a_{\mathbf v}\mathbf v=\tau\right\}.
\end{align*}
The first term is the $S_n$ coefficient.  The correction counts sets of
distinct nonzero vectors summing to $\tau$, with either $n$ such vectors or
$n-1$ such vectors and one zero vector.

\subsection{Stable agreement with \texorpdfstring{$S_n$}{S-n}}
Any labeling of total degree $|\tau|$ has at most $|\tau|$ nonzero labels.  If
$|\tau|\le n-2$, at least two labels are zero.  Their transposition is odd and
lies in the stabilizer, so no $S_n$-orbit can split.  Therefore
\[
 \boxed{[m_\tau]\smg{A_n}=[m_\tau]\SMG_{S_n}
 \quad\text{whenever }|\tau|\le n-2.}
\]
Outside that range the distinct-label term is the exact correction, not merely
an error estimate.

\subsection{Verification method and evidence}
The notebook enumerates even permutations, constructs their permutation
matrices, builds the full matrices on each symmetric tensor block, and averages
them to obtain the $A_n$ Reynolds projector.  Its rank is compared independently
with the character average, the even-cycle-type formula, and the $S_n$ orbit
count plus a distinct-label dynamic-programming correction.

\begin{center}
\small
\begin{tabular}{@{}lrrrrrr@{}}
\toprule
Case & $\dim W_\tau$ & rank & cycle & $S_n$ & corr. & total\\
\midrule
$A_3,\tau=(2)$ & 6 & 2 & 2 & 2 & 0 & 2\\
$A_3,\tau=(3)$ & 10 & 4 & 4 & 3 & 1 & 4\\
$A_4,\tau=(2,1)$ & 40 & 4 & 4 & 4 & 0 & 4\\
$A_5,\tau=(2,2)$ & 225 & 9 & 9 & 9 & 0 & 9\\
$A_5,\tau=(3,1)$ & 175 & 7 & 7 & 7 & 0 & 7\\
\bottomrule
\end{tabular}
\end{center}
The $A_3,\tau=(3)$ row deliberately lies outside the stable range and confirms
that one $S_3$ orbit splits, increasing the invariant dimension from three to
four.  The remaining rows test multiblock targets and larger projectors.
Idempotence residuals are at most $1.1\times10^{-15}$.

\paragraph{Conclusion.}
The sign-weighted cycle formula and orbit-splitting formula agree exactly, and
both match direct matrix invariant dimensions in all representative cases.
\endgroup


\clearpage
\section[Selected representations of S-n]{Selected representations of $S_n$}\label{proof:symmetric-representations}
\begingroup
\providecommand{\Sym}{\operatorname{Sym}}
\renewcommand{\Sym}{\operatorname{Sym}}
\providecommand{\tr}{\operatorname{Tr}}
\renewcommand{\tr}{\operatorname{Tr}}
\setcounter{theorem}{0}
\setcounter{proposition}{0}
\setcounter{lemma}{0}
\setcounter{corollary}{0}
\setcounter{remark}{0}
\begin{abstract}
We verify the $\sigma$-MGF formulas for the sign and standard representations
of $S_n$, the symmetric/exterior/tensor character constructions, and the
induced representation $\operatorname{Ind}_{C_3}^{S_3}\omega$.  Matrices are
constructed element by element.  Reynolds ranks, character averages, cycle
formulas, and decorated coset-cycle determinants agree in all small cases.
\end{abstract}

\subsection{Sign representation}
For $n\geq2$, exactly half of $S_n$ has sign $+1$ and half has sign $-1$.
Therefore
\begin{equation}
 \boxed{\displaystyle
 \MGF_{S_n,\mathrm{sgn}}(t_\bullet)=\frac12\left[
   \prod_i(1-t_i)^{-1}+\prod_i(1+t_i)^{-1}\right].}
 \label{symmetric-group-representations-proof-and-verification:eq:sign}
\end{equation}
The coefficient of $m_\tau$ is $1$ when $|\tau|$ is even and $0$ when it is
odd.  This follows either by expanding \eqref{symmetric-group-representations-proof-and-verification:eq:sign}, or because
$\bigotimes_a\Sym^{\tau_a}(\mathrm{sgn})=\mathrm{sgn}^{|\tau|}$.

\subsection{The standard representation}
Let $M_n=\mathbb C^n$ be the natural permutation representation and
$V_n=S^{(n-1,1)}$.  Since $M_n=\mathbf1\oplus V_n$, for every permutation
$g$,
\[
 \prod_i\det(I-t_i g|_{V_n})^{-1}
 =\prod_i(1-t_i)\prod_i\det(I-t_i g|_{M_n})^{-1}.
\]
If $g$ has cycle type $\lambda$, then
$\det(I-tg|_{M_n})=\prod_{d\in\lambda}(1-t^d)$.  Hence the exact finite-$n$
formula is
\begin{equation}
 \boxed{\displaystyle
 \MGF_{S_n,V_n}(t_1,\ldots,t_s)=
 \sum_{\lambda\vdash n}\frac1{z_\lambda}
 \prod_{i=1}^s\left[(1-t_i)
    \prod_{d\in\lambda}(1-t_i^d)^{-1}\right].}
 \label{symmetric-group-representations-proof-and-verification:eq:standardfinite}
\end{equation}

The natural-action coefficient counts $S_n$-orbits on tuples of multisets.
Multiplication by $\prod_i(1-t_i)$ gives the independent coefficient rule
\[
 [t_1^{\tau_1}\cdots t_s^{\tau_s}]\MGF_{S_n,V_n}
 =\sum_{J\subseteq[s]}(-1)^{|J|}
 N_n(\tau-\mathbf1_J),
\]
where $N_n(\alpha)$ is the corresponding vector-partition count and terms with
negative coordinates vanish.  Once total degree is at most $n$, the familiar
stable form follows:
\[
 \MGF_{S_n,V_n}\equiv
 \operatorname{Exp}_\sigma(h_2+h_3+\cdots)\pmod{\deg>n}.
\]

\paragraph{Concrete matrices.}
The code uses the basis $e_i-e_n$ for $1\leq i<n$.  If $P_g$ is the natural
permutation matrix and $B=[e_1-e_n\ \cdots\ e_{n-1}-e_n]$, the standard matrix
is the coordinate matrix of $P_gB$.  This produces exact integer matrices of
dimension $n-1$.  Although this basis is not orthonormal, the rank and
idempotence of the Reynolds operator are basis-invariant.

\subsection{Symmetric, exterior, and tensor constructions}
For any representation $V$ and $g\in G$, the cycle-index expansions give
\begin{align*}
 \chi_{\Sym^kV}(g)
 &=\sum_{\mu\vdash k}\frac1{z_\mu}
   \prod_{r\geq1}\chi_V(g^r)^{m_r(\mu)},\\
 \chi_{\wedge^kV}(g)
 &=\sum_{\mu\vdash k}\frac{(-1)^{k-\ell(\mu)}}{z_\mu}
   \prod_{r\geq1}\chi_V(g^r)^{m_r(\mu)},\\
 \chi_{V^{\otimes k}}(g)&=\chi_V(g)^k.
\end{align*}
For $k=2$ these reduce to
\[
 \chi_{\Sym^2V}(g)=\frac{\chi(g)^2+\chi(g^2)}2,
 \qquad
 \chi_{\wedge^2V}(g)=\frac{\chi(g)^2-\chi(g^2)}2.
\]
The implementation constructs the symmetric-square matrices on an occupancy
basis, the exterior square in the two-dimensional $S_3$ case, and tensor
characters.  Maximum discrepancies are below $10^{-15}$.

\subsection{An induced representation and the coset-cycle proof}
Let $H\leq G$, let $W$ be an $H$-representation, and choose representatives
$x$ for $G/H$.  If $g$ has a coset cycle $C$ of length $\ell_C$ beginning at
$xH$, define $h_C=x^{-1}g^{\ell_C}x\in H$.

\begin{proposition}
For $V=\operatorname{Ind}_H^GW$,
\[
 \det(I-t\rho_V(g))=
 \prod_{C\subseteq G/H}\det(I-t^{\ell_C}\rho_W(h_C)).
\]
\end{proposition}

\begin{proof}
Along one coset cycle, $g$ cyclically moves $\ell_C$ copies of $W$.  After one
lap the map on the starting copy is $\rho_W(h_C)$.  An eigenvalue $\alpha$ of
this monodromy produces all roots of $z^{\ell_C}=\alpha$, whose determinant
factor is $1-\alpha t^{\ell_C}$.  Multiply over eigenvalues and cycles.
\end{proof}

The test takes $G=S_3$, $H=A_3=C_3$, and
$\omega((123))=e^{2\pi i/3}$.  Two left-coset representatives give explicit
$2\times2$ monomial matrices.  The resulting induced representation is the
standard irreducible.  Its direct determinant average is compared with the
coset-cycle product, including the complex monodromy scalars.

\newpage
\subsection{Representative results}
\begin{center}
\small
\begin{tabular}{@{}lrrrrr@{}}
\toprule
case & rep. dim. & target dim. & projector & formula & max. error\\
\midrule
$S_4$ sign, $(2)$ & 1 & 1 & 1 & 1 & $2.3\cdot10^{-16}$\\
$S_4$ sign, $(3)$ & 1 & 1 & 0 & 0 & $2.3\cdot10^{-16}$\\
$S_3$ standard, $(2,1)$ & 2 & 6 & 1 & 1 & $3.5\cdot10^{-16}$\\
$S_4$ standard, $(2,2)$ & 3 & 36 & 3 & 3 & $2.5\cdot10^{-15}$\\
$S_5$ standard, $(2,1)$ & 4 & 40 & 1 & 1 & $1.3\cdot10^{-15}$\\
$\operatorname{Ind}_{C_3}^{S_3}\omega$, $(2,1)$ & 2 & 6 & 1 & 1 & $7.1\cdot10^{-16}$\\
\bottomrule
\end{tabular}
\end{center}
The determinant tests use $(t_1,t_2)=(0.11,-0.06)$ for sign/standard and
$(0.09,-0.04)$ for the induced example.  These tests exercise the entire
rational expression, while the projector tests certify selected exact
coefficients.

\subsection*{Reproducibility}
Open
\path{marimo_notebooks/symmetric_representations.py}
with marimo.  The shared pytest suite reruns every row and the derived-character
checks.
\endgroup


\clearpage
\section{Specht, Foulkes, and free-Lie modules}\label{proof:new-dist:symmetric_modules}

\resultbox{\textbf{Scope.} This section collects the non-permutation
$S_n$-modules that were formerly grouped with spin covers and rooted trees.
The common Frobenius--Molien formula is followed by fixed-tail Specht limits,
Foulkes modules, and multilinear free-Lie modules.  Exact finite formulas are
kept separate from the low-degree coefficients that obstruct raw limits.}

\subsection{The common target and generalized Molien theorem}

Let $G$ be finite, let $\rho:G\to GL(V)$ be a complex representation with
character $\chi_V$, and let $\tau=(\tau_1,\ldots,\tau_s)$.  Put
\[
 W_\tau(V)=\bigotimes_{j=1}^s\Sym^{\tau_j}V,
 \qquad
 a_\tau(G,V)=\dim W_\tau(V)^G.
\]
The $\sigma$-MGF is the symmetric series
\[
 \MGF_{G,V}=\sum_\tau a_\tau(G,V)m_\tau.
\]

\begin{theorem}[Imported generalized Molien identity]\label{nd:symmetric_spin_lie_tree:thm:molien}
For every finite group and finite-dimensional complex representation,
\begin{equation}\label{nd:symmetric_spin_lie_tree:eq:molien}
 \boxed{
 \MGF_{G,V}
 =\frac1{|G|}\sum_{g\in G}
 \exp\!\left(\sum_{r\ge1}\frac{p_r}{r}\chi_V(g^r)\right)
 =\sum_\tau\dim W_\tau(V)^Gm_\tau.}
\end{equation}
Equivalently, in an alphabet $\mathbf t=(t_1,t_2,\ldots)$,
\begin{equation}\label{nd:symmetric_spin_lie_tree:eq:det-molien}
 \MGF_{G,V}(\mathbf t)
 =\frac1{|G|}\sum_{g\in G}\prod_i
 \det(I-t_i\rho(g))^{-1}.
\end{equation}
\end{theorem}

\begin{proof}[Source]
This is Tools~\ref{tool:coefficient} and~\ref{tool:molien}, together with the
standard logarithmic determinant expansion. It is quoted here only to fix the
notation used in the representation-specific formulas below.
\end{proof}

Thus a character table without power maps is not sufficient: the exponent
uses $\chi_V(g^r)$ for every $r$.  For a permutation $g$ of cycle type
$\mu=1^{m_1}2^{m_2}\cdots$, the $r$th power has type $\mu^{[r]}$ satisfying
\begin{equation}\label{nd:symmetric_spin_lie_tree:eq:power-map}
 \boxed{m_j(\mu^{[r]})=
 \sum_{\substack{d\mid r\\(j,r/d)=1}}d\,m_{jd}(\mu).}
\end{equation}
Indeed, a cycle of length $jd$ splits under the $r$th power into $d$ cycles
of length $j$ precisely when $d=(jd,r)$, which is equivalent to the displayed
divisibility and coprimality conditions.

\subsection{Ordinary symmetric-group modules}

\subsubsection{Frobenius master formula}

For an $S_n$-module $V_n$ with Frobenius characteristic $f_n=\ch(V_n)$,
the standard identity $\chi_{V_n}(\mu)=z_\mu[p_\mu]f_n$ and
Theorem~\ref{nd:symmetric_spin_lie_tree:thm:molien} give
\begin{equation}\label{nd:symmetric_spin_lie_tree:eq:frob-master}
 \boxed{
 \MGF_{V_n}
 =\sum_{\mu\vdash n}\frac1{z_\mu}
 \exp\!\left(
 \sum_{r\ge1}\frac{p_r}{r}
 z_{\mu^{[r]}}[p_{\mu^{[r]}}]f_n
 \right).}
\end{equation}
For a Specht module $S^\lambda$, this specializes to
\begin{equation}\label{nd:symmetric_spin_lie_tree:eq:specht}
 \boxed{
 \MGF_{S_n,S^\lambda}
 =\sum_{\mu\vdash n}\frac1{z_\mu}
 \exp\!\left(\sum_{r\ge1}\frac{p_r}{r}
 \chi^\lambda(\mu^{[r]})\right).}
\end{equation}
The formula is exact for every $n$ and $\lambda$; character evaluation can be
performed by Murnaghan--Nakayama or any equivalent character algorithm.

Every irreducible $S_n$-module is self-dual and admits an invariant symmetric
bilinear form.  Therefore, for a nontrivial irreducible,
\[
 [m_{(1)}]=0,
 \qquad [m_{(2)}]=[m_{(1,1)}]=1.
\]
The fourth multilinear coefficient already sees Kronecker complexity:
\[
 [m_{(1,1,1)}]=g(\lambda,\lambda,\lambda),
 \qquad
 [m_{(1,1,1,1)}]=\sum_{\nu\vdash n}g(\lambda,\lambda,\nu)^2.
\]
Consequently \eqref{nd:symmetric_spin_lie_tree:eq:specht} is a complete exact formula, but not a claim
that arbitrary Kronecker multiplicities have a simpler closed form.

\subsubsection{Hooks, two rows, and the matrix audit}

Let $C_j(\sigma)$ count $j$-cycles.  Since
$S^{(n-k,1^k)}\cong\bigwedge^kS^{(n-1,1)}$, its character is
\begin{equation}\label{nd:symmetric_spin_lie_tree:eq:hook}
 \boxed{Q_{1^k}(C)=[u^k]\frac1{1+u}
 \prod_{j\ge1}(1-(-u)^j)^{C_j}.}
\end{equation}
This follows by deleting the trivial eigenline from the natural permutation
module and taking the generating function of exterior-power characters.

For $n\ge2k$, let
\[
 A_k(C)=[u^k]\prod_{j\ge1}(1+u^j)^{C_j}.
\]
This is the number of fixed $k$-subsets.  Young's rule gives
\[
 \C\!\left[\binom{[n]}k\right]
 \cong\bigoplus_{j=0}^kS^{(n-j,j)},
\]
and subtraction of consecutive characters proves
\begin{equation}\label{nd:symmetric_spin_lie_tree:eq:two-row}
 \boxed{\chi^{(n-k,k)}=A_k-A_{k-1}.}
\end{equation}

The verifier constructs the standard module on the integral basis
$e_i-e_n$, takes exterior minors for the hook cases, and realizes
$S^{(n-2,2)}$ as the vertex-sum-zero subspace of the two-subset permutation
module.  It checks the traces of powers $1\le r\le4$ for every matrix against
\eqref{nd:symmetric_spin_lie_tree:eq:hook} or \eqref{nd:symmetric_spin_lie_tree:eq:two-row}.  It then constructs
$W_\tau(V)$ for $\tau=(1),(2),(1,1),(3)$, finds the common generator-fixed
space, and compares its dimension with the full Molien average.

\begin{center}
\begin{tabular}{@{}lrrrrr@{}}\toprule
module&dim.&$[m_1]$&$[m_2]$&$[m_{11}]$&$[m_3]$\\\midrule
$S^{(3,1)}$ of $S_4$&3&0&1&1&1\\
$S^{(3,1,1)}$ of $S_5$&6&0&1&1&0\\
$S^{(2,2)}$ of $S_4$&2&0&1&1&1\\
$S^{(3,2)}$ of $S_5$&5&0&1&1&1\\\bottomrule
\end{tabular}
\end{center}
All $1{,}152$ character-power comparisons and all $16$ fixed-space/Molien
comparisons pass.  Numerical rank residuals are at most $1.3\times10^{-15}$;
the matrices and character calculations themselves have integral entries.

\subsubsection{Fixed-tail asymptotics and the boundary of the claim}

Fix a partition $\lambda$ and put
$\lambda[n]=(n-|\lambda|,\lambda_1,\lambda_2,\ldots)$. For
$n\ge |\lambda|+\lambda_1$, the partition is defined and its irreducible
character is given by a fixed character polynomial
$Q_\lambda(C_1,C_2,\ldots)$.  Let independent
$Z_j\sim\operatorname{Poisson}(1/j)$ and define
\begin{equation}\label{nd:symmetric_spin_lie_tree:eq:poisson-power}
 Z_j^{[r]}=\sum_{\substack{d\mid r\\(j,r/d)=1}}dZ_{jd}.
\end{equation}
Then, coefficientwise in bounded total degree,
\begin{equation}\label{nd:symmetric_spin_lie_tree:eq:fixed-tail-limit}
 \boxed{
 \MGF_{\lambda,\infty}
 =\E\exp\!\left(
 \sum_{r\ge1}\frac{p_r}{r}
 Q_\lambda(Z_1^{[r]},Z_2^{[r]},\ldots)
 \right).}
\end{equation}

\begin{proof}
In total degree at most $D$, the exponential uses only $r\le D$ and finitely
many cycle variables in each character polynomial.  Formula
\eqref{nd:symmetric_spin_lie_tree:eq:power-map} expresses the cycle counts of $\sigma^r$ in those of
$\sigma$.  The finite collection of cycle counts of a uniform permutation
converges jointly, with all fixed moments, to the independent $Z_j$.
Substitution and moment convergence prove \eqref{nd:symmetric_spin_lie_tree:eq:fixed-tail-limit}.
\end{proof}

This is a polynomial Poisson-chaos law, not generally an ordinary Poisson
law.  It applies to fixed hooks and fixed two-row tails through
\eqref{nd:symmetric_spin_lie_tree:eq:hook}--\eqref{nd:symmetric_spin_lie_tree:eq:two-row}.  It does not apply to rectangles or
diagrams with two growing dimensions.  For those shapes \eqref{nd:symmetric_spin_lie_tree:eq:specht}
remains exact, while normalized character asymptotics alone do not imply a
raw $\sigma$-MGF limit; boundedness of the fourth coefficient is already a
necessary test.

\subsection{Foulkes permutation modules}

Let
\[
 H^{(a^b)}=\Ind_{S_a\wr S_b}^{S_{ab}}\mathbf1.
\]
It is the permutation module on partitions of $[ab]$ into $b$ unlabeled
blocks of size $a$, and its Frobenius characteristic is
\[
 \ch H^{(a^b)}=h_b[h_a].
\]
Substitution in \eqref{nd:symmetric_spin_lie_tree:eq:frob-master} proves
\begin{equation}\label{nd:symmetric_spin_lie_tree:eq:foulkes}
 \boxed{
 \MGF_{H^{(a^b)}}=
 \sum_{\mu\vdash ab}\frac1{z_\mu}
 \exp\!\left(
 \sum_{r\ge1}\frac{p_r}{r}
 z_{\mu^{[r]}}[p_{\mu^{[r]}}]h_b[h_a]
 \right).}
\end{equation}

The pair coefficient has a concrete orbit interpretation.  Given two block
partitions, label their blocks temporarily and set
$A_{ij}=|B_i\cap C_j|$.  Every row and column sum is $a$.  Relabeling the two
block systems permutes rows and columns, and two pairs lie in the same
$S_{ab}$-orbit exactly when the tables are so related.  Therefore
\begin{equation}\label{nd:symmetric_spin_lie_tree:eq:contingency}
 \boxed{[m_{(1,1)}]\MGF_{H^{(a^b)}}
 =\#\{A\in\mathbb Z_{\ge0}^{b\times b}:
 \text{all margins are }a\}/(S_b\times S_b).}
\end{equation}

\begin{theorem}[Foulkes pair-growth obstruction]
Let $c_{a,b}$ denote the coefficient in
\eqref{nd:symmetric_spin_lie_tree:eq:contingency}.
\begin{enumerate}[label=(\roman*)]
\item For fixed $a\ge2$, $c_{a,b}\ge p(b)$, where $p(b)$ is the partition
number.
\item For fixed $b\ge2$, $c_{a,b}\ge\lfloor a/2\rfloor+1$.
\end{enumerate}
Consequently the raw Foulkes $\sigma$-MGF has no finite coefficientwise limit
in either standard regime $b\to\infty$ with $a\ge2$ fixed or
$a\to\infty$ with $b\ge2$ fixed.
\end{theorem}
\begin{proof}
For a partition $b=s_1+\cdots+s_r$, take a block diagonal matrix with a
$1\times1$ block $[a]$ for $s_i=1$ and, for $s_i\ge2$, the block
$(a-1)I_{s_i}+P_{s_i}$, where $P_{s_i}$ is a cyclic permutation matrix. The
sizes of the connected components of the weighted bipartite support are
invariant under independent row and column permutations and recover the
partition of $b$. This proves (i).

For (ii), use the $2\times2$ block
$\left(\begin{smallmatrix}j&a-j\\a-j&j\end{smallmatrix}\right)$ for
$0\le j\le\lfloor a/2\rfloor$, together with $b-2$ blocks $[a]$. The unique
nontrivial weighted component and its edge weights distinguish these values
of $j$ up to the already imposed symmetry $j\leftrightarrow a-j$.
\end{proof}

The verification lists the block partitions, constructs every zero-one
permutation matrix, and independently expands
\[
 h_b[h_a]=\sum_{\nu\vdash b}\frac1{z_\nu}
 \prod_{r\in\nu}\left(\sum_{\lambda\vdash a}
 \frac{p_{r\lambda}}{z_\lambda}\right).
\]
Thus every matrix trace is compared with a power-sum coefficient, not with a
second traversal of the same action.  The results are
\begin{center}
\begin{tabular}{@{}lrrrrrr@{}}\toprule
module&dim.&$[m_1]$&$[m_2]$&$[m_{11}]$&$[m_3]$&$[m_{21}]$\\\midrule
$H^{(2^2)}$ of $S_4$&3&1&2&2&3&4\\
$H^{(2^3)}$ of $S_6$&15&1&3&3&8&12\\\bottomrule
\end{tabular}
\end{center}
All $744$ character comparisons and ten direct-orbit/Molien comparisons pass;
the contingency-table counts are respectively $2$ and $3$. The preceding
theorem, rather than the finite computations, proves divergence in the two
standard Foulkes regimes.

\subsection{The multilinear free-Lie representation}

The degree-$n$ multilinear free-Lie module has Frobenius characteristic
\begin{equation}\label{nd:symmetric_spin_lie_tree:eq:witt}
 \boxed{L_n=\ch\Lie_n=\frac1n\sum_{d\mid n}\mu(d)p_d^{n/d}.}
\end{equation}
Formula \eqref{nd:symmetric_spin_lie_tree:eq:frob-master} therefore yields
\begin{equation}\label{nd:symmetric_spin_lie_tree:eq:lie-smg}
 \boxed{
 \MGF_{\Lie_n}=
 \sum_{\mu\vdash n}\frac1{z_\mu}
 \exp\!\left(
 \sum_{r\ge1}\frac{p_r}{r}
 z_{\mu^{[r]}}[p_{\mu^{[r]}}]L_n
 \right).}
\end{equation}
Power sums are orthogonal with
$\langle p_\lambda,p_\lambda\rangle=z_\lambda$.  Distinct divisors in
\eqref{nd:symmetric_spin_lie_tree:eq:witt} give distinct cycle types, hence
\begin{equation}\label{nd:symmetric_spin_lie_tree:eq:lie-norm}
 \boxed{[m_{(1,1)}]\MGF_{\Lie_n}
 =\langle L_n,L_n\rangle
 =\frac1{n^2}\sum_{d\mid n}\mu(d)^2d^{n/d}(n/d)!.}
\end{equation}
The $d=1$ term gives the lower bound $n!/n^2$, so the raw series cannot have
a coefficientwise limit as $n\to\infty$.  A sign twist, such as that appearing
in top partition-lattice homology, leaves this norm unchanged.

For an independent matrix construction, fix $x_n$ and use the bracket basis
\[
 [\cdots[[x_n,x_{\sigma(1)}],x_{\sigma(2)}],\ldots,
 x_{\sigma(n-1)}],\qquad \sigma\in S_{n-1}.
\]
Expanding brackets in the multilinear associative word basis gives integral
matrices for relabeling by $S_n$.  For $n=3,4$, all $30$ matrix characters
match \eqref{nd:symmetric_spin_lie_tree:eq:witt}.  Direct fixed spaces and Molien averages give
\begin{center}
\begin{tabular}{@{}lrrrrr@{}}\toprule
module&dim.&$[m_1]$&$[m_2]$&$[m_{11}]$&$[m_3]$\\\midrule
$\Lie_3$&2&0&1&1&1\\
$\Lie_4$&6&0&2&2&2\\\bottomrule
\end{tabular}
\end{center}
The pair values $1,2$ agree exactly with \eqref{nd:symmetric_spin_lie_tree:eq:lie-norm}.


\subsection{Verification and asymptotic scope}

The matrix/fixed-space route constructs integral Specht and free-Lie matrices,
while the formula route uses power maps, character polynomials, power-sum
plethysm, and the Witt characteristic.  Foulkes coefficients are also checked
by literal orbit enumeration.  The relevant conclusions are:
\begin{center}
\small
\begin{tabularx}{\linewidth}{@{}>{\raggedright\arraybackslash}p{0.27\linewidth}
>{\raggedright\arraybackslash}p{0.29\linewidth}X@{}}\toprule
family&exact finite formula&asymptotic status\\\midrule
fixed-tail Specht&\eqref{nd:symmetric_spin_lie_tree:eq:specht},
\eqref{nd:symmetric_spin_lie_tree:eq:fixed-tail-limit}&explicit polynomial Poisson chaos\\
hooks and two rows&\eqref{nd:symmetric_spin_lie_tree:eq:hook},
\eqref{nd:symmetric_spin_lie_tree:eq:two-row}&explicit in the fixed-tail regime\\
growing diagrams&\eqref{nd:symmetric_spin_lie_tree:eq:specht}&normalized character theory alone does not control raw moments\\
Foulkes&\eqref{nd:symmetric_spin_lie_tree:eq:foulkes},
\eqref{nd:symmetric_spin_lie_tree:eq:contingency}&pair-orbit growth obstructs standard raw limits\\
$\Lie_n$&\eqref{nd:symmetric_spin_lie_tree:eq:lie-smg},
\eqref{nd:symmetric_spin_lie_tree:eq:lie-norm}&$[m_{11}]\ge n!/n^2$, hence no raw coefficientwise limit\\\bottomrule
\end{tabularx}
\end{center}

\clearpage
\section{Subset actions and the universal permutation method}\label{proof:subset-actions}
\begingroup
\providecommand{\Sym}{\operatorname{Sym}}
\renewcommand{\Sym}{\operatorname{Sym}}
\providecommand{\Fix}{\operatorname{Fix}}
\renewcommand{\Fix}{\operatorname{Fix}}
\setcounter{theorem}{0}
\setcounter{proposition}{0}
\setcounter{lemma}{0}
\setcounter{corollary}{0}
\setcounter{remark}{0}
\begin{abstract}
We prove and test the fixed-point, Mobius-inversion, determinant, and orbit
formulas for the permutation representation on $k$-subsets.  Every induced
permutation matrix is constructed for small $S_n$ and $A_n$.  Direct Reynolds
ranks agree with the independent cycle formula in target spaces up to dimension
$550$.  The symmetric-group quadratic coefficient $k+1$ is checked at its
boundary, and the alternating-group orbit-splitting boundary is tested
separately.
\end{abstract}

\subsection{Imported permutation tools}
For $V=\C[\Omega]$, Tools~\ref{tool:molien} and
\ref{tool:permutation-orbits} give
\begin{equation}\label{permutation-actions-proof-and-verification:eq:cycle}
\MGF_{G,V}(t_\bullet)=\frac1{|G|}\sum_{g\in G}
\prod_{i,d\ge1}(1-t_i^d)^{-c_d^\Omega(g)},
\end{equation}
\begin{equation}\label{permutation-actions-proof-and-verification:eq:orbit}
[m_\tau]\MGF_{G,V}=\#\!\left(G\backslash
\prod_a\Sym^{\tau_a}(\Omega)\right).
\end{equation}
These identities are quoted, not reproved. The new calculation begins with
the conversion from fixed points of powers to induced cycles.

\subsection{Recovering cycles from fixed points}
For any permutation $q$ of finite order,
\[
 \#\Fix(q^r)=\sum_{d\mid r}d\,c_d(q).
\]
Mobius inversion gives
\begin{equation}
 \boxed{\displaystyle
 c_d(q)=\frac1d\sum_{r\mid d}\mu(d/r)\#\Fix(q^r).}
 \label{permutation-actions-proof-and-verification:eq:mobius}
\end{equation}
This identity is exact: the code checks divisibility by $d$ before accepting a
cycle count.

Now take $G=S_n$ or $A_n$ and
$\Omega=\binom{[n]}k$.  A $k$-subset fixed by $\sigma^r$ is a union of cycles
of $\sigma^r$.  If $X_j(\sigma^r)$ is its number of $j$-cycles, then
\begin{equation}
 \boxed{\displaystyle
 \#\Fix_{\binom{[n]}k}(\sigma^r)
 =[z^k]\prod_{j\geq1}(1+z^j)^{X_j(\sigma^r)}.}
 \label{permutation-actions-proof-and-verification:eq:fixed}
\end{equation}
Equations \eqref{permutation-actions-proof-and-verification:eq:cycle}, \eqref{permutation-actions-proof-and-verification:eq:mobius}, and \eqref{permutation-actions-proof-and-verification:eq:fixed} are the
proposed exact finite formula.

\subsection{Stability and the first nontrivial coefficients}
For fixed $k$ and $\tau$, a tuple of $|\tau|$ many $k$-subsets uses at most
$k|\tau|$ labels.  Deleting isolated labels identifies the stable orbits with
finite colored $k$-uniform multihypergraphs.  Consequently
\[
 [m_\tau]\MGF_{S_n,\mathbb C[\binom{[n]}k]}
 \quad\hbox{stabilizes for}\quad n\geq k|\tau|.
\]
For $A_n$, the same coefficients agree once two unused labels are available,
so $n\geq k|\tau|+2$ is sufficient: compose any odd relabeling with a
transposition of the unused labels.

For $S_n$, when $n\geq2k$, an ordered or unordered pair of $k$-subsets is
classified by $|A\cap B|\in\{0,1,\ldots,k\}$.  Thus
\begin{equation}
 [m_{(2)}]\MGF=[m_{(1,1)}]\MGF=k+1
 \quad\text{for $S_n$ and $n\geq2k$.}
 \label{permutation-actions-proof-and-verification:eq:quadratic}
\end{equation}
For $A_n$, the same conclusion is guaranteed once two unused labels are
available, so $n\geq2k+2$ is sufficient.  At $n=2k$ and $n=2k+1$, an $S_n$
orbit may split and must be checked separately.  For example, $A_4$ acting on
2-subsets has three orbits on unordered pairs but four orbits on ordered
pairs: the ordered adjacent-pair class splits into two $A_4$-orbits.  Thus
\[
 [m_{(2)}]\MGF=3,\qquad [m_{(1,1)}]\MGF=4
 \quad\text{for $(n,k)=(4,2)$.}
\]
The regression suite checks the symmetric boundary and the alternating cases
$A_4,A_5,A_6$ separately.

\subsection{Four independent computational layers}
\begin{enumerate}
\item Enumerate $S_n$ (or its even subgroup) as image-tuples and construct its
      action on the lexicographically ordered $k$-subsets.
\item Build the full Reynolds projector on
      $\bigotimes_a\Sym^{\tau_a}\mathbb C[\binom{[n]}k]$ and take its numerical
      rank.  Its idempotence residual is recorded.
\item Without examining the induced matrix cycles, compute
      \eqref{permutation-actions-proof-and-verification:eq:fixed}, apply \eqref{permutation-actions-proof-and-verification:eq:mobius}, and extract the coefficient of
      \eqref{permutation-actions-proof-and-verification:eq:cycle} using integer dynamic programming.
\item Evaluate both full determinant expressions at
      $(t_1,t_2)=(0.08,-0.03)$, safely inside the common disk of convergence.
\end{enumerate}
The symmetric-power occupancy basis has dimension
$\binom{\dim V+d-1}{d}$.  The chosen cases therefore exercise nontrivial
spaces while keeping full projectors feasible.

\subsection{Representative results}
\begin{center}
\small
\begin{tabular}{@{}lrrrrr@{}}
\toprule
case & rep. dim. & target dim. & projector & formula & max. error\\
\midrule
$S_4$ on 2-subsets, $(2)$ & 6 & 21 & 3 & 3 & $2.4\cdot10^{-16}$\\
$S_5$ on 2-subsets, $(1,1)$ & 10 & 100 & 3 & 3 & $2.6\cdot10^{-16}$\\
$S_5$ on 2-subsets, $(2,1)$ & 10 & 550 & 11 & 11 & $2.2\cdot10^{-15}$\\
$A_5$ on 2-subsets, $(2)$ & 10 & 55 & 3 & 3 & $2.6\cdot10^{-16}$\\
$A_4$ on 2-subsets, $(1,1)$ & 6 & 36 & 4 & 4 & $<10^{-14}$\\
$A_5$ on 2-subsets, $(1,1)$ & 10 & 100 & 3 & 3 & $<10^{-14}$\\
$A_6$ on 2-subsets, $(1,1)$ & 15 & 225 & 3 & 3 & $<10^{-14}$\\
\bottomrule
\end{tabular}
\end{center}
The character average equals the same integer in every row.  The $550$
dimensional case is a useful stress test: the exact cycle calculation predicts
$11$, and the independently assembled projector has rank $11$.

\subsection*{Reproducibility}
The interactive file is
\path{marimo_notebooks/permutation_actions.py}.  It
exposes $n$, $k$, $\tau$,
and the choice $S_n/A_n$, with conservative control ranges.  The shared pytest
suite also checks the stability boundary.
\endgroup


\clearpage
\section{Johnson-scheme permutation representations}\label{proof:johnson}
\begingroup
\providecommand{\Sym}{\operatorname{Sym}}
\renewcommand{\Sym}{\operatorname{Sym}}
\providecommand{\Fix}{\operatorname{Fix}}
\renewcommand{\Fix}{\operatorname{Fix}}
\providecommand{\Aut}{\operatorname{Aut}}
\renewcommand{\Aut}{\operatorname{Aut}}
\providecommand{\cS}{\SMG}
\renewcommand{\cS}{\SMG}
\setcounter{theorem}{0}
\setcounter{proposition}{0}
\setcounter{lemma}{0}
\setcounter{corollary}{0}
\setcounter{remark}{0}
\begin{abstract}
For the permutation representation on the $k$-subsets of $[v]$, we prove the
cycle-factor and orbit formulas, derive the fixed-subset formula, treat the
middle-dimensional complementation coset, and verify all coefficients through
total degree three in three representative finite cases.  The computation uses
the actual permutation matrices and does not assume the desired cycle formula.
\end{abstract}

\subsection{Imported permutation tools}
Let $c_\ell(g)$ count the $\ell$-cycles of $g$ on $X$.  Tools~\ref{tool:molien}
and~\ref{tool:permutation-orbits} give
\begin{align}
 \smgv{G}{\mathbb C[X]}
 &=\frac1{|G|}\sum_{g\in G}
 \prod_{i\ge1}\prod_{\ell\ge1}(1-t_i^\ell)^{-c_\ell(g)},
 \label{johnson-proof-and-verification:eq:master}\\
 [m_\tau]\,\smgv{G}{\mathbb C[X]}
 &=\#\bigl(G\backslash
 (\Sym^{\tau_1}X\times\cdots\times\Sym^{\tau_s}X)\bigr).
 \label{johnson-proof-and-verification:eq:orbit}
\end{align}
The Johnson-specific work begins with the fixed-subset calculation below.

\subsection{Johnson fixed points and complementation}
Put $X=\binom{[v]}k$.  If $\sigma\in S_v$ has $a_j$ cycles of length $j$,
then a $j$-cycle splits under $\sigma^r$ into $\gcd(j,r)$ cycles of length
$j/\gcd(j,r)$.  A fixed subset is exactly a union of these cycles; hence
\begin{equation}\label{johnson-proof-and-verification:eq:johnsonfixed}
 F_r^J(\sigma)=[u^k]\prod_{j\ge1}
 \left(1+u^{j/\gcd(j,r)}\right)^{a_j\gcd(j,r)}.
\end{equation}
The vertex-cycle counts follow from
\begin{equation}\label{johnson-proof-and-verification:eq:mobius}
 c_\ell(\sigma)=\frac1\ell\sum_{d\mid\ell}
 \mu(\ell/d)F_d^J(\sigma).
\end{equation}

When $v=2k$ and $k>1$, complementation $T(A)=A^c$ commutes with $S_{2k}$
and supplies the additional coset in the full Johnson-scheme automorphism
group.  For odd $r$, a subset fixed by $\sigma^rT$ must alternate membership
around every cycle of $\sigma^r$.  Thus
\begin{equation}\label{johnson-proof-and-verification:eq:coset}
F_r(\sigma T)=\begin{cases}
F_r^J(\sigma),&r\ \text{even},\\
2^{c(\sigma^r)},&r\ \text{odd and every cycle of $\sigma^r$ is even},\\
0,&\text{otherwise}.
\end{cases}
\end{equation}
Every alternating choice has size exactly half of $v$, so no further size
constraint is missing.  Equations \eqref{johnson-proof-and-verification:eq:master}--\eqref{johnson-proof-and-verification:eq:coset} prove
the claimed formula, including the middle case.

\subsection{Verification design}
For each group element, the code explicitly builds its $|X|\times|X|$
permutation matrix $P$.  It computes $h_d(P)=\operatorname{tr}(\Sym^dP)$ from
the matrix traces by Newton's recurrence
\[
 d h_d(P)=\sum_{r=1}^d\operatorname{tr}(P^r)h_{d-r}(P),\qquad h_0=1.
\]
It averages $\prod_j h_{\tau_j}(P)$ over the group.  Independently, it
extracts the same coefficient from the cycle factors in \eqref{johnson-proof-and-verification:eq:master}
and counts the orbits in \eqref{johnson-proof-and-verification:eq:orbit} by explicit set traversal.  Finally,
it evaluates \eqref{johnson-proof-and-verification:eq:johnsonfixed} or \eqref{johnson-proof-and-verification:eq:coset} for every required
power and checks that \eqref{johnson-proof-and-verification:eq:mobius} reproduces every actual cycle count.

\subsection{Exact results}
The entries list the common value in the order
$\tau=(1),(2),(1,1),(3),(2,1),(1,1,1)$.
\begin{center}
\begin{tabular}{@{}lrrl@{}}
\toprule
Action & $|G|$ & $|X|$ & common coefficient vector\\
\midrule
$\Aut J(4,2)$ (including $T$) & 48 & 6 & $(1,3,3,5,8,11)$\\
$S_5$ on $2$-subsets & 120 & 10 & $(1,3,3,7,11,15)$\\
$S_5$ on $1$-subsets & 120 & 5 & $(1,2,2,3,4,5)$\\
\bottomrule
\end{tabular}
\end{center}
All 18 triple comparisons are exact, and all fixed-point reconstructions pass.

At $(t_1,t_2)=(0.037,0.061)$ the direct determinant average and cycle product
for $\Aut J(4,2)$ coincide at
\[
1.1237420509116285\qquad\text{(reported error $0$)}.
\]
The pair values also give $[m_{(2)}]=[m_{(1,1)}]=\min(k,v-k)+1$, as predicted.

\subsection{Reproduction}
From the repository root run
\begin{verbatim}
marimo run marimo_notebooks/johnson_scheme.py
\end{verbatim}
or run the family module directly with Python.  The notebook displays every
matrix/cycle/orbit comparison rather than only the summary above.
\endgroup


\clearpage
\section{Hamming-scheme permutation representations}\label{proof:hamming}
\begingroup
\providecommand{\Sym}{\operatorname{Sym}}
\renewcommand{\Sym}{\operatorname{Sym}}
\providecommand{\cS}{\SMG}
\renewcommand{\cS}{\SMG}
\setcounter{theorem}{0}
\setcounter{proposition}{0}
\setcounter{lemma}{0}
\setcounter{corollary}{0}
\setcounter{remark}{0}
\begin{abstract}
We prove and test the generalized Molien formula for the full Hamming graph
automorphism group $S_q\wr S_D=S_q^D\rtimes S_D$ acting on $Q^D$.  The tests
construct every wreath-product element, not merely pure coordinate
permutations, and compare matrix, cycle, and orbit calculations through total
degree three.
\end{abstract}

\subsection{Imported permutation tools}
Tools~\ref{tool:molien} and~\ref{tool:permutation-orbits} give
\begin{equation}\label{hamming-proof-and-verification:eq:master}
\cS_{G,X}(\mathbf t)=\frac1{|G|}\sum_{g\in G}
\prod_{i,\ell}(1-t_i^\ell)^{-c_\ell(g)},
\end{equation}
\begin{equation}\label{hamming-proof-and-verification:eq:orbit}
[m_\tau]\cS_{G,X}=\#\!\left(G\backslash
\prod_j\Sym^{\tau_j}(X)\right).
\end{equation}
They are quoted here only to fix notation. The family-specific step is the
fixed-word formula for a full Hamming wreath-product element.

\subsection{Fixed words in the full wreath product}
Write $g=(\sigma_1,\ldots,\sigma_D;\pi)$.  For a fixed power $g^r$, follow a
coordinate around a cycle $C=(i_1,\ldots,i_s)$ of its coordinate permutation.
The successive alphabet relabelings compose to a permutation $h_{C,r}\in S_q$.
A fixed word is determined on $C$ by one symbol fixed by $h_{C,r}$, and
different coordinate cycles are independent.  Therefore
\begin{equation}\label{hamming-proof-and-verification:eq:fixed}
F_r^H(g)=\prod_{C\in\operatorname{Cyc}(\pi^r)}
|\operatorname{Fix}_Q(h_{C,r})|.
\end{equation}
The exact vertex-cycle counts are then
\begin{equation}\label{hamming-proof-and-verification:eq:mobius}
c_\ell^H(g)=\frac1\ell\sum_{d\mid\ell}\mu(\ell/d)F_d^H(g).
\end{equation}
Substitution in \eqref{hamming-proof-and-verification:eq:master} proves the proposed Hamming formula.  For
pure coordinate permutations the label products are identities, reducing
\eqref{hamming-proof-and-verification:eq:fixed} to $q^{c(\pi^r)}$.

Two ordered pairs of words are equivalent exactly when they have equal
Hamming distance.  Hence vertex transitivity and the distances $0,\ldots,D$
give
\[
[m_{(1)}]=1,\qquad [m_{(2)}]=[m_{(1,1)}]=D+1.
\]

\subsection{Independent computational checks}
The implementation explicitly constructs every word permutation and its
$q^D\times q^D$ matrix $P$.  It computes complete symmetric characters from
actual matrix traces,
\[
d h_d(P)=\sum_{r=1}^d\operatorname{tr}(P^r)h_{d-r}(P),\quad h_0=1,
\]
and averages $\prod_jh_{\tau_j}(P)$.  A second route extracts coefficients
from \eqref{hamming-proof-and-verification:eq:master}; a third explicitly traverses the orbits in
\eqref{hamming-proof-and-verification:eq:orbit}.  Separately, the coordinate maps and alphabet labels are
powered by composition, \eqref{hamming-proof-and-verification:eq:fixed} is evaluated, and \eqref{hamming-proof-and-verification:eq:mobius}
is required to reproduce the actual vertex cycles for every group element.

\subsection{Exact results}
Values are ordered by
$\tau=(1),(2),(1,1),(3),(2,1),(1,1,1)$.
\begin{center}
\begin{tabular}{@{}lrrl@{}}
\toprule
Action & $|G|$ & $|X|$ & common coefficient vector\\
\midrule
$S_2\wr S_2$ on $Q^2$ & 8 & 4 & $(1,3,3,4,7,10)$\\
$S_2\wr S_3$ on $Q^3$ & 48 & 8 & $(1,4,4,7,13,20)$\\
$S_3\wr S_2$ on $Q^2$ & 72 & 9 & $(1,3,3,7,11,15)$\\
\bottomrule
\end{tabular}
\end{center}
All 18 matrix/cycle/orbit comparisons and every fixed-word reconstruction
pass.  For $(q,D)=(3,2)$ at $(t_1,t_2)=(0.037,0.061)$, the determinant and
cycle averages are $1.1254339873575003$ and $1.1254339873575000$; the absolute
error is $2.22\times10^{-16}$.

\subsection{Reproduction}
Run
\begin{verbatim}
marimo run marimo_notebooks/hamming_scheme.py
\end{verbatim}
from the repository root.  The table in the notebook exposes every exact
comparison and the global fixed-point reconstruction result.
\endgroup


\clearpage
\section{Grassmann-scheme actions}\label{proof:grassmann}
\begingroup
\providecommand{\Sym}{\operatorname{Sym}}
\renewcommand{\Sym}{\operatorname{Sym}}
\providecommand{\Gr}{\operatorname{Gr}}
\renewcommand{\Gr}{\operatorname{Gr}}
\providecommand{\cS}{\SMG}
\renewcommand{\cS}{\SMG}
\setcounter{theorem}{0}
\setcounter{proposition}{0}
\setcounter{lemma}{0}
\setcounter{corollary}{0}
\setcounter{remark}{0}
\begin{abstract}
For $GL_N(q)$ acting on the $k$-subspaces of $\mathbb F_q^N$, we prove the
permutation Molien formula and verify its invariant-subspace input by direct
prime-field matrix construction.  Three representative cases test points,
hyperplanes, and a nontrivial scalar kernel.
\end{abstract}

\subsection{General formula}
Let $X=\Gr_q(N,k)$ and $V=\mathbb C[X]$.  If $c_\ell(g)$ is the number of
$\ell$-cycles of $g$ on $X$, the cycle-block determinant identity gives
\begin{equation}\label{grassmann-proof-and-verification:eq:master}
\cS_{G,X}(\mathbf t)=\frac1{|G|}\sum_{g\in G}
\prod_{i,\ell}(1-t_i^\ell)^{-c_\ell(g)}.
\end{equation}
The monomial basis of $\Sym^d(V)$ is indexed by $d$-multisets of subspaces,
so Burnside's lemma also proves
\begin{equation}\label{grassmann-proof-and-verification:eq:orbit}
[m_\tau]\cS_{G,X}=\#\left(G\backslash\prod_j\Sym^{\tau_j}(X)\right).
\end{equation}

For a linear matrix $g$, define
\begin{equation}\label{grassmann-proof-and-verification:eq:fixed}
F_r^{\rm Gr}(g)=\#\{U\leq\mathbb F_q^N:\dim U=k,\ g^rU=U\}.
\end{equation}
Every $\ell$-cycle contributes its $\ell$ vertices to $F_r$ exactly when
$\ell\mid r$.  Möbius inversion therefore yields
\begin{equation}\label{grassmann-proof-and-verification:eq:mobius}
c_\ell^{\rm Gr}(g)=\frac1\ell\sum_{d\mid\ell}
\mu(\ell/d)F_d^{\rm Gr}(g).
\end{equation}
Together, \eqref{grassmann-proof-and-verification:eq:master}, \eqref{grassmann-proof-and-verification:eq:fixed}, and \eqref{grassmann-proof-and-verification:eq:mobius} prove the
claimed exact formula.  Interpreting the ambient space as an
$\mathbb F_q[x]$-module with $x$ acting as $g^r$ identifies
\eqref{grassmann-proof-and-verification:eq:fixed} with the stated invariant-submodule problem.

Pair orbits are classified by $\dim(U\cap W)$, and so, with
$s=\min(k,N-k)$,
\[
[m_{(1)}]=1,\qquad [m_{(2)}]=[m_{(1,1)}]=s+1.
\]

\subsection{Verification strategy}
All matrices in $GL_N(q)$ are enumerated over the prime field.  Every
$k$-subspace is generated and stored in unique reduced-row-echelon form.
Matrix multiplication then gives both the induced vertex permutation and the
independent invariance test in \eqref{grassmann-proof-and-verification:eq:fixed}.  Scalar matrices are deduplicated
before averaging, so the average is over the actual permutation image; this is
equivalent to averaging over $GL_N(q)$ because the scalar kernel has constant
size.

The target coefficient is computed from actual permutation matrices using
Newton's trace recurrence for $h_d(P)=\operatorname{tr}(\Sym^dP)$.  The same
coefficient is independently extracted from \eqref{grassmann-proof-and-verification:eq:master} and counted as
the explicit orbit number in \eqref{grassmann-proof-and-verification:eq:orbit}.  Finally, invariant-subspace
counts for powers of every matrix representative must reconstruct every cycle
through \eqref{grassmann-proof-and-verification:eq:mobius}.

\subsection{Exact results}
The coefficient vector is ordered by
$\tau=(1),(2),(1,1),(3),(2,1),(1,1,1)$.
\begin{center}
\begin{tabular}{@{}lrrrl@{}}
\toprule
Action & $|GL|$ & $|G\text{ image}|$ & $|X|$ & common vector\\
\midrule
$\Gr_2(3,1)$ & 168 & 168 & 7 & $(1,2,2,4,5,6)$\\
$\Gr_2(3,2)$ & 168 & 168 & 7 & $(1,2,2,4,5,6)$\\
$\Gr_3(2,1)$ & 48 & 24 & 4 & $(1,2,2,3,4,5)$\\
\bottomrule
\end{tabular}
\end{center}
All 18 exact triple comparisons and all invariant-subspace reconstructions
pass.  The final row explicitly checks the scalar kernel.  In the first case,
the direct determinant and cycle averages at $(0.037,0.061)$ both equal
$1.1152638116688691$ (reported error $0$).

\subsection{Reproduction}
Run
\begin{verbatim}
marimo run marimo_notebooks/grassmann_scheme.py
\end{verbatim}
from the repository root.  The implementation is intentionally restricted to
prime fields for transparency; the theorem itself is not so restricted.
\endgroup


\clearpage
\section{Polar-scheme permutation representations}\label{proof:polar-schemes}
\begingroup
\providecommand{\Sym}{\operatorname{Sym}}
\renewcommand{\Sym}{\operatorname{Sym}}
\providecommand{\cS}{\SMG}
\renewcommand{\cS}{\SMG}
\setcounter{theorem}{0}
\setcounter{proposition}{0}
\setcounter{lemma}{0}
\setcounter{corollary}{0}
\setcounter{remark}{0}
\begin{abstract}
We prove the fixed-polar-subspace version of the permutation Molien formula
and test two nontrivial actions of the explicitly constructed matrix group
$Sp_4(2)$: isotropic points and maximal totally isotropic $2$-spaces.  The
latter is a rank-two dual-polar scheme.
\end{abstract}

\subsection{Cycle and orbit formulas}
Let a finite isometry group $G$ act on a finite family $X$ of totally
isotropic or singular subspaces, and put $V=\mathbb C[X]$.  If $c_\ell(g)$
counts $\ell$-cycles on $X$, the determinant of an $\ell$-cycle block is
$1-t^\ell$, so
\begin{equation}\label{polar-proof-and-verification:eq:master}
\cS_{G,X}(\mathbf t)=\frac1{|G|}\sum_{g\in G}
\prod_{i,\ell}(1-t_i^\ell)^{-c_\ell(g)}.
\end{equation}
Since the monomial basis of $\Sym^d(V)$ is indexed by $d$-multisets in $X$,
Burnside's lemma independently gives
\begin{equation}\label{polar-proof-and-verification:eq:orbit}
[m_\tau]\cS_{G,X}=\#\left(G\backslash\prod_j\Sym^{\tau_j}(X)\right).
\end{equation}

Define the geometric fixed-point count
\begin{equation}\label{polar-proof-and-verification:eq:fixed}
F_r^{\rm pol}(g)=\#\{U\in X:g^rU=U\}.
\end{equation}
The universal relation $F_r=\sum_{\ell\mid r}\ell c_\ell$ and Möbius
inversion yield
\begin{equation}\label{polar-proof-and-verification:eq:mobius}
c_\ell^{\rm pol}(g)=\frac1\ell\sum_{d\mid\ell}
\mu(\ell/d)F_d^{\rm pol}(g).
\end{equation}
Equations \eqref{polar-proof-and-verification:eq:master}--\eqref{polar-proof-and-verification:eq:mobius} prove the proposed formula.
Relative to the Grassmann case, the only additional predicate in
\eqref{polar-proof-and-verification:eq:fixed} is total isotropy or singularity.

For maximal isotropic subspaces in Witt rank $r$, pair orbits are classified
by $r-\dim(U\cap W)$.  Consequently
\[
[m_{(1)}]=1,\qquad [m_{(2)}]=[m_{(1,1)}]=r+1.
\]

\subsection{Concrete symplectic construction}
On $\mathbb F_2^4$ use the nondegenerate alternating form
\[
\langle x,y\rangle=x_1y_2+x_2y_1+x_3y_4+x_4y_3.
\]
The code enumerates all $20{,}160$ matrices in $GL_4(2)$ and retains exactly
the 720 satisfying $\langle gx,gy\rangle=\langle x,y\rangle$ on basis
vectors.  It separately enumerates reduced-row-echelon subspaces.  There are
15 one-spaces (all isotropic for an alternating form) and 15 two-spaces on
which the form vanishes identically.  Thus both induced representations have
dimension 15, but their vertices are different geometric objects.

For every group matrix and every needed power, the program acts on the stored
bases and counts invariant members of $X$ directly.  Equation
\eqref{polar-proof-and-verification:eq:mobius} must reproduce the cycles of the separately constructed
vertex permutation.

\subsection{Three independent coefficient routes}
For the actual $15\times15$ permutation matrix $P$, complete symmetric
characters are computed without cycle data from
\[
d h_d(P)=\sum_{a=1}^d\operatorname{tr}(P^a)h_{d-a}(P),\qquad h_0=1.
\]
The matrix-group target is the average of $\prod_jh_{\tau_j}(P)$.  A second
calculation extracts the coefficient from \eqref{polar-proof-and-verification:eq:master}; a third explicitly
enumerates the configuration orbits in \eqref{polar-proof-and-verification:eq:orbit}.  Agreement of these
routes tests the representation, the proposed formula, and its geometric
interpretation separately.

\subsection{Exact results}
Values are ordered by
$\tau=(1),(2),(1,1),(3),(2,1),(1,1,1)$.
\begin{center}
\begin{tabular}{@{}lrrl@{}}
\toprule
Action & $|G|$ & $|X|$ & common coefficient vector\\
\midrule
$Sp_4(2)$ on isotropic points & 720 & 15 & $(1,3,3,8,12,16)$\\
$Sp_4(2)$ on maximal isotropic 2-spaces & 720 & 15 & $(1,3,3,8,12,16)$\\
\bottomrule
\end{tabular}
\end{center}
All 12 matrix/cycle/orbit comparisons and all geometric fixed-point
reconstructions pass.  The dual-polar pair value is $3=r+1$ for $r=2$.  For
the maximal action at $(t_1,t_2)=(0.037,0.061)$, both numerical averages equal
$1.1263214575673286$ (reported error $0$).

\subsection{Scope and reproduction}
The exact proof applies equally to alternating, Hermitian, and quadratic
forms.  The executable representative is symplectic because $Sp_4(2)$ is
large enough to exhibit all rank-two pair relations while remaining fully
enumerable and auditable.

Run
\begin{verbatim}
marimo run marimo_notebooks/polar_scheme.py
\end{verbatim}
from the repository root to inspect every comparison.
\endgroup


\clearpage
\section{Association-scheme permutation groups}\label{proof:new-dist:association_schemes}

\resultbox{\textbf{Canonical testing artifacts.}\par\smallskip Exact backend: \path{lambda_distributions/dists/association_scheme_sigma_mgfs/verification.py}.\par Regression: \path{lambda_distributions/dists/association_scheme_sigma_mgfs/test_verification.py}.\par Marimo: \path{marimo_notebooks/association_schemes.py}.}

\subsection*{Scope and outcome}

For every finite automorphism group $G\leq\Sym(\Omega)$, the imported
permutation tools show that the complete $\sigma$-MGF on $\C[\Omega]$ is determined by the
power-fixed-point data $F_r(g)=|\Fix_\Omega(g^r)|$.  We then prove explicit
fixed-point templates for homogeneous spaces, affine translation actions, and
wreath-product actions.  These specialize to Johnson, Hamming, Grassmann,
form, polar, building, graph, Doob, folded-cube, and halved-cube families.
Stable ranges and degree-two obstructions are kept separate from exact finite
formulas.  An accompanying marimo notebook constructs representative matrix
groups and compares direct matrix traces, cycle products, and literal
configuration orbits in 99 exact coefficient checks.  Every check passes.

\subsection{The universal power-fixed-point theorem}

Let a finite group $G$ act on a finite set $\Omega$, let $V=\C[\Omega]$, and
let $P_g$ denote the resulting permutation matrix.  For a composition or
partition $\tau=(\tau_1,\ldots,\tau_s)$ put
\[
 W_\tau(V)=\bigotimes_{j=1}^s\Sym^{\tau_j}V,
 \qquad a_\tau(G,\Omega)=\dim W_\tau(V)^G.
\]
Our convention is
\[
 \MGF_{G,\Omega}(\mathbf t)
 =\E_{g\in G}\ExpS(X_V(g)h_1)
 =\sum_\tau a_\tau(G,\Omega)m_\tau(\mathbf t).
\]
For $g\in G$ define
\[
 F_r(g)=|\Fix_\Omega(g^r)|,
 \qquad c_d(g)=\#\{\text{$d$-cycles of $g$ on $\Omega$}\}.
\]

\begin{theorem}[Association-scheme power-fixed-point theorem]
For every finite permutation action,
\begin{align}
 F_r(g)&=\sum_{d\mid r}d\,c_d(g),\label{nd:association_schemes:eq:fixed-from-cycles}\\
 c_d(g)&=\frac1d\sum_{r\mid d}\mu(d/r)F_r(g),\label{nd:association_schemes:eq:mobius}\\
 \MGF_{G,\Omega}(\mathbf t)
 &=\frac1{|G|}\sum_{g\in G}
   \prod_{i\geq1}\prod_{d\geq1}(1-t_i^d)^{-c_d(g)},\label{nd:association_schemes:eq:universal}\\
 [m_\tau]\MGF_{G,\Omega}
 &=\#\left(G\backslash\prod_j\Sym^{\tau_j}\Omega\right).
 \label{nd:association_schemes:eq:orbits}
\end{align}
Thus the complete $\sigma$-MGF is determined by the collection of numbers
$F_r(g)$.
\end{theorem}

\begin{proof}[Imported identities and the fixed-point reduction]
The sigma-MGF and orbit formulas are Tools~\ref{tool:molien} and
\ref{tool:permutation-orbits}.  A point on a $d$-cycle is fixed by $g^r$
exactly when $d\mid r$, proving
\eqref{nd:association_schemes:eq:fixed-from-cycles}; arithmetic M\"obius
inversion gives \eqref{nd:association_schemes:eq:mobius}.  Substituting the
recovered cycle counts into the imported permutation formula proves
\eqref{nd:association_schemes:eq:universal}.  No Reynolds--Molien or orbit
identity is reproved.
\end{proof}

\begin{remark}[What is tested directly]
The verification stores $P_g$ exactly by the location of the single nonzero
entry in each column.  It computes symmetric-power characters from the actual
matrix traces using Newton's recurrence
\[
 d\,h_d(P_g)=\sum_{r=1}^d\operatorname{Tr}(P_g^r)h_{d-r}(P_g),\qquad h_0=1.
\]
This matrix route, the determinant cycle-factor route, and literal
configuration orbits are separate computations of the target coefficient.
\end{remark}

\subsection{Three reusable fixed-point templates}

\subsubsection{Homogeneous and parabolic actions}

Let $\Omega=G/H$ with left multiplication.

\begin{proposition}[Coset fixed points]\label{nd:association_schemes:prop:coset}
For $x\in G$,
\begin{equation}\label{nd:association_schemes:eq:coset}
 \boxed{|\Fix_{G/H}(x)|
 =\frac{|C_G(x)|\,|x^G\cap H|}{|H|}.}
\end{equation}
\end{proposition}

\begin{proof}
The coset $yH$ is fixed precisely when $y^{-1}xy\in H$.  For each element of
$x^G\cap H$, exactly $|C_G(x)|$ choices of $y$ conjugate $x$ to it.  Dividing
the number of such $y$ by $|H|$ accounts for the representatives of a coset.
Apply this to $x=g^r$, then use \eqref{nd:association_schemes:eq:mobius}--\eqref{nd:association_schemes:eq:universal}.
\end{proof}

This covers Grassmannians and other parabolic quotients, dual-polar spaces,
generalized polygons from $BN$-pairs, and building residues.

\subsubsection{Affine translation actions}

Let $A$ be a finite additive group, $H\leq\Aut(A)$, and
$G=A\rtimes H$ act by $x\mapsto b+h(x)$.  Define
\[
 b_r=(I+h+\cdots+h^{r-1})b.
\]

\begin{proposition}[Affine kernel/image formula]\label{nd:association_schemes:prop:affine}
For $g=(b,h)$,
\begin{equation}\label{nd:association_schemes:eq:affine}
 \boxed{F_r(b,h)=
 \begin{cases}
 |\ker(I-h^r)|,&b_r\in\operatorname{im}(I-h^r),\\
 0,&b_r\notin\operatorname{im}(I-h^r).
 \end{cases}}
\end{equation}
\end{proposition}

\begin{proof}
Induction gives $g^r(x)=b_r+h^r(x)$.  Thus $x$ is fixed exactly when
$(I-h^r)x=b_r$.  A fiber of the homomorphism $I-h^r$ is empty outside its
image and is a coset of its kernel otherwise.
\end{proof}

\subsubsection{Product and wreath actions}

Let $H\curvearrowright\Gamma$ and let $H\wr S_d=H^d\rtimes S_d$ act on
$\Gamma^d$.  Write $g=(h_1,\ldots,h_d;\pi)$.  For a cycle
$C=(i_1\ldots i_\ell)$ of $\pi$, let
\[
 a_C=h_{i_\ell}\cdots h_{i_1},\qquad e_C(r)=\gcd(\ell,r),
 \qquad f_s(a)=|\Fix_\Gamma(a^s)|.
\]

\begin{proposition}[Marked-cycle product formula]\label{nd:association_schemes:prop:wreath}
\begin{equation}\label{nd:association_schemes:eq:wreath}
 \boxed{F_r(g)=\prod_{C\in\operatorname{Cyc}(\pi)}
 f_{r/e_C(r)}(a_C)^{e_C(r)}.}
\end{equation}
\end{proposition}

\begin{proof}
Under $\pi^r$, an $\ell$-cycle splits into $e=\gcd(\ell,r)$ coordinate
cycles.  Transport once around any one has monodromy conjugate to
$a_C^{r/e}$.  Each of the $e$ cycles admits $f_{r/e}(a_C)$ choices.  Different
cycles of $\pi$ use disjoint coordinates, so their counts multiply.
\end{proof}

\subsection{Johnson schemes}

Let $\Omega=\binom{[v]}k$ and $G=S_v$.  Suppose $\sigma$ has cycle type
$\lambda=1^{m_1}2^{m_2}\cdots$.  A $j$-cycle of $\sigma$ splits under
$\sigma^r$ into $e=\gcd(j,r)$ cycles of length $j/e$.  A fixed $k$-subset is
a union of these cycles, hence
\begin{equation}\label{nd:association_schemes:eq:johnson-fixed}
 \boxed{F_r^{J(v,k)}(\lambda)=[u^k]\prod_{j\geq1}
 \left(1+u^{j/\gcd(j,r)}\right)^{m_j\gcd(j,r)}.}
\end{equation}
With $c_d^J(\lambda)=d^{-1}\sum_{r\mid d}\mu(d/r)F_r^J(\lambda)$ and
$z_\lambda=\prod_jj^{m_j}m_j!$,
\begin{equation}\label{nd:association_schemes:eq:johnson-smg}
 \boxed{\MGF_{J(v,k)}=\sum_{\lambda\vdash v}\frac1{z_\lambda}
 \prod_{i,d\geq1}(1-t_i^d)^{-c_d^J(\lambda)}.}
\end{equation}

For $D=|\tau|$, a configuration contributing to $[m_\tau]$ uses at most
$kD$ underlying points.  Relabelling those points proves
\begin{equation}\label{nd:association_schemes:eq:johnson-stable}
 [m_\tau]\MGF_{J(v,k)}\text{ is constant for }v\geq kD.
\end{equation}
The stable objects are isomorphism classes of $\tau$-colored $k$-uniform
multihypergraphs without isolated vertices.  Ordered pairs are classified by
intersection size, so
\begin{equation}\label{nd:association_schemes:eq:johnson-pair}
 \boxed{[m_{(1,1)}]\MGF_{J(v,k)}=\min(k,v-k)+1.}
\end{equation}
Thus no coefficientwise limit exists when $\min(k,v-k)\to\infty$.

\subsection{Hamming and Doob product schemes}

For $\Omega=[q]^d$ and $G=S_q\wr S_d$, Proposition~\ref{nd:association_schemes:prop:wreath} gives
\begin{equation}\label{nd:association_schemes:eq:hamming-fixed}
 F_r(g)=\prod_C\left|\Fix_{[q]}
 \left(a_C^{r/e_C(r)}\right)\right|^{e_C(r)}.
\end{equation}
Consequently
\begin{equation}\label{nd:association_schemes:eq:hamming-smg}
 \boxed{\MGF_{H(d,q)}=\frac1{(q!)^dd!}\sum_{\pi\in S_d}
 \sum_{h_1,\ldots,h_d\in S_q}\prod_{i,\ell\geq1}
 (1-t_i^\ell)^{-\frac1\ell\sum_{r\mid\ell}\mu(\ell/r)F_r(g)}.}
\end{equation}

For fixed $d$ and $D=|\tau|$, once $q\geq D$ the orbit of the $D$ symbols in
each coordinate is exactly its equality partition in the partition lattice
$\Pi_D$.  If $K_\tau=\prod_jS_{\tau_j}$ permutes repeated objects within the
colored multiset blocks, then
\begin{equation}\label{nd:association_schemes:eq:hamming-stable}
 \boxed{[m_\tau]\MGF_{H(d,q)}
 =\#\bigl((K_\tau\times S_d)\backslash\Pi_D^d\bigr),\qquad q\geq D.}
\end{equation}
On the other hand, ordered word pairs are classified by Hamming distance:
\begin{equation}\label{nd:association_schemes:eq:hamming-pair}
 \boxed{[m_{(1,1)}]\MGF_{H(d,q)}=d+1.}
\end{equation}

For the Doob graph $D(m,n)=\mathrm{Sh}^{\square m}\square K_4^{\square n}$,
take
\[
 G_{m,n}=(\Aut(\mathrm{Sh})\wr S_m)\times(S_4\wr S_n).
\]
Fixed points multiply between the two Cartesian blocks.  The Shrikhande
vertex action has four ordered-pair types and $S_4\curvearrowright[4]$ has
two.  Permuting equal factors leaves multisets of types, so
\begin{equation}\label{nd:association_schemes:eq:doob-pair}
 \boxed{[m_{(1,1)}]\MGF_{D(m,n)}=\binom{m+3}{3}(n+1).}
\end{equation}
In contrast, $H(2m+n,4)$ has pair coefficient $2m+n+1$ even though the two
families have the same distance-regular parameters.

\subsection{Grassmann, polar, polygon, and building actions}

\subsubsection{Grassmann schemes}

Let $\Omega=\Gr_k(\F_q^n)$ and take the projective image of $\GL_n(q)$.  Directly,
\begin{equation}\label{nd:association_schemes:eq:grass-fixed}
 F_r(g)=\#\{U\leq\F_q^n:\dim U=k,\ g^rU=U\}.
\end{equation}
Equivalently, with $P=P_{k,n-k}$, Proposition~\ref{nd:association_schemes:prop:coset} gives
\[
 F_r(g)=\frac{|C_G(g^r)|\,|(g^r)^G\cap P|}{|P|}.
\]
Invariant-subspace counts may be computed from the rational canonical form
of $g^r$.  Ordered pairs are classified by $\dim(U\cap W)$:
\begin{equation}\label{nd:association_schemes:eq:grass-pair}
 \boxed{[m_{(1,1)}]\MGF_{\Gr_k(n,q)}=\min(k,n-k)+1.}
\end{equation}
The span of $D$ participating $k$-spaces has dimension at most $kD$, and an
isomorphism between supported configurations extends to the ambient space.
Thus
\begin{equation}\label{nd:association_schemes:eq:grass-stable}
 [m_\tau]\MGF_{\Gr_k(n,q)}\text{ stabilizes for }n\geq kD
 \quad(q,k,\tau\text{ fixed}).
\end{equation}
Fixed-rank field-size limits can fail: four points on a projective line carry
a cross-ratio parameter.

\subsubsection{Dual-polar spaces and generalized polygons}

Let $V$ be a finite formed space of Witt rank $d$ and let $\Omega=G/P$ be the
maximal isotropic or singular subspaces.  Then
\begin{equation}\label{nd:association_schemes:eq:dualpolar-fixed}
 F_r(g)=\frac{|C_G(g^r)|\,|(g^r)^G\cap P|}{|P|}.
\end{equation}
Pair orbits are indexed by intersection codimension, hence
\begin{equation}\label{nd:association_schemes:eq:dualpolar-pair}
 \boxed{[m_{(1,1)}]\MGF_{\mathrm{DP}(d,q)}=d+1.}
\end{equation}

The same coset formula applies to a flag-transitive generalized polygon with
point set $G/P$.  If the point action is distance-transitive, its pair orbits
are indexed by distance.  Therefore the pair coefficients for generalized
quadrangles, hexagons, and octagons are respectively $3,4,5$.  Without
flag/distance transitivity, the order parameters alone do not imply these
values and do not determine the full series.

\subsubsection{Buildings}

Suppose $G$ has a $BN$-pair with Weyl group $W$.  For chambers $G/B$, Bruhat
decomposition gives $B\backslash G/B\cong W$, so
\begin{equation}\label{nd:association_schemes:eq:building-chamber}
 \boxed{[m_{(1,1)}]\MGF_{G,G/B}=|W|.}
\end{equation}
For residues of type $J$,
\begin{equation}\label{nd:association_schemes:eq:building-residue}
 \boxed{[m_{(1,1)}]\MGF_{G,G/P_J}=|W_J\backslash W/W_J|.}
\end{equation}
These conclusions require the natural strong or Weyl transitivity.  Higher
coefficients need not be controlled by $W$ because triples of flags can carry
field-valued moduli.

\subsection{Translation schemes of forms}

All families in this section use Proposition~\ref{nd:association_schemes:prop:affine}.  The target
representation is the complex permutation module on the additive form space;
the finite-field matrices are used to construct its zero-one permutation
matrices.

\subsubsection{Bilinear forms}

Let $A=M_{m\times n}(\F_q)$ and let $(P,Q)$ act linearly by
$L(X)=PXQ^{-1}$.  For $g=(B,P,Q)$ define
$B_r=(I+L+\cdots+L^{r-1})B$.  Then
\begin{equation}\label{nd:association_schemes:eq:bilinear-fixed}
 \boxed{F_r(g)=
 \begin{cases}
 q^{\dim\ker(I-L^r)},&B_r\in\operatorname{im}(I-L^r),\\
 0,&\text{otherwise}.
 \end{cases}}
\end{equation}
After translating one member of a pair to zero, differences are classified by
rank:
\begin{equation}\label{nd:association_schemes:eq:bilinear-pair}
 \boxed{[m_{(1,1)}]\MGF_{\mathrm{Bil}(m,n,q)}=\min(m,n)+1.}
\end{equation}
For fixed $m,q,\tau$, translating one of $D=|\tau|$ matrices to zero leaves
combined row support of dimension at most $m(D-1)$.  Thus the coefficient is
stable once $n\geq m(D-1)$.

\subsubsection{Alternating, Hermitian, and quadratic forms}

For $A=\operatorname{Alt}_n(\F_q)$ use $L(X)=PXP^T$; for Hermitian matrices
over $\F_{q^2}$ use $L(X)=PXP^{\sigma T}$.  Formula \eqref{nd:association_schemes:eq:bilinear-fixed}
applies verbatim with the appropriate linear operator.  Rank classifies the
one-form orbits in these two cases, while alternating ranks are even.  Hence
\begin{align}
 [m_{(1,1)}]\MGF_{\mathrm{Alt}(n,q)}&=\left\lfloor\frac n2\right\rfloor+1,
 \label{nd:association_schemes:eq:alt-pair}\\
 [m_{(1,1)}]\MGF_{\mathrm{Herm}(n,q)}&=n+1.
 \label{nd:association_schemes:eq:herm-pair}
\end{align}

Let $A=\mathcal Q(n,q)$ be the additive space of quadratic forms, with
$P\cdot Q=Q\circ P^{-1}$.  Again \eqref{nd:association_schemes:eq:affine} is exact.  For odd $q$,
each positive rank has two congruence types and rank zero has one, yielding
\begin{equation}\label{nd:association_schemes:eq:quadratic-pair}
 \boxed{[m_{(1,1)}]\MGF_{\mathcal Q(n,q)}=2n+1\qquad(q\text{ odd}).}
\end{equation}
Even characteristic requires the usual rank/type refinement but no change to
the affine fixed-point formula.  The displayed pair coefficients already rule
out rank-growing coefficientwise limits.  Field-size limits can also fail in
higher degree: after translating a triple of bilinear forms and making one
difference invertible, the other retains similarity invariants such as
eigenvalues.

\subsection{Affine-polar, folded-cube, and halved-cube actions}

For an affine-polar action take the additive formed space $A=V$ and the
relevant isometry, similitude, or semisimilitude group $H$.  Equation
\eqref{nd:association_schemes:eq:affine} is the complete exact answer.  For $D=|\tau|$, translate one
of the $D$ points to zero.  The other $D-1$ differences span a formed subspace
of dimension at most $D-1$.  Its restricted form and labelled vectors
determine the orbit; Witt extension embeds equivalent supported configurations
into all sufficiently large ambient spaces.  Therefore
\begin{equation}\label{nd:association_schemes:eq:affpolar-stable}
 [m_\tau]\MGF_{\mathrm{AffPol}(N,q)}\text{ is eventually constant as }N\to\infty
 \quad(q,\tau\text{ fixed}).
\end{equation}

For the folded cube, let
\[
 A_n=\F_2^n/\langle\mathbf1\rangle,\qquad G=A_n\rtimes S_n.
\]
For $g=(b,\sigma)$, equation \eqref{nd:association_schemes:eq:affine} becomes
\begin{equation}\label{nd:association_schemes:eq:folded-fixed}
 F_r(g)=\begin{cases}
 2^{\dim\ker(I-\sigma^r\mid A_n)},&b_r\in\operatorname{im}(I-\sigma^r),\\
 0,&\text{otherwise}.
 \end{cases}
\end{equation}
Weights $w$ and $n-w$ represent the same difference orbit, so
\begin{equation}\label{nd:association_schemes:eq:folded-pair}
 \boxed{[m_{(1,1)}]\MGF_{\mathrm{FoldCube}_n}=\left\lfloor\frac n2\right\rfloor+1.}
\end{equation}

For the halved cube use
$E_n=\{x\in\F_2^n:\sum_i x_i=0\}$ and $G=E_n\rtimes S_n$.
Formula \eqref{nd:association_schemes:eq:folded-fixed} holds with $A_n$ replaced by $E_n$.
Difference weights are the even integers from zero to the largest even value
at most $n$, giving
\begin{equation}\label{nd:association_schemes:eq:halved-pair}
 \boxed{[m_{(1,1)}]\MGF_{\frac12Q_n}=\left\lfloor\frac n2\right\rfloor+1.}
\end{equation}

\subsection{Distance-regular and strongly regular graphs}

For any finite graph $\Gamma$, equations
\eqref{nd:association_schemes:eq:mobius}--\eqref{nd:association_schemes:eq:universal} apply with
$G=\Aut(\Gamma)$ and $\Omega=V(\Gamma)$.  If the action is
distance-transitive of diameter $D$, then
\begin{equation}\label{nd:association_schemes:eq:distance-pair}
 [m_{(1,1)}]\MGF_\Gamma=D+1.
\end{equation}
More generally, the pair coefficient is the permutation rank of the vertex
action.

The $4\times4$ rook graph and the Shrikhande graph both have strongly regular
parameters $(16,6,2,2)$, but their full vertex actions have ranks three and
four.  In the Cayley model on $\mathbb Z_4^2$, the Shrikhande vertex
stabilizer has orbits
\[
 \{0\},\quad \{\text{six neighbors}\},\quad
 \{(2,0),(0,2),(2,2)\},\quad \{\text{remaining six vertices}\}.
\]
Therefore the $\sigma$-MGF is not determined by the intersection array,
spectrum, or strongly regular parameters.  At the other extreme, a graph with
trivial automorphism group on $N$ vertices has
\[
 \MGF_\Gamma=\prod_{i\geq1}(1-t_i)^{-N}.
\]

\subsection{Asymptotic classification}

Let $D=|\tau|$. The proof mechanism behind the positive stability results is
bounded support together with an extension theorem: a fixed-degree
configuration lives in a bounded subobject and supported isomorphisms extend
to ambient automorphisms. Most rank-growth failures below are already forced
by an unbounded degree-two coefficient. The fixed-rank Grassmann
$q\to\infty$ regime is different: its pair coefficient is constant, while
cross-ratio moduli can first obstruct convergence in higher degree.

\begin{center}
\small
\renewcommand{\arraystretch}{1.12}
\begin{tabularx}{\linewidth}{@{}>{\raggedright\arraybackslash}p{0.47\linewidth}X@{}}
\toprule
Family and regime & Coefficientwise behavior \\
\midrule
$J(v,k)$, fixed $k$, $v\to\infty$ & Stable for $v\geq kD$ \\
$J(v,k)$, $\min(k,v-k)\to\infty$ & No limit; pair coefficient diverges \\
$H(d,q)$, fixed $d$, $q\to\infty$ & Stable for $q\geq D$ \\
$H(d,q)$, fixed $q$, $d\to\infty$ & No limit; pair coefficient $d+1$ \\
$\Gr_k(n,q)$, fixed $q,k$, $n\to\infty$ & Stable for $n\geq kD$ \\
Grassmann, fixed rank, $q\to\infty$ & Can fail through cross-ratio moduli \\
Bilinear, fixed $m,q$, $n\to\infty$ & Stable for $n\geq m(D-1)$ \\
Alternating/Hermitian/quadratic, rank $\to\infty$ & No limit in degree two \\
Dual-polar, Witt rank $\to\infty$ & No limit; pair coefficient $d+1$ \\
Affine-polar, fixed $q$, rank $\to\infty$ & Eventually stable for fixed $\tau$ \\
Building chambers, rank $\to\infty$ & No limit when $|W|\to\infty$ \\
Bounded residues in classical towers, fixed $q$ & Stable by bounded support \\
Doob/folded/halved dimension $\to\infty$ & No limit in degree two \\
\bottomrule
\end{tabularx}
\end{center}

\subsection{Exact computational verification}

\subsubsection{Construction and comparison strategy}

The companion module \texttt{verification.py} uses only exact finite
arithmetic for the structural checks.
\begin{enumerate}
\item Enumerate the stated finite matrix group (or a faithful permutation
image) and the finite set $\Omega$.
\item Construct $P_g$ by applying each group element to every point of
$\Omega$.  The dimension is exactly $|\Omega|$.
\item Compute $F_r(g)$ directly from $P_g^r$ and independently from the
appropriate coset, affine, invariant-subspace, or wreath formula.
\item Recover $c_d(g)$ by M\"obius inversion and compare with literal cycle
decomposition.\par
\item Compute $[m_\tau]\MGF$ from Newton matrix traces, from cycle-factor
coefficient extraction, and from direct orbits on
$\prod_j\Sym^{\tau_j}\Omega$.
\end{enumerate}

The marimo notebook presents these checks by group family.  Existing Johnson,
Hamming, Grassmann, and polar constructors are imported from the canonical
association-scheme package in this workspace; the new notebook adds the
affine-form, building, graph, Doob, folded, and halved constructors.  No
combined anthology PDF is modified.

\subsubsection{Representative exact cases}

\begin{center}
\small
\renewcommand{\arraystretch}{1.1}
\begin{longtable}{@{}p{0.38\linewidth}rrrrl@{}}
\toprule
Action & $|\Omega|$ & $|H|$ & $|G|$ & $[m_{11}]$ & Audit \\
\midrule
\endfirsthead
\toprule
Action & $|\Omega|$ & $|H|$ & $|G|$ & $[m_{11}]$ & Audit \\
\midrule
\endhead
$S_v$ on $k$-subsets: $(4,2),(5,2),(5,1)$ & $6,10,5$ & -- & $48,120,120$ & $3,3,2$ & exact \\
$S_q\wr S_d$: $(2,2),(2,3),(3,2)$ & $4,8,9$ & -- & $8,48,72$ & $3,4,3$ & exact \\
$PGL_3(2)$ on $\Gr_1,\Gr_2$; $PGL_2(3)$ on $\Gr_1$ & $7,7,4$ & -- & $168,168,24$ & $2,2,2$ & exact \\
$Sp_4(2)$ on polar points / Lagrangians & $15$ & -- & $720$ & $3$ & exact \\
$GL_3(2)$ on complete flags & $21$ & $|B|=8$ & $168$ & $6$ & coset \\
\midrule
$M_2(\F_2)\rtimes(GL_2(2)\times GL_2(2))$ & $16$ & $36$ & $576$ & $3$ & affine \\
$\operatorname{Alt}_3(2)\rtimes GL_3(2)$ & $8$ & $168$ & $1344$ & $2$ & affine \\
$\operatorname{Herm}_2(4/2)\rtimes GL_2(4)$ image & $16$ & $60$ & $960$ & $3$ & affine \\
$\mathcal Q(2,3)\rtimes GL_2(3)$ image & $27$ & $24$ & $648$ & $5$ & affine \\
$\F_2^4\rtimes O_4^+(2)$ & $16$ & $72$ & $1152$ & $3$ & affine polar \\
Folded and halved $4$-cubes & $8$ & $24$ & $192$ & $3$ & affine \\
\midrule
Shrikhande full vertex action & $16$ & $12$ & $192$ & $4$ & graph \\
$4\times4$ rook graph & $16$ & -- & $1152$ & $3$ & graph \\
Doob $D(1,1)$ chosen product action & $64$ & -- & $4608$ & $8$ & product \\
\bottomrule
\end{longtable}
\end{center}

For the supplemental affine and graph cases, the reported low-degree values
are:
\[
\begin{array}{c|cccccccccc}
\text{case}&\mathrm{Bil}&\mathrm{Alt}&\mathrm{Herm}&\mathcal Q&\mathrm{AffPol}
&\mathrm{Fold}&\tfrac12Q&\mathrm{Sh}&\mathrm{Rook}&D(1,1)\\ \hline
[m_{11}]&3&2&3&5&3&3&3&4&3&8
\end{array}
\]
In every entry, the matrix-trace average, cycle formula, and direct orbit count
agree.  The quadratic example also illustrates why $[m_{(2)}]$ need not equal
$[m_{(1,1)}]$: its values are $4$ and $5$, respectively, because forgetting
the two colors identifies a difference with its negative.

\resultbox{\textbf{Verification outcome.}  The suite completes 99 exact
coefficient comparisons.  Every displayed coefficient agrees across all
implemented routes; every family-specific power-fixed-point formula
reconstructs the literal cycle counts; the $GL_3(2)/B$ example verifies the
centralizer/class-intersection formula element by element; and the Doob
$D(1,1)$ example verifies the product fixed-point identity for powers
$1\leq r\leq6$.}

\subsection{Scope and cautions}

\begin{itemize}
\item The universal theorem is valid for every finite action.  An association
scheme is useful structure for calculating $F_r(g)$; it is not an additional
hypothesis needed for the theorem.
\item The automorphism group must be specified.  Different subgroups of
$\Aut(\Omega)$ generally give different averages and different orbit counts.
\item Parameters such as an intersection array or strongly regular tuple do
not determine the automorphism action, as the rook/Shrikhande comparison
shows.
\item Semilinear, complement, or exceptional automorphisms are incorporated
by averaging over their extra cosets.  Omitting such a coset computes the
series for the smaller stated group, not for the full automorphism group.
\item Exact finite formulas, stable formulas, and coefficientwise limits are
distinct claims.  This note states a stable range only where bounded support
and extension give one directly.
\end{itemize}

\clearpage
\section{The hyperoctahedral group on the binary cube}\label{proof:hyperoctahedral-cube}
\begingroup
\providecommand{\F}{\mathbb F}
\renewcommand{\F}{\mathbb F}
\providecommand{\C}{\mathbb C}
\renewcommand{\C}{\mathbb C}
\providecommand{\E}{\mathbb E}
\renewcommand{\E}{\mathbb E}
\providecommand{\Prob}{\mathbb P}
\renewcommand{\Prob}{\mathbb P}
\providecommand{\Fix}{\operatorname{Fix}}
\renewcommand{\Fix}{\operatorname{Fix}}
\providecommand{\Tr}{\operatorname{Tr}}
\renewcommand{\Tr}{\operatorname{Tr}}
\providecommand{\Sym}{\operatorname{Sym}}
\renewcommand{\Sym}{\operatorname{Sym}}
\providecommand{\Exp}{\operatorname{Exp}}
\renewcommand{\Exp}{\operatorname{Exp}}
\providecommand{\stirling}[2]{\genfrac{[}{]}{0pt}{}{#1}{#2}}
\renewcommand{\stirling}[2]{\genfrac{[}{]}{0pt}{}{#1}{#2}}
\setcounter{theorem}{0}
\setcounter{proposition}{0}
\setcounter{lemma}{0}
\setcounter{corollary}{0}
\setcounter{remark}{0}
\begin{abstract}
Let $B_n=C_2^n\rtimes S_n$ act affinely on $X_n=\F_2^n$ and let
$V_n=\C[X_n]$.  We prove the proposed fixed-point distribution and moment
formulas, derive the exact signed-cycle-type expression for the full
$\sigma$-MGF, and describe an independent exhaustive verification.  The
program constructs every $2^n\times2^n$ permutation matrix, computes the
target directly, traverses the relevant monomial-basis orbits, and compares
both results with the signed-cycle formula for $1\leq n\leq5$.  All checks
pass; the largest case has $|B_5|=3840$ matrices of size $32\times32$.
\end{abstract}

\subsection{The action and the imported target}
Write elements of $B_n$ as $g=(b,\pi)$, where $b\in\F_2^n$ and
$\pi\in S_n$, and set
\[
g\cdot x=b+\pi x,\qquad x\in X_n.
\]
On the basis $\{e_x:x\in X_n\}$ of $V_n$, the corresponding matrix is
\begin{equation}\label{hyperoctahedral-cube-proof-and-verification:eq:matrix}
P_g e_x=e_{g\cdot x}.
\end{equation}
Thus $P_g$ is a $2^n\times2^n$ permutation matrix.  Tools~\ref{tool:molien}
and~\ref{tool:permutation-orbits} specialize to
\begin{align}
\smgv{B_n}{V_n}
&=\E_{g\in B_n}\prod_{d\geq1}Q_d(\mathbf t)^{c_d(g)},
\qquad Q_d(\mathbf t)=\prod_{a\geq1}(1-t_a^d)^{-1},
\label{hyperoctahedral-cube-proof-and-verification:eq:molien}\\
[m_\tau]\,\smgv{B_n}{V_n}
&=\#\left(B_n\backslash\prod_{j=1}^s\Sym^{\tau_j}(X_n)\right).
\label{hyperoctahedral-cube-proof-and-verification:eq:orbit-coeff}
\end{align}
The new argument is the conditional fixed-point law of the affine cube action.

\subsection{The fixed-point law}
Define
\[
F_n(g)=\Tr(P_g)=\#\Fix_{X_n}(g).
\]

\begin{theorem}[Exact conditional law]\label{hyperoctahedral-cube-proof-and-verification:thm:fixed-law}
For uniform $b\in\F_2^n$ and fixed $\pi\in S_n$, let $C(\pi)$ be the number
of coordinate cycles of $\pi$.  Then
\[
F_n(b,\pi)=
\begin{cases}
2^{C(\pi)},&\text{with probability }2^{-C(\pi)},\\
0,&\text{with probability }1-2^{-C(\pi)}.
\end{cases}
\]
\end{theorem}

\begin{proof}
On a coordinate cycle $(i_1,\ldots,i_\ell)$, the fixed equation
$x=b+\pi x$ becomes a cyclic system
\[
x_{i_1}=b_{i_1}+x_{i_\ell},\quad
x_{i_2}=b_{i_2}+x_{i_1},\quad\ldots\quad
x_{i_\ell}=b_{i_\ell}+x_{i_{\ell-1}}.
\]
Adding in $\F_2$ shows that it is solvable only if
$\sum_{j=1}^{\ell}b_{i_j}=0$.  This condition is also sufficient: after one
coordinate is chosen, the remaining coordinates are forced.  Consequently
each coordinate cycle imposes one independent parity condition on $b$ and,
when solvable, contributes two solutions.  There are $C(\pi)$ independent
conditions and $2^{C(\pi)}$ solutions.
\end{proof}

Let $\stirling{n}{k}$ be the unsigned Stirling number of the first kind.  Since
there are $\stirling{n}{k}$ permutations with $k$ cycles, Theorem
\ref{hyperoctahedral-cube-proof-and-verification:thm:fixed-law} yields
\begin{equation}\label{hyperoctahedral-cube-proof-and-verification:eq:distribution}
\Prob(F_n=2^k)=\frac{\stirling{n}{k}}{n!\,2^k},
\qquad 1\leq k\leq n.
\end{equation}
The remaining probability is at zero.  Using
$\sum_k\stirling{n}{k}x^k=x(x+1)\cdots(x+n-1)$ gives
\begin{equation}\label{hyperoctahedral-cube-proof-and-verification:eq:positive-probability}
\Prob(F_n>0)
=\frac{(1/2)^{\overline n}}{n!}
=\frac{\Gamma(n+\tfrac12)}{\sqrt\pi\,\Gamma(n+1)}
=\frac{\binom{2n}{n}}{4^n}
\sim\frac1{\sqrt{\pi n}}.
\end{equation}
Hence $F_n\to0$ in probability.

For every integer $m\geq1$, conditioning on $\pi$ instead gives
\[
\E_b[F_n^m\mid\pi]
=2^{-C(\pi)}2^{mC(\pi)}=2^{(m-1)C(\pi)}.
\]
Another application of the Stirling generating polynomial proves
\begin{equation}\label{hyperoctahedral-cube-proof-and-verification:eq:moments}
\boxed{\quad
\E[F_n^m]
=\frac{(2^{m-1})^{\overline n}}{n!}
=\binom{n+2^{m-1}-1}{n}.
\quad}
\end{equation}
In particular $\E F_n=1$, $\E F_n^2=n+1$,
$\E F_n^3=\binom{n+3}{3}$, and $\operatorname{Var}(F_n)=n$.  Thus convergence
in probability to zero coexists with divergent moments of every order
$m\geq2$.

There is also a direct Burnside proof of \eqref{hyperoctahedral-cube-proof-and-verification:eq:moments}.  Represent an
ordered $m$-tuple of cube vertices by an $m\times n$ binary matrix.  Coordinate
permutations permute its columns, while a flip in one cube coordinate adds
$\mathbf1=(1,\ldots,1)$ to the corresponding column.  Each column is therefore
remembered only as an element of
$\F_2^m/\langle\mathbf1\rangle$, a set of size $2^{m-1}$.  The orbit is a
multiset of $n$ such column types, so
\[
\#(B_n\backslash X_n^m)=\binom{n+2^{m-1}-1}{n}=\E[F_n^m].
\]
Combining this with \eqref{hyperoctahedral-cube-proof-and-verification:eq:orbit-coeff} gives
\begin{equation}\label{hyperoctahedral-cube-proof-and-verification:eq:moment-coeff}
[m_{(1^m)}]\mathcal H_n=\binom{n+2^{m-1}-1}{n}.
\end{equation}

\subsection{Signed cycle types and the full cycle index}
A signed coordinate-cycle type is a pair of sequences
$(a_\ell,b_\ell)_{\ell\geq1}$, where $a_\ell$ counts length-$\ell$
coordinate cycles with total flip parity zero and $b_\ell$ counts those with
total flip parity one.  They satisfy
\[
\sum_{\ell\geq1}\ell(a_\ell+b_\ell)=n.
\]
For a uniform element of $B_n$, the probability of this type is
\begin{equation}\label{hyperoctahedral-cube-proof-and-verification:eq:type-probability}
w(a,b)=\prod_{\ell\geq1}
\frac1{(2\ell)^{a_\ell+b_\ell}a_\ell!b_\ell!}.
\end{equation}
This follows from the usual $S_n$ cycle-type probability and the fact that
the accumulated parity on every coordinate cycle is an independent fair bit.

Let
\[
R_r(a,b)=\#\Fix_{X_n}(g^r)
\]
for any element of signed type $(a,b)$.  A positive length-$\ell$ cycle
contributes $2^{\gcd(\ell,r)}$ fixed choices.  For a negative cycle, write the
affine map as $T(x)=Rx+e_0$, with $R$ a cyclic rotation.  If
$d=\gcd(\ell,r)$, the permutation $R^r$ has $d$ coordinate orbits.  On each
orbit, the translation in $T^r$ has parity $r/d$.  The fixed equation is
therefore solvable exactly when $r/d$ is even; when solvable it has $2^d$
solutions.  Independence between coordinate cycles proves
\begin{equation}\label{hyperoctahedral-cube-proof-and-verification:eq:fixed-powers}
R_r(a,b)=
\begin{cases}
\displaystyle
2^{\sum_{\ell\geq1}(a_\ell+b_\ell)\gcd(\ell,r)},
&\displaystyle
\frac r{\gcd(\ell,r)}\text{ is even for every }\ell\text{ with }b_\ell>0,
\\[8pt]
0,&\text{otherwise}.
\end{cases}
\end{equation}

Because $R_r=\sum_{d\mid r}d\,c_d$, Möbius inversion recovers the exact
vertex-cycle counts:
\begin{equation}\label{hyperoctahedral-cube-proof-and-verification:eq:mobius}
c_d(a,b)=\frac1d\sum_{r\mid d}\mu(d/r)R_r(a,b).
\end{equation}
Substituting \eqref{hyperoctahedral-cube-proof-and-verification:eq:type-probability}--\eqref{hyperoctahedral-cube-proof-and-verification:eq:mobius} into
\eqref{hyperoctahedral-cube-proof-and-verification:eq:molien} proves the full finite-$n$ formula
\begin{equation}\label{hyperoctahedral-cube-proof-and-verification:eq:full-formula}
\boxed{
\mathcal H_n(\mathbf t)=
\sum_{\substack{(a_\ell),(b_\ell)\\
\sum_\ell\ell(a_\ell+b_\ell)=n}}
w(a,b)\prod_{d\geq1}Q_d(\mathbf t)^{c_d(a,b)}.
}
\end{equation}
This is exact, not an asymptotic or a truncation.

At total degree two, the orbit interpretation gives a transparent check.
The group is transitive on vertices, and pairs of vertices are classified by
their Hamming distance $0,1,\ldots,n$.  Therefore
\begin{equation}\label{hyperoctahedral-cube-proof-and-verification:eq:degree-two}
\mathcal H_n
=1+m_{(1)}+(n+1)\bigl(m_{(2)}+m_{(1,1)}\bigr)+\cdots.
\end{equation}
In particular, the coefficients do not stabilize as $n\to\infty$.

\subsection{Independent verification strategy}
The companion program avoids using the statement being tested as its only
computational route.

\paragraph{Literal matrix computation.}
For every one of the $2^n n!$ pairs $(b,\pi)$, the program builds the vertex
permutation and the matrix \eqref{hyperoctahedral-cube-proof-and-verification:eq:matrix}.  It checks that the matrices are
distinct permutation matrices.  For a partition $\tau$, it computes complete
symmetric characters from actual matrix traces using Newton's recurrence
\begin{equation}\label{hyperoctahedral-cube-proof-and-verification:eq:newton}
d\,h_d(P_g)=\sum_{r=1}^d\Tr(P_g^r)h_{d-r}(P_g),
\qquad h_0=1,
\end{equation}
and averages $\prod_jh_{\tau_j}(P_g)$ exactly.

\paragraph{Direct orbit computation.}
The monomial basis in \eqref{hyperoctahedral-cube-proof-and-verification:eq:orbit-coeff} is constructed literally.  A
breadth-first traversal using the $n$ coordinate flips and $n-1$ adjacent
coordinate transpositions counts its orbits.  This computation does not use
traces, determinants, signed types, or the proposed formula.

\paragraph{Signed-type computation.}
All nonnegative solutions $(a_\ell,b_\ell)$ of the size constraint are
generated without enumerating group elements.  Their rational weights
\eqref{hyperoctahedral-cube-proof-and-verification:eq:type-probability} are checked to sum to one.  Equations
\eqref{hyperoctahedral-cube-proof-and-verification:eq:fixed-powers} and \eqref{hyperoctahedral-cube-proof-and-verification:eq:mobius} then produce vertex cycles and
the coefficient predicted by \eqref{hyperoctahedral-cube-proof-and-verification:eq:full-formula}.

The tests also compare the entire convergent function, not only a finite
Taylor truncation.  At $(t_1,t_2)=(0.037,0.061)$ they evaluate
\[
\frac1{|B_n|}\sum_g
\frac1{\det(I-t_1P_g)\det(I-t_2P_g)}
\]
by dense determinants, by actual vertex cycles, and by the signed-type sum.

\subsection{Exact results}
The following coefficient vectors are ordered as
\[
\tau=(1),(2),(1,1),(3),(2,1),(1,1,1).
\]
Every entry agrees across literal matrices, direct orbit traversal, and the
signed-type formula.
\begin{center}
\begin{tabular}{@{}rrrrl@{}}
\toprule
$n$ & $|B_n|$ & $\dim V_n$ & $\Prob(F_n>0)$ & common coefficient vector\\
\midrule
1 & 2    & 2  & $1/2$   & $(1,2,2,2,3,4)$\\
2 & 8    & 4  & $3/8$   & $(1,3,3,4,7,10)$\\
3 & 48   & 8  & $5/16$  & $(1,4,4,7,13,20)$\\
4 & 384  & 16 & $35/128$& $(1,5,5,11,22,35)$\\
5 & 3840 & 32 & $63/256$& $(1,6,6,16,34,56)$\\
\bottomrule
\end{tabular}
\end{center}

For $n=1,\ldots,5$, the program additionally verifies, element by element,
the conditional parity law and reconstructs every vertex-cycle count from
\eqref{hyperoctahedral-cube-proof-and-verification:eq:fixed-powers}--\eqref{hyperoctahedral-cube-proof-and-verification:eq:mobius}.  It checks moments $m=1,2,3,4$,
giving 20 exact moment comparisons, and the six displayed coefficients,
giving 30 three-way coefficient comparisons.  In the largest numerical case,
the matrix determinant average is
\[
1.1620117926716296,
\]
and its maximum discrepancy from the cycle and signed-type evaluations is
$1.74\times10^{-14}$, consistent with floating-point roundoff.  All integer
and rational comparisons are exact.

The range $n\leq5$ is a practical exhaustive choice: it includes all 3840
elements at $n=5$ and orbit spaces as large as $|X_5|^3=32768$, while the
group order grows as $2^n n!$.  The proof and signed-type algorithm remain
valid for every $n$; only brute-force matrix enumeration is deliberately
bounded.

\subsection{Reproduction}
From the repository root, run
\begin{verbatim}
marimo run marimo_notebooks/hyperoctahedral_cube.py
\end{verbatim}
The notebook allows inspection of individual matrices and signed types and
shows every comparison table.  The noninteractive regression test is
\begin{verbatim}
pytest -q tests/test_hyperoctahedral_cube.py
\end{verbatim}

\subsection{Conclusion}
The proposed formulas are correct.  A random cube symmetry has no fixed
vertex with probability tending to one, but its rare nonzero fixed sets are
large enough that $\E F_n=1$ and all higher moments diverge.  The signed-cycle
power formula determines every vertex cycle and hence the complete
$\sigma$-MGF through \eqref{hyperoctahedral-cube-proof-and-verification:eq:full-formula}.  Exhaustive matrix, orbit, and
signed-type computations agree throughout the representative range
$1\leq n\leq5$.
\endgroup


\clearpage
\section{Socle, product-type, and primitive-extension permutation actions}\label{proof:new-dist:primitive_actions}

\resultbox{\textbf{Canonical testing artifacts.}\par\smallskip Exact backend: \path{lambda_distributions/dists/primitive_action_sigma_mgfs/verification.py}.\par Regression: \path{lambda_distributions/dists/primitive_action_sigma_mgfs/test_verification.py}.\par Marimo: \path{marimo_notebooks/primitive_actions.py}.}

\subsection*{Scope and outcome}

For a finite group $G$ acting on a finite set $\Omega$, we study the honest
complex matrix representation $V=\C[\Omega]$ and its multigraded
$\sigma$-MGF.  A universal Reynolds--Molien theorem expresses the series in
terms of the cycle counts of represented permutation matrices.  We prove the
specialized fixed-power formulas for almost-simple coset actions, the socle
actions underlying simple and compound diagonal type, wreath products in
product action, regular twisted-wreath actions, and holomorph-simple and
holomorph-compound actions. A full primitive extension is obtained only after
averaging the additional top-permutation and outer-automorphism cosets.
We also prove the exact moment and rare-event statements attached to these
formulas.  An executable marimo notebook constructs representative matrix
groups of dimensions $5,9,27,36,60$, and $3600$ and checks fixed points,
M\"obius-recovered cycles, Molien coefficients, and direct configuration
orbits.  The audit comprises $43{,}642$ exact checks, all passing.

\subsection{The common representation and universal formula}

Let $G$ act on a finite set $\Omega$, and let $V=\C[\Omega]$.  For a
composition or partition $\tau=(\tau_1,\ldots,\tau_s)$ put
\[
 W_\tau(V)=\bigotimes_{j=1}^s\Sym^{\tau_j}V,
 \qquad
 a_\tau(G,\Omega)=\dim W_\tau(V)^G.
\]
Our normalization is
\[
 \MGF_{G,\Omega}(\mathbf t)=\sum_\tau a_\tau(G,\Omega)m_\tau(\mathbf t).
\]
For $g\in G$, define
\[
 F_r(g)=|\Fix_\Omega(g^r)|,
 \qquad
 c_d(g)=\#\{\text{$d$-cycles of $g$ on $\Omega$}\}.
\]

\begin{theorem}[universal permutation theorem]\label{nd:primitive_actions:thm:universal}
For every finite permutation action,
\begin{align}
 F_r(g)&=\sum_{d\mid r}d\,c_d(g),
 &
 c_d(g)&=\frac1d\sum_{r\mid d}\mu(d/r)F_r(g),\label{nd:primitive_actions:eq:mobius}\\
 \MGF_{G,\Omega}(\mathbf t)
 &=\frac1{|G|}\sum_{g\in G}\prod_{i\ge1}\prod_{d\ge1}
 (1-t_i^d)^{-c_d(g)},\label{nd:primitive_actions:eq:universal-smg}\\
 [m_\tau]\MGF_{G,\Omega}
 &=\#\left(G\backslash\prod_j\Sym^{\tau_j}\Omega\right).
 \label{nd:primitive_actions:eq:orbit-interpretation}
\end{align}
In particular,
\begin{equation}\label{nd:primitive_actions:eq:fixed-moment}
 [m_{(1^s)}]\MGF_{G,\Omega}
 =\E_{g\in G}F_1(g)^s
 =\#(G\backslash\Omega^s).
\end{equation}
\end{theorem}

\begin{proof}[Imported part and family-independent reduction]
The sigma-MGF and orbit identities are Tools~\ref{tool:molien} and
\ref{tool:permutation-orbits}.  A point in a $d$-cycle of $g$ is fixed by
$g^r$ precisely when $d\mid r$, giving the first relation in
\eqref{nd:primitive_actions:eq:mobius}; ordinary M\"obius inversion gives the
second.  Substitution into the imported permutation formula proves the
remaining displayed identities.
\end{proof}

\begin{remark}[matrix construction]
The verification code stores a permutation matrix $P_g$ by the tuple
$p_g(x)$ defined by $P_ge_x=e_{p_g(x)}$.  This is a lossless sparse matrix
representation with exactly one nonzero entry per column.  Composition of
these tuples is matrix multiplication.  Thus the degree-$3600$ cases are
fully constructed matrix representations; only wasteful dense zero storage
is omitted.
\end{remark}

\subsection{Almost-simple coset actions}

Let $T\le G\le\operatorname{Aut}(T)$ and let $G$ act transitively on
$\Omega=G/M$, where $M<G$ is maximal.  For $y\in G$, let $y^G$ denote its
conjugacy class.

\begin{proposition}[class-power formula]\label{nd:primitive_actions:prop:almost-simple}
For $g\in G$ and $r\ge1$,
\begin{equation}\label{nd:primitive_actions:eq:almost-fixed}
 \boxed{F_r(g)=\frac{|C_G(g^r)|\,|(g^r)^G\cap M|}{|M|}.}
\end{equation}
Consequently,
\begin{equation}\label{nd:primitive_actions:eq:almost-smg}
 \boxed{\MGF_{G,G/M}=
 \sum_{[g]\subseteq G}\frac1{|C_G(g)|}
 \prod_{i,d}(1-t_i^d)^{-c_d(g)},}
\end{equation}
where
\[
 c_d(g)=\frac1d\sum_{r\mid d}\mu(d/r)
 \frac{|C_G(g^r)|\,|(g^r)^G\cap M|}{|M|}.
\]
\end{proposition}

\begin{proof}
The coset $xM$ is fixed by $g^r$ exactly when $x^{-1}g^rx\in M$.  Every
element of $(g^r)^G\cap M$ has $|C_G(g^r)|$ conjugators, while every fixed
coset has $|M|$ representatives.  Double counting proves
\eqref{nd:primitive_actions:eq:almost-fixed}.  Theorem~\ref{nd:primitive_actions:thm:universal} and grouping the group
average by conjugacy classes prove \eqref{nd:primitive_actions:eq:almost-smg}.
\end{proof}

There is no universal almost-simple limit.  Natural symmetric and alternating
actions exhibit the familiar Poisson behavior, while bounded-flag and
projective regimes depend on which field or rank parameter grows.  The exact
finite formula above remains valid in every case; asymptotics require
family-specific control of class powers and subgroup intersections.

\subsection{Socle diagonal actions}

Let $T$ be nonabelian simple, $N=T^k$, $D=\operatorname{diag}(T)$, and
$\Omega=N/D$.  For a conjugacy class $C\subseteq T$, set
$c_C=|C_T(x)|$ for $x\in C$.

\begin{theorem}[socle diagonal fixed powers and moments]\label{nd:primitive_actions:thm:diagonal}
For $g=(g_1,\ldots,g_k)\in T^k$,
\begin{equation}\label{nd:primitive_actions:eq:diagonal-fixed}
 \boxed{F_r(g)=
 \begin{cases}
 c_C^{k-1},&g_1^r,\ldots,g_k^r\text{ all lie in the same class }C,\\
 0,&\text{otherwise}.
 \end{cases}}
\end{equation}
Equations \eqref{nd:primitive_actions:eq:mobius} and \eqref{nd:primitive_actions:eq:universal-smg} therefore give the
exact $\sigma$-MGF.  If $X_k(g)=F_1(g)$, then for every $m\ge1$,
\begin{equation}\label{nd:primitive_actions:eq:diagonal-moment}
 \boxed{\E X_k^m=\sum_{C\in\Cl(T)}c_C^{(m-1)k-m}.}
\end{equation}
In particular,
\begin{equation}\label{nd:primitive_actions:eq:diagonal-pair}
 [m_{(1,1)}]\MGF=\sum_Cc_C^{k-2},
 \qquad
 \Pr(X_k>0)=\sum_Cc_C^{-k}.
\end{equation}
\end{theorem}

\begin{proof}
A coset $xD$ is fixed by $g^r$ exactly when
$x^{-1}g^rx=(y,\ldots,y)$ for some $y\in T$.  If the component powers are
not conjugate, there is no solution.  If all lie in $C$, each $y\in C$ has
$c_C$ possible conjugators in every coordinate.  Dividing the
$|C|c_C^k=|T|c_C^{k-1}$ solutions by $|D|=|T|$ proves
\eqref{nd:primitive_actions:eq:diagonal-fixed}.  A random component lies in $C$ with probability
$|C|/|T|=1/c_C$.  Thus the event that all $k$ components lie in $C$ has
probability $c_C^{-k}$ and the fixed-point value there is $c_C^{k-1}$.
Summing its $m$th power proves \eqref{nd:primitive_actions:eq:diagonal-moment}.
\end{proof}

The identity class alone contributes $|T|^{k-2}$ to the pair coefficient.
Therefore the unnormalized series has no finite coefficientwise limit as
$k\to\infty$.  At the same time $\sum_Cc_C^{-1}=1$, so $\E X_k=1$, while
$\Pr(X_k>0)\to0$ and all higher moments grow.  This is an exact rare-event
law: $X_k\to0$ in probability but not in moments.

\begin{remark}[Full primitive diagonal-type groups]
The theorem above is the exact series for the socle $N=T^k$ acting on
$N/\operatorname{diag}(T)$. If a larger primitive group $G$ contains $N$
normally, the full $G$-series is the normalized average over every coset of
$N$ in $G$. Top permutations and outer automorphisms act on the common-class
condition and may merge or split $N$-orbits. No average over $N$ alone is
labelled the full primitive-group series.
\end{remark}

\subsection{Compound diagonal actions}

Partition $b\ell$ factors into $b$ blocks $B_j$ of size $\ell$ and set
\[
 N=T^{b\ell},\qquad
 D=\prod_{j=1}^b\operatorname{diag}_{B_j}(T),\qquad
 \Omega=N/D\cong\prod_{j=1}^b T^\ell/\operatorname{diag}(T).
\]
Write $g=(g_{j,s})$.  Applying Theorem~\ref{nd:primitive_actions:thm:diagonal} independently in
each block proves
\begin{equation}\label{nd:primitive_actions:eq:compound-fixed}
 \boxed{F_r(g)=\prod_{j=1}^b
 \begin{cases}
 c_{C_j}^{\ell-1},&g_{j,1}^r,\ldots,g_{j,\ell}^r\in C_j,\\
 0,&\text{otherwise}.
 \end{cases}}
\end{equation}
Together with Theorem~\ref{nd:primitive_actions:thm:universal}, this is the exact compound-diagonal
$\sigma$-MGF.

Put $R_\ell(T)=\sum_Cc_C^{\ell-2}$, the rank of one diagonal block.  Without
block permutations, pair orbits are independent, so
\[
 [m_{(1,1)}]=R_\ell(T)^b.
\]
If the full $S_b$ permutes the blocks, a pair orbit is instead a multiset of
$b$ single-block orbital types.  Hence
\begin{equation}\label{nd:primitive_actions:eq:compound-pair}
 \boxed{[m_{(1,1)}]=\binom{b+R_\ell(T)-1}{R_\ell(T)-1}.}
\end{equation}
Thus block permutations reduce exponential divergence in $b$ to polynomial
divergence.

\subsection{Wreath products in product action}

Let $H$ act transitively on a finite set $\Omega$, and let
$W_k=H^k\rtimes S_k$ act on $\Omega^k$.  Write
$g=(h_1,\ldots,h_k;\pi)$.  For a cycle $C$ of $\pi$, let $a_C$ be the
product of its labels around the cycle; a different starting point conjugates
$a_C$, which does not change any fixed-point number.  Put
$f_s(a)=|\Fix_\Omega(a^s)|$ and $d_C(r)=\gcd(|C|,r)$.

\begin{theorem}[product-action fixed powers]\label{nd:primitive_actions:thm:product-fixed}
For every $r\ge1$,
\begin{equation}\label{nd:primitive_actions:eq:product-fixed}
 \boxed{F_r(g)=\prod_{C\in\Cyc(\pi)}
 f_{r/d_C(r)}(a_C)^{d_C(r)}.}
\end{equation}
Consequently \eqref{nd:primitive_actions:eq:mobius} and \eqref{nd:primitive_actions:eq:universal-smg} give the exact
product-action $\sigma$-MGF.
\end{theorem}

\begin{proof}
On a coordinate cycle $C$ of length $q$, the permutation $\pi^r$ splits into
$d=\gcd(q,r)$ cycles.  Following one new cycle to its starting coordinate
applies an element conjugate to $a_C^{r/d}$.  Each of the $d$ cycles can be
assigned independently a point fixed by that element, producing the factor
$f_{r/d}(a_C)^d$.  Different original cycles use disjoint coordinates and
multiply.
\end{proof}

\begin{theorem}[exact product-action moments]\label{nd:primitive_actions:thm:product-moment}
Let $R_m(H,\Omega)=\#(H\backslash\Omega^m)$.  Then
\begin{equation}\label{nd:primitive_actions:eq:product-moment}
 \boxed{\E_{g\in H\wr S_k}|\Fix_{\Omega^k}(g)|^m
 =\binom{k+R_m(H,\Omega)-1}{R_m(H,\Omega)-1}.}
\end{equation}
\end{theorem}

\begin{proof}
By \eqref{nd:primitive_actions:eq:fixed-moment}, the left side counts orbits on
$(\Omega^k)^m$.  Regard an $m$-tuple of words as $k$ columns in $\Omega^m$.
The base group independently replaces each column by one of the $R_m$
$H$-orbit types; the top $S_k$ forgets column order.  The remaining datum is
a multiset of size $k$ chosen from $R_m$ types, counted by the displayed
binomial coefficient.
\end{proof}

For $m=2$, if $H$ has permutation rank $R_2=r_H$, then
\[
 [m_{(1,1)}]\MGF=\binom{k+r_H-1}{r_H-1}.
\]
Every nontrivial transitive action has $r_H\ge2$, so a finite coefficientwise
limit is impossible as $k\to\infty$.

Let $\delta$ be the derangement proportion of $H$.  Conditional on $\pi$,
the cycle products $a_C$ are independent uniform elements of $H$.  A fixed
word exists exactly when no cycle product is a derangement, hence
\[
 \Pr(F_1>0\mid\pi)=(1-\delta)^{K(\pi)}.
\]
Using the cycle-count generating identity for $S_k$ gives
\begin{equation}\label{nd:primitive_actions:eq:product-survival}
 \boxed{\Pr(F_1>0)=\frac{(1-\delta)(2-\delta)\cdots(k-\delta)}{k!}
 \sim\frac{k^{-\delta}}{\Gamma(1-\delta)}.}
\end{equation}
For finitely many pairs $(q,C)$, the number of $q$-cycles whose cycle product
lies in a conjugacy class $C\subseteq H$ converges jointly to independent
Poisson variables of means $|C|/(q|H|)$.  The input is therefore marked
Poisson, but $F_1$ is a killed multiplicative functional of that input, not
ordinarily a compound-Poisson sum.

\subsection{Twisted-wreath regular actions}

Let $N=T^k$ and $G=N\rtimes P$, where $P$ acts on $N$ by automorphisms.
Use the affine action on the regular set $\Omega=N$,
\[
 x^{(n,p)}=n\,p(x).
\]
For $g=(n,p)$ write
\[
 g^r=(n_r,p^r),\qquad
 n_r=n\,p(n)\cdots p^{r-1}(n),
\]
and define the nonabelian Lang map $L_\alpha(x)=x\alpha(x)^{-1}$.

\begin{theorem}[Lang-map formula and spike law]\label{nd:primitive_actions:thm:twisted}
For every $r\ge1$,
\begin{equation}\label{nd:primitive_actions:eq:twisted-fixed}
 \boxed{F_r(n,p)=
 \begin{cases}
 |\Fix_N(p^r)|,&n_r\in\im L_{p^r},\\
 0,&\text{otherwise}.
 \end{cases}}
\end{equation}
Conditionally on $p$ and with $n$ uniform in $N$,
\begin{equation}\label{nd:primitive_actions:eq:spike-law}
 \boxed{F_1(n,p)=
 \begin{cases}
 |\Fix_N(p)|,&\text{with probability }|\Fix_N(p)|^{-1},\\
 0,&\text{otherwise}.
 \end{cases}}
\end{equation}
Consequently
\begin{equation}\label{nd:primitive_actions:eq:twisted-moment}
 \boxed{\E F_1^m=\frac1{|P|}\sum_{p\in P}|\Fix_N(p)|^{m-1}
 =\#(P\backslash N^{m-1}),}
\end{equation}
and $[m_{(1,1)}]\MGF=\#(P\backslash N)$.
\end{theorem}

\begin{proof}
The fixed-point equation is $x=n_rp^r(x)$, equivalently
$L_{p^r}(x)=n_r$.  If $L_\alpha(x)=L_\alpha(y)$, then
$z=y^{-1}x$ satisfies $z=\alpha(z)$; conversely, right multiplication by
any $z\in\Fix_N(\alpha)$ preserves the Lang value.  Thus every nonempty
fiber of $L_\alpha$ is a right coset of $\Fix_N(\alpha)$.  This proves
\eqref{nd:primitive_actions:eq:twisted-fixed}.  Its image has size
$|N|/|\Fix_N(p)|$, so a uniform $n$ belongs to it with probability
$|\Fix_N(p)|^{-1}$, proving \eqref{nd:primitive_actions:eq:spike-law}.  Averaging the $m$th
power gives the first equality in \eqref{nd:primitive_actions:eq:twisted-moment}; Burnside's lemma
for the action of $P$ on $N^{m-1}$ gives the second.
\end{proof}

If $p$ permutes factors with automorphism labels, and $\beta_C$ is the twist
product around a coordinate cycle $C$, then direct propagation around each
cycle gives
\begin{equation}\label{nd:primitive_actions:eq:twisted-cycle}
 \boxed{|\Fix_N(p)|=\prod_C|\Fix_T(\beta_C)|.}
\end{equation}
For the untwisted benchmark $P=S_k$, this becomes
$|\Fix_N(p)|=|T|^{K(p)}$, so
\begin{align}
 \E F_1^m
 &=\binom{k+|T|^{m-1}-1}{|T|^{m-1}-1},\label{nd:primitive_actions:eq:twisted-benchmark-moment}\\
 \Pr(F_1>0)
 &=\frac{(1/|T|)(1+1/|T|)\cdots(k-1+1/|T|)}{k!}
 \sim\frac{k^{1/|T|-1}}{\Gamma(1/|T|)}.
 \label{nd:primitive_actions:eq:twisted-benchmark-survival}
\end{align}
This retains the same Bernoulli-spike mechanism as genuine primitive
twisted-wreath groups.

\subsection{Holomorph-simple and holomorph-compound actions}

\subsubsection{Holomorph-simple left-right action}

Let $T\times T$ act on $\Omega=T$ by
\[
 x^{(a,b)}=a^{-1}xb.
\]
Then $x$ is fixed by $(a,b)^r$ precisely when $a^r=xb^rx^{-1}$.

\begin{proposition}[holomorph-simple formulas]\label{nd:primitive_actions:prop:hs}
\begin{equation}\label{nd:primitive_actions:eq:hs-fixed}
 \boxed{F_r(a,b)=
 \begin{cases}|C_T(a^r)|,&a^r\sim b^r,\\0,&\text{otherwise}.
 \end{cases}}
\end{equation}
If $X(a,b)=F_1(a,b)$, then
\begin{equation}\label{nd:primitive_actions:eq:hs-moment}
 \boxed{\E X^m=\sum_{C\in\Cl(T)}c_C^{m-2},
 \qquad [m_{(1,1)}]\MGF=|\Cl(T)|.}
\end{equation}
\end{proposition}

\begin{proof}
The fixed-point equation proves \eqref{nd:primitive_actions:eq:hs-fixed}; if it is solvable, the
set of conjugators is a centralizer coset.  For a class $C$, the probability
that independent uniform $a,b$ both lie in $C$ is
$(|C|/|T|)^2=c_C^{-2}$, and the fixed-point value is $c_C$.  Summing its
$m$th power proves \eqref{nd:primitive_actions:eq:hs-moment}.
\end{proof}

Thus a family in which $|\Cl(T)|$ grows has no coefficientwise
holomorph-simple limit.  Adding outer automorphisms replaces ordinary
conjugacy by the corresponding twisted-conjugacy equation.

\subsubsection{Holomorph-compound product action}

Let $(T^k\times T^k)\rtimes S_k$ act on $T^k$ coordinatewise by left-right
multiplication, with $S_k$ permuting coordinates.  This is exactly the
product action of the base group $H=T\times T\curvearrowright T$.  The
base rank is $R_2(H,T)=|\Cl(T)|$, so
\begin{equation}\label{nd:primitive_actions:eq:hc-pair}
 \boxed{[m_{(1,1)}]\MGF=
 \binom{k+|\Cl(T)|-1}{|\Cl(T)|-1}.}
\end{equation}
All fixed-power and higher-moment formulas are those of
Theorems~\ref{nd:primitive_actions:thm:product-fixed} and \ref{nd:primitive_actions:thm:product-moment} with this base
action.

\subsection{Exact verification strategy and results}

\subsubsection{Independent routes}

The executable package uses four logically distinct routes.
\begin{enumerate}
\item \textbf{Matrix route.}  It constructs every represented permutation
matrix in the exhaustive cases and computes $F_r$ from the cycles of its
actual permutation.
\item \textbf{Structural route.}  It computes the right side of the relevant
class-intersection, common-conjugacy-class, cycle-product, Lang-map, or
left-right power formula without consulting the represented matrix cycles.
\item \textbf{Molien route.}  For cycle counts $c_d(g)$ it extracts
$[t^q]\prod_d(1-t^d)^{-c_d(g)}$ exactly and averages products of these
coefficients over the group.
\item \textbf{Orbit route.}  It constructs products of multisets or ordered
tuples and counts their orbits directly.  This does not use determinant or
character formulas.
\end{enumerate}
The M\"obius formula is also checked by reconstructing every cycle count from
fixed points of powers.

\subsubsection{Representative cases}

\begin{center}
\small
\begin{tabularx}{\linewidth}{@{}p{0.22\linewidth}p{0.34\linewidth}rX@{}}
\toprule
Family & Constructed action & Matrix dimension & Exact comparison \\
\midrule
Almost simple
& $A_5\curvearrowright A_5/A_4$
& $5$
& All $60$ elements, powers $1$--$6$, five $\tau$ values \\
Simple diagonal / HS
& $A_5^2\curvearrowright A_5$ by left-right multiplication
& $60$
& All $3600$ elements, powers $1$--$6$, direct pair orbits \\
Simple diagonal
& $A_5^3\curvearrowright A_5^2$
& $3600$
& Representative fixed powers; all simultaneous-conjugacy orbits \\
Compound diagonal
& Two $A_5^2/\operatorname{diag}(A_5)$ blocks, with and without $S_2$
& $3600$
& Direct stabilizer orbits: $25$ and $15$ \\
Product action
& $S_3\wr S_k\curvearrowright\{0,1,2\}^k$, $k=2,3$
& $9,27$
& All $72$ and $1296$ elements, powers $1$--$6$, moments $2,3$ \\
Twisted wreath
& $S_3^2\rtimes S_2\curvearrowright S_3^2$ affinely
& $36$
& All $72$ elements, powers $1$--$6$, moments $1,2,3$ \\
Holomorph compound
& $((S_3\times S_3)^2\rtimes S_2)\curvearrowright S_3^2$
& $36$
& All $2592$ elements, five powers, direct pair orbits \\
\bottomrule
\end{tabularx}
\end{center}

The smaller $S_3$ wreath benchmarks are intentional.  The fixed-point
identities are valid for arbitrary finite base groups, while the smaller base
permits exhaustive enumeration of every matrix and target orbit.  The
$A_5$ cases separately test the hypotheses that require a nonabelian simple
group.

\subsubsection{Numerical results, all exact}

For $A_5$, the centralizer sizes are $3,4,5,5,60$.  Hence the direct
diagonal pair coefficients are
\[
 k=2:\quad 5,
 \qquad
 k=3:\quad 3+4+5+5+60=77.
\]
For two diagonal $k=2$ blocks, direct orbit traversal gives $5^2=25$ without
block swaps and $\binom{2+5-1}{5-1}=15$ with the full $S_2$.

For $S_3\wr S_k$ on $\{0,1,2\}^k$, the base has $R_2=2$ and $R_3=5$.
Direct orbit counts, Burnside averages, and Theorem~\ref{nd:primitive_actions:thm:product-moment}
all give
\[
\begin{array}{c|cc}
k &[m_{(1,1)}]&[m_{(1,1,1)}]\\ \hline
2&3&15\\
3&4&35
\end{array}
\]
For the affine benchmark $S_3^2\rtimes S_2$, the first three fixed-point
moments are exactly $1,21,666$, and the direct orbit count
$|S_2\backslash S_3^2|$ is $21$.  For the holomorph-compound benchmark, the
direct pair-orbit count, Burnside average, and formula
$\binom{2+3-1}{3-1}$ are all $6$.

The final check count is
\[
 360+21600+16+432+7776+432+12960+66=43{,}642.
\]
The summands are respectively: almost-simple class powers; degree-$60$
diagonal/HS powers; degree-$3600$ diagonal samples; product-action $k=2$;
product-action $k=3$; twisted-wreath powers; holomorph-compound powers; and a
separate universal M\"obius audit on all $A_5$ and $S_3$ natural-action
matrices.  Every equality is integral; no floating-point tolerance is used.

\subsection{Asymptotic classification and cautions}

\begin{center}
\small
\begin{tabularx}{\linewidth}{@{}p{0.19\linewidth}p{0.32\linewidth}X@{}}
\toprule
Type & Exact structural input & Typical asymptotic \\
\midrule
Almost simple & $C_G(g^r)$ and $(g^r)^G\cap M$ & Family-dependent \\
Simple diagonal & Common power-conjugacy class & Exponential moment explosion \\
Compound diagonal & Blockwise diagonal formula & Polynomial after block permutation \\
Product action & Marked cycle products in $H$ & $F_1\to0$; polynomially divergent moments \\
Twisted wreath & Lang solvability and $\Fix_N(p^r)$ & Bernoulli spikes \\
Holomorph simple & Conjugacy of left/right powers & Pair coefficient $|\Cl(T)|$ \\
Holomorph compound & Product action over HS & Polynomial moment explosion \\
\bottomrule
\end{tabularx}
\end{center}

\resultbox{\textbf{Limit class.}  These families typically combine
marked-Poisson cycle inputs with multiplicative rare-event observables:
$F_1\to0$ in probability, $\E F_1=1$, and higher moments diverge.  Calling
this an ordinary compound-Poisson $\Lambda$-limit would be misleading.
The precise statement is convergence in probability of the scalar observable
together with explicit divergence of selected moments; there is no raw
moment/$\sigma$-MGF limit.}

Three cautions delimit the statement.  First, all representations here are
the finite permutation modules $\C[\Omega]$; no canonical embedding of an
abstract finite group into complex matrices is being assumed beyond the
stated action.  Second, coefficientwise nonconvergence follows as soon as one
fixed coefficient diverges, but convergence in probability of $F_1$ is a
different notion.  Third, genuine primitive twisted-wreath groups impose
additional structure on $T$ and $P$; the exhaustive $S_3$ computation is a
benchmark for the universal affine/Lang identity, not a claim that this small
benchmark itself satisfies every primitivity hypothesis.

\subsection{Reproducibility}

The files accompanying this note are:
\begin{itemize}
\item \texttt{verification.py}: exact group constructors, sparse permutation
matrices, fixed-power formulas, coefficient extraction, and orbit traversal;
\item \texttt{test\_verification.py}: automated regression tests for every
family; and
\item \path{marimo_notebooks/primitive_actions.py}: a marimo interface
presenting the checks in the
same group order as this document.
\end{itemize}
From the repository root, run
\begin{verbatim}
.venv/bin/python -m lambda_distributions.dists.primitive_action_sigma_mgfs.verification
.venv/bin/pytest -q lambda_distributions/dists/primitive_action_sigma_mgfs/test_verification.py
.venv/bin/marimo run marimo_notebooks/primitive_actions.py
\end{verbatim}

\begin{thebibliography}{9}
\bibitem{nd:primitive_actions:product-embedding}
\emph{Embedding permutation groups into wreath products in product action},
\url{https://arxiv.org/abs/1108.3611}.

\bibitem{nd:primitive_actions:derangements}
\emph{Derangements in simple and primitive groups},
\url{https://arxiv.org/abs/math/0208022}.

\bibitem{nd:primitive_actions:diagonal-bases}
\emph{Base sizes of primitive groups of diagonal type},
\url{https://arxiv.org/abs/2303.14290}.

\bibitem{nd:primitive_actions:wreath-derangements}
\emph{Derangements in wreath products of permutation groups},
\url{https://arxiv.org/abs/2207.01215}.

\bibitem{nd:primitive_actions:twisted-bases}
\emph{Bases of twisted wreath products},
\url{https://arxiv.org/abs/2102.02190}.
\end{thebibliography}

\clearpage
\part{Finite affine, Lie-type, and exceptional finite families}
\section{Power characters, parabolic cycles, and bounded-support stability}\label{proof:power-bounded}
\begingroup
\definecolor{ink}{HTML}{17212B}
\definecolor{navy}{HTML}{17365D}
\definecolor{teal}{HTML}{147D78}
\definecolor{pale}{HTML}{EEF4F7}
\providecommand{\Sym}{\operatorname{Sym}}
\renewcommand{\Sym}{\operatorname{Sym}}
\providecommand{\Tr}{\operatorname{Tr}}
\renewcommand{\Tr}{\operatorname{Tr}}
\providecommand{\Fix}{\operatorname{Fix}}
\renewcommand{\Fix}{\operatorname{Fix}}
\providecommand{\ExpS}{\operatorname{Exp}_{\sigma}}
\renewcommand{\ExpS}{\operatorname{Exp}_{\sigma}}
\providecommand{\MGF}{\SMG}
\renewcommand{\MGF}{\SMG}
\providecommand{\F}{\mathbb F}
\renewcommand{\F}{\mathbb F}
\providecommand{\C}{\mathbb C}
\renewcommand{\C}{\mathbb C}
\providecommand{\E}{\mathbb E}
\renewcommand{\E}{\mathbb E}
\providecommand{\resultbox}[1]{\begin{center}\fcolorbox{navy}{pale}{\parbox{0.93\linewidth}{#1}}\end{center}}
\renewcommand{\resultbox}[1]{\begin{center}\fcolorbox{navy}{pale}{\parbox{0.93\linewidth}{#1}}\end{center}}
\setcounter{theorem}{0}
\setcounter{proposition}{0}
\setcounter{lemma}{0}
\setcounter{corollary}{0}
\setcounter{remark}{0}
\begin{abstract}
For an arbitrary finite group and honest complex representation we derive the
sigma-MGF from character values on all powers.  For permutation actions this
becomes an exact fixed-point and cycle formula; for coset actions it becomes a
centralizer--subgroup-intersection formula, and for conjugation it becomes a
centralizer formula.  We then prove coefficientwise fixed-field stability for
bounded flags in the classical growing-rank towers.  A marimo notebook checks
the linear, parabolic, conjugation, and stability statements by explicit
matrices or generator orbits.  The proof is general; finite computations are
reported only as independent exact verification.
\end{abstract}

\resultbox{\textbf{Status.} Equations \eqref{power-character-bounded-flag-stability-proof-and-verification:eq:molien}--\eqref{power-character-bounded-flag-stability-proof-and-verification:eq:cosetcycles}
and the type-$A$ threshold \eqref{power-character-bounded-flag-stability-proof-and-verification:eq:typeAstable} are exact theorems.  The
formed-space statement is coefficientwise stability with a finite
tower-dependent threshold, not a claimed universal numerical threshold.
Complete flags are excluded: their degree-two coefficient already grows with
the Weyl group.}

\subsection{Common target and universal character-power formula}

Let $G$ be finite and let $\rho:G\to GL(V)$ be an honest finite-dimensional
complex representation with character $\chi_V$.  For a partition
$\tau=(\tau_1,\ldots,\tau_s)$ set
\[
 W_\tau(V)=\bigotimes_{j=1}^s\Sym^{\tau_j}V,
 \qquad |\tau|=\sum_j\tau_j,
\]
and define
\[
 \MGF_{G,V}=\E_{g\in G}\!\left[\ExpS(X_V(g)h_1)\right]
 =\sum_\tau\dim W_\tau(V)^G\,m_\tau.
\]

\begin{theorem}[Universal power-character formula]\label{power-character-bounded-flag-stability-proof-and-verification:thm:power}
One has the exact identities
\begin{align}
\MGF_{G,V}
 &=\frac1{|G|}\sum_{g\in G}\prod_{i\ge1}
       \det(I-t_i\rho(g))^{-1},\label{power-character-bounded-flag-stability-proof-and-verification:eq:molien}\\
 &=\frac1{|G|}\sum_{g\in G}\exp\!\left(
       \sum_{r\ge1}\frac{\chi_V(g^r)}r p_r\right),\label{power-character-bounded-flag-stability-proof-and-verification:eq:power}\\
 &=\sum_{[g]\subseteq G}\frac1{|C_G(g)|}\exp\!\left(
       \sum_{r\ge1}\frac{\chi_V(g^r)}r p_r\right).\label{power-character-bounded-flag-stability-proof-and-verification:eq:classpower}
\end{align}
Thus the series is determined by conjugacy classes, class power maps, and the
character values.  Its coefficient is
\begin{equation}\label{power-character-bounded-flag-stability-proof-and-verification:eq:coefficient}
[m_\tau]\MGF_{G,V}=\frac1{|G|}\sum_{g\in G}\prod_{j=1}^s
\left(\sum_{\lambda\vdash\tau_j}\frac1{z_\lambda}
\prod_{r\ge1}\chi_V(g^r)^{m_r(\lambda)}\right).
\end{equation}
\end{theorem}

\begin{proof}[Reduction to the imported tools]
Equation~\eqref{power-character-bounded-flag-stability-proof-and-verification:eq:molien}
is Tool~\ref{tool:molien}.  Applying the standard logarithmic determinant
change of coordinates and using
$\Tr(\rho(g)^r)=\chi_V(g^r)$ gives
\eqref{power-character-bounded-flag-stability-proof-and-verification:eq:power}.
Grouping the resulting average by conjugacy classes gives
\eqref{power-character-bounded-flag-stability-proof-and-verification:eq:classpower}; expanding each symmetric-power character in power sums gives
\eqref{power-character-bounded-flag-stability-proof-and-verification:eq:coefficient}.
No invariant-dimension or Molien identity is reproved here.
\end{proof}

\begin{remark}[Character specializations]
For a genuine Steinberg or a specified genuine Weil representation, insert
its character into \eqref{power-character-bounded-flag-stability-proof-and-verification:eq:classpower}.  For a Deligne--Lusztig virtual
character $R_T^G(\theta)$ the same expression is interpreted in the completed
Grothendieck lambda-ring; its coefficients are virtual multiplicities and
need not be nonnegative.  These are specializations of the theorem, not
claims that all such character values collapse to one uniform product.
\end{remark}

\subsection{Permutation and parabolic quotient groups}

Suppose $G$ acts on a finite set $X$, and use the honest complex permutation
module $V=\C[X]$.  Let $c_r(g;X)$ count the $r$-cycles and put
$F_d(g;X)=|\Fix_X(g^d)|$.

\begin{proposition}[Permutation-cycle formula]\label{power-character-bounded-flag-stability-proof-and-verification:prop:permutation}
For every finite action,
\begin{align}
c_r(g;X)&=\frac1r\sum_{d\mid r}\mu(r/d)F_d(g;X),\label{power-character-bounded-flag-stability-proof-and-verification:eq:mobius}\\
\MGF_{G,X}&=\frac1{|G|}\sum_{g\in G}\prod_{i,r\ge1}
 (1-t_i^r)^{-c_r(g;X)},\label{power-character-bounded-flag-stability-proof-and-verification:eq:permutation}\\
[m_\tau]\MGF_{G,X}&=\#\left(G\backslash
 \prod_j\Sym^{\tau_j}X\right).\label{power-character-bounded-flag-stability-proof-and-verification:eq:orbits}
\end{align}
\end{proposition}

\begin{proof}
An $r$-cycle has determinant factor $1-t^r$ and contributes all of its $r$
points to $F_d$ precisely when $r\mid d$.  Hence
$F_d=\sum_{r\mid d}r c_r$, and M"obius inversion gives
\eqref{power-character-bounded-flag-stability-proof-and-verification:eq:mobius}.  The determinant factors give \eqref{power-character-bounded-flag-stability-proof-and-verification:eq:permutation}.
The monomial basis of $W_\tau(\C[X])$ is the displayed product of multiset
spaces, so its invariant dimension is the number of orbits, proving
\eqref{power-character-bounded-flag-stability-proof-and-verification:eq:orbits}.
\end{proof}

\begin{proposition}[Coset and parabolic cycles]\label{power-character-bounded-flag-stability-proof-and-verification:prop:coset}
For $X=G/H$,
\begin{align}
F_d(g;G/H)&=\frac{|C_G(g^d)|}{|H|}\,|(g^d)^G\cap H|,\label{power-character-bounded-flag-stability-proof-and-verification:eq:cosetfixed}\\
c_r(g;G/H)&=\frac1r\sum_{d\mid r}\mu(r/d)
 \frac{|C_G(g^d)|}{|H|}\,|(g^d)^G\cap H|.\label{power-character-bounded-flag-stability-proof-and-verification:eq:cosetcycles}
\end{align}
Substitution in \eqref{power-character-bounded-flag-stability-proof-and-verification:eq:permutation} gives the complete exact sigma-MGF.
In particular this applies to $G=\mathbf G^F$ and every $F$-stable parabolic
$P$, including projective points, Grassmannians, polar Grassmannians, and
fixed-shape partial flags.
\end{proposition}

\begin{proof}
The coset $xH$ is fixed by $g^d$ exactly when $x^{-1}g^dx\in H$.  Each element
of $(g^d)^G\cap H$ has $|C_G(g^d)|$ conjugators $x$.  Divide by $|H|$ for the
number of representatives of a coset and then apply \eqref{power-character-bounded-flag-stability-proof-and-verification:eq:mobius}.
\end{proof}

\subsection{Conjugation groups}

For conjugation on $G$ itself,
\[
 F_d(g;G_{\mathrm{conj}})=|C_G(g^d)|.
\]
Therefore
\begin{equation}\label{power-character-bounded-flag-stability-proof-and-verification:eq:conjcycles}
c_r^{\mathrm{conj}}(g)=\frac1r\sum_{d\mid r}\mu(r/d)|C_G(g^d)|,
\qquad
\MGF_{G,\mathrm{conj}}=\frac1{|G|}\sum_g\prod_{i,r}
(1-t_i^r)^{-c_r^{\mathrm{conj}}(g)}.
\end{equation}
On a fixed conjugacy class $C$, replace the fixed-point count by
$|C\cap C_G(g^d)|$.  Hence the conjugation series is determined by
centralizers, class intersections, and power maps.

\subsection{Fixed-field stability for bounded flags}

Fix $q$ and a flag shape
$\mathbf d=(d_1<\cdots<d_s)$, with $d=d_s$ independent of ambient rank.  Let
$X_n=\operatorname{Flag}_{\mathbf d}(\F_q^n)$.

\begin{theorem}[Type-$A$ bounded-support stability]\label{power-character-bounded-flag-stability-proof-and-verification:thm:typeA}
For every partition $\tau$ of total degree $D$,
\begin{equation}\label{power-character-bounded-flag-stability-proof-and-verification:eq:typeAstable}
[m_\tau]\MGF_{GL_n(q),X_n}\quad\text{is independent of $n$ for }n\ge dD.
\end{equation}
The same holds for $PGL_n(q)$.  For $SL_n(q)$ and $PSL_n(q)$ the safe range is
$n>dD$.  The stable coefficient is
\begin{equation}\label{power-character-bounded-flag-stability-proof-and-verification:eq:stablecoefficient}
a_{\tau,\mathbf d}(q)=\sum_{r=0}^{dD}\#\left(
GL_r(q)\backslash\mathcal C_{\tau,\mathbf d}^{\mathrm{span}}(r)\right),
\end{equation}
where the configuration space consists of $\tau$-colored flag multisets whose
top subspaces span $\F_q^r$.
\end{theorem}

\begin{proof}
A degree-$D$ configuration uses $D$ flags with multiplicity.  The sum of their
top subspaces, called the support, has dimension at most $dD$.  Two
configurations are $GL_n(q)$-conjugate exactly when their supported labelled
configurations are linearly isomorphic, because every isomorphism between the
supports extends across complements.  Once $n\ge dD$, every possible support
already embeds, so no new orbit type can appear.  Sorting by the support
dimension gives \eqref{power-character-bounded-flag-stability-proof-and-verification:eq:stablecoefficient}.  Scalar matrices act trivially,
so the $PGL$ orbits agree.  If $n>dD$, a nonzero complementary direction can
correct the determinant of any extension while fixing the target support;
this gives the stated $SL$ and $PSL$ range.
\end{proof}

\begin{proposition}[Formed-space towers]\label{power-character-bounded-flag-stability-proof-and-verification:prop:formed}
Fix one symplectic, unitary, odd-orthogonal, or one of the two
even-orthogonal towers, and let $X_n$ be the totally isotropic or singular
flags of fixed shape.  For every $\tau$ there is a finite threshold
$N_{\mathcal T}(\mathbf d,\tau)$ after which
$[m_\tau]\MGF_{G_n,X_n}$ is constant.
\end{proposition}

\begin{proof}
The span of a degree-$D$ configuration again has bounded dimension.  Its orbit
is determined by the labelled flags together with the restricted Hermitian,
alternating, or quadratic form and the multiplicities.  In characteristic two
the quadratic form, not merely its polar form, is retained.  There are only
finitely many such bounded-dimensional supported types over a fixed field.
Witt extension embeds and identifies each type once the chosen tower is large
enough; taking the maximum of these finitely many embedding ranks gives the
threshold.  Determinant or spinor-norm correction may require enlarging the
unused nondegenerate complement before passing to special, derived, or
projective groups.
\end{proof}

\begin{remark}[Why complete flags are excluded]
For $G/B$, Bruhat decomposition gives
$[m_{(1,1)}]\MGF_{G,G/B}=|B\backslash G/B|=|W|$.  This grows with rank, so
bounded shape is an essential hypothesis.
\end{remark}

\subsection{Representation scope for finite Lie-type groups}

Projective, subspace, flag, and conjugacy actions are honest complex
permutation representations.  Formula \eqref{power-character-bounded-flag-stability-proof-and-verification:eq:classpower} also applies to
every specified genuine complex character, including Steinberg and chosen
Weil representations.  The defining-characteristic action on
$\mathfrak g(\F_q)$ requires care:
\begin{itemize}
\item as a finite set with conjugation action it is an honest permutation
representation, with
$F_d(g)=q^{\dim_{\F_q}C_{\mathfrak g}(g^d)}$;
\item as a linear module over $\F_q$ it is not automatically a complex
representation, so a cross-characteristic realization or lift must be
specified before applying \eqref{power-character-bounded-flag-stability-proof-and-verification:eq:power}.
\end{itemize}

\subsection{Group-separated exact verification}

The executable companion is
\path{marimo_notebooks/power_character_stability.py}.
It calls a plain Python verification module, so the same checks run without a
notebook server.

\subsubsection{Linear group: $S_3\cong GL_2(2)$ and its Steinberg module}

Every one of the six integral $2\times2$ matrices of the standard $S_3$
representation is constructed.  For each $\tau$, the program builds the
matrices of every $\Sym^{\tau_j}V$, takes their Kronecker product, averages the
actual matrices, and computes the Reynolds rank.  Independently it uses only
$\chi(g^r)$ and Newton's recurrence to evaluate \eqref{power-character-bounded-flag-stability-proof-and-verification:eq:coefficient}.  In
the order
\[
(1),(2),(1,1),(3),(2,1),(1,1,1),
\]
both routes give
\[
(0,1,1,1,1,1).
\]
The largest projector-idempotence error is $5.56\times10^{-17}$.

\subsubsection{Parabolic group: $GL_3(2)$ on $G/P$}

The program enumerates all $168$ invertible $3\times3$ matrices over $\F_2$
and their action on the seven projective points.  It separately computes all
six conjugacy classes, the point stabilizer $P$ of order $24$, centralizers,
and class intersections with $P$.  Formula \eqref{power-character-bounded-flag-stability-proof-and-verification:eq:cosetfixed} reproduces
every needed fixed-point count and reconstructs the cycles of all $168$
elements.  Actual permutation-matrix averages, cycle factors, and direct
multiset-orbit traversal agree on
\[
(1,2,2,4,5,6).
\]
There are $715$ direct fixed-point equalities and $168$ full cycle
reconstructions.

\subsubsection{Conjugation group: $S_3$ on itself}

The six conjugation permutations are built from the group law.  Direct fixed
points agree with $|C_G(g^d)|$ for every required power, and all six cycle
decompositions are reconstructed.  Matrix averages, \eqref{power-character-bounded-flag-stability-proof-and-verification:eq:conjcycles},
and direct orbits agree on
\[
(3,8,11,17,32,49).
\]

\subsubsection{Stable type-$A$ groups}

Transvections generate the actions without enumerating the large ambient
groups.  Direct configuration-orbit traversal gives:
\begin{center}
\begin{tabular}{@{}lrrrr@{}}
\toprule
$\tau$ for projective points & $n=1$ & $n=2$ & $n=3$ & $n=4$\\
\midrule
$(1)$       &1&1&1&1\\
$(2)$       &1&2&2&2\\
$(1,1)$     &1&2&2&2\\
$(3)$       &1&3&4&4\\
$(2,1)$     &1&4&5&5\\
$(1,1,1)$   &1&5&6&6\\
\bottomrule
\end{tabular}
\end{center}
The onset is exactly consistent with $n\ge|\tau|$ when $d=1$.  For
$GL_n(2)$ on $\operatorname{Gr}_2(n)$, both $\tau=(2)$ and $(1,1)$ give three
orbits at $n=4$ and again at $n=5$, matching the threshold $n\ge2|\tau|=4$
and the intersection-dimension classification.

\resultbox{\textbf{Verification result.} All 26 reported coefficient or
stability rows pass.  These computations test the formulas at small
representative dimensions with distinct failure modes; the all-rank claims
rest on the proofs above.}

\subsection{Reproduction}

From the repository root, run
\begin{verbatim}
python3 -m for_this_guy.power_character_bounded_flag_stability.verification
.venv/bin/marimo run marimo_notebooks/power_character_stability.py
\end{verbatim}

\begin{thebibliography}{9}
\bibitem{power-character-bounded-flag-stability-proof-and-verification:Howe}
S. Howe, \emph{Random matrix statistics and zeroes of $L$-functions via
probability in lambda-rings}, \href{https://arxiv.org/abs/2412.19295}{arXiv:2412.19295}.

\bibitem{power-character-bounded-flag-stability-proof-and-verification:Huang}
J. Huang, \emph{Some extensions of Witt's theorem},
\href{https://arxiv.org/abs/0705.3839}{arXiv:0705.3839}.

\bibitem{power-character-bounded-flag-stability-proof-and-verification:Thomas}
T. Thomas, \emph{The character of the Weil representation},
\href{https://arxiv.org/abs/math/0610644}{arXiv:math/0610644}.

\bibitem{power-character-bounded-flag-stability-proof-and-verification:MWZ}
P. Magyar, J. Weyman, and A. Zelevinsky, \emph{Multiple flag varieties of
finite type}, \href{https://arxiv.org/abs/math/9805067}{arXiv:math/9805067}.
\end{thebibliography}
\endgroup


\clearpage
\section{$\sigma$-MGF for finite affine groups}\label{proof:affine}
\begingroup
\definecolor{navy}{RGB}{24,55,95}
\definecolor{softblue}{RGB}{235,242,250}
\providecommand{\F}{\mathbb F}
\renewcommand{\F}{\mathbb F}
\providecommand{\AGL}{\operatorname{AGL}}
\renewcommand{\AGL}{\operatorname{AGL}}
\providecommand{\GL}{\operatorname{GL}}
\renewcommand{\GL}{\operatorname{GL}}
\providecommand{\Mult}{\operatorname{Mult}}
\renewcommand{\Mult}{\operatorname{Mult}}
\providecommand{\Sym}{\operatorname{Sym}}
\renewcommand{\Sym}{\operatorname{Sym}}
\providecommand{\Fix}{\operatorname{Fix}}
\renewcommand{\Fix}{\operatorname{Fix}}
\providecommand{\im}{\operatorname{im}}
\renewcommand{\im}{\operatorname{im}}
\providecommand{\cA}{\mathcal A}
\renewcommand{\cA}{\mathcal A}
\setcounter{theorem}{0}
\setcounter{proposition}{0}
\setcounter{lemma}{0}
\setcounter{corollary}{0}
\setcounter{remark}{0}
\begin{abstract}
Let $\AGL_n(q)$ act on $S=\F_q^n$ and hence on the permutation module
$\mathbb C[S]$.  This note proves the cycle-index form of its generalized
Molien series, the tuple-of-multisets orbit interpretation, the stable range
$n\ge |\tau|-1$, and the Gaussian-binomial formula for fixed-point moments.
The companion Marimo notebook constructs the affine permutations and their
$q^n\times q^n$ permutation matrices.  It then compares literal orbit
enumeration with coefficient extraction, checks the linear fixed-point and
Möbius formulas, and tests the moment formula.  All 18 reported comparisons
passed exactly, including $\AGL_3(2)$ of order $1344$ and $\AGL_2(3)$ of
order $432$.
\end{abstract}

\subsection{The statement being tested}
For a finite group $G$ acting on a finite set $S$, put
\[
 \MGF_G(\mathbf t)=\frac1{|G|}\sum_{g\in G}
 \prod_{a\ge1}\det(I-t_ag\mid\mathbb C[S])^{-1}.
\]
If $\tau=(\tau_1,\ldots,\tau_\ell)$ is a partition, the coefficient of the
monomial symmetric function $m_\tau$ is
\begin{equation}\label{affine-groups-proof-and-verification:eq:orbit-target}
 [m_\tau]\MGF_G
 =\#\left(G\backslash\prod_{j=1}^{\ell}\Mult_{\tau_j}(S)\right).
\end{equation}
Our affine group is
\[
 G_{n,q}=\AGL_n(q)=\F_q^n\rtimes\GL_n(q),\qquad
 (A,b)x=Ax+b.
\]
The exact proposed formula is
\begin{equation}\label{affine-groups-proof-and-verification:eq:affine-cycle-index}
 \boxed{\cA_{n,q}(\mathbf t)=\frac1{q^n|\GL_n(q)|}
 \sum_{A\in\GL_n(q)}\sum_{b\in\F_q^n}
 \prod_{r\ge1}Q_r(\mathbf t)^{c_r(A,b)},}
 \qquad
 Q_r(\mathbf t)=\prod_{a\ge1}(1-t_a^r)^{-1}.
\end{equation}
Here $c_r(A,b)$ is the number of $r$-cycles of the affine permutation.

\subsection{Imported orbit and cycle reductions}
Equation~\eqref{affine-groups-proof-and-verification:eq:orbit-target} is
Tool~\ref{tool:permutation-orbits}, and
\eqref{affine-groups-proof-and-verification:eq:affine-cycle-index} is the
permutation specialization of Tool~\ref{tool:molien}.  Order the $q^n$ points
of $\F_q^n$; the verifier constructs the corresponding matrices by
\[
 (P_{(A,b)})_{y,x}=\mathbf 1_{\{y=Ax+b\}}.
\]
The affine-specific work below is therefore the fixed-point calculation,
M\"obius recovery of cycles, stability, and moment enumeration.

There is also a useful coefficient extractor that avoids symbolic
determinants.  The number of size-$d$ multisets fixed by $g$ is
\begin{equation}\label{affine-groups-proof-and-verification:eq:fixed-multiset}
 H_d(g)=[z^d]\prod_{r\ge1}(1-z^r)^{-c_r(g)}.
\end{equation}
Therefore
\begin{equation}\label{affine-groups-proof-and-verification:eq:burnside-coefficient}
 [m_\tau]\cA_{n,q}=\frac1{|G_{n,q}|}\sum_{g\in G_{n,q}}
 \prod_{j=1}^{\ell}H_{\tau_j}(g).
\end{equation}
The implementation extracts \eqref{affine-groups-proof-and-verification:eq:fixed-multiset} by integer dynamic
programming.  Separately, it explicitly moves every multiset tuple by every
affine permutation and counts canonical orbit representatives.  Thus the two
numbers compared in the main test do not share an orbit-counting routine.

\subsection{Linear fixed-point formula and Möbius recovery}
For $g=(A,b)$,
\[
 g^d(x)=A^dx+b_d,\qquad b_d=(I+A+\cdots+A^{d-1})b.
\]
Thus $(I-A^d)x=b_d$.  Elementary linear algebra over $\F_q$ gives
\begin{equation}\label{affine-groups-proof-and-verification:eq:linear-fix}
 F_d(A,b):=\#\Fix(g^d)=
 \begin{cases}
 q^{\dim\ker(I-A^d)},&b_d\in\im(I-A^d),\\
 0,&b_d\notin\im(I-A^d).
 \end{cases}
\end{equation}
Since every $r$-cycle contributes its $r$ points to $\Fix(g^d)$ exactly when
$r\mid d$,
\[
 F_d=\sum_{r\mid d}r c_r,
 \qquad
 \boxed{c_r=\frac1r\sum_{d\mid r}\mu(r/d)F_d.}
\]
The checker computes ranks of both the coefficient matrix and its augmented
matrix modulo $q$, compares \eqref{affine-groups-proof-and-verification:eq:linear-fix} with fixed points of the
explicit permutation power, and then reconstructs every cycle count by
Möbius inversion.  It does so for every element of $\AGL_2(2)$ and
$\AGL_2(3)$ and for every required power through $q^n$.

\subsection{Stable coefficients}
Let $k=|\tau|$.  A tuple of multisets has at most $k$ support points; attach
to each support point its nonzero multiplicity vector.  Its affine span has
dimension at most $k-1$.

Suppose two such labeled configurations have affinely isomorphic spans.  An
affine isomorphism between their spans extends to an element of
$\AGL_n(q)$: choose origins, extend the induced linear isomorphism of
direction subspaces to an element of $\GL_n(q)$, and restore the translation.
The converse follows by restriction.  Hence the orbit is determined by the
intrinsic labeled affine configuration, not the ambient space, once
$n\ge k-1$.  Therefore
\[
 \boxed{[m_\tau]\cA_{n,q}\text{ is independent of }n
 \text{ for }n\ge |\tau|-1.}
\]
The computed value for $\tau=(1,1,1)$ is $5$ for both $\AGL_2(2)$ and
$\AGL_3(2)$, directly exercising this threshold.

\subsection{Fixed-point moments and the two limits}
Let $F(g)=\#\Fix(g)$.  Burnside gives
\[
 \mathbb E[F^k]=\#\bigl(\AGL_n(q)\backslash(\F_q^n)^k\bigr).
\]
Translate the first point to zero and form the $n\times(k-1)$ matrix with
columns $x_i-x_1$.  Left multiplication by $\GL_n(q)$ preserves, and is
classified by, its row space in $\F_q^{k-1}$.  It follows that
\begin{equation}\label{affine-groups-proof-and-verification:eq:moments}
 \boxed{\mathbb E[F^k]=\sum_{r=0}^{\min(n,k-1)}
 \genfrac{[}{]}{0pt}{}{k-1}{r}_q.}
\end{equation}
In particular, $\mathbb E[F^3]=q+2$ for $n=1$ and $q+3$ for $n\ge2$.
For fixed $q$, the moment stabilizes but is not the mean-one Poisson moment:
when $q>2$ the third moment already differs from $5$; for $q=2$, the fourth
moment is $16$, not $15$.  For fixed $n\ge1$, the third moment diverges as
$q\to\infty$, so the unrenormalized $\sigma$-MGF has no
coefficientwise finite limit.

\subsection{Exact verification evidence}
All arithmetic in the finite fields, cycle extraction, Burnside averages,
and Gaussian binomials is exact.  No floating-point eigenvalues are used.
The determinant identity is independently recovered from traces of the
constructed permutation matrices via Newton identities.

\begin{center}
\small
\begin{tabular}{@{}lcrrrr@{}}
\toprule
Group & $\tau$ & $|G|$ & configurations & literal orbits & formula\\
\midrule
$\AGL_1(2)$ & $(3)$ & 2 & 4 & 2 & 2\\
$\AGL_1(2)$ & $(2,1)$ & 2 & 6 & 3 & 3\\
$\AGL_2(2)$ & $(3)$ & 24 & 20 & 3 & 3\\
$\AGL_2(2)$ & $(2,1)$ & 24 & 40 & 4 & 4\\
$\AGL_2(2)$ & $(1,1,1)$ & 24 & 64 & 5 & 5\\
$\AGL_3(2)$ & $(1,1,1)$ & 1344 & 512 & 5 & 5\\
$\AGL_1(3)$ & $(3)$ & 6 & 10 & 3 & 3\\
$\AGL_2(3)$ & $(2,1)$ & 432 & 405 & 5 & 5\\
$\AGL_1(5)$ & $(1,1,1)$ & 20 & 125 & 7 & 7\\
\bottomrule
\end{tabular}
\end{center}

\begin{center}
\small
\begin{tabular}{@{}lcrr@{}}
\toprule
Group & moment & direct average & formula \eqref{affine-groups-proof-and-verification:eq:moments}\\
\midrule
$\AGL_1(2)$ & $\mathbb E F^3$ & 4 & 4\\
$\AGL_2(2)$ & $\mathbb E F^3$ & 5 & 5\\
$\AGL_3(2)$ & $\mathbb E F^4$ & 16 & 16\\
$\AGL_1(3)$ & $\mathbb E F^3$ & 5 & 5\\
$\AGL_2(3)$ & $\mathbb E F^3$ & 6 & 6\\
$\AGL_1(5)$ & $\mathbb E F^3$ & 7 & 7\\
\bottomrule
\end{tabular}
\end{center}

\paragraph{Scope and reproducibility.}
The mathematical proof applies to every prime power $q$.  The transparent
enumerator supplied here implements prime fields $\F_q$ only; extension-field
arithmetic is the sole missing software layer for nonprime prime powers.  The
Marimo notebook exposes the group, partition, and moment as controls and
prints an actual permutation matrix.  The adjacent standalone module runs
the complete fixed suite and exits nonzero on any discrepancy.

\paragraph{Conclusion.}
The cycle-index formula, the linear fixed-point/Möbius formulas, the orbit
interpretation, the stable range, and the Gaussian-binomial moments agree
with explicit finite computations across the representative dimensions and
field sizes tested.
\endgroup


\clearpage
\section{Finite linear-group permutation actions}\label{proof:finite-linear}
\begingroup
\definecolor{navy}{RGB}{24,55,95}
\definecolor{green}{RGB}{30,110,70}
\providecommand{\Sym}{\operatorname{Sym}}
\renewcommand{\Sym}{\operatorname{Sym}}
\providecommand{\Fix}{\operatorname{Fix}}
\renewcommand{\Fix}{\operatorname{Fix}}
\providecommand{\F}{\mathbb F}
\renewcommand{\F}{\mathbb F}
\providecommand{\Sg}{\SMG}
\renewcommand{\Sg}{\SMG}
\setcounter{theorem}{0}
\setcounter{proposition}{0}
\setcounter{lemma}{0}
\setcounter{corollary}{0}
\setcounter{remark}{0}
\begin{abstract}
For a finite group acting on a set $S$, we prove the determinant/cycle formula,
its fixed-point reconstruction, and its interpretation as an orbit count on
tuples of multisets.  We then specialize to nonzero-vector, projective-point,
and complete-flag actions of finite linear groups.  Exact finite-field matrix
enumerations verify all coefficients through degree three for
$PSL_3(2)$, the degree-two vector claims for $GL_2(3)$, and the Bruhat
obstruction in ranks two and three.  The computation also exposes a sharp
boundary phenomenon for $SL_2(3)$.
\end{abstract}

\subsection{Imported finite permutation formula}
For $V=\mathbb C[S]$, Tools~\ref{tool:molien} and
\ref{tool:permutation-orbits} give
\begin{align}
 \smgv{G}{V}
 &=\frac1{|G|}\sum_{g\in G}
 \prod_{i\geq1}\prod_{d\geq1}(1-t_i^d)^{-c_d(g;S)},\\
 [m_\tau]\,\smgv{G}{V}
 &=\#\left(G\backslash\prod_b\Sym^{\tau_b}S\right).
\end{align}
The new calculations begin with the fixed-point counts for the vector,
projective-point, and flag actions.

\subsubsection*{Fixed points determine the cycles}
A point on a $d$-cycle is fixed by $g^r$ exactly when $d\mid r$.  Thus
\[
 \Fix_S(g^r)=\sum_{d\mid r}d\,c_d(g;S).
\]
M\"obius inversion gives
\begin{equation}
 \boxed{c_d(g;S)=\frac1d\sum_{r\mid d}
 \mu(d/r)\Fix_S(g^r).}
\end{equation}
For $S=G/P$, $\Fix_S(g^r)$ is the value at $g^r$ of the permutation
character $\operatorname{Ind}_P^G\mathbf1$.  Grouping (2) by conjugacy class
then yields the stated centralizer-weighted formula.

\subsubsection*{Direct orbit interpretation}
A monomial basis vector of $\Sym^r\mathbb C[S]$ is a multiset of $r$ points
of $S$.  Consequently a basis of $W_\tau$ is a tuple of such multisets, of
respective sizes $\tau_1,\ldots,\tau_s$.  Since $G$ permutes this basis, the
sum of the basis vectors in each orbit is one invariant basis vector.  Hence
\begin{equation}
 \boxed{[m_\tau]\Sg_{G,S}=\left|\{G\text{-orbits on tuples of multisets of
 sizes }\tau_1,\ldots,\tau_s\}\right|.}
\end{equation}
Equation (4) supplies the independent direct route used in the computation.

\subsection{Linear and projective specializations}
For $S=\F_q^n\setminus\{0\}$ and $g\in GL_n(q)$,
\begin{equation}
 \Fix_S(g^r)=q^{\dim\ker(g^r-I)}-1.
\end{equation}
For projective points, a fixed line is an eigenline whose eigenvalue lies in
$\F_q^\times$.  The eigenspaces for distinct eigenvalues contribute disjoint
sets of lines, so
\begin{equation}
 \Fix_{\mathbb P^{n-1}(\F_q)}(g^r)=
 \sum_{\lambda\in\F_q^\times}
 \frac{q^{\dim\ker(g^r-\lambda I)}-1}{q-1}.
\end{equation}
Substitution of (5) or (6) into (3) and then (2) proves the proposed exact
formulas.

Over $\F_2$, every projective line has one nonzero vector and
$GL_n(2)=SL_n(2)=PGL_n(2)=PSL_n(2)$.  Thus (5) directly gives the displayed
$PSL_n(2)$ formula.

\subsubsection*{Stable range}
A configuration of total degree $D=|\tau|$ involves at most $D$ projective
points and its span has dimension at most $D$.  When $n\geq D$, two such
configurations are $GL_n(q)$-equivalent exactly when the induced colored
configurations on their spans are linearly isomorphic: any isomorphism of the
spans extends after choosing complements.  This proves coefficientwise
stability for $GL/PGL$ projective actions.  The same argument for nonzero
vectors is immediate.  For $SL/PSL$, the determinant of an extension can be
corrected on an unused one-dimensional summand, which requires the strict
range $n>D$.

For projective points the stable expansion through degree three is
\begin{equation}
 1+m_{(1)}+2m_{(2)}+2m_{(1,1)}
 +4m_{(3)}+5m_{(2,1)}+6m_{(1,1,1)}+O(\deg4).
\end{equation}
For three distinct points, collinear and noncollinear triples give two
different orbits; this is the extra degree-three term relative to the
symmetric-group/Poisson pattern.

For nonzero vectors, ordered pairs consist of one independent orbit and one
dependent orbit $y=ax$ for each $a\in\F_q^\times$, giving
\[
 [m_{(1,1)}]\Sg=q.
\]
For an unordered pair, the dependent ratios other than $1$ are identified by
$a\leftrightarrow a^{-1}$.  Together with the repeated and independent types,
this gives
\[
 [m_{(2)}]\Sg=2+\left\lfloor\frac{q-1}{2}\right\rfloor.
\]

\subsection{Complete flags and failure of rank stability}
For $S=G/B$, orbitals are the double cosets $B\backslash G/B$.  Bruhat
decomposition identifies these with the Weyl group, whence
\begin{equation}
 [m_{(1,1)}]\Sg_{G,G/B}=|W|.
\end{equation}
In type $A_{n-1}$ this is $n!$, so even the degree-two coefficient cannot
stabilize as rank grows.

\subsection{Exact verification method}
The accompanying marimo notebook performs the following independent steps.
\begin{enumerate}
 \item Enumerate every finite-field matrix and retain the required determinant
 or invertibility condition.
 \item Induce its permutation on nonzero vectors, normalized projective
 representatives, or incident point--hyperplane flags.
 \item Count orbits directly on tuples of multisets, without using character
 or cycle formulas.
 \item Extract the same coefficient from the actual cycle counts in (2).
 \item Reconstruct every cycle count from fixed points of powers using (3).
 \item At $(t_1,t_2)=(0.07,0.11)$, compare the direct matrix-determinant
 average with the cycle-product average.
\end{enumerate}
The last two numerical averages both equal $1.247556190416$ for the
$PSL_3(2)$ projective action, with difference below displayed precision.

\begin{center}
\small
\begin{tabular}{@{}lcrrr@{}}
\toprule
Action & $\tau$ & direct orbits & cycle formula & claimed\\
\midrule
$PSL_3(2)$ on $\mathbb P^2(\F_2)$ & $(1)$ & 1&1&1\\
 & $(2)$ &2&2&2\\
 & $(1,1)$ &2&2&2\\
 & $(3)$ &4&4&4\\
 & $(2,1)$ &5&5&5\\
 & $(1,1,1)$ &6&6&6\\
$PSL_2(2)$ on $\mathbb P^1(\F_2)$ & $(3)$ &3&3&3\\
$GL_2(3)$ on nonzero vectors & $(2)$ &3&3&3\\
 & $(1,1)$ &3&3&3\\
$SL_2(3)$ on nonzero vectors & $(1,1)$ &4&4&4\\
$PGL_2(3)$ on projective points & $(1,1)$ &2&2&2\\
$PSL_2(3)$ on projective points & $(1,1)$ &2&2&2\\
$GL_2(2)$ on complete flags & $(1,1)$ &2&2&2\\
$GL_3(2)$ on complete flags & $(1,1)$ &6&6&6\\
\bottomrule
\end{tabular}
\end{center}

The constructed orders are $|GL_3(2)|=168$, $|GL_2(3)|=48$, and
$|SL_2(3)|=24$; their projective images in degree two have orders 24 and 12.
All displayed tests pass exactly.

\paragraph{Sharpness found by verification.}
The $SL_2(3)$ value $4$ at $\tau=(1,1)$ differs from the $GL_2(3)$ value
$3$.  Here $n=|\tau|=2$, so there is no unused complement on which to correct
the determinant.  Independent ordered pairs split into two $SL_2(3)$ orbits
according to their determinant.  This confirms that the proposed $SL/GL$
agreement must retain the strict hypothesis $n>|\tau|$.

\subsection{Scope of the conclusion}
The matrix, orbit, fixed-point, and determinant routes agree for every
representative case above.  These computations verify the exact mechanism and
the low-degree stable claims; the proofs of (1)--(8), rather than finite
sampling, establish the formulas for all finite parameters in their stated
ranges.
\endgroup


\clearpage
\section{Finite polar-group permutation actions}\label{proof:finite-polar}
\begingroup
\definecolor{navy}{RGB}{24,55,95}
\providecommand{\Sym}{\operatorname{Sym}}
\renewcommand{\Sym}{\operatorname{Sym}}
\providecommand{\Fix}{\operatorname{Fix}}
\renewcommand{\Fix}{\operatorname{Fix}}
\providecommand{\F}{\mathbb F}
\renewcommand{\F}{\mathbb F}
\providecommand{\Sg}{\SMG}
\renewcommand{\Sg}{\SMG}
\setcounter{theorem}{0}
\setcounter{proposition}{0}
\setcounter{lemma}{0}
\setcounter{corollary}{0}
\setcounter{remark}{0}
\begin{abstract}
This note specializes the exact permutation $\sigma$-MGF to isotropic and
singular point actions of finite unitary, symplectic, and orthogonal groups.
Witt extension proves the stable three-orbit classification of pairs.  The
companion marimo notebook constructs $Sp_4(2)$, $O_4^+(2)$, and $U_4(2)$ as
finite-field matrix groups, induces their point actions, and independently
compares direct orbit counts, cycle products, fixed-point reconstruction, and
matrix determinants.  Every tested coefficient agrees exactly.
\end{abstract}

\subsection{Exact formula for a polar action}
Let a finite classical group $G$ act on a finite set $S$ of isotropic or
singular one-spaces.  For $g\in G$, let $c_d(g;S)$ denote the number of
$d$-cycles of its action on $S$.  The permutation-matrix eigenvalues on a
$d$-cycle are the $d$th roots of unity, so
\begin{equation}
 \boxed{\Sg_{G,S}(t_\bullet)=\frac1{|G|}\sum_{g\in G}
 \prod_{i,d\geq1}(1-t_i^d)^{-c_d(g;S)}.}
\end{equation}
Moreover,
\begin{equation}
 \Fix_S(g^r)=\sum_{d\mid r}d\,c_d(g;S),\qquad
 c_d(g;S)=\frac1d\sum_{r\mid d}\mu(d/r)\Fix_S(g^r).
\end{equation}
Thus one may compute (1) purely from the number of isotropic or singular
eigenlines of every power $g^r$.  Equivalently, for a parabolic point
stabilizer $P$, this fixed-point count is
$(\operatorname{Ind}_P^G\mathbf1)(g^r)$.

For $\tau=(\tau_1,\ldots,\tau_s)$, expanding the determinants gives
\begin{equation}
 [m_\tau]\Sg_{G,S}=\dim\left(
 \bigotimes_{b=1}^s\Sym^{\tau_b}\mathbb C[S]\right)^G.
\end{equation}
The tensor product has a basis consisting of tuples of multisets of points.
Because $G$ permutes this basis, (3) is exactly the number of $G$-orbits on
such tuples.  This orbit count is the independent direct computation used
below.

\subsection{Why the degree-two coefficient is three}
For two isotropic or singular points in a nondegenerate polar space of
sufficient Witt rank, there are three possible geometric types:
\begin{enumerate}
 \item the two points coincide;
 \item they are distinct and perpendicular;
 \item they are distinct and nonperpendicular.
\end{enumerate}
Each relation is preserved by every isometry.  Conversely, an isometry
between either pair of generated subspaces extends to the ambient space by
Witt's extension theorem.  Hence each type is one orbit.  This applies to
both ordered pairs and unordered 2-multisets, proving
\begin{equation}
 \boxed{[m_{(2)}]\Sg=[m_{(1,1)}]\Sg=3.}
\end{equation}
It follows that the stable expansion begins
\[
 \Sg=1+m_{(1)}+3m_{(2)}+3m_{(1,1)}+O(\deg3).
\]
The extra perpendicularity relation already distinguishes these limits from
the Poisson/symmetric-group series in degree two.

\paragraph{Proposition (coefficientwise stability within a fixed polar tower).}
Fix the field, one classical type, and a total degree $D$.  For each tuple of
multisets of singular or isotropic points of total degree $D$, retain the
labelled span together with the restricted form and the multiplicities of the
labelled one-spaces; in characteristic-two orthogonal cases retain the
restricted quadratic form rather than only its polar bilinear form.  Two
configurations in the same ambient nondegenerate polar space are conjugate
precisely when these labelled formed spaces are isometric, because an
isometry of their spans extends by Witt's lemma.

Only finitely many labelled formed spaces of dimension at most $D$ exist over
the fixed finite field.  Let $m_0(D)$ be the maximum, over the types that occur
somewhere in the chosen tower, of the least Witt rank in which the type
embeds.  Every coefficient of total degree $D$ is therefore constant for Witt
rank at least $m_0(D)$.  Orthogonal plus, minus, and odd-dimensional towers
are treated separately.  This proves existence of a stable threshold but
does not claim that the displayed bound is sharp.  The degree-two value three
was proved directly above and does not depend on this proposition.

\subsection{The three exact matrix constructions}
All arithmetic in the program is exact integer arithmetic.

\paragraph{$Sp_4(2)$.}
We enumerate $GL_4(2)$ and retain exactly the matrices preserving
\[
 B(x,y)=x_1y_2+x_2y_1+x_3y_4+x_4y_3.
\]
The resulting group has order $720$.  In characteristic two every one-space
is isotropic, giving a permutation action on the 15 nonzero vectors (equally,
the 15 projective points).

\paragraph{$O_4^+(2)$.}
We retain the matrices preserving the hyperbolic quadratic form
\[
 Q(x)=x_1x_2+x_3x_4.
\]
The resulting group has order $72$ and acts on its 9 nonzero singular vectors.

\paragraph{$U_4(2)$.}
We implement $\F_4=\F_2[\alpha]/(\alpha^2+\alpha+1)$ and the Hermitian form
\[
 h(x,y)=\sum_{j=1}^4\overline{x_j}y_j,\qquad \overline z=z^2.
\]
Every ordered orthonormal basis is enumerated recursively, producing exactly
\[
 |U_4(2)|=2^6\prod_{i=1}^4(2^i-(-1)^i)=77{,}760
\]
matrices.  Normalizing each nonzero vector by its first nonzero coordinate
gives 45 isotropic projective points.  The scalar center has order three, so
the faithful induced permutation image has order $25{,}920$.

\subsection{Verification strategy and results}
For each action, the notebook:
\begin{enumerate}
 \item constructs the full finite-field matrix group and induced permutation;
 \item directly enumerates orbits on ordered point pairs and 2-multisets;
 \item extracts the same coefficients from the cycle product in (1);
 \item reconstructs cycle counts from fixed points of powers using (2); and
 \item numerically compares direct permutation-matrix determinants with the
 cycle products at $(t_1,t_2)=(0.07,0.11)$.
\end{enumerate}
The regression module also compares a degree-three labelled-formed-space
orbit count in two consecutive ranks of each implemented tower.  This tests
the restricted-form encoding; it is not used as a proof of the proposition.

\begin{center}
\begin{tabular}{@{}lcrrr@{}}
\toprule
Action & $\tau$ & direct orbits & cycle formula & claimed\\
\midrule
$Sp_4(2)$ on 15 isotropic points & $(2)$ &3&3&3\\
 & $(1,1)$ &3&3&3\\
$O_4^+(2)$ on 9 singular points & $(2)$ &3&3&3\\
 & $(1,1)$ &3&3&3\\
$U_4(2)$ on 45 isotropic points & $(2)$ &3&3&3\\
 & $(1,1)$ &3&3&3\\
\bottomrule
\end{tabular}
\end{center}

Every equality in the table is exact.  All cycle counts in the tested
$Sp_4(2)$ and projective $U_4(2)$ actions are also recovered exactly from
fixed points of powers.  For the symplectic action, the direct determinant
average and cycle-product average both equal $1.304867127368$, with zero
difference at the displayed floating-point precision.

\subsection{Representations requiring additional choices}
The Weil-representation paragraph is governed by the general identity
\[
 \det(I-t\rho(g))^{-1}=\exp\left(
 \sum_{r\geq1}\frac{\chi_\rho(g^r)}{r}t^r\right),
\]
but it does not specify one numerical representation.  A reproducible matrix
test requires a choice of $q$ and $n$, an additive character, a linearization,
and whether $\rho$ is the full Weil representation or a constituent.  Even
characteristic requires a different construction.  Accordingly, no arbitrary
choice was silently inserted into the verification.

Likewise, the exceptional/Suzuki/Ree formula is an exact template for a
specified pair $(G(q),P(q))$, not a single representation.  A numerical test
requires the concrete group, field size, and parabolic action.  Once supplied,
the same fixed-point, cycle, determinant, and orbit pipeline applies without
change.

\paragraph{Conclusion.}
The three classical polar families realize exactly the predicted three pair
types in small but stable-rank examples.  The agreement of full matrix-group
enumeration, direct configuration orbits, cycle extraction, M\"obius
reconstruction, and determinants validates both the formula and the proposed
verification method.
\endgroup


\clearpage
\section{Finite local-ring permutation actions}\label{proof:new-dist:local_ring_actions}

\resultbox{\textbf{Canonical testing artifacts.}\par\smallskip Exact backend: \path{lambda_distributions/dists/local_ring_permutation_actions/verification.py}.\par Regression: \path{lambda_distributions/dists/local_ring_permutation_actions/test_verification.py}.\par Marimo: \path{marimo_notebooks/local_ring_actions.py}.}

\subsection*{Scope and outcome}

Let $R_a=\calO/\pi^a\calO$ be a finite chain ring with residue field
$\F_q$.  A subgroup of $M_n(R_a)$ is not canonically a complex matrix group;
the representations studied here are instead the honest complex permutation
modules $\C[X]$ on primitive vectors, free projective points, free summands,
flags, formed-space submodules, and congruence adjoint quotients.  We prove the
fixed-point, cycle, and $\sigma$-MGF formulas and the low-degree orbit formulas.
The fixed-level rank-stability statements are conditional on the explicit
generating-tuple and formed-space extension properties stated below.  An executable marimo notebook independently
constructs representative groups over $\mathbb Z/4$ and $\mathbb Z/9$ and
compares direct configuration orbits with Smith-kernel, homogeneous-space,
M\"obius-cycle, and adjoint-kernel calculations.  All reported checks pass
exactly.

\subsection{Representation convention and universal permutation theorem}

Fix a finite group $H$ acting on a finite set $X$, and put $V=\C[X]$.  For
$h\in H$, write
\[
 F_d(h;X)=|\Fix_X(h^d)|,
 \qquad c_r(h;X)=\#\{\text{$r$-cycles of $h$ on $X$}\}.
\]
A point on an $r$-cycle is fixed by $h^d$ exactly when $r\mid d$, so
\[
 F_d(h;X)=\sum_{r\mid d}r c_r(h;X).
\]
M\"obius inversion proves
\begin{equation}\label{nd:local_ring_actions:eq:mobius}
 \boxed{c_r(h;X)=\frac1r\sum_{d\mid r}\mu(r/d)F_d(h;X).}
\end{equation}

The permutation matrix of an $r$-cycle has all $r$th roots of unity as its
eigenvalues.  Its characteristic factor is therefore $1-t^r$.  Reynolds
averaging and the symmetric-power identity give the exact series
\begin{equation}\label{nd:local_ring_actions:eq:universal}
 \boxed{\MGF_{H,X}(t_1,t_2,\ldots)=\frac1{|H|}\sum_{h\in H}
 \prod_{i\ge1}\prod_{r\ge1}(1-t_i^r)^{-c_r(h;X)}.}
\end{equation}
Indeed, for a partition or composition $\tau=(\tau_1,\ldots,\tau_s)$,
\[
 [m_\tau]\MGF_{H,X}
 =\dim\left(\bigotimes_{j=1}^s\Sym^{\tau_j}\C[X]\right)^H.
\]
A monomial basis for $\Sym^d\C[X]$ is indexed by $d$-element multisets in
$X$.  Orbit sums are a basis of the invariants, proving the independent
interpretation
\begin{equation}\label{nd:local_ring_actions:eq:orbits}
 \boxed{[m_\tau]\MGF_{H,X}
 =\#\left(H\backslash\prod_j\Sym^{\tau_j}X\right).}
\end{equation}
Equations \eqref{nd:local_ring_actions:eq:mobius}--\eqref{nd:local_ring_actions:eq:orbits} reduce every family in this
note to a fixed-point problem on a finite set.

\subsection[GLn over Ra on primitive vectors]{$\GL_n(R_a)$ on primitive vectors}

Let
\[
 R_a=\calO/\pi^a\calO,\qquad R_a/\pi R_a\simeq\F_q,
 \qquad X^{\rm prim}_{n,a}=R_a^n\setminus\pi R_a^n.
\]
For $A\in M_n(R_a)$, let $0\le\nu_1(A),\ldots,\nu_n(A)\le a$ be its
truncated Smith exponents; a zero Smith entry receives exponent $a$.  Define
\begin{equation}\label{nd:local_ring_actions:eq:Kb}
 K_b(A)=q^{\sum_{j=1}^n\min(\nu_j(A),b)}.
\end{equation}
Smith reduction shows that $K_b(A)$ is the size of the kernel of $A$ after
reduction to $R_b$.

\subsubsection{Fixed points and exact series}
A primitive vector fixed by $h^d$ lies in $\ker(h^d-I)$ but not in
$\pi R_a^n$.  Multiplication by $\pi$ identifies $R_{a-1}^n$ with
$\pi R_a^n$, and under this identification the second kernel has size
$K_{a-1}(h^d-I)$.  Hence
\begin{equation}\label{nd:local_ring_actions:eq:primitive-fixed}
 \boxed{F_d^{\rm prim}(h)=K_a(h^d-I)-K_{a-1}(h^d-I).}
\end{equation}
Substitution in \eqref{nd:local_ring_actions:eq:mobius} and \eqref{nd:local_ring_actions:eq:universal} proves
\begin{equation}\label{nd:local_ring_actions:eq:primitive-smg}
 \boxed{\MGF^{\rm prim}_{n,a}=\frac1{|\GL_n(R_a)|}\sum_h
 \prod_{i,r\ge1}(1-t_i^r)^{-\frac1r\sum_{d\mid r}\mu(r/d)
 [K_a(h^d-I)-K_{a-1}(h^d-I)]}.}
\end{equation}
The identity element in \eqref{nd:local_ring_actions:eq:primitive-fixed} also gives
$|X^{\rm prim}_{n,a}|=q^{an}-q^{(a-1)n}$.

\subsubsection{Ordered-pair coefficient}
For $n\ge2$, transitivity allows the first primitive vector to be fixed as
$e_1$.  Write the second as $w=xe_1+y$.  Its component $y$ in a chosen
complement has valuation $r\in\{0,\ldots,a\}$, with $r=a$ meaning $y=0$.
The stabilizer of $e_1$ moves $y$ to $\pi^r e_2$ and changes $x$ by an
element of $\pi^rR_a$.  For $r=0$ there is one orbit.  For $1\le r\le a$,
primitivity forces the residue of $x$ modulo $\pi^r$ to be a unit, giving
$q^r-q^{r-1}$ orbits.  Consequently
\begin{equation}\label{nd:local_ring_actions:eq:primitive-pair}
 \boxed{[m_{(1,1)}]\MGF^{\rm prim}_{n,a}
 =1+\sum_{r=1}^a(q^r-q^{r-1})=q^a.}
\end{equation}

\subsection[GLn over Ra on projective points]{$\GL_n(R_a)$ on projective points}

Let $X^{\rm proj}_{n,a}=\mathbf P^{n-1}(R_a)$ be the free rank-one direct
summands.  If $L=R_av$ is fixed by $h^d$, then $h^dv=uv$ for a unique
$u\in R_a^\times$.  Conversely each primitive $u$-eigenvector generates a
fixed line, and every line has $|R_a^\times|$ primitive generators.  Thus
\begin{equation}\label{nd:local_ring_actions:eq:projective-fixed}
 \boxed{F_d^{\rm proj}(h)=\frac1{|R_a^\times|}\sum_{u\in R_a^\times}
 \left[K_a(h^d-uI)-K_{a-1}(h^d-uI)\right].}
\end{equation}
Together with \eqref{nd:local_ring_actions:eq:mobius} this yields
\begin{equation}\label{nd:local_ring_actions:eq:projective-smg}
 \boxed{\MGF^{\rm proj}_{n,a}=\frac1{|\GL_n(R_a)|}\sum_h
 \prod_{i,r\ge1}(1-t_i^r)^{-\frac1r\sum_{d\mid r}\mu(r/d)
 F_d^{\rm proj}(h)}.}
\end{equation}
Counting primitive generators gives
\[
 |\mathbf P^{n-1}(R_a)|
 =q^{(a-1)(n-1)}\frac{q^n-1}{q-1}.
\]

For $n\ge2$, after fixing the first line, every second line has a unique
relative-position parameter
\[
 R_a(e_1+\pi^r e_2),\qquad 0\le r\le a,
\]
up to the stabilizer.  The parameter is unchanged when the two lines are
interchanged.  It classifies both ordered pairs and two-element multisets,
including equality at $r=a$.  Therefore
\begin{equation}\label{nd:local_ring_actions:eq:projective-pair}
 \boxed{[m_{(1,1)}]\MGF^{\rm proj}_{n,a}
 =[m_{(2)}]\MGF^{\rm proj}_{n,a}=a+1.}
\end{equation}
For $n=2$ this is also the complete-flag action.  At $a=2$ its three
orbitals already exceed the field-level Weyl count $|S_2|=2$, concretely
exhibiting the extra valuation parameter in local-ring Bruhat theory.

\subsection{Free summands, flags, and homogeneous spaces}

Let $H=\GL_n(R_a)$ and let $P_{k,n-k}$ stabilize the standard free
rank-$k$ summand.  For any finite homogeneous space $H/P$,
\begin{equation}\label{nd:local_ring_actions:eq:coset-fixed}
 |\Fix_{H/P}(x)|=\frac{|C_H(x)|\,|x^H\cap P|}{|P|}.
\end{equation}
To prove this, count $y\in H$ satisfying $y^{-1}xy\in P$.  Every fixed coset
has exactly $|P|$ such representatives.  On the other hand, each element of
$x^H\cap P$ has exactly $|C_H(x)|$ conjugators.

For $X^{(k)}_{n,a}=\Gr_k^{\rm free}(R_a^n)=H/P_{k,n-k}$, equations
\eqref{nd:local_ring_actions:eq:coset-fixed}, \eqref{nd:local_ring_actions:eq:mobius}, and \eqref{nd:local_ring_actions:eq:universal} give
\begin{align}
 F_d^{(k)}(h)
 &=\frac{|C_H(h^d)|\,|(h^d)^H\cap P_{k,n-k}|}{|P_{k,n-k}|},\label{nd:local_ring_actions:eq:grass-fixed}\\
 \MGF^{(k)}_{n,a}
 &=\frac1{|H|}\sum_h\prod_{i,r\ge1}
 (1-t_i^r)^{-\frac1r\sum_{d\mid r}\mu(r/d)F_d^{(k)}(h)}.
 \label{nd:local_ring_actions:eq:grass-smg}
\end{align}
Reduction modulo $\pi$ and graph lifting over a fixed complement yield
\[
 |X^{(k)}_{n,a}|=q^{(a-1)k(n-k)}\qbinom{n}{k}.
\]

\subsubsection{Pair coefficient for free summands}
Assume $n\ge2k$ and fix one summand $U$.  For a second summand $W$, apply
Smith reduction to the projection $W\to R_a^n/U$.  Its uniquely determined
exponents satisfy
$0\le\lambda_1\le\cdots\le\lambda_k\le a$.  Changes of bases in $W$, in
$U$, and in a complementary $k$-summand put the pair in the form
\[
 U=\langle e_1,\ldots,e_k\rangle,\qquad
 W=\langle e_i+\pi^{\lambda_i}e_{k+i}:1\le i\le k\rangle.
\]
Parabolic shears remove the remaining components, so the Smith exponents are
also a complete invariant.  There are as many nondecreasing $k$-tuples from
$\{0,\ldots,a\}$ as multisets of size $k$ from an $(a+1)$-element set:
\begin{equation}\label{nd:local_ring_actions:eq:grass-pair}
 \boxed{[m_{(1,1)}]\MGF^{(k)}_{n,a}
 =[m_{(2)}]\MGF^{(k)}_{n,a}=\binom{a+k}{k}.}
\end{equation}
At $a=1$ this reduces to the $k+1$ field-level intersection types.

For a partial flag shape $\mathbf d=(d_1<\cdots<d_s)$, put
$X_{\mathbf d,n,a}=H/P_{\mathbf d,n}$.  The same proof gives
\begin{equation}\label{nd:local_ring_actions:eq:flag}
 F_d(h;X_{\mathbf d,n,a})
 =\frac{|C_H(h^d)|\,|(h^d)^H\cap P_{\mathbf d,n}|}{|P_{\mathbf d,n}|},
\end{equation}
followed by \eqref{nd:local_ring_actions:eq:mobius} and \eqref{nd:local_ring_actions:eq:universal}.  Unlike the field
case, $P_{\mathbf d,n}\backslash H/P_{\mathbf d,n}$ generally retains
valuation data and is not indexed only by the Weyl group.

\subsection{Conditional fixed-level stability as rank grows}

Fix $a$, take $Q\le\GL_d(R_a)$, and define
\[
 X_n(Q)=\operatorname{SplitInj}(R_a^d,R_a^n)/Q.
\]
Primitive vectors, projective points, rank-$k$ summands, and fixed-shape flags
arise from $d=1,Q=1$; $d=1,Q=R_a^\times$; $d=k,Q=\GL_k(R_a)$; and
$d=d_s,Q=P_{\mathbf d}(R_a)$, respectively.

Let $D=|\tau|$.  Choose a split frame for each of the $D$ objects and
concatenate the resulting blocks into an $n\times dD$ matrix $A$.  Let $L(A)$
be its row submodule in $R_a^{dD}$.  Every block projection $L(A)\to R_a^d$
is surjective.  Write
\[
 \mathcal L_{\tau,d}(R_a)=\{L\le R_a^{dD}:\text{every block projection is
 surjective}\},\qquad
 \Gamma_{\tau,Q}=Q^D\rtimes\prod_jS_{\tau_j}.
\]

Two matrices with the same row submodule are two generating $n$-tuples of
the same finite $R_a$-module. We isolate the exact hypothesis used by the
orbit argument:
\begin{description}
\item[Generating-tuple extension property $\mathrm{GT}(n,L)$.]
The group $\GL_n(R_a)$ is transitive on the surjections
$R_a^n\twoheadrightarrow L$.
\end{description}
Every submodule of $R_a^{dD}$ needs at most $dD$ generators. Assuming
$\mathrm{GT}(n,L)$ for every $L\le R_a^{dD}$, the block changes and
within-color reorderings are precisely $\Gamma_{\tau,Q}$. Hence
\begin{equation}\label{nd:local_ring_actions:eq:stable}
 \boxed{[m_\tau]\MGF_{\GL_n(R_a),X_n(Q)}
 =\#\bigl(\Gamma_{\tau,Q}\backslash\mathcal L_{\tau,d}(R_a)\bigr),
 \qquad n\ge dD.}
\end{equation}
The right side is independent of $n$. Under the same extension property, the
$\SL_n(R_a)$ coefficient agrees in the uniform sufficient range $n>dD$: an
unused coordinate corrects the determinant without moving the configuration.
The conclusion then passes to projective quotients after factoring the action
kernel. A proof attempt and the remaining module-theoretic lemma are recorded
in the final provisional part.

\clearpage
\subsection{Symplectic and orthogonal towers}
\thispagestyle{fancy}

Let $G_n$ be a finite symplectic or orthogonal group over $R_a$ and let $X_n$
be a homogeneous orbit of split totally isotropic or singular free submodules
or flags of fixed shape.  The exact statement needs no new averaging:
\begin{equation}\label{nd:local_ring_actions:eq:formed}
 \boxed{\MGF_{G_n,X_n}=\frac1{|G_n|}\sum_{g\in G_n}\prod_{i,r\ge1}
 (1-t_i^r)^{-\frac1r\sum_{d\mid r}\mu(r/d)|\Fix_{X_n}(g^d)|}.}
\end{equation}
When $X_n=G_n/P_n$, equation \eqref{nd:local_ring_actions:eq:coset-fixed} computes its fixed
points. For fixed $a$, shape, and coefficient degree, each configuration is
supported on a formed submodule of bounded rank. Coefficient stability follows
provided the relevant split unimodular Witt-extension theorem holds over
$R_a$ and every supported isometry extends in the stated tower. Under that
hypothesis, Theorem~\ref{thm:bounded-support-orbit-stability} applies. In
residue characteristic two, the quadratic form itself (not only its polar
form) must be included in the support data. The unverified extension input is
isolated in the final provisional part.

\subsection{Reductive group schemes and adjoint congruence quotients}

Let $\mathcal G$ be a reductive group scheme, $G_a=\mathcal G(R_a)$, and let
$\rho_a$ be any honest complex representation with character $\chi_a$.
Expanding each inverse determinant logarithm gives the class-compressed master
formula
\begin{equation}\label{nd:local_ring_actions:eq:character}
 \boxed{\MGF_{G_a,\rho_a}=\sum_{[g]\subseteq G_a}\frac1{|C_{G_a}(g)|}
 \exp\!\left(\sum_{r\ge1}\frac{\chi_a(g^r)}r p_r\right).}
\end{equation}
If $X_a=G_a/P_a$ is a single parabolic orbit, combine
\eqref{nd:local_ring_actions:eq:coset-fixed} with \eqref{nd:local_ring_actions:eq:universal}.  If
$(\mathcal G/\mathcal P)(R_a)$ has several $G_a$-orbits, the formula must be
applied orbit by orbit; silently treating the set as one coset space would be
incorrect.

Now let $c\le a$ and
$\mathfrak g_c=\Lie(\mathcal G)(R_c)$.  Let $G_a$ act on this finite set by
reduction modulo $\pi^c$ and the adjoint action.  Then
\begin{equation}\label{nd:local_ring_actions:eq:adjoint-fixed}
 \boxed{F_d^{\Ad,c}(g)=\left|\ker\bigl(\Ad(\bar g^d)-I:
 \mathfrak g_c\to\mathfrak g_c\bigr)\right|.}
\end{equation}
If the Smith exponents of the displayed operator are
$\eta_1,\ldots,\eta_N\in\{0,\ldots,c\}$, the kernel has size
$q^{\eta_1+\cdots+\eta_N}$.  Therefore
\begin{equation}\label{nd:local_ring_actions:eq:adjoint-smg}
 \boxed{\MGF_{G_a,\mathfrak g_c}=\frac1{|G_a|}\sum_g\prod_{i,r\ge1}
 (1-t_i^r)^{-\frac1r\sum_{d\mid r}\mu(r/d)
 q^{\sum_j\eta_j(g^d)}}.}
\end{equation}

\subsection{Level growth and the correct asymptotic object}

The rank-stable degree-two coefficients still grow with the level:
\[
 \begin{array}{c|c}
 \text{action}&[m_{(1,1)}]\text{ in stable rank}\\ \hline
 \text{primitive vectors}&q^a\\
 \text{projective points}&a+1\\
 \text{rank-$k$ free summands}&\binom{a+k}{k}.
 \end{array}
\]
Thus no unnormalized coefficientwise limit exists as $a\to\infty$ for these
actions.  The natural replacement is the level-generating series.  With the
formal level-zero convention,
\begin{equation}\label{nd:local_ring_actions:eq:level-series}
 \sum_{a\ge0}q^a u^a=\frac1{1-qu},\qquad
 \sum_{a\ge0}(a+1)u^a=\frac1{(1-u)^2},\qquad
 \sum_{a\ge0}\binom{a+k}{k}u^a=\frac1{(1-u)^{k+1}}.
\end{equation}
For actual rings indexed from $a=1$, subtract the constant term.  The growth
is exponential for primitive pairs and polynomial of degree $1$ or $k$ for
projective points or rank-$k$ summands.

\subsection{Independent computational verification}

The accompanying executable files are
\begin{center}
\path{verification.py},\qquad
\path{marimo_notebooks/local_ring_actions.py},\qquad
\path{test_verification.py}.
\end{center}
The marimo notebook is separated into four group/action sections.  It performs
the following independent checks.
\begin{enumerate}
\item Enumerate matrices over $\mathbb Z/p^a$ and retain exactly those whose
determinant is a unit; retain determinant one for $\Sp_2=\SL_2$.
\item Construct primitive vectors and free rank-one summands as literal finite
sets, then induce the matrix permutations.
\item Count orbits directly on tuples of multisets, without using characters,
fixed-point formulas, or cycle indices.
\item Separately compute the cycle-index coefficient in
\eqref{nd:local_ring_actions:eq:universal}; reconstruct cycles from fixed points of powers using
\eqref{nd:local_ring_actions:eq:mobius}.
\item Recover truncated Smith exponents from exact kernel sizes over every
$R_b$ and compare \eqref{nd:local_ring_actions:eq:primitive-fixed} and
\eqref{nd:local_ring_actions:eq:projective-fixed} with literal fixed points.
\item Check \eqref{nd:local_ring_actions:eq:coset-fixed} from centralizers and conjugacy-class
intersections with the projective parabolic.
\item Build all rank-two free summands of $(\mathbb Z/4)^4$ as unique graph
lifts of planes in $\F_2^4$, then compute stabilizer orbits from parabolic
generators.
\item Construct the adjoint action by explicit conjugation and compare every
fixed-point count with \eqref{nd:local_ring_actions:eq:adjoint-fixed}.
\end{enumerate}

\subsubsection{Exact result ledger}

\begin{center}
\small
\begin{tabular}{@{}lrrrr@{}}
\toprule
Action & $|G|$ & $|X|$ & $\tau$ & direct/formula\\
\midrule
$\GL_2(\mathbb Z/4)$, primitive & 96&12&$(1,1)$&$4/4$\\
$\GL_2(\mathbb Z/4)$, projective &96&6&$(2)$&$3/3$\\
&&&$(1,1)$&$3/3$\\
$\GL_2(\mathbb Z/9)$, primitive &3888&72&$(1,1)$&$9/9$\\
$\GL_2(\mathbb Z/9)$, projective &3888&12&$(2)$&$3/3$\\
&&&$(1,1)$&$3/3$\\
$\Sp_2(\mathbb Z/4)$, isotropic lines &48&6&$(1,1)$&$3/3$\\
$\Sp_4(2)$, isotropic points &720&15&$(2),(1,1)$&$3/3$ each\\
$O_4^+(2)$, singular points &72&9&$(2),(1,1)$&$3/3$ each\\
$\GL_2(\mathbb Z/4)$, $M_2(\F_2)$ adjoint &96&16&$(1),(2),(1,1)$
&$6/6,34/34,56/56$\\
\bottomrule
\end{tabular}
\end{center}

The $\GL_2(\mathbb Z/4)$ Smith, homogeneous-space, and M\"obius checks cover
all $96$ matrices and every relevant power.  For $\GL_2(\mathbb Z/9)$ the
orbit and cycle averages use all $3888$ matrices, while the more expensive
pointwise comparison uses $32$ deterministic representatives.  The rank-two
summand construction contains exactly
\[
 2^{(2-1)2(4-2)}\left.\qbinom{4}{2}\right|_{q=2}=16\cdot35=560
\]
summands.  The stabilizer has exactly six orbits, with Smith types
\[
 (0,0),(0,1),(0,2),(1,1),(1,2),(2,2),
\]
which is the complete list and agrees with $\binom{2+2}{2}=6$.

\resultbox{\textbf{Verification outcome.}  Ten newly tabulated coefficient
comparisons pass exactly.  All fixed-point/cycle reconstructions pass in the
stated exhaustive or representative ranges.  Four automated tests complete
successfully, and the marimo notebook passes its static notebook check.}

\subsection{Scope and reproducibility}

The proof applies to finite chain rings with residue cardinality $q$.  The
executable constructors specialize to $R_a=\mathbb Z/p^a\mathbb Z$, so their
residue cardinality is $q=p$; this is sufficient to test non-field level
behavior at two distinct primes.  The computations prove no untested
all-parameter statement: the universal Reynolds, Smith, coset-counting, and
bounded-support arguments above provide those proofs.

From the workspace root, run
\begin{verbatim}
marimo run marimo_notebooks/local_ring_actions.py
pytest -q lambda_distributions/dists/local_ring_permutation_actions/test_verification.py
\end{verbatim}
The standalone PDF generated from this source is designed to be included
later, without modification, by the proof anthology's \verb|\proofnote|
mechanism.

\begin{thebibliography}{9}
\bibitem{nd:local_ring_actions:grassmann-residue}
Y.~Guo, J.~Guo, and L.~Huang,
\emph{Vector spaces and Grassmann graphs over residue class rings},
\href{https://arxiv.org/abs/1705.04610}{arXiv:1705.04610}.

\bibitem{nd:local_ring_actions:bruhat-local}
U.~Onn, A.~Prasad, and L.~Vaserstein,
\emph{A note on Bruhat decomposition of $GL(n)$ over local principal ideal
rings},
\href{https://arxiv.org/abs/math/0506094}{arXiv:math/0506094}.

\bibitem{nd:local_ring_actions:witt-extension}
J.-P.~Tignol,
\emph{Witt's extension theorem for quadratic spaces over semiperfect rings},
\href{https://arxiv.org/abs/1408.0522}{arXiv:1408.0522}.
\end{thebibliography}

\clearpage
\section[PSL(2,q) on the projective line]{$PSL_2(q)$ on the projective line}\label{proof:psl2}
\begingroup
\definecolor{ink}{HTML}{172337}
\definecolor{accent}{HTML}{8B4A35}
\definecolor{pale}{HTML}{FAF1ED}
\definecolor{passgreen}{HTML}{147A4A}
\providecommand{\Sym}{\operatorname{Sym}}
\renewcommand{\Sym}{\operatorname{Sym}}
\providecommand{\Exp}{\operatorname{Exp}_{\sigma}}
\renewcommand{\Exp}{\operatorname{Exp}_{\sigma}}
\providecommand{\PASS}{\textcolor{passgreen}{\textbf{PASS}}}
\renewcommand{\PASS}{\textcolor{passgreen}{\textbf{PASS}}}
\providecommand{\summarybox}[1]{\begin{center}\fcolorbox{accent}{pale}{\parbox{0.91\linewidth}{#1}}\end{center}}
\renewcommand{\summarybox}[1]{\begin{center}\fcolorbox{accent}{pale}{\parbox{0.91\linewidth}{#1}}\end{center}}
\setcounter{theorem}{0}
\setcounter{proposition}{0}
\setcounter{lemma}{0}
\setcounter{corollary}{0}
\setcounter{remark}{0}
\summarybox{\textbf{Outcome.} Formula (3.3) in the supplied statement is correct.  Prime-field groups $q=2,3,5,7$ were constructed from all determinant-one $2\times2$ matrices and then quotiented by their projective action.  All cycle-type multiplicities, $48$ coefficients through degree four, and direct orbit counts agree.  Determinant-level error is at most $6.67\times10^{-16}$.}

\subsection{Representation and coefficient meaning}
Let $q=p^f$, $d=\gcd(2,q-1)$, and $G=PSL_2(q)$, so $|G|=q(q^2-1)/d$.  Let $V=\mathbb C[\mathbf P^1(\mathbf F_q)]$, of dimension $q+1$.  For a permutation with $c_m$ cycles of length $m$,
\[
 \det(I-tg)^{-1}=\prod_m(1-t^m)^{-c_m},\qquad
 \prod_j\det(I-t_jg)^{-1}=\Exp\left(\sum_m c_mp_m\right). \tag{1}
\]
Tools~\ref{tool:coefficient} and~\ref{tool:permutation-orbits} identify the
coefficient at $\tau$ with
$\dim(\bigotimes_s\Sym^{\tau_s}V)^G$ and with the number of $G$-orbits on
tuples of multisets having sizes $\tau_s$.  The latter gives a genuinely
independent combinatorial check of the class calculation.

\subsection{Four element types}
Put $L_-=(q-1)/d$ and $L_+=(q+1)/d$.
\begin{center}\begin{tabularx}{\linewidth}{@{}lXl@{}}\toprule
type & cycle structure on $\mathbf P^1$ & multiplicity\\\midrule
identity & $1^{q+1}$ & $1$\\
nonidentity unipotent & $1^1p^{q/p}$ & $q^2-1$\\
split semisimple, order $m>1$ & $1^2m^{(q-1)/m}$ & $\frac{q(q+1)}2\varphi(m)$, $m\mid L_-$\\
nonsplit semisimple, order $m>1$ & $m^{(q+1)/m}$ & $\frac{q(q-1)}2\varphi(m)$, $m\mid L_+$\\\bottomrule
\end{tabularx}\end{center}
The unipotent form $x\mapsto x+1$ fixes infinity and partitions the affine line into $q/p$ cycles.  A split element has two eigenlines; a nonsplit element has none.  Each nontrivial semisimple element determines its eigenline pair, hence its unique torus.  The split and nonsplit torus counts are $q(q+1)/2$ and $q(q-1)/2$ respectively, giving the displayed multiplicities.

\subsection{Exact formula}
Substituting the four rows into (1) proves
\begin{align}
\MGF_{G,V}=\frac1{|G|}\Bigg[&\Exp((q+1)p_1)
+(q^2-1)\Exp\left(p_1+\frac qp p_p\right)\notag\\
&+\frac{q(q+1)}2\sum_{\substack{m\mid L_-\\m>1}}\varphi(m)
 \Exp\left(2p_1+\frac{q-1}{m}p_m\right)\notag\\
&+\frac{q(q-1)}2\sum_{\substack{m\mid L_+\\m>1}}\varphi(m)
 \Exp\left(\frac{q+1}{m}p_m\right)\Bigg]. \tag{2}
\end{align}
The multiplicities sum to $|G|$, an immediate global consistency check.

\subsection{Poisson comparison and the exact boundary}
Two-transitivity forces the same orbit counts as the stable symmetric-group action through total degree two:
\[
 [m_{(1)}]=1,\qquad [m_{(2)}]=[m_{(1,1)}]=2.
\]
For ordered triples, equality patterns contribute $1+3$ orbits, while distinct triples contribute
\[
 \frac{(q+1)q(q-1)}{|G|}=d.
\]
Thus
\begin{equation}[m_{(1,1,1)}]=4+d=\begin{cases}5&q\text{ even},\\6&q\text{ odd}.\end{cases} \tag{3}\end{equation}
At degree four, equality patterns with $1,2,3,4$ distinct values give
\begin{equation}[m_{(1,1,1,1)}]=1+7+6d+d(q-2)=8+d(q+4). \tag{4}\end{equation}
The last term is the projective cross-ratio parameter.  Consequently the raw sigma--MGF has no finite coefficientwise limit as $q\to\infty$.

\subsection{Verification strategy}
For each prime $q$ tested, the implementation enumerates every
$(a,b,c,d)\in\mathbf F_q^4$ with $ad-bc=1$, induces the fractional-linear
permutation on $\mathbf F_q\cup\{\infty\}$, and removes scalar-kernel
duplicates.  It compares the actual cycle-type histogram with the four-type
theorem, extracts every coefficient through degree four, independently counts
orbits on tuples of multisets, and finally compares determinant averages of
the permutation matrices at $(t_1,t_2)=(0.07,0.11)$.
The construction is intentionally transparent and is presently implemented for prime fields.  The proof and formula cover all prime powers; testing $f>1$ would require a small finite-field layer, not a change to the group-theoretic argument.
\endgroup


\clearpage
\section[Delta(3n-squared) and Delta(6n-squared) monomial SU(3) groups]{$\Delta(3n^2)$ and $\Delta(6n^2)$ monomial $SU(3)$ groups}\label{proof:delta-su3}
\begingroup
\definecolor{ink}{HTML}{172337}
\definecolor{accent}{HTML}{594A9A}
\definecolor{pale}{HTML}{F1EFF8}
\definecolor{passgreen}{HTML}{147A4A}
\providecommand{\Sym}{\operatorname{Sym}}
\renewcommand{\Sym}{\operatorname{Sym}}
\providecommand{\Exp}{\operatorname{Exp}_{\sigma}}
\renewcommand{\Exp}{\operatorname{Exp}_{\sigma}}
\providecommand{\PASS}{\textcolor{passgreen}{\textbf{PASS}}}
\renewcommand{\PASS}{\textcolor{passgreen}{\textbf{PASS}}}
\providecommand{\summarybox}[1]{\begin{center}\fcolorbox{accent}{pale}{\parbox{0.91\linewidth}{#1}}\end{center}}
\renewcommand{\summarybox}[1]{\begin{center}\fcolorbox{accent}{pale}{\parbox{0.91\linewidth}{#1}}\end{center}}
\setcounter{theorem}{0}
\setcounter{proposition}{0}
\setcounter{lemma}{0}
\setcounter{corollary}{0}
\setcounter{remark}{0}
\summarybox{\textbf{Outcome.} Both proposed formulas are correct.  The standard $3$-dimensional matrices were built for $n=2,3,4,5$ (orders $12$--$150$).  All $240$ coefficients through total degree six agree with the exact $Q_n/R_n$ formulas.  Eight determinant-level comparisons have maximum error $8.89\times10^{-16}$.}

\subsection{Groups and imported sigma-MGF target}
Put $\eta=e^{2\pi i/n}$ and
\[
 N_n=\{D(a,b)=\operatorname{diag}(\eta^a,\eta^b,\eta^{-a-b}):a,b\in\mathbb Z/n\}.
\]
The coefficient and determinant identities are Tools~\ref{tool:coefficient}
and~\ref{tool:molien}; the new content is the coset spectral calculation.

Let $P$ be the permutation matrix of a $3$-cycle.  Then
$\Delta(3n^2)=N_n\rtimes\langle P\rangle$.  For $\pi\in S_3$ use the determinant-one lift $\widetilde P_\pi=\operatorname{sgn}(\pi)P_\pi$; then
$\Delta(6n^2)=N_n\rtimes\widetilde S_3$.

\subsection{Coset proof}
Define
\[
 \mathcal A_n=\frac1{n^2}\sum_{a,b}\prod_j
 \frac1{(1-\eta^at_j)(1-\eta^bt_j)(1-\eta^{-a-b}t_j)}. \tag{2}
\]
For $DP$ or $DP^2$, the three entries around the monomial $3$-cycle multiply to $\det D=1$, so $\det(I-tDP)=1-t^3$.  Therefore
\begin{equation}
 \boxed{\MGF_{\Delta(3n^2)}=\tfrac13\mathcal A_n+\tfrac23\Exp(p_3).} \tag{3}
\end{equation}

For a transposition coset in the determinant-one lift, let $\alpha$ be the product around the $2$-cycle and $\beta$ the fixed-coordinate entry.  Since the underlying permutation has determinant $-1$, $-\alpha\beta=1$ and $\beta=-\alpha^{-1}$.  Thus
\[
 \det(I-tg)=(1-\alpha t^2)(1+\alpha^{-1}t).
\]
As $D(a,b)$ varies, $\alpha$ is uniform in $\mu_n$.  Put
\[
 \mathcal B_n=\frac1n\sum_{a\bmod n}\prod_j
 \frac1{(1-\eta^at_j^2)(1+\eta^{-a}t_j)}. \tag{4}
\]
The identity, $3$-cycle, and transposition proportions in $S_3$ are $1/6$, $1/3$, and $1/2$.  Hence
\begin{equation}
 \boxed{\MGF_{\Delta(6n^2)}=\tfrac16\mathcal A_n+\tfrac13\Exp(p_3)+\tfrac12\mathcal B_n.} \tag{5}
\end{equation}

\subsection{Exact coefficient formulas}
For $\tau=(\tau_1,\ldots,\tau_\ell)$, let $Q_n(\tau)$ count rows
$(x_s,y_s,z_s)$ with $x_s+y_s+z_s=\tau_s$ and
$\sum_sx_s\equiv\sum_sy_s\equiv\sum_sz_s\pmod n$.  Root-of-unity orthogonality gives $[m_\tau]\mathcal A_n=Q_n(\tau)$.

Let
\[
 R_n(\tau)=\sum_{\substack{x_s+2y_s=\tau_s\\ \sum y_s\equiv\sum x_s\ (n)}}(-1)^{\sum x_s}.
\]
This follows by expanding $(1+\alpha^{-1}t)^{-1}$ and $(1-\alpha t^2)^{-1}$.  Finally, $[m_\tau]\Exp(p_3)$ is $1$ exactly when every $\tau_s$ is divisible by $3$.  Therefore
\begin{align}
 [m_\tau]\MGF_{\Delta(3n^2)}&=\tfrac13Q_n(\tau)+\tfrac23\mathbf1_{3\mid\tau_s\ \forall s}, \tag{6}\\
 [m_\tau]\MGF_{\Delta(6n^2)}&=\tfrac16Q_n(\tau)+\tfrac13\mathbf1_{3\mid\tau_s\ \forall s}+\tfrac12R_n(\tau). \tag{7}
\end{align}
Their integrality is automatic from (1), even though (6)--(7) display rational weights.

\subsection{Verification design and results}
For each $n$, all matrices $D(a,b)P^k$ and $D(a,b)\widetilde P_\pi$ are constructed.  The code checks group size and determinant one.  It then compares: (i) eigenvalue-based symmetric-power character averages; (ii) exact enumeration of $Q_n$ and $R_n$; and (iii) the direct determinant average at $(0.07,0.11)$ against (3) or (5).
\begin{center}\begin{tabular}{@{}rrrrrrc@{}}\toprule
$n$ & $|\Delta(3n^2)|$ & $|\Delta(6n^2)|$ & degree & checks/family & max. error & result\\\midrule
2&12&24&$0$--$6$&30&$4.45\times10^{-16}$&\PASS\\
3&27&54&$0$--$6$&30&$6.67\times10^{-16}$&\PASS\\
4&48&96&$0$--$6$&30&$6.67\times10^{-16}$&\PASS\\
5&75&150&$0$--$6$&30&$8.89\times10^{-16}$&\PASS\\\bottomrule
\end{tabular}\end{center}
There are two families for each row, hence $2\cdot4\cdot30=240$ exact coefficient checks.  As useful spot checks, for $n=2$ the coefficients at $\tau=(3),(2,1),(1,1,1)$ are $(1,1,2)$ for $\Delta(12)$ and $(0,0,1)$ for $\Delta(24)$.

\subsection{Novel diagnostic}
The coset decomposition is also a debugging tool.  A mismatch isolated to $\mathcal B_n$ almost always signals the missed minus sign in the determinant-one lift of a transposition.  A mismatch isolated to $\Exp(p_3)$ signals an incorrect monomial-cycle product.  This three-channel comparison is stronger than a single ordinary Molien specialization.

\subsection*{Reproduction}
\begin{verbatim}
.venv/bin/python -m \
  for_this_guy.generalized_molien_extensions.delta_su3.verification
\end{verbatim}
\endgroup


\clearpage
\section{Code--monomial matrix groups}\label{proof:code-monomial}
\begingroup
\definecolor{ink}{HTML}{172337}
\definecolor{accent}{HTML}{255C73}
\definecolor{pale}{HTML}{EEF5F7}
\definecolor{passgreen}{HTML}{147A4A}
\providecommand{\Sym}{\operatorname{Sym}}
\renewcommand{\Sym}{\operatorname{Sym}}
\providecommand{\Cyc}{\operatorname{Cyc}}
\renewcommand{\Cyc}{\operatorname{Cyc}}
\providecommand{\PASS}{\textcolor{passgreen}{\textbf{PASS}}}
\renewcommand{\PASS}{\textcolor{passgreen}{\textbf{PASS}}}
\providecommand{\summarybox}[1]{\begin{center}\fcolorbox{accent}{pale}{\parbox{0.91\linewidth}{#1}}\end{center}}
\renewcommand{\summarybox}[1]{\begin{center}\fcolorbox{accent}{pale}{\parbox{0.91\linewidth}{#1}}\end{center}}
\setcounter{theorem}{0}
\setcounter{proposition}{0}
\setcounter{lemma}{0}
\setcounter{corollary}{0}
\setcounter{remark}{0}
\summarybox{\textbf{Outcome.} The proposed weighted-cycle formula and the dual-code orbit formula are correct under the stated prime-alphabet convention.  Explicit matrices for three groups of orders $8$, $24$, and $18$ were used to check all $57$ coefficients through total degree five.  All agree.  Three independent determinant evaluations agree to machine precision (maximum error $0$).}

\subsection{Imported coefficient identity}
For the code--monomial module $V$, put $\MGF_{G_C,V}=\smgv{G_C}{V}$.
The determinant and coefficient identities are imported; the new theorem is
the cycle-decorated dual-code filter.

\subsection{Code--monomial theorem}
Let $r$ be prime, $C\leq\mathbf F_r^n$, and let
\[
H\leq\operatorname{Aut}_{\rm perm}(C):=
\{\pi\in S_n:\pi(C)=C\}
\]
be a coordinate-permutation automorphism group.  Put
$\zeta=e^{2\pi i/r}$ and define
\[
 D_c=\operatorname{diag}(\zeta^{c_1},\ldots,\zeta^{c_n}),\qquad G_C=D_C\rtimes H.
\]
Use $P_\pi e_i=e_{\pi(i)}$.  On a coordinate cycle $O$ of $\pi$, the weighted shift $D_cP_\pi$ returns after $|O|$ steps with scalar $\zeta^{c(O)}$, where $c(O)=\sum_{i\in O}c_i$.  Therefore
\begin{equation}
 \det(I-tD_cP_\pi)=\prod_{O\in\Cyc(\pi)}(1-\zeta^{c(O)}t^{|O|}). \tag{2}
\end{equation}
Substitution into the master average proves
\begin{equation}
 \boxed{\MGF_{G_C,V}(\mathbf t)=\frac1{|C||H|}\sum_{\pi\in H}\sum_{c\in C}
 \prod_{j,O}(1-\zeta^{c(O)}t_j^{|O|})^{-1}.} \tag{3}
\end{equation}

\subsubsection{Why the coefficient is an orbit count}
A monomial basis of $\bigotimes_s\Sym^{\tau_s}V$ is indexed by arrays $A=(a_{si})\geq0$ with $\sum_i a_{si}=\tau_s$.  Set
\[
 w(A)_i=\sum_s a_{si}\pmod r.
\]
The diagonal matrix $D_c$ multiplies the basis vector by $\zeta^{c\cdot w(A)}$.  Character orthogonality says it survives all diagonal matrices precisely when $w(A)\in C^\perp$.  The remaining group $H$ permutes columns, so invariants in this permutation module are constant on $H$-orbits.  Hence
\begin{equation}
 \boxed{[m_\tau]\MGF_{G_C,V}
 =\#\bigl(H\backslash\{A:\sum_i a_{si}=\tau_s,\ w(A)\in C^\perp\}\bigr).} \tag{4}
\end{equation}
This also proves integrality without relying on cancellation in (3).

The proof uses only coordinate permutations and therefore does not cover the
full monomial or semilinear automorphism group of a code.  If coordinate
scalings or field automorphisms are included, their characters must be
incorporated into the diagonal filter and the column-orbit argument must be
replaced by the corresponding semilinear action.  The verifier rejects an
automorphism entry that is not a coordinate permutation.

\subsection{Independent verification method}
The implementation uses routes that share the group input but not the coefficient algorithm:
\begin{enumerate}[leftmargin=*]
\item Enumerate every complex matrix $D_cP_\pi$.  Obtain $\chi_{\Sym^dV}(g)$ from the eigenvalue series $\prod_i(1-t\lambda_i)^{-1}$ and average products of characters.
\item Enumerate exponent arrays, impose $c\cdot w(A)=0$ for every $c\in C$, and count column-permutation orbits.  No eigenvalues occur in this route.
\item At $(t_1,t_2)=(0.07,0.11)$ compare the direct matrix determinant average with (3), evaluated solely from weighted coordinate cycles.
\end{enumerate}

\subsection{Representative cases and results}
Every partition of every total degree $0$ through $5$ was checked (19 coefficients per group).
\begin{center}\begin{tabularx}{\linewidth}{@{}Xrrrrc@{}}\toprule
Case & $\dim V$ & $|C|$ & $|H|$ & checks & result\\\midrule
$W(B_2)$, full binary code &2&4&2&19&\PASS\\
$W(D_3)$, even binary code &3&4&6&19&\PASS\\
ternary repetition code $\rtimes S_3$ &3&3&6&19&\PASS\\\bottomrule
\end{tabularx}\end{center}
\begin{center}\begin{tabular}{@{}lrrrr@{}}\toprule
Case & $(2)$ & $(3)$ & $(2,1)$ & $(1,1,1)$\\\midrule
$W(B_2)$ &1&0&0&0\\
$W(D_3)$ &1&1&1&1\\
ternary repetition &0&3&4&5\\\bottomrule
\end{tabular}\end{center}
The numerical generalized-Molien values by both routes were
\[1.025494342880870,\qquad 1.028638382342067,\qquad
1.010621973731615.\]

\subsection{Scope and reusable tool}
For $\mathbf F_{p^f}$ with $f>1$, the notation $c\mapsto\zeta^c$ is not a faithful additive embedding of the whole additive group into one cyclic root group.  Formula (4) remains valid after replacing $C^\perp$ by the annihilator of the chosen diagonal-character group.  The reusable computational idea is to test (3) numerically while testing (4) exactly: agreement localizes errors to either the weighted-cycle determinant or the character-orthogonality filter.
\endgroup


\clearpage
\subsection{Extended equivariant code-monomial theory}\label{proof:new-dist:code_monomial}

\resultbox{\textbf{Canonical testing artifacts.}\par\smallskip Exact backend: \path{lambda_distributions/dists/code_monomial_sigma_mgfs/verification.py}.\par Regression: \path{lambda_distributions/dists/code_monomial_sigma_mgfs/test_verification.py}.\par Marimo: \path{marimo_notebooks/code_monomial.py}.}

\subsubsection*{Scope and outcome}

Let $C\leq(\Z/r\Z)^N$ be an additive code and let a coordinate-permutation
group $P\leq S_N$ preserve $C$.  We construct the honest unitary monomial
group $G(C,P)=D_C\rtimes P\leq U(N)$ and prove two equivalent formulas for
its complete generalized Molien, or $\sigma$-MGF, series: a weighted-cycle
average over $C\times P$ and a root-filtered sum over dual codewords that are
constant on permutation cycles.  We also prove the exact coefficient-orbit
interpretation, the complete-weight-enumerator specialization for $P=1$, and
the resulting coefficientwise convergence criterion and orbit-growth
obstruction.  An accompanying marimo notebook constructs five groups of
dimensions $2$ through $4$ and orders $4,8,24,18,64$.  All $95$ coefficients
through total degree five agree exactly; three independent evaluations of the
full series agree to $4.45\times10^{-16}$.

\subsubsection{Representation and target quantity}

Fix an integer $r\geq2$ and write $R=\Z/r\Z$ and
$\zeta=e^{2\pi i/r}$.  Let $C\leq R^N$ be an additive subgroup.  For
$c=(c_1,\ldots,c_N)\in C$, define
\[
 D(c)=\operatorname{diag}(\zeta^{c_1},\ldots,\zeta^{c_N}),
 \qquad D_C=\{D(c):c\in C\}.
\]
Let $P\leq S_N$ preserve $C$ under coordinate permutation, and use the
convention $P_\pi e_j=e_{\pi(j)}$.  Then conjugation by $P_\pi$ permutes the
diagonal entries of $D(c)$ and preserves $D_C$.  Hence
\[
 G(C,P)=\{D(c)P_\pi:c\in C,\ \pi\in P\}=D_C\rtimes P\leq U(N).
\]
The parametrization by $C\times P$ is injective, so $|G(C,P)|=|C||P|$.

For any finite $G\leq U(V)$ define
\begin{equation}\label{nd:code_monomial:eq:target}
 \MGF_{G,V}(\mathbf t)
 =\frac1{|G|}\sum_{g\in G}\prod_{i\geq1}\det(I-t_i g)^{-1}.
\end{equation}
For a partition $\tau=(\tau_1,\ldots,\tau_s)$, Tool~\ref{tool:coefficient}
gives
\begin{equation}\label{nd:code_monomial:eq:coefficient-target}
 [m_\tau]\,\smgv{G}{V}
 =\dim\left(\bigotimes_{\alpha=1}^s\Sym^{\tau_\alpha}V\right)^G.
\end{equation}
The new work is the dual-code and cycle-decorated evaluation of this imported
coefficient.

\subsubsection{Weighted-cycle and equivariant dual-code formulas}

For $a\in R$, define the formal series
\begin{equation}\label{nd:code_monomial:eq:H-def}
 H_a(\mathbf t)=
 \sum_{\substack{\nu_1,\nu_2,\ldots\geq0\\
                  \sum_i\nu_i\equiv a\pmod r}}
 \prod_i t_i^{\nu_i}.
\end{equation}
Character orthogonality on $R$ gives the root filter
\begin{equation}\label{nd:code_monomial:eq:H-filter}
 \boxed{H_a(\mathbf t)=\frac1r\sum_{s=0}^{r-1}
 \zeta^{-as}\prod_i(1-\zeta^s t_i)^{-1}.}
\end{equation}
For a positive integer $\ell$, put
$H_a^{[\ell]}(\mathbf t)=H_a(t_1^\ell,t_2^\ell,\ldots)$, and define
\[
 C^\perp=\{u\in R^N:\langle c,u\rangle=0\text{ in }R
                         \text{ for every }c\in C\}.
\]

\begin{theorem}[Equivariant dual-code $\sigma$-MGF]\label{nd:code_monomial:thm:main}
For $G(C,P)$ on $V=\C^N$,
\begin{align}
 \MGF_{C,P}(\mathbf t)
 &=\frac1{|C||P|}\sum_{\pi\in P}\sum_{c\in C}
   \prod_{i\geq1}\prod_{Q\in\Cyc(\pi)}
   \left(1-\zeta^{\sum_{j\in Q}c_j}t_i^{|Q|}\right)^{-1},
   \label{nd:code_monomial:eq:weighted-cycle}\\
 &=\frac1{|P|}\sum_{\pi\in P}
   \sum_{\substack{u\in C^\perp\\
                    u\text{ constant on every cycle of }\pi}}
   \prod_{Q\in\Cyc(\pi)}H_{u_Q}^{[|Q|]}(\mathbf t),
   \label{nd:code_monomial:eq:dual-cycle}
\end{align}
where $u_Q$ is the common value of $u$ on $Q$.
\end{theorem}

\begin{proof}
Consider a cycle $Q=(j_1\,j_2\,\cdots\,j_\ell)$ of $\pi$.  On its
coordinate subspace, $D(c)P_\pi$ is a weighted cyclic shift.  Its $\ell$th
power is scalar multiplication by
$\zeta^{c_{j_1}+\cdots+c_{j_\ell}}$, and therefore
\[
 \det(I-tD(c)P_\pi\mid\C^Q)
 =1-\zeta^{\sum_{j\in Q}c_j}t^\ell.
\]
Multiplication over the disjoint cycles, over the variables $t_i$, and then
averaging over $C\times P$ proves \eqref{nd:code_monomial:eq:weighted-cycle}.

For a fixed cycle $Q$, expand its factor as
\[
 \prod_i\left(1-\zeta^{\sum_{j\in Q}c_j}t_i^{|Q|}\right)^{-1}
 =\sum_{\nu_1,\nu_2,\ldots\geq0}
  \zeta^{(\sum_i\nu_i)(\sum_{j\in Q}c_j)}
  \prod_i t_i^{|Q|\nu_i}.
\]
Assign to every $j\in Q$ the residue
$u_j=\sum_i\nu_i\pmod r$.  The resulting word $u$ is constant on every
cycle of $\pi$, and the total phase is $\zeta^{\langle c,u\rangle}$.
Orthogonality on the finite additive group $C$ gives
\[
 \frac1{|C|}\sum_{c\in C}\zeta^{\langle c,u\rangle}
 =\begin{cases}1,&u\in C^\perp,\\0,&u\notin C^\perp.\end{cases}
\]
After this filter, the sum of the allowed exponents on $Q$ is precisely
$H_{u_Q}^{[|Q|]}$.  Multiplying over cycles proves
\eqref{nd:code_monomial:eq:dual-cycle}.  This orthogonality argument is valid over
$\Z/r\Z$ for composite as well as prime $r$.
\end{proof}

\subsubsection{Pure diagonal groups and complete weight enumerators}

When $P=1$, every coordinate is a one-cycle.  Define the complete weight
enumerator of a code $A\leq R^N$ by
\[
 \cwe_A(X_0,\ldots,X_{r-1})
 =\sum_{u\in A}\prod_{j=1}^N X_{u_j}.
\]

\begin{corollary}[Pure diagonal specialization]\label{nd:code_monomial:cor:cwe}
For the diagonal group $D_C$,
\begin{equation}\label{nd:code_monomial:eq:cwe}
 \boxed{\MGF_{D_C}=\cwe_{C^\perp}(H_0,\ldots,H_{r-1}).}
\end{equation}
For $r=2$ this becomes
\[
 H_0=\tfrac12\left[\prod_i(1-t_i)^{-1}+\prod_i(1+t_i)^{-1}\right],
 \quad
 H_1=\tfrac12\left[\prod_i(1-t_i)^{-1}-\prod_i(1+t_i)^{-1}\right],
\]
so $\MGF_{D_C}=W_{C^\perp}(H_0,H_1)$ for the ordinary binary weight
enumerator $W$.
\end{corollary}

\begin{proof}
Set $P=1$ in \eqref{nd:code_monomial:eq:dual-cycle}.  The contribution of a fixed
$u\in C^\perp$ is $\prod_jH_{u_j}$, exactly the corresponding monomial in
the complete weight enumerator.
\end{proof}

This is the first verification family below.  It isolates the diagonal
character filter before any coordinate-orbit quotient is imposed.

\subsubsection{Exact coefficient formula}

For $\tau=(\tau_1,\ldots,\tau_s)$, let $\mathcal E_\tau(C^\perp)$ be the
set of matrices $A=(a_{\alpha j})\in\Z_{\geq0}^{s\times N}$ satisfying
\[
 \sum_{j=1}^N a_{\alpha j}=\tau_\alpha\quad(1\leq\alpha\leq s),
 \qquad
 \left(\sum_\alpha a_{\alpha1},\ldots,
       \sum_\alpha a_{\alpha N}\right)\bmod r\in C^\perp.
\]
The group $P$ acts on these arrays by permuting columns.

\begin{theorem}[Dual-code configuration orbits]\label{nd:code_monomial:thm:orbits}
For every partition $\tau$,
\begin{equation}\label{nd:code_monomial:eq:orbit-count}
 \boxed{[m_\tau]\MGF_{C,P}
 =\#\bigl(P\backslash\mathcal E_\tau(C^\perp)\bigr).}
\end{equation}
\end{theorem}

\begin{proof}
A monomial basis of
$\bigotimes_\alpha\Sym^{\tau_\alpha}\C^N$ is indexed by the exponent
arrays with the stated row sums.  The matrix $D(c)$ multiplies the basis
vector indexed by $A$ by
\[
 \zeta^{\sum_jc_j\sum_\alpha a_{\alpha j}}.
\]
It is fixed by every $D(c)$ exactly when the vector of total coordinate
exponents lies in $C^\perp$ modulo $r$.  Thus the $D_C$-invariant monomial
basis is indexed by $\mathcal E_\tau(C^\perp)$.  The quotient group $P$
permutes this basis by permuting columns, and its orbit sums form a basis of
the invariant space.  Equation \eqref{nd:code_monomial:eq:coefficient-target} completes the
proof.
\end{proof}

\subsubsection{Coefficientwise asymptotics}

Consider a sequence $(C_N,P_N)$ with a fixed modulus $r$.  The exact orbit
formula immediately gives
\begin{equation}\label{nd:code_monomial:eq:criterion}
 \MGF_{C_N,P_N}\text{ converges coefficientwise}
 \quad\Longleftrightarrow\quad
 \#\bigl(P_N\backslash\mathcal E_\tau(C_N^\perp)\bigr)
 \text{ is eventually constant for every fixed }\tau.
\end{equation}
Here ordinary convergence of a nonnegative integer sequence is equivalent to
eventual constancy.

There is a code-independent obstruction.  In each copy of $\Sym^r\C^N$,
the vectors $x_j^r$ are fixed by every $D(c)$ because $\zeta^{rc_j}=1$.
Consequently, for $\tau=(r,\ldots,r)$ with $s$ parts,
\begin{equation}\label{nd:code_monomial:eq:obstruction}
 \boxed{[m_{(r^s)}]\MGF_{C_N,P_N}
 \geq\#\bigl(P_N\backslash[N]^s\bigr).}
\end{equation}
Thus bounded orbit counts on every fixed Cartesian power of the coordinate
set are necessary for a raw coefficientwise limit.  Two immediate examples
are useful:
\begin{enumerate}
\item If $P_N=1$, then $[m_{(r)}]\MGF\geq N$, so every length-growing pure
diagonal family diverges in a fixed degree.
\item If $P_N$ is the regular cyclic shift group, then it has $N$ orbits on
ordered coordinate pairs, so $[m_{(r,r)}]\MGF\geq N$.
\end{enumerate}
Weight-enumerator convergence alone cannot control these coordinate-orbit
counts.  Conversely, \eqref{nd:code_monomial:eq:criterion} identifies the extra information
needed: bounded-degree \emph{equivariant} dual-code configurations.

\subsubsection{Independent verification strategy}

The executable verifier uses four logically separated calculations.
\begin{enumerate}
\item \textbf{Matrix construction.}  It creates every dense complex matrix
$D(c)P_\pi$, checks the expected group cardinality and distinctness, and
checks $g^*g=I$.
\item \textbf{Direct full-series value.}  At
$(t_1,t_2)=(0.07,0.11)$ it evaluates \eqref{nd:code_monomial:eq:target} using
\texttt{numpy.linalg.det} on the constructed matrices.
\item \textbf{Two full-series comparators.}  It evaluates
\eqref{nd:code_monomial:eq:weighted-cycle} from codeword cycle sums and
\eqref{nd:code_monomial:eq:dual-cycle} by enumerating $C^\perp$, testing cycle constancy, and
evaluating \eqref{nd:code_monomial:eq:H-filter}.
\item \textbf{Coefficient comparison.}  For the matrix route it obtains
$\operatorname{Tr}(\Sym^d g)$ from traces of actual matrix powers using
Newton's recurrence
\[
 d\,h_d(g)=\sum_{k=1}^d\operatorname{Tr}(g^k)h_{d-k}(g),\qquad h_0=1,
\]
then averages products over the matrix group.  The independent exact route
enumerates exponent arrays, imposes the dual-code condition, and counts
$P$-orbits as in Theorem~\ref{nd:code_monomial:thm:orbits}.
\end{enumerate}
The proof is general; these computations are regression tests, not a finite
substitute for the proof.

\subsubsection{Verification organized by group}

Every partition of every total degree from $0$ through $5$ was tested: $19$
coefficients for each of five groups, or $95$ exact comparisons in total.

\paragraph{Pure diagonal group}

For $C=(\Z/2\Z)^2$ and $P=1$, the group has order $4$ on $\C^2$.  This
case separately tests Corollary~\ref{nd:code_monomial:cor:cwe}; in particular the degree-two
coefficient is $2$ before the coordinate swap identifies the two squares.

\paragraph{Binary semidirect groups}

The full binary code of length two with $P=S_2$ produces the order-$8$
signed permutation group $W(B_2)$.  The even binary code of length three with
$P=S_3$ produces the order-$24$ group $W(D_3)$.  These test nontrivial cycle
types and dual codes of different dimensions.

\paragraph{Ternary semidirect group}

For the ternary repetition code
$C=\{(a,a,a):a\in\Z/3\Z\}$ and $P=S_3$, the constructed group has order
$18$ on $\C^3$.  Its nonzero coefficients occur in a substantially different
degree pattern from the binary groups.

\paragraph{Composite cyclic-alphabet group}

For
\[
 C=\{(a,b,a,b):a,b\in\Z/4\Z\}\leq(\Z/4\Z)^4
\]
and the cyclic coordinate group generated by $(1\,2\,3\,4)$, the constructed
group has order $64$ on $\C^4$.  This case directly tests the composite
modulus asserted in Theorem~\ref{nd:code_monomial:thm:main}.

\begin{center}
\begin{tabularx}{\linewidth}{@{}Xrrrrrc@{}}
\toprule
Group & $r$ & $N$ & $|C|$ & $|P|$ & $|G|$ & result\\
\midrule
$D_{(\Z/2)^2}$, pure diagonal &2&2&4&1&4&\textcolor{teal}{\textbf{PASS}}\\
$W(B_2)$ &2&2&4&2&8&\textcolor{teal}{\textbf{PASS}}\\
$W(D_3)$ &2&3&4&6&24&\textcolor{teal}{\textbf{PASS}}\\
ternary repetition $\rtimes S_3$ &3&3&3&6&18&\textcolor{teal}{\textbf{PASS}}\\
quaternary pair code $\rtimes C_4$ &4&4&16&4&64&\textcolor{teal}{\textbf{PASS}}\\
\bottomrule
\end{tabularx}
\end{center}

Selected coefficients are shown below; all remaining partitions through
degree five are checked by the notebook and test suite.
\begin{center}
\begin{tabularx}{\linewidth}{@{}Xrrrrrr@{}}
\toprule
Group & $(1)$ & $(2)$ & $(3)$ & $(2,1)$ & $(1,1,1)$ & $(4)$\\
\midrule
$D_{(\Z/2)^2}$ &0&2&0&0&0&3\\
$W(B_2)$ &0&1&0&0&0&2\\
$W(D_3)$ &0&1&1&1&1&2\\
ternary repetition $\rtimes S_3$ &0&0&3&4&5&0\\
quaternary pair $\rtimes C_4$ &0&0&0&0&0&3\\
\bottomrule
\end{tabularx}
\end{center}

At the full-series test point the direct matrix values are respectively
\[
 1.0507589401246002,\quad
 1.0254943428808700,\quad
 1.0286383823420668,\quad
 1.0106219737316153,\quad
 1.0013332933694528.
\]
Both formula routes reproduce every value.  Across all groups, the maximum
full-series discrepancy is $4.45\times10^{-16}$, the maximum coefficient
rounding residual is $2.23\times10^{-14}$, and the maximum unitarity residual
is $2.23\times10^{-16}$.

\subsubsection{Scope and reproduction}

The literal scalar model uses the cyclic additive alphabet $\Z/r\Z$.  For a
noncyclic alphabet such as the additive group of $\mathbf F_{p^s}$, there is
no faithful map of the entire alphabet into a single cyclic root group by
$a\mapsto\zeta^a$.  One must instead choose the appropriate family of
additive characters, character module, or Clifford--Weil representation.
Likewise, a monomial automorphism involving coordinate scalings, or a
semilinear automorphism involving a field automorphism, is not merely a
coordinate permutation in $P$ and needs a correspondingly generalized
cycle calculation.

It is also useful to separate three levels of code information.  For a pure
diagonal binary group, the ordinary Hamming weight enumerator of $C^\perp$
determines the answer through Corollary~\ref{nd:code_monomial:cor:cwe}.  For $r>2$, the
ordinary weight enumerator forgets the symbols and must be replaced by the
complete weight enumerator.  Once $P$ is nontrivial, even the complete weight
enumerator forgets how the coordinate group acts.  The exact required datum
is the equivariant collection of orbit counts in
\eqref{nd:code_monomial:eq:orbit-count}, equivalently a cycle-indexed complete-weight
enumerator.  Thus ordinary weight-enumerator convergence does not by itself
imply $\sigma$-MGF convergence.

The finite audit deliberately exercises the following boundaries:
\begin{enumerate}
\item $P=1$ versus nontrivial $P$, so the complete-weight and equivariant
formulas are tested separately;
\item prime moduli $2,3$ versus the composite modulus $4$;
\item full, even, repetition, and two-generator pair codes;
\item identity, symmetric, and cyclic coordinate groups.
\end{enumerate}
It does not claim to test every asymptotic code family mentioned in broader
surveys.  Equations \eqref{nd:code_monomial:eq:criterion} and \eqref{nd:code_monomial:eq:obstruction} are the
proved reusable tests for such a family once its codes and coordinate groups
are specified.

From the repository root, run
\begin{verbatim}
.venv/bin/python -m lambda_distributions.dists.code_monomial_sigma_mgfs.verification
.venv/bin/pytest -q lambda_distributions/dists/code_monomial_sigma_mgfs/test_verification.py
.venv/bin/marimo edit marimo_notebooks/code_monomial.py
\end{verbatim}
The source and PDF in this directory are standalone artifacts intended for
later inclusion.  No combined PDF is edited or rebuilt.

\clearpage
\part{Braid, Pauli, extraspecial, Clifford, and Barnes--Wall groups}
\section{Scalar cyclic finite braid images}\label{proof:braid-cyclic}
\begingroup
\definecolor{navy}{HTML}{17365D}
\definecolor{teal}{HTML}{147D78}
\providecommand{\Sgm}{\SMG}
\renewcommand{\Sgm}{\SMG}
\setcounter{theorem}{0}
\setcounter{proposition}{0}
\setcounter{lemma}{0}
\setcounter{corollary}{0}
\setcounter{remark}{0}
\begin{abstract}
Using the imported sigma--MGF identities, we compute the scalar cyclic
quotients of $B_3$ in arbitrary representation dimension.  These examples
isolate the scalar-center selection rule.  Explicit Reynolds-projector ranks
agree with both the elementwise power-trace formula and its conjugacy-class
grouping in all tested dimensions $D=1,2,3,4$.
\end{abstract}

\subsection{The representation and the proposed quantity}
Fix $m\geq2$, $D\geq1$, and $\zeta=e^{2\pi i/m}$.  Define
\[
 \rho(\sigma_1)=\rho(\sigma_2)=\zeta I_D.
\]
The matrices are unitary and satisfy
$\rho(\sigma_1)\rho(\sigma_2)\rho(\sigma_1)
=\rho(\sigma_2)\rho(\sigma_1)\rho(\sigma_2)$, so this is a
$D$-dimensional representation of $B_3$.  Its image is $C_m$.

For a tuple $\tau=(\tau_1,\ldots,\tau_\ell)$ let
\[
 W_\tau=\bigotimes_{a=1}^{\ell}\operatorname{Sym}^{\tau_a}(\mathbb C^D),
 \qquad |\tau|=\sum_a\tau_a.
\]
The coefficient under test is
\[
 [m_\tau]\Sgm_{C_m}(\mathbf t),\qquad
 \Sgm_{C_m}(\mathbf t)=\frac1m\sum_{j=0}^{m-1}
 \prod_{a\geq1}\det(I-t_a\zeta^jI_D)^{-1}.
\]

\subsection{Scalar specialization of the imported formula}
Tools~\ref{tool:coefficient} and~\ref{tool:molien} supply the common
finite-group identities.  In the present representation, $\zeta^jI_D$ acts
on all of $W_\tau$ by
$\zeta^{j|\tau|}$.  Hence the Reynolds operator is either zero or identity:
\[
 \boxed{[m_\tau]\Sgm_{C_m}=\mathbf1_{m\mid|\tau|}
 \prod_a\binom{D+\tau_a-1}{\tau_a}.}
\]
This proves the scalar selection rule and also determines the nonzero value,
not merely its support.

\subsection{Independent verification strategy}
The program first closes the matrix group from the displayed generator and
asserts order $m$, unitarity, and the braid relation.  It then constructs an
orthonormal occupation basis of every $\operatorname{Sym}^{r}\mathbb C^D$
inside $(\mathbb C^D)^{\otimes r}$.  In that basis it forms
\[
 P_\tau=\frac1m\sum_{g\in C_m}\bigotimes_a\operatorname{Sym}^{\tau_a}(g)
\]
and computes its rank.  This is the direct target calculation.

The comparator never uses this projector.  It obtains the complete
homogeneous characters from Newton's recurrence
\[
 r h_r(g)=\sum_{k=1}^r\operatorname{Tr}(g^k)h_{r-k}(g),\qquad h_0=1,
\]
then averages $\prod_a h_{\tau_a}(g)$ elementwise and, separately, by
matrix conjugacy classes.

\begin{center}
\begin{tabular}{@{}rrrrrr@{}}\toprule
$m$&$D$&$\tau$&$\dim W_\tau$&$\operatorname{rank}P_\tau$&formula\\\midrule
3&1&$(1)$&1&0&0\\
3&1&$(3)$&1&1&1\\
3&2&$(2,1)$&6&6&6\\
4&2&$(4)$&5&5&5\\
4&3&$(2,2)$&36&36&36\\
5&4&$(5)$&56&56&56\\\bottomrule
\end{tabular}
\end{center}
The maximum projector idempotence/Hermiticity error in this sweep is below
$6\times10^{-15}$.

\subsection{Reproduction and scope}
From the repository root run
\begin{verbatim}
.venv/bin/python for_this_guy/finite_braid_images/
  cyclic_scalar/verification.py
.venv/bin/marimo edit marimo_notebooks/braid_cyclic_scalar.py
\end{verbatim}
after joining each wrapped command.  The proof is valid for every $m,D,\tau$;
the table is a small numerical regression suite, not evidence substituted for
the proof.
\endgroup


\clearpage
\section[The permutation braid image B3 onto S3]{The permutation braid image $B_3\twoheadrightarrow S_3$}\label{proof:braid-s3}
\begingroup
\definecolor{navy}{HTML}{17365D}
\definecolor{teal}{HTML}{147D78}
\providecommand{\Sgm}{\SMG}
\renewcommand{\Sgm}{\SMG}
\setcounter{theorem}{0}
\setcounter{proposition}{0}
\setcounter{lemma}{0}
\setcounter{corollary}{0}
\setcounter{remark}{0}
\begin{abstract}
We construct the adjacent-transposition quotient of $B_3$ in the
three-dimensional permutation representation.  We derive its sigma--MGF
coefficient formula from the three spectral types of $S_3$, and independently
recover the same integers as ranks of explicit Reynolds projectors on tensor
products of symmetric powers.
\end{abstract}

\subsection{Matrix group and braid relation}
Let
\[
 A=\begin{pmatrix}0&1&0\\1&0&0\\0&0&1\end{pmatrix},\qquad
 B=\begin{pmatrix}1&0&0\\0&0&1\\0&1&0\end{pmatrix}.
\]
These unitary matrices obey $ABA=BAB$.  Sending
$\sigma_1\mapsto A$ and $\sigma_2\mapsto B$ therefore defines a
representation of $B_3$.  Matrix closure gives the six permutation matrices,
so the image is $S_3$ acting on $V=\mathbb C^3$.

\subsection{Imported target and the three-class formula}
For $W_\tau=\bigotimes_a\operatorname{Sym}^{\tau_a}V$,
Tools~\ref{tool:coefficient} and~\ref{tool:molien} give
\[
 [m_\tau]\,\smgv{S_3}{V}=\dim(W_\tau^{S_3})
 =\frac16\sum_{g\in S_3}\prod_a h_{\tau_a}(g).
\]
The representation-specific input is that the three conjugacy classes have
eigenvalues
\[
 (1,1,1),\qquad(1,1,-1),\qquad(1,\omega,\omega^2),
 \quad \omega^3=1.
\]
Their generating series are respectively
\[
 \frac1{(1-t)^3},\qquad \frac1{(1-t)(1-t^2)},\qquad\frac1{1-t^3}.
\]
Consequently, with
\[
 A_r=\binom{r+2}{2},\qquad B_r=\left\lfloor\frac r2\right\rfloor+1,
 \qquad C_r=\mathbf1_{3\mid r},
\]
the proposed class formula becomes the completely explicit identity
\[
 \boxed{[m_\tau]\Sgm_{S_3}=\frac16\left(
 \prod_a A_{\tau_a}+3\prod_a B_{\tau_a}+2\prod_a C_{\tau_a}\right).}
\]
This is also the character-and-power-map formula: Newton's recurrence
$r h_r=\sum_{k=1}^r\operatorname{Tr}(g^k)h_{r-k}$ is obtained by taking
the logarithm of the determinant generating series.

\subsection{Direct versus formula verification}
The direct route embeds the normalized occupation basis of
$\operatorname{Sym}^rV$ into $V^{\otimes r}$, restricts $g^{\otimes r}$
to that basis, forms the tensor product over the parts of $\tau$, and ranks
the averaged matrix.  The two formula routes compute the power-trace
recurrence over all six matrices and over the three constructed conjugacy
classes.  Thus the direct route does not diagonalize a Reynolds projector by
using its character in advance.

\begin{center}
\begin{tabular}{@{}rrrrr@{}}\toprule
$\tau$&$\dim W_\tau$&projector rank&six elements&three classes\\\midrule
$(1)$&3&1&1&1\\
$(2)$&6&2&2&2\\
$(1,1)$&9&2&2&2\\
$(3)$&10&3&3&3\\
$(2,1)$&18&4&4&4\\
$(2,2)$&36&8&8&8\\\bottomrule
\end{tabular}
\end{center}
All matrices close to exactly six elements, exactly three conjugacy classes
are recovered, and the largest projector error is below $8\times10^{-16}$.

\subsection{Reproduction}
From the repository root run
\begin{verbatim}
.venv/bin/python for_this_guy/finite_braid_images/
  s3_permutation/verification.py
.venv/bin/marimo edit marimo_notebooks/braid_s3_permutation.py
\end{verbatim}
after joining wrapped lines.  This example verifies the universal finite-image
claim in a non-scalar braid quotient; it is not asserted to be a particular
fusion-space representation of a chosen modular category.
\endgroup


\clearpage
\section{The finite Ising/Jones braid image}\label{proof:braid-ising}
\begingroup
\definecolor{navy}{HTML}{17365D}
\definecolor{teal}{HTML}{147D78}
\providecommand{\Sgm}{\SMG}
\renewcommand{\Sgm}{\SMG}
\setcounter{theorem}{0}
\setcounter{proposition}{0}
\setcounter{lemma}{0}
\setcounter{corollary}{0}
\setcounter{remark}{0}
\begin{abstract}
We test the finite-image formulas on the two-dimensional Ising/Jones braid
representation.  The displayed braid generators close to a 192-element
unitary lift with scalar center of order eight.  We verify its one-sided
sigma--MGF and the full-twist selection rule.  We then pass to conjugation on
$\operatorname{End}(\mathbb C^2)$, obtaining the canonical 24-element
projective image in dimension four, and verify the projective formula there.
\end{abstract}

\subsection{The finite braid image}
Put
\[
 H=\frac1{\sqrt2}\begin{pmatrix}1&1\\1&-1\end{pmatrix},\qquad
 S=\begin{pmatrix}1&0\\0&i\end{pmatrix},\qquad q=e^{-\pi i/8},
\]
and define
\[
 \rho(\sigma_1)=qS,\qquad \rho(\sigma_2)=qHSH.
\]
Direct multiplication gives the braid relation.  Breadth-first matrix closure
gives $|G|=192$; quotienting its eight scalar matrices gives a projective image
of order $24$.  The $B_3$ full twist is
\[
 \rho((\sigma_1\sigma_2)^3)=e^{-\pi i/4}I.
\]
It therefore acts on
$W_\tau=\bigotimes_a\operatorname{Sym}^{\tau_a}\mathbb C^2$ by
$e^{-\pi i|\tau|/4}$, proving the necessary selection rule
\[
 \boxed{W_\tau^G=0\quad\text{unless}\quad8\mid|\tau|.}
\]
The rule is necessary, not sufficient, as the degree-eight computations show.

\subsection{Imported finite-image tools}
Once the 192-element represented image is fixed, its one-sided $\sigma$-MGF
and coefficients are Tools~\ref{tool:molien} and~\ref{tool:coefficient}. The
new issue is projective: scalar changes alter the one-sided series, whereas
adjoint and balanced constructions are lift-independent.

\subsection{Projective correction}
Changing a lift $g$ to $zg$ changes its eigenvalues, so the one-sided series
is not a projective invariant.  Conjugation on $\operatorname{End}(V)$ is
unchanged:
\[
 \operatorname{Ad}_{zg}(A)=(zg)A(zg)^{-1}=gAg^{-1}.
\]
Thus the projectively intrinsic series is the ordinary finite Molien series
of the adjoint representation,
\[
 \Sgm^{\mathrm{proj}}(\mathbf t)=\frac1{|\overline G|}
 \sum_{\bar g\in\overline G}\prod_a
 \det(I-t_a\operatorname{Ad}_{\bar g})^{-1}.
\]
If $g$ is any lift, then
\[
 \operatorname{Tr}(\operatorname{Ad}_{g}^{k})
 =\operatorname{Tr}(g^k)\overline{\operatorname{Tr}(g^k)}
 =|\chi(g^k)|^2,
\]
which proves the proposed projective power-trace formula and its independence
of lift.

\subsection{Independent matrix checks}
For each $\tau$, the program explicitly restricts tensor powers to normalized
symmetric occupation bases and ranks the averaged Reynolds matrix.  Separately,
it computes $h_r(g)$ from traces of powers using Newton's recurrence and
averages $\prod_a h_{\tau_a}(g)$ first over elements and then over matrix
conjugacy classes.  The lift has 32 classes; the adjoint image has five.

\begin{center}
\begin{tabular}{@{}llrrrr@{}}\toprule
representation&$\tau$&$\dim W_\tau$&rank&elements&classes\\\midrule
lift on $\mathbb C^2$&$(1)$&2&0&0&0\\
&$(4)$&5&0&0&0\\
&$(8)$&9&0&0&0\\
&$(4,4)$&25&1&1&1\\\addlinespace
adjoint on $\operatorname{End}(\mathbb C^2)$&$(1)$&4&1&1&1\\
&$(2)$&10&2&2&2\\
&$(1,1)$&16&2&2&2\\
&$(3)$&20&2&2&2\\\bottomrule
\end{tabular}
\end{center}
All braid, unitarity, closure, integrality, and equality assertions pass.  The
largest projector error is below $4\times10^{-15}$.

\subsection{Reproduction and interpretation}
Run
\begin{verbatim}
.venv/bin/python for_this_guy/finite_braid_images/
  ising_clifford/verification.py
.venv/bin/marimo edit marimo_notebooks/braid_ising_clifford.py
\end{verbatim}
after joining wrapped lines.  The computation verifies the stated formulas
for the displayed $B_3$ Ising representation.  It does not claim that every
root-of-unity Jones representation has finite image; finiteness remains an
essential hypothesis.
\endgroup


\clearpage
\section{Specified braid images and conditional quantum-image closures}\label{proof:new-dist:braid_tqft}

\resultbox{\textbf{Canonical testing artifacts.}\par\smallskip Exact backend: \path{lambda_distributions/dists/braid_tqft_quantum_images_sigma_mgfs/verification.py}.\par Regression: \path{lambda_distributions/dists/braid_tqft_quantum_images_sigma_mgfs/test_verification.py}.\par Marimo: \path{marimo_notebooks/braid_tqft.py}.}

\subsection*{Scope and outcome}

A finite braid or TQFT image has a uniform measure; an infinite image does
not, and its canonical unitary substitute is normalized Haar measure on the
compact closure.  We import the generalized Molien identity in both settings,
record its class-power form, and test each mathematically distinct mechanism
on a specified representation.  The program constructs a $192$-matrix Ising braid lift, the
$27$-matrix qutrit Heisenberg group, a $96$-matrix rank-one Weil lift, the
$24$-element projective qubit Clifford group, and regular and affine
phase-space permutation groups through degree $25$.  All 35 finite formula
comparisons and all 1,440 affine fixed-power checks pass.  Sixteen seeded Haar
diagnostics corroborate the exact $U$, $SU$, $O$, and $USp$ conclusions.  The
generic Jones, mapping-class, and Hadamard statements are recorded with their
necessary image/closure hypotheses; no unspecified representation is treated
as though it determined a matrix group.

\proofstatus{Mixed}{The finite-image and fixed compact-closure formulas are
exact. Generic TQFT image determination and closure classification are not
proved here; those conditional research directions are moved to the final
provisional part.}

\subsection{The probability distinction and the common target}

Let $V$ be a $d$-dimensional complex unitary module.  For a finite represented
image $I\leq U(V)$, define
\begin{equation}\label{nd:braid_tqft:eq:finite-def}
\smgv{I}{V}(\mathbf t)
=\frac1{|I|}\sum_{x\in I}\prod_{a\geq1}
\det(I-t_a x)^{-1}.
\end{equation}
If a unitary representation has infinite image, there is no uniform
probability measure on that countable image.  Put instead
\[
K=\overline{\rho(\Gamma)}\leq U(V)
\]
and use normalized Haar measure:
\begin{equation}\label{nd:braid_tqft:eq:compact-def}
\mathbb E_{k\in K}\!\left[\ExpS(X_{K,V}h_1)\right](\mathbf t)
=\int_K\prod_{a\geq1}\det(I-t_a k)^{-1}\,dk.
\end{equation}
This agrees with \eqref{nd:braid_tqft:eq:finite-def} for finite $K$ and is unchanged when a
dense subgroup is replaced by its closure.  Thus ``braid representation'' or
``TQFT representation'' is not enough data: one must specify the represented
image, a finite central lift when the action is projective, or the compact
closure.

For a partition or composition $\tau=(\tau_1,\ldots,\tau_s)$, set
\[
W_\tau(V)=\bigotimes_{j=1}^s\Sym^{\tau_j}V,
\qquad |\tau|=\sum_j\tau_j.
\]
The coefficient tested throughout is
\[
[m_\tau]\MGF_{G,V}=\dim W_\tau(V)^G.
\]

\subsection{Universal finite-image and compact-closure theorem}

\begin{theorem}[Generalized Molien/class-power identity]\label{nd:braid_tqft:thm:universal}
For a finite unitary image $I$ with character $\chi_V$,
\begin{align}
\smgv{I}{V}
&=\frac1{|I|}\sum_{x\in I}
\exp\!\left(\sum_{r\geq1}\frac{p_r}{r}\chi_V(x^r)\right)\label{nd:braid_tqft:eq:element-power}\\
&=\sum_{[x]\subseteq I}\frac1{|C_I(x)|}
\exp\!\left(\sum_{r\geq1}\frac{p_r}{r}\chi_V(x^r)\right),\label{nd:braid_tqft:eq:class-power}
\end{align}
where $p_r=\sum_a t_a^r$.  Moreover
\begin{equation}\label{nd:braid_tqft:eq:coeff}
\boxed{[m_\tau]\,\smgv{I}{V}=\dim W_\tau(V)^I.}
\end{equation}
The same statements hold for compact $K$ after replacing the finite average
by normalized Haar integration.
\end{theorem}

\begin{proof}[Source]
The coefficient and determinant statements are Tools~\ref{tool:coefficient}
and~\ref{tool:molien}.  The power-character expression is the standard
logarithmic determinant change of coordinates, and the class form merely
groups the imported finite average by conjugacy class.  The compact statement
is the same imported identity with normalized Haar measure.
\end{proof}

The theorem shows that the exact finite answer is determined by conjugacy
classes, power maps, and the represented character.  It does \emph{not} say
that the abstract source group or a TQFT label determines those data.

\subsubsection{Projective representations and scalar selection}

Suppose a projective image becomes honest on a finite central extension
\[
1\longrightarrow Z\longrightarrow\widetilde I\longrightarrow I
\longrightarrow1.
\]
Theorem~\ref{nd:braid_tqft:thm:universal} applies to $\widetilde I$.  If $z\in Z$ acts as
$\xi(z)I$, then $z$ acts on $W_\tau(V)$ as $\xi(z)^{|\tau|}$.  Therefore
\begin{equation}\label{nd:braid_tqft:eq:central-rule}
[m_\tau]\MGF_{\widetilde I,V}=0
\quad\text{unless}\quad
\xi(z)^{|\tau|}=1\ \text{for every }z\in Z.
\end{equation}
This is a necessary condition, not a promise of a nonzero invariant in every
allowed degree.  It also explains why the one-sided series depends on the
chosen finite lift.

\subsection{Finite Ising braid image}

Let
\[
H=\frac1{\sqrt2}\begin{pmatrix}1&1\\1&-1\end{pmatrix},\qquad
R=e^{-\pi i/8}\begin{pmatrix}1&0\\0&i\end{pmatrix},
\]
and put
\begin{equation}\label{nd:braid_tqft:eq:ising-generators}
\rho(\sigma_1)=R,\qquad \rho(\sigma_2)=HRH.
\end{equation}
Direct multiplication gives
\[
\rho(\sigma_1)\rho(\sigma_2)\rho(\sigma_1)
=\rho(\sigma_2)\rho(\sigma_1)\rho(\sigma_2),
\]
so \eqref{nd:braid_tqft:eq:ising-generators} is an honest two-dimensional representation of
$B_3$.  The verifier starts from these matrices and closes under
multiplication.  It obtains $192$ matrices, $24$ matrices modulo scalar phase,
and exactly eight scalar matrices $\mu_8I$.  Thus \eqref{nd:braid_tqft:eq:central-rule}
forces $8\mid |\tau|$.

For each matrix $A$, the direct route diagonalizes $A$ and expands
$\det(I-tA)^{-1}$.  The independent power route computes $\Tr(A^r)$ by matrix
multiplication and uses Newton's recurrence
\begin{equation}\label{nd:braid_tqft:eq:newton}
h_0(A)=1,\qquad h_n(A)=\frac1n\sum_{r=1}^n
\Tr(A^r)h_{n-r}(A).
\end{equation}
The literal matrix averages agree:
\begin{center}
\begin{tabular}{@{}rrrrrrrr@{}}
\toprule
$\tau$ & $(1)$ & $(2)$ & $(4)$ & $(8)$ & $(7,1)$ & $(4,4)$ & $(1^8)$\\
\midrule
direct matrices &0&0&0&0&1&1&7\\
power formula   &0&0&0&0&1&1&7\\
\bottomrule
\end{tabular}
\end{center}
At $(t_1,t_2,t_3)=(0.023,-0.031,0.041)$, direct determinants and the
truncated power-character exponential agree to $2.2\times10^{-18}$; the
braid-relation residual is $3.0\times10^{-16}$.  The power series is truncated
at $r=60$ only for this nonzero numerical evaluation.  All displayed
coefficients use the exact finite recurrence \eqref{nd:braid_tqft:eq:newton} through the
required degree.

\begin{remark}
This example validates a particular finite Ising/Clifford braid image.  It
does not imply that generic Jones, Hecke, Temperley--Lieb, or BMW sectors have
finite image.  Density results for broad Jones--Wenzl families show why the
finite/closure distinction is substantive \cite{nd:braid_tqft:flw}.
\end{remark}

\subsection{Monomial finite image: the qutrit Heisenberg group}

Suppose a represented finite group is monomial in a chosen basis.  Write
$x=D(\alpha)P_\pi$.  For a cycle $C=(j_1\cdots j_\ell)$ of $\pi$, define
$\alpha_C=\prod_{j\in C}\alpha_j$.  Reordering the basis cycle by cycle gives
\begin{equation}\label{nd:braid_tqft:eq:monomial-det}
\det(I-tx)=\prod_{C\in\operatorname{Cyc}(\pi)}
(1-\alpha_Ct^{|C|}),
\end{equation}
and hence
\begin{equation}\label{nd:braid_tqft:eq:monomial-mgf}
\boxed{\MGF_{I,V}=\frac1{|I|}\sum_{x\in I}
\prod_{a,C}(1-\alpha_C(x)t_a^{|C|})^{-1}.}
\end{equation}

The concrete test is the qutrit Heisenberg group
\[
H_3=\{\omega^cX^aZ^b:a,b,c\in\F_3\},\qquad
Xe_j=e_{j+1},\quad Ze_j=\omega^je_j,
\]
where $\omega=e^{2\pi i/3}$.  All $27$ dense $3\times3$ matrices are built.
The code separately reads the nonzero entry in each column, traverses the
permutation cycles, multiplies their phases, and evaluates
\eqref{nd:braid_tqft:eq:monomial-mgf}.  The dense and cycle-phase multivariate averages
agree to $1.8\times10^{-20}$.

Here the central phase rule and the two spectral sectors simplify to
\begin{equation}\label{nd:braid_tqft:eq:h3formula}
[m_\tau]\MGF_{H_3,V}=
\begin{cases}
\displaystyle \frac19\prod_j\binom{\tau_j+2}{\tau_j}
+\frac89\prod_j\mathbf1_{3\mid\tau_j},&3\mid|\tau|,\\[6pt]
0,&3\nmid|\tau|.
\end{cases}
\end{equation}
Indeed, the three central matrices contribute the first term after the phase
filter, while every noncentral matrix has spectrum a scalar multiple of all
three cube roots and contributes the second term.

\begin{center}
\begin{tabular}{@{}lrrrrrrr@{}}
\toprule
$\tau$&$(1)$&$(2)$&$(3)$&$(2,1)$&$(4)$&$(3,3)$&$(6)$\\
\midrule
direct matrices&0&0&2&2&0&12&4\\
formula \eqref{nd:braid_tqft:eq:h3formula}&0&0&2&2&0&12&4\\
\bottomrule
\end{tabular}
\end{center}
This is a representative finite Gaussian/Heisenberg core of the type that
occurs in metaplectic and parafermionic constructions.  Identifying a
particular physical braid sector with this group requires its own generators;
finiteness alone does not imply a universal asymptotic law.

\subsection{Rank-one Weil lift and the phase issue}

Over $\F_3$, define the Fourier and chirp matrices
\begin{equation}\label{nd:braid_tqft:eq:weil-generators}
F_{xy}=3^{-1/2}\omega^{xy},\qquad
P_{xx}=\omega^{x^2/2},
\end{equation}
where $1/2=2$ in $\F_3$.  Simultaneously track the generators
\[
S=\begin{pmatrix}0&-1\\1&0\end{pmatrix},\qquad
T=\begin{pmatrix}1&0\\1&1\end{pmatrix}
\quad\text{in }SL_2(\F_3).
\]
Projective matrix closure gives all $24$ elements of $SL_2(\F_3)$; retaining
the generated scalar phases gives a $96$-matrix lift.

The finite-field Weil character has the form
\begin{equation}\label{nd:braid_tqft:eq:weil-character}
\chi_\omega(\widetilde A)=\gamma(\widetilde A)
|\ker(A-I)|^{1/2},\qquad |\gamma(\widetilde A)|=1,
\end{equation}
with a lift-dependent Weil-index phase \cite{nd:braid_tqft:thomas}.  Consequently
\begin{equation}\label{nd:braid_tqft:eq:weil-magnitude}
|\chi_\omega(\widetilde A)|^2=|\ker(A-I)|.
\end{equation}
The verifier checks \eqref{nd:braid_tqft:eq:weil-magnitude} for every $A\in SL_2(\F_3)$.
The distribution is
\begin{center}
\begin{tabular}{@{}crrr@{}}
\toprule
$|\ker(A-I)|$&1&3&9\\
number of matrices&15&8&1\\
observed $|\Tr\omega(A)|^2$&1&3&9\\
\bottomrule
\end{tabular}
\end{center}

For the full lift, the phase-sensitive class-power formula is
\begin{equation}\label{nd:braid_tqft:eq:weil-class}
\MGF^{\mathrm{Weil}}=
\sum_{[\widetilde A]}\frac1{|C(\widetilde A)|}
\exp\!\left(\sum_{r\geq1}\frac{p_r}{r}
\chi_\omega(\widetilde A^r)\right).
\end{equation}
The direct and power-character coefficient averages give
\begin{center}
\begin{tabular}{@{}lrrrrr@{}}
\toprule
$\tau$&$(1)$&$(2)$&$(4)$&$(3,1)$&$(4,4)$\\
\midrule
direct lift&0&0&1&2&6\\
formula \eqref{nd:braid_tqft:eq:weil-class}&0&0&1&2&6\\
\bottomrule
\end{tabular}
\end{center}
The full numerical MGF error is $2.3\times10^{-16}$.

Equation \eqref{nd:braid_tqft:eq:weil-magnitude} controls balanced trace moments, but it does
not determine the holomorphic series \eqref{nd:braid_tqft:eq:weil-class}: the latter also
needs the joint phases of $A,A^2,\ldots$.  This is exactly why fixed-space
magnitude stabilization cannot by itself prove convergence of the full Weil
sigma--MGF.

The same rank-one construction may be viewed as a finite quotient of a torus
mapping-class representation.  Higher-genus or composite-level formulas
remain exact class sums only after the relevant lift and Weil character have
been specified; Chinese-remainder structure alone does not remove the local
phase data.

\subsection{Balanced series and unitary designs}

Scalar phases often annihilate the one-sided series.  The appropriate
design-sensitive object is
\begin{equation}\label{nd:braid_tqft:eq:balanced}
\BMGF_{G,V}(\mathbf s,\mathbf t)=\frac1{|G|}\sum_{g\in G}
\prod_i\det(I-s_ig)^{-1}\prod_j\det(I-t_j\overline g)^{-1}.
\end{equation}
Character expansion and Reynolds averaging prove
\begin{equation}\label{nd:braid_tqft:eq:balanced-coeff}
[m_\alpha(\mathbf s)m_\beta(\mathbf t)]\BMGF_{G,V}
=\dim\left[
\left(\bigotimes_i\Sym^{\alpha_i}V\right)\otimes
\left(\bigotimes_j\Sym^{\beta_j}V^*\right)
\right]^G.
\end{equation}
Only $|\alpha|=|\beta|$ is intrinsically projective.  For
$\alpha=\beta=(1^k)$, the coefficient becomes the frame potential
\begin{equation}\label{nd:braid_tqft:eq:frame}
\Phi_k(G)=\frac1{|G|}\sum_{g\in G}|\Tr g|^{2k}
=\dim\operatorname{End}_G(V^{\otimes k}).
\end{equation}

If $G$ is a unitary $t$-design, every balanced coefficient with
$|\alpha|,|\beta|\leq t$ equals its Haar $U(d)$ value.  Conversely, equality
of \eqref{nd:braid_tqft:eq:frame} with Haar at degree $k$ forces equality of the nested
invariant spaces and hence equality of the $k$th moment projector.  By
Schur--Weyl duality,
\[
\Phi_k(U(d))=\sum_{\substack{\lambda\vdash k\\\ell(\lambda)\leq d}}
(f^\lambda)^2.
\]

Closing $H$ and $S=\operatorname{diag}(1,i)$ modulo scalar phase gives the
$24$-element projective one-qubit Clifford group.  Direct matrix averaging
and the hook-length formula give
\begin{center}
\begin{tabular}{@{}crrrr@{}}
\toprule
$k$&1&2&3&4\\
\midrule
$\Phi_k(\mathrm{Cliff}_1)$&1&2&5&15\\
$\Phi_k(U(2))$&1&2&5&14\\
\bottomrule
\end{tabular}
\end{center}
Thus the tested group agrees through degree three and first differs at degree
four, consistent with the general multiqubit design theorem and its extra
fourth-order stabilizer invariant \cite{nd:braid_tqft:clifford4}.

\subsection{Weyl--Heisenberg regular and affine phase-space actions}

\subsubsection{Regular action and the prime SIC formula}

Let $A_d\cong(\mathbb Z/d\mathbb Z)^2$ act regularly on itself, viewed as the
$d^2$ projective outcomes in a Weyl--Heisenberg covariant SIC setting.  An
element $a$ of order $o(a)$ has $d^2/o(a)$ cycles, so
\begin{equation}\label{nd:braid_tqft:eq:regular-general}
\MGF_{A_d,\mathrm{Reg}}=
\frac1{d^2}\sum_{a\in A_d}\prod_i
(1-t_i^{o(a)})^{-d^2/o(a)}.
\end{equation}
For prime $p$, every nonidentity element has order $p$:
\begin{equation}\label{nd:braid_tqft:eq:regular-prime}
\boxed{\MGF_{A_p,\mathrm{Reg}}=
\frac1{p^2}\left[
\prod_i(1-t_i)^{-p^2}+(p^2-1)\prod_i(1-t_i^p)^{-p}
\right].}
\end{equation}
In particular
\begin{equation}\label{nd:braid_tqft:eq:regular-pair}
[m_{(1,1)}]\MGF_{A_p,\mathrm{Reg}}=p^2,
\end{equation}
because only the identity fixes a point and the pair coefficient is the
average squared fixed-point count.

The verifier constructs all permutation matrices for $p=2,3,5$.  Selected
exact comparisons are
\begin{center}
\begin{tabular}{@{}crrrr@{}}
\toprule
$p$&permutation degree&$[m_{11}]$&$[m_2]$&$[m_{p,p}]$\\
\midrule
2&4&4&4&28\\
3&9&9&5&3,033\\
5&25&25&13&564,110,025\\
\bottomrule
\end{tabular}
\end{center}
Every entry agrees with coefficient extraction from \eqref{nd:braid_tqft:eq:regular-prime};
full nonzero determinant evaluations agree to $2.3\times10^{-16}$.  The growth
in \eqref{nd:braid_tqft:eq:regular-pair} proves that the raw regular-action series has no
coefficientwise $p\to\infty$ limit.  The contextual link between prime
Weyl--Heisenberg SICs and Clifford normalizers is discussed in \cite{nd:braid_tqft:zhu-sic}.

\subsubsection{Affine symmetry groups}

Let $\Gamma=A\rtimes K$ act on the additive space $A$.  Write
$g=(b,A_0)$ with action $x\mapsto A_0x+b$ and put
\[
b_r=(I+A_0+\cdots+A_0^{r-1})b.
\]
Then
\[
g^r(x)=A_0^rx+b_r,
\]
so the fixed-point equation is $(I-A_0^r)x=b_r$.  Therefore
\begin{equation}\label{nd:braid_tqft:eq:affine-fixed}
\boxed{
|\Fix(g^r)|=
\begin{cases}
|\ker(I-A_0^r)|,&b_r\in\operatorname{im}(I-A_0^r),\\
0,&\text{otherwise}.
\end{cases}}
\end{equation}
Fixed points of powers recover cycle counts by M\"obius inversion, hence
\eqref{nd:braid_tqft:eq:affine-fixed} determines the full permutation sigma--MGF.

Ordered-pair orbits are classified by the difference of the two points, so
\begin{equation}\label{nd:braid_tqft:eq:affine-pair}
[m_{(1,1)}]\MGF_{A\rtimes K,A}
=\#(K\backslash A)=1+\#(K\backslash(A\setminus\{0\})).
\end{equation}
For $K=SL_2(\F_p)$, the nonzero vectors form one orbit and the answer is $2$.
The code constructs the full affine permutation groups for $p=2,3$:
\begin{center}
\begin{tabular}{@{}crrrr@{}}
\toprule
$p$&degree&group order&fixed-power checks&pair coefficient\\
\midrule
2&4&24&144&2\\
3&9&216&1,296&2\\
\bottomrule
\end{tabular}
\end{center}
All $1,440$ comparisons of direct fixed points with
\eqref{nd:braid_tqft:eq:affine-fixed}, for powers $1\leq r\leq6$, pass exactly.

\subsection{Compact closures and stable classical laws}

If the represented image is infinite, Theorem~\ref{nd:braid_tqft:thm:universal} applies to
its compact closure $K$.  Four cases must be distinguished.

\subsubsection{Unitary and special-unitary closure}

If $K=U(d)$ in the defining representation, the scalar circle acts on every
positive-degree $W_\tau(V)$ with weight $|\tau|$.  Hence
\begin{equation}\label{nd:braid_tqft:eq:u-trivial}
\boxed{\MGF_{U(d),\C^d}=1.}
\end{equation}
For $SU(d)$, its center $\mu_dI$ implies that a coefficient vanishes unless
$d\mid|\tau|$.  Consequently, for every fixed positive $\tau$, the
coefficient is exactly zero once $d>|\tau|$, and
\begin{equation}\label{nd:braid_tqft:eq:su-limit}
\MGF_{SU(d),\C^d}\longrightarrow1
\quad\text{coefficientwise.}
\end{equation}
The divisibility condition is only necessary.  For example,
$[m_{(1^d)}]=1$ because the determinant tensor is invariant, whereas
$[m_{(d)}]=0$ for $d>1$ in the defining module.

\subsubsection{Orthogonal and compact symplectic closure}

The first fundamental theorems for $O(d)$ and $USp(2n)$ say that stable tensor
invariants are generated by pair contractions with, respectively, a
symmetric or alternating bilinear form.  Repackaging these contractions by
symmetric powers yields the stable identities
\begin{equation}\label{nd:braid_tqft:eq:stable-classical}
\boxed{
\MGF_{O(\infty)}=\ExpS(h_2),\qquad
\MGF_{USp(\infty)}=\ExpS(e_2).}
\end{equation}
This is coefficientwise stability, not equality in every degree at every
finite rank; contraction diagrams can acquire relations outside the stable
range.

The proof of \eqref{nd:braid_tqft:eq:u-trivial}--\eqref{nd:braid_tqft:eq:stable-classical} is algebraic.
As an independent diagnostic, the notebook draws seeded Haar matrices by
complex Gaussian QR for $U$, determinant-normalized QR for $SU$, real QR for
$O$, and quaternionic Gram--Schmidt for $USp$.  It averages the same
symmetric-power characters used in the finite tests.  Representative rows
from $3,000$ samples each are:
\begin{center}
\begin{tabular}{@{}llrrr@{}}
\toprule
group&$\tau$&estimate&target&standard error\\
\midrule
$U(4)$&$(1,1)$&$-0.0279-0.0160i$&0&0.0258\\
$SU(2)$&$(1,1)$&1.0064&1&0.0182\\
$SU(3)$&$(1,1,1)$&$0.9930-0.0013i$&1&0.0412\\
$O(6)$&$(2)$&0.9956&1&0.0182\\
$O(6)$&$(1,1,1,1)$&2.7695&3&0.1503\\
$O(6)$&$(2,2)$&2.1018&2&0.0731\\
$USp(4)$&$(2)$&$-0.0209$&0&0.0178\\
$USp(4)$&$(1,1,1,1)$&3.1935&3&0.1554\\
$USp(4)$&$(2,2)$&1.0291&1&0.0437\\
\bottomrule
\end{tabular}
\end{center}
All sixteen diagnostics lie within the stated five-standard-error/absolute
tolerance.  Worst unitarity residuals are below $6.1\times10^{-15}$.  These
samples test the matrix construction and low-degree consequences; they are
not substituted for the Reynolds and invariant-theory proofs.

\subsection{How the remaining family labels fit}

The universal theorem organizes the broader claims without assigning a
formula to an unspecified representation.

\begin{itemize}
\item \textbf{Jones, Hecke, Temperley--Lieb, and BMW sectors.}  For a specified
finite image, use \eqref{nd:braid_tqft:eq:class-power}.  For an infinite image, use
\eqref{nd:braid_tqft:eq:compact-def}.  If the closure is eventually $U(d_n)$ or $SU(d_n)$,
the one-sided coefficientwise limit is $1$; orthogonal and symplectic closures
give \eqref{nd:braid_tqft:eq:stable-classical}.  Fusion-space dimension growth alone does
not determine the closure.
\item \textbf{Property (F) and finite mapping-class images.}  Property (F)
supplies finiteness, not a universal class table.  One still needs the lift,
power maps, and represented character.  Twisted quantum-double and
metaplectic examples often allow monomial or Weil compression; the
Heisenberg and rank-one Weil sections demonstrate the two reductions.
\item \textbf{Dense mapping-class images.}  Projective density must first be
translated to the honest compact closure relevant to the chosen lift.  Once
that closure and its central scalars are identified, the compact theorem
applies.
\item \textbf{Hessian and Heisenberg-normalizer groups.}  The defining
unitary action uses a Weil--Heisenberg class sum; the phase-space action uses
\eqref{nd:braid_tqft:eq:affine-fixed}.  Their asymptotics depend on fixed-space phases or on
the orbit structure of the complement.
\item \textbf{Hadamard-generated groups.}  Hadamard generation alone does not
imply finiteness.  In the finite qubit example, $H$ and $HS$ are both complex
Hadamard and generate the same lift as $H,S$, since $H(HS)=S$.  Infinite
examples require the closure-Haar formula.
\end{itemize}

The general obstruction to a raw coefficientwise limit is already visible at
balanced fourth tensor degree:
\[
\dim\operatorname{End}_{I_n}(V_n^{\otimes2}).
\]
If this diverges, the corresponding coefficient diverges.  Conversely,
bounded low moments do not determine the full holomorphic series when
power-character phases fail to stabilize.

\subsection{Verification design and reproducibility}

The program uses structurally different routes:
\begin{enumerate}
\item \textbf{Matrix closure.}  Products of the displayed generators are
stored until no new matrix appears.  Projective closures normalize a single
nonzero entry only when the coefficient being tested is phase neutral.
\item \textbf{Direct determinants.}  Every finite represented matrix is
inserted into \eqref{nd:braid_tqft:eq:finite-def}.
\item \textbf{Power characters.}  Matrix powers and \eqref{nd:braid_tqft:eq:newton} compute
coefficients without diagonalizing; the numerical whole-function check uses
\eqref{nd:braid_tqft:eq:element-power} through power $60$.
\item \textbf{Monomial cycles.}  Nonzero column entries determine permutation
cycles and accumulated phases, independently checking
\eqref{nd:braid_tqft:eq:monomial-det}.
\item \textbf{Finite-field tracking.}  Weil words are followed simultaneously
in $SL_2(\F_3)$; affine fixed points are obtained from direct permutations and
from \eqref{nd:braid_tqft:eq:affine-fixed}.
\item \textbf{Orbit/design comparators.}  Pair orbits and the hook-length Haar
frame potential avoid the determinant formula.
\item \textbf{Haar diagnostics.}  Seeded continuous samples include standard
errors and group residuals and are explicitly reported as numerical evidence.
\end{enumerate}

Run the regression test from the repository root:
\begin{verbatim}
.venv/bin/python -m pytest -q \
  lambda_distributions/dists/braid_tqft_quantum_images_sigma_mgfs/test_verification.py
\end{verbatim}
Open the interactive laboratory with:
\begin{verbatim}
.venv/bin/marimo run \
  marimo_notebooks/braid_tqft.py
\end{verbatim}
The suite reports \textcolor{teal}{\textbf{PASS}}: 35 finite comparisons, 1,440 affine fixed-power
checks, and 16 compact-closure diagnostics.

\begin{thebibliography}{9}
\bibitem{nd:braid_tqft:howe}
S.~Howe, \emph{Random matrix statistics and zeroes of $L$-functions via
probability in $\lambda$-rings},
\url{https://arxiv.org/abs/2412.19295}.

\bibitem{nd:braid_tqft:flw}
M.~H.~Freedman, M.~J.~Larsen, and Z.~Wang,
\emph{The two-eigenvalue problem and density of Jones representation of braid
groups}, \url{https://arxiv.org/abs/math/0103200}.

\bibitem{nd:braid_tqft:thomas}
T.~Thomas, \emph{The Character of the Weil Representation},
\url{https://arxiv.org/abs/math/0610644}.

\bibitem{nd:braid_tqft:clifford4}
H.~Zhu, R.~Kueng, M.~Grassl, and D.~Gross,
\emph{The Clifford group fails gracefully to be a unitary 4-design},
\url{https://arxiv.org/abs/1609.08172}.

\bibitem{nd:braid_tqft:zhu-sic}
H.~Zhu, \emph{SIC--POVMs and Clifford groups in prime dimensions},
\url{https://arxiv.org/abs/1003.3591}.
\end{thebibliography}

\clearpage
\section{Central extensions and projective representations}\label{proof:new-dist:central_extensions}

\resultbox{\textbf{Canonical testing artifacts.}\par\smallskip Exact backend: \path{lambda_distributions/dists/central_extensions_projective_representations/verification.py}.\par Regression: \path{lambda_distributions/dists/central_extensions_projective_representations/test_verification.py}.\par Marimo: \path{marimo_notebooks/central_extensions.py}.}

\subsection*{Scope and outcome}

For a finite or compact central extension acting with one scalar central
character, we prove that the one-sided $\sigma$-MGF is obtained
by a homogeneous root-of-unity projection.  A continuous scalar center kills
every positive degree.  Since a projective matrix has no canonical
one-sided eigenvalue distribution, we construct and prove lift independence
of the charge-zero balanced series.  We then specialize the theory to Schur
and pin covers, finite Weil representations, extraspecial groups, Clifford
groups, sporadic covers, and determinant-one kernels.  The companion marimo
notebook constructs nine representative matrix groups in dimensions $2$,
$3$, $5$, and $9$.  It performs 180 coefficient comparisons through total
degree six, five full determinant-average checks, and nine balanced-moment
checks.  Every comparison passes; the maximum numerical discrepancy is
$6.28\times10^{-12}$.

\subsection{The generalized Molien target}

Let $K$ be a finite or compact group and let
$\rho:K\to U(V)$ be a finite-dimensional unitary representation.  For a
partition $\tau=(\tau_1,\ldots,\tau_\ell)$ set
\[
 W_\tau(V)=\bigotimes_{i=1}^{\ell}\Sym^{\tau_i}V,
 \qquad |\tau|=\sum_i\tau_i.
\]
The one-sided $\sigma$-MGF is
\begin{equation}\label{nd:central_extensions:eq:molien}
 \MGF_{K,V}(\mathbf t)
 =\int_K\prod_{i\geq1}\det(I-t_i\rho(g))^{-1}\,dg.
\end{equation}

\begin{theorem}[Imported generalized Molien identity]\label{nd:central_extensions:thm:molien}
With normalized counting measure in the finite case and normalized Haar
measure in the compact case,
\[
 \MGF_{K,V}(\mathbf t)
 =\sum_\tau\dim W_\tau(V)^K\,m_\tau(\mathbf t).
\]
Equivalently,
\begin{equation}\label{nd:central_extensions:eq:trace-exp}
 \prod_i\det(I-t_i\rho(g))^{-1}
 =\exp\!\left(\sum_{k\geq1}\frac{p_k(\mathbf t)}{k}
                   \Tr(\rho(g)^k)\right).
\end{equation}
\end{theorem}

\begin{proof}[Source]
This is Tools~\ref{tool:coefficient} and~\ref{tool:molien}; it is retained
with its equation labels because later central-character formulas refer to
them.
\end{proof}

\subsection{Central-character projection}

Consider a central extension
\begin{equation}\label{nd:central_extensions:eq:extension}
 1\longrightarrow Z\longrightarrow\widetilde G
 \xrightarrow{\pi}G\longrightarrow1
\end{equation}
and a representation $\rho:\widetilde G\to U(V)$ for which
$\rho(z)=\chi(z)I_V$.  Write $C=\chi(Z)\leq U(1)$.

\begin{theorem}[Finite central-character projector]\label{nd:central_extensions:thm:projector}
Suppose $C=\mu_r$.  Then
\begin{equation}\label{nd:central_extensions:eq:selection}
 [m_\tau]\MGF_{\widetilde G,V}=0\qquad\text{unless }r\mid|\tau|.
\end{equation}
For any set-theoretic section $g\mapsto\widetilde g$,
\begin{equation}\label{nd:central_extensions:eq:fibre-average}
 \boxed{\MGF_{\widetilde G,V}(\mathbf t)
 =\frac1{|G|}\sum_{g\in G}\frac1r\sum_{\zeta\in\mu_r}
 \prod_i\det(I-\zeta t_i\rho(\widetilde g))^{-1}.}
\end{equation}
Thus adjoining scalar $r$th roots applies
\begin{equation}\label{nd:central_extensions:eq:Pr}
 \PP_rF(\mathbf t)=\frac1r\sum_{\zeta\in\mu_r}F(\zeta\mathbf t),
\end{equation}
which retains precisely the homogeneous degrees divisible by $r$.
\end{theorem}

\begin{proof}
The center acts on $W_\tau(V)$ by the scalar $\chi(z)^{|\tau|}$.
A fixed vector can exist only when this scalar is one for all $z$, proving
\eqref{nd:central_extensions:eq:selection}.  For the integral proof, disintegrate normalized
measure over the fibres of $\pi$.  The pushforward of normalized measure on
$Z$ under $\chi$ is uniform measure on $C$.  In total degree $d$, the fibre
average is
\[
 \frac1r\sum_{\zeta\in\mu_r}\zeta^d
 =\begin{cases}1,&r\mid d,\\0,&r\nmid d.\end{cases}
\]
This proves \eqref{nd:central_extensions:eq:fibre-average} and \eqref{nd:central_extensions:eq:Pr}, independently of
the chosen section and even when $\chi$ has a kernel.
\end{proof}

\begin{corollary}[Continuous scalar saturation]\label{nd:central_extensions:cor:u1}
If $C=U(1)$, then
\[
 \boxed{\MGF_{\widetilde G,V}=1.}
\]
\end{corollary}

\begin{proof}
The scalar average in total degree $d$ is
$\int_{U(1)}z^d\,dz$, equal to zero for every $d>0$ and one for $d=0$.
\end{proof}

This explains the trivial one-sided series of $U(d)$ in its defining
module and of a full unitary normalizer containing every scalar matrix.

\subsection{The canonical projective series}

A projective matrix $[U_g]\in PU(V)$ has no canonical eigenvalue multiset:
another lift $\zeta_gU_g$ multiplies every eigenvalue by $\zeta_g$.
Consequently a projective representation does not canonically define
\eqref{nd:central_extensions:eq:molien}.

For partitions $\alpha,\beta$, define the charge-zero balanced series
\begin{equation}\label{nd:central_extensions:eq:balanced-def}
 \boxed{\BSMG^{\mathrm{proj}}_{G,[V]}(\mathbf s,\mathbf t)
 =\sum_{|\alpha|=|\beta|}
 \dim\!\left(W_\alpha(V)\otimes W_\beta(V)^*\right)^{\widetilde G}
 m_\alpha(\mathbf s)m_\beta(\mathbf t).}
\end{equation}

\begin{theorem}[Lift-independent projective Molien formula]
For any choice of unitary lift $U_g$ of each projective matrix,
\begin{equation}\label{nd:central_extensions:eq:balanced-ct}
 \boxed{\BSMG^{\mathrm{proj}}_{G,[V]}
 =\CT_u\int_G
 \prod_i\det(I-u s_iU_g)^{-1}
 \prod_j\det(I-u^{-1}t_jU_g^{-*})^{-1}\,dg.}
\end{equation}
The answer is unchanged by arbitrary, element-dependent replacements
$U_g\mapsto\zeta_gU_g$.  It therefore depends only on the projective image,
not on a Schur cover, cocycle representative, scalar phases, or a
determinant-one lift.
\end{theorem}

\begin{proof}
The coefficient indexed by $(\alpha,\beta)$ in the integrand acquires
$u^{|\alpha|-|\beta|}$.  The constant term imposes equality of the two
degrees.  Under $U_g\mapsto\zeta_gU_g$, that coefficient is multiplied by
$\zeta_g^{|\alpha|-|\beta|}=1$.  Haar averaging then gives the invariant
dimension exactly as in Theorem~\ref{nd:central_extensions:thm:molien}.
\end{proof}

For a finite scalar center $\mu_r$, the full two-sided series has sectors
$|\alpha|-|\beta|\equiv0\pmod r$.  Exact equality, rather than congruence,
is the part intrinsic to the projective representation.  The multilinear
diagonal coefficients are the frame potentials
\begin{equation}\label{nd:central_extensions:eq:frame}
 [m_{(1^k)}(\mathbf s)m_{(1^k)}(\mathbf t)]\BSMG^{\mathrm{proj}}
 =\int_G|\Tr U_g|^{2k}\,dg.
\end{equation}

\subsection{Finite groups: class compression and power maps}

If $\widetilde G$ is finite with character $\chi_V$, grouping
\eqref{nd:central_extensions:eq:trace-exp} by conjugacy classes gives
\begin{equation}\label{nd:central_extensions:eq:class-formula}
 \boxed{\MGF_{\widetilde G,V}
 =\frac1{|\widetilde G|}\sum_{C}|C|
 \exp\!\left(\sum_{k\geq1}\frac{p_k(\mathbf t)}{k}
                         \chi_V(c_C^k)\right).}
\end{equation}
Thus an exact calculation needs the represented character, class sizes, and
class power maps.  The projective balanced version is
\begin{equation}\label{nd:central_extensions:eq:balanced-class}
 \CT_u\frac1{|\widetilde G|}\sum_C|C|\exp\!\left[
 \sum_{k\geq1}\frac{u^kp_k(\mathbf s)\chi_V(c_C^k)
 +u^{-k}p_k(\mathbf t)\overline{\chi_V(c_C^k)}}{k}\right].
\end{equation}
Equations \eqref{nd:central_extensions:eq:class-formula} and \eqref{nd:central_extensions:eq:balanced-class} are the
practical route for large Schur covers, spin Weyl covers, and sporadic
covers whose dense matrices would be wasteful.

\subsection{Schur and spin covers: the binary octahedral test}

For a genuine spin module of a double cover $2.S_n$ or $2.A_n$, the central
involution acts by $-I$.  Theorem~\ref{nd:central_extensions:thm:projector} immediately gives
\begin{equation}\label{nd:central_extensions:eq:spin-parity}
 [m_\tau]\MGF_{2.S_n,V_n}=[m_\tau]\MGF_{2.A_n,V_n}=0
 \qquad(|\tau|\text{ odd}).
\end{equation}

For the standard basic-spin module, Clifford algebra identifies
$\End(V_n)$ with all or one parity of the exterior algebra of the reflection
module $E_n$.  Its exterior powers are distinct hook modules.  With the
standard associate convention this gives
\begin{equation}\label{nd:central_extensions:eq:basic-spin-fourth}
 \frac1{|2.S_n|}\sum_g|\chi_{V_n}(g)|^4
 =\begin{cases}n,&n\text{ odd},\\n/2,&n\text{ even}.
 \end{cases}
\end{equation}
The fourth balanced moment therefore grows linearly although $\dim V_n$
grows exponentially; it does not approach the Haar value $2$.

The notebook uses the first nontrivial concrete model relevant here:
the $48$ unit quaternions of the binary octahedral group.  Under
\[
 a+bi+cj+dk\longmapsto
 \begin{pmatrix}a+bi&c+di\\-c+di&a-bi\end{pmatrix}\in SU(2),
\]
they form a double cover $2.S_4\to S_4$ and give the two-dimensional
basic-spin representation.  The list consists of
\[
 \{\pm1,\pm i,\pm j,\pm k\},\qquad
 \tfrac12(\pm1\pm i\pm j\pm k),\qquad
 \tfrac1{\sqrt2}(\pm e_a\pm e_b)\ (a<b).
\]
All signs are independent in the displayed families.  Dense multiplication
checks all $48^2$ products.  The fourth trace moment is $2$, agreeing with
\eqref{nd:central_extensions:eq:basic-spin-fourth} at $n=4$.

\begin{center}
\begin{tabular}{@{}c c c c c@{}}
\toprule
$k$ & direct $|\Tr|^{2k}$ & phase-twisted lifts & Haar $U(2)$ & conclusion\\
\midrule
1 & 1 & 1 & 1 & agree\\
2 & 2 & 2 & 2 & agree\\
3 & 5 & 5 & 5 & agree\\
4 & 15 & 15 & 14 & first separation\\
\bottomrule
\end{tabular}
\end{center}

The phase-twisted column multiplies the $j$th matrix by an unrelated
$37$th root of unity.  This is a direct test that projective balanced
quantities do not depend on a coherent linear lift.  The non-multilinear
test
\[
 (\alpha,\beta)=((2,1),(1,1,1))
\]
also remains $3$ after phase twisting, checking charge-zero invariance beyond
frame potentials.

\subsection{Pin covers of Weyl groups}

Let $W\leq O(E)$ be a finite Weyl or Coxeter group,
$\widetilde W\leq\operatorname{Pin}(E)$ its pin cover, and $S$ a genuine
spinor.  The central involution gives the parity rule
\[
 [m_\tau]\MGF_{\widetilde W,S}=0\qquad(|\tau|\text{ odd}).
\]
Moreover
\begin{equation}\label{nd:central_extensions:eq:cliff-end}
 \End(S)\cong\Cl(E)\ \text{or}\ \Cl^0(E)
 \cong\bigwedge^\bullet E\ \text{or}\ \bigwedge^{\mathrm{even}}E
\end{equation}
as $W$-modules.  Hence
\begin{equation}\label{nd:central_extensions:eq:weyl-moment}
 \boxed{\frac1{|\widetilde W|}\sum_{\widetilde w}
 |\Tr S(\widetilde w)|^{2k}
 =\dim\left(\End(S)^{\otimes k}\right)^W.}
\end{equation}
For $k=2$, if $\End(S)=\bigoplus_\lambda m_\lambda U_\lambda$, the answer
is $\sum_\lambda m_\lambda^2$.  Distinct exterior powers in the classical
towers therefore produce a rank-growing exterior-algebra law, not a Haar
unitary law.  The notebook's $Q_8$ and binary octahedral constructions are
the rank-small pin/spin test cases; higher ranks are more efficiently tested
by \eqref{nd:central_extensions:eq:cliff-end} or class power maps.

\subsection{Finite Weil representations}

Assume either that $q$ is odd, or more generally that a Weyl operator basis
$\{D_v:v\in W_n\}$ of $\End(\mathcal H_n)$ and unitary lifts $U_g$ have
been chosen with exact covariance. Let $W_n=\F_q^{2n}$ and let
$\mathcal H_n$ be the full finite Schr\"odinger space,
$\dim\mathcal H_n=q^n$. The required covariance is
\[
 U_gD_vU_g^{-1}=D_{gv}
\]
after a compatible normalization.  Since the $D_v$ form a basis of
$\End(\mathcal H_n)$,
\begin{equation}\label{nd:central_extensions:eq:weil-end}
 \End(\mathcal H_n)\cong\C[W_n]
\end{equation}
as an $\Sp(2n,q)$-module.  Burnside's orbit lemma gives
\begin{equation}\label{nd:central_extensions:eq:weil-orbits}
 \boxed{\frac1{|\Sp(2n,q)|}\sum_g|\Tr U_g|^{2k}
 =\#\bigl(\Sp(2n,q)\backslash W_n^k\bigr).}
\end{equation}

\proofstatus{Exact under covariance}{The orbit formula applies to the full
Schr\"odinger module. It is not inferred from the name ``Weil
representation'' alone, and it must be modified for an irreducible even or
odd constituent.}

There are two orbits on $W_n$: zero and nonzero, so the second trace moment
is $2$.  For ordered pairs and $n\geq2$, the orbit types are
\begin{itemize}
\item $(0,0)$, $(0,w\ne0)$, and $(v\ne0,0)$;
\item $w=av$, one orbit for each $a\in\F_q^\times$;
\item independent orthogonal pairs;
\item independent pairs with $\langle v,w\rangle=c$, one orbit for each
      $c\in\F_q^\times$.
\end{itemize}
Thus the stable fourth moment is $3+(q-1)+1+(q-1)=2q+2$.
Witt extension proves that every fixed-$k$ orbit count stabilizes once the
ambient symplectic rank is large enough.

\begin{remark}[Rank-one correction]
When $n=1$, an orthogonal pair in a two-dimensional symplectic space is
dependent.  The independent-orthogonal orbit is absent, so
\begin{equation}\label{nd:central_extensions:eq:rank-one-weil}
 \mathbb E|\Tr U_g|^4=2q+1\qquad(n=1).
\end{equation}
This hypothesis is essential when testing a reasonable small dimension.
\end{remark}

For odd prime $q$, the notebook independently constructs the projective Weil
matrices from
\[
 F_{xy}=q^{-1/2}e^{2\pi ixy/q},\qquad
 P_{xx}=e^{2\pi i(2^{-1}x^2)/q}.
\]
Projective closure yields $|SL(2,3)|=24$ matrices in dimension $3$ and
$|SL(2,5)|=120$ matrices in dimension $5$.  Separately, the code enumerates
the $SL(2,q)$-orbits on $(\F_q^2)^k$.

\begin{center}
\begin{tabular}{@{}c c c c c@{}}
\toprule
$q$ & projective order & $k$ & matrix average & orbit count\\
\midrule
3 & 24  & 1 & 2 & 2\\
3 & 24  & 2 & 7 & 7\\
5 & 120 & 1 & 2 & 2\\
5 & 120 & 2 & 11 & 11\\
\bottomrule
\end{tabular}
\end{center}

For an irreducible even or odd Weil constituent, one must insert the
corresponding idempotent in \eqref{nd:central_extensions:eq:weil-end}; the full-module orbit count
must not be copied unchanged.  Likewise, over finite fields where a
metaplectic extension splits, the actual scalar image, rather than the name
of the cover, determines the selection rule.

\subsection{Odd extraspecial groups and their central products}

Let $E_n(p)$ be an extraspecial group of order $p^{1+2n}$ and exponent $p$,
where $p$ is odd.  Its Schr\"odinger module has dimension $N=p^n$.  In the
explicit model used by the notebook, for $a,b\in\F_p^n$ and $c\in\F_p$,
\begin{equation}\label{nd:central_extensions:eq:heis-matrix}
 \rho(c,a,b)e_x=\omega^{c+b\cdot x}e_{x+a},
 \qquad \omega=e^{2\pi i/p}.
\end{equation}
These $pN^2$ matrices are unitary and faithful.

\begin{theorem}[Odd extraspecial formula]\label{nd:central_extensions:thm:extra-odd}
For the module \eqref{nd:central_extensions:eq:heis-matrix},
\begin{equation}\label{nd:central_extensions:eq:extra-series}
 \boxed{\MGF_{E_n(p),V}
 =N^{-2}\PP_p\prod_i(1-t_i)^{-N}
 +(1-N^{-2})\prod_i(1-t_i^p)^{-N/p}.}
\end{equation}
Consequently,
\begin{align}\label{nd:central_extensions:eq:extra-coeff}
 [m_\tau]\MGF_{E_n(p),V}
 =\mathbf1_{p\mid|\tau|}\biggl[&N^{-2}
 \prod_i\binom{N+\tau_i-1}{\tau_i}\\
 &+(1-N^{-2})\mathbf1_{p\mid\tau_i\ \forall i}
 \prod_i\binom{N/p+\tau_i/p-1}{\tau_i/p}\biggr].\nonumber
\end{align}
If scalar $s$th roots are adjoined with $p\mid s$, then
\begin{equation}\label{nd:central_extensions:eq:central-product}
 \MGF_{\mu_sE_n(p),V}=\PP_s\MGF_{E_n(p),V}.
\end{equation}
\end{theorem}

\begin{proof}
The center consists of $p$ scalar matrices $\omega^cI$, a fraction $N^{-2}$
of the group.  Averaging their determinant factors gives the first term of
\eqref{nd:central_extensions:eq:extra-series}.  For a noncentral element, the weighted shift
\eqref{nd:central_extensions:eq:heis-matrix} has all $p$th roots of unity as eigenvalues, each with
multiplicity $N/p$.  Therefore
$\det(I-tg)=(1-t^p)^{N/p}$, giving the second term.  Extracting coefficients
with $(1-t)^{-a}=\sum_{d\geq0}\binom{a+d-1}{d}t^d$ proves
\eqref{nd:central_extensions:eq:extra-coeff}.  Equation \eqref{nd:central_extensions:eq:central-product} is
Theorem~\ref{nd:central_extensions:thm:projector}.
\end{proof}

The notebook constructs $E_1(3)$, $E_2(3)$, and $E_1(5)$, of respective
orders $27$, $243$, and $125$ and dimensions $3$, $9$, and $5$.  Every
coefficient through total degree six agrees with \eqref{nd:central_extensions:eq:extra-coeff}.
At $(t_1,t_2)=(0.07,0.11)$ the full determinant averages are:
\begin{center}
\begin{tabular}{@{}c c c c@{}}
\toprule
group & dense matrix average & formula & absolute error\\
\midrule
$E_1(3)$ & 1.006152412632117 & 1.006152412632117 & $7.8\times10^{-18}$\\
$E_2(3)$ & 1.016020200753344 & 1.016020200753347 & $2.7\times10^{-15}$\\
$E_1(5)$ & 1.000511444770786 & 1.000511444770785 & $1.6\times10^{-15}$\\
\bottomrule
\end{tabular}
\end{center}

The trace is zero off the center and equals $N\omega^c$ on the central
element $\omega^cI$.  Hence
\begin{equation}\label{nd:central_extensions:eq:trace-spike}
 \mathbb E[T^a\overline T^{\,b}]
 =\mathbf1_{p\mid(a-b)}N^{a+b-2},
\end{equation}
so $\mathbb E|T|^2=1$ while
$\mathbb E|T|^{2k}=N^{2k-2}$ for $k\geq2$.  Thus $T\to0$ in probability
as $n\to\infty$, but its higher balanced moments diverge: the limit is a
rare central spike, not a Gaussian.

For fixed $\tau$ and $d=|\tau|$, the central term of
\eqref{nd:central_extensions:eq:extra-coeff} is asymptotic to
$N^{d-2}/\prod_i\tau_i!$.  When all $\tau_i$ are divisible by $p$, the
noncentral term is asymptotic to
$(N/p)^{d/p}/\prod_i(\tau_i/p)!$.  Raw coefficients therefore generally
have no finite coefficientwise limit.

\subsection{Extraspecial \texorpdfstring{$2$}{2}-groups: separate plus and minus sectors}

Let $E_n^\varepsilon$ be extraspecial of order $2^{1+2n}$, sign
$\varepsilon\in\{+1,-1\}$, and let $N=2^n$.  Central elements again give
the root-filtered sector.  Noncentral elements whose squares are $+I$ have
eigenvalues $\pm1$ equally often; those whose squares are $-I$ have
eigenvalues $\pm i$ equally often.  Counting the two quadratic-form types
gives
\begin{equation}\label{nd:central_extensions:eq:extra-two}
 \boxed{\begin{aligned}
 \MGF_{E_n^\varepsilon}={}&
 2^{-2n}\PP_2\prod_i(1-t_i)^{-N}\\
 &+\left(\frac12+\varepsilon2^{-n-1}-2^{-2n}\right)
       \prod_i(1-t_i^2)^{-N/2}\\
 &+\left(\frac12-\varepsilon2^{-n-1}\right)
       \prod_i(1+t_i^2)^{-N/2}.
 \end{aligned}}
\end{equation}
The final two sectors can cancel in low degree; they may not be combined by
discarding the sign.

The rank-one plus group is $D_8$, represented by $\pm I,\pm X,\pm Z,\pm XZ$.
The rank-one minus group is $Q_8$, represented by
$\pm I,\pm iX,\pm iZ,\pm(-XZ)$.  Both $8$-element sets pass exact closure.
All coefficients through degree six agree with \eqref{nd:central_extensions:eq:extra-two}.  At the
same determinant test point, the errors are below $3.6\times10^{-18}$.

\subsection{Clifford groups and determinant-one lifts}

The full unitary Clifford normalizer contains $U(1)I$, so
Corollary~\ref{nd:central_extensions:cor:u1} gives
\[
 \MGF_{\mathrm{Cliff}_n}=1
\]
for the one-sided series.  A finite phase lift with scalar image $\mu_s$
instead applies $\PP_s$ and has lift-dependent coefficients in degrees
divisible by $s$.

Let $K_n=\mathrm{Cliff}_n\cap SU(N)$ with $N=2^n$.  It contains
$\{\zeta I:\zeta^N=1\}\cong\mu_N$, and therefore
\begin{equation}\label{nd:central_extensions:eq:cliff-selection}
 [m_\tau]\MGF_{K_n}=0\qquad\text{unless }N\mid|\tau|.
\end{equation}
For every fixed positive degree, eventually $N>|\tau|$, whence
\begin{equation}\label{nd:central_extensions:eq:cliff-limit}
 \MGF_{K_n}\longrightarrow1\qquad\text{coefficientwise}.
\end{equation}

In the one-qubit case, $K_1$ is the binary octahedral group already tested.
Adjoining $\mu_4$ produces $96$ distinct unitary matrices; all coefficients
through total degree six equal the binary-octahedral coefficient multiplied
by $\mathbf1_{4\mid|\tau|}$.  This directly tests the finite-lift rule.

Balanced quantities are the same for the full normalizer, any finite scalar
lift, $K_n$, and the projective quotient.  Multiqubit Clifford groups are
unitary $3$-designs, so their balanced series agrees with Haar unitary
through balanced degree three and differs at degree four.  The one-qubit
row in the table above gives the concrete values $1,2,5,15$ versus Haar
$1,2,5,14$.

More generally, suppose $\delta:\widetilde G\to\mu_q$ is a character and
$K=\ker\delta$.  Define
\[
 \MGF_{\widetilde G}^{(j)}
 =\frac1{|\widetilde G|}\sum_{g\in\widetilde G}\delta(g)^j
 \prod_i\det(I-t_i\rho(g))^{-1}.
\]
Character orthogonality gives the exact determinant-one formula
\begin{equation}\label{nd:central_extensions:eq:kernel-fourier}
 \boxed{\MGF_{K,V}=\sum_{j=0}^{q-1}\MGF_{\widetilde G}^{(j)}.}
\end{equation}
Indeed, $\sum_j\delta(g)^j$ equals $q$ on $K$ and zero off $K$, while
$|\widetilde G|=q|K|$ when $\delta$ is onto.  Passing to determinant one
therefore adds determinant-twisted sectors; it does not merely ``remove
scalars.''

\subsection{Sporadic covers: exact data workflow}

For a genuine irreducible module of a sporadic cover $2.G$, $3.G$, or $6.G$,
let $r$ be the order of its central scalar character.  Then
\[
 [m_\tau]\MGF_{\widetilde G,V}=0\qquad\text{unless }r\mid|\tau|,
\]
and the complete exact answer is \eqref{nd:central_extensions:eq:class-formula}.  An auditable
calculation imports three items from the represented covering group:
\begin{enumerate}
\item conjugacy-class sizes;
\item the genuine character values on those classes;
\item every class power map needed to the chosen truncation degree.
\end{enumerate}
The verification then compares a class-compressed sum against an element
enumeration when matrices are available, or against an independent character
table implementation otherwise.  There is no intrinsic rank asymptotic for
an isolated sporadic group.  Useful comparisons are the first separating
degree, balanced design strength, and equality of charge-zero series across
different lifts.

No external sporadic character table was supplied with this workspace, so
the notebook deliberately does not invent a sporadic numerical example.  The
binary octahedral and finite Weil cases exercise the same class/power and
phase mechanisms with explicit matrices; inserting certified ATLAS data in
\eqref{nd:central_extensions:eq:class-formula} is the correct sporadic verification strategy.

\subsection{Verification method and reproducibility}

The companion files are
\begin{itemize}
\item \texttt{verification.py}: constructions and comparison routes;
\item \path{marimo_notebooks/central_extensions.py}: the marimo interface,
grouped by group family;
\item \texttt{test\_verification.py}: a headless regression test.
\end{itemize}

For each dense matrix $g$, the code computes
$\Tr(\Sym^d g)$ from the recurrence
\begin{equation}\label{nd:central_extensions:eq:newton}
 h_0(g)=1,\qquad
 d\,h_d(g)=\sum_{k=1}^{d}\Tr(g^k)h_{d-k}(g).
\end{equation}
Thus
\[
 [m_\tau]\MGF=\frac1{|G|}\sum_{g\in G}\prod_i h_{\tau_i}(g)
\]
is obtained directly from actual matrix powers.  It is compared with
\eqref{nd:central_extensions:eq:extra-coeff}, coefficient extraction from \eqref{nd:central_extensions:eq:extra-two}, or
the scalar projector.  This is independent of the determinant-sector
comparison, which evaluates the full rational functions at
$(0.07,0.11)$.

For the projective tests, frame potentials are computed from the generated
matrices while the Weil comparator counts finite symplectic orbits.  The
binary octahedral matrices are also phase-twisted element by element.  Matrix
audits check expected order, distinctness, identity, unitarity, and exact
closure where a quadratic scan is modest.  The $243$-element $E_2(3)$ case
uses the exact Heisenberg multiplication law from \eqref{nd:central_extensions:eq:heis-matrix}
instead of a redundant $243^2$ dense closure scan.

\begin{center}
\begin{tabularx}{\linewidth}{@{}l r r r X@{}}
\toprule
family & groups & dimensions & checks & independent comparator\\
\midrule
odd extraspecial & 3 & $3,9,5$ & $90+3$ & spectral sectors and binomial formula\\
extraspecial $2$ & 2 & $2,2$ & $60+2$ & signed $\pm1/\pm i$ sectors\\
Schur/scalar lift & 2 & $2,2$ & 30 & homogeneous $\mu_4$ projector\\
projective spin/Clifford & 1 image & 2 & 5 & phase changes and Haar values\\
finite Weil & 2 & $3,5$ & 4 & symplectic orbit enumeration\\
\midrule
total & 9 constructions & $2$--$9$ & 194 & all pass\\
\bottomrule
\end{tabularx}
\end{center}

Here ``194'' means 180 coefficient, five determinant, and nine balanced
comparisons.  The maximum residual, $6.28\times10^{-12}$, comes from
floating-point roots of unity in an integer-valued coefficient; all stated
integer answers are within the declared $2\times10^{-9}$ tolerance.

\subsection{Additional Schur-cover audit for genuine modules}

Let
\[
 1\longrightarrow\{1,z\}\longrightarrow\widetilde S_n
 \longrightarrow S_n\longrightarrow1
\]
be a Schur double cover, and let $D$ be a genuine module, so $z$ acts by
$-I$.  Theorem~\ref{nd:symmetric_spin_lie_tree:thm:molien}, grouped by conjugacy classes, gives
\begin{equation}\label{nd:symmetric_spin_lie_tree:eq:schur-cover}
 \boxed{
 \MGF_{\widetilde S_n,D}
 =\sum_{[\widetilde g]}
 \frac1{|C_{\widetilde S_n}(\widetilde g)|}
 \exp\!\left(\sum_{r\ge1}\frac{p_r}{r}
 \chi_D(\widetilde g^r)\right).}
\end{equation}
Pairing the summands $g$ and $zg$ shows immediately that
\begin{equation}\label{nd:symmetric_spin_lie_tree:eq:central-parity}
 \boxed{[m_\tau]\MGF_{\widetilde S_n,D}=0
 \quad\text{when }|\tau|\text{ is odd}.}
\end{equation}

If $D\otimes D\cong\bigoplus_\eta m_\eta U_\eta$ and the required
self-duality identifies fourth invariants with
$\operatorname{End}_{\widetilde S_n}(D\otimes D)$, Schur's lemma gives
\begin{equation}\label{nd:symmetric_spin_lie_tree:eq:spin-fourth}
 \dim(D^{\otimes4})^{\widetilde S_n}=\sum_\eta m_\eta^2.
\end{equation}
Without that self-duality, the always-valid statement pairs multiplicities
of dual constituents instead.  This qualification is essential when the
chosen genuine constituent is not self-dual.

For the audit, Hermitian Clifford matrices $\gamma_i$ satisfy
$\gamma_i\gamma_j+\gamma_j\gamma_i=2\delta_{ij}I$.  The lifts
\[
 u_i=\frac{\gamma_i-\gamma_{i+1}}{\sqrt2},\qquad 1\le i<4,
\]
generate the Pin preimage of $S_4$, an order-$48$ matrix group containing
$-I$.  The represented module has dimension four.  Direct generator-fixed
spaces and \eqref{nd:symmetric_spin_lie_tree:eq:schur-cover} give
\begin{center}
\begin{tabular}{@{}crrrrr@{}}\toprule
$\tau$&$(1)$&$(2)$&$(1,1)$&$(3)$&$(1,1,1,1)$\\\midrule
coefficient&0&1&2&0&16\\\bottomrule
\end{tabular}
\end{center}
All five comparisons pass, the odd rows vanish as predicted, and the maximum
Clifford-relation residual is zero at machine precision.  This test verifies
the universal formula and central selection rule for a genuine matrix module;
it does not purport to tabulate every strict-partition spin character.


\subsection{Asymptotic taxonomy and conclusions}

\begin{center}
\small
\begin{tabularx}{\linewidth}{@{}l X X@{}}
\toprule
family & one-sided behavior & balanced/projective behavior\\
\midrule
$2.S_n,2.A_n$ spin & odd degrees vanish & basic-spin fourth moment is linear in $n$\\
spin Weyl covers & parity projection & exterior-algebra centralizer law\\
sporadic covers & degree modulo central character & exact finite class formula; no intrinsic tower\\
finite Weil modules & lift-dependent scalar rule & fixed-$k$ symplectic orbit law stabilizes\\
extraspecial $p$-groups & coefficients generally diverge & rare central spikes; high moments diverge\\
full Clifford normalizer & exactly $1$ & Haar agreement through design degree\\
determinant-one Clifford & coefficientwise limit $1$ & same law as the projective quotient\\
scalar quotient & not canonically defined & charge zero is canonical\\
\bottomrule
\end{tabularx}
\end{center}

Central extensions therefore produce three distinct mechanisms:
\[
\boxed{
\begin{array}{ll}
\textbf{finite scalar center:}&\text{congruence projection on total degree};\\[1mm]
\textbf{growing or continuous scalar center:}&\text{one-sided collapse to }1;\\[1mm]
\textbf{projective representation:}&\text{information survives canonically at charge zero.}
\end{array}}
\]
The natural asymptotic objects are consequently projective or charge-graded
$\Lambda$-distributions, not phase-dependent one-sided distributions on a
scalar quotient.

\clearpage
\section{Pauli--Heisenberg groups}\label{proof:pauli-heisenberg}
\begingroup
\definecolor{navy}{HTML}{17365D}
\definecolor{teal}{HTML}{147D78}
\providecommand{\Sgm}{\SMG}
\renewcommand{\Sgm}{\SMG}
\providecommand{\E}{\mathbb E}
\renewcommand{\E}{\mathbb E}
\providecommand{\F}{\mathbb F}
\renewcommand{\F}{\mathbb F}
\setcounter{theorem}{0}
\setcounter{proposition}{0}
\setcounter{lemma}{0}
\setcounter{corollary}{0}
\setcounter{remark}{0}
\begin{abstract}
For odd characteristic we prove the two-spectrum formula for the defining
Schrödinger representation and check it by constructing every matrix in
$P_{3,1}$, $P_{5,1}$, $P_{9,1}$, and $H_3(\F_9)$.  The tests evaluate the
target invariant dimension directly, not from the spectral classification
used in the proof.  All 17 representative coefficient comparisons pass.
\end{abstract}

\subsection{Statement}
Let $q=p^f$ with $p$ odd, $V=\mathbb C[\F_q^n]$, and $N=q^n$.  Fix a
nontrivial additive character $\psi:\F_q\to\mu_p$.  With
\[
X(a)e_x=e_{x+a},\qquad Z(b)e_x=\psi(b\cdot x)e_x,
\qquad W(a,b)=\psi\!\left(\tfrac12a\cdot b\right)X(a)Z(b),
\]
put $P_{q,n}=\{\zeta W(a,b):\zeta\in\mu_p\}$.  For a partition
$\tau=(\tau_1,\ldots,\tau_\ell)$ define
\[
A_\tau=\prod_i\binom{N+\tau_i-1}{\tau_i},\qquad
B_\tau=\begin{cases}
\displaystyle\prod_i\binom{N/p+\tau_i/p-1}{\tau_i/p},&p\mid\tau_i\ \forall i,\\
0,&\text{otherwise.}
\end{cases}
\]
Then the claimed generalized Molien coefficient is
\begin{equation}\label{pauli-heisenberg-proof-and-verification:eq:coefficient}
[m_\tau]\Sgm_{P_{q,n}}=
\begin{cases}
0,&p\nmid |\tau|,\\
q^{-2n}A_\tau+(1-q^{-2n})B_\tau,&p\mid |\tau|.
\end{cases}
\end{equation}
Equivalently,
\begin{equation}\label{pauli-heisenberg-proof-and-verification:eq:series}
\Sgm_{P_{q,n}}(\mathbf t)=
\frac{q^{-2n}}p\sum_{\zeta\in\mu_p}\prod_j(1-\zeta t_j)^{-N}
+(1-q^{-2n})\prod_j(1-t_j^p)^{-N/p}.
\end{equation}

\subsection{Proof}
There are $p$ central matrices $\zeta I_N$, a fraction $q^{-2n}$ of the
group.  Their average is the first term of \eqref{pauli-heisenberg-proof-and-verification:eq:series}.  Now take
$(a,b)\ne(0,0)$.  For $1\le k<p$, $W(a,b)^k$ has trace zero: if $ka\ne0$
it is a translation with zero diagonal; otherwise its diagonal trace is a
nontrivial additive-character sum.  Also $W(a,b)^p=I$.  Fourier inversion of
the multiplicities of the $p$th roots of unity therefore gives multiplicity
$N/p$ for every root.  Hence
\[
\det(I-tW(a,b))=(1-t^p)^{N/p}.
\]
A scalar in $\mu_p$ only permutes this spectrum, proving \eqref{pauli-heisenberg-proof-and-verification:eq:series}.
Expanding the central factors gives $A_\tau$ and forces $p\mid|\tau|$;
expanding the noncentral factors gives $B_\tau$.  This proves
\eqref{pauli-heisenberg-proof-and-verification:eq:coefficient}.

For the finite Heisenberg group $H_{2n+1}(\F_q)$, the Schrödinger matrix is
$\rho_\psi(a,b,z)=\psi(z)W(a,b)$.  When $f>1$, $\ker\psi$ has size $q/p$,
so each matrix in $P_{q,n}$ has exactly $q/p$ preimages.  Uniform averages
over the abstract group and its image consequently agree.

\subsection{Independent verification method}
For every constructed matrix $g$, the code computes its eigenvalues
$\lambda_1,\ldots,\lambda_N$ numerically and obtains $h_d(\lambda)$ from
\[
\prod_{r=1}^N(1-\lambda_r t)^{-1}=\sum_{d\ge0}h_d(\lambda)t^d.
\]
It then evaluates the target directly as
\[
D_\tau=\frac1{|G|}\sum_{g\in G}\prod_i h_{\tau_i}(\operatorname{spec}g).
\]
This is the character of $\bigotimes_i\operatorname{Sym}^{\tau_i}V$ averaged
over $G$, hence the rank of its Reynolds projector.  Only after computing
$D_\tau$ does the program evaluate \eqref{pauli-heisenberg-proof-and-verification:eq:coefficient}.

The $q=9$ construction uses $\F_9=\F_3[u]/(u^2+1)$ and
$\psi(x)=\exp(2\pi i\operatorname{Tr}_{\F_9/\F_3}(x)/3)$.  Thus the test
includes a genuine prime-power field.  The $H_3(\F_9)$ run enumerates all
729 abstract triples, retaining the threefold central kernel rather than
silently replacing the group by its image.

\subsection{Representative results}
\begin{center}
\begin{tabular}{@{}lrrrr@{}}\toprule
group & $\dim V$ & matrices averaged & $\tau$ & direct $=$ formula\\\midrule
$P_{3,1}$ &3&27&$(3)$&2\\
$P_{3,1}$ &3&27&$(3,3)$&12\\
$P_{5,1}$ &5&125&$(4,1)$&14\\
$P_{5,1}$ &5&125&$(5,5)$&636\\
$P_{9,1}$ &9&243&$(3)$&5\\
$P_{9,1}$ &9&243&$(3,3)$&345\\
$H_3(\F_9)$ &9&729&$(2,1)$&5\\
$H_3(\F_9)$ &9&729&$(3,3)$&345\\\bottomrule
\end{tabular}
\end{center}
All 17 built-in cases pass to tolerance $2\times10^{-8}$ and round to
nonnegative integers.  Characteristic two is intentionally excluded here:
the displayed half-phase normalization is an odd-characteristic statement;
the phase-extended qubit lift is handled in its own verification.

\subsection{Reproduction}
Run
\begin{verbatim}
.venv/bin/python for_this_guy/quantum_matrix_groups/
  pauli_heisenberg/verification.py
.venv/bin/marimo edit marimo_notebooks/pauli_heisenberg.py
\end{verbatim}
The line breaks after the directory slash are for display; join each command
onto one shell line.
\endgroup


\clearpage
\section{Scalar-extended Pauli groups}\label{proof:scalar-pauli}
\begingroup
\definecolor{navy}{HTML}{17365D}
\definecolor{teal}{HTML}{147D78}
\definecolor{red}{HTML}{9B2C2C}
\providecommand{\Sgm}{\SMG}
\renewcommand{\Sgm}{\SMG}
\setcounter{theorem}{0}
\setcounter{proposition}{0}
\setcounter{lemma}{0}
\setcounter{corollary}{0}
\setcounter{remark}{0}
\begin{abstract}
Building on Theorem A(3) of the 14 July paper, we prove the coefficient form of the scalar-extended Pauli formula and test
the complete matrix groups $\mu_4P_{2,1}$, $\mu_4P_{2,2}$, and
$\mu_6P_{3,1}$.  All 13 representative checks pass.  The calculation also
locates a sign error in the supplied qubit-recovery prose: equation (13)
specializes with a positive, projected noncentral sector, not a minus sign.
\end{abstract}

\subsection{Formula and proof}
Let $q=p^f$, $N=q^n$, and let $s$ be divisible by $p$.  For the finite lift
$\widetilde P_{q,n}^{(s)}=\mu_sP_{q,n}$, its distinct matrices may be indexed
by $\alpha\in\mu_s$ and $v\in\mathbb F_q^{2n}$, so its order is $sq^{2n}$.
For odd $p$ use the usual phase-normalized Weyl operators.  For qubits we use
$\{I,X,Z,iXZ\}$ on each tensor factor and adjoin $\mu_4$.

If $v=0$, the matrix is $\alpha I_N$.  If $v\ne0$, the $p$ spectral branches
of the phase-normalized $W(v)$ occur with multiplicity $N/p$, and therefore
\[
\det(I-t\alpha W(v))=(1-\alpha^p t^p)^{N/p}.
\]
Averaging these two sectors proves
\begin{equation}\label{scalar-extended-pauli-proof-and-verification:eq:scalar}
\Sgm_{\widetilde P}(\mathbf t)=
q^{-2n}\frac1s\sum_{\alpha\in\mu_s}\prod_j(1-\alpha t_j)^{-N}
+(1-q^{-2n})\frac1s\sum_{\alpha\in\mu_s}
\prod_j(1-\alpha^pt_j^p)^{-N/p}.
\end{equation}
Root-of-unity orthogonality forces $s\mid|\tau|$ in both sectors.  Thus,
with
\[
A_\tau=\prod_i\binom{N+\tau_i-1}{\tau_i},\qquad
B_\tau=\mathbf1_{p\mid\tau_i\ \forall i}
\prod_i\binom{N/p+\tau_i/p-1}{\tau_i/p},
\]
the coefficient formula is
\begin{equation}\label{scalar-extended-pauli-proof-and-verification:eq:coeff}
[m_\tau]\Sgm_{\widetilde P}=\begin{cases}
0,&s\nmid|\tau|,\\
q^{-2n}A_\tau+(1-q^{-2n})B_\tau,&s\mid|\tau|.
\end{cases}
\end{equation}

\subsection{Qubit specialization and correction}
For $p=q=2$, $s=4$, equation \eqref{scalar-extended-pauli-proof-and-verification:eq:scalar} gives
\[
4^{-n}\Pi_4\prod_j(1-t_j)^{-2^n}
+(1-4^{-n})\frac14\sum_{\alpha\in\mu_4}
\prod_j(1-\alpha^2t_j^2)^{-2^{n-1}}.
\]
The second summand is an average of determinant reciprocals and has positive
coefficients in allowed multidegrees.  In particular, for one qubit,
\[
[t^4]\Sgm=\tfrac14\binom{5}{4}+\tfrac34\binom{2}{2}=2.
\]
The supplied “qubit recovery” paragraph instead writes a minus before the
noncentral term, which would give $1/2$, not even an integer invariant
dimension.  Therefore that minus sign cannot be correct.  Equation (13)
itself, interpreted as the root-of-unity average above, is consistent with
the direct matrices and with \eqref{scalar-extended-pauli-proof-and-verification:eq:coeff}.

\subsection{Verification strategy and results}
For each matrix the verifier computes the coefficients $h_d$ of
$\det(I-tg)^{-1}$ from its eigenvalues, then independently averages
$\prod_i h_{\tau_i}(g)$ over every group matrix.  It compares this character
average with \eqref{scalar-extended-pauli-proof-and-verification:eq:coeff}.  Representative results are:
\begin{center}
\begin{tabular}{@{}lrrr@{}}\toprule
group & $\tau$ & direct & formula\\\midrule
$\mu_4P_{2,1}$ & $(2)$&0&0\\
$\mu_4P_{2,1}$ & $(4)$&2&2\\
$\mu_4P_{2,1}$ & $(2,2)$&3&3\\
$\mu_4P_{2,2}$ & $(4)$&5&5\\
$\mu_4P_{2,2}$ & $(4,4)$&85&85\\
$\mu_6P_{3,1}$ & $(3)$&0&0\\
$\mu_6P_{3,1}$ & $(6)$&4&4\\
$\mu_6P_{3,1}$ & $(5,1)$&7&7\\\bottomrule
\end{tabular}
\end{center}
All 13 built-in comparisons pass to $2\times10^{-8}$.  The cases $(2)$ for
$s=4$ and $(3)$ for $s=6$ directly exercise the stronger scalar selection
rule, while $(3,1)$ and $(5,1)$ test mixed symmetric powers.

\subsection{Reproduction}
\begin{verbatim}
.venv/bin/python for_this_guy/quantum_matrix_groups/
  scalar_extended_pauli/verification.py
.venv/bin/marimo edit marimo_notebooks/scalar_extended_pauli.py
\end{verbatim}
Join each displayed command onto one shell line.
\endgroup


\clearpage
\section{Extraspecial groups}\label{proof:extraspecial}
\begingroup
\definecolor{navy}{HTML}{17365D}
\definecolor{teal}{HTML}{147D78}
\providecommand{\Sgm}{\SMG}
\renewcommand{\Sgm}{\SMG}
\setcounter{theorem}{0}
\setcounter{proposition}{0}
\setcounter{lemma}{0}
\setcounter{corollary}{0}
\setcounter{remark}{0}
\begin{abstract}
We verify the proposed spectral formulas on all matrices of the two
extraspecial groups of order $3^3$ and on the faithful representations of
$D_8$ and $Q_8$.  These examples are minimal but discriminating: they detect
both the exponent-$p$ versus exponent-$p^2$ split and the plus/minus Arf type.
All 20 coefficient comparisons pass.
\end{abstract}

\subsection{Odd $p$: what the character does not see}
Let $E$ be extraspecial of order $p^{1+2n}$ and let $\rho$ be a faithful
irreducible representation of dimension $N=p^n$, with
$\rho(z)=\omega I_N$.  Central matrices contribute
\[
p^{-2n}\Pi_p\prod_j(1-t_j)^{-N}.
\]
In the exponent-$p$ case every noncentral element has all $p$th roots once
per block of size $p$, hence
\begin{equation}\label{extraspecial-groups-proof-and-verification:eq:ep}
\Sgm_{E,\exp p}=p^{-2n}\Pi_p\prod_j(1-t_j)^{-N}
+(1-p^{-2n})\prod_j(1-t_j^p)^{-N/p}.
\end{equation}

In the exponent-$p^2$ case, $g^p=z^{\ell(gZ)}$ for a nonzero linear
functional $\ell:E/Z(E)\to\mathbb F_p$.  Set
$A_a(\mathbf t)=\prod_j(1-\omega^at_j^p)^{-N/p}$.  Counting fibers of
$\ell$ gives
\begin{equation}\label{extraspecial-groups-proof-and-verification:eq:ep2}
\Sgm_{E,\exp p^2}=p^{-2n}\Pi_p\prod_j(1-t_j)^{-N}
+(p^{-1}-p^{-2n})A_0+p^{-1}\sum_{a=1}^{p-1}A_a.
\end{equation}
The coefficients in \eqref{extraspecial-groups-proof-and-verification:eq:ep2} sum to one.  The last sum is the spectral
information missed by the ordinary character, which vanishes off the center.

For the direct exponent-$9$ test we use
$E=C_9\rtimes C_3=\langle x,y:y^{-1}xy=x^4\rangle$ and the matrices
\[
\rho(x)=\operatorname{diag}(\eta,\eta^4,\eta^7),\qquad
\rho(y)e_j=e_{j-1},\qquad \eta=e^{2\pi i/9}.
\]
The exponent-$3$ comparison is the $3$-dimensional Heisenberg representation.

\subsection{Extraspecial 2-groups}
For sign $\varepsilon=+1$ (plus type) or $-1$ (minus type), put $N=2^n$ and
\[
A_+=\prod_j(1-t_j^2)^{-N/2},\qquad
A_-=\prod_j(1+t_j^2)^{-N/2}.
\]
The quadratic form on $E/Z(E)$ has
$2^{2n-1}+\varepsilon2^{n-1}$ zeros and
$2^{2n-1}-\varepsilon2^{n-1}$ ones.  Removing the identity coset from the
first sector, restoring its two central lifts, and dividing by $|E|$ proves
\begin{align}\label{extraspecial-groups-proof-and-verification:eq:e2}
\Sgm_{E^\varepsilon}={}&2^{-2n}\Pi_2\prod_j(1-t_j)^{-N}\notag\\
&+(\tfrac12+\varepsilon2^{-n-1}-2^{-2n})A_+
+(\tfrac12-\varepsilon2^{-n-1})A_-.
\end{align}
For $n=1$ the plus and minus representations are $D_8$ and $Q_8$.

\subsection{Direct coefficient test}
The verifier constructs every representation matrix, computes
$h_d(\operatorname{spec}g)$ from $\det(I-tg)^{-1}$, and averages
$\prod_i h_{\tau_i}$ over the group.  It separately expands
\eqref{extraspecial-groups-proof-and-verification:eq:ep}, \eqref{extraspecial-groups-proof-and-verification:eq:ep2}, or \eqref{extraspecial-groups-proof-and-verification:eq:e2}.  Selected outcomes are:
\begin{center}
\begin{tabular}{@{}lrrr@{}}\toprule
group & $\tau$ & direct & formula\\\midrule
exponent $3$ & $(3)$&2&2\\
exponent $9$ & $(3)$&1&1\\
exponent $3$ & $(3,3)$&12&12\\
exponent $9$ & $(3,3)$&11&11\\
$D_8$ & $(2)$&1&1\\
$Q_8$ & $(2)$&0&0\\
$D_8$ & $(1,1)$&1&1\\
$Q_8$ & $(1,1)$&1&1\\
$D_8$ & $(2,2)$&3&3\\
$Q_8$ & $(2,2)$&3&3\\\bottomrule
\end{tabular}
\end{center}
The unequal rows demonstrate that the tests would fail if the spectral
sectors were collapsed.  All 20 built-in cases pass to $2\times10^{-8}$.

\subsection{Reproduction}
\begin{verbatim}
.venv/bin/python for_this_guy/quantum_matrix_groups/
  extraspecial_groups/verification.py
.venv/bin/marimo edit marimo_notebooks/extraspecial_groups.py
\end{verbatim}
Join each displayed command onto one shell line.
\endgroup


\clearpage
\section{Clifford groups and balanced design invariants}\label{proof:clifford}
\begingroup
\definecolor{navy}{HTML}{17365D}
\definecolor{teal}{HTML}{147D78}
\definecolor{soft}{HTML}{F2F6F8}
\providecommand{\Sgm}{\SMG}
\renewcommand{\Sgm}{\SMG}
\providecommand{\Bmgf}{\mathcal B}
\renewcommand{\Bmgf}{\mathcal B}
\setcounter{theorem}{0}
\setcounter{proposition}{0}
\setcounter{lemma}{0}
\setcounter{corollary}{0}
\setcounter{remark}{0}
\begin{abstract}
The ordinary one-sided sigma--MGF does not detect unitary-design order.  We
prove the correct two-sided coefficient statement, isolate a further phase
subtlety, and verify the design-detecting coefficient by closing explicit
projective Clifford matrix groups in dimensions $2$ and $4$.  The computed
frame potentials agree with Haar through degree three and differ at degree
four: $(1,2,5,15)$ versus $(1,2,5,14)$ in dimension two, and
$(1,2,6,29)$ versus $(1,2,6,24)$ in dimension four.
\end{abstract}

\subsection{One-sided series and the scalar obstruction}
For a finite lift $\widetilde G\leq U(V)$, put
\[
\Sgm_{\widetilde G,V}(\mathbf t)=\frac1{|\widetilde G|}\sum_{g\in\widetilde G}
 \prod_{i\geq1}\det(I-t_i g)^{-1}.
\]
For a compact group $K\leq U(V)$, define instead
\[
\Sgm_{K,V}(\mathbf t)=\int_K\prod_{i\geq1}\det(I-t_i k)^{-1}\,dk
\]
with normalized Haar measure.
Expanding each determinant gives
\[
[m_\tau]\Sgm_{\widetilde G,V}=\dim\left(\bigotimes_j
\operatorname{Sym}^{\tau_j}V\right)^{\widetilde G}.
\]
If the compact group contains the full scalar circle, every positive-degree
one-sided coefficient vanishes.  Hence
\[
\boxed{\Sgm_{K,V}=1.}
\]
for $U(2^n)$ and for the full normalizer $N_{U(2^n)}(P_n)$.  A finite scalar
lift instead obeys only the congruence imposed by its finite scalar center and
may have nonzero lift-dependent coefficients.

\subsection{The two-sided series, correctly phase neutralized}
For a specified finite lift $\widetilde G\leq U(V)$ define
\[
\Bmgf_{\widetilde G}(\mathbf s,\mathbf t)
=\frac1{|\widetilde G|}\sum_{g\in\widetilde G}
 \prod_i\det(I-s_i g)^{-1}\prod_j\det(I-t_j\bar g)^{-1}.
\]
Character expansion and averaging prove
\begin{equation}\label{clifford-groups-proof-and-verification:eq:two-sided}
[m_\alpha(\mathbf s)m_\beta(\mathbf t)]\Bmgf_{\widetilde G}
=\dim\left(\operatorname{Sym}^{\alpha}V\otimes
\operatorname{Sym}^{\beta}V^*\right)^{\widetilde G}.
\end{equation}

There is a useful refinement.  Replacing a representative $g$ by $zg$
multiplies the coefficient integrand by $z^{|\alpha|-|\beta|}$.  Thus the
entire two-sided series is not intrinsically projective.  Its
\emph{phase-neutral projection}
\begin{equation}\label{clifford-groups-proof-and-verification:eq:neutral}
\Bmgf_G^0=\sum_{|\alpha|=|\beta|}
[m_\alpha m_\beta]\Bmgf_{\widetilde G}\,m_\alpha m_\beta
\end{equation}
is independent of representatives.  Equivalently it is the constant term in
$z$ after $(\mathbf s,\mathbf t)\mapsto(z\mathbf s,z^{-1}\mathbf t)$.
For the full normalizer, continuous scalar integration performs exactly this
projection.  All design coefficients used below lie in \eqref{clifford-groups-proof-and-verification:eq:neutral}.

\subsection{The design-detecting coefficient}
For $(1^k)=(1,\ldots,1)$, one has
$\operatorname{Sym}^{(1^k)}V=V^{\otimes k}$.  Hence \eqref{clifford-groups-proof-and-verification:eq:two-sided}
gives
\begin{align}
[m_{(1^k)}(\mathbf s)m_{(1^k)}(\mathbf t)]\Bmgf_G^0
&=\dim\operatorname{End}_G(V^{\otimes k})\\
&=\frac1{|G|}\sum_{g\in G}|\operatorname{Tr}g|^{2k}.
\label{clifford-groups-proof-and-verification:eq:frame}
\end{align}
The last quantity is the $k$th unitary frame potential.

Let $P_G^{(k)}$ be the average of
$g^{\otimes k}\otimes\bar g^{\otimes k}$.  It is the orthogonal projection
onto the corresponding $G$-fixed space.  Since $G\subseteq U(d)$, the Haar
fixed space is contained in the $G$-fixed space.  Equality of the coefficient
in \eqref{clifford-groups-proof-and-verification:eq:frame} with the Haar coefficient is equality of the two ranks;
the nested spaces, and therefore their orthogonal projections, are equal.
This proves
\[
\boxed{G\text{ is a unitary }k\text{-design}\iff
\Phi_k(G)=\Phi_k(U(d)).}
\]
Schur--Weyl duality supplies the exact Haar reference
\[
\Phi_k(U(d))=\sum_{\substack{\lambda\vdash k\\\ell(\lambda)\leq d}}
(f^\lambda)^2,
\]
which reduces to $k!$ when $d\geq k$.

\subsection{Independent matrix verification}
The program constructs
\[
H=2^{-1/2}\begin{pmatrix}1&1\\1&-1\end{pmatrix},\qquad
S=\begin{pmatrix}1&0\\0&i\end{pmatrix}
\]
and, for two qubits, their tensor placements together with CNOT.  Matrices
are multiplied until closure, with equality taken modulo global phase.  The
resulting orders are asserted to be $24$ and $11{,}520$.  No class table or
design theorem is used in the numerical average.  The Haar comparator is
computed independently by the hook-length formula in the Schur--Weyl sum.

\begin{center}
\begin{tabular}{@{}ccrrrr@{}}\toprule
$n$&$d$&$|\overline{\mathrm{Cliff}}_n|$&$k$&direct $\Phi_k$&Haar\\\midrule
1&2&24&1&1&1\\
 &&&2&2&2\\
 &&&3&5&5\\
 &&&4&15&14\\\addlinespace
2&4&11,520&1&1&1\\
 &&&2&2&2\\
 &&&3&6&6\\
 &&&4&29&24\\\bottomrule
\end{tabular}
\end{center}
Thus the first total degree at which the canonical phase-neutral coefficient
can disagree is four, and the frame coefficient does disagree there.
This is precisely the Clifford $3$-design/$4$-design transition established
in \cite{clifford-groups-proof-and-verification:zhu3design,clifford-groups-proof-and-verification:graceful}.
The notebook establishes the displayed frame-potential values only for the
one- and two-qubit projective groups through $k=4$.  The all-qubit 3-design
and failure-of-4-design statements are external theorems supplied by those
references.

As a separate regression test, closing the ordinary one-qubit lift generated
by $H,S$ gives $192$ matrices and scalar center $\mu_8I$.  Its one-sided
coefficients for $\tau=(1),(2),(4)$ vanish, while those for $(8)$ and $(4,4)$
are $1$ and $2$.  This confirms both lift dependence and the scalar selection
rule $8\mid|\tau|$.

\subsection{Reproduction and scope}
Run
\begin{verbatim}
.venv/bin/python for_this_guy/quantum_matrix_groups/
  clifford_groups/verification.py
.venv/bin/marimo edit marimo_notebooks/clifford_groups.py
\end{verbatim}
after joining each wrapped command onto one line.  The sweep verifies the
multiqubit statement for $n=1,2$ and $k\leq4$; the cited theorems supply the
all-$n$ result.  It makes no inference about odd-prime qudit Clifford groups.

\begin{thebibliography}{9}
\bibitem{clifford-groups-proof-and-verification:zhu3design} H.~Zhu, ``Multiqubit Clifford groups are unitary
3-designs,'' \url{https://arxiv.org/abs/1510.02619}.
\bibitem{clifford-groups-proof-and-verification:graceful} H.~Zhu, R.~Kueng, M.~Grassl, and D.~Gross,
``The Clifford group fails gracefully to be a unitary 4-design,''
\url{https://arxiv.org/abs/1609.08172}.
\end{thebibliography}
\endgroup


\clearpage
\section{The real extraspecial Pauli layer}\label{proof:real-pauli}
\begingroup
\definecolor{navy}{HTML}{17365D}
\definecolor{teal}{HTML}{147D78}
\providecommand{\MGF}{\SMG}
\renewcommand{\MGF}{\SMG}
\setcounter{theorem}{0}
\setcounter{proposition}{0}
\setcounter{lemma}{0}
\setcounter{corollary}{0}
\setcounter{remark}{0}
\begin{abstract}
For the plus-type extraspecial real Pauli group
$E_m\cong 2_+^{1+2m}$ on $V_m=\mathbb C^{2^m}$, we derive the complete
four-sector determinant average.  An independent program constructs every
matrix $\pm X^aZ^b$ and checks 24 one-row and mixed coefficients in
dimensions $2,4,8$.  All comparisons are exact integers and all pass.
\end{abstract}

\subsection{Representation and coefficient target}
Put $N=2^m$.  On the basis $\{e_x:x\in\mathbb F_2^m\}$ define
\[
 X(a)e_x=e_{x+a},\qquad Z(b)e_x=(-1)^{b\cdot x}e_x.
\]
Then
\[
 E_m=\{\pm X(a)Z(b):a,b\in\mathbb F_2^m\},\qquad |E_m|=2N^2.
\]
For a partition $\tau=(\tau_1,\ldots,\tau_r)$, determinant expansion gives
\begin{equation}\label{extraspecial-real-pauli-proof-and-verification:eq:target}
 [m_\tau]\MGF_{E_m,m}
 =\frac1{|E_m|}\sum_{g\in E_m}\prod_i
 h_{\tau_i}(\operatorname{spec}g),
\end{equation}
where $\sum_{d\geq0}h_d(\operatorname{spec}g)t^d=
\det(I-tg)^{-1}$.  This is exactly
$\dim(\bigotimes_i\operatorname{Sym}^{\tau_i}V_m)^{E_m}$.

\subsection{The four spectral sectors}
The central elements $I,-I$ have determinants $(1-t)^N,(1+t)^N$.
For $g=\pm X(a)Z(b)$ one has
\[
 g^2=(-1)^{a\cdot b}I.
\]
Off the center, the trace is zero.  Hence an involution has $N/2$ eigenvalues
of each sign, while an element squaring to $-I$ has $N/2$ eigenvalues each
of $i$ and $-i$.  Their determinants are respectively
$(1-t^2)^{N/2}$ and $(1+t^2)^{N/2}$.

The quadratic form $q(a,b)=a\cdot b$ is plus type.  It has
$(N^2+N)/2$ zeros, including $(0,0)$.  Accounting for both central lifts
therefore gives
\[
 \#\{g\notin Z(E_m):g^2=I\}=N^2+N-2,
 \qquad \#\{g:g^2=-I\}=N^2-N.
\]
The counts sum with the two central elements to $2N^2$.  Averaging proves
\begin{equation}\label{extraspecial-real-pauli-proof-and-verification:eq:formula}
\boxed{\begin{aligned}
\MGF_{E_m,m}(\mathbf t)=\frac1{2N^2}\Big[&
\prod_j(1-t_j)^{-N}+\prod_j(1+t_j)^{-N}\\
&+(N^2+N-2)\prod_j(1-t_j^2)^{-N/2}\\
&+(N^2-N)\prod_j(1+t_j^2)^{-N/2}\Big].
\end{aligned}}
\end{equation}
Because $-I\in E_m$, the first two terms cancel in odd total degree and the
other two contain only even powers.  Thus every odd-total-degree coefficient
vanishes.

\subsection{Independent coefficient extraction}
The verifier does not classify matrices after construction.  For each
explicit matrix it computes the eigenvalues, obtains $h_d$ recursively from
the product generating function, forms the integrand in \eqref{extraspecial-real-pauli-proof-and-verification:eq:target},
and averages.  The formula side separately extracts coefficients from the
four factors in \eqref{extraspecial-real-pauli-proof-and-verification:eq:formula} using binomial coefficients.
Before coefficient extraction, every matrix is independently checked to have
real entries and to satisfy $g^{\mathsf T}g=I$.

\begin{center}
\begin{tabular}{@{}c@{\quad}c@{\quad}c@{\quad}c@{\qquad}r@{\qquad}r@{}}\toprule
$m$&$N$&$|E_m|$&$\tau$&direct&formula\\\midrule
1&2&8&$(3,1)$&2&2\\
1&2&8&$(2,2)$&3&3\\
2&4&32&$(4)$&5&5\\
2&4&32&$(4,2)$&23&23\\
3&8&128&$(3,1)$&15&15\\
3&8&128&$(6)$&29&29\\
3&8&128&$(4,2)$&190&190\\\bottomrule
\end{tabular}
\end{center}
The full sweep uses eight partitions at each rank.  It includes the odd
case $(1)$ and the mixed cases $(1,1)$, $(3,1)$, $(2,2)$, and $(4,2)$.
All 24 pass.

\subsection{Reproduction}
\begin{verbatim}
.venv/bin/python for_this_guy/barnes_wall_lattice_automorphism_groups/
  extraspecial_real_pauli/verification.py
.venv/bin/marimo edit marimo_notebooks/
  extraspecial_real_pauli.py
\end{verbatim}
Join each wrapped command onto one shell line.  The computation is exhaustive
at the stated ranks; no random sampling or precomputed character table is
used.
\endgroup


\clearpage
\section{The real Clifford layer}\label{proof:real-clifford}
\begingroup
\definecolor{navy}{HTML}{17365D}
\definecolor{teal}{HTML}{147D78}
\providecommand{\MGF}{\SMG}
\renewcommand{\MGF}{\SMG}
\setcounter{theorem}{0}
\setcounter{proposition}{0}
\setcounter{lemma}{0}
\setcounter{corollary}{0}
\setcounter{remark}{0}
\begin{abstract}
We prove the exact reduction of the real Clifford sigma--MGF to coset
averages over $O^+(2m,2)$ and state precisely what the self-dual-code theorem
does and does not imply for mixed coefficients.  Exhaustive matrix closures
in dimensions $2$ and $4$ recover group orders $16$ and $2304$, quotient
orders $2$ and $72$, lift independence, and 16 direct coefficient matches.
\end{abstract}

\subsection{Exact quotient formula}
Let $E_m\triangleleft\mathcal C_m=N_{O(V_m)}(E_m)$, with
\[
 \mathcal C_m/E_m\cong O^+(2m,2).
\]
For $A$ in the quotient and any lift $\widetilde A$, define
\[
 \Phi_m(A;\mathbf t)=\frac1{|E_m|}\sum_{e\in E_m}
 \prod_j\det(I-t_je\widetilde A)^{-1}.
\]
Replacing $\widetilde A$ by $e_0\widetilde A$ permutes $E_m$, so $\Phi_m$
is well defined.  Partitioning $\mathcal C_m$ into right cosets proves
\begin{equation}\label{real-clifford-proof-and-verification:eq:coset}
\boxed{\MGF_{\mathcal C_m,m}(\mathbf t)=
\frac1{|O^+(2m,2)|}\sum_{A\in O^+(2m,2)}\Phi_m(A;\mathbf t).}
\end{equation}
Coefficient extraction gives
\[
[m_\tau]\MGF_{\mathcal C_m,m}=
\dim\left(\bigotimes_i\operatorname{Sym}^{\tau_i}V_m\right)^{\mathcal C_m}.
\]

\subsection{One-row code theorem and its boundary}
For a binary self-dual code $C\leq\mathbb F_2^{2k}$, its genus-$m$
complete weight enumerator is
\[
 \operatorname{cwe}_m(C)=\sum_{(c^{(1)},\ldots,c^{(m)})\in C^m}
 \prod_{j=1}^{2k}x_{(c_j^{(1)},\ldots,c_j^{(m)})}.
\]
The Clifford invariant theorem identifies the degree-$2k$ invariant space
with the span of these enumerators.  In the stable range $m\geq k-1$, they
are independent and therefore
\begin{equation}\label{real-clifford-proof-and-verification:eq:codes}
 [m_{(2k)}]\MGF_{\mathcal C_m,m}
 =\#\{\text{binary self-dual codes of length }2k\}/\!\cong.
\end{equation}
Consequently the stable values in degrees $2,4,6,8$ are $1,1,1,2$.
The identification with genus-$m$ complete weight enumerators and their
independence in this range are external inputs from the
Nebe--Rains--Sloane invariant theorem cited below.  This note proves the coset
formula and verifies low-rank examples; it does not reprove that theorem.
The degree-eight value is guaranteed by \eqref{real-clifford-proof-and-verification:eq:codes} only for $m\geq3$;
our smaller closures also happen to give two, but we report that as a direct
low-rank computation rather than misuse the stable-range hypothesis.

For a general partition $\tau$, block-symmetrizing the full code tensor
produces explicit invariants.  The ordinary invariant-ring theorem does not
by itself show that these span every mixed fixed space.  Thus \eqref{real-clifford-proof-and-verification:eq:coset}
is exact for all $\tau$, while a code-indexed equality for all mixed $\tau$
would require a separate multigraded theorem.

\subsection{Explicit matrices and exhaustive checks}
We generate the real Pauli matrices together with tensor placements of
\[
 H=2^{-1/2}\begin{pmatrix}1&1\\1&-1\end{pmatrix}
\]
and, for $m=2$, CNOT.  Breadth-first closure yields
\[
|\mathcal C_1|=16=8\cdot2,
\qquad |\mathcal C_2|=2304=32\cdot72.
\]
Every closed matrix is checked to be real orthogonal.  Conjugation by every
selected coset representative sends the Pauli group to itself.  The program
then partitions the closed group into Pauli cosets, averages each coset, and
compares \eqref{real-clifford-proof-and-verification:eq:coset} with an independent direct average.  Each lift is
also multiplied by a varying Pauli matrix, and every corresponding pair of
coset contributions is compared before the quotient average is taken.

\begin{center}
\begin{tabular}{@{}cccrrr@{}}\toprule
$m$&$|\mathcal C_m|$&$|O^+|$&$\tau$&direct&coset\\\midrule
1&16&2&$(2)$&1&1\\
1&16&2&$(8)$&2&2\\
1&16&2&$(2,2)$&2&2\\
2&2304&72&$(4)$&1&1\\
2&2304&72&$(8)$&2&2\\
2&2304&72&$(4,2)$&2&2\\\bottomrule
\end{tabular}
\end{center}
The full sweep tests eight partitions at each rank.  Stable-code references
are applied only when $m\geq k-1$; all 16 rows pass.

\subsection{Reproduction and references}
\begin{verbatim}
.venv/bin/python for_this_guy/barnes_wall_lattice_automorphism_groups/
  real_clifford/verification.py
.venv/bin/marimo edit marimo_notebooks/
  real_clifford.py
\end{verbatim}
Join wrapped commands onto one line.  For the invariant theorem, see
G.~Nebe, E.~Rains, and N.~J.~A.~Sloane,
\emph{The invariants of the Clifford groups},
\url{https://arxiv.org/abs/math/0001038}.
\endgroup


\clearpage
\section{Barnes--Wall lattice automorphism groups}\label{proof:barnes-wall}
\begingroup
\definecolor{navy}{HTML}{17365D}
\definecolor{teal}{HTML}{147D78}
\providecommand{\MGF}{\SMG}
\renewcommand{\MGF}{\SMG}
\setcounter{theorem}{0}
\setcounter{proposition}{0}
\setcounter{lemma}{0}
\setcounter{corollary}{0}
\setcounter{remark}{0}
\begin{abstract}
Away from the exceptional rank $m=3$, the Barnes--Wall automorphism group is
an index-two subgroup of the real Clifford group.  We prove the resulting
twisted sigma--MGF identity and verify it exhaustively in the concrete models
$BW_2=\mathbb Z^2$ and $BW_4=D_4$.  The constructed groups have orders $8$
and $1152$ and agree exactly with the Clifford matrices preserving the
specified lattice bases.  All 16 coefficient comparisons pass.
\end{abstract}

\subsection{The index-two identity}
Assume $m\ne3$ and write
\[
B_m=\operatorname{Aut}(BW_{2^m})=\ker\varepsilon_m
\quad\text{inside}\quad \mathcal C_m,
\]
where $\varepsilon_m:\mathcal C_m\to\{\pm1\}$ is the nontrivial quotient
character.  Define the twisted series
\[
\MGF^{\varepsilon_m}_{\mathcal C_m,m}(\mathbf t)
=\frac1{|\mathcal C_m|}\sum_{g\in\mathcal C_m}\varepsilon_m(g)
\prod_j\det(I-t_jg)^{-1}.
\]
Since $\mathbf1_{B_m}(g)=(1+\varepsilon_m(g))/2$ and
$|B_m|=|\mathcal C_m|/2$, direct substitution proves
\begin{equation}\label{barnes-wall-proof-and-verification:eq:twist}
\boxed{\MGF_{B_m,m}=\MGF_{\mathcal C_m,m}
+\MGF^{\varepsilon_m}_{\mathcal C_m,m}.}
\end{equation}
For $W_\tau=\bigotimes_i\operatorname{Sym}^{\tau_i}V_m$, coefficient
extraction equivalently gives
\[
\dim W_\tau^{B_m}=\dim W_\tau^{\mathcal C_m}
+\dim\operatorname{Hom}_{\mathcal C_m}(\varepsilon_m,W_\tau).
\]
The second nonnegative summand is precisely the lattice-sensitive sector.

The structural equality
$\operatorname{Aut}(BW_{2^m})=\ker\varepsilon_m\leq\mathcal C_m$ and the
exceptional behavior at $m=3$ are external inputs from the cited
Barnes--Wall/Clifford-group literature.  Equation \eqref{barnes-wall-proof-and-verification:eq:twist} is proved
conditionally from this index-two statement, and the verifier checks the
concrete ranks $m=1,2$.

\subsection{Concrete lattice models}
For $m=1$ we use $BW_2=\mathbb Z^2$.  Its automorphism group is the signed
permutation group of order $8$, equal to $E_1$, while the real Clifford group
has order $16$.

For $m=2$ we use
\[
 BW_4=D_4=\{x\in\mathbb Z^4:\textstyle\sum_i x_i\equiv0\pmod2\}.
\]
A basis is given by the columns of
\[
L=\begin{pmatrix}
1&0&0&0\\-1&1&0&0\\0&-1&1&1\\0&0&-1&1
\end{pmatrix}.
\]
Starting with $E_2$, tensor Hadamard $H\otimes H$, CNOT, and controlled-$Z$,
matrix closure gives a subgroup of order $1152$.  Independently, for each of
the $2304$ matrices $g\in\mathcal C_2$ the verifier tests whether
$L^{-1}gL\in GL(4,\mathbb Z)$.  Exactly the same $1152$ matrices pass.  Thus
the subgroup identification uses lattice preservation, not only an expected
order.  Every matrix in both closed groups is also checked to have real
entries and satisfy $g^{\mathsf T}g=I$.

\subsection{Direct versus twisted averages}
For each partition $\tau$, the direct side averages
$\prod_i h_{\tau_i}(\operatorname{spec}g)$ over $B_m$.  On the right side,
the same integrand is averaged first with weight $1$ and then with weight
$\varepsilon_m(g)$ over all of $\mathcal C_m$.  Membership in the independently
identified lattice stabilizer determines the sign.

\begin{center}
\begin{tabular}{@{}cccrrrrr@{}}\toprule
$m$&$|B_m|$&$\tau$&direct&Clifford&twisted&sum\\\midrule
1&8&$(4)$&2&1&1&2\\
1&8&$(8)$&3&2&1&3\\
1&8&$(4,2)$&4&2&2&4\\
2&1152&$(4)$&1&1&0&1\\
2&1152&$(6)$&2&1&1&2\\
2&1152&$(8)$&3&2&1&3\\
2&1152&$(4,2)$&3&2&1&3\\\bottomrule
\end{tabular}
\end{center}
These rows exhibit a real distinction: the first tested one-row twisted
contribution occurs in degree $4$ for $B_1$ but only degree $6$ for $B_2$.
The full sweep tests eight partitions at each rank, including mixed powers;
all 16 instances of \eqref{barnes-wall-proof-and-verification:eq:twist} pass.

\subsection{Scope and reproduction}
No $m=3$ conclusion is drawn: that rank is explicitly exceptional to the
index-two statement.  Ranks $m\geq4$ are covered algebraically by
\eqref{barnes-wall-proof-and-verification:eq:twist}, but exhaustive closure is not a reasonable verification
strategy because the groups grow rapidly.
\begin{verbatim}
.venv/bin/python for_this_guy/barnes_wall_lattice_automorphism_groups/
  barnes_wall/verification.py
.venv/bin/marimo edit marimo_notebooks/barnes_wall.py
\end{verbatim}
Join each wrapped command onto one line.  Background on the group chain and
the exceptional rank appears in G.~Nebe, E.~Rains, and N.~J.~A.~Sloane,
\emph{The invariants of the Clifford groups},
\url{https://arxiv.org/abs/math/0001038}.
\endgroup


\clearpage
\section{Lattice automorphism groups}\label{proof:new-dist:lattices}
\proofstatus{Conditional on certified class data}{The universal class/power or
frame-shape formula is exact. Concrete numerical tables in this section depend
on external class sizes, power maps, character values, or frame shapes and are
not independently derived by the Molien argument.}


\resultbox{\textbf{Canonical testing artifacts.}\par\smallskip Exact backend: \path{lambda_distributions/dists/lattice_sigma_mgfs/verification.py}.\par Regression: \path{lambda_distributions/dists/lattice_sigma_mgfs/test_verification.py}.\par Marimo: \path{marimo_notebooks/lattices.py}.}

\subsection*{Scope and outcome}

For the canonical complex representation of the full automorphism group of a
positive-definite integral lattice, we prove the exact frame-shape
$\sigma$-MGF and its coefficient interpretation by Reynolds projection.  We
then prove the cycle formulas and coefficientwise limits for type-$A$ root
lattices, signed-permutation lattices, and repeated orthogonal sums; give the
finite class reductions for Niemeier, Coxeter--Todd, extremal, and modular
lattices; exhibit explicit rank-$2,3,4$ lattices with automorphism group
exactly $\{\pm I\}$; and prove the contrasting exponential moment growth for
tensor powers.  A reproducible marimo notebook constructs all tractable
matrix groups and compares direct matrix averages, recovered frame shapes,
family cycle formulas, and independently computed common fixed spaces.

\subsection{Canonical representation and the frame-shape theorem}

Let $L$ be a positive-definite integral lattice of rank $d$, put
\[
 G_L=\Aut(L),\qquad V_L=L\otimes_{\Z}\C,
\]
and let $G_L$ act through its canonical lattice matrices.  Positivity makes
$G_L$ finite.  If $\tau=(\tau_1,\ldots,\tau_s)$ is a partition, define
\[
 W_\tau(V_L)=\bigotimes_{j=1}^s\Sym^{\tau_j}V_L,
 \qquad a_\tau(L)=\dim W_\tau(V_L)^{G_L}.
\]
The symmetric series $\MGF_L=\sum_\tau a_\tau(L)m_\tau$ is the lattice
$\sigma$-MGF.

Every $g\in G_L$ has finite order.  Its characteristic polynomial can be
written uniquely in frame-shape form
\begin{equation}\label{nd:lattices:eq:frame}
 \det(xI-g)=\prod_{m\ge1}(x^m-1)^{a_m(g)},
 \qquad \sum_m m a_m(g)=d.
\end{equation}
The exponents are integers and may be negative; cyclotomic cancellation makes
the product a polynomial.  Equivalently,
\begin{equation}\label{nd:lattices:eq:frame-det}
 \det(I-tg)=\prod_{m\ge1}(1-t^m)^{a_m(g)}.
\end{equation}

\begin{theorem}[frame-shape $\sigma$-MGF]\label{nd:lattices:thm:frame}
For every positive-definite integral lattice,
\begin{align}
 \MGF_L(\mathbf t)
 &=\frac1{|G_L|}\sum_{g\in G_L}
   \prod_{i\ge1}\prod_{m\ge1}(1-t_i^m)^{-a_m(g)},
   \label{nd:lattices:eq:element-frame}\\
 &=\sum_{[g]\subseteq G_L}\frac1{|C_{G_L}(g)|}
   \prod_{i,m}(1-t_i^m)^{-a_m(g)},\label{nd:lattices:eq:class-frame}\\
 &=\sum_{[g]}\frac1{|C_{G_L}(g)|}
   \exp\!\left(\sum_{r\ge1}\frac{p_r}{r}\Tr(g^r\mid V_L)\right).
   \label{nd:lattices:eq:trace-frame}
\end{align}
Moreover,
\begin{equation}\label{nd:lattices:eq:coefficient}
 \boxed{[m_\tau]\MGF_L=\dim W_\tau(V_L)^{G_L}.}
\end{equation}
\end{theorem}

\begin{proof}
For any complex matrix $A$,
\begin{equation}\label{nd:lattices:eq:complete}
 \sum_{q\ge0}\Tr(\Sym^qA)t^q
 =\det(I-tA)^{-1}
 =\exp\!\left(\sum_{r\ge1}\frac{t^r}{r}\Tr(A^r)\right).
\end{equation}
Multiply over the variables $t_i$ and average over $G_L$.  The coefficient of
$t_1^{\tau_1}\cdots t_s^{\tau_s}$ is the trace on $W_\tau(V_L)$ of the
Reynolds idempotent
\[
 R_\tau=\frac1{|G_L|}\sum_{g\in G_L}\rho_{W_\tau}(g).
\]
Its image is exactly $W_\tau(V_L)^{G_L}$, so its trace proves
\eqref{nd:lattices:eq:coefficient}.  Equation~\eqref{nd:lattices:eq:frame-det} gives
\eqref{nd:lattices:eq:element-frame}; grouping elements into conjugacy classes gives
\eqref{nd:lattices:eq:class-frame}; and the exponential identity in
\eqref{nd:lattices:eq:complete} gives \eqref{nd:lattices:eq:trace-frame}.
\end{proof}

The frame exponents and power traces determine one another:
\begin{equation}\label{nd:lattices:eq:trace-frame-inversion}
 \Tr(g^r)=\sum_{m\mid r}m a_m(g),\qquad
 a_m(g)=\frac1m\sum_{e\mid m}\mu(m/e)\Tr(g^e).
\end{equation}
Indeed, the sum of the $r$th powers of all $m$th roots of unity is $m$ if
$m\mid r$ and zero otherwise.  Möbius inversion proves the second formula.
This is also the exact recovery algorithm used by the verifier.

\subsection{Universal parity and low-degree diagnostics}

The central automorphism $-I$ belongs to every $G_L$.  It acts on
$W_\tau(V_L)$ by $(-1)^{|\tau|}$; therefore
\begin{equation}\label{nd:lattices:eq:parity}
 \boxed{[m_\tau]\MGF_L=0\quad\text{if }|\tau|\text{ is odd}.}
\end{equation}
Because the lattice form identifies $V_L$ with its dual,
\begin{equation}\label{nd:lattices:eq:low-degree}
 [m_{(1,1)}]\MGF_L=\dim\End_{G_L}(V_L),
 \qquad [m_{(2)}]\MGF_L\ge1.
\end{equation}
If $V_L$ is complex irreducible, Schur's lemma and uniqueness of the invariant
symmetric form give equality to one in both positions.

The fourth multilinear coefficient is already more sensitive.  If
$V_L\otimes V_L=\bigoplus_\rho b_\rho\rho$, then
\begin{equation}\label{nd:lattices:eq:fourth}
 [m_{(1,1,1,1)}]\MGF_L
 =\dim\End_{G_L}(V_L\otimes V_L)=\sum_\rho b_\rho^2.
\end{equation}
Thus irreducibility in degree one does not control a growing family in degree
four.

More generally, Newton's identity gives the exact moment formula
\begin{equation}\label{nd:lattices:eq:moment}
 [m_\tau]\MGF_L
 =\E_{g\in G_L}\prod_j\left[
 \sum_{\lambda\vdash\tau_j}\frac1{z_\lambda}
 \prod_{r\ge1}\Tr(g^r)^{m_r(\lambda)}\right].
\end{equation}
A coefficient of total degree $D$ depends only on the joint moments of
$\Tr(g),\ldots,\Tr(g^D)$.  This observation is the correct bridge from exact
finite formulas to coefficientwise asymptotics.

\subsection{Type-A root lattices}

For
\[
 A_{n-1}=\{(x_1,\ldots,x_n)\in\Z^n:\textstyle\sum_i x_i=0\},
 \qquad n\ge3,
\]
the full automorphism group is $\{\pm1\}\times S_n$ on the standard
hyperplane representation.  If $\sigma$ has $c_m(\sigma)$ cycles of length
$m$, deleting the constant eigenline from the permutation representation
gives
\begin{equation}\label{nd:lattices:eq:a-det}
 \det(I-t\varepsilon\sigma\mid V_{A_{n-1}})
 =\frac{\prod_m(1-(\varepsilon t)^m)^{c_m(\sigma)}}{1-\varepsilon t}.
\end{equation}
Let $m_m(\lambda)$ denote the multiplicity of $m$ in a partition $\lambda$,
and $z_\lambda=\prod_m m^{m_m(\lambda)}m_m(\lambda)!$.  Averaging
\eqref{nd:lattices:eq:a-det} by cycle type proves
\begin{equation}\label{nd:lattices:eq:a-formula}
 \boxed{
 \MGF_{A_{n-1}}(\mathbf t)=\frac12\sum_{\varepsilon=\pm1}
 \sum_{\lambda\vdash n}\frac1{z_\lambda}
 \prod_i\left[(1-\varepsilon t_i)
 \prod_{m\ge1}\bigl(1-(\varepsilon t_i)^m\bigr)^{-m_m(\lambda)}\right].}
\end{equation}

\subsubsection*{Large-rank limit}
For a uniform permutation, the finite collection of cycle counts converges
with all fixed moments to independent $Z_m\sim\operatorname{Poisson}(1/m)$.
Applying the Poisson probability-generating function to
\eqref{nd:lattices:eq:a-formula} yields, coefficientwise,
\begin{equation}\label{nd:lattices:eq:a-limit}
 \boxed{
 \MGF_{A_\infty}=\frac12\sum_{\varepsilon=\pm1}
 \prod_i(1-\varepsilon t_i)
 \exp\!\left[\sum_{m\ge1}\frac1m\left(
 \prod_i(1-(\varepsilon t_i)^m)^{-1}-1\right)\right].}
\end{equation}
Only finitely many cycle moments enter a bounded-degree coefficient, so the
usual finite-dimensional cycle-count convergence is sufficient.

\subsection{Signed-permutation lattices: types B, C, and D}

The cubic lattice $\Z^n$ has the signed-permutation group
$C_2^n\rtimes S_n$.  The same natural matrix group is the Weyl group in types
$B_n$ and $C_n$, and it is the full automorphism group of $D_n$ for $n\ge5$.
The exceptional triality of $D_4$ must be handled by its own finite class sum.

For a permutation cycle of length $\ell$, let $\delta\in\{\pm1\}$ be the
product of its coordinate signs.  Its determinant is $1-\delta t^\ell$.
After averaging the independent sign product, put
\begin{equation}\label{nd:lattices:eq:b-ell}
 B_\ell(\mathbf t)=\frac12\left[
 \prod_i(1-t_i^\ell)^{-1}+\prod_i(1+t_i^\ell)^{-1}\right].
\end{equation}
The cycle index of $S_n$ now gives
\begin{align}
 \MGF_{B_n}(\mathbf t)
 &=\sum_{\lambda\vdash n}\frac1{z_\lambda}
   \prod_{\ell\ge1}B_\ell(\mathbf t)^{m_\ell(\lambda)},
   \label{nd:lattices:eq:b-formula}\\
 \sum_{n\ge0}u^n\MGF_{B_n}(\mathbf t)
 &=\exp\!\left(\sum_{\ell\ge1}\frac{u^\ell}{\ell}B_\ell(\mathbf t)\right).
 \label{nd:lattices:eq:b-egf}
\end{align}
Since $B_\ell=1+$ terms of positive $\mathbf t$-degree, removing the factor
$(1-u)^{-1}$ and taking coefficients proves the stable limit
\begin{equation}\label{nd:lattices:eq:b-limit}
 \boxed{\MGF_{B_\infty}=\exp\!\left(
 \sum_{\ell\ge1}\frac{B_\ell(\mathbf t)-1}{\ell}\right).}
\end{equation}

\subsection{Repeated orthogonal sums and the stable law}

Fix a lattice $L$ with $G=\Aut(L)$ and let $G\wr S_k$ act on
$V_L^{\oplus k}$ by independent lattice automorphisms and component
permutations.  Define
\[
 A_\ell(\mathbf t)=\MGF_L(t_1^\ell,t_2^\ell,\ldots).
\]

\begin{theorem}[marked compound-Poisson law]\label{nd:lattices:thm:wreath-sum}
\begin{align}
 \sum_{k\ge0}u^k\MGF_{G\wr S_k,V_L^{\oplus k}}
 &=\exp\!\left(\sum_{\ell\ge1}\frac{u^\ell}{\ell}A_\ell(\mathbf t)\right),
 \label{nd:lattices:eq:wreath-sum}\\
 \lim_{k\to\infty}\MGF_{G\wr S_k,V_L^{\oplus k}}
 &=\exp\!\left[\sum_{\ell\ge1}\frac{
 \MGF_L(t_1^\ell,t_2^\ell,\ldots)-1}{\ell}\right]
 \label{nd:lattices:eq:wreath-limit}
\end{align}
coefficientwise in $\mathbf t$.
\end{theorem}

\begin{proof}
For a top permutation cycle $C$ of length $\ell$, multiply the $G$-labels
around $C$ to obtain its cycle product $g_C$.  On the associated block,
\[
 \det(I-tx\mid C)=\det(I-t^\ell g_C\mid V_L).
\]
For uniformly independent labels, $g_C$ is uniform in $G$, independently for
distinct cycles.  The exponential cycle-index formula proves
\eqref{nd:lattices:eq:wreath-sum}.  Next write
\[
 \exp\!\left(\sum_{\ell\ge1}\frac{u^\ell A_\ell}{\ell}\right)
 =\frac1{1-u}\exp\!\left(
 \sum_{\ell\ge1}\frac{u^\ell(A_\ell-1)}{\ell}\right).
\]
In each fixed $\mathbf t$-degree, the second factor is a polynomial in $u$.
Its coefficient sequence therefore stabilizes after multiplication by
$(1-u)^{-1}$, and evaluation at $u=1$ proves \eqref{nd:lattices:eq:wreath-limit}.
\end{proof}

This theorem applies to repeated copies of any fixed root lattice, the
Coxeter--Todd lattice, or any other fixed component.  It is a statement about
the natural wreath subgroup; an orthogonal sum can have additional
automorphisms when accidental isometries occur.

\subsection{Isolated and finite lattice families}

\subsubsection*{Rootful Niemeier lattices}
Let $N$ be rootful with root sublattice $R$, Weyl group
$W(R)\triangleleft\Aut(N)$, and umbral quotient
$G^N=\Aut(N)/W(R)$.  Choose a lift $\widetilde\gamma$ of each
$\gamma\in G^N$.  The exact finite reduction is
\begin{equation}\label{nd:lattices:eq:niemeier}
 \boxed{\MGF_N=\frac1{|G^N|}\sum_{\gamma\in G^N}\frac1{|W(R)|}
 \sum_{w\in W(R)}\prod_i\det(I-t_iw\widetilde\gamma)^{-1}.}
\end{equation}
Changing a lift multiplies it by an element of $W(R)$ and merely permutes the
inner sum, so \eqref{nd:lattices:eq:niemeier} is lift-independent.  If a lift cycles
isomorphic root components with length $\ell$ and total diagram twist
$\delta_C$, that component cycle reduces to
\begin{equation}\label{nd:lattices:eq:niemeier-component}
 \frac1{|W(R_a)|}\sum_{x\in W(R_a)}
 \prod_i\det(I-t_i^\ell x\delta_C)^{-1}.
\end{equation}
Thus all 23 rootful cases are finite Weyl/quotient class calculations.  They
are a rank-$24$ collection, not a rank-asymptotic tower.

\subsubsection*{Coxeter--Todd, extremal, modular, and neighbor constructions}
For the rank-$12$ Coxeter--Todd lattice $K_{12}$, and for every specified
extremal or modular lattice, the exact answer is simply
\begin{equation}\label{nd:lattices:eq:isolated}
 \MGF_L=\sum_{[g]\subseteq\Aut(L)}\frac1{|C(g)|}
 \prod_{i,m}(1-t_i^m)^{-a_m(g)}.
\end{equation}
This is complete once the matrix classes or weighted frame-shape table are
supplied.  Lattice extremality, modularity, strong perfection, a neighbor
relation, or an inclusion $L_n\subset L_{n+1}$ does not by itself impose
compatible automorphism representations.  Consequently no universal
coefficientwise limit follows from those geometric labels alone.

The notebook does not fabricate missing class data for these large isolated
groups.  It verifies the determinant and frame-shape mechanism on exhaustive
small matrix groups; a numerical evaluation of \eqref{nd:lattices:eq:niemeier} or
\eqref{nd:lattices:eq:isolated} requires the corresponding complete external class table.

\subsection{Explicit lattices with only scalar symmetry}

The smallest possible full automorphism group is $\{\pm I\}$.  If $L$ has
rank $d$ and this full group, then
\begin{equation}\label{nd:lattices:eq:scalar-series}
 \boxed{\MGF_L=\frac12\left[
 \prod_i(1-t_i)^{-d}+\prod_i(1+t_i)^{-d}\right].}
\end{equation}
In particular,
\begin{equation}\label{nd:lattices:eq:scalar-low}
 [m_{(1,1)}]=d^2,\qquad [m_{(2)}]=\frac{d(d+1)}2.
\end{equation}
The verifier uses the following concrete positive-definite integral Gram
matrices:
\begin{align*}
 Q_2&=\begin{pmatrix}6&-1\\-1&8\end{pmatrix},\\
 Q_3&=\begin{pmatrix}6&1&2\\1&7&-1\\2&-1&10\end{pmatrix},\\
 Q_4&=\begin{pmatrix}6&1&1&2\\1&7&1&2\\1&1&11&1\\2&2&1&13\end{pmatrix}.
\end{align*}

\begin{proposition}\label{nd:lattices:prop:scalar-grams}
For $d=2,3,4$, the lattice with Gram matrix $Q_d$ has full automorphism group
$\{\pm I_d\}$.
\end{proposition}

\begin{proof}
Each matrix is strictly diagonally dominant with positive diagonal, hence
positive definite.  Its Gershgorin lower bound $\delta_d$ satisfies
$v^TQ_dv\ge\delta_d\|v\|_2^2$.  If $A^TQ_dA=Q_d$ and $v_j$ is column $j$ of
$A$, then $v_j^TQ_dv_j=(Q_d)_{jj}$.  Therefore every coordinate of every
column has absolute value at most
\[
 B_d=\left\lfloor\sqrt{\max_j(Q_d)_{jj}/\delta_d}\right\rfloor.
\]
The values $(\delta_d,B_d)$ are $(5,1),(3,1),(2,2)$.  Exhaustive exact
enumeration of the finite boxes $[-B_d,B_d]^d$, retaining precisely the
column tuples with pairings $v_i^TQ_dv_j=(Q_d)_{ij}$, returns two matrices in
each rank: $I_d$ and $-I_d$.  This column-pairing condition is exactly
$A^TQ_dA=Q_d$, so the enumeration is exhaustive.
\end{proof}

Equations~\eqref{nd:lattices:eq:scalar-low} show immediate degree-two divergence in any
rank-growing tower whose full automorphism groups stay scalar.  This is the
basic obstruction to a universal law for arbitrary lattice families.

\subsection{Tensor products of lattice representations}

Let $G=\Aut(L)$ act on $V=V_L$ with character $\chi$.  The local group $G^k$
acts on $V^{\otimes k}$.  Since
\[
 \Tr\bigl((g_1\otimes\cdots\otimes g_k)^r\bigr)
 =\prod_{j=1}^k\chi(g_j^r),
\]
Theorem~\ref{nd:lattices:thm:frame} gives
\begin{equation}\label{nd:lattices:eq:tensor-local}
 \boxed{\MGF_{G^k,V^{\otimes k}}
 =\frac1{|G|^k}\sum_{g_1,\ldots,g_k}
 \exp\!\left[\sum_{r\ge1}\frac{p_r}{r}
 \prod_{j=1}^k\chi(g_j^r)\right].}
\end{equation}
Reordering tensor factors in the multilinear coefficient gives
\begin{equation}\label{nd:lattices:eq:tensor-multilinear}
 \boxed{[m_{(1^s)}]\MGF_{G^k,V^{\otimes k}}
 =\left(\dim(V^{\otimes s})^G\right)^k.}
\end{equation}
If $V$ is irreducible orthogonal of dimension greater than one, then
$\dim(V^{\otimes2})^G=1$, while
\[
 \dim(V^{\otimes4})^G=\dim\End_G(V\otimes V)\ge2.
\]
The last inequality holds because $V\otimes V$ contains the trivial module
and at least one nontrivial constituent.  Hence
\begin{equation}\label{nd:lattices:eq:tensor-growth}
 [m_{(1,1,1,1)}]\MGF_{G^k,V^{\otimes k}}\ge2^k.
\end{equation}
Tensor towers therefore usually exhibit exponential moment divergence,
unlike the repeated-sum law in Theorem~\ref{nd:lattices:thm:wreath-sum}.

If factor permutations are included, write
$x=(g_1,\ldots,g_k;\pi)\in G\wr S_k$.  For a cycle $C$ of $\pi$, let $a_C$
be the ordered cycle product and $d_C(r)=\gcd(|C|,r)$.  The $r$th power splits
$C$ into $d_C(r)$ cycles, each carrying a conjugate of
$a_C^{r/d_C(r)}$.  Contracting indices around those cycles proves
\begin{equation}\label{nd:lattices:eq:tensor-wreath-trace}
 \boxed{\Tr(x^r\mid V^{\otimes k})
 =\prod_C\chi\!\left(a_C^{r/d_C(r)}\right)^{d_C(r)}.}
\end{equation}
Substitution into \eqref{nd:lattices:eq:trace-frame} is the exact tensor-wreath
$\sigma$-MGF.

\subsection{Exact computational verification}

The companion marimo notebook calls a pure-Python exact verifier.  No
floating-point eigenvalue or numerical determinant enters a pass condition.
The cases are grouped by the formulas they test.

\subsubsection*{Matrix constructions}
\begin{itemize}
\item \textbf{Type $A$.}  For $n=3,4,5$, every $\pm\sigma$ is represented on
the integral basis $e_1-e_n,\ldots,e_{n-1}-e_n$.  The matrices preserve the
Gram matrix with diagonal $2$ and off-diagonal $1$.
\item \textbf{Signed permutations.}  Every signed permutation matrix is
enumerated in dimensions $2,3,4$ and checked against the identity Gram
matrix.  These are structural tests of \eqref{nd:lattices:eq:b-formula}; the dimension-$4$
case is not claimed to be the full $D_4$ automorphism group.
\item \textbf{Scalar lattices.}  All automorphisms of $Q_2,Q_3,Q_4$ are
enumerated by the exhaustive bound in Proposition~\ref{nd:lattices:prop:scalar-grams}.
\item \textbf{Repeated sums.}  All $288$ block-monomial matrices of
$\Aut(A_2)\wr S_2$ on $A_2\perp A_2$ are constructed.
\item \textbf{Tensor products.}  Kronecker matrices construct
$\Aut(A_2)^2$ and $\Aut(A_2)\wr S_2$ on $V_{A_2}^{\otimes2}$.  Their abstract
orders are $144$ and $288$, while the matrix-image orders are $72$ and $144$
because $(-I,-I)$ acts trivially.
\end{itemize}

\subsubsection*{Four independent checks}
For each concrete matrix $A$, the direct route computes
$T_r=\Tr(A^r)$ and then
\begin{equation}\label{nd:lattices:eq:newton-audit}
 h_0(A)=1,\qquad qh_q(A)=\sum_{r=1}^qT_rh_{q-r}(A).
\end{equation}
It averages $\prod_jh_{\tau_j}(A)$, which is the Reynolds trace on
$W_\tau$.  Independently, the verifier recovers $a_m(A)$ by
\eqref{nd:lattices:eq:trace-frame-inversion}, checks the polynomial identity
\[
 \det(I-tA)=\prod_m(1-t^m)^{a_m(A)},
\]
including all negative exponents by cross-multiplication, and expands the
frame product.  A third route uses only the family formulas: ordinary cycles,
signed cycle products, marked wreath cycles, or the tensor trace
\eqref{nd:lattices:eq:tensor-wreath-trace}.

Finally, for $\tau=(1),(2),(1,1),(3)$, the code explicitly constructs the
integer matrices on each $\Sym^{\tau_j}V$, forms their tensor products, and
solves all generator-fixed equations over $\mathbb Q$.  This common-kernel
dimension is compared with the Reynolds trace.  Thus the fixed-space check is
not merely another expansion of the averaged determinant.

\subsubsection*{Representative cases and results}
\begin{center}
\small
\begin{tabular}{@{}llrr@{}}\toprule
family & matrix image & dimension & image order\\\midrule
type $A$ & $\Aut(A_2),\Aut(A_3),\Aut(A_4)$ & $2,3,4$ & $12,48,240$\\
signed & $C_2\wr S_2,C_2\wr S_3,C_2\wr S_4$ & $2,3,4$ & $8,48,384$\\
scalar & $\Aut(Q_2),\Aut(Q_3),\Aut(Q_4)$ & $2,3,4$ & $2,2,2$\\
orthogonal sum & $\Aut(A_2)\wr S_2$ & $4$ & $288$\\
local tensor & $\rho(\Aut(A_2)^2)$ & $4$ & $72$\\
tensor wreath & $\rho(\Aut(A_2)\wr S_2)$ & $4$ & $144$\\\bottomrule
\end{tabular}
\end{center}

Selected degree-four outputs illustrate the different asymptotics:
\begin{center}
\small
\begin{tabular}{@{}lrrrr@{}}\toprule
matrix image & $[m_2]$ & $[m_{11}]$ & $[m_{22}]$ & $[m_{1111}]$\\\midrule
$\Aut(A_2)$ &1&1&2&3\\
$\Aut(A_3)$ &1&1&3&4\\
$\Aut(A_4)$ &1&1&3&4\\
$C_2\wr S_n$, $n=2,3,4$ &1&1&3&4\\
$\Aut(Q_2)$ &3&4&9&16\\
$\Aut(Q_3)$ &6&9&36&81\\
$\Aut(Q_4)$ &10&16&100&256\\
$\Aut(A_2)\wr S_2$ on $V^{\oplus2}$ &1&1&4&6\\
$\rho(\Aut(A_2)^2)$ on $V^{\otimes2}$ &1&1&5&9\\
$\rho(\Aut(A_2)\wr S_2)$ on $V^{\otimes2}$ &1&1&4&6\\\bottomrule
\end{tabular}
\end{center}
All odd total-degree values vanish, as required by \eqref{nd:lattices:eq:parity}.  For
the local tensor case, the value $9=3^2$ at $m_{1111}$ verifies
\eqref{nd:lattices:eq:tensor-multilinear}; adjoining the factor swap takes invariants in
the symmetric square of a three-dimensional space and gives
$\binom{3+1}{2}=6$.

The degree-four audit contains twelve partitions per case and three recorded
coefficient routes, for $12\cdot12\cdot3=432$ values.  The fixed-space audit
has $12\cdot4=48$ comparisons.  There are $1{,}250$ distinct matrices; six
frame checks per matrix (one determinant identity and five complete
characters through degree four) give $7{,}500$ elementwise checks.

\subsection{Asymptotic classification and scope}

\begin{center}
\small
\begin{tabularx}{\linewidth}{@{}>{\raggedright\arraybackslash}p{0.25\linewidth}
>{\raggedright\arraybackslash}p{0.29\linewidth}X@{}}\toprule
family & exact formula & asymptotic conclusion\\\midrule
$A_{n-1}$ & \eqref{nd:lattices:eq:a-formula} & Poisson-cycle limit \eqref{nd:lattices:eq:a-limit}\\
$\Z^n$, types $B/C$, $D_n$ for $n\ge5$ & \eqref{nd:lattices:eq:b-formula} & stable signed-cycle law \eqref{nd:lattices:eq:b-limit}\\
$R^{\perp k}$ with component permutations & \eqref{nd:lattices:eq:wreath-sum} & compound-Poisson limit \eqref{nd:lattices:eq:wreath-limit}\\
rootful Niemeier & \eqref{nd:lattices:eq:niemeier} & finite rank-$24$ collection\\
$K_{12}$ and isolated extremal/modular lattices & \eqref{nd:lattices:eq:isolated} & sequence-dependent; repeated sums use \eqref{nd:lattices:eq:wreath-limit}\\
$\Aut(L)=\{\pm I\}$ & \eqref{nd:lattices:eq:scalar-series} & degree-two divergence\\
$V_L^{\otimes k}$ under $G^k$ & \eqref{nd:lattices:eq:tensor-local} & degree-four exponential growth\\
neighbor or laminated towers & \eqref{nd:lattices:eq:class-frame} & requires compatible group data\\\bottomrule
\end{tabularx}
\end{center}

The structural distinction is sharp: orthogonal repeated sums with component
permutations yield stable marked compound-Poisson frame-shape laws, whereas
tensor powers and low-symmetry towers force divergent invariant moments.
The universal object is therefore the weighted frame-shape distribution
\[
 \boxed{\MGF_L=\E_{g\in\Aut(L)}
 \prod_{i,m}(1-t_i^m)^{-a_m(g)},}
\]
together with whatever compatibility a specified lattice tower supplies.

\subsection{Reproducibility}

From the repository root, run
\begin{verbatim}
python3 -m lambda_distributions.dists.lattice_sigma_mgfs.verification
python3 -m pytest -q lambda_distributions/dists/lattice_sigma_mgfs/test_verification.py
marimo edit marimo_notebooks/lattices.py
\end{verbatim}
The LaTeX source is standalone and can be rebuilt in its folder with
\begin{verbatim}
latexmk -pdf -interaction=nonstopmode -halt-on-error \
  lattice_sigma_mgfs_proof_and_verification.tex
\end{verbatim}
No command in this workflow modifies
\texttt{proved\_matrix\_groups/all\_new\_matrix\_groups\_proofs.pdf}.

\clearpage
\part{Compact classical, exceptional, and Schur-functor families}
\section{Classical compact groups and weighted tori}\label{proof:classical-compact}
\begingroup
\providecommand{\Sym}{\operatorname{Sym}}
\renewcommand{\Sym}{\operatorname{Sym}}
\providecommand{\CT}{\operatorname{CT}}
\renewcommand{\CT}{\operatorname{CT}}
\providecommand{\SU}{\operatorname{SU}}
\renewcommand{\SU}{\operatorname{SU}}
\providecommand{\SO}{\operatorname{SO}}
\renewcommand{\SO}{\operatorname{SO}}
\providecommand{\OO}{\operatorname{O}}
\renewcommand{\OO}{\operatorname{O}}
\providecommand{\EE}{\mathbb E}
\renewcommand{\EE}{\mathbb E}
\setcounter{theorem}{0}
\setcounter{proposition}{0}
\setcounter{lemma}{0}
\setcounter{corollary}{0}
\setcounter{remark}{0}
\begin{abstract}
We prove the coefficient formulas proposed for the defining representations of
$\SU(n)$, $\SO(n)$, $\OO(n)$, and arbitrary compact tori.  We then describe an
independent matrix calculation: connected-group invariants are the common
kernel of the derived representation, while torus integrals are replaced by
provably alias-free finite diagonal grids.  The accompanying marimo notebooks
perform 100 exact integer comparisons.  All pass.  We also state two necessary
corrections explicitly: orthogonal Schur shapes must have length at most $n$,
and all formulas are interpreted coefficientwise (or after finite truncation).
\end{abstract}

\subsection{Imported compact-group identity}
For each compact group below, $\MGF_{G,V}=\smgv{G}{V}$ and the coefficient
formula are Tools~\ref{tool:molien} and~\ref{tool:coefficient}. We begin with
the branching rules distinguishing $SU(n)$, $SO(n)$, $O(n)$, and weighted
tori.

\subsection{The $\SU(n)$ coefficient formula}
As a polynomial $GL_n$-module,
\[
 \operatorname{ch}W_\tau=h_\tau
 =\sum_{\substack{\lambda\vdash |\tau|\\\ell(\lambda)\le n}}
 K_{\lambda,\tau}s_\lambda.
\]
The restriction of $S_\lambda V$ to $\SU(n)$ is trivial precisely when its
highest-weight differences vanish, namely when
$\lambda=(q^n)=(q,\ldots,q)$.  In that case
$S_{(q^n)}V\cong (\det V)^q$, which is trivial on $\SU(n)$.  Hence
\begin{theorem}
For the defining representation of $\SU(n)$,
\[
 \boxed{\dim W_\tau(V)^{\SU(n)}=
 \begin{cases}
 K_{(q^n),\tau},&|\tau|=nq,\\
 0,&n\nmid |\tau|.
 \end{cases}}
\]
\end{theorem}
The first correction to the $U(n)$ answer occurs at degree $n$ and includes
$\tau=(1^n)$, represented by the determinant tensor.

\subsection{The $\SO(n)$ and $\OO(n)$ formulas}
The orthogonal branching criterion says
\[
 \dim(S_\lambda V)^{\OO(n)}=1
 \quad\Longleftrightarrow\quad
 \lambda=2\delta,
\]
provided $\ell(\lambda)\le n$; otherwise $S_\lambda(\mathbb C^n)=0$.
Because $\SO(n)$ has index two in $\OO(n)$, an $\SO(n)$-fixed line is either
trivial or determinant-isotypic under $\OO(n)$.  Tensoring by determinant adds
one box to each of exactly $n$ rows.  Therefore
\begin{theorem}
For $n\ge2$ and the defining representation,
\begin{align*}
 \dim W_\tau(V)^{\OO(n)}
 &=\sum_{\substack{\delta\\\ell(2\delta)\le n}}K_{2\delta,\tau},\\
 \dim W_\tau(V)^{\SO(n)}
 &=\sum_{\substack{\delta\\\ell(2\delta)\le n}}K_{2\delta,\tau}
 +\sum_{\substack{\delta\\\ell(\delta)\le n}}
 K_{2\delta+(1^n),\tau}.
\end{align*}
\end{theorem}
The length restriction is essential in finite dimension.  The orientation
correction begins at degree $n$, occurs only in degrees congruent to $n$ modulo
two, and contains the Levi-Civita tensor at $\tau=(1^n)$.

\subsection{Weighted compact tori}
Let $T^r$ act diagonally on $\mathbb C^d$ with integer weights
$w_1,\ldots,w_d\in\mathbb Z^r$.  Orthogonality of torus characters gives
\[
 \boxed{\MGF_{T^r,V}
 =\CT_{z_1,\ldots,z_r}
 \prod_{i\ge1}\prod_{j=1}^d(1-t_i z^{w_j})^{-1}.}
\]
Consequently, $\dim W_\tau(V)^{T^r}$ is the number of arrays
$a_{ij}\in\mathbb Z_{\ge0}$ satisfying
\[
 \sum_j a_{ij}=\tau_i,
 \qquad
 \sum_{i,j}a_{ij}w_j=0.
\]

\begin{lemma}[Alias-free finite grid]
Suppose every tested monomial has each weight coordinate in $[-B,B]$.
For any integer $M>2B$, averaging over the explicit subgroup
$(\mu_M)^r\subset T^r$ selects exactly the zero integer weight.
\end{lemma}
\begin{proof}
The finite average selects weights congruent to zero modulo $M$.  The only
multiple of $M$ in $[-B,B]$ is zero.  Thus the finite matrix average equals the
continuous Haar constant term throughout the tested truncation.
\end{proof}

\subsection{Independent direct matrix computation}
The notebook does not recover the left side from Schur decomposition.  In the
normalized occupancy basis $|a_1,\ldots,a_n\rangle$ of $\Sym^dV$, a matrix unit
$E_{uv}$ acts infinitesimally by
\[
 d\Sym^d(E_{uv})|a\rangle=
 \begin{cases}
 \sqrt{a_v(a_u+1)}\,|a+e_u-e_v\rangle,&u\ne v,\\
 a_u|a\rangle,&u=v.
 \end{cases}
\]
On $W_\tau$ the derived action is the Kronecker sum over tensor factors.
The common nullity for a basis of $\mathfrak{sl}_n(\mathbb C)$ gives the
$\SU(n)$ invariants; skew matrix units give $\SO(n)$ invariants.  Adding the
constraint $\rho(\operatorname{diag}(-1,1,\ldots,1))-I$ gives $\OO(n)$
invariants.  This is valid because invariants of a connected matrix group are
exactly the vectors killed by its derived Lie algebra.

Kostka numbers on the formula side are counted separately by backtracking over
semistandard tableaux.  Thus the two columns share only the mathematical claim
being tested, not an implementation route.

\subsection{Test matrix and outcome}
\begin{center}
\begin{tabular}{@{}llllr@{}}
\toprule
Family & Dimensions & Tested degrees & Direct mechanism & Checks\\
\midrule
$\SU(n)$ & $n=2,3$ & $0$--$4$ representative $\tau$ & Lie-kernel nullity & 15\\
$\SO(n)$ & $n=2,3,4$ & through orientation degree & Lie-kernel nullity & 22\\
$\OO(n)$ & $n=2,3,4$ & through orientation degree & kernel plus reflection & 20\\
$T^r$ & $r=1,2$; $d=2,3$ & all $\tau$ through degree $4$ or $5$ & finite grid average & 43\\
\midrule
Total & & & & \textbf{100}\\
\bottomrule
\end{tabular}
\end{center}
All 100 comparisons pass.  The largest directly formed continuous-group
target has dimension 256.  Torus grids use the conservative bound
$M=2D\max_{j,k}|(w_j)_k|+1$, so the no-aliasing lemma applies.

\subsection{Reproduction}
From the repository root, run the following modules with \texttt{python3 -m}:
\begin{itemize}
\small
\item \path{for_this_guy.matrix_group_formula_verification.classical_compact.verification}
\item \path{for_this_guy.matrix_group_formula_verification.tori.verification}
\end{itemize}
Open the corresponding \path{marimo_notebooks/classical_compact.py} and
\path{marimo_notebooks/tori.py} files with \texttt{marimo edit} for the interactive
tables.
The verification ranges are deliberately small enough to inspect and large
enough to include determinant/orientation corrections.  These computations
test the stated dimensions; they do not replace the proofs above or establish
unlisted exceptional-group simplifications.
\endgroup


\clearpage
\section[The compact G2 sigma-MGF on V7]{The compact $G_2$ sigma-MGF on $V_7$}\label{proof:g2}
\begingroup
\definecolor{navy}{RGB}{24,55,95}
\providecommand{\Sg}{\SMG}
\renewcommand{\Sg}{\SMG}
\providecommand{\CT}{\operatorname{CT}}
\renewcommand{\CT}{\operatorname{CT}}
\setcounter{theorem}{0}
\setcounter{proposition}{0}
\setcounter{lemma}{0}
\setcounter{corollary}{0}
\setcounter{remark}{0}
\begin{abstract}
We prove the compact Molien--Weyl formula for the minimal seven-dimensional
$G_2$-module, construct its maximal-torus matrices from the weights, and
verify every Schur coefficient through degree three exactly.  The result is
$\Sg_{G_2,7}=1+h_2+e_3+O_{\ge4}$, and the ordinary specialization is
$(1-t^2)^{-1}$.
\end{abstract}

\subsection{Imported compact identity}
The compact Molien and coefficient formulas are imported. Weyl integration
reduces $\smgv{G_2}{V_7}$ to the maximal-torus constant term below.

\subsection{Specialization to the minimal module}
The weights of $V_7$ are zero together with the six short roots.  Since
$|W(G_2)|=12$,
\[
 \boxed{\Sg_{G_2,7}=\frac1{12}\CT_z\left[D_{G_2}(z)
 \prod_a\left((1-t_a)^{-1}\prod_{\beta\in\Phi_{\rm short}}
 (1-t_a z^\beta)^{-1}\right)\right].}
\]
The invariant metric lies in $\operatorname{Sym}^2V_7$, while the octonionic
associative form lies in $\bigwedge^3V_7$.  Their uniqueness and the absence
of other invariants below degree four give
\[
 \boxed{\Sg_{G_2,7}=1+h_2+e_3+O_{\ge4}.}
\]
A standard invariant-theory theorem for the minimal $G_2$ module states that
$\mathbb C[V_7]^{G_2}=\mathbb C[q]$, where $q$ is the invariant quadratic
norm.  Using that external theorem,
$\Sg_{G_2,7}(t,0,\ldots)=(1-t^2)^{-1}$.  The exact Weyl computation in this
note independently verifies the coefficients through degree three.

\subsection{Exact matrix verification}
The notebook generates the root system from the $G_2$ Cartan matrix and forms
\[
 M(\theta)=\operatorname{diag}\bigl(e^{i\langle\mu,\theta\rangle}\bigr)_{
 \mu\in\{0\}\cup\Phi_{\rm short}}\in U(7).
\]
It compares $\prod_a\det(I-t_aM)^{-1}$ directly with the product over the
seven diagonal eigenvalues.  For integration it evaluates Schur characters by
Jacobi--Trudi and uses the one-sided Weyl denominator
\[
 \prod_{\alpha>0}(1-z^{-\alpha})
 =\sum_{w\in W}\det(w)z^{w\rho-\rho}.
\]
Thus every coefficient is an exact integer, with no quadrature tolerance.
In the requested monomial basis the exact degree-at-most-three signature is
\[
([m_{(2)}],[m_{(1,1)}],[m_{(3)}],[m_{(2,1)}],[m_{(1,1,1)}])
=(1,1,0,0,1).
\]

\begin{center}\small
\begin{tabular}{@{}c|rrrrrrr@{}}\toprule
$\lambda$&$\varnothing$&$(1)$&$(2)$&$(1,1)$&$(3)$&$(2,1)$&$(1,1,1)$\\\midrule
proposed&1&0&1&0&0&0&1\\
exact Weyl check&1&0&1&0&0&0&1\\\bottomrule
\end{tabular}
\end{center}
The sampled torus matrix has unitarity residual $3.2\times10^{-16}$ and the
two integrand evaluations differ by less than $9\times10^{-16}$.  All seven
coefficient checks pass.

\paragraph{Reproducibility.}
Run \texttt{marimo\_notebooks/lie\_group\_g2.py} with marimo.  The shared exact-arithmetic
module is in the parent folder.

\subsection*{External invariant-theory source}
See J.~F.~Adams, \emph{Lectures on Exceptional Lie Groups}, Chicago
Lectures in Mathematics (1996), for the minimal $G_2$ module and its invariant
metric.  The all-degree one-copy invariant-ring statement is an external
input; the low-degree Weyl checks above are internal to this artifact.
\endgroup


\clearpage
\section[The compact Spin(7) sigma-MGF on S8]{The compact $\operatorname{Spin}(7)$ sigma-MGF on $S_8$}\label{proof:spin7}
\begingroup
\definecolor{navy}{RGB}{24,55,95}
\providecommand{\Sg}{\SMG}
\renewcommand{\Sg}{\SMG}
\providecommand{\CT}{\operatorname{CT}}
\renewcommand{\CT}{\operatorname{CT}}
\setcounter{theorem}{0}
\setcounter{proposition}{0}
\setcounter{lemma}{0}
\setcounter{corollary}{0}
\setcounter{remark}{0}
\begin{abstract}
For the complexification of the real eight-dimensional spin representation,
we prove the compact Molien--Weyl formula, construct explicit maximal-torus
matrices, and verify all coefficients through total degree four by exact
constant terms.  In the Schur basis the answer is
$1+h_2+h_4+s_{(2,2)}+e_4+O_{\ge6}$.  The coefficient of
$m_{(1,1,1,1)}$ is therefore $4$, with the extra invariant supplied by the
Cayley alternating four-form.
\end{abstract}

\subsection{Imported coefficient formula}
Tools~\ref{tool:coefficient} and~\ref{tool:molien} apply to the compact spin
module. The new calculation begins with the $B_3$ weights and Weyl constant
term.

\subsection{Matrix model and Weyl constant term}
Use type $B_3$ with orthonormal coordinates $e_1,e_2,e_3$.  The weights of
$S_8$ are
\[
 \operatorname{Wt}(S_8)=
 \left\{\tfrac12(\epsilon_1e_1+\epsilon_2e_2+\epsilon_3e_3):
 \epsilon_i\in\{\pm1\}\right\}.
\]
For $\theta\in\mathbb R^3$ this constructs the explicit unitary matrix
\[
 M(\theta)=\operatorname{diag}_{\epsilon\in\{\pm1\}^3}
 \exp\!\left(\frac{i}{2}\sum_{j=1}^3\epsilon_j\theta_j\right)\in U(8).
\]
These matrices form the image of a maximal torus of $\operatorname{Spin}(7)$.
Together with the $18$ roots and $|W(B_3)|=48$, the Weyl integration formula
turns the group integral into the exact constant term
\[
 \boxed{\Sg_{\operatorname{Spin}(7),8}
 =\frac1{48}\CT_z\left[
 \prod_{\alpha>0}(1-z^\alpha)(1-z^{-\alpha})
 \prod_{a,\mu\in\operatorname{Wt}(S_8)}(1-t_a z^\mu)^{-1}
 \right].}
\]
Thus it is unnecessary to approximate Haar measure by random matrices.

\subsection{Predicted low-degree tensors}
The central element $-1\in\operatorname{Spin}(7)$ acts by $-I$ on $S_8$,
so all odd-total-degree coefficients vanish.  The tensor square is
\[
 S_8\otimes S_8=\mathbf1\oplus\mathbf7\oplus\mathbf{21}\oplus\mathbf{35},
 \qquad
 \operatorname{Sym}^2S_8=\mathbf1\oplus\mathbf{35},\quad
 \bigwedge^2S_8=\mathbf7\oplus\mathbf{21}.
\]
It is multiplicity-free, hence
\[
 \dim(S_8^{\otimes4})^{\operatorname{Spin}(7)}
 =\dim\operatorname{End}_{\operatorname{Spin}(7)}(S_8\otimes S_8)=4.
\]
Three tensors are the metric pairings.  The fourth is the Cayley form in
$\bigwedge^4S_8^*$.  In Schur notation this yields
\[
 \boxed{\Sg_{\operatorname{Spin}(7),8}
 =1+h_2+h_4+s_{(2,2)}+e_4+O_{\ge6}.}
\]
Changing to the monomial basis gives, through degree four,
\[
\begin{array}{c|rrrrrrr}
\tau&(2)&(1,1)&(4)&(3,1)&(2,2)&(2,1,1)&(1,1,1,1)\\ \hline
[m_\tau]&1&1&1&1&2&2&4
\end{array}
\]
Every coefficient of total degree one or three is zero.

\subsection{Independent exact verification}
The marimo notebook generates the $B_3$ root system and the eight-weight
orbit from the Cartan matrix.  It then performs two distinct checks.
First, for a concrete $M(\theta)$ it compares
\[
 \prod_a\det(I-t_aM(\theta))^{-1}
 \quad\text{with}\quad
 \prod_{a,\mu}(1-t_a e^{i\langle\mu,\theta\rangle})^{-1}.
\]
Second, it forms the characters $h_{\tau_1}\cdots h_{\tau_r}$ directly from
the weights and applies the one-sided Weyl-denominator identity
\[
 \prod_{\alpha>0}(1-z^{-\alpha})
 =\sum_{w\in W}\det(w)z^{w\rho-\rho}.
\]
The resulting constant terms are exact integers and do not use the quoted
tensor-square decomposition as input.

The exact output is
\[
\begin{array}{c|rrrrrrr}
\tau&(2)&(1,1)&(4)&(3,1)&(2,2)&(2,1,1)&(1,1,1,1)\\ \hline
\text{proposed}&1&1&1&1&2&2&4\\
\text{Weyl check}&1&1&1&1&2&2&4
\end{array}
\]
with every omitted partition of degree at most four returning zero.  In the
Schur basis the nonzero multiplicities are
\[
 [s_{(2)}]=[s_{(4)}]=[s_{(2,2)}]=[s_{(1,1,1,1)}]=1.
\]

\paragraph{Reproducibility.}
Run the marimo notebook
\path{marimo_notebooks/lie_group_spin7.py}.
Its shared
exact-arithmetic core is in the parent folder.  The test suite also executes
this degree-four case noninteractively.
\endgroup


\clearpage
\section{Compact Spin groups and a Pin component reduction}\label{proof:new-dist:spin_pin}

\resultbox{\textbf{Canonical testing artifacts.}\par\smallskip Exact backend: \path{lambda_distributions/dists/verification.py}.\par Regression: \path{lambda_distributions/dists/test_verification.py}.\par Marimo: \path{marimo_notebooks/spin_pin.py}.}

\subsection*{Scope and outcome}

We prove the exact compact Molien--Weyl formulas for the spin representation
of $\operatorname{Spin}(2n+1)$ and the two half-spin representations of
$\operatorname{Spin}(2n)$.  Clifford algebra decompositions give the decisive
fourth multilinear coefficients $n+1$ and, respectively,
$n/2+1$ or $(n-1)/2$.  They grow with rank, so the unnormalized spinor
sigma-MGFs have no coefficientwise large-rank limit.  We also state the exact
two-component Pin formula, distinguish the $\operatorname{Pin}^{+}$ and
$\operatorname{Pin}^{-}$ lifts, and record the mixed spinor--vector and
adjoint consequences.  An accompanying marimo notebook independently
constructs Weyl characters and Clifford matrices.  All 76 exact and
matrix-level checks pass through rank five where applicable.

\subsection{Imported compact coefficient formula}
Let $G$ be compact and $W$ a finite-dimensional unitary module. The following
is Tools~\ref{tool:molien} and~\ref{tool:coefficient}, followed by Weyl
integration; it is stated with labels because later spin formulas use it.
\begin{theorem}[Imported compact Molien--Weyl formula]
\label{nd:spin_pin:thm:molien}
\begin{align}
\MGF_{G,W}(\mathbf t)
&=\int_G\prod_{i\ge1}\det(I-t_i\rho(g))^{-1}\,dg,
\label{nd:spin_pin:eq:det-molien}\\
&=\int_G\exp\!\left(\sum_{r\ge1}\frac{p_r(\mathbf t)}r
\chi_W(g^r)\right)dg.\label{nd:spin_pin:eq:power-molien}
\end{align}
If $T$ is a maximal torus with positive roots $\Phi^+$ and Weyl group $W_G$,
\begin{equation}\label{nd:spin_pin:eq:weyl-master}
\boxed{\MGF_{G,W}=\frac1{|W_G|}\CT\!\left[
\prod_{\alpha\in\Phi^+}(1-z^\alpha)(1-z^{-\alpha})
\exp\!\left(\sum_{r\ge1}\frac{p_r}{r}\chi_W(z^r)\right)\right].}
\end{equation}
\end{theorem}
For a $W_G$-invariant Laurent polynomial $f$, the Weyl denominator gives the
exact finite lookup
\begin{equation}\label{nd:spin_pin:eq:alternant-lookup}
\frac1{|W_G|}\CT\!\left[
\prod_{\alpha>0}(1-z^\alpha)(1-z^{-\alpha})f(z)\right]
=\sum_{w\in W_G}\det(w)\,[z^{\rho_\Phi-w\rho_\Phi}]f(z).
\end{equation}
No Reynolds or Cauchy proof is repeated here. The new content starts with the
Clifford-theoretic weights, all-rank tensor-square decompositions, and the Pin
component reduction.

\subsection{\texorpdfstring{Type $B_n$: $\operatorname{Spin}(2n+1)$ on its spin module}%
{Type Bn: Spin(2n+1) on its spin module}}

Let $V=V_{2n+1}$ be the vector representation and let $\Delta_n$ be the
$2^n$-dimensional complex spin representation.  In orthonormal torus
coordinates the weights are
\[
 \frac12(\epsilon_1e_1+\cdots+\epsilon_ne_n),
 \qquad \epsilon_j\in\{\pm1\},
\]
each once.  Hence
\begin{equation}\label{nd:spin_pin:eq:b-character}
 \chi_{\Delta_n}(z)=\prod_{j=1}^n(z_j^{1/2}+z_j^{-1/2}).
\end{equation}
Since $|W(B_n)|=2^nn!$, Theorem~\ref{nd:spin_pin:thm:molien} gives
\begin{equation}\label{nd:spin_pin:eq:b-mgf}
 \boxed{\MGF_{\operatorname{Spin}(2n+1),\Delta_n}
 =\frac1{2^nn!}\CT\!\left[\mathcal D_{B_n}(z)
 \exp\!\left(\sum_{r\ge1}\frac{p_r}{r}
 \prod_{j=1}^n(z_j^{r/2}+z_j^{-r/2})\right)\right].}
\end{equation}
Equivalently, if
\[
 H_k^{\Delta_n}(z)=[u^k]\exp\!\left(
 \sum_{r\ge1}\frac{u^r}{r}\chi_{\Delta_n}(z^r)\right),
\]
then every requested coefficient is the finite calculation
\begin{equation}\label{nd:spin_pin:eq:b-coeff}
 [m_\tau]\MGF
 =\frac1{2^nn!}\CT\!\left[
 \mathcal D_{B_n}(z)\prod_jH_{\tau_j}^{\Delta_n}(z)\right].
\end{equation}

The nontrivial element in the kernel of
$\operatorname{Spin}(2n+1)\to SO(2n+1)$ acts as $-I$ on $\Delta_n$.
It follows immediately that
\begin{equation}\label{nd:spin_pin:eq:b-parity}
 [m_\tau]\MGF=0\qquad\text{if $|\tau|$ is odd}.
\end{equation}
The spin module is self-dual.  The Clifford-periodic invariant bilinear form
is symmetric for $n\equiv0,3\pmod4$ and alternating for
$n\equiv1,2\pmod4$.  Therefore
\begin{equation}\label{nd:spin_pin:eq:b-quadratic}
 [m_{(1,1)}]\MGF=1,
 \qquad
 [m_{(2)}]\MGF=\begin{cases}
 1,&n\equiv0,3\pmod4,\\0,&n\equiv1,2\pmod4.
 \end{cases}
\end{equation}
One proof of the symmetry sign starts with the alternating form on the
$SU(2)\cong\operatorname{Spin}(3)$ module and tensors the Pauli-matrix
Clifford construction.  At each two-dimensional Clifford extension the
transpose sign changes in the periodic pattern $--++$.

\begin{theorem}[odd-spin fourth moment]\label{nd:spin_pin:thm:b-fourth}
For every $n\ge1$,
\begin{equation}\label{nd:spin_pin:eq:b-fourth}
 \boxed{\dim(\Delta_n^{\otimes4})^{\operatorname{Spin}(2n+1)}=n+1.}
\end{equation}
\end{theorem}

\begin{proof}
The even Clifford algebra acts faithfully on $\Delta_n$, and the symbol map
from Clifford products to exterior forms is Spin-equivariant.  Hodge duality
then gives the multiplicity-free decomposition
\begin{equation}\label{nd:spin_pin:eq:b-tensor-square}
 \Delta_n\otimes\Delta_n
 \cong\End(\Delta_n)
 \cong\bigoplus_{k=0}^{n}\Lambda^kV_{2n+1}.
\end{equation}
The summands are pairwise nonisomorphic and self-dual.  Since $\Delta_n$ is
self-dual,
\[
 \dim(\Delta_n^{\otimes4})^G
 =\dim\End_G(\Delta_n\otimes\Delta_n)=n+1.
\]
\end{proof}

\subsection{\texorpdfstring{Type $D_n$: half-spin modules of $\operatorname{Spin}(2n)$}%
{Type Dn: half-spin modules of Spin(2n)}}

Let $\Delta_n^+$ and $\Delta_n^-$ be the two half-spin modules, each of
dimension $2^{n-1}$.  Set
\[
 A_n(z)=\prod_{j=1}^n(z_j^{1/2}+z_j^{-1/2}),\qquad
 B_n(z)=\prod_{j=1}^n(z_j^{1/2}-z_j^{-1/2}).
\]
Separating spin weights by the parity of their minus signs gives
\begin{equation}\label{nd:spin_pin:eq:d-character}
 \chi_{\Delta_n^\pm}(z)=\frac12(A_n(z)\pm B_n(z)).
\end{equation}
Because $|W(D_n)|=2^{n-1}n!$,
\begin{equation}\label{nd:spin_pin:eq:d-mgf}
 \boxed{\MGF_{\operatorname{Spin}(2n),\Delta_n^\pm}
 =\frac1{2^{n-1}n!}\CT\!\left[\mathcal D_{D_n}(z)
 \exp\!\left(\sum_{r\ge1}\frac{p_r}{r}
 \chi_{\Delta_n^\pm}(z^r)\right)\right].}
\end{equation}

The half-spin central character has order two for even $n$ and order four
for odd $n$.  Thus
\begin{equation}\label{nd:spin_pin:eq:d-selection}
 [m_\tau]\MGF=0\quad
 \begin{cases}
 \text{if $|\tau|$ is odd},&n\text{ even},\\
 \text{unless $4\mid|\tau|$},&n\text{ odd}.
 \end{cases}
\end{equation}
Moreover,
\[
 (\Delta_n^+)^*\cong
 \begin{cases}\Delta_n^+,&n\text{ even},\\
 \Delta_n^-,&n\text{ odd},\end{cases}
\]
and, in the self-dual case, the form is symmetric for $n\equiv0\pmod4$
and alternating for $n\equiv2\pmod4$.  Consequently
\begin{equation}\label{nd:spin_pin:eq:d-quadratic}
 [m_{(1,1)}]\MGF=\mathbf1_{2\mid n},
 \qquad
 [m_{(2)}]\MGF=\mathbf1_{4\mid n}.
\end{equation}

\begin{theorem}[half-spin fourth moments]\label{nd:spin_pin:thm:d-fourth}
For $n\ge2$,
\begin{equation}\label{nd:spin_pin:eq:d-fourth}
 \boxed{\dim((\Delta_n^+)^{\otimes4})^{\operatorname{Spin}(2n)}
 =\begin{cases}
 n/2+1,&n\text{ even},\\[2pt]
 (n-1)/2,&n\text{ odd}.
 \end{cases}}
\end{equation}
The balanced count is
\begin{equation}\label{nd:spin_pin:eq:d-balanced}
 \boxed{\dim\left((\Delta_n^+\otimes\Delta_n^{+*})^{\otimes2}
 \right)^{\operatorname{Spin}(2n)}=\left\lfloor\frac n2\right\rfloor+1.}
\end{equation}
\end{theorem}

\begin{proof}
The Clifford symbol map and the two Hodge eigenspaces in middle degree give
\begin{equation}\label{nd:spin_pin:eq:d-tensor-square}
 \Delta_n^+\otimes\Delta_n^+
 \cong
 \bigoplus_{\substack{0\le k<n\\k\equiv n\ ({\rm mod}\ 2)}}
 \Lambda^kV_{2n}\ \oplus\ \Lambda^n_{\varepsilon}V_{2n},
\end{equation}
where $\Lambda^n_{\varepsilon}$ is one of the two middle-degree irreducibles.
The tensor square of $\Delta_n^-$ has the same lower-degree summands and the
opposite middle constituent.

If $n$ is even, $\Delta_n^+$ is self-dual and
\eqref{nd:spin_pin:eq:d-tensor-square} has $n/2+1$ pairwise nonisomorphic summands.
The invariant fourth power is its endomorphism algebra, proving the first
line of \eqref{nd:spin_pin:eq:d-fourth}.  If $n$ is odd, invariants are
\[
 \Hom_G(\Delta_n^-\otimes\Delta_n^-,
         \Delta_n^+\otimes\Delta_n^+).
\]
Only the $(n-1)/2$ lower exterior powers occur on both sides; the middle
constituents are opposite.  This proves the second line.

Finally, $\Delta_n^+\otimes\Delta_n^{+*}=\End(\Delta_n^+)$ is
multiplicity-free.  Its exterior-form constituents are the even degrees below
$n$, together with one middle summand when $n$ is even.  Their number is
$\lfloor n/2\rfloor+1$, proving \eqref{nd:spin_pin:eq:d-balanced}.
\end{proof}

The special case $n=3$ is a useful boundary check:
$\operatorname{Spin}(6)\cong SU(4)$ and $\Delta_3^+$ is the defining
four-dimensional representation.  Its first one-copy invariant is the
degree-four determinant, so the fourth multilinear coefficient is one, while
both quadratic coefficients vanish.

\subsection{Asymptotic obstruction}

\begin{corollary}\label{nd:spin_pin:cor:no-limit}
Neither the unnormalized odd-spin sequence
$\MGF_{\operatorname{Spin}(2n+1),\Delta_n}$ nor either half-spin sequence has
a finite coefficientwise limit in the completed symmetric-function ring.
\end{corollary}

\begin{proof}
The fixed coefficient $[m_{(1,1,1,1)}]$ equals $n+1$ in type $B_n$ and is
asymptotic to $n/2$ in type $D_n$.  It diverges along every sequence with
$n\to\infty$.  Coefficientwise convergence would require this scalar
coefficient to converge, which is impossible.
\end{proof}

The conclusion is not contradicted by $\dim\Delta_n=2^n$.  Increasing rank
opens a new exterior-power channel in \eqref{nd:spin_pin:eq:b-tensor-square} or every two
ranks in \eqref{nd:spin_pin:eq:d-tensor-square}.  A normalized, logarithmic, or cumulant
renormalization may still have a limit, but the raw sigma-MGF does not.

\subsection{Pin groups: the disconnected component formula}

Fix a concrete complex pinor representation $\Pi_N^\varepsilon$ of
$\operatorname{Pin}^\varepsilon(N)$, $\varepsilon\in\{+,-\}$, and choose
$u_\varepsilon$ in the odd component.  Since
\[
 \operatorname{Pin}^\varepsilon(N)
 =\operatorname{Spin}(N)\sqcup
 u_\varepsilon\operatorname{Spin}(N),
\]
normalized Haar measure assigns mass $1/2$ to each component.  Therefore
\begin{align}\label{nd:spin_pin:eq:pin-components}
 \MGF_{\operatorname{Pin}^\varepsilon(N),\Pi_N^\varepsilon}
 ={}&\frac12\int_{\operatorname{Spin}(N)}
 \exp\!\left(\sum_{r\ge1}\frac{p_r}{r}
 \chi_{\Pi_N^\varepsilon}(g^r)\right)dg\\
 &+\frac12\int_{\operatorname{Spin}(N)}
 \exp\!\left(\sum_{r\ge1}\frac{p_r}{r}
 \chi_{\Pi_N^\varepsilon}((u_\varepsilon g)^r)\right)dg.\notag
\end{align}
This is an exact identity, but the second integrand depends on the chosen Pin
extension, not only on the restricted Spin character.

For $N=2n$, the restriction is $\Delta_n^+\oplus\Delta_n^-$.  An odd
Clifford generator anticommutes with chirality and is therefore block
off-diagonal.  Every odd power remains block off-diagonal, so
\begin{equation}\label{nd:spin_pin:eq:pin-odd-trace}
 \chi_{\Pi_{2n}^\varepsilon}(x^{2r+1})=0
 \qquad(x\notin\operatorname{Spin}(2n)).
\end{equation}
Only the $p_{2r}$ terms remain in the odd-component exponential.  The two Pin
groups must not be identified: in the concrete Clifford models used here,
\begin{equation}\label{nd:spin_pin:eq:pin-square}
 u_+=\gamma_1,\quad u_- = i\gamma_1,
 \qquad u_+^2=I,\quad u_-^2=-I.
\end{equation}
Thus their even power characters can differ.

\subsection{Mixed spinor--vector and adjoint modules}

Character calculus immediately extends the exact constant term.  For
$W_{a,b}=\Delta^{\otimes a}\otimes V^{\otimes b}$,
\begin{equation}\label{nd:spin_pin:eq:mixed-character}
 \chi_{W_{a,b}}(z)=\chi_\Delta(z)^a\chi_V(z)^b,
\end{equation}
where
\[
 \chi_{V_{2n+1}}=1+\sum_j(z_j+z_j^{-1}),\qquad
 \chi_{V_{2n}}=\sum_j(z_j+z_j^{-1}).
\]
Substitution of \eqref{nd:spin_pin:eq:mixed-character} into \eqref{nd:spin_pin:eq:weyl-master} is the
exact sigma-MGF at every finite rank.  For a direct sum $\Delta\oplus V$, the
character product is replaced by $\chi_\Delta+\chi_V$.

The adjoint representation factors through $SO(N)$ and is
$\mathfrak{so}_N\cong\Lambda^2V_N$.  Its torus character is
\begin{equation}\label{nd:spin_pin:eq:adjoint-character}
 \chi_{\rm Ad}(z)=\operatorname{rank}(G)+\sum_{\alpha\in\Phi}z^\alpha,
\end{equation}
so \eqref{nd:spin_pin:eq:weyl-master} again gives an exact finite-rank formula.  The
exterior-square identity gives the directly testable trace formula
\begin{equation}\label{nd:spin_pin:eq:adjoint-trace}
 \boxed{\Tr(\operatorname{Ad}(g)^r)=\frac12\left(
 \Tr(V(g)^r)^2-\Tr(V(g)^{2r})\right).}
\end{equation}
Unlike the genuine spinors, $\Lambda^2V_N$ has a coefficientwise stable
orthogonal regime.  For fixed $\tau$ and the safe range $N>2|\tau|$,
\begin{equation}\label{nd:spin_pin:eq:adjoint-stable}
 [m_\tau]\MGF_{{\rm Ad},\infty}
 =\left\langle\ExpS(h_2),h_\tau[e_2]\right\rangle.
\end{equation}
The strict inequality is a sufficient stable range, not a claim of
sharpness in every degree.

\subsection{Finite restrictions and spin covers}

If $H\le\operatorname{Spin}(N)$ is finite and $W$ is a restricted spinor
module, Tool~\ref{tool:molien}, followed by the power-trace change of
coordinates and class grouping, gives
\begin{equation}\label{nd:spin_pin:eq:finite-restriction}
 \smgv{H}{W}=\frac1{|H|}\sum_{h\in H}
 \exp\!\left(\sum_{r\ge1}\frac{p_r}{r}\chi_W(h^r)\right)
 =\sum_{[h]}\frac1{|C_H(h)|}
 \exp\!\left(\sum_{r\ge1}\frac{p_r}{r}\chi_W(h^r)\right).
\end{equation}
The same formula applies to the inverse image $\widetilde W$ of a finite real
reflection group in Pin and to a Schur double cover $\widetilde S_n$ with a
chosen genuine spin character.  The cover and power maps are essential:
ordinary conjugacy classes can split upstairs.  No representation-independent
large-$n$ law exists for $\widetilde S_n$ until a sequence of strict-partition
labels is specified.

\subsection{Independent computational verification}

\subsubsection{Exact Weyl route}

For type $B_n$ the verifier enumerates all signed permutations; for type
$D_n$ it retains those with an even number of sign changes.  In doubled
coordinates,
\[
 2\rho_{B_n}=(2n-1,2n-3,\ldots,1),\qquad
 2\rho_{D_n}=(2n-2,2n-4,\ldots,0).
\]
For each $\tau$, dynamic programming expands
$\prod_jh_{\tau_j}$ from the actual spin weights.  The finite lookup
\eqref{nd:spin_pin:eq:alternant-lookup} then returns the invariant multiplicity.  The
decompositions in Theorems~\ref{nd:spin_pin:thm:b-fourth} and \ref{nd:spin_pin:thm:d-fourth} are not
used by this route.

\begin{center}
\begin{tabular}{@{}c@{\quad}r@{\quad}rrrr@{\quad}c@{}}
\toprule
Group&$\dim W$&$[m_{(2)}]$&$[m_{(1,1)}]$&$[m_{(1^4)}]$&formula&result\\
\midrule
$\operatorname{Spin}(3)$&2&0&1&2&$n+1$&\textcolor{teal}{\textbf{PASS}}\\
$\operatorname{Spin}(5)$&4&0&1&3&$n+1$&\textcolor{teal}{\textbf{PASS}}\\
$\operatorname{Spin}(7)$&8&1&1&4&$n+1$&\textcolor{teal}{\textbf{PASS}}\\
$\operatorname{Spin}(9)$&16&1&1&5&$n+1$&\textcolor{teal}{\textbf{PASS}}\\
$\operatorname{Spin}(11)$&32&0&1&6&$n+1$&\textcolor{teal}{\textbf{PASS}}\\
\bottomrule
\end{tabular}
\end{center}

\begin{center}
\begin{tabular}{@{}c@{\quad}r@{\quad}rrrr@{\quad}c@{}}
\toprule
Group&$\dim\Delta^+$&$[m_{(2)}]$&$[m_{(1,1)}]$&$[m_{(1^4)}]$&balanced fourth&result\\
\midrule
$\operatorname{Spin}(4)$&2&0&1&2&2&\textcolor{teal}{\textbf{PASS}}\\
$\operatorname{Spin}(6)$&4&0&0&1&2&\textcolor{teal}{\textbf{PASS}}\\
$\operatorname{Spin}(8)$&8&1&1&3&3&\textcolor{teal}{\textbf{PASS}}\\
$\operatorname{Spin}(10)$&16&0&0&2&3&\textcolor{teal}{\textbf{PASS}}\\
\bottomrule
\end{tabular}
\end{center}

All odd-total-degree type-$B$ checks vanish.  The type-$D$ suite separately
checks the order-two/order-four selection rule.  In total, 45 exact spin,
half-spin, and balanced coefficient comparisons pass.

\subsubsection{Direct Clifford-matrix route}

Let $X,Y,Z$ be the Pauli matrices.  The implementation uses the
Jordan--Wigner gamma matrices
\[
 \gamma_{2j+1}=Z^{\otimes j}\otimes X\otimes I^{\otimes(n-j-1)},\qquad
 \gamma_{2j+2}=Z^{\otimes j}\otimes Y\otimes I^{\otimes(n-j-1)},
\]
with $Z^{\otimes n}$ appended in odd dimension.  They satisfy
\[
 \gamma_a\gamma_b+\gamma_b\gamma_a=2\delta_{ab}I.
\]
Concrete represented Spin elements are products
\begin{equation}\label{nd:spin_pin:eq:matrix-spin}
 \rho(g)=\prod_{a<b}\exp\!\left(
 \frac{\theta_{ab}}2\gamma_a\gamma_b\right).
\end{equation}
Because $(\gamma_a\gamma_b)^2=-I$, each exponential in
\eqref{nd:spin_pin:eq:matrix-spin} is evaluated exactly as a sine--cosine matrix
polynomial.

On a tensor square, the code constructs the positive quadratic Casimir
\[
 C=-\sum_{a<b}\left(J_{ab}\otimes I+I\otimes J_{ab}\right)^2,
 \qquad J_{ab}=\tfrac12\gamma_a\gamma_b.
\]
Its eigenspace dimensions provide an independent matrix signature of the
irreducible summands.

\begin{center}\small
\begin{tabularx}{\linewidth}{@{}lrrXc@{}}
\toprule
Module&matrix dimension&Casimir blocks&observed block dimensions&result\\
\midrule
$\operatorname{Spin}(3)$ spin&2&2&$1,3$&\textcolor{teal}{\textbf{PASS}}\\
$\operatorname{Spin}(5)$ spin&4&3&$1,5,10$&\textcolor{teal}{\textbf{PASS}}\\
$\operatorname{Spin}(7)$ spin&8&4&$1,7,21,35$&\textcolor{teal}{\textbf{PASS}}\\
$\operatorname{Spin}(9)$ spin&16&5&$1,9,36,84,126$&\textcolor{teal}{\textbf{PASS}}\\
$\operatorname{Spin}(4)$ half-spin&2&2&$1,3$&\textcolor{teal}{\textbf{PASS}}\\
$\operatorname{Spin}(8)$ half-spin&8&3&$1,28,35$&\textcolor{teal}{\textbf{PASS}}\\
\bottomrule
\end{tabularx}
\end{center}

The type-$B$ dimensions are exactly
$\binom{2n+1}{0},\ldots,\binom{2n+1}{n}$, as predicted by
\eqref{nd:spin_pin:eq:b-tensor-square}.  Every Clifford anticommutator residual is zero
in floating arithmetic for these Pauli matrices.

\subsubsection{Pin, mixed, and adjoint checks}

For $N=2,4,6$, the verifier constructs $u_+=\gamma_1$ and
$u_-=i\gamma_1$, a nontrivial Spin torus element, and their product in the
odd component.  All six cases have the correct square, anticommute with
chirality, and have zero traces in powers $1,3,5$.  The largest unitarity
residual is $5.7\times10^{-16}$.

Mixed spinor--vector Kronecker matrices in types $B_3$ and $D_3$ agree with
the product-weight construction and \eqref{nd:spin_pin:eq:mixed-character}.  Explicit
adjoint torus matrices verify \eqref{nd:spin_pin:eq:adjoint-trace} for $r=1,2,3,4$ with
maximum residual $7.2\times10^{-15}$.  Exact Weyl integration in types
$B_2,B_3,D_4$ also recovers one invariant metric and one invariant
multilinear structure tensor.  These are 25 additional checks.  Together
with the six Casimir rows, the complete report contains 76 passing checks.

\subsection{Reproduction and scope}

From the repository root, run
\begin{verbatim}
.venv/bin/python lambda_distributions/dists/verification.py
.venv/bin/python -m pytest -q lambda_distributions/dists/test_verification.py
.venv/bin/marimo edit marimo_notebooks/spin_pin.py
\end{verbatim}
The first command prints the six verification subtotals and fails with a
nonzero exit status if any comparison fails.  The notebook exposes the same
suite with a rank selector and group-separated tables.

The exact coefficient computations establish the tested low-degree cases;
the all-rank claims are proved by Reynolds/Weyl integration and the Clifford
decompositions.  The matrix computations are an independent finite-rank
audit, not a substitute for those proofs.  For Pin, the document proves the
component reduction and the even-dimensional odd-trace simplification.  A
fully expanded Pin sigma-MGF still requires the power characters of the
chosen $\operatorname{Pin}^{\pm}$ extension.

\begin{thebibliography}{9}
\bibitem{nd:spin_pin:GKRS}
B.~Gross, B.~Kostant, P.~Ramond, and S.~Sternberg,
\emph{The Weyl character formula, the half-spin representations, and equal
rank subgroups}, Proc. Natl. Acad. Sci. USA 95 (1998),
\href{https://arxiv.org/abs/math/9808133}{arXiv:math/9808133}.

\bibitem{nd:spin_pin:Koike}
K.~Koike,
\emph{Spin representations and centralizer algebras for the spinor groups},
\href{https://arxiv.org/abs/math/0502397}{arXiv:math/0502397}.

\bibitem{nd:spin_pin:Wenzl}
H.~Wenzl,
\emph{On centralizer algebras for spin representations},
\href{https://arxiv.org/abs/1107.4183}{arXiv:1107.4183}.

\bibitem{nd:spin_pin:LM}
H.~B.~Lawson, Jr. and M.-L.~Michelsohn,
\emph{Spin Geometry}, Princeton Mathematical Series 38,
Princeton University Press, 1989.
\end{thebibliography}

\clearpage
\section[The compact F4 sigma-MGF on V26]{The compact $F_4$ sigma-MGF on $V_{26}$}\label{proof:f4}
\begingroup
\definecolor{navy}{RGB}{24,55,95}
\providecommand{\Sg}{\SMG}
\renewcommand{\Sg}{\SMG}
\providecommand{\CT}{\operatorname{CT}}
\renewcommand{\CT}{\operatorname{CT}}
\setcounter{theorem}{0}
\setcounter{proposition}{0}
\setcounter{lemma}{0}
\setcounter{corollary}{0}
\setcounter{remark}{0}
\begin{abstract}
For the trace-zero Albert module we prove the compact Molien--Weyl expression,
construct the 26-dimensional torus matrices, and exactly verify
$\Sg_{F_4,26}=1+h_2+h_3+O_{\ge4}$.  The one-copy Hilbert series is
$((1-t^2)(1-t^3))^{-1}$.
\end{abstract}

\subsection{Weyl specialization}
By the imported compact Molien identity and Weyl integration, the weights of
$V_{26}$ give
\[
\boxed{\Sg_{F_4,26}=\smgv{F_4}{V_{26}}
=\frac1{1152}\CT_z\!\left[D_{F_4}(z)\prod_a(1-t_a)^{-2}
\prod_{\beta\in\Phi_{\rm short}}(1-t_a z^\beta)^{-1}\right].}
\]
Only this weight specialization is family-specific.

\subsection{Invariant tensors}
The trace-zero Albert algebra carries the restricted trace form and the
restricted determinant.  These are respectively a unique symmetric quadratic
and a unique symmetric cubic invariant.  No other Schur type occurs below
degree four, so
\[
 \boxed{\Sg_{F_4,26}=1+h_2+h_3+O_{\ge4}.}
\]
The standard invariant theory of the trace-zero Albert algebra gives
$\mathbb C[V_{26}]^{F_4}=\mathbb C[q,c]$, with generators of degrees two and
three.  Therefore
$\Sg(t,0,\ldots)=((1-t^2)(1-t^3))^{-1}$.  This all-degree invariant-ring
description is an external input; the notebook verifies only the displayed
low-degree terms.

\subsection{Verification method and evidence}
Starting with the $F_4$ Cartan matrix, the notebook reflects a short simple
root to obtain 24 short roots and appends two zero weights.  It constructs
\[
 M(\theta)=\operatorname{diag}(e^{i\langle\mu,\theta\rangle})_{\mu}\in U(26)
\]
and checks the determinant integrand against its eigenvalue product.  Schur
characters are computed independently by Jacobi--Trudi.  Their Haar integrals
are exact: for a character $f$ invariant under $W$,
\[
 \frac1{|W|}\CT(D_{F_4}f)=\CT\left(f\prod_{\alpha>0}(1-z^{-\alpha})\right),
\]
and the Weyl-denominator identity supplies every coefficient on the right.
The direct monomial-basis check returns
\[
([m_{(2)}],[m_{(1,1)}],[m_{(3)}],[m_{(2,1)}],[m_{(1,1,1)}])
=(1,1,1,1,1).
\]

\begin{center}\small
\begin{tabular}{@{}c|rrrrrrr@{}}\toprule
$\lambda$&$\varnothing$&$(1)$&$(2)$&$(1,1)$&$(3)$&$(2,1)$&$(1,1,1)$\\\midrule
proposed&1&0&1&0&1&0&0\\
exact Weyl check&1&0&1&0&1&0&0\\\bottomrule
\end{tabular}
\end{center}
The generated root count is 48, the matrix dimension is 26, the unitarity
residual is $2.8\times10^{-16}$, and the two integrand evaluations agree to
$1.3\times10^{-14}$.  Every degree-at-most-three coefficient passes.

\paragraph{Reproducibility.} Run \texttt{marimo\_notebooks/lie\_group\_f4.py} with marimo.
\subsection*{External invariant-theory source}
See T.~A.~Springer and F.~D.~Veldkamp, \emph{Octonions, Jordan Algebras and
Exceptional Groups}, Springer Monographs in Mathematics (2000), for the
Albert-algebra realization and its norm invariants.
\endgroup


\clearpage
\section[The compact E6 sigma-MGF on V27]{The compact $E_6$ sigma-MGF on $V_{27}$}\label{proof:e6}
\begingroup
\definecolor{navy}{RGB}{24,55,95}
\providecommand{\Sg}{\SMG}
\renewcommand{\Sg}{\SMG}
\providecommand{\CT}{\operatorname{CT}}
\renewcommand{\CT}{\operatorname{CT}}
\setcounter{theorem}{0}
\setcounter{proposition}{0}
\setcounter{lemma}{0}
\setcounter{corollary}{0}
\setcounter{remark}{0}
\begin{abstract}
For simply connected compact $E_6$ and its minuscule 27-dimensional module,
we prove the constant-term formula, the mod-three selection rule, and the
initial expansion $1+h_3+O_{\ge6}$.  An exact 27-weight computation verifies
all Schur coefficients through degree three.
\end{abstract}

\subsection{Weyl specialization}
Because the 27 weights form $W(E_6)\omega_1$, the imported compact Molien
identity and Weyl integration give
\[
\boxed{\Sg_{E_6,27}=\smgv{E_6}{V_{27}}
=\frac1{51840}\CT_z\!\left[D_{E_6}(z)
\prod_a\prod_{\mu\in W\omega_1}(1-t_a z^\mu)^{-1}\right].}
\]
The new points are the central selection rule and exceptional cubic.

\subsection{Selection rule and one-copy series}
The generator of the center $Z(E_6)\cong\mathbb Z/3$ acts on $V_{27}$ by a
primitive cube root.  It acts on every degree-$d$ tensor construction by its
$d$th power.  Consequently
\[
 [m_\tau]\Sg_{E_6,27}=0\quad\text{unless}\quad 3\mid |\tau|.
\]
The exceptional determinant is the unique symmetric cubic invariant.  Thus
\[
 \boxed{\Sg_{E_6,27}=1+h_3+O_{\ge6}},\qquad
 \boxed{\Sg_{E_6,27}(t,0,\ldots)=\frac1{1-t^3}}.
\]
The module belongs to simply connected $E_6$ and does not descend to the
adjoint quotient $E_6/Z_3$.

The all-degree statement $\mathbb C[V_{27}]^{E_6}=\mathbb C[\det]$ is an
external invariant-theory input for the minimal 27-dimensional module.  The
notebook verifies the cubic invariant and the displayed coefficients but does
not independently prove generation in all degrees.

\subsection{Exact verification}
The code generates the 72 roots, obtains $\omega_1$ by rationally solving
$A^Tc=e_1$, and closes its Weyl orbit under simple reflections.  The resulting
27 weights define the explicit matrix
$M(\theta)=\operatorname{diag}(e^{i\langle\mu,\theta\rangle})\in U(27)$.
The determinant formula and eigenvalue product are evaluated independently.

For coefficients, Jacobi--Trudi produces $\chi_{\mathbb S_\lambda V}$.
The one-sided Weyl denominator
\[
 \prod_{\alpha>0}(1-z^{-\alpha})=
 \sum_{w\in W}\det(w)z^{w\rho-\rho}
\]
makes its integral an exact signed weight sum.  Orbit membership of
$\rho-\mu$ is tested by reflecting it to the dominant chamber.
The monomial-basis check gives zero in degrees one and two and
\[
([m_{(3)}],[m_{(2,1)}],[m_{(1,1,1)}])=(1,1,1).
\]

\begin{center}\small
\begin{tabular}{@{}c|rrrrrrr@{}}\toprule
$\lambda$&$\varnothing$&$(1)$&$(2)$&$(1,1)$&$(3)$&$(2,1)$&$(1,1,1)$\\\midrule
proposed&1&0&0&0&1&0&0\\
exact Weyl check&1&0&0&0&1&0&0\\\bottomrule
\end{tabular}
\end{center}
The orbit contains exactly 27 distinct weights.  Unitarity holds to
$4.3\times10^{-16}$, the integrand residual is below $1.8\times10^{-15}$,
and all coefficient checks pass.

\paragraph{Reproducibility.} Run \texttt{marimo\_notebooks/lie\_group\_e6.py} with marimo.
\subsection*{External invariant-theory source}
See T.~A.~Springer and F.~D.~Veldkamp, \emph{Octonions, Jordan Algebras and
Exceptional Groups}, Springer Monographs in Mathematics (2000), for the
Albert determinant and the $E_6$ structure group.
\endgroup


\clearpage
\section[The compact E7 sigma-MGF on V56]{The compact $E_7$ sigma-MGF on $V_{56}$}\label{proof:e7}
\begingroup
\definecolor{navy}{RGB}{24,55,95}
\providecommand{\Sg}{\SMG}
\renewcommand{\Sg}{\SMG}
\providecommand{\CT}{\operatorname{CT}}
\renewcommand{\CT}{\operatorname{CT}}
\setcounter{theorem}{0}
\setcounter{proposition}{0}
\setcounter{lemma}{0}
\setcounter{corollary}{0}
\setcounter{remark}{0}
\begin{abstract}
For simply connected compact $E_7$ on its minuscule 56, we prove the
Molien--Weyl expression and parity rule, exactly verify the full quadratic
expansion with explicit torus matrices, and state the degree-four expansion
using cited invariant theory.  We also give a sparse quartic verification
algorithm; the present artifact does not report that computation as executed.
\end{abstract}

\subsection{Weyl specialization and parity}
The imported compact Molien identity and Weyl integration yield
\[
\boxed{\Sg_{E_7,56}=\smgv{E_7}{V_{56}}
=\frac1{2903040}\CT_z\!\left[D_{E_7}(z)
\prod_a\prod_{\mu\in W\omega_7}(1-t_a z^\mu)^{-1}\right].}
\]
The nontrivial central element acts as $-I$, so odd total degree vanishes.
This is the simply connected group; the 56 does not descend to $E_7/Z_2$.

\subsection{Low-degree tensors}
Let $\omega\in\bigwedge^2V^*$ be the invariant symplectic form and let
$q\in\operatorname{Sym}^4V^*$ be the Cartan--Freudenthal quartic.  At degree
four, the three pairwise contractions made from two copies of $\omega$
decompose under $S_4$ as the Schur types $(2,2)$ and $(1,1,1,1)$.  The quartic
$q$ supplies type $(4)$.  The one-copy statement
$\mathbb C[V_{56}]^{E_7}=\mathbb C[q]$ and the exhaustion of degree-four
tensor invariants by the symplectic pairings and $q$ are external
invariant-theory inputs; see Brown \cite{e7-proof-and-verification:brown}.  Using those inputs,
\[
 \boxed{\Sg_{E_7,56}=1+e_2+h_4+s_{(2,2)}+e_4+O_{\ge6}},
 \qquad \boxed{\Sg(t,0,\ldots)=\frac1{1-t^4}}.
\]
The one-variable specialization kills $e_2,s_{(2,2)},e_4$ and retains the
quartic generator.
In particular $([m_{(2)}],[m_{(1,1)}])=(0,1)$, while the predicted degree-four
row is
\[
([m_{(4)}],[m_{(3,1)}],[m_{(2,2)}],[m_{(2,1,1)}],[m_{(1,1,1,1)}])
=(1,1,2,2,4).
\]

\subsection{Executed quadratic verification and optional quartic confirmation}
The notebook generates the 126 roots and the 56-weight orbit, constructs
$M(\theta)=\operatorname{diag}(e^{i\langle\mu,\theta\rangle})$, and verifies
the determinant integrand.  Jacobi--Trudi characters are integrated exactly
using
\[
 \frac1{|W|}\CT(Df)=\CT\left(f\prod_{\alpha>0}(1-z^{-\alpha})\right).
\]
The complete quadratic computation is
\begin{center}\small
\begin{tabular}{@{}c|rrrr@{}}\toprule
$\lambda$&$\varnothing$&$(1)$&$(2)$&$(1,1)$\\\midrule
proposed&1&0&0&1\\
exact Weyl check&1&0&0&1\\\bottomrule
\end{tabular}
\end{center}
All entries pass.  The matrix unitarity residual is $4.2\times10^{-16}$ and
the determinant/eigenvalue residual is $8.6\times10^{-14}$.

\paragraph{Status of the degree-four row.}
The row follows from the cited invariant-theory decomposition.  The notebook
output displayed here verifies only degrees at most two.  A fully mechanical
confirmation may be obtained as follows, but is not executed in this
artifact.  Do not form dense matrices on
$V^{\otimes4}$ (dimension $56^4$).  Instead: (i) generate the weight-sparse
root-operator action on each of $\operatorname{Sym}^4V$, $\mathbb S_{(2,2)}V$,
and $\bigwedge^4V$; (ii) retain only zero-weight basis vectors; (iii) solve the
sparse raising-operator nullspace; and (iv) check dimensions $1,1,1$.
Equivalently, stream the three Schur characters through the same exact Weyl
functional.  An executed result must report the zero-weight dimensions,
exact raising-kernel ranks, arithmetic mode, and output hash.

\paragraph{Reproducibility.} Run \texttt{marimo\_notebooks/lie\_group\_e7.py} with marimo.
\begin{thebibliography}{9}
\bibitem{e7-proof-and-verification:brown} R.~B.~Brown, ``Groups of type $E_7$,'' \emph{J. Reine Angew.
Math.} \textbf{236} (1969), 79--102.
\bibitem{e7-proof-and-verification:adams} J.~F.~Adams, \emph{Lectures on Exceptional Lie Groups},
Chicago Lectures in Mathematics (1996).
\end{thebibliography}
\endgroup


\clearpage
\section[The compact E8 sigma-MGF on the adjoint module]{The compact $E_8$ sigma-MGF on the adjoint module}\label{proof:e8}
\begingroup
\definecolor{navy}{RGB}{24,55,95}
\providecommand{\Sg}{\SMG}
\renewcommand{\Sg}{\SMG}
\providecommand{\CT}{\operatorname{CT}}
\renewcommand{\CT}{\operatorname{CT}}
\setcounter{theorem}{0}
\setcounter{proposition}{0}
\setcounter{lemma}{0}
\setcounter{corollary}{0}
\setcounter{remark}{0}
\begin{abstract}
We prove the exact constant-term formula for the 248-dimensional adjoint
module, verify its root and zero-weight data and its quadratic invariant
exactly, and prove the alternating cubic row algebraically.  We also
record the correct one-copy Chevalley series and explain why it does not
determine the multivariate sigma-MGF.
\end{abstract}

\subsection{Exact formula}
For every unitary compact-group representation,
\[
 \Sg_{K,V}=\int_K\prod_a\det(I-t_a\rho(g))^{-1}dg
 =\sum_\lambda\dim(\mathbb S_\lambda V)^K s_\lambda(\mathbf t).
\]
Equivalently,
$[m_\tau]\Sg_{K,V}=\dim(\bigotimes_i\operatorname{Sym}^{\tau_i}V)^K$.
The adjoint $E_8$ weights are the 240 roots and zero with multiplicity eight.
Weyl integration therefore proves
\[
 \boxed{\Sg_{E_8,248}=\frac1{696729600}\CT_z\left[D_{E_8}(z)
 \prod_a(1-t_a)^{-8}\prod_{\alpha\in\Phi(E_8)}
 (1-t_a z^\alpha)^{-1}\right].}
\]

\subsection{Universal adjoint tensors}
For a compact simple Lie algebra, the Killing form $\kappa$ is a unique
symmetric quadratic invariant.  In degree three,
\[
 c(x,y,z)=\kappa(x,[y,z])
\]
is invariant, nonzero, and fully alternating.  Moreover
\[
(\bigwedge^3\mathfrak g)^G\cong
\operatorname{Hom}_G(\bigwedge^2\mathfrak g,\mathfrak g),
\]
and for the simple Lie algebra $\mathfrak e_8$ the bracket spans this
Hom-space.  Chevalley restriction gives basic invariant degrees
$2,8,12,14,18,20,24,30$, so $(\operatorname{Sym}^3\mathfrak g)^G=0$.
Since the only invariant in $\mathfrak g^{\otimes3}$ is alternating, the
Schur type $(2,1)$ also has multiplicity zero.  Consequently
\[
 \boxed{\Sg_{E_8,248}=1+h_2+e_3+O_{\ge4}.}
\]
Chevalley restriction gives the one-copy series
\[
 \boxed{\Sg(t,0,\ldots)=\prod_{d\in\{2,8,12,14,18,20,24,30\}}
 (1-t^d)^{-1}.}
\]
This specialization cannot see $e_3$ and therefore does not determine the
full sigma-MGF.
Thus the degree-at-most-three monomial signature is
\[
([m_{(2)}],[m_{(1,1)}],[m_{(3)}],[m_{(2,1)}],[m_{(1,1,1)}])
=(1,1,0,0,1).
\]

\subsection{Executed quadratic check and optional sparse cubic confirmation}
The notebook generates the $E_8$ root orbit (240 vectors), appends eight zero
weights, and constructs a $248$ by $248$ diagonal unitary torus matrix.  It
checks the determinant integrand directly and integrates the quadratic Schur
characters exactly with the Weyl-denominator identity.
\begin{center}\small
\begin{tabular}{@{}c|rrrr@{}}\toprule
$\lambda$&$\varnothing$&$(1)$&$(2)$&$(1,1)$\\\midrule
proposed&1&0&1&0\\
exact Weyl check&1&0&1&0\\\bottomrule
\end{tabular}
\end{center}
The computed dimension is 248, with 240 roots and eight zero weights; all
quadratic entries pass.  Floating-point work is confined to the illustrative
torus matrix, while invariant multiplicities use rational weights and exact
integer sums.

The cubic multiplicities
$(\operatorname{Sym}^3,\mathbb S_{(2,1)},\bigwedge^3)=(0,0,1)$ are therefore
proved algebraically above.  For an optional mechanical confirmation,
construct only the zero-weight part of
$\bigwedge^3\mathfrak e_8$ and impose the eight simple-root raising operators.
The expected nullspace is one-dimensional and represented by $c(x,y,z)$; the
analogous expected nullities on $\operatorname{Sym}^3$ and
$\mathbb S_{(2,1)}$ are zero.  This sparse route avoids materializing spaces
of dimensions on the order of $\binom{248}{3}$.  It is not reported as an
executed computation in this artifact.

\paragraph{Reproducibility.} Run \texttt{marimo\_notebooks/lie\_group\_e8.py} with marimo.
\endgroup


\clearpage
\section[The adjoint SU(n) sigma-MGF]{The adjoint $SU(n)$ sigma-MGF}\label{proof:su-adjoint}
\begingroup
\definecolor{navy}{RGB}{24,55,95}
\providecommand{\Sg}{\SMG}
\renewcommand{\Sg}{\SMG}
\providecommand{\CT}{\operatorname{CT}}
\renewcommand{\CT}{\operatorname{CT}}
\setcounter{theorem}{0}
\setcounter{proposition}{0}
\setcounter{lemma}{0}
\setcounter{corollary}{0}
\setcounter{remark}{0}
\begin{abstract}
We derive the compact Molien--Weyl formula for the adjoint representation of
$SU(n)$ and test the manageable cases $n=2,3,4$.  Exact Schur-character
integration confirms the quadratic Killing form, the alternating bracket
tensor, and the symmetric cubic precisely when $n\ge3$.
\end{abstract}

\subsection{Formula}
The adjoint weights of $SU(n)$ are zero with multiplicity $n-1$ and
$e_i-e_j$ for $i\ne j$.  The general Haar/Cauchy proof yields
\[
 \boxed{\Sg_{SU(n),\mathrm{Ad}}=
 \frac1{n!}\CT_z\left[D_{A_{n-1}}(z)
 \prod_a(1-t_a)^{-(n-1)}
 \prod_{i\ne j}(1-t_a z_i/z_j)^{-1}\right].}
\]
Indeed, expanding the determinant gives
$\sum_\lambda s_\lambda(\mathbf t)\chi_{\mathbb S_\lambda\mathfrak{sl}_n}$;
Haar integration retains exactly the invariant multiplicity.

\subsection{Low-degree structure}
The Killing form supplies $h_2$, and
$\kappa(X,[Y,Z])$ supplies $e_3$ for every $n\ge2$.  For $n\ge3$ there is also
the symmetric invariant
\[
 (X,Y,Z)\longmapsto\operatorname{Tr}(XYZ+XZY),
\]
which supplies $h_3$.  For $SU(2)$ it vanishes.  Thus, through degree three,
\[
 \Sg_{SU(2),\mathrm{Ad}}=1+h_2+e_3+O_{\ge4},
\]
and, for $n\ge3$,
\[
 \Sg_{SU(n),\mathrm{Ad}}=1+h_2+h_3+e_3+O_{\ge4}.
\]
Chevalley restriction gives
$\Sg(t,0,\ldots)=\prod_{d=2}^n(1-t^d)^{-1}$.

\subsection{Direct and exact verification}
For each $n=2,3,4$, the notebook generates the $A_{n-1}$ roots, appends the
correct zero weights, and constructs the diagonal adjoint torus matrix.
It compares the determinant with the product over eigenvalues.  It then
computes every partition of degrees zero through three by Jacobi--Trudi and
applies the exact Weyl functional
\[
 \frac1{|W|}\CT(Df)=\CT\left(f\prod_{\alpha>0}(1-z^{-\alpha})\right).
\]

\begin{center}\small
\begin{tabular}{@{}c|c|rrrr@{}}\toprule
$n$&$\dim\mathfrak{sl}_n$&$[h_2]$&$[h_3]$&$[e_3]$&all other $d\le3$\\\midrule
2&3&1&0&1&0\\
3&8&1&1&1&0\\
4&15&1&1&1&0\\\bottomrule
\end{tabular}
\end{center}
All 21 partition-by-case comparisons pass.  The tested representation
dimensions $3,8,15$ are large enough to distinguish the exceptional $n=2$
behavior from the stable $n\ge3$ cubic while remaining directly inspectable.

\paragraph{Reproducibility.} Run
\path{marimo_notebooks/lie_group_su_adjoint.py}
with marimo.
\endgroup


\clearpage
\section[Schur-functor representations of U(n)]{Schur-functor representations of $U(n)$}\label{proof:schur-unitary}
\begingroup
\definecolor{ink}{HTML}{17212B}
\definecolor{accent}{HTML}{315B7D}
\definecolor{pale}{HTML}{EDF4F8}
\definecolor{passgreen}{HTML}{16794A}
\providecommand{\Sym}{\operatorname{Sym}}
\renewcommand{\Sym}{\operatorname{Sym}}
\providecommand{\Tr}{\operatorname{Tr}}
\renewcommand{\Tr}{\operatorname{Tr}}
\providecommand{\Exp}{\operatorname{Exp}_{\sigma}}
\renewcommand{\Exp}{\operatorname{Exp}_{\sigma}}
\providecommand{\summarybox}[1]{\begin{center}\fcolorbox{accent}{pale}{\parbox{0.91\linewidth}{#1}}\end{center}}
\renewcommand{\summarybox}[1]{\begin{center}\fcolorbox{accent}{pale}{\parbox{0.91\linewidth}{#1}}\end{center}}
\setcounter{theorem}{0}
\setcounter{proposition}{0}
\setcounter{lemma}{0}
\setcounter{corollary}{0}
\setcounter{remark}{0}
\summarybox{\textbf{Outcome.} The finite-dimensional plethystic trace formula is correct. For five pairs $(n,\lambda)$, including $\lambda=(2,1)$, 40 direct Schur-matrix comparisons pass with maximum error $2.98\times10^{-15}$. The important qualification is exact: for nonempty $\lambda$, the purely holomorphic $U(n)$ sigma-MGF is $1$. The nontrivial statement is the joint law with the conjugate. Three seeded Haar diagnostics also pass.}

\subsection{Statement with the unitary correction}
Let $W_n=\mathbb C^n$, let $U\in U(n)$, fix $\lambda\vdash d$, and put $V_{\lambda,n}=S_\lambda(W_n)$. If $X(U)$ is the eigenvalue alphabet of $U$, then the eigenvalue alphabet of $S_\lambda(U)$ is the plethysm $s_\lambda[X(U)]$. Hence, for every $r\ge1$,
\begin{equation}\label{unitary-schur-functor-proof-and-verification:eq:trace}
 \boxed{\Tr(S_\lambda(U)^r)=p_r[s_\lambda][X(U)]=s_\lambda[X(U^r)].}
\end{equation}
For a partition $\tau$, the coefficient integrand of the sigma-MGF is
\begin{equation}\label{unitary-schur-functor-proof-and-verification:eq:int}
 \prod_i h_{\tau_i}[s_\lambda][X(U)].
\end{equation}
Associativity of plethysm gives $h_\tau[s_\lambda[X]]=(h_\tau[s_\lambda])[X]$, so Haar integration yields
\begin{equation}
 [m_\tau]\MGF^U_{\lambda,n}
 =\dim\left(\bigotimes_i\Sym^{\tau_i}S_\lambda(W_n)\right)^{U(n)}.
\end{equation}

\textbf{Proposition (central-character obstruction).} If $d>0$, then
\begin{equation}\label{unitary-schur-functor-proof-and-verification:eq:hol}
 \boxed{\MGF^U_{\lambda,n}=1.}
\end{equation}
Indeed, $zI$ acts on $S_\lambda(W_n)$ by $z^d$ and on the tensor product above by $z^{d|\tau|}$. For $|\tau|>0$ this character is not identically one on $U(1)$, so its fixed space is zero. This is an exact finite-$n$ proof, not an asymptotic statement.

The nontrivial object uses a conjugate alphabet. In the stable range its coefficient is
\begin{equation}
 [m_\tau\bar m_\nu]\mathcal J^U_{\lambda,\infty}
 =\left\langle \Exp(h_1\bar h_1),\,h_\tau[s_\lambda]\,\overline{h_\nu[s_\lambda]}\right\rangle.
\end{equation}
This is the plethystic pushforward of the defining-representation joint law.

\textbf{Stable-range qualification.}  Let $|\lambda|=d$ and consider a
balanced coefficient indexed by outer partitions $\alpha,\beta$, with
$d|\alpha|=d|\beta|=:m$.  Embed its coefficient space into
$W_n^{\otimes m}\otimes(W_n^*)^{\otimes m}$.  The first fundamental theorem
for $GL_n$ says that contraction tensors indexed by $S_m$ span the invariants,
and the second fundamental theorem says that they are linearly independent
when $n\geq m$.  Therefore the stable plethystic formula is valid whenever
\[
n\geq m=d|\alpha|=d|\beta|.
\]
If the two degrees are unequal, the coefficient is exactly zero for every $n$
by the scalar-center argument.

\subsection{Explicit polynomial-Gaussian law}
The Frobenius formula gives
\begin{equation}
 s_\lambda=\sum_{\alpha\vdash d}\frac{\chi^\lambda(\alpha)}{z_\alpha}p_\alpha,
 \qquad z_\alpha=\prod_j j^{m_j(\alpha)}m_j(\alpha)!.
\end{equation}
Substituting $p_k[X(U^r)]=\Tr(U^{rk})$ in \eqref{unitary-schur-functor-proof-and-verification:eq:trace} gives the exact identity
\begin{equation}\label{unitary-schur-functor-proof-and-verification:eq:frob}
 \Tr(S_\lambda(U)^r)=\sum_{\alpha\vdash d}\frac{\chi^\lambda(\alpha)}{z_\alpha}
 \prod_{a\in\alpha}\Tr(U^{ra}).
\end{equation}
For fixed powers, $\Tr(U^k)$ converges jointly in moments to independent circular complex Gaussians $Z_k$ satisfying $\mathbb E Z_k=\mathbb E Z_k^2=0$ and $\mathbb E|Z_k|^2=k$. Polynomial substitution in \eqref{unitary-schur-functor-proof-and-verification:eq:frob} therefore gives
\begin{equation}
 \Tr(S_\lambda(U)^r)\longrightarrow
 Q^U_{\lambda,r}=\sum_{\alpha\vdash d}\frac{\chi^\lambda(\alpha)}{z_\alpha}\prod_{a\in\alpha}Z_{ra}.
\end{equation}
For example, $Q^U_{(2),1}=\tfrac12(Z_1^2+Z_2)$ and $Q^U_{(1,1),1}=\tfrac12(Z_1^2-Z_2)$. Pure holomorphic moments vanish, while mixed moments generally do not.

\subsection{Independent finite-dimensional verification}
For each test the code draws a seeded Haar unitary by complex Gaussian QR factorization. It constructs a standard Young symmetrizer $c_\lambda$ on $W_n^{\otimes d}$, obtains an orthonormal basis $B$ of its image by singular-value decomposition, and forms the actual matrix
\begin{equation}
 D_\lambda(U)=B^*U^{\otimes d}B.
\end{equation}
This route does not use symmetric-function trace identities. The comparison route computes $p_k=\Tr(U^{rk})$, obtains $h_j$ from Newton's recurrence
\begin{equation}
 jh_j=\sum_{k=1}^j p_kh_{j-k},\qquad h_0=1,
\end{equation}
and evaluates $s_\lambda=\det(h_{\lambda_i-i+j})$ by Jacobi--Trudi. It compares both \eqref{unitary-schur-functor-proof-and-verification:eq:trace} for $r=1,2,3$ and \eqref{unitary-schur-functor-proof-and-verification:eq:int} for $\tau=(),(1),(2),(1,1),(3)$.

\begin{center}\begin{tabular}{@{}llllr@{}}\toprule
$n$ & tested $\lambda$ & direct checks & largest purpose & status\\\midrule
2 & $(1),(2)$ & trace and sigma integrands & low-rank boundary & PASS\\
3 & $(1,1),(3),(2,1)$ & trace and sigma integrands & mixed Young symmetry & PASS\\\midrule
\multicolumn{4}{r}{40 comparisons; maximum absolute error} & $2.98\times10^{-15}$\\\bottomrule
\end{tabular}\end{center}

\clearpage
The seeded $U(8)$ Haar run uses 1,200 matrices. It estimates $\mathbb E\chi_{(2)}=0.0222+0.0343i$, $\mathbb E\chi_{(2)}^2=0.0322-0.0404i$, and $\mathbb E\chi_{(1,1)}=-0.0059+0.0148i$, all within $4.5$ estimated standard errors of the exact target zero. These statistical rows check the implementation but do not replace the central-character proof.

\subsection{Scope and reproducibility}
The exact trace formula holds whenever $S_\lambda(\mathbb C^n)$ is defined ($\ell(\lambda)\le n$). Coefficient stabilization is a separate assertion: the balanced unitary threshold is $n\geq d|\alpha|=d|\beta|$, while the orthogonal/symplectic Brauer bound is stated in their respective notes.  Unitary holomorphic vanishing is exact for every allowed $n$. Run the verification module below from the repository root, or open \path{marimo_notebooks/schur_unitary.py} with marimo.
\begin{center}\footnotesize\texttt{python -m for\_this\_guy.schur\_functor\_classical\_groups.unitary.verification}\end{center}

\subsection*{References}
S. Howe, \href{https://arxiv.org/abs/2412.19295}{\emph{Random matrix statistics and zeroes of $L$-functions via probability in $\lambda$-rings}}.\\
G. Lehrer and R. Zhang, \href{https://arxiv.org/abs/1207.5889}{\emph{The Brauer Category and Invariant Theory}}.
\endgroup


\clearpage
\section[Schur-functor representations of O(n)]{Schur-functor representations of $O(n)$}\label{proof:schur-orthogonal}
\begingroup
\definecolor{ink}{HTML}{17212B}
\definecolor{accent}{HTML}{315B7D}
\definecolor{pale}{HTML}{EDF4F8}
\providecommand{\Sym}{\operatorname{Sym}}
\renewcommand{\Sym}{\operatorname{Sym}}
\providecommand{\Tr}{\operatorname{Tr}}
\renewcommand{\Tr}{\operatorname{Tr}}
\providecommand{\Exp}{\operatorname{Exp}_{\sigma}}
\renewcommand{\Exp}{\operatorname{Exp}_{\sigma}}
\providecommand{\summarybox}[1]{\begin{center}\fcolorbox{accent}{pale}{\parbox{0.91\linewidth}{#1}}\end{center}}
\renewcommand{\summarybox}[1]{\begin{center}\fcolorbox{accent}{pale}{\parbox{0.91\linewidth}{#1}}\end{center}}
\setcounter{theorem}{0}
\setcounter{proposition}{0}
\setcounter{lemma}{0}
\setcounter{corollary}{0}
\setcounter{remark}{0}
\summarybox{\textbf{Outcome.} The orthogonal formula is correct as a coefficientwise stable plethystic pushforward. The finite trace identity is exact before taking a limit. Five representative $(n,\lambda)$ cases give 40 direct comparisons with maximum error $5.33\times10^{-15}$. Haar $O(9)$ tests recover the invariant quadratic form in $\Sym^2W$, its absence from $\Lambda^2W$, and the odd-degree parity obstruction.}

\subsection{Stable sigma-MGF and its proof}
Let $W_n=\mathbb C^n$ with the defining action of the compact group $O(n)$, fix $\lambda\vdash d$, and put $V_{\lambda,n}=S_\lambda(W_n)$. For a partition $\tau$,
\begin{equation}\label{orthogonal-schur-functor-proof-and-verification:eq:coeff}
 [m_\tau]\MGF^O_{\lambda,n}
 =\dim\left(\bigotimes_i\Sym^{\tau_i}S_\lambda(W_n)\right)^{O(n)}
 =\mu_{O(n)}\bigl(h_\tau[s_\lambda]\bigr).
\end{equation}
The second equality is simply associativity of plethysm:
$h_\tau[s_\lambda[X]]=(h_\tau[s_\lambda])[X]$.
The defining-representation stable sigma-MGF is $\Exp(h_2)$. Applying its coefficient functional to the fixed symmetric function $h_\tau[s_\lambda]$ proves
\begin{equation}\label{orthogonal-schur-functor-proof-and-verification:eq:push}
 \boxed{[m_\tau]\MGF^O_{\lambda,\infty}
 =\left\langle\Exp(h_2),h_\tau[s_\lambda]\right\rangle.}
\end{equation}
No new limiting invariant theory is needed after the defining case. A conservative sufficient stable range is $n\ge d|\tau|$: the relevant expression has tensor degree $d|\tau|$, where Brauer pairings are independent in that range.

Since $-I\in O(n)$ acts on $S_\lambda(W_n)$ by $(-1)^d$, it acts on the space in \eqref{orthogonal-schur-functor-proof-and-verification:eq:coeff} by $(-1)^{d|\tau|}$. Therefore
\begin{equation}\label{orthogonal-schur-functor-proof-and-verification:eq:parity}
 \boxed{[m_\tau]\MGF^O_{\lambda,n}=0\quad\text{if }d|\tau|\text{ is odd}.}
\end{equation}
This parity conclusion is exact at finite $n$.

\subsection{Exact traces and the limiting law}
For every $g\in O(n)$ and $r\ge1$, functoriality gives $S_\lambda(g)^r=S_\lambda(g^r)$, whence
\begin{equation}\label{orthogonal-schur-functor-proof-and-verification:eq:trace}
 \boxed{\Tr(S_\lambda(g)^r)=s_\lambda[X(g^r)]
 =\sum_{\alpha\vdash d}\frac{\chi^\lambda(\alpha)}{z_\alpha}
 \prod_{a\in\alpha}\Tr(g^{ra}).}
\end{equation}
This formula is exact for every allowed $n$; stability is needed only when replacing traces by Gaussian variables.

The joint moment convergence of finitely many orthogonal trace powers is an
external classical-group trace theorem (Diaconis--Evans, with the
normalization below).  Let independent $A_k\sim N(0,1)$ and set
\begin{equation}
 Z_k^O=\sqrt{k}\,A_k+\mathbf1_{2\mid k}.
\end{equation}
Joint moment convergence of fixed orthogonal trace powers and polynomial substitution in \eqref{orthogonal-schur-functor-proof-and-verification:eq:trace} give
\begin{equation}
 \Tr(S_\lambda(g)^r)\longrightarrow Q^O_{\lambda,r}
 =\sum_{\alpha\vdash d}\frac{\chi^\lambda(\alpha)}{z_\alpha}\prod_{a\in\alpha}Z_{ra}^O.
\end{equation}
For $\lambda=(2)$, $Q^O_{(2),1}=\tfrac12((Z_1^O)^2+Z_2^O)$ has mean $1$, recording the invariant quadratic form. For $\lambda=(1,1)$, the minus sign makes the mean zero.

\subsection{Matrix construction and formula cross-check}
The verification draws seeded Haar orthogonal matrices by real Gaussian QR with sign correction. A standard Young symmetrizer $c_\lambda$ is built explicitly on $W_n^{\otimes d}$. If $B$ is an orthonormal basis of $\operatorname{im}c_\lambda$, the direct representation matrix is
\begin{equation}
 D_\lambda(g)=B^*g^{\otimes d}B.
\end{equation}
The formula cross-check never uses this matrix: it computes the defining traces, uses Newton's recurrence $jh_j=\sum_{k=1}^jp_kh_{j-k}$, and evaluates the Jacobi--Trudi determinant $s_\lambda=\det(h_{\lambda_i-i+j})$. Tests compare $\Tr(D_\lambda(g)^r)$ for $r=1,2,3$. They also compare
\begin{equation}
 \prod_i h_{\tau_i}[D_\lambda(g)]
 \quad\text{against}\quad
 \prod_i h_{\tau_i}[s_\lambda][X(g)]
\end{equation}
for $\tau=(),(1),(2),(1,1),(3)$, so the actual sigma-MGF integrands are covered.

\begin{center}\begin{tabular}{@{}llllr@{}}\toprule
$n$ & tested $\lambda$ & checks & purpose & status\\\midrule
3 & $(1),(2),(3)$ & trace and sigma integrands & defining/symmetric powers & PASS\\
4 & $(1,1),(2,1)$ & trace and sigma integrands & alternating/mixed symmetry & PASS\\\midrule
\multicolumn{4}{r}{40 comparisons; maximum absolute error} & $5.33\times10^{-15}$\\\bottomrule
\end{tabular}\end{center}

The seeded Haar diagnostic uses 1,200 matrices in $O(9)$. It gives means $0.9979$ for $\chi_{(2)}$ (target $1$), $-0.0289$ for $\chi_{(1,1)}$ (target $0$), and $0.0116$ for $\chi_{(3)}$ (target $0$). Each lies inside $4.5$ estimated standard errors. Group residuals are at most $3.24\times10^{-15}$.

\subsection{Reproducibility and interpretation}
Run the module below from the repository root, or open \path{marimo_notebooks/schur_orthogonal.py} with marimo. Exact rows are the proof-driven regression test. Haar rows are seeded numerical diagnostics and are reported with standard errors; they are not presented as proofs of convergence.
\begin{center}\footnotesize\texttt{python -m for\_this\_guy.schur\_functor\_classical\_groups.orthogonal.verification}\end{center}

\subsection*{References}
S. Howe, \href{https://arxiv.org/abs/2412.19295}{\emph{Random matrix statistics and zeroes of $L$-functions via probability in $\lambda$-rings}}.\\
G. Lehrer and R. Zhang, \href{https://arxiv.org/abs/1207.5889}{\emph{The Brauer Category and Invariant Theory}}.
\\P. Diaconis and S. N. Evans, \emph{Linear functionals of eigenvalues of
random matrices}, Trans. Amer. Math. Soc. 353 (2001), 2615--2633.
\endgroup


\clearpage
\section[Schur-functor representations of USp(2n)]{Schur-functor representations of $USp(2n)$}\label{proof:schur-symplectic}
\begingroup
\definecolor{ink}{HTML}{17212B}
\definecolor{accent}{HTML}{315B7D}
\definecolor{pale}{HTML}{EDF4F8}
\providecommand{\Sym}{\operatorname{Sym}}
\renewcommand{\Sym}{\operatorname{Sym}}
\providecommand{\Tr}{\operatorname{Tr}}
\renewcommand{\Tr}{\operatorname{Tr}}
\providecommand{\Exp}{\operatorname{Exp}_{\sigma}}
\renewcommand{\Exp}{\operatorname{Exp}_{\sigma}}
\providecommand{\summarybox}[1]{\begin{center}\fcolorbox{accent}{pale}{\parbox{0.91\linewidth}{#1}}\end{center}}
\renewcommand{\summarybox}[1]{\begin{center}\fcolorbox{accent}{pale}{\parbox{0.91\linewidth}{#1}}\end{center}}
\setcounter{theorem}{0}
\setcounter{proposition}{0}
\setcounter{lemma}{0}
\setcounter{corollary}{0}
\setcounter{remark}{0}
\summarybox{\textbf{Outcome.} Interpreting $Sp(2n)$ as the compact group $USp(2n)$, the stable formula is the plethystic pushforward of $\Exp(e_2)$. Five representative cases in ambient dimensions 2 and 4 give 40 direct comparisons with maximum error $3.70\times10^{-15}$. A separate Haar $USp(8)$ run detects the invariant alternating form in $\Lambda^2W$ and verifies odd-degree parity.}

\subsection{Convention and stable sigma-MGF}
Here $Sp(2n)$ means the compact matrix group
\begin{equation}
 USp(2n)=\{g\in U(2n):g^TJg=J\},\qquad
 J=\begin{pmatrix}0&I_n\\-I_n&0\end{pmatrix}.
\end{equation}
This convention matters: the complex algebraic group is not a Haar probability space, while its compact form is.

Let $W_n=\mathbb C^{2n}$, fix $\lambda\vdash d$, and put $V_{\lambda,n}=S_\lambda(W_n)$. For each partition $\tau$,
\begin{equation}\label{symplectic-schur-functor-proof-and-verification:eq:coeff}
 [m_\tau]\MGF^{Sp}_{\lambda,n}
 =\dim\left(\bigotimes_i\Sym^{\tau_i}S_\lambda(W_n)\right)^{USp(2n)}
 =\mu_{Sp(2n)}\bigl(h_\tau[s_\lambda]\bigr).
\end{equation}
Associativity of plethysm proves the last equality. Since the defining-representation stable sigma-MGF is $\Exp(e_2)$, applying its coefficient functional to the fixed function $h_\tau[s_\lambda]$ proves
\begin{equation}\label{symplectic-schur-functor-proof-and-verification:eq:push}
 \boxed{[m_\tau]\MGF^{Sp}_{\lambda,\infty}
 =\left\langle\Exp(e_2),h_\tau[s_\lambda]\right\rangle.}
\end{equation}
A conservative sufficient stable range is $\dim W_n=2n\ge d|\tau|$, obtained by embedding in tensor degree $d|\tau|$ and applying symplectic Brauer invariant theory.

The central element $-I$ acts by $(-1)^d$ on $S_\lambda(W_n)$. Thus the exact finite-dimensional parity rule is
\begin{equation}\label{symplectic-schur-functor-proof-and-verification:eq:parity}
 \boxed{[m_\tau]\MGF^{Sp}_{\lambda,n}=0\quad\text{if }d|\tau|\text{ is odd}.}
\end{equation}

\subsection{Exact trace formula and Gaussian substitution}
Functoriality and the Frobenius expansion of $s_\lambda$ give, for every $g\in USp(2n)$,
\begin{equation}\label{symplectic-schur-functor-proof-and-verification:eq:trace}
 \boxed{\Tr(S_\lambda(g)^r)=s_\lambda[X(g^r)]
 =\sum_{\alpha\vdash d}\frac{\chi^\lambda(\alpha)}{z_\alpha}
 \prod_{a\in\alpha}\Tr(g^{ra}).}
\end{equation}
This is exact whenever $\ell(\lambda)\le2n$.

The joint moment convergence of finitely many symplectic trace powers is an
external classical-group trace theorem (Diaconis--Evans, with the
normalization below).  Let independent $A_k\sim N(0,1)$ and set
\begin{equation}
 Z_k^{Sp}=\sqrt{k}\,A_k-\mathbf1_{2\mid k}.
\end{equation}
Joint moment convergence of fixed trace powers, followed by polynomial substitution in \eqref{symplectic-schur-functor-proof-and-verification:eq:trace}, gives
\begin{equation}
 \Tr(S_\lambda(g)^r)\longrightarrow Q^{Sp}_{\lambda,r}
 =\sum_{\alpha\vdash d}\frac{\chi^\lambda(\alpha)}{z_\alpha}\prod_{a\in\alpha}Z_{ra}^{Sp}.
\end{equation}
For $\lambda=(1,1)$,
$Q^{Sp}_{(1,1),1}=\tfrac12((Z_1^{Sp})^2-Z_2^{Sp})$ has mean $1$, representing the invariant symplectic form. For $\lambda=(2)$ the corresponding plus sign yields mean zero.

\subsection{Finite matrix construction and formula cross-check}
The sampler performs quaternionic Gram--Schmidt. If $v$ is accepted, it also inserts $-J\bar v$; with columns ordered in pairs of blocks this gives both $g^*g=I$ and $g^TJg=J$. The largest observed group residual in all reported sampling is $6.67\times10^{-15}$.

For the direct route, the code constructs a Young symmetrizer $c_\lambda$ on $W_n^{\otimes d}$, computes an orthonormal basis $B$ of its image, and forms
\begin{equation}
 D_\lambda(g)=B^*g^{\otimes d}B.
\end{equation}
For the formula route, only $\Tr(g^k)$ is retained. Newton's recurrence constructs the $h_j$, and Jacobi--Trudi constructs $s_\lambda$. The comparison covers $\Tr(D_\lambda(g)^r)$ for $r=1,2,3$ and the sigma-MGF integrands $\prod_i h_{\tau_i}[D_\lambda(g)]$ for $\tau=(),(1),(2),(1,1),(3)$.

\begin{center}\small\begin{tabular}{@{}llllr@{}}\toprule
ambient dimension & tested $\lambda$ & checks & purpose & status\\\midrule
2 & $(1)$ & trace and sigma integrands & rank-one boundary & PASS\\
4 & $(2),(1,1),(3),(2,1)$ & trace and sigma integrands & all degree $\le3$ types & PASS\\\midrule
\multicolumn{4}{r}{40 comparisons; maximum absolute error} & $3.70\times10^{-15}$\\\bottomrule
\end{tabular}\end{center}

The seeded Haar diagnostic uses 1,200 matrices in $USp(8)$. It estimates means $0.0031$ for $\chi_{(2)}$ (target $0$), $0.9917$ for $\chi_{(1,1)}$ (target $1$), and $0.0818$ for $\chi_{(3)}$ (target $0$). All lie within $4.5$ estimated standard errors. The last target also exercises \eqref{symplectic-schur-functor-proof-and-verification:eq:parity}.

\subsection{Reproducibility and interpretation}
Run the module below from the repository root, or open \path{marimo_notebooks/schur_symplectic.py} with marimo. The exact matrix/character identities certify the finite trace formula.  The seeded Haar rows are diagnostics only; they neither prove Haar integration identities nor the $n\to\infty$ convergence theorem.
\begin{center}\footnotesize\texttt{python -m for\_this\_guy.schur\_functor\_classical\_groups.symplectic.verification}\end{center}

\subsection*{References}
S. Howe, \href{https://arxiv.org/abs/2412.19295}{\emph{Random matrix statistics and zeroes of $L$-functions via probability in $\lambda$-rings}}.\\
G. Lehrer and R. Zhang, \href{https://arxiv.org/abs/1207.5889}{\emph{The Brauer Category and Invariant Theory}}.
\\P. Diaconis and S. N. Evans, \emph{Linear functionals of eigenvalues of
random matrices}, Trans. Amer. Math. Soc. 353 (2001), 2615--2633.
\endgroup


\clearpage
\section{Balanced sigma-MGFs for local unitary tensor products}\label{proof:local-unitary}
\begingroup
\definecolor{ink}{HTML}{17212B}
\definecolor{accent}{HTML}{315B7D}
\definecolor{pale}{HTML}{EDF4F8}
\definecolor{passgreen}{HTML}{16794A}
\providecommand{\Sym}{\operatorname{Sym}}
\renewcommand{\Sym}{\operatorname{Sym}}
\providecommand{\End}{\operatorname{End}}
\renewcommand{\End}{\operatorname{End}}
\providecommand{\Tr}{\operatorname{Tr}}
\renewcommand{\Tr}{\operatorname{Tr}}
\providecommand{\Hom}{\operatorname{Hom}}
\renewcommand{\Hom}{\operatorname{Hom}}
\providecommand{\Exp}{\operatorname{Exp}_{\sigma}}
\renewcommand{\Exp}{\operatorname{Exp}_{\sigma}}
\providecommand{\SMG}{\operatorname{SMG}}
\renewcommand{\SMG}{\operatorname{SMG}}
\providecommand{\PASS}{\textcolor{passgreen}{\textbf{PASS}}}
\renewcommand{\PASS}{\textcolor{passgreen}{\textbf{PASS}}}
\providecommand{\summarybox}[1]{\begin{center}\fcolorbox{accent}{pale}{\parbox{0.91\linewidth}{#1}}\end{center}}
\renewcommand{\summarybox}[1]{\begin{center}\fcolorbox{accent}{pale}{\parbox{0.91\linewidth}{#1}}\end{center}}
\setcounter{theorem}{0}
\setcounter{proposition}{0}
\setcounter{lemma}{0}
\setcounter{corollary}{0}
\setcounter{remark}{0}
\summarybox{\textbf{Outcome.} For
$G=\prod_{a=1}^kU(d_a)$ on $V=\bigotimes_{a=1}^k\mathbb C^{d_a}$,
the ordinary one-sided sigma-MGF is exactly $1$. The nontrivial object is the
balanced two-sided series involving both $V$ and $V^*$. Its stable power-sum
coefficient is $z_\lambda^{k-2}$, its bipartite single-state coefficient is a
restricted partition number, and its multilinear coefficient is an exact
product of Schur--Weyl centralizer dimensions. A reproducible marimo notebook
constructs the Kronecker-product matrices and compares direct Haar estimates
with these formulas. All 46 numerical comparisons, 4 construction checks, and
5 scalar-center witnesses pass.}

\subsection{Representation and the correction at the outset}

Put $V_a=\mathbb C^{d_a}$ and
\[
 G=U(d_1)\times\cdots\times U(d_k),\qquad
 V=V_1\otimes\cdots\otimes V_k.
\]
The representation is the concrete matrix homomorphism
\begin{equation}
 \rho(g_1,\ldots,g_k)=g_1\otimes\cdots\otimes g_k
 \in U(D),\qquad D=\prod_{a=1}^kd_a.
\end{equation}
For a partition $\tau=(\tau_1,\ldots,\tau_\ell)$, write
\[
 \Sym^\tau V=\bigotimes_{j=1}^{\ell}\Sym^{\tau_j}V,
 \qquad |\tau|=\sum_j\tau_j.
\]
The ordinary coefficient is
\[
 [m_\tau]\SMG_{G,V}=\dim(\Sym^\tau V)^G.
\]

\textbf{Proposition 1 (one-sided collapse).}
\begin{equation}
 \boxed{\SMG_{G,V}=1.}
\end{equation}

\textit{Proof.}
For $z\in U(1)$ the group contains
$(zI_{d_1},I_{d_2},\ldots,I_{d_k})$, whose matrix in (1) is $zI_V$.
It acts on $\Sym^\tau V$ as $z^{|\tau|}I$. If $|\tau|>0$, one can choose
$z$ with $z^{|\tau|}\ne1$, so a fixed vector must be zero. Only degree zero
remains. \hfill$\square$

The same center explains why local-unitary entanglement invariants use a state
$\psi$ and its conjugate $\bar\psi$: the weights $z$ and $z^{-1}$ cancel only
when the two total degrees agree.

\subsection{The balanced series and Molien formula}

Introduce independent alphabets $\mathbf s$ and $\mathbf t$ and define
\begin{equation}
 \mathcal B_{G,V}(\mathbf s,\mathbf t)
 =\int_G \Exp\!\left(X_V(g)h_1(\mathbf s)\right)
          \Exp\!\left(X_{V^*}(g)h_1(\mathbf t)\right)\,dg.
\end{equation}
Coefficient extraction gives
\begin{align}
 \mathcal B_{G,V}
 &=\sum_{\alpha,\beta}b_{\alpha,\beta}
          m_\alpha(\mathbf s)m_\beta(\mathbf t),\\
 b_{\alpha,\beta}
 &=\dim\left(\Sym^\alpha V\otimes\Sym^\beta V^*\right)^G
  =\dim\Hom_G(\Sym^\beta V,\Sym^\alpha V).
\end{align}
The center now gives the balanced selection rule
\begin{equation}
 \boxed{b_{\alpha,\beta}=0\quad\text{unless}\quad
 |\alpha|=|\beta|.}
\end{equation}

For finitely many variables, (3) is the balanced Molien integral
\begin{equation}
 \mathcal B_{G,V}
 =\int_G\prod_i\det(I-s_i\rho(g))^{-1}
         \prod_j\det(I-t_j\rho(g)^{-1})^{-1}\,dg.
\end{equation}
If $x^{(a)}_1,\ldots,x^{(a)}_{d_a}$ are the eigenvalues of $g_a$, then the
eigenvalues of $\rho(g)$ are $x^{(1)}_{i_1}\cdots x^{(k)}_{i_k}$.
Consequently (7), together with the Weyl integration formula for each local
factor, is an exact constant-term prescription in every dimension.

\subsection{Stable power-sum formula}

For every positive integer $r$, Kronecker products satisfy the directly
testable identity
\begin{equation}
 \Tr(\rho(g)^r)=\prod_{a=1}^k\Tr(g_a^r).
\end{equation}
Thus, for a partition $\lambda$,
$p_\lambda(\rho(g))=\prod_a p_\lambda(g_a)$, and product Haar measure gives
\begin{equation}
 \mathbb E_G\!\left[p_\lambda(\rho(g))
 \overline{p_\mu(\rho(g))}\right]
 =\prod_{a=1}^k\mathbb E_{U(d_a)}
 \left[p_\lambda(g_a)\overline{p_\mu(g_a)}\right].
\end{equation}
Let $z_\lambda=\prod_i i^{m_i(\lambda)}m_i(\lambda)!$. Schur orthogonality,
or equivalently the stable Hall inner product, gives
\begin{equation}
 \mathbb E_{U(d)}[p_\lambda(U)\overline{p_\mu(U)}]
 =\delta_{\lambda\mu}z_\lambda
 \quad\text{when }d\ge |\lambda|,|\mu|.
\end{equation}
Combining (9) and (10), if $N\le\min_a d_a$ then
\begin{equation}
 \mathbb E_G[p_\lambda(\rho)\overline{p_\mu(\rho)}]
 =\delta_{\lambda\mu}z_\lambda^k
 \quad (|\lambda|,|\mu|\le N).
\end{equation}
Each determinant exponential contributes a factor $p_\lambda/z_\lambda$.
Therefore through bidegree $(N,N)$,
\begin{equation}
 \boxed{\mathcal B_{G,V}\equiv
 \sum_{|\lambda|\le N}z_\lambda^{k-2}
 p_\lambda(\mathbf s)p_\lambda(\mathbf t).}
\end{equation}
For $k=1,2,3$, the coefficients are respectively
$z_\lambda^{-1},1,z_\lambda$.

\subsection{Exact single-state coefficients}

For partitions $\lambda^{(1)},\ldots,\lambda^{(k)}\vdash m$, define the
generalized Kronecker coefficient
\[
 g(\lambda^{(1)},\ldots,\lambda^{(k)})
 =\dim\left([\lambda^{(1)}]\otimes\cdots\otimes[\lambda^{(k)}]\right)^{S_m}.
\]
Applying Schur--Weyl duality to each $V_a^{\otimes m}$ and then taking
diagonal $S_m$-invariants proves
\begin{equation}
 \Sym^m(V_1\otimes\cdots\otimes V_k)
 \cong\bigoplus_{\lambda^{(1)},\ldots,\lambda^{(k)}\vdash m}
 g(\lambda^{(1)},\ldots,\lambda^{(k)})
 \bigotimes_{a=1}^kS_{\lambda^{(a)}}V_a,
\end{equation}
where terms with $\ell(\lambda^{(a)})>d_a$ vanish. Schur orthogonality now
gives the all-dimensional identity
\begin{equation}
 \boxed{b_{(m),(m)}=
 \sum_{\substack{\lambda^{(a)}\vdash m\\
                  \ell(\lambda^{(a)})\le d_a}}
 g(\lambda^{(1)},\ldots,\lambda^{(k)})^2.}
\end{equation}

For $k=2$, $g(\lambda,\mu)=\delta_{\lambda\mu}$. Hence
\begin{equation}
 \boxed{b_{(m),(m)}=
 \#\{\lambda\vdash m:\ell(\lambda)\le\min(d_1,d_2)\}.}
\end{equation}
This is $p(m)$ when both dimensions are at least $m$. In stable range, setting
each alphabet in (12) to one variable yields the equivalent formula
\begin{equation}
 \boxed{b_{(m),(m)}=\sum_{\lambda\vdash m}z_\lambda^{k-2}.}
\end{equation}

For completeness, general $b_{\alpha,\beta}$ also has an exact finite formula.
Decompose each $\Sym^{\alpha_j}V$ by (13), then combine local Schur functors
with Littlewood--Richardson coefficients. If the resulting multiplicity of
$\bigotimes_aS_{\mu^{(a)}}V_a$ is $M_\alpha(\mu^{(1)},\ldots,\mu^{(k)})$,
then
\begin{equation}
 b_{\alpha,\beta}=\sum_{\substack{\mu^{(1)},\ldots,\mu^{(k)}\\
                          \ell(\mu^{(a)})\le d_a}}
 M_\alpha(\mu^{(1)},\ldots,\mu^{(k)})
 M_\beta(\mu^{(1)},\ldots,\mu^{(k)}).
\end{equation}
This separates the two combinatorial ingredients: generalized Kronecker
coefficients enter the symmetric powers of a tensor product, while
Littlewood--Richardson coefficients enter products of symmetric powers.

\subsection{Exact multilinear coefficients}

For $\alpha=\beta=(1^r)$, $\Sym^{(1^r)}V=V^{\otimes r}$, so
\begin{equation}
 b_{(1^r),(1^r)}=\dim\End_G(V^{\otimes r}).
\end{equation}
Because the local factors act independently,
\[
 \End_G(V^{\otimes r})\cong
 \bigotimes_{a=1}^k\End_{U(d_a)}(V_a^{\otimes r}).
\]
If $f^\lambda$ is the number of standard Young tableaux of shape $\lambda$,
Schur--Weyl duality and the hook-length formula give
\begin{equation}
 \boxed{b_{(1^r),(1^r)}=
 \prod_{a=1}^k\sum_{\substack{\lambda\vdash r\\
                    \ell(\lambda)\le d_a}}(f^\lambda)^2.}
\end{equation}
When every $d_a\ge r$, the inner sum is $|S_r|=r!$, and therefore
\begin{equation}
 \boxed{b_{(1^r),(1^r)}=(r!)^k.}
\end{equation}
This is the dimension of the stable family of contraction diagrams indexed by
$k$ independent permutations in $S_r$.

\subsection{Verification design}

The accompanying implementation uses two independent routes.

\textbf{Exact route.}
Integer partitions are generated recursively. The hook-length formula computes
$f^\lambda$, which is substituted into (19). Equation (15) is evaluated by
restricted partition counting, and (16) by direct computation of $z_\lambda$.
No sampled matrices enter these tables.

\textbf{Direct matrix route.}
For every local factor, a complex Gaussian matrix $Z$ is QR-factorized and the
diagonal phases of $R$ are removed. The resulting $Q$ is Haar distributed on
$U(d_a)$. The code explicitly forms
$\rho=Q_1\otimes\cdots\otimes Q_k$, checks $\rho^*\rho=I$, and evaluates its
power traces. Newton's recurrence
\begin{equation}
 mh_m(\rho)=\sum_{j=1}^m\Tr(\rho^j)h_{m-j}(\rho)
\end{equation}
then produces characters of $\Sym^mV$ directly. The sampled target quantities
are
\begin{align*}
 &\mathbb E\prod_jh_{\tau_j}(\rho) &&\text{(ordinary one-sided)},\\
 &\mathbb E[p_\lambda(\rho)\overline{p_\mu(\rho)}]
     &&\text{(balanced power-sum)},\\
 &\mathbb E|h_m(\rho)|^2 &&\text{(single-state)},\\
 &\mathbb E|\Tr\rho|^{2r} &&\text{(multilinear)}.
\end{align*}
The last two are character inner products, hence exactly the invariant
multiplicities in (14) and (19). Numerical integration is corroboration, not a
replacement for the proofs above.

\clearpage
\subsection{Results}

Each case used 6,000 deterministic-seed Haar samples. A numerical row passes
when the absolute error is at most six empirical standard errors plus a small
fixed rounding margin. Matrix unitarity and trace factorization were also
checked at tolerance $10^{-11}$.

\begin{center}
\begin{tabularx}{\linewidth}{@{}Xlrrrr@{}}
\toprule
Representation & Target & Direct estimate & Formula & Std. error & Result\\
\midrule
$U(3)$ & $|h_3|^2$ & 1.0367 & 1 & 0.0253 & \PASS\\
$U(3)$ & $|\Tr\rho|^4$ & 2.0318 & 2 & 0.0593 & \PASS\\
$U(2)\times U(3)$ & $|h_2|^2$ & 1.9533 & 2 & 0.0546 & \PASS\\
$U(2)\times U(3)$ & $|p_{(2)}|^2$ & 3.8911 & 4 & 0.0638 & \PASS\\
$U(2)\times U(2)$ & $|h_3|^2$ & 2.0797 & 2 & 0.1219 & \PASS\\
$U(2)^3$ & $|h_2|^2$ & 4.0815 & 4 & 0.1641 & \PASS\\
$U(2)^3$ & $|\Tr\rho|^4$ & 8.2364 & 8 & 0.6379 & \PASS\\
$U(2)^3$ & $p_{(2)}\overline{p_{(1,1)}}$ & $0.0032$ & 0 & 0.0284 & \PASS\\
\bottomrule
\end{tabularx}
\end{center}

The exact calculations make the stable-range boundary visible:
\begin{center}
\begin{tabular}{@{}llrrl@{}}
\toprule
Dimensions & Sector & Degree & Exact & Stable expression\\
\midrule
$(2,2)$ & single-state & 3 & 2 & $p(3)=3$ (not applicable)\\
$(2,2)$ & multilinear & 3 & 25 & $(3!)^2=36$ (not applicable)\\
$(2,3)$ & multilinear & 3 & 30 & $(3!)^2=36$ (not applicable)\\
$(2,2,2)$ & multilinear & 3 & 125 & $(3!)^3=216$ (not applicable)\\
\bottomrule
\end{tabular}
\end{center}
Thus the dimension restrictions are not cosmetic: applying a stable formula
outside $r,m\le\min_a d_a$ produces wrong answers.

The scalar-center test separately constructs, for each degree $1\le n\le5$,
$z=e^{2\pi i/(n+1)}$ and the group matrix
$\rho(zI_2,I_3)=zI_6$. Its action on degree $n$ is $z^n\ne1$ in every case,
giving five direct matrix witnesses for Proposition~1.

\subsection{Reproduction and scope}

The folder \texttt{for\_this\_guy/local\_unitary\_tensor\_products} contains:
\begin{itemize}[leftmargin=1.5em]
 \item \path{marimo_notebooks/local_unitary_tensor_products.py}, the interactive marimo report;
 \item \texttt{verification.py}, the direct and exact implementation;
 \item \texttt{test\_verification.py}, deterministic regression tests;
 \item this \LaTeX{} source and its compiled PDF.
\end{itemize}
The notebook exposes the sample count and reports every tested coefficient, not
only the representative rows printed here. Its global $U(3)$ case is explicitly
labeled as a comparator; it is not conflated with the local tensor-product
groups. All local-unitary cases remain grouped in this single family because
they are specializations of the same representation (1) and the same proof.
\endgroup


\clearpage
\section{Unitary conjugation on endomorphisms}\label{proof:unitary-conjugation}
\begingroup
\definecolor{ink}{HTML}{18212A}
\definecolor{accent}{HTML}{245A9C}
\definecolor{soft}{HTML}{EEF5FB}
\providecommand{\C}{\mathbb C}
\renewcommand{\C}{\mathbb C}
\providecommand{\End}{\operatorname{End}}
\renewcommand{\End}{\operatorname{End}}
\providecommand{\Sym}{\operatorname{Sym}}
\renewcommand{\Sym}{\operatorname{Sym}}
\providecommand{\tr}{\operatorname{Tr}}
\renewcommand{\tr}{\operatorname{Tr}}
\providecommand{\Ad}{\operatorname{Ad}}
\renewcommand{\Ad}{\operatorname{Ad}}
\providecommand{\Scal}{\SMG}
\renewcommand{\Scal}{\SMG}
\providecommand{\Sfun}{\mathbb S}
\renewcommand{\Sfun}{\mathbb S}
\providecommand{\lr}[2]{c_{#1}^{#2}}
\renewcommand{\lr}[2]{c_{#1}^{#2}}
\setcounter{theorem}{0}
\setcounter{proposition}{0}
\setcounter{lemma}{0}
\setcounter{corollary}{0}
\setcounter{remark}{0}
\begin{abstract}
For the conjugation representation of $U(n)$ on
$W=\End(\C^n)$, this note proves the trace form of the generalized Molien
integral, an exact finite-rank Littlewood--Richardson coefficient formula, and
the stable description by colored necklaces.  The formulas are checked by an
independent matrix computation.  The accompanying marimo notebook constructs
unitary generators and their $n^2$-dimensional conjugation matrices, builds the
full derived action on tensor products of symmetric powers, and computes its
common kernel.  All 18 representative comparisons pass for $1\leq n\leq4$,
total degrees through $3$, and full target dimensions through $405$.  The
checks include the nonstable value $5$ at $n=2$, $\tau=(1,1,1)$, versus the
stable value $6$.
\end{abstract}

\subsection{Imported target}
For a compact group $G$ acting on $W$, Tools~\ref{tool:molien} and
\ref{tool:coefficient} give
\begin{equation}\label{unitary-conjugation-proof-and-verification:eq:definition}
\Scal_{G,W}(\mathbf t)=\smgv{G}{W}
=\int_G\prod_{a\ge1}\det(I-t_a\rho_W(g))^{-1}\,dg,
\end{equation}
\begin{equation}\label{unitary-conjugation-proof-and-verification:eq:target}
[m_\tau]\Scal_{G,W}=\dim\left(\bigotimes_a\Sym^{\tau_a}W\right)^G.
\end{equation}
Writing $P_k=\sum_a t_a^k$, the logarithmic determinant expansion gives
\begin{equation}\label{unitary-conjugation-proof-and-verification:eq:exponential}
\boxed{\Scal_{G,W}(\mathbf t)=\int_G\exp\!\left(
\sum_{k\ge1}\frac{P_k}{k}\Tr(\rho_W(g)^k)\right)dg.}
\end{equation}
These identities are quoted; the new calculation begins with the conjugation
trace identity and Littlewood--Richardson decomposition.

\subsection{The conjugation representation}

Take $G=U(n)$, $V=\C^n$, and $W=\End(V)=V\otimes V^*$.  The action is
$A\mapsto gAg^{-1}$.  If $g$ has eigenvalues $z_1,\ldots,z_n$, its
eigenvalues on $W$ are $z_i z_j^{-1}$.  Hence
\begin{equation}\label{unitary-conjugation-proof-and-verification:eq:trace}
  \tr(\Ad(g)^k)
  =\sum_{i,j}z_i^kz_j^{-k}
  =\tr(g^k)\tr(g^{-k})
  =|\tr(g^k)|^2.
\end{equation}
Substitution into \eqref{unitary-conjugation-proof-and-verification:eq:exponential} proves the exact finite-rank formula
\begin{equation}\label{unitary-conjugation-proof-and-verification:eq:finite-integral}
  \boxed{
  \Scal_{U(n),\End(V)}(\mathbf t)
  =\int_{U(n)}\exp\left(
    \sum_{k\geq1}\frac{P_k}{k}|\tr(g^k)|^2
  \right)dg.}
\end{equation}

\subsection{Exact finite-rank coefficients}

\begin{theorem}[Finite-rank formula]\label{unitary-conjugation-proof-and-verification:thm:finite}
For every partition $\tau=(\tau_1,\ldots,\tau_s)$,
\begin{equation}\label{unitary-conjugation-proof-and-verification:eq:lr}
  \boxed{
  [m_\tau]\Scal_{U(n),\End(V)}
  =\sum_{\substack{
      \lambda^{(a)}\vdash\tau_a\\
      \ell(\lambda^{(a)})\leq n}}
    \ \sum_{\substack{\nu\vdash |\tau|\\\ell(\nu)\leq n}}
    \left(c_{\lambda^{(1)},\ldots,\lambda^{(s)}}^\nu\right)^2.}
\end{equation}
\end{theorem}

\begin{proof}
Cauchy decomposition gives, as a $GL(V)\times GL(V)$ module,
\[
  \Sym^r(V\otimes V^*)
  \cong\bigoplus_{\substack{\lambda\vdash r\\\ell(\lambda)\leq n}}
  \Sfun_\lambda V\otimes\Sfun_\lambda V^*.
\]
Tensor this decomposition over $a=1,\ldots,s$.  For a fixed tuple
$(\lambda^{(1)},\ldots,\lambda^{(s)})$, the Littlewood--Richardson rule gives
\[
  \bigotimes_a\Sfun_{\lambda^{(a)}}V
  \cong\bigoplus_{\substack{\nu\\\ell(\nu)\leq n}}
  \left(\Sfun_\nu V\right)^{\oplus
  c_{\lambda^{(1)},\ldots,\lambda^{(s)}}^\nu}.
\]
The associated dual factor has the same multiplicities.  Schur's lemma says
that the invariant dimension of this summand is the sum of the squares of
those multiplicities.  Summing over the partition tuple proves
\eqref{unitary-conjugation-proof-and-verification:eq:lr}.
\end{proof}

\subsection{Stable contractions and colored necklaces}

Let $d=|\tau|$ and
$H_\tau=S_{\tau_1}\times\cdots\times S_{\tau_s}\leq S_d$.  Since
\[
  \bigotimes_a\Sym^{\tau_a}W=(W^{\otimes d})^{H_\tau},
  \qquad
  W^{\otimes d}=V^{\otimes d}\otimes(V^*)^{\otimes d},
\]
the first fundamental theorem for $GL(V)$ supplies invariant contraction
tensors $T_\sigma$, indexed by $\sigma\in S_d$.  When $n\geq d$ these tensors
are linearly independent, hence form a basis.  A block permutation
$h\in H_\tau$ acts by
\[
  hT_\sigma=T_{h\sigma h^{-1}}.
\]
Therefore the fixed dimension is the number of orbits of this permutation
basis.

\begin{corollary}[Stable coefficient]\label{unitary-conjugation-proof-and-verification:cor:stable}
If $n\geq d=|\tau|$, then
\begin{equation}\label{unitary-conjugation-proof-and-verification:eq:orbits}
  \boxed{
  [m_\tau]\Scal_{U(n),\End(V)}
  =\#(S_d/H_\tau\text{-conjugacy}).}
\end{equation}
Consequently all coefficients of total degree at most $n$ have stabilized.
\end{corollary}

Give the $d$ letters colors so that exactly $\tau_a$ letters have color $a$.
Conjugating by $H_\tau$ relabels letters within each color and preserves only
the cyclic color word attached to each cycle of a permutation.  Thus an orbit
is precisely a multiset of nonempty colored necklaces.  If $\mathcal N$ is
the set of all nonempty necklaces and
\[
  \operatorname{wt}(N)=\prod_a t_a^{\#\{\text{letters of color }a\text{ in }N\}},
\]
the multiset construction gives the stable Euler product
\begin{equation}\label{unitary-conjugation-proof-and-verification:eq:necklace}
  \boxed{
  \Scal_{U(\infty),\End}(\mathbf t)
  =\prod_{N\in\mathcal N}\frac{1}{1-\operatorname{wt}(N)}.}
\end{equation}
This proves the proposed colored-cycle interpretation and fixes its meaning
coefficientwise.

\subsection{Independent matrix verification}

The notebook does not use \eqref{unitary-conjugation-proof-and-verification:eq:lr} to obtain its direct target.  It
constructs the group representation and then solves the invariant problem in
the full target space.

\subsubsection*{Group and conjugation matrices}
For column-major vectorization,
\begin{equation}\label{unitary-conjugation-proof-and-verification:eq:admatrix}
  \operatorname{vec}(gAg^*)=(\overline g\otimes g)\operatorname{vec}(A).
\end{equation}
The code constructs $n^2$ explicit unitary generators by exponentiating the
standard skew-Hermitian basis
\[
  iE_{jj},\qquad E_{ij}-E_{ji},\qquad i(E_{ij}+E_{ji})\quad(i<j).
\]
It checks unitarity, the homomorphism identity
$\Ad(gh)=\Ad(g)\Ad(h)$, and \eqref{unitary-conjugation-proof-and-verification:eq:trace} for powers $1\leq k\leq4$.

\subsubsection*{Direct invariant dimension}
Differentiating \eqref{unitary-conjugation-proof-and-verification:eq:admatrix} gives the $n^2\times n^2$ matrix
\begin{equation}\label{unitary-conjugation-proof-and-verification:eq:derived-ad}
  d\Ad(X)=I\otimes X-X^{\mathsf T}\otimes I,
  \qquad d\Ad(X)(A)=XA-AX.
\end{equation}
If $B=(b_{ij})$ acts on a $q$-dimensional space and
$x^\alpha=x_1^{\alpha_1}\cdots x_q^{\alpha_q}$ is a monomial basis element
of $\Sym^r$, its derived action is implemented directly as
\begin{equation}\label{unitary-conjugation-proof-and-verification:eq:sym-action}
  d\Sym^r(B)x^\alpha
  =\sum_{i,j}\alpha_j b_{ij}x^{\alpha-e_j+e_i}.
\end{equation}
On a tensor product, the code sums \eqref{unitary-conjugation-proof-and-verification:eq:sym-action} over the tensor
factors.  It forms the positive semidefinite Gram matrix
\[
  Q=\sum_{X\text{ in the basis of }\mathfrak u(n)}
  d\rho_\tau(X)^*d\rho_\tau(X).
\]
Its nullity is the common-kernel dimension.  Because $U(n)$ is connected,
this common kernel is exactly the $U(n)$-fixed space.  This is a deterministic
floating-point calculation, not a Haar-sampling estimate.  Each reported row
includes the scale-dependent rank threshold, largest null eigenvalue,
smallest positive eigenvalue, spectral-gap ratio, and invariance residual.  A
row without a gap ratio of at least $10^4$ is labelled inconclusive rather
than rounded to a nullity.

\subsubsection*{Three comparison routes}
The finite formula is computed by enumerating LR tableaux with all length
cutoffs retained.  In stable cases, the notebook separately (i) enumerates
$H_\tau$-conjugacy orbits in $S_d$ and (ii) expands the necklace Euler product
through multidegree $\tau$.  Agreement of all applicable columns therefore
tests the representation construction, the finite formula, the stability
threshold, and the necklace bijection.

\begin{table}[ht]
\centering
\small
\setlength{\tabcolsep}{5.3pt}
\begin{tabular}{@{}r c r r r c@{}}
\toprule
$n$ & $\tau$ & full target dim. & direct nullity & finite LR & stable orbit/Euler \\
\midrule
1 & $(1)$       & 1   & 1 & 1 & 1 \\
1 & $(2)$       & 1   & 1 & 1 & -- \\
1 & $(1,1)$     & 1   & 1 & 1 & -- \\
1 & $(3)$       & 1   & 1 & 1 & -- \\
1 & $(2,1)$     & 1   & 1 & 1 & -- \\
2 & $(1)$       & 4   & 1 & 1 & 1 \\
2 & $(2)$       & 10  & 2 & 2 & 2 \\
2 & $(1,1)$     & 16  & 2 & 2 & 2 \\
2 & $(3)$       & 20  & 2 & 2 & -- \\
2 & $(2,1)$     & 40  & 3 & 3 & -- \\
2 & $(1,1,1)$   & 64  & 5 & 5 & -- \\
3 & $(1)$       & 9   & 1 & 1 & 1 \\
3 & $(2)$       & 45  & 2 & 2 & 2 \\
3 & $(1,1)$     & 81  & 2 & 2 & 2 \\
3 & $(3)$       & 165 & 3 & 3 & 3 \\
3 & $(2,1)$     & 405 & 4 & 4 & 4 \\
4 & $(1)$       & 16  & 1 & 1 & 1 \\
4 & $(2)$       & 136 & 2 & 2 & 2 \\
\bottomrule
\end{tabular}
\caption{All 18 representative comparisons pass.  A dash means $n<|\tau|$,
so the stable formula is not asserted.}
\label{unitary-conjugation-proof-and-verification:tab:checks}
\end{table}

For $n=1,2,3,4$, the maximum errors in unitarity, the homomorphism identity,
and the trace identity are respectively below
$3.2\times10^{-16}$, $1.5\times10^{-16}$, and $1.8\times10^{-15}$.
The test suite also checks the stable values
\[
  \tau=(3),(2,1),(1,1,1),(2,2)
  \quad\longmapsto\quad 3,4,6,10
\]
by both orbit enumeration and necklace coefficient extraction.

\Needspace{11\baselineskip}
\subsection{Reproduction and scope}

The noninteractive companion files are grouped under
\path{proved_matrix_groups/packages/unitary_conjugation_endomorphisms/}.
From the repository root, run
\begin{verbatim}
python3 proved_matrix_groups/packages/unitary_conjugation_endomorphisms/
  unitary_conjugation_core.py
python3 -m pytest proved_matrix_groups/packages/
  unitary_conjugation_endomorphisms/test_verification.py
marimo edit marimo_notebooks/
  unitary_conjugation.py
\end{verbatim}
Join each wrapped command onto one shell line.
The notebook's interactive direct calculation uses a
safety cap of $800$ on the full target dimension; the combinatorial formulas
can be explored farther.  The proof itself has no dimension restriction
beyond $n\geq1$ for the finite formula and the explicitly stated stable
condition $n\geq|\tau|$.
\endgroup


\clearpage
\section{Proctor's odd symplectic group and compact replacements}\label{proof:odd-symplectic}
\begingroup
\definecolor{ink}{HTML}{17212B}
\definecolor{accent}{HTML}{315B7D}
\definecolor{pale}{HTML}{EDF4F8}
\definecolor{passgreen}{HTML}{16794A}
\providecommand{\Sfun}{\SMG}
\renewcommand{\Sfun}{\SMG}
\providecommand{\Sym}{\operatorname{Sym}}
\renewcommand{\Sym}{\operatorname{Sym}}
\providecommand{\Sp}{\operatorname{Sp}}
\renewcommand{\Sp}{\operatorname{Sp}}
\providecommand{\rad}{\operatorname{rad}}
\renewcommand{\rad}{\operatorname{rad}}
\providecommand{\Spec}{\operatorname{Spec}}
\renewcommand{\Spec}{\operatorname{Spec}}
\providecommand{\PASS}{\textcolor{passgreen}{\textbf{PASS}}}
\renewcommand{\PASS}{\textcolor{passgreen}{\textbf{PASS}}}
\providecommand{\summarybox}[1]{\begin{center}\fcolorbox{accent}{pale}{\parbox{0.91\linewidth}{#1}}\end{center}}
\renewcommand{\summarybox}[1]{\begin{center}\fcolorbox{accent}{pale}{\parbox{0.91\linewidth}{#1}}\end{center}}
\setcounter{theorem}{0}
\setcounter{proposition}{0}
\setcounter{lemma}{0}
\setcounter{corollary}{0}
\setcounter{remark}{0}
\summarybox{\textbf{Outcome.} The proposed normalized Haar expectation for the
full complex odd symplectic group is not defined: the group contains both a
vector group and $\mathbb C^\times$. Three valid replacements were proved and
tested. Exact constant-term calculations verify the two compact formulas
through total degree six. The full algebraic invariant series is exactly $1$.
For finite fields, four matrix groups of orders $24$, $432$, $12{,}000$, and
$11{,}520$ were enumerated in dimensions three and five; all 38 tested
coefficients agree between direct permutation-matrix averages and the proposed
cycle formula. Selected coefficients also agree with explicit orbit counts.}

\subsection{The obstruction comes before computation}

Let $V=W\oplus L$ over $\mathbb C$, where $\dim W=2n$, $L=\mathbb C e_0$,
and $\omega$ is a nondegenerate alternating form on $W$. Extend it by
\[
 \Omega((w,s),(w',s'))=\omega(w,w'),\qquad \rad\Omega=L.
\]
Consider the full stabilizer
\[
 G=\{g\in GL(V):\Omega(gv,gw)=\Omega(v,w)\}.
\]
Because the radical is intrinsic, $gL=L$. In the decomposition $W\oplus L$,
every element therefore has the form
\begin{equation}
 g=\begin{pmatrix}A&0\\ \ell&a\end{pmatrix},\qquad
 A\in\Sp(W),\quad \ell\in W^*,\quad a\in\mathbb C^\times.
\end{equation}
Conversely, every matrix in (1) preserves $\Omega$. Thus
\begin{equation}
 G\cong W^*\rtimes\bigl(\Sp_{2n}(\mathbb C)\times\mathbb C^\times\bigr).
\end{equation}
The vector group $W^*$ is the unipotent radical, so $G$ is nonreductive; both
$W^*$ and $\mathbb C^\times$ are noncompact. A locally compact group has a
normalized Haar probability exactly when it is compact. Consequently
\[
 \frac{1}{\operatorname{vol}(G)}\int_G f(g)\,dg
\]
is not a defined normalization. This is a mathematical obstruction, not a
numerical convergence problem. A finite box cutoff would depend on coordinates
and cutoff shape and is not a group-invariant substitute.

There is a useful convention warning. Some authors impose $g(e_0)=e_0$. We
write this subgroup as
\begin{equation}
 G^1=W^*\rtimes\Sp_{2n}(\mathbb C),
\end{equation}
which is obtained by setting $a=1$. The full stabilizer $G$ and $G^1$ give
different answers below and should not be conflated.

\subsection{A structural identity and what it sees}

Block triangularity gives, for every scalar $t$,
\begin{equation}
 \det(I-tg)=(1-at)\det(I-tA),\qquad
 \Spec(g)=\Spec(A)\cup\{a\}.
\end{equation}
The shear coordinate $\ell$ can change Jordan structure but is invisible to
the characteristic polynomial. Therefore any probability measure $\nu$ for
which the coefficientwise moments exist satisfies
\begin{equation}
 \int_G\prod_i\det(I-t_i g)^{-1}\,d\nu(g)
 =\int_{\Sp_{2n}(\mathbb C)\times\mathbb C^\times}
 \prod_i\frac{\det(I-t_iA)^{-1}}{1-at_i}\,d\bar\nu(A,a),
\end{equation}
where $\bar\nu$ is the pushforward to the Levi quotient. This gives a useful
diagnostic for heat-kernel or user-chosen noncompact probabilities: the answer
depends on the chosen measure, but never on its conditional distribution in
the unipotent coordinate.

The code tests (4) modulo $q$ on deterministic samples from every constructed
finite group. It is important that this structural test is separate from the
permutation-representation test in Section~\ref{proctor-odd-symplectic-proof-and-verification:sec:finite}.

\subsection{Two compact replacements}

A maximal compact subgroup of the full stabilizer is
\[
 K=\Sp(n)\times U(1),\qquad V|_K=W\oplus\chi,
\]
where $\chi(z)=z$ on the radical line. Write
\[
 \Sfun_{H,X}(\mathbf t)=\int_H\prod_i\det(I-t_i h|_X)^{-1}\,dh
\]
for a compact group with normalized Haar probability. Direct-sum
multiplicativity and product Haar measure give
\[
 \Sfun_{K,V}=\Sfun_{\Sp(n),W}\Sfun_{U(1),\chi}.
\]
Since
\[
 \prod_i(1-zt_i)^{-1}=\sum_{d\ge0}z^d h_d(\mathbf t),
 \qquad \int_{U(1)}z^d\,dz=\delta_{d0},
\]
we obtain the exact identity
\begin{equation}
 \boxed{\Sfun_{K,V}=\Sfun_{\Sp(n),W}.}
\end{equation}
The radical line disappears because it has a nontrivial $U(1)$ weight.

For the fixed-radical group $G^1$, a maximal compact subgroup is $\Sp(n)$ and
$V|_{\Sp(n)}=W\oplus\mathbf 1$. The extra eigenvalue is now $1$, so
\begin{equation}
 \boxed{\Sfun_{\Sp(n),W\oplus\mathbf1}
 =\Sfun_{\Sp(n),W}\prod_i(1-t_i)^{-1}.}
\end{equation}

\subsubsection*{Exact rank-one check}

For $n=1$, restrict to the maximal torus of $SU(2)=\Sp(1)$ with weights
$z,z^{-1}$. Weyl integration is the exact constant-term operation
\[
 \int_{SU(2)}f=\operatorname{CT}_z\bigl((1-z^2)f(z)\bigr).
\]
The verification represents every complete homogeneous character as a Laurent
polynomial. It extracts both the $z$ and $U(1)$ constant terms and tests all 30
partitions through total degree six. Representative results are:
\begin{center}
\begin{tabular}{@{}lrrrr@{}}
\toprule
$\tau$ & $K$ direct & $\Sp(1),W$ & $G^1$ compact direct & formula (7)\\
\midrule
$()$      &1&1&1&1\\
$(1)$     &0&0&1&1\\
$(2)$     &0&0&1&1\\
$(1,1)$   &1&1&2&2\\
$(2,1)$   &0&0&2&2\\
$(2,2)$   &1&1&3&3\\
\bottomrule
\end{tabular}
\end{center}
All 30 comparisons pass exactly, with integer arithmetic.

\subsection{The full algebraic invariant series}

Even though Haar averaging fails, the covariant invariant series is meaningful:
\begin{equation}
 \Sfun^{\mathrm{inv}}_{G,V}
 =\sum_\tau\dim\left(\bigotimes_j\Sym^{\tau_j}V\right)^Gm_\tau.
\end{equation}

\textbf{Theorem.} For the full stabilizer $G$ in (1),
\begin{equation}
 \boxed{\Sfun^{\mathrm{inv}}_{G,V}=1.}
\end{equation}

\textit{Proof.}
Embed $\bigotimes_j\Sym^{\tau_j}V$ into $V^{\otimes d}$, where
$d=|\tau|$. The subgroup $\operatorname{diag}(I_W,a)$ acts with weight $a^r$
on tensors having $r$ radical-line factors. Its invariant subspace is thus
contained in $W^{\otimes d}$.

For $\ell\in W^*$, differentiate the action of
\[
 u_\ell=\begin{pmatrix}I_W&0\\ \ell&1\end{pmatrix}.
\]
On $x\in W^{\otimes d}$ the derivative is a sum of $d$ contractions, one for
each tensor position, and the output of the $j$th contraction has its unique
$e_0$ in position $j$. These outputs lie in distinct direct summands. Hence
invariance forces every positional contraction against every $\ell$ to vanish.
Already at the first position, the contractions by a basis of $W^*$ recover
all coordinates of $x$, so $x=0$. Thus there are no positive-degree
invariants. \hfill$\square$

This proof gives a compact computational certificate: after the torus weight
filter, stack the first-position contractions against a basis of $W^*$. The
result is a coordinate permutation of rank $(2n)^d$. The implementation records
full rank for $n=1,2$ and $1\le d\le4$ (eight checks), so every recorded kernel
has dimension zero. This certificate is stronger and cheaper than building
large Reynolds matrices for a nonreductive group.

\subsection{Finite odd symplectic groups and a genuine complex representation}
\label{proctor-odd-symplectic-proof-and-verification:sec:finite}

Let $E=W\oplus L$ over a prime field $\mathbb F_q$. The same block proof gives
\begin{equation}
 G^{\mathrm{odd}}_{2n+1}(q)
 \cong W^*\rtimes\bigl(\Sp_{2n}(q)\times\mathbb F_q^\times\bigr),
\end{equation}
and hence
\begin{align}
 |G^{\mathrm{odd}}_{2n+1}(q)|
 &=q^{2n}(q-1)|\Sp_{2n}(q)|,\notag\\
 |\Sp_{2n}(q)|&=q^{n^2}\prod_{j=1}^n(q^{2j}-1).
\end{align}

Matrices with entries in $\mathbb F_q$ are not automatically complex matrices.
We therefore use the honest complex permutation representation on
$S=\mathbb P(E)$. If $P_g$ is its permutation matrix and $c_r(g)$ is the
number of $r$-cycles, then
\begin{equation}
 \det(I-tP_g)^{-1}=\prod_{r\ge1}(1-t^r)^{-c_r(g)}.
\end{equation}
Consequently the proposed finite formula is
\begin{equation}
 \boxed{\Sfun_{G(q),\mathbb C[S]}(\mathbf t)
 =\frac1{|G(q)|}\sum_{g\in G(q)}
 \prod_i\prod_{r\ge1}(1-t_i^r)^{-c_r(g)}.}
\end{equation}
For a partition $\tau$, Burnside's lemma also gives
\begin{equation}
 [m_\tau]\Sfun_{G(q),\mathbb C[S]}
 =\#\left(G(q)\backslash\prod_j\operatorname{MSet}_{\tau_j}(S)\right).
\end{equation}

There are exactly two orbits on $\mathbb P(E)$: the radical point and all
nonradical points. The latter assertion follows because $\Sp(W)$ is transitive
on nonzero vectors of $W$, while $\ell(w)$ adjusts the radical coordinate.
Therefore
\begin{equation}
 \boxed{[m_{(1)}]\Sfun_{G(q),\mathbb C[\mathbb P(E)]}=2.}
\end{equation}

\subsection{Independent verification algorithms}

For each $(n,q)$, the program performs the following steps.
\begin{enumerate}[leftmargin=*,itemsep=3pt]
\item Enumerate every $2n\times2n$ matrix $A$ satisfying
$A^{\mathsf T}JA=J$, then every triple $(A,\ell,a)$ in (10). Check the observed
orders against (11), preservation of the degenerate form, and deterministic
closure samples.
\item Enumerate normalized representatives of every projective point and
construct the permutation $P_g$ for every group matrix.
\item \textbf{Direct matrix route:} compute $p_k=\operatorname{tr}(P_g^k)$
from explicit permutation powers and use Newton's recurrence
\[
 d h_d=\sum_{k=1}^d p_kh_{d-k}.
\]
Average $\prod_jh_{\tau_j}$ over all matrices.
\item \textbf{Cycle route:} independently find the cycles and expand the
right side of (12) by dynamic programming. Average the resulting products and
compare every coefficient with the direct route.
\item \textbf{Orbit route:} in the two smaller cases, enumerate the orbits on
multisets for $\tau=(1),(2),(1,1)$ and compare with (14).
\end{enumerate}

The direct route does not call the cycle-product routine. This separation
matters: otherwise (13) would only be tested against a rearrangement of itself.

\subsection{Exact results}

\begin{center}
\begin{tabularx}{\linewidth}{@{}rrrrrrX@{}}
\toprule
$n$ & $q$ & $\dim E$ & $|\Sp_{2n}(q)|$ & $|G(q)|$ & $|\mathbb P(E)|$ & tested degrees\\
\midrule
1&2&3&6&24&7&all partitions through 4\\
1&3&3&24&432&13&all partitions through 4\\
1&5&3&120&12,000&31&all partitions through 3\\
2&2&5&720&11,520&31&all partitions through 3\\
\bottomrule
\end{tabularx}
\end{center}
All four form, order, radical-point, determinant, closure, and coefficient
checks pass. In total, 38 matrix-average coefficients match 38 cycle-formula
coefficients exactly. Selected values illustrate both stability and genuine
dimension dependence:
\begin{center}
\begin{tabular}{@{}ccrrrrrr@{}}
\toprule
$(n,q)$ & $m_{(1)}$ & $m_{(2)}$ & $m_{(1,1)}$ & $m_{(3)}$ & $m_{(2,1)}$ & $m_{(1,1,1)}$\\
\midrule
$(1,2)$&2&5&6&11&18&25\\
$(1,3)$&2&5&6&12&19&28\\
$(1,5)$&2&5&6&14&22&30\\
$(2,2)$&2&6&7&19&31&43\\
\bottomrule
\end{tabular}
\end{center}
For $(n,q)=(1,2)$ and $(1,3)$, explicit orbit enumeration gives respectively
$2$, $5$, and $6$ orbits for $\tau=(1),(2),(1,1)$, exactly matching both
averaging routes.

\subsection{Interpretation, limitations, and new tools}

The results separate four notions that are easy to blur:
\begin{itemize}[leftmargin=*,itemsep=3pt]
\item The full complex-group Haar expression is undefined, and no computation
can validate it without first choosing a new measure.
\item Maximal-compact averaging is canonical up to conjugacy, but it forgets
the unipotent radical. For the full stabilizer it also removes the radical
line by $U(1)$ weight; for $G^1$ it retains that line as a trivial summand.
\item The full covariant invariant series is well-defined but collapses to $1$
in the defining representation. Interesting odd symplectic representation
theory therefore requires other modules, mixed $V,V^*$ constructions, or
polynomial-function invariants.
\item Finite-field permutation actions give exact, nontrivial, reproducible
$\Lambda$-distributions, but they are different representations of different
finite groups and should not be presented as complex matrix realizations of
the defining modular matrices.
\end{itemize}

Two reusable tools emerge. First, the \emph{Levi visibility test} (4)--(5)
immediately decides which nonreductive coordinates an eigenvalue statistic can
detect. Second, the \emph{torus-filter plus derivative certificate} proves
invariant vanishing using small full-rank contraction maps instead of a
nonexistent Reynolds average.

\subsection{Reproducibility}

From the repository root, run
\begin{verbatim}
.venv/bin/python \
  for_this_guy/proctor_odd_symplectic_group/verification.py
\end{verbatim}
and open the interactive report with
\begin{verbatim}
.venv/bin/marimo edit \
  marimo_notebooks/proctor_odd_symplectic.py
\end{verbatim}
The verification uses only exact integer and rational arithmetic. Exhaustive
group enumeration is intentionally limited to prime fields and to the four
displayed cases; extending to larger ranks should replace all-matrix search by
standard symplectic generators.

\begin{thebibliography}{9}
\bibitem{proctor-odd-symplectic-proof-and-verification:Okada}
S.~Okada, \emph{A bialternant formula for odd symplectic characters and its
application}, 2019, \url{https://arxiv.org/abs/1905.12964}.
\bibitem{proctor-odd-symplectic-proof-and-verification:Burnside}
W.~Burnside, \emph{Theory of Groups of Finite Order}, second edition, 1911.
\bibitem{proctor-odd-symplectic-proof-and-verification:Haar}
A.~Haar, \emph{Der Massbegriff in der Theorie der kontinuierlichen Gruppen},
Annals of Mathematics 34 (1933), 147--169.
\end{thebibliography}
\endgroup


\clearpage
\part{Polyhedral, sporadic, and special 24-dimensional groups}
\section{Finite rotation groups and baseline matrix-group formulas}\label{proof:finite-matrix-groups}
\begingroup
\providecommand{\Sym}{\operatorname{Sym}}
\renewcommand{\Sym}{\operatorname{Sym}}
\providecommand{\EE}{\mathbb E}
\renewcommand{\EE}{\mathbb E}
\setcounter{theorem}{0}
\setcounter{proposition}{0}
\setcounter{lemma}{0}
\setcounter{corollary}{0}
\setcounter{remark}{0}
\begin{abstract}
We verify finite-group instances of the generalized Molien formula by full
matrix enumeration.  The examples cover the tetrahedral, octahedral, and
icosahedral rotation representations; the full $H_3$ reflection group;
restricted monomial groups including $G(r,p,n)$; an exterior-power
representation; and a regular representation.  In total, 492 coefficient
comparisons pass, together with independent Reynolds-projector checks.  A
dual-code cycle counter for monomial groups is included as a reusable tool.
\end{abstract}

\subsection{Imported finite-group tools}
The universal determinant and coefficient identities are imported
Tools~\ref{tool:molien} and~\ref{tool:coefficient}. This section is a spectral
compression and verification study, beginning with rotation-angle classes.

\subsection{Polyhedral rotation groups}
A rotation through angle $\theta$ in three dimensions has spectrum
$1,e^{i\theta},e^{-i\theta}$ and hence
\[
 \det(I-tR_\theta)^{-1}
 =\frac1{(1-t)(1-2\cos(\theta)t+t^2)}.
\]
Notice the multiplication by $t$ in the middle term; this corrects a
typographical ambiguity in the supplied expression.  Writing $\Phi_\theta$
for the product of this factor over all variables $t_i$, class enumeration
gives
\begin{align*}
 \MGF_{A_4}^{\rm rot}
 &=\frac1{12}(\Phi_0+8\Phi_{2\pi/3}+3\Phi_\pi),\\
 \MGF_{S_4}^{\rm rot}
 &=\frac1{24}(\Phi_0+8\Phi_{2\pi/3}+6\Phi_{\pi/2}+9\Phi_\pi),\\
 \MGF_{A_5}^{\rm rot}
 &=\frac1{60}(\Phi_0+20\Phi_{2\pi/3}+15\Phi_\pi
 +12\Phi_{2\pi/5}+12\Phi_{4\pi/5}).
\end{align*}
The notebook constructs the first two groups as determinant-one signed
permutation matrices (restricting to a tetrahedron for $A_4$).  It constructs
$A_5$ by enumerating all orthogonal maps preserving the twelve icosahedron
vertices.  No class labels are used on the direct side.

Adjoining the negatives of the 60 icosahedral rotations gives the order-120
$H_3$ reflection group.  Its ten spectral classes are checked against complete
matrix averaging.  Independently, the one-variable invariant series
\[
 \frac1{(1-t^2)(1-t^6)(1-t^{10})}
\]
matches the directly averaged dimensions through degree ten.

\subsection{A dual-code cycle formula for monomial groups}
Let $A\le(\mathbb Z/r\mathbb Z)^n$ be a linear code stable under a permutation
group $H\le S_n$.  Put
\[
 G=D(A)\rtimes H,
 \qquad D(x)=\operatorname{diag}(\zeta^{x_1},\ldots,\zeta^{x_n}),
 \quad \zeta=e^{2\pi i/r}.
\]
For a permutation $\sigma$, expand the determinant factor cycle by cycle.
For every cycle $C$ and every variable index $a$, choose a multiplicity
$m_{C,a}\ge0$ satisfying
\[
 \sum_C |C|m_{C,a}=\tau_a.
\]
The diagonal phase exponent is constant on $C$ and equals
$e_j=\sum_a m_{C,a}\pmod r$.  Character orthogonality over $A$ gives
\[
 \frac1{|A|}\sum_{x\in A}\zeta^{\langle x,e\rangle}
 =\mathbf 1_{\{e\in A^\perp\}}.
\]
Thus the coefficient is obtained by counting these cycle assignments whose
exponent vector lies in the dual code, then averaging over $H$.  This removes
the diagonal loop entirely and is useful well beyond the test cases.

For the complex reflection group $G(r,p,n)$, the code is
$A=\{x:\sum_jx_j\equiv0\pmod p\}$.  The notebook also tests a specialized
counter: for each $q=0,\ldots,p-1$, take label vectors
\[
 u\in\prod_a\{0,\ldots,\tau_a\},
 \qquad \sum_a u_a\equiv q\,r/p\pmod r,
\]
and count multisets of exactly $n$ labels with componentwise sum $\tau$.
Summing these counts over $q$ yields the coefficient.  Dynamic programming
implements this finite multiset count independently of matrix averaging.

\subsection{Alternate representations}
If $V$ has eigenvalues $\lambda_1,\ldots,\lambda_d$, then $\bigwedge^kV$ has
eigenvalues $\prod_{j\in J}\lambda_j$ over $k$-subsets $J$.  The notebook
constructs $\bigwedge^2$ of the three-dimensional standard reflection
representation of $S_4$ using matrix minors, while the formula side uses the
spectral products.  All group elements are enumerated.

For the regular representation, an element $g$ of order $o(g)$ acts as
$|G|/o(g)$ cycles of length $o(g)$.  Hence
\[
 \boxed{\MGF_{G,\mathrm{Reg}}
 =\frac1{|G|}\sum_{g\in G}\prod_i
 (1-t_i^{o(g)})^{-|G|/o(g)}.}
\]
For $S_3$, the formula uses the order distribution
$N_1=1,N_2=3,N_3=2$ and is compared with six explicit $6\times6$ left-regular
permutation matrices.

\subsection{Test matrix and outcome}
\begin{center}
\begin{tabular}{@{}p{0.34\linewidth}p{0.38\linewidth}r@{}}
\toprule
Cases & Independent comparison & Checks\\
\midrule
$A_4,S_4,A_5$ rotations & all matrices versus angle classes & 135\\
Full $H_3$ reflection group & 120 matrices versus ten spectra; degrees $2,6,10$ & 139\\
$D(A)\rtimes C_3$ and five $G(r,p,n)$ cases & matrix average versus dual-code or multiset counter & 180\\
$\bigwedge^2$ standard $S_4$ and regular $S_3$ & explicit matrices versus transformed spectra/order counts & 38\\
\midrule
Total & & \textbf{492}\\
\bottomrule
\end{tabular}
\end{center}
Every comparison passes.  Additional Reynolds projectors satisfy
$\|R^2-R\|_F<9\times10^{-16}$.  Representative ranks are: rank one on
$\Sym^2(\mathbb C^3)$ for $H_3$; rank one on
$\Sym^2(\mathbb C^3)\otimes\mathbb C^3$ for $G(2,2,3)$; and rank five on
$\Sym^2$ of the six-dimensional regular $S_3$ representation.

\subsection{Reproduction and scope}
From the repository root, run these modules with \texttt{python3 -m}:
\begin{itemize}
\small
\item \path{for_this_guy.matrix_group_formula_verification.h3_reflection.verification}
\item \path{for_this_guy.matrix_group_formula_verification.restricted_monomial.verification}
\item \path{for_this_guy.matrix_group_formula_verification.representation_variants.verification}
\end{itemize}
Their canonical marimo apps are cataloged in
\path{marimo_notebooks/README.md}.  The verification establishes the formulas for the displayed finite cases and
truncations.  The common finite-group identity is imported as
Tools~\ref{tool:coefficient}--\ref{tool:molien}; these tests do not claim
elementary simplifications for every finite group of Lie type, spin
representation, or exceptional representation mentioned in the broader
proposal.
\endgroup


\clearpage
\section[Three representations associated with A5]{Three representations associated with $A_5$}\label{proof:a5-representations}
\begingroup
\definecolor{ink}{HTML}{25213A}
\definecolor{accent}{HTML}{7A3E65}
\definecolor{pale}{HTML}{F8EFF5}
\definecolor{passgreen}{HTML}{16794A}
\providecommand{\Sfun}{\SMG}
\renewcommand{\Sfun}{\SMG}
\providecommand{\Sym}{\operatorname{Sym}}
\renewcommand{\Sym}{\operatorname{Sym}}
\providecommand{\PASS}{\textcolor{passgreen}{\textbf{PASS}}}
\renewcommand{\PASS}{\textcolor{passgreen}{\textbf{PASS}}}
\providecommand{\summarybox}[1]{\begin{center}\fcolorbox{accent}{pale}{\parbox{0.91\linewidth}{#1}}\end{center}}
\renewcommand{\summarybox}[1]{\begin{center}\fcolorbox{accent}{pale}{\parbox{0.91\linewidth}{#1}}\end{center}}
\setcounter{theorem}{0}
\setcounter{proposition}{0}
\setcounter{lemma}{0}
\setcounter{corollary}{0}
\setcounter{remark}{0}
\summarybox{\textbf{Outcome.} All 45 multivariate coefficients indexed by
partitions through total degree seven agree in each of the five-dimensional
permutation, four-dimensional deleted permutation, and three-dimensional
rotation representations.  The explicit groups all have order 60.  Separate
Reynolds projectors confirm the low-degree invariant dimensions.}

\subsection{Imported target}
Tools~\ref{tool:coefficient} and~\ref{tool:molien} identify the desired
series as $\smgv{G}{V}$ and its $m_\tau$ coefficient as
$\dim(\bigotimes_a\Sym^{\tau_a}V)^G$.  This section therefore starts with the
three representation-specific class spectra; it does not rederive the common
identity.

\subsection{Three class-spectrum formulas}

Write $F_{(\lambda_1,\ldots,\lambda_n)}(\mathbf t)
=\prod_{a,i}(1-\lambda_i t_a)^{-1}$,
$\omega=e^{2\pi i/3}$, and $\xi=e^{2\pi i/5}$.  The formulas tested are
\begin{align}
\Sfun_{\rm perm}=\frac1{60}\big[&F_{(1,1,1,1,1)}
+15F_{(1,1,1,-1,-1)}+20F_{(1,1,1,\omega,\omega^{-1})}\notag\\
&+24F_{(1,\xi,\xi^2,\xi^3,\xi^4)}\big],                  \tag{3}\\
\Sfun_{\rm del}=\frac1{60}\big[&F_{(1,1,1,1)}
+15F_{(1,1,-1,-1)}+20F_{(1,1,\omega,\omega^{-1})}\notag\\
&+24F_{(\xi,\xi^2,\xi^3,\xi^4)}\big],                   \tag{4}\\
\Sfun_{\rm rot}=\frac1{60}\big[&F_{(1,1,1)}
+15F_{(1,-1,-1)}+20F_{(1,\omega,\omega^{-1})}\notag\\
&+12F_{(1,\xi,\xi^{-1})}+12F_{(1,\xi^2,\xi^{-2})}\big]. \tag{5}
\end{align}

For (3), a permutation of cycle type $(1^5)$, $(2^2,1)$, $(3,1^2)$,
or $(5)$ has precisely the displayed roots of unity as eigenvalues.  Removing
the always-fixed all-ones line deletes one eigenvalue $1$ and proves (4).
For (5), a rotation through angle $\theta$ has spectrum
$(1,e^{i\theta},e^{-i\theta})$; the rotation counts are
$1,15,20,12,12$ for angles $0,\pi,2\pi/3,2\pi/5,4\pi/5$.

The split of the fifth turns is essential in (5): the two classes have
different 3D spectra.  They may be combined in (3)--(4) because a 5-cycle and
its nontrivial powers have the same full permutation spectrum.

\subsection{Independent matrix construction}

For the first representation, all 60 even permutations of five symbols are
converted directly to $5\times5$ permutation matrices.  For the deleted
representation, let $B$ be an orthonormal basis for
\[
\{x\in\mathbb C^5:x_1+\cdots+x_5=0\};
\]
the matrices are $B^*P_\pi B$.  Neither route uses the spectral tables above.

For the rotation representation, start from the twelve icosahedron vertices
\[
(0,\pm1,\pm\varphi),\quad(\pm1,\pm\varphi,0),\quad
(\pm\varphi,0,\pm1),\qquad \varphi=(1+\sqrt5)/2.
\]
Mapping a fixed independent vertex frame to every compatible target frame and
retaining the orientation-preserving isometries reconstructs 60 distinct
$3\times3$ matrices.

\subsection{Verification protocol}

For every partition $\tau$ through total degree seven, the direct route
computes $h_{\tau_a}$ recursively from the eigenvalues of every explicit
matrix and averages $\prod_a h_{\tau_a}$.  The formula route uses only the
weighted spectra in (3), (4), or (5).  It never receives the matrices.

As a third route, the code explicitly constructs each $\Sym^2(V)$ inside
$V^{\otimes2}$ using normalized occupancy states, averages the resulting
matrices, and checks projector rank, trace, and idempotence.  The same check is
also performed on $V\otimes V$.

\subsection{Results}

\begin{center}
\begin{tabularx}{\linewidth}{@{}Xrrrrc@{}}
\toprule
Representation & $|G|$ & $\dim V$ & degree & checks & result\\
\midrule
$A_5$ permutation & 60 & 5 & $0$--$7$ & 45 & \PASS\\
$A_5$ deleted permutation & 60 & 4 & $0$--$7$ & 45 & \PASS\\
$A_5$ rotation & 60 & 3 & $0$--$7$ & 45 & \PASS\\
\bottomrule
\end{tabularx}
\end{center}

\begin{center}
\begin{tabular}{@{}lrrrr@{}}
\toprule
Representation & $[m_{(1)}]$ & $[m_{(2)}]$ & $[m_{(1,1)}]$ & $[m_{(1,1,1)}]$\\
\midrule
permutation & 1 & 2 & 2 & 5\\
deleted permutation & 0 & 1 & 1 & 1\\
rotation & 0 & 1 & 1 & 1\\
\bottomrule
\end{tabular}
\end{center}

The symmetric-square projector ranks are respectively $2,1,1$.  Their traces
agree with those integers to $3\times10^{-16}$, and the largest idempotence
error is $3.27\times10^{-16}$.  These checks recover the claims that the
permutation representation has a fixed line while the deleted and rotation
representations do not.

The value $[m_{(1,1,1)}]=1$ in the rotation case is the invariant orientation
tensor in $V^{\otimes3}$.  This is not contradicted by the absence of an
invariant cubic in $\Sym^3(V)$: $(1,1,1)$ denotes three separate degree-one
factors, not the partition $(3)$.

\subsection{Reproducibility and scope}

Run the command-line suite from the repository root:
{\scriptsize
\begin{verbatim}
.venv/bin/python -m for_this_guy.matrix_group_formula_verification.a5_representations.verification
\end{verbatim}
}
Open \texttt{marimo\_notebooks/a5\_representations.py} with marimo to change the
degree bound and inspect every row.

The finite computation is a transcription and convention test, not the source
of the all-degree theorem.  All-degree validity follows from (1)--(2) and the
complete class spectra; the numerical suite checks the matrix constructions,
class sizes, fifth-root choices, and coefficient extraction in representative
small dimensions.
\endgroup


\clearpage
\section{Binary polyhedral groups}\label{proof:binary-polyhedral}
\begingroup
\definecolor{ink}{HTML}{142B37}
\definecolor{accent}{HTML}{2B6681}
\definecolor{pale}{HTML}{EDF6FA}
\definecolor{passgreen}{HTML}{16794A}
\providecommand{\Sfun}{\SMG}
\renewcommand{\Sfun}{\SMG}
\providecommand{\Sym}{\operatorname{Sym}}
\renewcommand{\Sym}{\operatorname{Sym}}
\providecommand{\PASS}{\textcolor{passgreen}{\textbf{PASS}}}
\renewcommand{\PASS}{\textcolor{passgreen}{\textbf{PASS}}}
\providecommand{\summarybox}[1]{\begin{center}\fcolorbox{accent}{pale}{\parbox{0.91\linewidth}{#1}}\end{center}}
\renewcommand{\summarybox}[1]{\begin{center}\fcolorbox{accent}{pale}{\parbox{0.91\linewidth}{#1}}\end{center}}
\setcounter{theorem}{0}
\setcounter{proposition}{0}
\setcounter{lemma}{0}
\setcounter{corollary}{0}
\setcounter{remark}{0}
\summarybox{\textbf{Outcome.} Explicit $2\times2$ complex matrix groups of
orders $24$, $48$, and $120$ were reconstructed for $2T$, $2O$, and $2I$.
For each group, all 67 multivariate coefficients through total degree eight
agree with the proposed paired-lift formula.  Closure, unitarity, determinant
one, the central element $-I$, and the odd-degree filter all pass.}

\subsection{Imported target}
For each binary polyhedral matrix group, the $\sigma$-MGF and invariant
coefficients are imported Tools~\ref{tool:molien} and
\ref{tool:coefficient}. The new calculation pairs the two lifts of each
rotation class.

\subsection{Proof of the binary-lift formula}

Let $\widetilde\Gamma\leq SU(2)$ be the inverse image of a finite rotation
group $\Gamma\leq SO(3)$.  If $R\in\Gamma$ rotates through angle $\theta$,
its two lifts $q$ and $-q$ have spectra
\[
(e^{i\theta/2},e^{-i\theta/2}),\qquad
(-e^{i\theta/2},-e^{-i\theta/2}).
\]
Define
\begin{align*}
\Psi_\theta^+(\mathbf t)&=\prod_a
(1-2\cos(\theta/2)t_a+t_a^2)^{-1},\\
\Psi_\theta^-(\mathbf t)&=\prod_a
(1+2\cos(\theta/2)t_a+t_a^2)^{-1},\\
B_\theta&=\tfrac12(\Psi_\theta^++\Psi_\theta^-).
\end{align*}
Pair the two terms in (1) over each $R\in\Gamma$.  Because
$|\widetilde\Gamma|=2|\Gamma|$, this gives
\begin{equation}
\Sfun_{\widetilde\Gamma,\mathbb C^2}
=\frac1{|\Gamma|}\sum_\theta N_\Gamma(\theta)B_\theta.    \tag{3}
\end{equation}
This proves the proposed formula without requiring a separate conjugacy table
for the binary group.

Substituting the rotation counts yields
\begin{align}
\Sfun_{2T}&=\frac1{12}(B_0+8B_{2\pi/3}+3B_\pi),            \tag{4}\\
\Sfun_{2O}&=\frac1{24}(B_0+6B_{\pi/2}+8B_{2\pi/3}+9B_\pi),\tag{5}\\
\Sfun_{2I}&=\frac1{60}(B_0+15B_\pi+20B_{2\pi/3}
+12B_{2\pi/5}+12B_{4\pi/5}).                              \tag{6}
\end{align}

\subsection{The central parity filter}

The matrix $-I$ belongs to each binary group.  On
$W_\tau=\bigotimes_a\Sym^{\tau_a}(\mathbb C^2)$ it acts as
$(-1)^{|\tau|}$.  If $|\tau|$ is odd and $v$ is invariant, then
$v=(-I)v=-v$, so $v=0$.  Therefore
\begin{equation}
[m_\tau]\Sfun_{2T}=[m_\tau]\Sfun_{2O}=[m_\tau]\Sfun_{2I}=0
\quad\text{for odd }|\tau|.                                \tag{7}
\end{equation}
The same cancellation is visible termwise in the average defining
$B_\theta$.

\subsection{Independent matrix construction}

The parent rotation matrices are constructed geometrically: signed
permutation matrices give the octahedral group; the subgroup preserving the
four vertices $(1,1,1),(1,-1,-1),(-1,1,-1),(-1,-1,1)$ gives the tetrahedral
group; and compatible vertex-frame maps of the icosahedron give the
icosahedral group.

For each rotation matrix, a standard trace-and-diagonal algorithm recovers one
unit quaternion $q=w+xi+yj+zk$.  It is mapped to
\begin{equation}
\rho(q)=\begin{pmatrix}
w+iz&y+ix\\-y+ix&w-iz
\end{pmatrix}\in SU(2).                                    \tag{8}
\end{equation}
Both $\rho(q)$ and $-\rho(q)$ are retained.  Since the sign choice for an
individual lift is arbitrary but both signs are present, the resulting set is
the full preimage and is independent of those choices.

This construction does not use (4)--(6).  The implementation verifies all
pairwise products against the reconstructed set, as well as unitarity,
determinant one, and the presence of $-I$.

\subsection{Verification protocol and results}

For every partition $\tau$ through total degree eight, the direct route
recursively computes complete homogeneous characters from all explicit
matrices and averages their products.  The independent route uses only the
parent angle counts and paired half-angle spectra in (3).  All odd rows are
also checked separately against zero.

Finally, explicit matrices on $V\otimes V$ and $\Sym^2(V)$ are averaged into
Reynolds projectors.  The alternating determinant tensor gives a one-dimensional
invariant in $V\otimes V$, while $\Sym^2(V)$ has no invariant vector.

\begin{center}
\begin{tabularx}{\linewidth}{@{}Xrrrrc@{}}
\toprule
Group & order & dimension & degree & checks & result\\
\midrule
$2T$ & 24 & 2 & $0$--$8$ & 67 & \PASS\\
$2O$ & 48 & 2 & $0$--$8$ & 67 & \PASS\\
$2I$ & 120 & 2 & $0$--$8$ & 67 & \PASS\\
\bottomrule
\end{tabularx}
\end{center}

All three sets are closed under multiplication.  The largest observed
unitarity error is $3.17\times10^{-16}$ and the largest determinant error is
$2.27\times10^{-16}$.  Every odd-degree coefficient is exactly zero after
integer validation.

For each group, the $4\times4$ Reynolds projector on $V\otimes V$ has rank and
trace $1$; its largest idempotence error is $1.04\times10^{-16}$.  The
symmetric-square projector has rank $0$.  Consequently the computation
recovers the proposed low-degree distinction
\[
[m_{(2)}]=0,\qquad [m_{(1,1)}]=1.
\]

\subsection{Reproducibility and scope}

Run from the repository root:
{\scriptsize
\begin{verbatim}
.venv/bin/python -m for_this_guy.matrix_group_formula_verification.binary_polyhedral.verification
\end{verbatim}
}
Open
\path{marimo_notebooks/binary_polyhedral.py}
with marimo for interactive
degree selection and complete row tables.

The finite checks validate the matrix lift, angle convention, class counts,
central signs, and coefficient code in reasonable low dimensions.  The
all-degree conclusion is the formal proof (1)--(3); it does not depend on the
chosen cutoff or floating-point experiment.
\endgroup


\clearpage
\section{Solvable and finite-subgroup families}\label{proof:new-dist:solvable_finite}

\resultbox{\textbf{Canonical testing artifacts.}\par\smallskip Exact backend: \path{lambda_distributions/dists/solvable_finite_subgroup_sigma_mgfs/verification.py}.\par Regression: \path{lambda_distributions/dists/solvable_finite_subgroup_sigma_mgfs/test_verification.py}.\par Marimo: \path{marimo_notebooks/solvable_finite.py}.}

\subsection*{Scope and outcome}

A group name alone does not determine a $\sigma$-MGF: the series depends on a
specified complex representation.  We prove the formulas for two distinct
finite-lamplighter representations, split metacyclic induced modules, finite
affine and Frobenius permutation actions, unitriangular, Borel, and parabolic
actions, finite defining subgroups of $SU(3)$, $SU(4)$, $Sp(4)$, and $SO(4)$,
and symplectic-reflection wreath products.  The accompanying marimo notebook
constructs every representative matrix group and compares direct Reynolds
averages with weighted-cycle, fixed-power, class-spectrum, or wreath-recursion
formulas.  Dimensions range from $2$ to $27$; the audit performs $7{,}580$
fixed-power checks and $70$ coefficient comparisons, all passing.  The finite
theorems, computational checks, and asymptotic statements are kept logically
separate.

\subsection{The representation-dependent target}

Let $\rho:G\to GL(W)$ be a finite complex representation.  For a partition
or composition $\tau=(\tau_1,\ldots,\tau_s)$, put
\[
 W_\tau=\bigotimes_{j=1}^s\Sym^{\tau_j}W,
 \qquad
 a_\tau(G,W)=\dim W_\tau^G.
\]
The normalization used here is
\[
 \MGF_{G,W}(\mathbf t)=\sum_\tau a_\tau(G,W)m_\tau(\mathbf t).
\]

\begin{theorem}[Imported Reynolds--Molien identity]\label{nd:solvable_finite:thm:universal}
For every finite complex representation,
\begin{align}
 \MGF_{G,W}(\mathbf t)
 &=\frac1{|G|}\sum_{g\in G}\prod_{i\geq1}
       \det(I-t_i\rho(g))^{-1},\label{nd:solvable_finite:eq:molien}\\
 &=\frac1{|G|}\sum_{g\in G}
 \exp\!\left(\sum_{r\geq1}\frac{p_r(\mathbf t)}r\chi_W(g^r)\right),
 \label{nd:solvable_finite:eq:power-trace}\\
 [m_\tau]\MGF_{G,W}&=a_\tau(G,W).\label{nd:solvable_finite:eq:coefficient}
\end{align}
\end{theorem}

\begin{proof}[Source]
Equations~\eqref{nd:solvable_finite:eq:molien} and
\eqref{nd:solvable_finite:eq:coefficient} are Tools~\ref{tool:molien} and
\ref{tool:coefficient}; the power-trace form is the standard logarithmic
determinant expansion.
\end{proof}

If $G$ acts on a finite set $X$, we always mean the honest complex
permutation module $W=\C[X]$.  Write
\[
 F_r(g)=|\Fix_X(g^r)|,
 \qquad
 c_d(g)=\#\{\text{$d$-cycles of $g$ on $X$}\}.
\]

\begin{corollary}[permutation form]\label{nd:solvable_finite:cor:permutation}
For every finite permutation action,
\begin{align}
 F_r(g)&=\sum_{d\mid r}d c_d(g),&
 c_d(g)&=\frac1d\sum_{r\mid d}\mu(d/r)F_r(g),\label{nd:solvable_finite:eq:mobius}\\
 \MGF_{G,X}&=\frac1{|G|}\sum_{g\in G}
 \prod_{i,d\geq1}(1-t_i^d)^{-c_d(g)},\label{nd:solvable_finite:eq:perm-molien}\\
 [m_\tau]\MGF_{G,X}
 &=\#\left(G\backslash\prod_j\Sym^{\tau_j}X\right).\label{nd:solvable_finite:eq:orbit}
\end{align}
In particular, $[m_{(1^s)}]\MGF=\E F_1(g)^s=\#(G\backslash X^s)$.
\end{corollary}

\begin{proof}
A point in a $d$-cycle is fixed by $g^r$ exactly when $d\mid r$; M\"obius
inversion gives \eqref{nd:solvable_finite:eq:mobius}.  A $d$-cycle has determinant factor
$1-t^d$, which proves \eqref{nd:solvable_finite:eq:perm-molien}.  Monomials in
$\Sym^q\C[X]$ are indexed by $q$-element multisets, and orbit sums form a
basis for the invariant space, proving \eqref{nd:solvable_finite:eq:orbit}.
\end{proof}

\subsection{Finite lamplighter groups}

\subsubsection{Natural monomial representation}

Put $L_{r,n}=\mu_r^n\rtimes C_n$, where the generator $P$ cyclically shifts
the coordinates, and use the $n$-dimensional representation
\[
 (z_1,\ldots,z_n;P^s)\longmapsto
 \operatorname{diag}(z_1,\ldots,z_n)P^s.
\]
Define
\[
 A_{r,\ell}(\mathbf t)=\frac1r\sum_{\zeta^r=1}
 \prod_{i\geq1}(1-\zeta t_i^\ell)^{-1}.
\]
Root-of-unity filtering also gives
\[
 A_{r,\ell}=\sum_{q\geq0}h_{qr}(t_1^\ell,t_2^\ell,\ldots).
\]

\begin{theorem}[lamplighter monomial formula]\label{nd:solvable_finite:thm:lamp-monomial}
\[
 \boxed{\MGF^{\mathrm{mon}}_{L_{r,n}}
 =\frac1n\sum_{\ell\mid n}\varphi(\ell)
 A_{r,\ell}^{\,n/\ell}.}
\]
\end{theorem}

\begin{proof}
The shift $P^s$ has $d=\gcd(n,s)$ cycles, each of length
$\ell=n/d$.  On one coordinate cycle $C$, weighted permutation elimination
gives
\[
 \det(I-t\operatorname{diag}(z)P^s\mid_C)
 =1-\left(\prod_{j\in C}z_j\right)t^\ell.
\]
When the labels are averaged independently, every cycle product is an
independent uniform element of $\mu_r$.  Thus a shift of cycle length $\ell$
contributes $A_{r,\ell}^{n/\ell}$.  Exactly $\varphi(\ell)$ shifts have this
cycle length.
\end{proof}

There is no raw coefficientwise limit with $r$ fixed and $n\to\infty$.  The
identity-shift sector has weight $1/n$ and contributes $A_{r,1}^n/n$.
Choosing the $h_r$ term from two factors gives
\[
 [m_{(r,r)}]\MGF^{\mathrm{mon}}_{L_{r,n}}
 \geq \frac1n\binom n2[m_{(r,r)}]h_r^2
 =\frac{(r+1)(n-1)}2.
\]

\subsubsection{Affine action on lamp configurations}

The same abstract group acts on $X=C_r^n$ by $x\mapsto b+P^sx$.  For
$g=(b,P^s)$ set $b_k=(I+P^s+\cdots+P^{(k-1)s})b$.

\begin{proposition}[lamp fixed powers and pair coefficient]
\begin{align}
 F_k(g)&=\begin{cases}
 r^{\gcd(n,ks)},&b_k\in\im(I-P^{ks}),\\
 0,&\text{otherwise},
 \end{cases}\label{nd:solvable_finite:eq:lamp-affine-fixed}\\
 [m_{(1,1)}]\MGF^{\mathrm{aff}}_{L_{r,n}}
 &=\boxed{\frac1n\sum_{d\mid n}\varphi(d)r^{n/d}}.\label{nd:solvable_finite:eq:necklace}
\end{align}
\end{proposition}

\begin{proof}
Iteration gives $g^k(x)=b_k+P^{ks}x$.  Hence fixed points solve
$(I-P^{ks})x=b_k$.  A soluble fiber has the size of the kernel, namely
$r^{\gcd(n,ks)}$, proving \eqref{nd:solvable_finite:eq:lamp-affine-fixed}.  For ordered pairs,
translate the first entry to $0$.  The remaining difference is acted on only
by cyclic shifts.  Burnside's lemma gives the necklace count
\eqref{nd:solvable_finite:eq:necklace}.
\end{proof}

For fixed $r\geq2$, the necklace count is asymptotic to $r^n/n$; thus this
representation has exponential degree-two growth.  This does not contradict
Theorem~\ref{nd:solvable_finite:thm:lamp-monomial}: the two series use different representations.

\subsection{Split metacyclic induced representations}

Let
\[
 M(m,n,u)=\langle a,b:a^m=b^n=1,\ bab^{-1}=a^u\rangle,
 \qquad u^n\equiv1\pmod m,
\]
and suppose the chosen character has a full orbit of length $n$.  With
$\zeta=e^{2\pi i/m}$, define
\[
 \rho(a)e_j=\zeta^{u^{-j}}e_j,
 \qquad \rho(b)e_j=e_{j+1}.
\]
For $g=a^xb^y$, let $\mathcal C_y$ be the cycles of addition by $y$ on
$\mathbb Z/n\mathbb Z$, and put
\[
 \alpha_C(x,y)=\zeta^{x\sum_{j\in C}u^{-j}}.
\]

\begin{theorem}[metacyclic weighted-cycle formula]\label{nd:solvable_finite:thm:metacyclic}
\[
 \boxed{\MGF_{M(m,n,u),\rho}
 =\frac1{mn}\sum_{x=0}^{m-1}\sum_{y=0}^{n-1}
 \prod_{i\geq1}\prod_{C\in\mathcal C_y}
 (1-\alpha_C(x,y)t_i^{|C|})^{-1}.}
\]
\end{theorem}

\begin{proof}
The relations hold because shifting the diagonal weights sends
$u^{-j}$ to $u^{-(j-1)}=u\,u^{-j}$.  On a cycle $C$ of the permutation
$j\mapsto j+y$, the product of the diagonal labels is $\alpha_C(x,y)$.
Weighted permutation elimination gives the characteristic factor
$1-\alpha_C(x,y)t^{|C|}$.  Multiplication over cycles and Reynolds averaging
prove the result.
\end{proof}

If the character orbit has length $d<n$, the same proof applies to its
$d$-dimensional induced module.  A nonsplit relation $b^n=a^c$ inserts the
corresponding scalar in cycles crossing that relation; it must not be silently
omitted.  The arithmetic quantities $\gcd(n,y)$ and
$\sum_{j\in C}u^{-j}\pmod m$ prevent a group-independent asymptotic law.

\subsection{General affine and Frobenius permutation actions}

\subsubsection{Affine groups}

Let $V=\F_q^d$, $H\leq GL(V)$, and let $V\rtimes H$ act naturally on $V$.
For $g=(b,h)$ set $b_r=(I+h+\cdots+h^{r-1})b$.

\begin{theorem}[affine fixed powers and moments]\label{nd:solvable_finite:thm:affine}
\begin{align}
 F_r(b,h)&=\begin{cases}
 q^{\dim\ker(I-h^r)},&b_r\in\im(I-h^r),\\
 0,&\text{otherwise},
 \end{cases}\label{nd:solvable_finite:eq:affine-fixed}\\
 [m_{(1^s)}]\MGF_{V\rtimes H,V}
 &=\#(H\backslash V^{s-1}).\label{nd:solvable_finite:eq:affine-moment}
\end{align}
For $H=GL_d(q)$ this specializes to
\begin{equation}
 \boxed{[m_{(1^s)}]\MGF
 =\sum_{j=0}^{\min(d,s-1)}\qbinom{s-1}{j}.}\label{nd:solvable_finite:eq:gaussian}
\end{equation}
\end{theorem}

\begin{proof}
Iteration gives $g^r(x)=b_r+h^rx$, and solving the affine linear equation
gives \eqref{nd:solvable_finite:eq:affine-fixed}.  Translation sends the first point of an
$s$-tuple to zero and leaves the diagonal $H$-action on its $s-1$
differences, proving \eqref{nd:solvable_finite:eq:affine-moment}.  Regard an element of
$V^{s-1}$ as a $d$-by-$(s-1)$ matrix.  Two matrices lie in the same
$GL_d(q)$ orbit under left multiplication exactly when they have the same row
space.  The possible row spaces are the $j$-subspaces of $\F_q^{s-1}$ with
$j\leq d$, proving \eqref{nd:solvable_finite:eq:gaussian}.
\end{proof}

For fixed $q$ and fixed $s$, the right side stabilizes once $d\geq s-1$.
For general $H_d$, coefficientwise convergence is equivalent to stabilization
of the $H_d$-orbit counts on every fixed number of vectors.

\subsubsection{Finite Frobenius groups}

Let $G=K\rtimes H$ be a Frobenius group in its natural action on
$G/H\simeq K$, and put $N=|K|$.  Define
\[
 E_d=\prod_i(1-t_i^d)^{-N/d},\qquad
 R_d=\prod_i(1-t_i)^{-1}(1-t_i^d)^{-(N-1)/d}.
\]

\begin{theorem}[Frobenius cycle inventory]\label{nd:solvable_finite:thm:frobenius}
\begin{align}
 \MGF_{G,K}
 &=\boxed{\frac1{N|H|}\left[E_1+
 \sum_{1\ne k\in K}E_{\operatorname{ord}(k)}+
 N\sum_{1\ne h\in H}R_{\operatorname{ord}(h)}\right]},\label{nd:solvable_finite:eq:frob-molien}\\
 [m_{(1,1)}]\MGF_{G,K}
 &=\boxed{1+\frac{N-1}{|H|}}.\label{nd:solvable_finite:eq:frob-pair}
\end{align}
\end{theorem}

\begin{proof}
A nonidentity $k\in K$ is a fixed-point-free translation; if its order is
$d$, its cycles have type $d^{N/d}$.  Every element outside $K$ lies in a
unique conjugate of a complement and fixes one point.  Its nontrivial powers
fix the same point, so an element of order $d$ has cycle type
$1^1d^{(N-1)/d}$.  The $N$ complements have disjoint nonidentity parts,
which proves \eqref{nd:solvable_finite:eq:frob-molien}.  After translating the first member of
an ordered pair to $1\in K$, the complement acts freely on $K\setminus\{1\}$.
Its orbit count proves \eqref{nd:solvable_finite:eq:frob-pair}.
\end{proof}

With $|H|$ fixed and $N\to\infty$, the pair coefficient grows linearly.
Sharper transitivity can keep degree two bounded, but higher orbit parameters
still require separate control.

\subsection{Unitriangular, Borel, and parabolic actions}

Finite-field matrices are not complex matrices until a complex
representation is chosen.  Here the choices are permutation modules on
$V_n=\F_q^n$, $V_n^\times$, projective points, and flag varieties.

\begin{proposition}[vector and projective fixed powers]\label{nd:solvable_finite:prop:linear-fixed}
For $g\in GL_n(q)$,
\begin{align}
 F_r(g;V_n)&=q^{\dim\ker(g^r-I)},\label{nd:solvable_finite:eq:vector-fixed}\\
 F_r(g;V_n^\times)&=q^{\dim\ker(g^r-I)}-1,\label{nd:solvable_finite:eq:nonzero-fixed}\\
 F_r(g;\mathbb P^{n-1})
 &=\sum_{\lambda\in\F_q^\times}
 \frac{q^{\dim\ker(g^r-\lambda I)}-1}{q-1}.\label{nd:solvable_finite:eq:projective-fixed}
\end{align}
\end{proposition}

\begin{proof}
The first two statements count fixed vectors, with or without zero.  A
projective point is fixed exactly when its vectors form an eigenline.  The
nonzero vectors in distinct eigenspaces are disjoint, and every line has
$q-1$ representatives, giving \eqref{nd:solvable_finite:eq:projective-fixed}.
\end{proof}

For the upper-unitriangular and upper-triangular Borel groups, elementary row
operations give
\begin{align}
 \#(UT_n(q)\backslash V_n^\times)&=n(q-1),&
 \#(UT_n(q)\backslash\mathbb P^{n-1})&=n,\label{nd:solvable_finite:eq:ut-orbits}\\
 \#(B_n(q)\backslash V_n^\times)&=n,&
 \#(B_n(q)\backslash\mathbb P^{n-1})&=n.\label{nd:solvable_finite:eq:b-orbits}
\end{align}
For $UT_n(q)$ a nonzero vector is classified by the position and value of its
last nonzero coordinate; projectivization removes the value.  The diagonal
torus in $B_n(q)$ removes it before projectivization.  On all of $V_n$, add
one orbit for zero.  Consequently these rank-growing actions already diverge
in degree one.

Let $G$ be a finite reductive group and let $P,Q$ be parabolics with Weyl
subgroups $W_P,W_Q$.  Bruhat decomposition and the coset fixed-point count
give
\begin{align}
 [m_{(1)}]\MGF_{P,G/Q}&=|P\backslash G/Q|=|W_P\backslash W/W_Q|,
 \label{nd:solvable_finite:eq:parabolic-orbits}\\
 |\Fix_{G/Q}(p^r)|&=
 \frac{|C_G(p^r)|\,|(p^r)^G\cap Q|}{|Q|}.\label{nd:solvable_finite:eq:parabolic-fixed}
\end{align}
For $P=Q=B_n(q)$, \eqref{nd:solvable_finite:eq:parabolic-orbits} is $|W|=n!$.
Equation~\eqref{nd:solvable_finite:eq:parabolic-fixed} follows by counting the $x\in G$ for
which $x^{-1}p^rx\in Q$, then dividing by $|Q|$.

\subsection{Finite defining subgroups of unitary and symplectic groups}

\subsubsection{Class spectra and scalar selection}

For $G\leq U(d)$ let $\lambda_1(g),\ldots,\lambda_d(g)$ be the eigenvalues
of the defining matrix.

\begin{proposition}[finite unitary class formula]\label{nd:solvable_finite:prop:unitary-class}
\begin{align}
 \MGF_{G,\C^d}
 &=\boxed{\sum_{[g]}\frac1{|C_G(g)|}
 \prod_{i\geq1}\prod_{j=1}^d(1-\lambda_j(g)t_i)^{-1}},
 \label{nd:solvable_finite:eq:unitary-class}\\
 &=\sum_{[g]}\frac1{|C_G(g)|}
 \exp\!\left(\sum_{r\geq1}\frac{p_r}{r}\chi(g^r)\right).
 \label{nd:solvable_finite:eq:unitary-power}
\end{align}
Thus a character table determines the complete series only together with its
class power maps.  If $G$ contains a scalar $\zeta I$ of order $s$, then
\[
 [m_\tau]\MGF_G=0\quad\text{unless}\quad s\mid|\tau|.
 \label{nd:solvable_finite:eq:scalar-selection}
\]
\end{proposition}

\begin{proof}
Group the average in Theorem~\ref{nd:solvable_finite:thm:universal} by conjugacy classes to get
\eqref{nd:solvable_finite:eq:unitary-class}; the logarithmic determinant identity gives
\eqref{nd:solvable_finite:eq:unitary-power}.  The scalar acts on $W_\tau$ as
$\zeta^{|\tau|}$, so a fixed vector requires that scalar to be one.
\end{proof}

\subsubsection{General monomial theorem and the SU(3) type-C and type-D series}

Let $G=D\rtimes P\leq U(d)$ with $D$ diagonal and $P\leq S_d$.  For a
cycle $C$ of $\pi\in P$, write $z_C=\prod_{j\in C}z_j$.

\begin{theorem}[monomial class formula]\label{nd:solvable_finite:thm:monomial}
\[
 \boxed{\MGF_{D\rtimes P}
 =\frac1{|D||P|}\sum_{\pi\in P}\sum_{z\in D}
 \prod_{i\geq1}\prod_{C\in\Cyc(\pi)}
 (1-z_Ct_i^{|C|})^{-1}.}
\]
\end{theorem}

\begin{proof}
Weighted permutation elimination gives
$\det(I-tD(z)P_\pi)=\prod_C(1-z_Ct^{|C|})$.  Theorem~\ref{nd:solvable_finite:thm:universal}
then applies.  This includes the determinant-one monomial type-C and type-D
families in $SU(3)$.
\end{proof}

If $D_m$ is the $m$-torsion in a fixed compact diagonal torus $T$ normalized
by $P$, then in total degree at most $D$ only characters of bounded exponent
occur.  For all sufficiently large $m$, such a character is trivial on
$D_m$ exactly when it is trivial on $T$.  Hence
\[
 \MGF_{D_m\rtimes P}\longrightarrow\MGF_{T\rtimes P}
 \quad\text{coefficientwise}.
\]
For the full coordinate torus and the determinant-one torus, $m>D$ suffices.

\subsubsection{Two exceptional SU(3) class formulas}

Write
\[
 \mathcal D(\alpha,\beta,\gamma)
 =\prod_{i\geq1}\frac1{(1-\alpha t_i)(1-\beta t_i)(1-\gamma t_i)}.
\]
For the $3$-dimensional rotation representation of
$\Sigma(60)\cong A_5$,
\begin{align}
 \MGF_{\Sigma(60)}=\frac1{60}\bigl[&
 \mathcal D(1,1,1)+15\mathcal D(1,-1,-1)
 +20\mathcal D(1,\omega,\omega^2)\notag\\
 &+12\mathcal D(1,\zeta_5,\zeta_5^4)
 +12\mathcal D(1,\zeta_5^2,\zeta_5^3)\bigr].
 \label{nd:solvable_finite:eq:a5-su3}
\end{align}
For either conjugate $3$-dimensional representation of
$\Sigma(168)\cong PSL_2(7)$,
\begin{align}
 \MGF_{\Sigma(168)}=\frac1{168}\bigl[&
 \mathcal D(1,1,1)+21\mathcal D(1,-1,-1)
 +56\mathcal D(1,\omega,\omega^2)
 +42\mathcal D(1,i,-i)\notag\\
 &+24\mathcal D(\zeta_7,\zeta_7^2,\zeta_7^4)
 +24\mathcal D(\zeta_7^3,\zeta_7^5,\zeta_7^6)\bigr].
 \label{nd:solvable_finite:eq:psl-su3}
\end{align}
Both identities are immediate from Proposition~\ref{nd:solvable_finite:prop:unitary-class} and
the displayed class spectra.  Isolated exceptional groups have no intrinsic
rank asymptotic.  A sequence whose scalar-center order tends to infinity has
series tending coefficientwise to $1$ by \eqref{nd:solvable_finite:eq:scalar-selection}.

\subsubsection{SU(4), Sp(4), and SO(4)}

For $G\leq SU(4)$, Proposition~\ref{nd:solvable_finite:prop:unitary-class} with its defining
character $\chi_4$ is the complete answer.  For $G\leq Sp(4)$, the spectrum
is $\{\lambda,\lambda^{-1},\mu,\mu^{-1}\}$ and hence
\[
 \MGF_G=\sum_{[g]}\frac1{|C_G(g)|}\prod_i
 \frac1{(1-\lambda_gt_i)(1-\lambda_g^{-1}t_i)
 (1-\mu_gt_i)(1-\mu_g^{-1}t_i)}.
\]
If $-I\in G$, all odd-total-degree coefficients vanish.

For $SO(4)$ use $\operatorname{Spin}(4)\simeq SU(2)_L\times SU(2)_R$.
If $\widetilde G$ is the inverse image of $G$, the complexified vector
representation is the tensor product of the two defining planes, so
\[
 \boxed{\MGF_{G,\C^4}
 =\frac1{|\widetilde G|}\sum_{(a,b)\in\widetilde G}
 \exp\!\left(\sum_{r\geq1}\frac{p_r}{r}
 \chi_2(a^r)\chi_2(b^r)\right).}
 \label{nd:solvable_finite:eq:so4-spin}
\]
The kernel element $(-I,-I)$ acts trivially, so averaging over the spin cover
equals averaging over the represented image.

\subsection{Symplectic-reflection wreath products}

Let $K\leq SL_2(\C)$ act on its defining symplectic plane $E$, and let
$G_n=K\wr S_n$ act on $E^{\oplus n}$.  Define
\[
 A_\ell(\mathbf t)=\frac1{|K|}\sum_{k\in K}
 \prod_{i\geq1}\det(I-t_i^\ell k)^{-1}.
\]

\begin{theorem}[wreath exponential and stable law]\label{nd:solvable_finite:thm:wreath}
\begin{align}
 \sum_{n\geq0}u^n\MGF_{K\wr S_n,E^{\oplus n}}
 &=\boxed{\exp\!\left(\sum_{\ell\geq1}\frac{u^\ell}{\ell}A_\ell\right)},
 \label{nd:solvable_finite:eq:wreath-exponential}\\
 \lim_{n\to\infty}\MGF_{K\wr S_n,E^{\oplus n}}
 &=\boxed{\exp\!\left(\sum_{\ell\geq1}\frac{A_\ell-1}{\ell}\right)
 =\ExpS(\MGF_K-1)}.
 \label{nd:solvable_finite:eq:wreath-limit}
\end{align}
The limit is coefficientwise.
\end{theorem}

\begin{proof}
For a permutation cycle $C$ of length $\ell$, multiply its $K$-labels around
the cycle to obtain $k_C$.  Block elimination gives
$\det(I-tg_C)=\det(I-t^\ell k_C)$.  The cycle product is uniform in $K$, and
the exponential formula for labelled permutation cycles gives
\eqref{nd:solvable_finite:eq:wreath-exponential}.  Split $A_\ell=1+(A_\ell-1)$ to write the
right side as
\[
 \frac1{1-u}\exp\!\left(\sum_{\ell\geq1}
 \frac{u^\ell(A_\ell-1)}\ell\right).
\]
In a fixed positive $\mathbf t$-degree only finitely many $\ell$ occur, and
the second factor has bounded $u$-degree after that coefficient is extracted.
Multiplication by $(1-u)^{-1}$ makes the coefficients eventually constant,
proving \eqref{nd:solvable_finite:eq:wreath-limit}.
\end{proof}

This is a genuine marked compound-Poisson law, unlike the rare-sector
divergence of the lamplighter and fixed-complement Frobenius examples.

\subsection{Verification strategy and results}

\subsubsection{Independent computational routes}

The marimo notebook imports one pure verification module.  Every check uses
at least two of the following routes.

\begin{enumerate}
\item \textbf{Direct represented matrices.}  Dense complex matrices are used
for the monomial and defining representations.  A permutation matrix is
stored losslessly by its unique nonzero row in each column.
\item \textbf{Direct target quantity.}  For a matrix $A$, complete homogeneous
characters are computed from its eigenvalues, and
\[
 \frac1{|G|}\sum_g\prod_j h_{\tau_j}(\operatorname{Spec}\rho(g))
\]
is rounded only after its imaginary part and distance to the nearest integer
are checked.  Permutation coefficients are computed exactly from cycle
factors.  Ordered-pair and difference-tuple orbits are also enumerated
literally where applicable.
\item \textbf{Structural formula.}  Lamplighter and metacyclic cases use
root-filtered or label-product cycles; affine cases use image and kernel
tests; Frobenius cases use the predicted full cycle inventory; flag actions
use centralizers and class intersections; unitary cases use class spectra;
and wreath products use the coefficient recurrence obtained by differentiating
\eqref{nd:solvable_finite:eq:wreath-exponential}.
\item \textbf{Analytic spot check.}  The full determinant products are
evaluated at $(t_1,t_2)=(0.031,-0.047)$.  This checks all degrees at once and
is separate from the coefficient table.
\end{enumerate}

\subsubsection{Constructed groups and reasonable dimensions}

\begin{longtable}{@{}p{0.34\linewidth}rrrrp{0.15\linewidth}@{}}
\toprule
Representation & dim. & group/cover & fixed checks & coeff. checks & result\\
\midrule
\endhead
$L_{2,3},L_{3,2},L_{2,4}$, monomial & $2$--$4$ & $18$--$64$ & -- &15&\textcolor{teal}{\textbf{PASS}}\\
$L_{2,3},L_{2,4},L_{3,3}$, affine lamps &$8$--$27$&$24$--$81$&823&--&\textcolor{teal}{\textbf{PASS}}\\
$M(5,2,4),M(7,3,2)$, induced &$2$--$3$&$10,21$&relations&10&\textcolor{teal}{\textbf{PASS}}\\
$AGL(2,2),AGL(3,2)$, natural &$4,8$&$24,1344$&6634&5&\textcolor{teal}{\textbf{PASS}}\\
$\F_5\rtimes\F_5^\times$, $D_{14}$ &$5,7$&$20,14$&cycles&10&\textcolor{teal}{\textbf{PASS}}\\
$UT_3(2),B_2(3)$, vectors/projective &$7$--$8$&$8,12$&123&--&\textcolor{teal}{\textbf{PASS}}\\
$B_3(2)$ on complete flags &$21$&$8$&class powers&--&\textcolor{teal}{\textbf{PASS}}\\
$A_5,PSL_2(7),\Delta(12)$ in $SU(3)$ &$3$&$60,168,12$&relations&15&\textcolor{teal}{\textbf{PASS}}\\
$A_5$ in $SU(4)$ &$4$&$60$&classes&5&\textcolor{teal}{\textbf{PASS}}\\
$Q_8\wr S_2$ in $Sp(4)$ &$4$&$128$&center&5&\textcolor{teal}{\textbf{PASS}}\\
$(Q_8\times Q_8)/\{\pm(I,I)\}$ in $SO(4)$ &$4$&$64/32$&power traces&5&\textcolor{teal}{\textbf{PASS}}\\
\midrule
\textbf{Total}&&&\textbf{$7{,}580$}&\textbf{$70$}&\textcolor{teal}{\textbf{PASS}}\\
\bottomrule
\end{longtable}

For the $SU(3)$ exceptional cases the notebook constructs actual generators.
For $PSL_2(7)$, with $\eta=e^{2\pi i/7}$, it uses
\[
 A=\operatorname{diag}(\eta,\eta^2,\eta^4),\qquad
 B=\frac{i}{\sqrt7}
 \begin{pmatrix}
 \eta^2-\eta^5&\eta-\eta^6&\eta^4-\eta^3\\
 \eta-\eta^6&\eta^4-\eta^3&\eta^2-\eta^5\\
 \eta^4-\eta^3&\eta^2-\eta^5&\eta-\eta^6
 \end{pmatrix}.
\]
Closure produces $168$ distinct unitary determinant-one matrices with order
inventory $1,21,56,42,48$ for element orders $1,2,3,4,7$.  The maximum
unitarity and determinant errors are below $5.4\times10^{-14}$.

The $SO(4)$ test constructs all $64$ matrices $a\otimes b$ with
$a,b\in Q_8\leq SU(2)$.  Deduplication leaves $32$ represented matrices,
exactly detecting the kernel $\{(I,I),(-I,-I)\}$.  For powers $1\leq r\leq4$
it verifies
\[
 \operatorname{Tr}((a\otimes b)^r)
 =\operatorname{Tr}(a^r)\operatorname{Tr}(b^r)
\]
to machine precision, and the cover and image Molien values agree within
$6.7\times10^{-16}$.

\subsubsection{Representative coefficient values}

\begin{center}
\small
\begin{tabular}{@{}lrrrrr@{}}
\toprule
Representation & $m_1$ & $m_2$ & $m_{11}$ & $m_3$ & $m_{21}$\\
\midrule
$L_{2,3}$ monomial &0&1&1&0&0\\
$L_{3,2}$ monomial &0&0&0&1&1\\
$M(5,2,4)$ induced &0&1&1&0&0\\
$M(7,3,2)$ induced &0&0&0&1&1\\
$\F_5\rtimes\F_5^\times$ natural &1&2&2&3&5\\
$D_{14}$ natural &1&4&4&8&16\\
$A_5\leq SU(3)$ &0&1&1&0&0\\
$PSL_2(7)\leq SU(3)$ &0&0&0&0&0\\
$\Delta(12)\leq SU(3)$ &0&1&1&1&1\\
$A_5\leq SU(4)$ &0&1&1&1&1\\
$Q_8\wr S_2\leq Sp(4)$ &0&0&1&0&0\\
$(Q_8\times Q_8)/\pm\leq SO(4)$ &0&1&1&0&0\\
\bottomrule
\end{tabular}
\end{center}

The affine-lamp pair coefficients are $4,6,11$ for $(r,n)=(2,3),(2,4),(3,3)$,
matching literal pair orbits and the necklace formula.  For $AGL(2,2)$, the
moments $[m_{(1^s)}]$ for $s=2,3,4$ are $2,5,15$; for $AGL(3,2)$ the tested
values for $s=2,3$ are $2,5$.  In each case direct difference-tuple orbits,
affine Burnside averages, and Gaussian-binomial sums agree.

\subsection{Asymptotic summary and scope}

\begin{center}
\small
\begin{tabularx}{\linewidth}{@{}p{0.43\linewidth}X@{}}
\toprule
Family and representation & large-parameter behavior\\
\midrule
$C_r^n\rtimes C_n$, natural monomial
& $[m_{(r,r)}]\gg n$; no raw coefficientwise limit.\\
$C_r^n\rtimes C_n$, affine lamps
& $[m_{11}]\sim r^n/n$; no raw coefficientwise limit.\\
Split metacyclic induced modules
& arithmetic and subsequence dependent.\\
$V_d\rtimes GL_d(q)$, $q$ fixed
& every fixed $(1^s)$ moment stabilizes once $d\geq s-1$.\\
Frobenius groups with fixed complement
& $[m_{11}]=1+(|K|-1)/|H|$ grows linearly.\\
$UT_n(q),B_n(q)$ on nonzero/projective points
& $[m_1]$ grows linearly in $n$.\\
$B_n(q)$ on complete flags
& $[m_1]=n!$.\\
$SU(3)$ diagonal torsion series
& coefficientwise torus-normalizer limit as torsion becomes dense.\\
Exceptional fixed-dimensional unitary groups
& isolated exact class sums; no intrinsic rank limit.\\
$K\wr S_n\leq Sp(2n)$
& $\ExpS(\MGF_K-1)$, a genuine coefficientwise limit.\\
\bottomrule
\end{tabularx}
\end{center}

A full crystallographic group is generally infinite and noncompact; it has no
uniform probability measure and therefore no $\sigma$-MGF of the present
kind.  One must first choose a finite quotient, a compact quotient action, or
a nonuniform probability law.  A finite linear point group preserving a
complex lattice is covered by Proposition~\ref{nd:solvable_finite:prop:unitary-class}, or by
Theorem~\ref{nd:solvable_finite:thm:monomial} when it is monomial.

\resultbox{\textbf{Reproducibility.} Run
\texttt{python -m new\_dists.solvable\_finite\_subgroup\_sigma\_mgfs.verification}
for the noninteractive audit, or open
\path{marimo_notebooks/solvable_finite.py} with marimo for the
group-organized laboratory.  The source keeps the direct matrix
route and each family-specific formula in separate functions so that agreement
cannot result from comparing an expression with itself.}

\clearpage
\section[The H3 reflection group]{The $H_3$ reflection group}\label{proof:h3}
\proofstatus{Formula exact; reconstruction diagnostic}{The Molien formula
is exact once the represented group is specified. Floating-point closure and
rank tests do not certify the algebraic reflection group; an exact
$\mathbb Q(\sqrt5)$ construction or certified character table is required for
a fully exact computational certificate.}

\begingroup
\definecolor{ink}{HTML}{211B32}
\definecolor{accent}{HTML}{6D3B72}
\definecolor{pale}{HTML}{F6F0F7}
\definecolor{passgreen}{HTML}{16794A}
\providecommand{\Sfun}{\SMG}
\renewcommand{\Sfun}{\SMG}
\providecommand{\Sym}{\operatorname{Sym}}
\renewcommand{\Sym}{\operatorname{Sym}}
\providecommand{\PASS}{\textcolor{passgreen}{\textbf{PASS}}}
\renewcommand{\PASS}{\textcolor{passgreen}{\textbf{PASS}}}
\providecommand{\summarybox}[1]{\begin{center}\fcolorbox{accent}{pale}{\parbox{0.91\linewidth}{#1}}\end{center}}
\renewcommand{\summarybox}[1]{\begin{center}\fcolorbox{accent}{pale}{\parbox{0.91\linewidth}{#1}}\end{center}}
\setcounter{theorem}{0}
\setcounter{proposition}{0}
\setcounter{lemma}{0}
\setcounter{corollary}{0}
\setcounter{remark}{0}
\summarybox{\textbf{Outcome.}  Explicit rotation matrices verify the proposed
$A_4$, $S_4$, and $A_5$ angle formulas in 45 multivariate coefficients each.
The 120 matrices of $W(H_3)$ were reconstructed from the icosahedron; all 139
coefficients through total degree ten agree with the signed class-spectrum
formula, and the one-variable series agrees with
$((1-t^2)(1-t^6)(1-t^{10}))^{-1}$.}

\subsection{Imported target}
The target is $\smgv{G}{V}$ from Tools~\ref{tool:coefficient} and
\ref{tool:molien}.  The calculation below does not reprove those identities:
it constructs the $H_3$ matrices independently and compares their invariant
coefficients with a weighted spectral computation.

\subsection{Floating-point reconstruction diagnostic for $W(H_3)$}

Let $\varphi=(1+\sqrt5)/2$.  Use the twelve vertices
\begin{equation}
\Omega=\{(0,\pm1,\pm\varphi),(\pm1,\pm\varphi,0),
(\pm\varphi,0,\pm1)\}.                                      \tag{3}
\end{equation}
Choose an ordered, linearly independent base triple $B$ of vertices.  For
every ordered target triple $T$ having the same Gram matrix, form the unique
linear map $R=TB^{-1}$.  Retain $R$ when
\[
R^{\mathsf T}R=I,\qquad\det R=1,\qquad R\Omega=\Omega.
\]
This produces 60 distinct rotations.  Adjoining their negatives produces 120
orthogonal transformations.  This is the full symmetry group of the
icosahedron, which is the real reflection group $W(H_3)$.

This ``vertex-frame reconstruction'' is useful when no computer algebra class
table is available, but the present implementation uses floating-point
coordinates and therefore supplies a diagnostic rather than an exact group
certificate.  The implementation reports its equality tolerance, the minimum
separation between reconstructed matrices, the maximum closure-matching
distance, and the condition number of the chosen base triple; it then tests
every one of the $120^2$ products against the reconstructed set.  No closure
claim here is described as exact.

The regenerated diagnostic reports equality tolerance $5\times10^{-10}$,
minimum separation $1.6625$ between distinct reconstructed matrices, maximum
closure-matching distance $3.77\times10^{-16}$, base-triple condition number
$1.852$, and maximum orthogonality residual $4.16\times10^{-16}$.

\subsection{Polyhedral rotation groups}

For a three-dimensional rotation through angle $\theta$, define
\[
\Phi_\theta(\mathbf t)=\prod_{a\geq1}
\big((1-t_a)(1-2\cos\theta\,t_a+t_a^2)\big)^{-1}.
\]
The spectrum is $(1,e^{i\theta},e^{-i\theta})$, so direct substitution of the
rotation counts gives
\begin{align*}
\Sfun_{A_4,\mathrm{rot}}&=\frac1{12}
(\Phi_0+3\Phi_\pi+8\Phi_{2\pi/3}),\\
\Sfun_{S_4,\mathrm{rot}}&=\frac1{24}
(\Phi_0+6\Phi_{\pi/2}+8\Phi_{2\pi/3}+9\Phi_\pi),\\
\Sfun_{A_5,\mathrm{rot}}&=\frac1{60}
(\Phi_0+15\Phi_\pi+20\Phi_{2\pi/3}
+12\Phi_{2\pi/5}+12\Phi_{4\pi/5}).
\end{align*}
The direct route independently constructs the 12 tetrahedral matrices as the
determinant-one signed permutation matrices preserving a fixed tetrahedron,
the 24 octahedral matrices as all determinant-one signed permutation matrices,
and the 60 icosahedral matrices by the vertex-frame construction above.  For
each group, all 45 coefficients indexed by partitions through total degree
seven agree with the displayed angle formula.

\begin{center}
\begin{tabular}{@{}lrrc@{}}
\toprule
group & order & coefficient checks & result\\
\midrule
$A_4$ tetrahedral rotations & 12 & 45 & \PASS\\
$S_4$ octahedral rotations & 24 & 45 & \PASS\\
$A_5$ icosahedral rotations & 60 & 45 & \PASS\\
\bottomrule
\end{tabular}
\end{center}

\subsection{The explicit $H_3$ spectral formula}

Write
\[
A_z(\mathbf t)=\prod_{a\geq1}(1-zt_a)^{-1},\qquad
\omega=e^{2\pi i/3},\qquad \xi=e^{2\pi i/5}.
\]
The orientation-preserving subgroup has the following spectra and class sizes:
\begin{center}
\begin{tabular}{@{}clc@{}}
\toprule
size & eigenvalues & geometric type\\
\midrule
1  & $(1,1,1)$ & identity\\
15 & $(1,-1,-1)$ & half turns\\
20 & $(1,\omega,\omega^2)$ & third turns\\
12 & $(1,\xi,\xi^{-1})$ & first fifth-turn class\\
12 & $(1,\xi^2,\xi^{-2})$ & second fifth-turn class\\
\bottomrule
\end{tabular}
\end{center}
Multiplication by central inversion negates all three eigenvalues and supplies
the other five classes.  Grouping (1) by these ten spectra gives
\begin{align}
\Sfun_{H_3}=\frac1{120}\Big[&A_1^3+A_{-1}^3
+15(A_1A_{-1}^2+A_{-1}A_1^2)\notag\\
&+20(A_1A_\omega A_{\omega^2}
+A_{-1}A_{-\omega}A_{-\omega^2})\notag\\
&+12\sum_{j=1}^{2}(A_1A_{\xi^j}A_{\xi^{-j}}
+A_{-1}A_{-\xi^j}A_{-\xi^{-j}})\Big].                       \tag{4}
\end{align}
Equation (4) is exact in all multidegrees.  It is more informative than the
ordinary invariant degrees because mixed coefficients depend on the entire
class spectrum.

\subsubsection{Ordinary Molien specialization}

Setting only $t_1=t$ nonzero gives
\begin{equation}
\Sfun_{H_3}(t,0,\ldots)=
\frac1{(1-t^2)(1-t^6)(1-t^{10})}.                             \tag{5}
\end{equation}
The right side has degree-$d$ coefficient equal to the number of nonnegative
solutions of $2a+6b+10c=d$.  This elementary count is used as a third route,
independent of both explicit matrix spectra and the multivariate class
calculation.

\subsection{Verification protocol}

\begin{enumerate}[leftmargin=*,itemsep=3pt]
\item Reconstruct the 60 rotations from (3), adjoin negatives, and test order,
orthogonality, signed determinants, and closure.
\item For every partition $\tau$ through total degree ten, compute the direct
matrix average in (2) over all 120 matrices.
\item Compute the same coefficient using only the ten weighted spectra in (4).
The explicit matrices are not passed to this route.
\item For every $0\leq d\leq10$, compare the direct one-variable average with
the solution count for $2a+6b+10c=d$.
\item Construct the six-dimensional $\Sym^2(V)$ matrices and average them to a
Reynolds projector.
\end{enumerate}

\subsection{Results}

\begin{center}
\begin{tabularx}{\linewidth}{@{}XrX@{}}
\toprule
Check & observed & status\\
\midrule
number of reconstructed matrices & 120 & \PASS\\
closure products tested & 14,400 & \PASS\\
maximum $\|g^{\mathsf T}g-I\|_F$ & $4.16\times10^{-16}$ & \PASS\\
signed determinant condition & all matrices & \PASS\\
partitions through degree 10 & 139 coefficients & \PASS\\
ordinary Molien degrees $0$ through $10$ & 11 coefficients & \PASS\\
\bottomrule
\end{tabularx}
\end{center}

Some matched multivariate coefficients are
\begin{center}
\begin{tabular}{@{}lrrrrrrr@{}}
\toprule
$\tau$ & $(1)$ & $(2)$ & $(3)$ & $(4)$ & $(2,1)$ & $(2,2)$ & $(3,2)$\\
\midrule
dimension & 0 & 1 & 0 & 1 & 0 & 2 & 0\\
\bottomrule
\end{tabular}
\end{center}
The coefficients of (5) for $d=0,\ldots,10$ are
\[
1,0,1,0,1,0,2,0,2,0,3,
\]
and agree termwise with direct matrix averaging.

For $\tau=(2)$, the target $\Sym^2(V)$ has dimension six.  The explicit
Reynolds projector has
\[
\operatorname{rank}P=1,\qquad \operatorname{tr}P=0.9999999999999998,
\qquad \|P^2-P\|_F=1.26\times10^{-16}.
\]
The invariant is the expected quadratic form.

\subsection{Reproducibility and limitations}

Run from the repository root:
{\small
\begin{verbatim}
.venv/bin/python -m for_this_guy.matrix_group_formula_verification.h3_reflection.verification
\end{verbatim}
}
The companion marimo notebook exposes the maximum tested degree and displays
all coefficient rows.

The group reconstruction uses double-precision values of $\varphi$, not an
exact quadratic-number field.  This does not affect the theoretical proof of
(4), but it makes the group diagnostics numerical.  To control this, vertices
and matrices are canonicalized to nine decimal places, every product is tested,
and all averaged characters must lie within $5\times10^{-7}$ of an integer.
The actual orthogonality and projector residuals are around $10^{-16}$.

The finite checks are diagnostic evidence; the all-degree validity of (4)
comes from the master identity and the complete class spectra.  The computation
specifically verifies that the class weights, signs, fifth-root classes, and
matrix convention have been transcribed consistently.
\endgroup


\clearpage
\section[The M24 permutation representation]{The $M_{24}$ permutation representation}\label{proof:m24}
\proofstatus{Conditional on certified class data}{The universal class/power or
frame-shape formula is exact. Concrete numerical tables in this section depend
on external class sizes, power maps, character values, or frame shapes and are
not independently derived by the Molien argument.}

\begingroup
\definecolor{navy}{HTML}{17365D}
\definecolor{teal}{HTML}{147D78}
\providecommand{\Sgm}{\SMG}
\renewcommand{\Sgm}{\SMG}
\setcounter{theorem}{0}
\setcounter{proposition}{0}
\setcounter{lemma}{0}
\setcounter{corollary}{0}
\setcounter{remark}{0}
\begin{abstract}
We prove the cycle-shape formula for the sigma--MGF of the natural
$24$-dimensional permutation representation of $M_{24}$.  Two published
ATLAS permutations are converted into $24\times24$ matrices.  Independent
generator-orbit enumeration is compared with a full conjugacy-class sum.
Agreement holds in every tested case.  The computation also sharpens the
5-transitivity statement: equality with $S_{24}$ holds through total degree
five, while $[m_{(6)}]$ is $12$ rather than $11$ and
$[m_{(1^6)}]$ is $204$ rather than $203$.
\end{abstract}

\subsection{Exact formula}
Let $g$ have cycle shape $\prod_{r\geq1}r^{c_r(g)}$.  On an $r$-cycle the
eigenvalues are all $r$th roots of unity, so
\[
\det(I-tP_r)=\prod_{j=0}^{r-1}(1-t\zeta_r^j)=1-t^r.
\]
Multiplying over cycles and averaging the generalized Molien integrand gives
\begin{equation}\label{m24-permutation-proof-and-verification:eq:molien}
\boxed{
\Sgm_{M_{24},\mathrm{perm}}(\mathbf t)
=\sum_C\frac{|C|}{|M_{24}|}
\prod_{i\geq1}\prod_{r\geq1}(1-t_i^r)^{-c_r(C)}.}
\end{equation}
This is an all-degree identity, not an asymptotic expansion.

The coefficient of $m_\tau$ is
\[
\dim\left(\bigotimes_j\operatorname{Sym}^{\tau_j}\mathbb C^{24}\right)^{M_{24}}.
\]
Each symmetric power has a basis indexed by multisets.  A permutation-group
fixed vector is constant on each orbit of the resulting basis.  Therefore
\begin{equation}\label{m24-permutation-proof-and-verification:eq:orbits}
[m_\tau]\Sgm_{M_{24},\mathrm{perm}}
=\#\left(M_{24}\backslash\prod_j
\operatorname{MSet}_{\tau_j}(\{1,\ldots,24\})\right).
\end{equation}
Equations \eqref{m24-permutation-proof-and-verification:eq:molien} and \eqref{m24-permutation-proof-and-verification:eq:orbits} are mathematically
equivalent, but they produce genuinely different algorithms: weighted class
expansion versus graph traversal under generators.  Their agreement is the
central cross-check.

\subsection{Why equality with $S_{24}$ stops after degree five}
An $S_{24}$ orbit of a product of multisets is determined by the multiplicity
profile carried by its support points.  If $|\tau|\leq5$, that support has at
most five points.  The 5-transitivity of $M_{24}$ maps any ordered listing of
one support to the corresponding listing of another, so every $S_{24}$ orbit
remains a single $M_{24}$ orbit.  Consequently
\[
\Sgm_{M_{24},\mathrm{perm}}\equiv
\Sgm_{S_{24},\mathrm{perm}}\pmod{\deg>5}.
\]
For one symmetric power, the $S_{24}$ coefficient is the partition number
$p(k)$; for $(1^k)$ it is the Bell number $B_k$ (provided $k\leq24$).
The claim is not extended above degree five.  Indeed, generator traversal
finds two orbits on six-element subsets.  This splits the all-distinct
equality pattern and creates the first extra orbit.

\subsection{Constructed matrices and checks}
The verifier uses the ATLAS standard generators
\begin{align*}
a={}&(1,4)(2,7)(3,17)(5,13)(6,9)(8,15)(10,19)(11,18)\\
&\quad(12,21)(14,16)(20,24)(22,23),\\
b={}&(1,4,6)(2,21,14)(3,9,15)(5,18,10)(13,17,16)(19,24,23).
\end{align*}
They are converted to permutation matrices by $P_ge_j=e_{g(j)}$.  The code
asserts the standard-generator fingerprints
\[
o(a)=2,\quad o(b)=3,\quad o(ab)=23,\quad
o(abababbababbabb)=4,
\]
and checks that the aggregated class sizes sum to
$|M_{24}|=244{,}823{,}040$.

For a class shape, the coefficient of $t^k$ in
$\prod_r(1-t^r)^{-c_r}$ is obtained with exact integer convolution.  In
parallel, breadth-first traversal acts with $a^{\pm1},b^{\pm1}$ on every
multiset through degree six and every ordered tuple through degree four.  At
degree six the tensor calculation uses equality patterns together with the
directly enumerated two orbits of six-element subsets.

\begin{center}
\begin{tabular}{@{}lrrr@{}}\toprule
coefficient & direct orbits & class formula & $S_{24}$\\\midrule
$m_{(1)}$&1&1&1\\
$m_{(2)}$&2&2&2\\
$m_{(3)}$&3&3&3\\
$m_{(4)}$&5&5&5\\
$m_{(5)}$&7&7&7\\
$m_{(6)}$&12&12&11\\\addlinespace
$m_{(1^2)}$&2&2&2\\
$m_{(1^3)}$&5&5&5\\
$m_{(1^4)}$&15&15&15\\
$m_{(1^6)}$&204&204&203\\\bottomrule
\end{tabular}
\end{center}

The degree-one invariant is the all-ones line.  The degree-six discrepancies
are a useful regression test: a program that silently substitutes
5-transitivity beyond its valid range returns the wrong values $11$ and
$203$.

\subsection{Reproduction and provenance}
Run
\begin{verbatim}
.venv/bin/python for_this_guy/sporadic_matrix_groups/
  m24_permutation/verification.py
.venv/bin/marimo edit marimo_notebooks/m24_permutation.py
\end{verbatim}
after joining wrapped lines.  The full cycle-shape/class-size table is stored
as exact integers in \path{sporadic_matrix_groups/common.py}; classes with
the same permutation cycle shape are aggregated because \eqref{m24-permutation-proof-and-verification:eq:molien}
cannot distinguish them.

\begin{center}\scriptsize
\begin{tabular}{@{}lrl@{}}\toprule
class & aggregated size & permutation cycle shape\\\midrule
$1A$&1&$1^{24}$\\
$2A$&11,385&$1^8 2^8$\\
$2B$&31,878&$2^{12}$\\
$3A$&226,688&$1^6 3^6$\\
$3B$&485,760&$3^8$\\
$4A$&637,560&$2^4 4^4$\\
$4B$&1,912,680&$1^4 2^2 4^4$\\
$4C$&2,550,240&$4^6$\\
$5A$&4,080,384&$1^4 5^4$\\
$6A$&10,200,960&$1^2 2^2 3^2 6^2$\\
$6B$&10,200,960&$6^4$\\
$7AB$&11,658,240&$1^3 7^3$\\
$8A$&15,301,440&$1^2 2 4 8^2$\\
$10A$&12,241,152&$2^2 10^2$\\
$11A$&22,256,640&$1^2 11^2$\\
$12A$&20,401,920&$2 4 6 12$\\
$12B$&20,401,920&$12^2$\\
$14AB$&34,974,720&$1 2 7 14$\\
$15AB$&32,643,072&$1 3 5 15$\\
$21AB$&23,316,480&$3 21$\\
$23AB$&21,288,960&$1 23$\\\bottomrule
\end{tabular}
\end{center}
The sizes sum to $244{,}823{,}040$.  Source: Online ATLAS of Finite Group
Representations, retrieved 18 July 2026.  The SHA-256 of the canonical
serialized table is
\begin{center}\footnotesize
\texttt{103ba576555f01bea74df73f8e2f788b02bee4e8b20df3ed198eb005d16de044}.
\end{center}
The SHA-256 of the software data file is
\begin{center}\footnotesize
\texttt{3fbc8e5cde4ac5670696ecc1c5f692a2afb2065469a83ce486aa8b28d046d579}.
\end{center}

Generator-orbit enumeration is independent of the class table.  The class-sum
calculation is an exact cross-check conditional on the externally supplied
ATLAS data; the complete table and checksums make that input auditable from
this PDF.

Generator data and the group order are from the online ATLAS of Finite Group
Representations, \url{https://brauer.maths.qmul.ac.uk/Atlas/spor/M24/}.
The calculation verifies the stated representation and coefficients without
enumerating all $244$ million group elements.
\endgroup


\clearpage
\section[The Co0 Leech-space representation]{The $Co_0$ Leech-space representation}\label{proof:co0}
\proofstatus{Conditional on certified class data}{The universal class/power or
frame-shape formula is exact. Concrete numerical tables in this section depend
on external class sizes, power maps, character values, or frame shapes and are
not independently derived by the Molien argument.}

\begingroup
\definecolor{navy}{HTML}{17365D}
\definecolor{teal}{HTML}{147D78}
\providecommand{\Sgm}{\SMG}
\renewcommand{\Sgm}{\SMG}
\setcounter{theorem}{0}
\setcounter{proposition}{0}
\setcounter{lemma}{0}
\setcounter{corollary}{0}
\setcounter{remark}{0}
\begin{abstract}
We prove the frame-shape class formula for the natural $24$-dimensional
complexified Leech-lattice representation of $Co_0$ and the vanishing of all
odd total-degree coefficients.  Since direct enumeration of $Co_0$ is not a
reasonable notebook computation, we test the determinant mechanism at the
correct dimension on the explicit signed-coordinate subgroup
$\{\pm P_g:g\in M_{24}\}$.  Direct invariant counts and exact class expansion
agree through degree six.  The distinction between this stress test and a
full $Co_0$ class-table evaluation is kept explicit.
\end{abstract}

\subsection{Frame-shape formula}
Let $V=\Lambda_{\rm Leech}\otimes_{\mathbb Z}\mathbb C$.  If a conjugacy-class
representative $g\in Co_0$ has frame shape
$\pi_g=\prod_{r\geq1}r^{k_r(g)}$, by definition
\[
\det(xI-g)=\prod_{r\geq1}(x^r-1)^{k_r(g)},\qquad
\sum_r r k_r(g)=24.
\]
Substitute $x=t^{-1}$ and multiply by $t^{24}$ to obtain
\[
\det(I-tg)=\prod_{r\geq1}(1-t^r)^{k_r(g)}.
\]
The generalized Molien formula, grouped by conjugacy class, is therefore
\begin{equation}\label{co0-leech-proof-and-verification:eq:co0}
\boxed{\Sgm_{Co_0,V}(\mathbf t)=
\sum_{C\subseteq Co_0}\frac{|C|}{|Co_0|}
\prod_{i,r}(1-t_i^r)^{-k_r(C)}.}
\end{equation}
Negative frame exponents cause no difficulty: each class contribution is a
rational function whose product is the characteristic determinant.

\subsection{Central parity theorem}
The central element $-I$ belongs to $Co_0$.  On
$\operatorname{Sym}^{\tau}V=\bigotimes_j\operatorname{Sym}^{\tau_j}V$ it
acts by $(-1)^{|\tau|}$.  A fixed vector in odd total degree would satisfy
$v=-v$, hence vanish.  Thus
\begin{equation}\label{co0-leech-proof-and-verification:eq:parity}
\boxed{[m_\tau]\Sgm_{Co_0,V}=0\quad\text{whenever }|\tau|\text{ is odd}.}
\end{equation}
In particular $[m_{(1)}]=0$, unlike the $M_{24}$ permutation representation,
whose all-ones vector is fixed.  The two 24-dimensional series therefore
differ in their first positive degree.

This proof is representation theoretic and does not require any frame-shape
table.  As a data-integrity check on a full table, the class contributions in
every odd degree must cancel exactly.

\subsection{A tractable subgroup at the true dimension}
In the standard coordinate construction of the Leech lattice, its monomial
automorphism group contains $M_{24}$ coordinate permutations and the central
sign.  We use
\[
H=\{\pm P_g:g\in M_{24}\}\leq Co_0.
\]
The notebook explicitly constructs three $24\times24$ integer orthogonal
generators $P_a,P_b,-I$, using the ATLAS permutations $a,b$ documented in the
companion $M_{24}$ note.  This is large enough to exercise the actual ambient
dimension while retaining an exact aggregated class computation.

If $g$ has cycle counts $c_r(g)$, then
\[
\det(I-tP_g)=\prod_r(1-t^r)^{c_r(g)},\qquad
\det(I+tP_g)=\prod_r(1-(-t)^r)^{c_r(g)}.
\]
Averaging the two signs projects onto even total degree.  Hence
\begin{equation}\label{co0-leech-proof-and-verification:eq:signed}
\Sgm_H(t)=\frac1{2|M_{24}|}\sum_{g\in M_{24}}
\left\{\prod_r(1-t^r)^{-c_r(g)}+
\prod_r(1-(-t)^r)^{-c_r(g)}\right\}.
\end{equation}

\subsection{Independent verification}
For even $k$, $-I$ acts trivially on $\operatorname{Sym}^kV$, so the $H$-fixed
space has a basis indexed by $M_{24}$ orbits of degree-$k$ multisets.  For odd
$k$, \eqref{co0-leech-proof-and-verification:eq:parity} makes the dimension zero.  This gives a direct
generator-orbit algorithm independent of the class expansion
\eqref{co0-leech-proof-and-verification:eq:signed}.  Exact results are:
\begin{center}
\begin{tabular}{@{}crrc@{}}\toprule
$k$ & direct invariant dimension & signed class formula & parity\\\midrule
1&0&0&odd\\
2&2&2&even\\
3&0&0&odd\\
4&5&5&even\\
5&0&0&odd\\
6&12&12&even\\\bottomrule
\end{tabular}
\end{center}
The degree-six value also exercises the first point at which the underlying
$M_{24}$ multiset orbits differ from those of $S_{24}$.

\subsection{What this does and does not verify}
The proof of \eqref{co0-leech-proof-and-verification:eq:co0} is exact for full $Co_0$, and the proof of
\eqref{co0-leech-proof-and-verification:eq:parity} establishes every odd-degree zero for full $Co_0$.  The
matrix sweep validates determinant expansion, sign cancellation, and direct
invariant counting in dimension $24$ on a genuine subgroup.  It does not
claim that the even coefficients $2,5,12$ are full-$Co_0$ coefficients; fixed
spaces can shrink when the acting group grows.  A full numerical evaluation
of \eqref{co0-leech-proof-and-verification:eq:co0} requires the complete $Co_0$ frame-shape and class-size
table, which is deliberately not reconstructed from partial data.

Run
\begin{verbatim}
.venv/bin/python for_this_guy/sporadic_matrix_groups/
  co0_leech/verification.py
.venv/bin/marimo edit marimo_notebooks/co0_leech.py
\end{verbatim}
after joining wrapped lines.  For background on the Leech representation and
Conway fixed lattices, see \url{https://arxiv.org/abs/1505.06420}.
\endgroup


\clearpage
\section{Sporadic groups and represented covers}\label{proof:new-dist:sporadic}
\proofstatus{Conditional on certified class data}{The universal class/power or
frame-shape formula is exact. Concrete numerical tables in this section depend
on external class sizes, power maps, character values, or frame shapes and are
not independently derived by the Molien argument.}


\resultbox{\textbf{Canonical testing artifacts.}\par\smallskip Exact backend: \path{lambda_distributions/dists/sporadic_group_sigma_mgfs/verification.py}.\par Regression: \path{lambda_distributions/dists/sporadic_group_sigma_mgfs/test_verification.py}.\par Marimo: \path{marimo_notebooks/sporadic.py}.}

\subsection{A representation is part of the data}

A sporadic abstract group does not by itself determine a sigma--MGF.  Fix a
finite-dimensional complex representation
\[
\rho:G\longrightarrow GL(V),\qquad \chi(g)=\operatorname{Tr}\rho(g).
\]
For formal variables $\mathbf t=(t_1,t_2,\ldots)$ and power sums
$p_r=\sum_i t_i^r$, define
\begin{equation}\label{nd:sporadic:eq:def}
\MGF_{G,V}(\mathbf t)=\frac1{|G|}\sum_{g\in G}
\prod_{i\geq1}\det(I-t_i\rho(g))^{-1}.
\end{equation}
Different representations of the same group generally give different
functions.  This note groups together two representations arising from the
same $M_{11}$ action: the full permutation module and its deleted constituent.

\subsection{Universal class-power theorem}

Let $x_1,\ldots,x_s$ represent the conjugacy classes of $G$, and put
$z_a=|C_G(x_a)|$.  Then
\begin{equation}\label{nd:sporadic:eq:universal}
\boxed{
\MGF_{G,V}(\mathbf t)
=\sum_{a=1}^s\frac1{z_a}
\exp\!\left(\sum_{r\geq1}\frac{p_r}{r}\chi(x_a^r)\right).}
\end{equation}
For a partition $\tau=(\tau_1,\ldots,\tau_\ell)$, with
$m_\tau=t_1^{\tau_1}\cdots t_\ell^{\tau_\ell}$ before symmetrization,
\begin{equation}\label{nd:sporadic:eq:coefficient}
\boxed{[m_\tau]\MGF_{G,V}
=\dim\left(\bigotimes_{j=1}^{\ell}\Sym^{\tau_j}V\right)^G.}
\end{equation}
Thus class centralizers, class power maps, and one ordinary character suffice
to determine the complete sigma--MGF.

\subsubsection*{Source of the universal identities}
Equation~\eqref{nd:sporadic:eq:coefficient} is Tool~\ref{tool:coefficient},
and the determinant form underlying
\eqref{nd:sporadic:eq:universal} is Tool~\ref{tool:molien}.  The displayed
class-power form is obtained only by the standard logarithmic determinant
change of coordinates and grouping the imported average by conjugacy class.
Thus the sporadic input is precisely the certified class sizes, power maps,
and character values.

\subsection{The \texorpdfstring{$M_{11}$}{M11} matrix group}

The online ATLAS gives the following standard generators in the natural
11-point action:
\begin{align*}
a&=(2,10)(4,11)(5,7)(8,9),\\
b&=(1,4,3,8)(2,5,6,9).
\end{align*}
The verifier closes these permutations under multiplication and inversion;
it obtains exactly $7{,}920$ elements.  With $P_ge_j=e_{g(j)}$, this constructs
all $11\times11$ zero-one matrices $P_g$ rather than merely storing abstract
class data.  The standard-generator fingerprint is
\[
o(a)=2,\qquad o(b)=4,\qquad o(ab)=11,\qquad
o(ababbabbb)=5.
\]

For the deleted representation, let
\[
W=\{x\in\mathbb C^{11}:x_1+\cdots+x_{11}=0\},\qquad
f_j=e_j-e_{11}\quad(1\leq j\leq10).
\]
The code constructs the restriction $D_g=P_g|_W$ in this basis.  Explicitly,
\begin{equation}\label{nd:sporadic:eq:deleted-matrix}
(D_g)_{ij}=\delta_{i,g(j)}-\delta_{i,g(11)}
\quad(1\leq i,j\leq10),
\end{equation}
where a delta is zero if its second index is $11$.  Every entry is integral.
All $7{,}920$ matrices are distinct in each representation, and sampled
products satisfy $P_gP_h=P_{gh}$ and $D_gD_h=D_{gh}$ exactly.

\subsection{Exact formulas for the two \texorpdfstring{$M_{11}$}{M11} representations}

If $\lambda=1^{c_1}2^{c_2}\cdots$ is a permutation cycle shape, set
\[
Z_\lambda(\mathbf t)=\prod_{i\geq1}\prod_{r\geq1}
(1-t_i^r)^{-c_r},\qquad
\widetilde Z_\lambda(\mathbf t)=\prod_{i\geq1}(1-t_i)Z_\lambda(\mathbf t).
\]
The second identity follows from
$\mathbb C^{11}\cong\mathbf1\oplus W$.  Direct enumeration gives the eight
aggregated spectra below; the two order-eight classes and two order-eleven
classes may be combined because their spectra in these representations agree.

\begin{center}
\begin{tabular}{@{}lrrl@{}}
\toprule
class row & size & centralizer & cycle shape\\
\midrule
$1A$ & 1 & 7,920 & $1^{11}$\\
$2A$ & 165 & 48 & $1^3 2^4$\\
$3A$ & 440 & 18 & $1^2 3^3$\\
$4A$ & 990 & 8 & $1^3 4^2$\\
$5A$ & 1,584 & 5 & $1\,5^2$\\
$6A$ & 1,320 & 6 & $2\,3\,6$\\
$8AB$ & 1,980 & $8$ in each class & $1\,2\,8$\\
$11AB$ & 1,440 & $11$ in each class & $11$\\
\bottomrule
\end{tabular}
\end{center}

Consequently the natural permutation formula is
\begin{equation}\label{nd:sporadic:eq:permformula}
\boxed{\begin{aligned}
\MGF_{M_{11},\mathbb C^{11}}=\frac1{7920}\big(&
Z_{1^{11}}+165Z_{1^3 2^4}+440Z_{1^2 3^3}
+990Z_{1^3 4^2}\\
&+1584Z_{1\,5^2}+1320Z_{2\,3\,6}
+1980Z_{1\,2\,8}+1440Z_{11}\big).
\end{aligned}}
\end{equation}
The deleted formula is the same weighted sum with every $Z$ replaced by
$\widetilde Z$:
\begin{equation}\label{nd:sporadic:eq:delformula}
\boxed{\MGF_{M_{11},W}
=\left(\prod_{i\geq1}(1-t_i)\right)
\MGF_{M_{11},\mathbb C^{11}}.}
\end{equation}
Equation \eqref{nd:sporadic:eq:delformula} is a relation between generalized Molien
integrands and hence between their averages; it does not assert a tensor
factorization of invariant spaces.

The same cycle data recover every character of a power:
\begin{equation}\label{nd:sporadic:eq:powertrace}
\chi_{\rm perm}(g^r)=\sum_{d\mid r}d\,c_d(g),\qquad
\chi_W(g^r)=\chi_{\rm perm}(g^r)-1.
\end{equation}
Substitution in \eqref{nd:sporadic:eq:universal} gives a third, power-character evaluation
of \eqref{nd:sporadic:eq:permformula}--\eqref{nd:sporadic:eq:delformula}.

\subsection{Verification strategy}

The program uses four routes with deliberately different inputs.
\begin{enumerate}[leftmargin=*,itemsep=3pt]
\item \textbf{Direct matrices.}  For every generated matrix $A$, traces of
$A^r$ are computed by integer matrix multiplication.  Complete symmetric
characters are recovered from Newton's exact recurrence
\[
h_0(A)=1,\qquad
h_d(A)=\frac1d\sum_{r=1}^d\operatorname{Tr}(A^r)h_{d-r}(A).
\]
The target coefficient is then the literal average
$|G|^{-1}\sum_g\prod_jh_{\tau_j}(\rho(g))$.
\item \textbf{Class determinants.}  Only the eight table rows are used to
expand $\prod_r(1-t^r)^{-c_r}$; the deleted factor $1-t$ is then applied.
\item \textbf{Power characters.}  Equation \eqref{nd:sporadic:eq:powertrace} and the
same Newton recurrence expand the exponential formula, without using the
determinant factorization in the implementation.
\item \textbf{Configuration orbits.}  For selected permutation coefficients,
the two generators traverse the finite set
\[
\prod_j\MSet_{\tau_j}(\{1,\ldots,11\}).
\]
The number of orbits must equal \eqref{nd:sporadic:eq:coefficient}.  This route uses
neither the class table nor matrix traces.
\end{enumerate}

Finally, at three nonzero variables the code evaluates every determinant
numerically and compares the full average with the closed class formula.
This catches mistakes that could be hidden by a short coefficient truncation.

\subsection{Results}

All 29 partitions of positive total degree at most six were tested in each
representation, giving 58 exact three-way comparisons.  Representative rows
are shown here.

\begin{center}
\begin{tabular}{@{}llrrrc@{}}
\toprule
representation & $\tau$ & matrix & determinant & power character & result\\
\midrule
permutation & $(1)$ & 1 & 1 & 1 & \textcolor{teal}{\textbf{PASS}}\\
permutation & $(2)$ & 2 & 2 & 2 & \textcolor{teal}{\textbf{PASS}}\\
permutation & $(2,1)$ & 4 & 4 & 4 & \textcolor{teal}{\textbf{PASS}}\\
permutation & $(2,2)$ & 9 & 9 & 9 & \textcolor{teal}{\textbf{PASS}}\\
permutation & $(1,1,1,1)$ & 15 & 15 & 15 & \textcolor{teal}{\textbf{PASS}}\\
permutation & $(6)$ & 13 & 13 & 13 & \textcolor{teal}{\textbf{PASS}}\\
\addlinespace
deleted & $(1)$ & 0 & 0 & 0 & \textcolor{teal}{\textbf{PASS}}\\
deleted & $(2)$ & 1 & 1 & 1 & \textcolor{teal}{\textbf{PASS}}\\
deleted & $(2,1)$ & 1 & 1 & 1 & \textcolor{teal}{\textbf{PASS}}\\
deleted & $(2,2)$ & 3 & 3 & 3 & \textcolor{teal}{\textbf{PASS}}\\
deleted & $(1,1,1,1)$ & 4 & 4 & 4 & \textcolor{teal}{\textbf{PASS}}\\
deleted & $(6)$ & 5 & 5 & 5 & \textcolor{teal}{\textbf{PASS}}\\
\bottomrule
\end{tabular}
\end{center}

The separate generator-orbit counts for
$\tau=(1),(2),(1,1),(3),(2,1),(2,2),(1,1,1,1)$ are respectively
\[
1,\ 2,\ 2,\ 3,\ 4,\ 9,\ 15,
\]
and agree with the class coefficients.  In particular, the last value is the
Bell number $B_4=15$, as expected from the 4-transitivity of this action.

At $(t_1,t_2,t_3)=(0.037,-0.051,0.083)$, the full determinant averages are
\begin{center}
\begin{tabular}{@{}lrrr@{}}
\toprule
representation & direct matrix average & class formula & absolute error\\
\midrule
permutation & 1.0866701548337285 & 1.0866701548337498 & $2.13\times10^{-14}$\\
deleted & 1.0085468522143926 & 1.0085468522144532 & $6.06\times10^{-14}$\\
\bottomrule
\end{tabular}
\end{center}
The discrepancies are at floating-point roundoff scale; all coefficient and
group-law checks are exact integers.

\subsection{Other sporadic constructions and scope}

\subsubsection{Other ordinary representations}
For any of the 26 sporadic simple groups, and for any chosen ordinary complex
character, equation \eqref{nd:sporadic:eq:universal} applies verbatim.  For a nontrivial
irreducible representation of a sporadic simple group, the kernel is normal
and therefore trivial.  This observation does not select a representation;
one must still state which character is intended.

\subsubsection{Central covers}
If $\widetilde G$ is a central cover and a central element $z$ acts on $V$ as
the scalar $\omega(z)$, it acts on
$\bigotimes_j\Sym^{\tau_j}V$ as $\omega(z)^{|\tau|}$.  Hence
\[
[m_\tau]\MGF_{\widetilde G,V}=0
\quad\text{unless}\quad \omega(z)^{|\tau|}=1
\quad\text{for every }z\in Z(\widetilde G).
\]
Projective representations of the simple quotient must therefore be treated
as honest representations of the appropriate cover.

\subsubsection{Lattice and conjugation actions}
For a lattice action, class eigenvalues or frame shapes may be substituted
directly into \eqref{nd:sporadic:eq:def}.  For conjugation on the basis $\mathbb C[G]$,
the fixed-point character is $\chi(g^r)=|C_G(g^r)|$; M\"obius inversion then
recovers the conjugation cycle counts.  These are separate representation
choices and are not silently combined with the $M_{11}$ action tested here.

\subsubsection{Asymptotics}
There is no intrinsic limiting sequence over the sporadic simple groups:
there are only 26.  For the one-variable coefficients of either fixed faithful
$M_{11}$ representation, ordinary rational-function asymptotics do apply.
The identity matrix supplies the unique pole of order $\dim V$ at $t=1$, up
to root-of-unity oscillations of lower polynomial order, so
\begin{align*}
[t^d]\MGF_{M_{11},\mathbb C^{11}}(t)
&=\frac1{7920}\binom{d+10}{10}+O(d^9),\\
[t^d]\MGF_{M_{11},W}(t)
&=\frac1{7920}\binom{d+9}{9}+O(d^8).
\end{align*}
The exact sequences are quasipolynomial because every eigenvalue is a root of
unity.

\subsection{Reproduction and provenance}

From the repository root, run
\begin{verbatim}
.venv/bin/python -m lambda_distributions.dists.sporadic_group_sigma_mgfs.verification
.venv/bin/python -m pytest -q \
  lambda_distributions/dists/sporadic_group_sigma_mgfs/test_verification.py
.venv/bin/marimo edit marimo_notebooks/sporadic.py
\end{verbatim}
The generator data, group order, standard-generator fingerprint, representation
degree, and class centralizers are recorded by the online
\href{https://brauer.maths.qmul.ac.uk/Atlas/v3/spor/M11/}{ATLAS page for
$M_{11}$}; the natural action and its cycle types are recorded on the
\href{https://brauer.maths.qmul.ac.uk/Atlas/v3/permrep/M11G1-p11B0}{ATLAS
11-point representation page}.  The program independently enumerates the
group and reproduces every aggregated class multiplicity used in
\eqref{nd:sporadic:eq:permformula}.

This is a standalone inclusion-ready proof.  It does not modify or rebuild the
combined proof PDF.

\clearpage
\part{Nonuniform measures, heat flow, and limiting laws}
\section{Nonuniform laws on cyclic groups}\label{proof:nonuniform-cyclic}
\begingroup
\definecolor{ink}{HTML}{19332D}
\definecolor{accent}{HTML}{187C68}
\definecolor{pale}{HTML}{ECF8F4}
\definecolor{passgreen}{HTML}{16794A}
\providecommand{\Sfun}{\SMG}
\renewcommand{\Sfun}{\SMG}
\providecommand{\Sym}{\operatorname{Sym}}
\renewcommand{\Sym}{\operatorname{Sym}}
\providecommand{\Tr}{\operatorname{Tr}}
\renewcommand{\Tr}{\operatorname{Tr}}
\providecommand{\PASS}{\textcolor{passgreen}{\textbf{PASS}}}
\renewcommand{\PASS}{\textcolor{passgreen}{\textbf{PASS}}}
\providecommand{\summarybox}[1]{\begin{center}\fcolorbox{accent}{pale}{\parbox{0.91\linewidth}{#1}}\end{center}}
\renewcommand{\summarybox}[1]{\begin{center}\fcolorbox{accent}{pale}{\parbox{0.91\linewidth}{#1}}\end{center}}
\setcounter{theorem}{0}
\setcounter{proposition}{0}
\setcounter{lemma}{0}
\setcounter{corollary}{0}
\setcounter{remark}{0}
\summarybox{\textbf{Outcome.}  The averaging-operator and convolution formulas
were checked on explicit diagonal models of $C_3,C_4,C_5$ in dimensions
$1,2,3$.  Arbitrary nonuniform measures and four walk times were tested for
every partition through total degree five.  All \textbf{288} comparisons pass;
the maximum absolute error is $2.91\times10^{-14}$.  Three additional
unexpanded determinant evaluations also pass.}

\subsection{The general coefficient theorem}
Let $G$ be finite, $\rho:G\to U(V)$, and let $\nu$ be a probability measure.
For finitely many active variables define
\begin{equation}
\Sfun_{\nu,V}(\mathbf t)=\sum_{g\in G}\nu(g)
 \prod_a\det(I-t_a\rho(g))^{-1}.                         \tag{1}
\end{equation}
The standard identity
$\det(I-tA)^{-1}=\sum_{r\geq0}\chi_{\Sym^r V}(A)t^r$
implies, by multiplying the series variable by variable, that
\begin{equation}
[\mathbf t^\tau]\Sfun_{\nu,V}
=\sum_g\nu(g)\chi_{W_\tau}(g)
=\Tr\!\left(\sum_g\nu(g)\rho_{W_\tau}(g)\right),
\qquad W_\tau=\bigotimes_a\Sym^{\tau_a}V.                \tag{2}
\end{equation}
This proves the proposed replacement of the Haar invariant dimension by the
trace of a nonuniform averaging operator.  Positivity of the measure does not
force the trace to be real or integral; the $C_5$ test below deliberately
exhibits a complex coefficient.

If $\nu=\kappa^{*r}$ and the convention is that a step product is
$X_1\cdots X_r$, then representation multiplicativity gives
\begin{equation}
\widehat\nu(W)=\widehat\kappa(W)^r.                        \tag{3}
\end{equation}
Equations (2)--(3) prove the convolution formula without a centrality
assumption.

\subsection{Specialization to \texorpdfstring{$C_m$}{C-m}}
Write $C_m=\{0,\ldots,m-1\}$ additively and
$\zeta_m=e^{2\pi i/m}$.  A diagonal representation is determined by weights
$w_1,\ldots,w_d\pmod m$:
\begin{equation}
\rho(k)=\operatorname{diag}(\zeta_m^{kw_1},\ldots,
\zeta_m^{kw_d}).                                           \tag{4}
\end{equation}
Let $M_\tau(q)$ denote the multiplicity of the character $k\mapsto
\zeta_m^{kq}$ in $W_\tau$.  Extracting weights from
\begin{equation}
\prod_{j=1}^d(1-xz^{w_j})^{-1}
=\sum_{r\geq0}\sum_q M_{(r)}(q)x^rz^q                    \tag{5}
\end{equation}
and convolving the multiplicities for the factors in $W_\tau$ gives the
finite Fourier formula
\begin{equation}
[\mathbf t^\tau]\Sfun_{\nu,V}
=\sum_{q\bmod m}M_\tau(q)\widehat\nu(q),
\qquad
\widehat\nu(q)=\sum_{k=0}^{m-1}\nu(k)\zeta_m^{kq}.        \tag{6}
\end{equation}
For a walk this becomes
\begin{equation}
[\mathbf t^\tau]\Sfun_{\kappa^{*r},V}
=\sum_qM_\tau(q)\widehat\kappa(q)^r.                     \tag{7}
\end{equation}
Thus the abstract matrix Fourier transform reduces to a transparent scalar
filter on each weight.

\subsection{Independent verification routes}
The direct route builds every matrix in (4), computes the power traces
$p_j=\Tr(\rho(k)^j)$, and obtains complete characters from Newton's recurrence
\begin{equation}
h_0=1,\qquad h_r=\frac1r\sum_{j=1}^r p_jh_{r-j}.            \tag{8}
\end{equation}
It then averages $\prod_a h_{\tau_a}$ using the fully enumerated
probability distribution.  The formula route never uses matrix traces: it
constructs the integer weight multiplicities in (5) and applies (6) or (7).
For walks, the direct distribution is obtained by repeated cyclic convolution,
whereas the formula route raises Fourier scalars to the $r$th power.

As a full target check, the code also evaluates (1) at
$(t_1,t_2)=(0.07,-0.04)$ using matrix determinants and compares it with the
product over the weights in (4).  These values stay safely inside every pole.

\subsection{Cases and results}
The representations and the nonuniform arbitrary measure are
\begin{center}
\begin{tabular}{@{}cccc@{}}
\toprule
Group & weights & dimension & $\nu(k)$ before normalization\\
\midrule
$C_3$ & $(1)$ & 1 & $1,2,3$\\
$C_4$ & $(1,-1)$ & 2 & $1,2,3,4$\\
$C_5$ & $(0,1,2)$ & 3 & $1,2,3,4,5$\\
\bottomrule
\end{tabular}
\end{center}
The walk step is $0.35\delta_0+0.40\delta_1+0.25\delta_{-1}$ and is tested
at $r=0,1,2,5$.  There are 19 partitions through total degree five, hence
$3\cdot(1+4)\cdot19=285$ coefficient comparisons plus three determinant
comparisons.

\begin{center}
\begin{tabularx}{\linewidth}{@{}lXrrc@{}}
\toprule
Case & coefficient & direct value & formula value & result\\
\midrule
$C_3$, arbitrary & $\tau=(3)$ & $1.0000000000$ & $1.0000000000$ & \PASS\\
$C_4$, walk $r=5$ & $\tau=(2,1)$ & $-0.0210525000$ & $-0.0210525000$ & \PASS\\
$C_5$, arbitrary & $\tau=(2,2)$ & $2.1666667+0.2293970i$ & $2.1666667+0.2293970i$ & \PASS\\
\bottomrule
\end{tabularx}
\end{center}
The complex $C_5$ row is a useful diagnostic: the implementation is genuinely
testing non-Haar averaging, not silently rounding to invariant dimensions.

\subsection{Reproducibility and scope}
Run \path{cyclic_groups/verification.py}, or open
\path{marimo_notebooks/nonuniform_cyclic.py} in marimo.  The notebook exposes the
degree cutoff and every comparison.  The finite computations verify all
coefficients in the stated range; equations (1)--(7) are all-degree proofs.
No sampling is used.
\endgroup


\clearpage
\section{Nonuniform laws on symmetric groups}\label{proof:nonuniform-symmetric}
\begingroup
\definecolor{ink}{HTML}{252E43}
\definecolor{accent}{HTML}{4B63B7}
\definecolor{pale}{HTML}{EFF2FF}
\definecolor{passgreen}{HTML}{16794A}
\providecommand{\Sfun}{\SMG}
\renewcommand{\Sfun}{\SMG}
\providecommand{\Tr}{\operatorname{Tr}}
\renewcommand{\Tr}{\operatorname{Tr}}
\providecommand{\PASS}{\textcolor{passgreen}{\textbf{PASS}}}
\renewcommand{\PASS}{\textcolor{passgreen}{\textbf{PASS}}}
\providecommand{\summarybox}[1]{\begin{center}\fcolorbox{accent}{pale}{\parbox{0.91\linewidth}{#1}}\end{center}}
\renewcommand{\summarybox}[1]{\begin{center}\fcolorbox{accent}{pale}{\parbox{0.91\linewidth}{#1}}\end{center}}
\setcounter{theorem}{0}
\setcounter{proposition}{0}
\setcounter{lemma}{0}
\setcounter{corollary}{0}
\setcounter{remark}{0}
\summarybox{\textbf{Outcome.}  Explicit permutation matrices for
$S_3,S_4,S_5$ were compared with the cycle-count formula under six families
of nonuniform measures.  A separate $S_3$ calculation decomposes every tested
tensor-symmetric character into all three irreducibles and verifies both
central and noncentral Fourier operators.  All \textbf{267} checks pass; the
maximum absolute error is $4.45\times10^{-15}$.}

\subsection{Permutation matrices and the cycle transform}
For the defining permutation representation $V=\mathbb C^n$, a $d$-cycle
contributes all $d$th roots of unity.  Consequently
\begin{equation}
\det(I-tP_\pi)=\prod_{d\geq1}(1-t^d)^{C_d(\pi)}.            \tag{1}
\end{equation}
With
\begin{equation}
Q_d(\mathbf t)=\prod_a(1-t_a^d)^{-1},                      \tag{2}
\end{equation}
direct substitution in the determinant definition proves
\begin{equation}
\Sfun_{\nu,n}(\mathbf t)=
\mathbb E_\nu\!\left[\prod_{d\geq1}Q_d(\mathbf t)^{C_d(\pi)}\right]. \tag{3}
\end{equation}
This identity is valid for every probability measure, whether or not it is a
class measure.  It is the joint cycle-count probability-generating function
evaluated at the transformed variables $Q_d$.

For coefficient extraction, each factor is expanded using
\begin{equation}
(1-t^d)^{-c}=\sum_{j\geq0}\binom{c+j-1}{j}t^{dj}.           \tag{4}
\end{equation}
The verification uses integer dynamic programming for (4), separately for
every permutation.  This route does not diagonalize a matrix.

\subsection{Measures tested}
All probabilities are constructed on the fully enumerated group:
\begin{enumerate}[leftmargin=*,itemsep=3pt]
\item the lazy random-transposition walk with
$\kappa=0.35\delta_e+0.65\operatorname{Unif}(C_2)$, after three steps;
\item the lazy random $3$-cycle walk with laziness $0.4$, after two steps;
\item the lazy adjacent-transposition walk with laziness $0.25$, after four steps;
\item the inverse $2$-riffle law
\begin{equation}
\mathbb P(\pi)=2^{-n}\binom{n+1-d(\pi)}{n},                \tag{5}
\end{equation}
where $d(\pi)$ is the number of descents;
\item Ewens measure with $\theta=1.7$,
$\mathbb P(\pi)\propto\theta^{K(\pi)}$;
\item the general cycle weight
$\mathbb P(\pi)\propto1.2^{C_1}0.7^{C_2}1.6^{C_3}$.
\end{enumerate}
For $S_3,S_4,S_5$ and every partition through total degree four, the direct
route averages characters computed from the explicit matrices using Newton's
identities.  The cycle route applies (4) to the cycle structure.  Thus the two
routes share the probability law but not the spectral computation.

The normalizing generating-function formulas follow immediately from the
exponential formula for labeled cycles.  For example,
\begin{equation}
\Sfun_{n,\theta}^{\mathrm{Ewens}}=
\frac{[z^n]\exp\!\left(\theta\sum_{d\geq1}Q_dz^d/d\right)}
{[z^n]\exp\!\left(\theta\sum_{d\geq1}z^d/d\right)}.        \tag{6}
\end{equation}
Equation (6) records the corrected typography: the normalization is
$(\theta)_n$, and the summation index is $d\geq1$.

\subsection{Fourier proof and an independent \texorpdfstring{$S_3$}{S3} test}
Decompose
\begin{equation}
W_\tau=\bigotimes_a\operatorname{Sym}^{\tau_a}V
\cong\bigoplus_{\lambda\vdash n}M_{\tau,\lambda}\otimes V_\lambda,
\qquad m_{\tau,\lambda}=\dim M_{\tau,\lambda}.            \tag{7}
\end{equation}
Linearity of trace gives, for an arbitrary measure,
\begin{equation}
[\mathbf t^\tau]\Sfun_{\nu,n}
=\sum_{\lambda\vdash n}m_{\tau,\lambda}
\Tr\widehat\nu(\lambda),
\qquad
\widehat\nu(\lambda)=\sum_\pi\nu(\pi)\rho_\lambda(\pi). \tag{8}
\end{equation}
For a convolution walk, $\widehat{\kappa^{*r}}(\lambda)
=\widehat\kappa(\lambda)^r$.  If $\kappa$ is central, Schur's lemma makes
this operator scalar.  In the lazy transposition case,
\begin{equation}
\widehat\kappa(\lambda)=\beta_\lambda I,
\qquad
\beta_\lambda=a+(1-a)\frac{\chi^\lambda((12))}{d_\lambda}. \tag{9}
\end{equation}
For adjacent transpositions the full matrix power must be retained.

The $S_3$ test constructs the trivial, sign, and two-dimensional standard
matrices explicitly.  Multiplicities in (7) are obtained independently by
the uniform character inner product, then (8) is evaluated for an arbitrary
law, lazy central and adjacent walks, and the riffle law.  For laziness
$a=0.35$, the central scalars on $(\mathbf1,\mathrm{sgn},\mathrm{std})$ are
\begin{equation}
(1,-0.3,0.35).                                               \tag{10}
\end{equation}
The adjacent operator on the standard representation instead has two distinct
eigenvalues $0.025$ and $0.675$, directly exposing the claimed matrix-valued
behavior.

\subsection{Results}
There are 12 partitions through total degree four.  The cycle calculation
therefore contributes $3\cdot6\cdot12=216$ comparisons.  The four $S_3$
Fourier laws contribute 48 more, and the three irreducible operator checks
bring the total to 267.
\begin{center}
\begin{tabularx}{\linewidth}{@{}lXrrc@{}}
\toprule
Case & coefficient & matrix/direct & proposed formula & result\\
\midrule
$S_3$, inverse $2$-riffle & $(3)$ & $5.750000$ & $5.750000$ & \PASS\\
$S_4$, Ewens $\theta=1.7$ & $(2,1)$ & $7.429557$ & $7.429557$ & \PASS\\
$S_5$, transposition walk & $(2,2)$ & $39.683180$ & $39.683180$ & \PASS\\
$S_3$, arbitrary Fourier law & $(2,1)$ & $1.904762$ & $1.904762$ & \PASS\\
$S_3$, adjacent walk Fourier law & $(2,2)$ & $9.671793$ & $9.671793$ & \PASS\\
\bottomrule
\end{tabularx}
\end{center}

\subsection{Warnings, corrections, and scope}
Without laziness, a transposition walk alternates parity and does not converge
to uniform measure on $S_n$.  A walk on odd-length cycles stays in $A_n$.
These are support obstructions, not failures of (3) or (8).

The supplied continuous-time class-walk expressions contained a stray comma
in the exponential; the intended factor is
$\exp(t(r_\lambda-1))$.  Continuous-time graph interchange processes were not
numerically tested here, but their operator formula follows from replacing
matrix powers in (8) by the matrix exponential of the generator.

Run
\begin{verbatim}
.venv/bin/python -m \
  for_this_guy.nonuniform_lambda_distributions.symmetric_groups.verification
\end{verbatim}
or open \path{marimo_notebooks/nonuniform_symmetric.py} in marimo.
\endgroup


\clearpage
\section[U(1) heat kernels and weight attenuation]{$U(1)$ heat kernels and weight attenuation}\label{proof:u1-heat}
\begingroup
\definecolor{ink}{HTML}{3A293A}
\definecolor{accent}{HTML}{A44E82}
\definecolor{pale}{HTML}{FFF0F8}
\definecolor{passgreen}{HTML}{16794A}
\providecommand{\Sfun}{\SMG}
\renewcommand{\Sfun}{\SMG}
\providecommand{\PASS}{\textcolor{passgreen}{\textbf{PASS}}}
\renewcommand{\PASS}{\textcolor{passgreen}{\textbf{PASS}}}
\providecommand{\summarybox}[1]{\begin{center}\fcolorbox{accent}{pale}{\parbox{0.91\linewidth}{#1}}\end{center}}
\renewcommand{\summarybox}[1]{\begin{center}\fcolorbox{accent}{pale}{\parbox{0.91\linewidth}{#1}}\end{center}}
\setcounter{theorem}{0}
\setcounter{proposition}{0}
\setcounter{lemma}{0}
\setcounter{corollary}{0}
\setcounter{remark}{0}
\summarybox{\textbf{Outcome.}  Three diagonal $U(1)$ representations of
dimensions $1,2,3$ were checked at five heat times for all partitions through
total degree five.  Direct quadrature against the heat kernel agrees with the
weight/Casimir formula in all \textbf{285} cases.  The maximum absolute error
is $5.73\times10^{-14}$.}

\subsection{Heat measure and normalization}
Identify $U(1)$ with $e^{i\theta}$, $0\leq\theta<2\pi$, and choose the
Laplacian normalization for which the character $e^{iq\theta}$ has Casimir
$q^2$.  Brownian motion with generator $\tfrac12\Delta$ has density relative
to normalized Haar measure
\begin{equation}
k_t(\theta)=\sum_{q\in\mathbb Z}e^{-q^2t/2}e^{-iq\theta}.   \tag{1}
\end{equation}
Orthogonality immediately yields
\begin{equation}
\int_{U(1)}e^{iw\theta}k_t(\theta)\,d\theta/(2\pi)
=e^{-w^2t/2}.                                               \tag{2}
\end{equation}
This fixes the otherwise convention-dependent factor in the Casimir exponent.

Let $V$ be the diagonal representation with integer weights
$a_1,\ldots,a_d$,
\begin{equation}
\rho(e^{i\theta})=\operatorname{diag}
(e^{ia_1\theta},\ldots,e^{ia_d\theta}).                    \tag{3}
\end{equation}
For $W_\tau=\bigotimes_j\operatorname{Sym}^{\tau_j}V$, write
\begin{equation}
W_\tau=\bigoplus_{w\in\mathbb Z}M_\tau(w)\,\mathbb C_w.  \tag{4}
\end{equation}
Combining (2)--(4) proves the coefficient formula
\begin{equation}
\boxed{
[\mathbf t^\tau]\Sfun_t
=\sum_{w\in\mathbb Z}M_\tau(w)e^{-w^2t/2}.}               \tag{5}
\end{equation}
This is the compact-group heat formula in its simplest exact form: each
irreducible weight is attenuated by its Casimir eigenvalue.

\subsection{Endpoints and interpretation}
At $t=0$, every exponential in (5) is one, hence
\begin{equation}
[\mathbf t^\tau]\Sfun_0=\dim W_\tau,
\qquad
\Sfun_0(\mathbf t)=\prod_j(1-t_j)^{-d}.                    \tag{6}
\end{equation}
This is evaluation at the identity matrix.  As $t\to\infty$, only weight zero
survives:
\begin{equation}
\lim_{t\to\infty}[\mathbf t^\tau]\Sfun_t=M_\tau(0)
=\dim W_\tau^{U(1)}.                                       \tag{7}
\end{equation}
Thus (5) gives a literal spectral interpolation between dimension and Haar
invariants.  For a purely positive defining weight there are no positive-degree
Haar invariants.  Mixed positive and negative weights, however, can cancel and
leave nontrivial weight-zero terms.

\subsection{Representations and verification method}
The tested weight lists are
\begin{center}
\begin{tabular}{@{}lccp{7.1cm}@{}}
\toprule
Name & weights & dimension & purpose\\
\midrule
defining & $(1)$ & 1 & isolates the scalar Casimir $r^2$\\
weight pair & $(1,-1)$ & 2 & creates many weight-zero invariants\\
mixed & $(0,1,-2)$ & 3 & combines a fixed line with asymmetric weights\\
\bottomrule
\end{tabular}
\end{center}
For each representation, $t=0,0.25,0.8,3,20$ and all 19 partitions through
total degree five are tested.

The direct route uses 512 equally spaced angles.  At each angle it constructs
the matrix (3), computes complete characters from the matrix power traces, and
integrates their product against the Fourier heat kernel (1), truncated at
$|q|\leq32$.  The formula route never evaluates an angle or matrix: it builds
the integer multiplicities $M_\tau(w)$ and applies (5).

\subsubsection*{Why the quadrature is especially strong}
This is not a generic Monte Carlo comparison.  Through total degree five the
largest character frequency is at most $5\max_j|a_j|=10$.  Multiplication by
the truncated kernel produces frequencies at most 42.  Since 512 exceeds
twice this bandwidth, the periodic trapezoidal rule integrates the entire
trigonometric polynomial exactly in exact arithmetic.  At the smallest
positive time, the first omitted heat coefficient is
$e^{-33^2/8}<10^{-59}$.  The observed errors are therefore floating-point
roundoff, not sampling noise or appreciable heat-kernel truncation.

At $t=0$ the heat density is a delta distribution, so the code evaluates the
identity matrix directly as in (6).  This avoids approximating a singular
endpoint with a finite Fourier sum.

\subsection{Results}
\begin{center}
\begin{tabularx}{\linewidth}{@{}lXrrc@{}}
\toprule
Representation and time & coefficient & quadrature & Casimir formula & result\\
\midrule
defining, $t=0.8$ & $(3)$ & $0.0273237224$ & $0.0273237224$ & \PASS\\
weight pair, $t=3$ & $(2,1)$ & $0.8925233825$ & $0.8925233825$ & \PASS\\
mixed, $t=20$ & $(2,2)$ & $5.0003631994$ & $5.0003631994$ & \PASS\\
\bottomrule
\end{tabularx}
\end{center}
The last value is already close to the integer Haar multiplicity five, while
the small positive correction records the surviving nonzero weights at finite
time.

\subsection{Scope and reproducibility}
The computation verifies the compact heat formula for $U(1)$ with several
representation dimensions and both one-variable and mixed coefficients.  It
does not by itself verify the branching rules proposed for $U(n)$, $SU(n)$,
$O(n)$, or $USp(2n)$.  Those require separate highest-weight and Casimir
normalizations.  The disconnected-group warning is also essential: Brownian
motion from the identity in $O(n)$ remains in $SO(n)$ unless component-changing
jumps are added.

Run \path{compact_u1_heat/verification.py}, or open
\path{marimo_notebooks/u1_heat.py} in marimo.
\endgroup


\clearpage
\section{Non-Haar Fourier-weighted sigma-MGFs}\label{proof:new-dist:non_haar_fourier}

\resultbox{\textbf{Canonical testing artifacts.}\par\smallskip Exact backend: \path{lambda_distributions/dists/non_haar_fourier_sigma_mgfs/verification.py}.\par Regression: \path{lambda_distributions/dists/non_haar_fourier_sigma_mgfs/test_verification.py}.\par Marimo: \path{marimo_notebooks/non_haar_fourier.py}.}

\subsection*{Scope and outcome}

For a probability law on a finite or compact matrix group, Haar averaging is
replaced by a Fourier-weighted operator rather than a Reynolds projection.
We prove the exact Schur and monomial formulas, the full matrix Fourier formula
for arbitrary random walks, its scalar character-ratio specialization for
central walks, and the Casimir-weighted heat formula.  Verification is split
by group family.  For $S_3$, every natural and deleted-representation matrix
is constructed and arbitrary convolution is compared with all irreducible
Fourier blocks.  For $U(1)$, direct heat-kernel quadrature of $2\times2$
matrices is compared with weight decomposition.  For Ewens laws on
$S_3,\ldots,S_6$, exact matrix enumeration is compared with the cycle-index
formula.  The accompanying marimo notebook exposes every comparison.

\subsection{The non-Haar master theorem}

Let $G$ be finite or compact, let $dg$ denote normalized Haar measure, and
let $\rho:G\to U(V)$ be a finite-dimensional unitary representation.  A
probability measure $\mu$ on $G$ need not be invariant.  For variables
$\mathbf t=(t_1,t_2,\ldots)$ define
\begin{equation}\label{nd:non_haar_fourier:eq:def}
 \MGF_{\mu,V}(\mathbf t)
 =\int_G\prod_{i\ge1}\det(I-t_i\rho(g))^{-1}\,d\mu(g).
\end{equation}
The statement is formal coefficientwise.  It is also analytic when only
finitely many variables are nonzero and $|t_i|<1$.

Write $p_r=\sum_i t_i^r$.  For a partition
$\tau=(\tau_1,\ldots,\tau_s)$ put
\begin{equation}\label{nd:non_haar_fourier:eq:wtau}
 W_\tau(V)=\bigotimes_{j=1}^s\Sym^{\tau_j}V.
\end{equation}

\begin{theorem}[non-Haar $\sigma$-MGF]\label{nd:non_haar_fourier:thm:master}
The following identities hold coefficientwise:
\begin{align}
 \MGF_{\mu,V}(\mathbf t)
 &=\int_G\exp\!\left(\sum_{r\ge1}\frac{p_r}{r}
     \chi_V(g^r)\right)d\mu(g),\label{nd:non_haar_fourier:eq:power}\\
 &=\sum_\lambda s_\lambda(\mathbf t)
     \int_G\chi_{S_\lambda V}(g)\,d\mu(g),\label{nd:non_haar_fourier:eq:schur}\\
 [m_\tau]\MGF_{\mu,V}
 &=\int_G\chi_{W_\tau(V)}(g)\,d\mu(g).\label{nd:non_haar_fourier:eq:monomial}
\end{align}
If
$S_\lambda V\cong\bigoplus_{\pi\in\widehat G}
m_\pi(S_\lambda V)\pi$, then
\begin{equation}\label{nd:non_haar_fourier:eq:fourier-mu}
 \boxed{\MGF_{\mu,V}
 =\sum_\lambda s_\lambda
   \sum_{\pi\in\widehat G}m_\pi(S_\lambda V)
        \Tr\widehat\mu(\pi),\qquad
 \widehat\mu(\pi)=\int_G\pi(g)\,d\mu(g).}
\end{equation}
The analogous monomial formula replaces $S_\lambda V$ by $W_\tau(V)$.
\end{theorem}

\begin{proof}
If $A$ has eigenvalues $\alpha_1,\ldots,\alpha_d$, then
\begin{equation}\label{nd:non_haar_fourier:eq:detexp}
 \det(I-tA)^{-1}
 =\prod_j(1-t\alpha_j)^{-1}
 =\exp\!\left(\sum_{r\ge1}\frac{t^r}{r}\Tr(A^r)\right).
\end{equation}
Apply this identity to $A=\rho(g)$ and multiply over $i$ to obtain
\eqref{nd:non_haar_fourier:eq:power}.  The Cauchy identity gives
\begin{equation}\label{nd:non_haar_fourier:eq:cauchy}
 \prod_i\det(I-t_i\rho(g))^{-1}
 =\sum_\lambda s_\lambda(\mathbf t)\chi_{S_\lambda V}(g),
\end{equation}
which proves \eqref{nd:non_haar_fourier:eq:schur}.  Alternatively, the coefficient of $t^q$ in
one determinant inverse is $\chi_{\Sym^qV}(g)$.  Multiplying the selected
coefficients proves \eqref{nd:non_haar_fourier:eq:monomial}.

Finally, characters add over direct sums.  On each irreducible summand $\pi$,
the integral of its character is $\Tr\int\pi(g)d\mu(g)$.  Summing with
multiplicity proves \eqref{nd:non_haar_fourier:eq:fourier-mu}.  No invariance of $\mu$ was used.
\end{proof}

\begin{remark}[what replaces Reynolds projection]
For Haar measure,
$\int_G\rho_W(g)dg$ is the idempotent projection onto $W^G$, so
$\int\chi_W=\dim W^G$.  For general $\mu$ the operator
$T_{\mu,W}=\int\rho_W(g)d\mu(g)$ is usually neither idempotent nor central.
The coefficient is its trace.  Thus a non-Haar law produces a
Fourier-weighted averaging operator, not an invariant dimension.
\end{remark}

\subsection{Random walks and the cutoff-visible spectrum}

Let $\nu$ be a step distribution and $\mu_k=\nu^{*k}$.  With the convention
in \eqref{nd:non_haar_fourier:eq:fourier-mu}, representation multiplicativity and Fubini give
\begin{equation}\label{nd:non_haar_fourier:eq:conv}
 \widehat{\nu^{*k}}(\pi)=\widehat\nu(\pi)^k.
\end{equation}

\begin{corollary}[arbitrary walk]\label{nd:non_haar_fourier:cor:walk}
For every partition $\tau$,
\begin{equation}\label{nd:non_haar_fourier:eq:walk}
 \boxed{[m_\tau]\MGF_{\nu^{*k},V}
 =\sum_{\pi\in\widehat G}m_\pi(W_\tau(V))
       \Tr\bigl(\widehat\nu(\pi)^k\bigr).}
\end{equation}
If $\nu$ is central, Schur's lemma gives
\begin{equation}\label{nd:non_haar_fourier:eq:theta}
 \widehat\nu(\pi)=\theta_\pi I_{d_\pi},\qquad
 \theta_\pi=\frac1{d_\pi}\int_G\chi_\pi(g)d\nu(g),
\end{equation}
and hence
\begin{equation}\label{nd:non_haar_fourier:eq:central}
 \boxed{[m_\tau]\MGF_{\nu^{*k},V}
 =\sum_\pi m_\pi(W_\tau(V))d_\pi\theta_\pi^k.}
\end{equation}
\end{corollary}

The trivial term in \eqref{nd:non_haar_fourier:eq:central} is the Haar coefficient.  Therefore
\begin{equation}\label{nd:non_haar_fourier:eq:deviation}
 [m_\tau](\MGF_{\nu^{*k},V}-\MGF_{\mathrm{Haar},V})
 =\sum_{\pi\ne\one}m_\pi(W_\tau(V))d_\pi\theta_\pi^k.
\end{equation}
For a fixed finite group this gives the exact leading exponential decay after
combining all terms with maximal $|\theta_\pi|$.  For a sequence
$(G_n,V_n,\nu_n)$ it also identifies the correct bounded-degree cutoff
criterion: only irreducibles occurring in the fixed $W_{n,\tau}$ are visible.
This is weaker than total-variation mixing of the full group-valued walk.

\begin{remark}[centrality is not optional]
Equation \eqref{nd:non_haar_fourier:eq:central} must not be used for a noncentral generator law.
The trace of $\widehat\nu(\pi)$ does not determine
$\Tr(\widehat\nu(\pi)^k)$.  Section~\ref{nd:non_haar_fourier:sec:s3} includes a concrete
$2\times2$ counterexample to that shortcut.
\end{remark}

\subsection{Compact heat kernels}

Let the Brownian generator be $\frac12\Delta_G$, and let $\kappa_\pi\ge0$
be the Casimir eigenvalue on $\pi$ for this normalization.  The heat
semigroup diagonalizes on matrix coefficients:
\begin{equation}\label{nd:non_haar_fourier:eq:heat-fourier}
 \widehat\mu_t(\pi)=e^{-t\kappa_\pi/2}I_{d_\pi}.
\end{equation}
Substitution in Theorem~\ref{nd:non_haar_fourier:thm:master} proves the following.

\begin{theorem}[heat formula]\label{nd:non_haar_fourier:thm:heat}
\begin{align}
 \MGF^{\mathrm{heat}}_{G,V}(t)
 &=\sum_\lambda s_\lambda
   \sum_\pi m_\pi(S_\lambda V)d_\pi e^{-t\kappa_\pi/2},
   \label{nd:non_haar_fourier:eq:heat-schur}\\
 [m_\tau]\MGF^{\mathrm{heat}}_{G,V}(t)
 &=\sum_\pi m_\pi(W_\tau(V))d_\pi e^{-t\kappa_\pi/2}.
   \label{nd:non_haar_fourier:eq:heat-monomial}
\end{align}
At $t=0$ the coefficient is $\dim W_\tau(V)$, and as $t\to\infty$ it tends
to $\dim W_\tau(V)^G$. If a positive-Casimir constituent occurs, let
$\gamma_\tau$ be the smallest positive Casimir appearing in $W_\tau(V)$.
Then
\begin{equation}\label{nd:non_haar_fourier:eq:heat-asymp}
 [m_\tau]\MGF_t=\dim W_\tau(V)^G
 +e^{-t\gamma_\tau/2}
   \sum_{\kappa_\pi=\gamma_\tau}m_\pi(W_\tau)d_\pi
 +o(e^{-t\gamma_\tau/2}).
\end{equation}
\end{theorem}

This theorem applies without conceptual change
If no positive-Casimir constituent occurs, the coefficient is already
constant and the remainder in the preceding asymptotic formula is zero.
 to classical and exceptional
compact groups once the representation and metric normalization are fixed.
For a growing-rank statement, that normalization is essential: rescaling the
metric rescales every $\kappa_\pi$ and therefore the asserted cutoff time.

\subsection{Finite matrix group: two walks on \texorpdfstring{$S_3$}{S3}}
\label{nd:non_haar_fourier:sec:s3}

\subsubsection{Constructed representations}
For $\sigma\in S_3$, let $P_\sigma$ be its $3\times3$ permutation matrix,
$P_\sigma e_j=e_{\sigma(j)}$.  The natural module decomposes as
\begin{equation}\label{nd:non_haar_fourier:eq:natdec}
 \C^3\cong\one\oplus V_{\mathrm{std}}.
\end{equation}
The verifier also constructs the deleted matrices on
$V_{\mathrm{std}}=\{x_1+x_2+x_3=0\}$ in the basis
$f_1=e_1-e_3,f_2=e_2-e_3$.  Their entries are
\begin{equation}\label{nd:non_haar_fourier:eq:deleted}
 (D_\sigma)_{ij}=\delta_{i,\sigma(j)}-\delta_{i,\sigma(3)},
 \qquad 1\le i,j\le2,
\end{equation}
where a delta is zero if its second index is $3$.  All six matrices in each
image are generated explicitly and multiplication is exact over $\mathbb Z$.

The three irreducible character rows, ordered by identity, transposition, and
three-cycle, are
\begin{center}
\begin{tabular}{@{}lrrrr@{}}
\toprule
$\pi$ & $d_\pi$ & $\chi_\pi(1)$ & $\chi_\pi((12))$ & $\chi_\pi((123))$\\
\midrule
trivial & 1 & 1 & 1 & 1\\
sign & 1 & 1 & $-1$ & 1\\
standard & 2 & 2 & 0 & $-1$\\
\bottomrule
\end{tabular}
\end{center}

\subsubsection{Central lazy-transposition law}
Set
\begin{equation}\label{nd:non_haar_fourier:eq:centralstep}
 \nu_c=\frac14\delta_e+\frac14\sum_{r\text{ a transposition}}\delta_r.
\end{equation}
It is central, and \eqref{nd:non_haar_fourier:eq:theta} gives
\begin{equation}\label{nd:non_haar_fourier:eq:s3theta}
 \theta_{\mathrm{triv}}=1,\qquad
 \theta_{\mathrm{sign}}=-\frac12,\qquad
 \theta_{\mathrm{std}}=\frac14.
\end{equation}
Thus every tested coefficient has the completely explicit form
\begin{equation}\label{nd:non_haar_fourier:eq:s3centralformula}
 m_{\mathrm{triv}}(W_\tau)
 +m_{\mathrm{sign}}(W_\tau)(-1/2)^k
 +2m_{\mathrm{std}}(W_\tau)(1/4)^k.
\end{equation}

\subsubsection{Noncentral two-generator law}
Set
\begin{equation}\label{nd:non_haar_fourier:eq:noncentralstep}
 \nu_{nc}=\frac12\delta_{(23)}+\frac12\delta_{(123)}.
\end{equation}
Using \eqref{nd:non_haar_fourier:eq:deleted}, its standard Fourier block is
\begin{equation}\label{nd:non_haar_fourier:eq:noncentralmatrix}
 \widehat\nu_{nc}(\mathrm{std})
 =\frac12D_{(23)}+\frac12D_{(123)}
 =\begin{pmatrix}0&-1/2\\0&-1/2\end{pmatrix}.
\end{equation}
This is not scalar.  The correct term in \eqref{nd:non_haar_fourier:eq:walk} is
$m_{\mathrm{std}}(W_\tau)\Tr(\widehat\nu_{nc}(\mathrm{std})^k)$;
a character-ratio substitution would discard matrix information.

\subsubsection{Exact verification strategy}
For each representation, partition $\tau$ through total degree four, and
$k\in\{0,1,2,3,5\}$, the program compares:
\begin{enumerate}
\item \textbf{Direct convolution and matrices.}  It forms $\nu^{*k}$ on all
six elements.  For each matrix $A$, it computes
\begin{equation}\label{nd:non_haar_fourier:eq:newton}
 h_0(A)=1,\qquad h_d(A)=\frac1d\sum_{r=1}^d
 \Tr(A^r)h_{d-r}(A),
\end{equation}
then averages $\prod_jh_{\tau_j}(A)$ literally.
\item \textbf{Full Fourier route.}  Multiplicities are computed independently
from
\begin{equation}\label{nd:non_haar_fourier:eq:multiplicity}
 m_\pi(W_\tau)=\frac1{6}\sum_{g\in S_3}
 \chi_{W_\tau}(g)\chi_\pi(g^{-1}),
\end{equation}
and inserted into \eqref{nd:non_haar_fourier:eq:walk} using exact rational matrix powers.
\item \textbf{Central scalar route.}  For $\nu_c$ only, the code separately
checks that every Fourier block is scalar and evaluates
\eqref{nd:non_haar_fourier:eq:s3centralformula}.
\end{enumerate}

\begin{center}
\begin{tabular}{@{}lllrrc@{}}
\toprule
walk & representation & $(k,\tau)$ & direct & Fourier & result\\
\midrule
central & natural & $(3,(1))$ & $33/32$ & $33/32$ & \textcolor{teal}{\textbf{PASS}}\\
central & natural & $(3,(2,2))$ & $63/8$ & $63/8$ & \textcolor{teal}{\textbf{PASS}}\\
central & deleted & $(5,(1))$ & $1/512$ & $1/512$ & \textcolor{teal}{\textbf{PASS}}\\
noncentral & natural & $(3,(1))$ & $7/8$ & $7/8$ & \textcolor{teal}{\textbf{PASS}}\\
noncentral & deleted & $(3,(2,1))$ & $3/4$ & $3/4$ & \textcolor{teal}{\textbf{PASS}}\\
noncentral & deleted & $(5,(2,2))$ & $61/32$ & $61/32$ & \textcolor{teal}{\textbf{PASS}}\\
\bottomrule
\end{tabular}
\end{center}
The displayed noncentral values are generated by the same suite as the full
table.  All 220 exact comparisons pass.

\subsection{Compact matrix group: \texorpdfstring{$U(1)$}{U(1)} heat flow}

Take
\begin{equation}\label{nd:non_haar_fourier:eq:u1rep}
 \rho(e^{i\phi})=\begin{pmatrix}e^{i\phi}&0\\0&e^{-i\phi}\end{pmatrix}
\end{equation}
on a two-dimensional representation with weights $+1,-1$.  For the standard
circle Laplacian, the weight-$q$ character has heat multiplier
$e^{-tq^2/2}$.  Moreover
\begin{equation}\label{nd:non_haar_fourier:eq:symweights}
 \chi_{\Sym^dV}(e^{i\phi})
 =\sum_{a=0}^d e^{i(d-2a)\phi}.
\end{equation}
If $c_\tau(q)$ is the coefficient of $z^q$ in
\begin{equation}\label{nd:non_haar_fourier:eq:weightpoly}
 \prod_{j=1}^{\ell(\tau)}
   (z^{\tau_j}+z^{\tau_j-2}+\cdots+z^{-\tau_j}),
\end{equation}
then Theorem~\ref{nd:non_haar_fourier:thm:heat} becomes the finite exact sum
\begin{equation}\label{nd:non_haar_fourier:eq:u1formula}
 \boxed{[m_\tau]\MGF_t
 =\sum_{q\in\mathbb Z}c_\tau(q)e^{-tq^2/2}.}
\end{equation}
At $t=0$ this is $\prod_j(\tau_j+1)$; at Haar time it tends to
$c_\tau(0)$.  The first nonzero $|q|$ in \eqref{nd:non_haar_fourier:eq:weightpoly} gives the
leading decay rate, an explicit instance of \eqref{nd:non_haar_fourier:eq:heat-asymp}.

\subsubsection{Independent numerical route}
The verifier does not use \eqref{nd:non_haar_fourier:eq:weightpoly} in its direct route.  On a
256-node periodic grid it constructs \eqref{nd:non_haar_fourier:eq:u1rep}, recovers every
$h_d(\rho(e^{i\phi}))$ from matrix traces using \eqref{nd:non_haar_fourier:eq:newton}, and
multiplies by
\begin{equation}\label{nd:non_haar_fourier:eq:heatkernel}
 K_t(\phi)=1+2\sum_{q=1}^{48}e^{-tq^2/2}\cos(q\phi).
\end{equation}
The trapezoidal average is compared with \eqref{nd:non_haar_fourier:eq:u1formula}.  The grid
resolves all retained Fourier modes; the omitted tail is far below the
reported tolerance at the smallest tested time $t=0.2$.

\begin{center}
\begin{tabular}{@{}crrrc@{}}
\toprule
$t$ & $\tau$ & direct quadrature & weight formula & error\\
\midrule
$0.2$ & $(1)$ & $1.809674836071920$ & $1.809674836071919$ & $6.7\times10^{-16}$\\
$0.2$ & $(2,2)$ & $6.085073220131871$ & $6.085073220131868$ & $3.6\times10^{-15}$\\
$0.7$ & $(1,1)$ & $2.493193927883212$ & $2.493193927883213$ & $1.3\times10^{-15}$\\
$2.0$ & $(5)$ & $0.736005701978834$ & $0.736005701978834$ & $4.4\times10^{-16}$\\
\bottomrule
\end{tabular}
\end{center}
All 54 comparisons through degree five pass.  The largest absolute error over
the complete suite is $1.5988\times10^{-14}$.  For all 18 partitions, the
$t=0$ dimension and $t\to\infty$ Haar-weight identities also pass.

\subsection{Ewens laws on symmetric matrix groups}

For $\sigma\in S_n$, let $K(\sigma)$ be its number of cycles and take
\begin{equation}\label{nd:non_haar_fourier:eq:ewenslaw}
 \mathbb P_{n,\theta}(\sigma)
 =\frac{\theta^{K(\sigma)}}{\theta^{(n)}},\qquad
 \theta^{(n)}=\theta(\theta+1)\cdots(\theta+n-1).
\end{equation}
The normalization follows from
$\sum_{\sigma\in S_n}\theta^{K(\sigma)}=\theta^{(n)}$.

If $m_j(\lambda)$ is the number of parts of size $j$ in $\lambda\vdash n$
and $z_\lambda=\prod_jj^{m_j(\lambda)}m_j(\lambda)!$, a permutation of cycle
shape $\lambda$ satisfies
\begin{equation}\label{nd:non_haar_fourier:eq:cycledet}
 \det(I-tP_\sigma)^{-1}
 =\prod_{j\ge1}(1-t^j)^{-m_j(\lambda)}.
\end{equation}

\begin{proposition}[exact Ewens formula]\label{nd:non_haar_fourier:prop:ewens}
For the natural permutation representation,
\begin{equation}\label{nd:non_haar_fourier:eq:ewens}
 \boxed{\MGF_{n,\theta}(\mathbf t)
 =\frac{n!}{\theta^{(n)}}
  \sum_{\lambda\vdash n}\frac{\theta^{\ell(\lambda)}}{z_\lambda}
  \prod_{i,j}(1-t_i^j)^{-m_j(\lambda)}.}
\end{equation}
\end{proposition}

\begin{proof}
There are $n!/z_\lambda$ permutations of shape $\lambda$.  They all have
$\ell(\lambda)$ cycles, hence the same Ewens weight, and
\eqref{nd:non_haar_fourier:eq:cycledet} gives their common matrix integrand.  Grouping the direct
element sum by shape proves \eqref{nd:non_haar_fourier:eq:ewens}.
\end{proof}

For fixed $j$, the Ewens cycle counts converge jointly to independent
Poisson variables of means $\theta/j$.  Consequently, coefficientwise in the
positive-degree variables,
\begin{equation}\label{nd:non_haar_fourier:eq:ewenslimit}
 \MGF_{\theta,\infty}
 =\exp\!\left[\sum_{j\ge1}\frac{\theta}{j}
   \left(\prod_i(1-t_i^j)^{-1}-1\right)\right].
\end{equation}
For a coefficient of bounded degree, only finitely many $j$ contribute, so
the formal interpretation is unambiguous.

\subsubsection{Exact verification strategy and results}
The direct route enumerates all $n!$ permutation matrices, computes their
symmetric characters from matrix powers, and applies \eqref{nd:non_haar_fourier:eq:ewenslaw}.
The independent route stores only cycle lengths, expands
\eqref{nd:non_haar_fourier:eq:cycledet}, and aggregates class sizes.  Both use rational
arithmetic for
\begin{equation*}
 n=3,4,5,6,\qquad \theta\in\{1/2,1,2\},
 \qquad |\tau|\le4.
\end{equation*}
This gives 132 coefficient comparisons.  At
$(t_1,t_2)=(1/10,-1/20)$, 12 further comparisons evaluate the full rational
determinant average without truncating in degree.  Every comparison passes.

\begin{center}
\begin{tabular}{@{}cccrrc@{}}
\toprule
$n$ & $\theta$ & $\tau$ & matrix enumeration & cycle formula & result\\
\midrule
3 & $1/2$ & $(1)$ & $3/5$ & $3/5$ & \textcolor{teal}{\textbf{PASS}}\\
3 & $1/2$ & $(2,2)$ & $4$ & $4$ & \textcolor{teal}{\textbf{PASS}}\\
3 & $2$ & $(1)$ & $3/2$ & $3/2$ & \textcolor{teal}{\textbf{PASS}}\\
6 & $1/2$ & $(2,1)$ & $326/231$ & $326/231$ & \textcolor{teal}{\textbf{PASS}}\\
6 & $1$ & $(2,2)$ & $9$ & $9$ & \textcolor{teal}{\textbf{PASS}}\\
6 & $2$ & $(1,1)$ & $32/7$ & $32/7$ & \textcolor{teal}{\textbf{PASS}}\\
\bottomrule
\end{tabular}
\end{center}

\subsection{Related measures and representation choices}

\subsubsection{Conjugacy-class walks}
If $\nu_K$ is uniform on a conjugacy class $K$, then
$\theta_\pi=\chi_\pi(K)/d_\pi$.  Therefore
\begin{equation}\label{nd:non_haar_fourier:eq:classwalk}
 [m_\tau]\MGF_{\nu_K^{*k},V}
 =\sum_\pi m_\pi(W_\tau)d_\pi
   \left(\frac{\chi_\pi(K)}{d_\pi}\right)^k.
\end{equation}
This includes random transpositions and central reflection walks.  If a
generator set is a union of classes, its weighted average character ratio is
used.  Noncentral elementary-generator walks require \eqref{nd:non_haar_fourier:eq:walk} instead.

\subsubsection{Character-weighted laws}
For an irreducible $\eta$, character orthogonality makes
$|\chi_\eta(g)|^2dg$ a probability measure.  Formula \eqref{nd:non_haar_fourier:eq:monomial}
becomes
\begin{equation}\label{nd:non_haar_fourier:eq:charweighted}
 [m_\tau]\MGF_{|\chi_\eta|^2,V}
 =\langle\chi_{W_\tau},\chi_\eta\chi_{\eta^*}\rangle_G
 =\dim\Hom_G(W_\tau,\eta\otimes\eta^*).
\end{equation}
This is a tensor-product multiplicity, not generally an invariant dimension.
Plancherel measure itself lives on $\widehat G$, so an additional construction
is needed before it defines a law on group elements.

\subsubsection{Finite-field and Clifford cautions}
A finite-field matrix action is not automatically a complex representation.
One must specify a complex permutation, Weil, or other cross-characteristic
representation before applying the theorem.  For scalar-rich unitary groups,
one-sided moments may vanish in many degrees; a balanced series involving
both $g$ and $g^*$ can be more informative.  These are changes of
representation or observable, not exceptions to Theorem~\ref{nd:non_haar_fourier:thm:master}.

\begingroup
\small
\setlength{\parskip}{3pt}
\subsection{Reproducibility and audit summary}

The package contains one pure-Python verifier and one marimo interface.  No
external character table, symbolic algebra service, or random seed is used.
Finite computations use \texttt{fractions.Fraction}; only the compact
quadrature uses floating point.

\begin{center}
\begin{tabularx}{\linewidth}{@{}lXrrc@{}}
\toprule
family & independent routes & cases & comparisons & result\\
\midrule
$S_3$ walks & convolution/matrix traces; full Fourier blocks; central scalars
& $2$ laws, $2$ reps & 220 & \textcolor{teal}{\textbf{PASS}}\\
$U(1)$ heat & matrix-integrand quadrature; weight/Casimir sum
& $3$ times, $18$ partitions & 54 & \textcolor{teal}{\textbf{PASS}}\\
Ewens $S_n$ & element matrices; cycle-shape aggregation
& $4$ ranks, $3$ parameters & 132 & \textcolor{teal}{\textbf{PASS}}\\
Ewens full series & element determinants; class determinants
& $12$ laws & 12 & \textcolor{teal}{\textbf{PASS}}\\
\bottomrule
\end{tabularx}
\end{center}

Run the tests from the repository root:
\begin{verbatim}
.venv/bin/python -m pytest -q \
  lambda_distributions/dists/non_haar_fourier_sigma_mgfs/test_verification.py
\end{verbatim}
Open the grouped interactive notebook with:
\begin{verbatim}
.venv/bin/marimo edit \
  marimo_notebooks/non_haar_fourier.py
\end{verbatim}
\resultbox{\textbf{Conclusion.}  The tested cases verify the intended
replacement precisely: Haar measure extracts the trivial representation,
whereas a non-Haar law weights every constituent by a Fourier trace.  Central
walks reduce to character ratios, noncentral walks retain full Fourier
matrices, heat flow supplies Casimir exponentials, and Ewens weighting is
captured exactly by cycle-shape aggregation.}
\endgroup

\clearpage
\section[Limiting laws for symmetric-group k-subset actions]{Limiting laws for symmetric-group $k$-subset actions}\label{proof:limit-subsets}
\begingroup
\definecolor{navy}{HTML}{17365D}
\definecolor{teal}{HTML}{147D78}
\providecommand{\Fix}{\operatorname{Fix}}
\renewcommand{\Fix}{\operatorname{Fix}}
\providecommand{\SMG}{\operatorname{SMG}}
\renewcommand{\SMG}{\operatorname{SMG}}
\setcounter{theorem}{0}
\setcounter{proposition}{0}
\setcounter{lemma}{0}
\setcounter{corollary}{0}
\setcounter{remark}{0}
\begin{abstract}
For fixed $k$, we prove the joint limit of all power traces of the permutation
representation of $S_n$ on $k$-subsets.  The limit is an explicit functional
of independent $Z_j\sim\mathrm{Poisson}(1/j)$ variables and is non-Poisson for
$k\ge2$.  An independent program constructs 1,584 representation matrices in
dimensions $6,10,15,20$ and performs 6,336 exact power-trace comparisons.
\end{abstract}

\subsection{The representation and exact formula}
Let $\Omega_{n,k}=\binom{[n]}k$ and let $\rho_{n,k}$ permute its standard
basis.  Since the trace of a permutation matrix counts fixed basis vectors,
\[
 p_r(X_{\rho}(\sigma))=\operatorname{Tr}(\rho(\sigma)^r)
 =|\Fix_{\Omega_{n,k}}(\sigma^r)|.
\]
Write $A_j(\sigma)$ for the number of $j$-cycles of $\sigma$.  A $j$-cycle
splits under the $r$th power into $d=\gcd(j,r)$ cycles of length $j/d$.
An invariant subset independently includes or excludes each resulting cycle,
so
\begin{equation}\label{symmetric-k-subsets-proof-and-verification:eq:fixed}
\boxed{F_{k,r}(\sigma)=[u^k]\prod_{j=1}^{kr}
\left(1+u^{j/\gcd(j,r)}\right)^{\gcd(j,r)A_j(\sigma)}.}
\end{equation}
The cutoff is valid: if $j>kr$, then $j/\gcd(j,r)>k$.

For a partition $\lambda$, direct averaging therefore gives
\begin{equation}\label{symmetric-k-subsets-proof-and-verification:eq:smg}
\SMG_{n,k}=\sum_\lambda\frac{\mathbb E_{S_n}
[\prod_aF_{k,\lambda_a}(\sigma)]}{z_\lambda}p_\lambda.
\end{equation}

\subsection{Coefficientwise limit}
Fix a finite set of exponents.  Every variable in \eqref{symmetric-k-subsets-proof-and-verification:eq:fixed} is a
polynomial in finitely many cycle counts.  We use the following elementary
factorial-moment lemma.  For fixed nonnegative integers $r_1,\ldots,r_R$,
\[
\mathbb E\prod_{j=1}^R(A_j)_{r_j}=\prod_{j=1}^Rj^{-r_j}
\quad\text{whenever}\quad n\geq\sum_{j=1}^Rjr_j.
\]
Indeed, the falling-factorial product counts ordered selections of pairwise
disjoint distinguished cycles, with $r_j$ cycles of length $j$.  Choosing
their supports and cyclic orders and then permuting the remaining labels gives
$n!\prod_jj^{-r_j}/n!$.  These are exactly the joint factorial moments of
independent $Z_j\sim\operatorname{Poisson}(1/j)$.  Hence the finite vector of
cycle counts converges jointly, and every factorial moment required by a fixed
coefficient is eventually equal to its limiting value.  Thus
\[
(A_1,\ldots,A_R)\Longrightarrow(Z_1,\ldots,Z_R),\qquad
Z_j\ \text{independent},\quad Z_j\sim\mathrm{Poisson}(1/j).
\]
Thus
\[
F_{k,r}(\sigma)\Longrightarrow F^{(\infty)}_{k,r}:=
[u^k]\prod_{j=1}^{kr}
(1+u^{j/\gcd(j,r)})^{\gcd(j,r)Z_j},
\]
jointly for every finite collection of $r$'s.  Moment stabilization permits
expectations in \eqref{symmetric-k-subsets-proof-and-verification:eq:smg} to pass to the limit, proving the coefficientwise
$\Lambda$-limit
\[
\boxed{\SMG^{(k)}_\infty=\sum_\lambda
\frac{\mathbb E[\prod_aF^{(\infty)}_{k,\lambda_a}]}{z_\lambda}p_\lambda.}
\]
For $r=1$, $F^{(\infty)}_{k,1}$ has mean one.  The action has rank $k+1$
once $n\ge2k$, hence its limiting second moment is $k+1$ and its variance is
$k$.  A Poisson variable of mean one has variance one; consequently the limit
is non-Poisson for $k\ge2$.  For example
\[
F^{(\infty)}_{2,1}=\binom{Z_1}{2}+Z_2.
\]

\subsection{Exact finite formula checks}
For each $(n,k)=(4,2),(5,2),(6,2),(6,3)$, the verifier enumerates all $n!$
permutations, induces their action on lexicographically ordered $k$-subsets,
and materializes the corresponding zero-one matrix.  Distinct byte keys prove
that the constructed matrices form a faithful copy of $S_n$.  It then compares
the integer matrix traces $\operatorname{Tr}(M^r)$, $1\le r\le4$, with
\eqref{symmetric-k-subsets-proof-and-verification:eq:fixed}.  Six joint products indexed by
$(1),(2),(3),(1,1),(2,1),(1,1,1)$ are averaged independently on both sides.

\begin{center}
\begin{tabular}{@{}rrrrrr@{}}\toprule
$n$&$k$&dimension&matrices&trace checks&$(\mathbb E F,\operatorname{Var}F)$\\\midrule
4&2&6&24&96&$(1,2)$\\
5&2&10&120&480&$(1,2)$\\
6&2&15&720&2880&$(1,2)$\\
6&3&20&720&2880&$(1,3)$\\\bottomrule
\end{tabular}
\end{center}
All 6,336 trace comparisons and all 24 joint-average comparisons pass exactly.
These exhaustive matrices verify the finite fixed-point identity, not the
asymptotic convergence; convergence is proved by the factorial-moment lemma
above.

\subsection{Reproduction and reference}
\begin{verbatim}
.venv/bin/python for_this_guy/limiting_lambda_claims/
  symmetric_k_subsets/verification.py
.venv/bin/marimo edit marimo_notebooks/symmetric_k_subsets.py
\end{verbatim}
Join each wrapped command onto one line.  For the fixed-point and moment
background, see P.~Diaconis, J.~Fulman, and R.~Guralnick,
\emph{On fixed points of permutations},
\url{https://arxiv.org/abs/0708.2643}.
\endgroup


\clearpage
\section{Rooted-tree automorphism groups}\label{proof:rooted-trees}
\begingroup
\definecolor{navy}{HTML}{17365D}
\definecolor{teal}{HTML}{147D78}
\providecommand{\MGF}{\SMG}
\renewcommand{\MGF}{\SMG}
\setcounter{theorem}{0}
\setcounter{proposition}{0}
\setcounter{lemma}{0}
\setcounter{corollary}{0}
\setcounter{remark}{0}
\begin{abstract}
We prove the exact recursive sigma--MGF law for automorphisms of any finite
rooted tree.  Explicit leaf-permutation matrix groups for complete binary
trees through height three and ternary trees through height two give 30 exact
mixed-coefficient comparisons.  The document separates this proved finite
recursion from the still-open general random-tree $\Lambda$-limit.
\end{abstract}

\subsection{Wreath-product lemma}
Let a finite permutation group $H$ act on $X$, and let $H\wr S_m$ act on
$X\times[m]$.  Define
\[
\MGF_H(\mathbf t)=\frac1{|H|}\sum_{h\in H}
\prod_i\det(I-t_i h)^{-1}.
\]
Write an element of the wreath product as $(h_1,\ldots,h_m;\sigma)$.  On a
top cycle $c=(i_1\cdots i_\ell)$, one circuit accumulates
$h_c=h_{i_\ell}\cdots h_{i_1}$ and block elimination gives
\[
\det(I-tg|_c)=\det(I-t^\ell h_c).
\]
Uniform independent $h_i$ make the products $h_c$ uniform and independent
over the top cycles.  Averaging over $S_m$ and using its cycle index proves
\begin{equation}\label{rooted-tree-proof-and-verification:eq:wreath}
\boxed{\sum_{m\ge0}u^m\MGF_{H\wr S_m}(\mathbf t)=
\exp\left(\sum_{\ell\ge1}\frac{u^\ell}{\ell}
\MGF_H(t_1^\ell,t_2^\ell,\ldots)\right).}
\end{equation}

\subsection{Arbitrary rooted trees}
If the root branches have isomorphism types $A$ with multiplicities $m_A$,
then internal automorphisms and permutations of equal branches give
\[
\operatorname{Aut}(T)\cong\prod_A(\operatorname{Aut}(A)\wr S_{m_A}).
\]
The leaf representation is the corresponding direct sum.  Determinants
multiply on direct sums, so \eqref{rooted-tree-proof-and-verification:eq:wreath} yields
\begin{equation}\label{rooted-tree-proof-and-verification:eq:tree}
\boxed{\MGF_T(\mathbf t)=\prod_A[u_A^{m_A}]
\exp\left(\sum_{\ell\ge1}\frac{u_A^\ell}{\ell}
\MGF_A(\mathbf t^\ell)\right).}
\end{equation}
For the complete $b$-ary tree, $G_h=G_{h-1}\wr S_b$, starting with the
trivial action on one leaf.

\subsection{Exact computational verification}
The verifier recursively enumerates every wreath-product permutation and
materializes its zero-one matrix.  Distinct matrix keys and the identity
\[
|G_h|=b!\,|G_{h-1}|^b
\]
independently check each closure.  For a matrix $M$, the direct symmetric
characters are obtained over exact rationals from Newton's identity
\[
d h_d(M)=\sum_{r=1}^d\operatorname{Tr}(M^r)h_{d-r}(M).
\]
The other side is computed from the exact cycle-index average in
\eqref{rooted-tree-proof-and-verification:eq:wreath}, with truncated multivariate polynomial arithmetic.  Thus
the two routes share neither matrices nor coefficient extraction.

\begin{center}
\begin{tabular}{@{}rrrrr@{}}\toprule
$b$&height&leaves&base order&group order\\\midrule
2&1&2&1&2\\
2&2&4&2&8\\
2&3&8&8&128\\
3&1&3&1&6\\
3&2&9&6&1296\\\bottomrule
\end{tabular}
\end{center}
For every row the program compares coefficients indexed by
$(1),(2),(3),(1,1),(2,1),(2,2)$; all 30 comparisons pass exactly.

\subsection{What is and is not proved}
Equation \eqref{rooted-tree-proof-and-verification:eq:tree} is an exact theorem for each finite tree.  For a
Galton--Watson tree it becomes a distributional recursion after conditioning
on branch-type counts.  Also,
\[
\log|\operatorname{Aut}(T)|=\sum_i\log|\operatorname{Aut}(T_i)|
+\sum_A\log(m_A!)
\]
is an additive tree functional; Gaussian-type asymptotics for this scalar
quantity are known.  Neither fact alone proves convergence of the full random
series $\MGF_T(\mathbf t)$.  A general limiting rooted-tree
$\Lambda$-distribution therefore remains open, and the verifier makes no such
claim.

\subsection{Reproduction and reference}
\begin{verbatim}
.venv/bin/python for_this_guy/limiting_lambda_claims/
  rooted_tree_automorphisms/verification.py
.venv/bin/marimo edit marimo_notebooks/rooted_tree_automorphisms.py
\end{verbatim}
Join wrapped commands onto one line.  For automorphism-count asymptotics, see
M.~Olsson and S.~Wagner, \emph{The distribution of the number of
automorphisms of random trees},
\url{https://arxiv.org/abs/2211.12801}.
\endgroup

\clearpage

\subsection{Level-transitive leaf actions and the universal no-limit obstruction}

Let $L\leq S_d$ be transitive, put $G_0=1$, and define
$G_{h+1}=G_h\wr L$ in the natural action on the $d^{h+1}$ leaves of the
rooted $d$-ary tree.  If $F_h$ is the number of fixed leaves of a uniform
element of $G_h$, and if $Y$ is the number of fixed letters of a uniform
element of $L$, then
\begin{equation}\label{nd:symmetric_spin_lie_tree:eq:branching}
 F_{h+1}\ \stackrel{d}{=}\ \sum_{j=1}^{Y}F_h^{(j)},
\end{equation}
where the summands are independent copies of $F_h$.  Transitivity and
Burnside's lemma give $\mathbb E Y=1$.  Writing
$r_L=\#(L\backslash[d]^2)=\mathbb E(Y^2)$ and conditioning in the preceding
recursion yields
\begin{equation}\label{nd:symmetric_spin_lie_tree:eq:second-moment}
 \mathbb E F_h=1,\qquad
 \mathbb E(F_h^2)=1+h(r_L-1).
\end{equation}
Consequently
\[
 [m_{(1,1)}]\MGF_{G_h,\C[X_h]}=1+h(r_L-1),
\]
so every nontrivial transitive local action fails to have a raw
coefficientwise limit.

There is also a recurrence-free form.  For any group acting transitively on
the leaves at height $h$, the depth of the lowest common ancestor of an
ordered pair is invariant.  Equality and the $h$ possible divergence depths
therefore give
\begin{equation}\label{nd:symmetric_spin_lie_tree:eq:universal-bound}
 [m_{(1,1)}]\MGF_{G_h,\C[X_h]}
 =\#(G_h\backslash X_h^2)\ge h+1.
\end{equation}
This proves failure of raw moment convergence for every level-transitive
family, not only for iterated full wreath products.  It does not rule out
convergence in probability or a normalized theory.

\subsection{Exact audit of the iterated examples}

Direct leaf permutations, truncated cycle-index recurrences, fixed-point
squares, and pair-orbit searches give the same values:
\begin{center}
\small
\begin{tabular}{@{}lrrrrr@{}}\toprule
local action&height&leaves&group order&$\mathbb E F_h^2$&pair orbits\\\midrule
$C_2$&1&2&2&2&2\\
$C_2$&2&4&8&3&3\\
$C_2$&3&8&128&4&4\\
$C_3$&1&3&3&3&3\\
$C_3$&2&9&81&5&5\\
$S_3$&1&3&6&2&2\\
$S_3$&2&9&1296&3&3\\\bottomrule
\end{tabular}
\end{center}
The six coefficients indexed by $(1),(2),(1,1),(3),(2,1),(2,2)$ agree in
every row.  The finite recursion for arbitrary rooted trees remains exact,
whereas an infinite-height random-tree law is model-dependent.

\clearpage
\clearpage
\appendix
\section{Self-contained build note}

This edition is a single-source artifact.  Every proof body is present in this
file, and compilation does not read any local TeX source, bibliography,
graphic, code, data, or manifest.  File names and commands printed inside the
proofs are human-readable provenance notes only; they are not build inputs.
The only required dependencies are the standard LaTeX packages declared in
the preamble.


\clearpage
\part{Provisional claims and unresolved proof obligations}

\section{Local-ring generating-tuple transitivity}
\proofstatus{Open lemma in this edition}{The exact finite formulas in
Section~\ref{proof:new-dist:local_ring_actions} do not depend on this lemma;
only the advertised uniform rank-stability bound does.}
Let $R_a$ be a finite chain ring and $L$ a finite $R_a$-module generated by
at most $n$ elements. The stability argument would follow from:
\begin{quote}
$\mathrm{GT}(n,L)$: any two surjections $R_a^n\twoheadrightarrow L$ differ
by precomposition with an element of $\GL_n(R_a)$.
\end{quote}
Reducing modulo the maximal ideal proves transitivity on minimal generating
sets over the residue field, and elementary column operations lift changes of
basis through $I+\pi M_n(R_a)$. The remaining issue is to prove that the
kernels of arbitrary nonminimal surjections can always be matched by such a
lift for every module $L\le R_a^{dD}$. Until this module-theoretic statement
is cited or proved, equation~\eqref{nd:local_ring_actions:eq:stable} is a
conditional theorem, not an unconditional all-chain-ring result.

\section{Formed-space extension over finite chain rings}
The bounded-support argument for symplectic and orthogonal towers requires a
precise theorem saying that an isometry between supported split, unimodular
formed submodules extends to the ambient group in the stated rank range. In
residue characteristic two the quadratic refinement must be part of the data.
A complete proof must specify the ring class, form type, nondegeneracy and
split hypotheses, and the exact unused-rank requirement. Without these
hypotheses the universal fixed-point formula remains exact, but the
eventual-stability conclusion is only conditional.

\section{Fixed-tail wreath character polynomials}
Equation~\eqref{nd:multiplicative_wreath:eq:fixed-tail-limit} requires a
wreath-product character-polynomial theorem with an explicit threshold
$N_{\boldsymbol\alpha}(D)$ that remains valid after the class-power maps used
for every $r\le D$. Once such a theorem is supplied, the marked Poisson
moment argument is finite-dimensional and rigorous. This edition does not
claim a universal threshold for arbitrary multipartition tails.

\section{Generic braid and TQFT image families}
For a specified finite represented image, or for a specified compact closure
with normalized Haar measure, the Molien and balanced formulas in
Section~\ref{proof:new-dist:braid_tqft} are exact. A generic Jones or TQFT
family additionally requires a theorem identifying the image or compact
closure and its represented character/power data. The concrete Ising,
Heisenberg, Weil, and Clifford examples do not by themselves establish such a
family theorem.

\section{Exact computational certificates still required}
The following upgrades are evidentiary rather than changes to the universal
formulas:
\begin{itemize}
\item reconstruct $H_3$ over $\mathbb Q(\sqrt5)$ or use a certified character
table;
\item attach versioned class sizes, power maps, frame shapes, and hashes for
$M_{24}$, $Co_0$, lattice, and sporadic examples;
\item for every numerical rank, report the largest singular value declared
zero and the smallest declared positive;
\item archive the scripts, lock files, exact outputs, and data tables referenced
by the reproduction commands.
\end{itemize}

\end{document}
